\documentclass[10pt]{amsart}
\usepackage{amsmath,amssymb,amsthm,mathtools,mathrsfs}
\usepackage[a4paper,hmargin=0.76in,vmargin=0.60in]{geometry}
\usepackage{booktabs}
\usepackage{array}
\usepackage{longtable}
\usepackage{tabularx}
\usepackage[hypertexnames=false]{hyperref}
\hypersetup{hidelinks}
\usepackage[final,expansion=false]{microtype}
\usepackage{url}
\emergencystretch=2.5em

\newtheorem{theorem}{Theorem}[section]
\newtheorem{proposition}[theorem]{Proposition}
\newtheorem{lemma}[theorem]{Lemma}
\newtheorem{corollary}[theorem]{Corollary}
\theoremstyle{definition}
\newtheorem{definition}[theorem]{Definition}
\newtheorem{problem}[theorem]{Problem}
\theoremstyle{remark}
\newtheorem{remark}[theorem]{Remark}

\newcommand{\C}{\mathbb C}
\newcommand{\R}{\mathbb R}
\newcommand{\N}{\mathbb N}
\newcommand{\Z}{\mathbb Z}
\newcommand{\Q}{\mathbb Q}
\newcommand{\Rel}{\operatorname{Re}}
\newcommand{\Iml}{\operatorname{Im}}
\newcommand{\supp}{\operatorname{supp}}
\newcommand{\TV}{\operatorname{TV}}
\newcommand{\Kp}{\mathcal K^{\mathrm{prim}}_{2T}}
\newcommand{\Kw}{\mathcal K_{2T}}
\newcommand{\PT}{\mathcal P_T}
\newcommand{\IT}{I_T}
\newcommand{\lamP}{\lambda^{\mathrm{prim}}_T}
\newcommand{\lamW}{\lambda^{\mathrm W}_T}
\newcommand{\lamA}{\lambda_{\mathrm A}(T)}
\newcommand{\sT}{\sigma_T}
\newcommand{\gh}[1]{\widehat g_{#1}}
\newcommand{\nuh}{\widehat\nu}
\newcommand{\wh}{\widehat w}
\newcommand{\Fh}{\widehat F}
\newcommand{\cF}{c_F}
\newcommand{\Lev}{\mathfrak E}
\newcommand{\Wcl}{\mathcal W_T}
\newcommand{\NT}{\mathcal N_T}
\newcommand{\Aconst}{\mathsf A_T}
\newcommand{\Ainf}{\mathcal A_\infty}
\newcommand{\PiT}{\Pi_T}
\newcommand{\mT}{\mathfrak m_T}
\newcommand{\laminf}{\lambda^{\infty}_{\min}}
\newcommand{\laminfP}{\lambda^{\infty,\mathrm P}_{\min}}
\newcolumntype{Y}{>{\raggedright\arraybackslash}X}

\providecommand{\anCP}{C_P}
\providecommand{\anPa}{\Phi_a}
\providecommand{\anPb}{\Phi_b}
\providecommand{\anPl}{\Phi}
\providecommand{\ancl}{\mathrm{cl}}
\providecommand{\aneu}{\mathrm{eu}_d}
\providecommand{\anrd}{\mathrm{rd}}
\providecommand{\auAz}{\mathcal A_0}
\providecommand{\auLat}{\mathcal L}
\providecommand{\auP}{\mathcal P}
\providecommand{\auQlb}{\overline{\Q}_\ell}
\providecommand{\auS}{\mathcal S}
\providecommand{\auT}{\mathcal T}
\providecommand{\audT}{d_{\auT}}
\providecommand{\augain}{\operatorname{gain}}
\providecommand{\auhmain}{h_{\mathrm{main}}}
\providecommand{\aurem}{\operatorname{rem}}
\providecommand{\ausmain}{s_{\mathrm{main}}}
\providecommand{\paI}{I_p^{\circ}}
\providecommand{\parsgn}{\operatorname{sgn}}
\providecommand{\pavp}{\varpi}
\providecommand{\pbdeg}{\operatorname{deg}}
\providecommand{\pbeq}{\mathsf q}
\providecommand{\pcF}{\mathcal F}
\providecommand{\pcG}{\mathsf G}
\providecommand{\pcRal}{\mathcal R}
\providecommand{\pcW}{\mathbb W}
\providecommand{\pcmc}{\operatorname{mc}}
\providecommand{\pcvol}{\operatorname{vol}}
\providecommand{\pgpr}{\operatorname{pr}}
\providecommand{\pgv}{v}
\providecommand{\pmP}{\mathsf P}
\providecommand{\psHS}{\mathcal H}
\providecommand{\psM}{\mathsf M}
\providecommand{\psvac}{\mathbf v}
\newenvironment{quotedthm}[2]{\par\medskip\noindent\begingroup
  \textbf{#1}#2.\ \itshape\ignorespaces}{\endgroup\par\medskip}
\newenvironment{quoteddef}[2]{\par\medskip\noindent\begingroup
  \textbf{#1}#2.\ \ignorespaces}{\endgroup\par\medskip}
\begin{document}
\addtolength{\abovedisplayskip}{-3.4pt}
\addtolength{\belowdisplayskip}{-3.4pt}
\addtolength{\belowdisplayshortskip}{-2pt}
\title[The Covolume Theorem: one theta lift]{The Covolume Theorem: one theta lift, Hecke $L$-values and test forms of least valuation}
\author{David Alejandro Trejo Pizzo}
\date{}
\begin{abstract}
For a theta lift in Weil's range from a unitary group to a definite unitary group over a CM
field, formed with a character twisted by a tamely ramified character~$\vartheta$,
we prove that the Gram determinant of the lift, multiplied by the determinant of an explicit
operator compressed to the space of lifts, is a constant times a power of a product of values of
Hecke $L$-functions twisted by~$\vartheta$. The operator is integral away from the conductor
of~$\vartheta$, and its determinant is bounded at each prime by the congruences between lifts and
other forms. We then bound the heights of the ratios of twisted to untwisted Hecke values: their
valuation at every prime is determined, for suitable twists, by the divisibility of Taylor
coefficients of theta functions with complex multiplication, and we show that this divisibility
is attained by an integral test form.
\end{abstract}
\maketitle
\setcounter{tocdepth}{1}\tableofcontents

\section{Introduction}\label{C2:intro}

The main paper~\cite{P46} of this series deduces the Riemann hypothesis for Dirichlet
$L$-functions from three theorems on integral algebraic modular forms on a chain of $m$ theta
lifts between unitary groups over a CM field $E$: the Short Forms Theorem, the Spanning Theorem
and the Covolume Theorem. In the proof of the third, the Gram determinant $Z(\vartheta)$ of the
chain formed with a twist $\vartheta$ is compared with $Z(1)$ at each prime~$\ell$, for a twist
depending on~$\ell$. By the factorisation proved in~\cite{P53} ({\cite[Proposition~14.37]{P46}}),
$Z(1)/Z(\vartheta)$ is, up to explicit ratios of Hecke $L$-values, the product of the ratios
$Z_j(1)/Z_j(\vartheta^{\,j})$ for $m$ single theta lifts $W_j\to V_j$. This paper is devoted to
one such lift. It proves the two propositions of~\cite{P46} that concern it:
{\cite[Proposition~14.41]{P46}}, which relates $Z_j(\vartheta)$ to values of Hecke $L$-functions,
and {\cite[Proposition~14.42]{P46}}, which bounds the heights of these values along a regular family
of twists.

We fix $1\le j\le m$ and use the notation of {\cite[\S\S14.1--14.3]{P46}}. The
space $W=W_j$ is hermitian of rank $r$ over $E$, with $k+j-1$ positive and $m-j$ negative signs
at every infinite place, and $V=V_j$ is positive definite of rank $r+t$, $t=5r$. A twist
$\vartheta$ is a unitary automorphic character of $\mathrm U(1)$, trivial at the infinite places
and tamely ramified at a finite set $Q_\vartheta$ of primes. The theta lift $\Theta_\vartheta$
for the pair $(\mathrm U(W),\mathrm U(V))$, formed with the character $\mu_j^0\vartheta$, maps
the space $e_0\mathcal A(W_j)$ of good forms of discrete series type on $\mathrm U(W)$, of
dimension $d_j$, to a space $\mathcal A(V;\vartheta)$ of algebraic modular forms on
$\mathrm U(V)$, and $Z_j(\vartheta)=\det(\Theta_\vartheta^*\Theta_\vartheta)$. The space
$\mathcal A(V;\vartheta)$ has an integral structure at every prime; the image
$\mathcal Z_\vartheta$ of $\Theta_\vartheta$ is not a direct summand for it, and the failure is
measured as follows. A \emph{false lift} is a form of $\mathcal A(V;\vartheta)$ of the type of
the lifts at the chosen places that satisfies the Hecke equation of a lift at one split place
and is orthogonal to $\mathcal Z_\vartheta$ (Definition~\ref{po:false}). The twist $\vartheta$
\emph{has denominators at most $\beta$} at a place datum $\iota$ if some endomorphism of
$\mathcal Z_\vartheta$ of determinant at least $\exp(-d_j\,r\,\beta)$ in absolute value at
$\iota$ carries to integral forms the orthogonal projections to $\mathcal Z_\vartheta$ of all
integral forms orthogonal to the false lifts (Definition~\ref{po:denom}).

Our first result is the following ({\cite[Proposition~14.41]{P46}}). \emph{There are a constant
$C_j\ne0$ and, for every twist $\vartheta$, an endomorphism $U_\vartheta$ of
$\mathcal A(V;\vartheta)$ such that
$\Delta_j(\vartheta)=\det(\operatorname{pr}_{\mathcal Z_\vartheta}U_\vartheta|_{\mathcal Z_\vartheta})$
is not zero and
\[
Z_j(\vartheta)\cdot\Delta_j(\vartheta)=C_j\cdot\Bigl(\frac{H_j(\vartheta)}{H_j(1)}\Bigr)^{d_j},
\qquad
H_j(\vartheta)=\prod_{i\in\{\ast,+,-\}}\Bigl(\bigl(\operatorname{Nm}\mathsf S_\vartheta\bigr)^{z_i^2}
\prod_{z=1}^{z_i}\bigl|L^{\mathsf S_\vartheta}\bigl(z,(\nu_i\vartheta^{s_i})^{-1}\bigr)\bigr|^2\Bigr),
\]
where the $L^{\mathsf S_\vartheta}(s,\nu^{-1})$ are Hecke $L$-functions of $E$ without the Euler
factors at the level, $(\nu_i,z_i)$ are the three pairs of \S\ref{po:ss-hecke}, and the signs
are $s_\ast=s_+=1$, $s_-=-1$. Moreover
$|\sigma\Delta_j(1)|\le\exp(3d_j\,f\,m\,b\,(\log b)^2)$ for every automorphism $\sigma$ of
$\mathbb C$, and $|\Delta_j(\vartheta)|_\iota\le\exp(d_j\,r\,\beta)$ whenever $\ell(\iota)$ does
not divide $\operatorname{Nm}Q_\vartheta$ and $\vartheta$ has denominators at most $\beta$ at
$\iota$.} Thus $Z_j(\vartheta)$ is not itself a product of Hecke $L$-values, but it becomes one
after multiplication by a number whose size at the trivial twist and whose valuations are
controlled. The characters $\nu_\ast=\mu_j^0\upsilon^{n'}$, $\nu_+=\upsilon^{-m}\chi_0^{\,j}$ and
$\nu_-=\upsilon^{-m/2}\chi_0^{-j}$ are ramified at the chosen places, and the arguments
$1\le z\le z_i$, with $z_\ast=(7m+k)/2$, $z_+=m$ and $z_-=m/2$, lie in the region $\Re s\ge1$;
the positivity of $H_j(\vartheta)$ is Lemma~\ref{po:Hpos}. For $z\ge2$ the Euler product converges
absolutely. For $z=1$ the statement is the case $s=1$ of Hecke's theorem that these $L$-functions
do not vanish on the line $\Re s=1$; we deduce it in the proof of the lemma from the analytic
continuation of Hecke $L$-functions~\cite{Hecke1920}, by the classical argument with the product
$\zeta_E^3L(s,\chi)^4L(s,\chi^2)$ if $\chi^2\ne1$ and by Landau's theorem on Dirichlet series
with non-negative coefficients if $\chi^2=1$.

The operator is $U_\vartheta=c\cdot\operatorname{glue}_\vartheta\circ\operatorname{glue}_\vartheta^*$
with a positive integer $c$, where the \emph{gluing map} $\operatorname{glue}_\vartheta$
(\S\ref{po:ss-glue}) sends tuples of forms on the unitary group of a definite subspace
$V^\circ\subset V$ to forms on $\mathrm U(V)$, by multiplication with integral forms on the group of the
complement, the \emph{first factors}, after a change of lattice at the chosen places. Three properties of this map give the theorem
(Proposition~\ref{po:glue}). First, the forms $\operatorname{glue}_\vartheta(\operatorname{glue}_\vartheta^*\Theta_\vartheta y)$
are orthogonal to the false lifts; this is a vanishing statement for periods of false lifts over
$\mathrm U(V^\circ)$ (Proposition~\ref{po:periods}, proved in \S\ref{ppe:sec}), and it is
what makes the compression to $\mathcal Z_\vartheta$ compatible with the definition of
denominators. Secondly, $\operatorname{glue}_\vartheta$ preserves integrality at the primes not dividing
$\operatorname{Nm}Q_\vartheta$, the map $\operatorname{glue}_1$ increases the square of no
conjugate norm by more than a factor $\exp(2fmb(\log b)^2)$, and an integral tight frame of the
space of values (Lemma~\ref{po:frame}), that is, a family of integral vectors $\eta$ with
$\sum\langle\eta,\xi\rangle\eta=H^2\xi$ for an integer $H$, writes $U_\vartheta$ as a sum of
operators $\Xi\mapsto c\,\langle g,\Xi\rangle\,g$ of rank one with integral forms $g$, so that $U_\vartheta$ preserves integrality at
these primes and all conjugates of $U_1$ are bounded; this is the origin of
the two bounds for $\Delta_j$. Thirdly, the composite
$\operatorname{glue}_\vartheta^*\circ\Theta_\vartheta$ is, up to a scalar $c'_\vartheta$ with
$|c'_\vartheta|^2H_j(\vartheta)$ independent of $\vartheta$, a map that does not depend on
$\vartheta$ at all: this is the statement of the \emph{three exchanges}
(Proposition~\ref{po:exchange}), in which the exchange identity of~\cite{P53} is applied three
times; the constants of the three applications produce the three factors of $H_j$. That the resulting map is injective on $e_0\mathcal A(W_j)$ is the statement that
\emph{products span} (Proposition~\ref{po:span}); its proof, in
Section~\ref{C2:g4}, uses a reproducing formula which expresses the
coefficients of holomorphic discrete series through theta lifts to a definite group
(Theorem~\ref{psr:R}). The delicate point of the whole construction is the choice of the first
factors: they must be integral at every prime not dividing $\operatorname{Nm}Q_\vartheta$,
short at the trivial twist, and compatible with the three exchanges
(Lemma~\ref{po:ff}, Theorem~\ref{pfq:ff} and Proposition~\ref{pfq:exchange}).

The second result concerns the numbers $H_j$. Recall from {\cite[\S14.2]{P46}} that a family
$y=(y_\iota)$ of complex numbers indexed by the place data $\iota=(\ell,\,\mathbb C\cong
\overline{\mathbb Q}_\ell)$ has \emph{height at most $c$} if, for a multiplier $H\in\mathbb C^\times$
all of whose conjugates are at most $e^\alpha$ in absolute value, $|Hy_\iota|_\iota\le\ell^{h_\ell}$
for all $\iota$ and $\alpha+\sum_{\ell\in T}h_\ell\log\ell\le c$ for all finite sets $T$ of primes.
A family $\theta=(\theta_\iota)$ of $m$-fold twists is \emph{regular}
({\cite[Definition~14.39]{P46}}) if the conductor of $\theta_\iota$ is prime to $\ell(\iota)$,
if certain ratios of twisted to untwisted Hecke values have height at most $x_1$, and if, with
exponents $\beta_\ell$ whose sum is at most the budget $\mathrm{bud}=f\,a\,(\log b)^3/(mk)$, the
twists $\theta_\iota^{\,j}$ have denominators at most $\beta_\ell$ for the $j$-th lift, and
the inner products of the projections of integral test forms to the lines of the lifts of
characters considered below are at most $e^{\beta_\ell}$ at $\iota$; such families exist by~\cite{P55}.
\emph{For a regular family $\theta$, the family $\iota\mapsto H_j(1)/H_j(\theta_\iota^{\,j})$ has
height at most}
\[
10\,m\,x_1+k\,(6\,\mathrm{bud}-h_{\mathrm{dir}})
\]
({\cite[Proposition~14.42]{P46}}). Here $h_{\mathrm{dir}}=\sum_{\sigma}\bigl(s_\infty(\kappa(\sigma))-
\sum_p2\log p\cdot s_p(\kappa(\sigma),\sigma)\bigr)$ is the direction amount of
{\cite[Definition~14.8]{P46}}: the sum is over the infinite places $\sigma$ of $E$, $\kappa$ is
the weight, $s_\infty(n)=\log n!+n\log(M/\Omega)$ is the archimedean size of a Taylor coefficient
of order $n$ of a theta function with complex multiplication by $E$, normalised by the period
$\Omega$, and the local prime sum $s_p(n,\sigma)$ is the $p$-adic divisibility of such a
coefficient guaranteed by {\cite[Proposition~14.10]{P46}}, which vanishes unless $p$ is not split in
the imaginary quadratic subfield~$K$. The bound says that, for the twists of a regular family,
the Hecke values are no more divisible than these Taylor coefficients are known to be. It is this
statement that makes the main terms of the upper and lower bounds of~\cite{P46} cancel.

The proof considers the $5m+k/2$ critical pairs $(\nu_i,z)$ that occur in $H_j$
(Lemma~\ref{po:Hratio}). For $5m$ of them the regularity of the family is used directly. The
remaining $k/2$ are \emph{main pairs} $p=(\nu,z)$ (Definition~\ref{po:main}), and for them the
Hecke value is the squared norm of a theta lift from $\mathrm U(1)$: if $\varphi_{p,\vartheta}$
denotes the normalised lift of $\nu\vartheta$ to the definite unitary group of rank $2z$, then
\[
\|\varphi_{p,\vartheta}\|^2=A_p\,\bigl(\operatorname{Nm}Q_\vartheta\bigr)^{2z-1}\,
L^{\mathsf S_\vartheta}\bigl(z,(\nu\vartheta)^{-1}\bigr)
\]
with $A_p>0$ independent of $\vartheta$ (Proposition~\ref{pex:rallis}); this follows from the
inner product formula of~\cite{P53}, hence from Weil's Siegel formula~\cite{WeilSiegel}. The
values of $\varphi_{p,\vartheta}$ are Taylor coefficients, of orders close to $\kappa(\sigma)$,
of theta functions with complex multiplication, and by the count their valuation at a
prime~$\ell$ is at least about $F(\ell)=\sum_\sigma s_\ell(\kappa(\sigma),\sigma)$. We prove two
facts. The first is archimedean: the constant family $\|\varphi_{p,1}\|^2$ has height at most
$\mathrm{bud}-\sum_\sigma s_\infty(\kappa(\sigma))$ (Theorem~\ref{psu:arch}), so that all conjugates
of the untwisted norm have the expected size; the passage to all conjugates uses the law of Galois
conjugation of the lifts of characters proved in~\cite{P53}. The second, which is the essential
point, is that the lower bound of the count is attained: there is $c_0\ne0$, all of
whose conjugates are at most $\exp(fmb/k+3x^{35})$, such that for every prime $\ell$, every
embedding and every twist $\vartheta$ of conductor prime to $\ell$ there is a test form $\Xi$,
integral at $\ell$ and orthogonal to the relevant false lifts, with
\[
\bigl|c_0\,\rho_{q_2}(\vartheta)\bigr|\cdot\ell^{-F(\ell)}\le\bigl|\langle\Xi,\varphi_{p,\vartheta}
\rangle\bigr|
\]
in the $\ell$-adic absolute value, where $\rho_{q_2}$ is the ratio of Hecke values of an
auxiliary pair (Theorem~\ref{psu:small}). Since the inner product of the projections of two test
forms to the line of $\varphi_{p,\vartheta}$ is $\langle\Xi,\varphi\rangle\langle\varphi,\Xi'\rangle/\|\varphi\|^2$, and regularity bounds it by $e^{\beta_\ell}$, the norm
$\|\varphi_{p,\theta_\iota^{\,j}}\|^2$ cannot have valuation much larger than $2F(\ell)$; summing
over $\ell$ gives the height $6\,\mathrm{bud}-h_{\mathrm{dir}}$ of a main pair
(Proposition~\ref{po:mainheight}).

The existence of the test form rests on the \emph{sharpness of the count}
(Theorems~\ref{psm:twelve} and~\ref{psm:attained}): at every prime $p$, some Taylor coefficient of
the prescribed orders of a theta function with complex multiplication, for a suitable indicator
datum, has valuation at most $F(p)$, up to a factor whose conjugates are bounded uniformly in
$p$. This is a statement about every prime, including $2$, the primes that ramify in $E$ and the
primes of the level, and the cases are treated separately according to the splitting of $p$ in
$K$ and to its size (\S\ref{psm:sec}). The passage from one Taylor coefficient to an integral form orthogonal
to the false lifts whose pairing with the lift is computed exactly occupies
\S\ref{psu:sec}.

In the course of the proof we also establish the statement (IM)
(Propositions~\ref{pex:IM} and~\ref{pex:IM1}): for a ramified critical pair $(\nu,z)$ with $z\ge2$
the suitably normalised norm $\|\varphi_{p,\vartheta}\|^2$, and for $z=1$ a number
$\Lambda(\vartheta)$ with $|\Lambda(\vartheta)|^2$ a constant times
$\operatorname{Nm}Q_\vartheta\,|L(1,(\nu\vartheta)^{-1})|^2$, is a finite linear combination of
values of $\vartheta$ whose coefficients are integral at the primes not dividing
$\operatorname{Nm}Q_\vartheta$. It is the input from theta lifts of the twisting theorem
of~\cite{P55}.

The notation and the data are those of~\cite{P46}. From~\cite{P49} we use the properties of the
Weil representation, the theta functions with complex multiplication with the count, and the
archimedean results; from~\cite{P51} estimates for the coefficients of discrete series. With
\cite{P53} the relation goes in both directions, and with~\cite{P55} in one; we describe both. From~\cite{P53}
we use the exchange identity, the inner product formula and the law of conjugation
({\cite[Theorem~9.1]{P53}}). In return \cite{P53} uses the archimedean facts on one lift of
Section~\ref{C2:g3} (Lemmas~\ref{plf:A12} and~\ref{plf:harm}, Proposition~\ref{plf:L4}), the standing facts on a lift (Lemma~\ref{po:facts}) and two lemmas of Sections~\ref{C2:g3} and~\ref{C2:g4}
(Lemmas~\ref{plf:KS} and~\ref{psr:Egrowth}), which are proved here without the exchange identity
and the inner product formula, and it uses the formula for the norm of the lift of a character
(Proposition~\ref{pex:rallis}, with Lemmas~\ref{pex:arch}, \ref{po:naut} and~\ref{psu:flat}) in the proof of the law of conjugation; that law enters this paper
only in Theorem~\ref{psu:arch}, after Proposition~\ref{pex:rallis} and these lemmas. No proof in this paper cites a result of~\cite{P55}. The existence of regular
families is proved there with the help of the statement (IM) above (Propositions~\ref{pex:IM} and~\ref{pex:IM1}), and of the equivariance
and the Hecke equation of a lift (Lemma~\ref{po:facts}); these do not involve regular families,
which enter this paper only as the hypothesis of {\cite[Proposition~14.42]{P46}}. So there is no circle
in either case. The statements quoted are reproduced in Appendix~\ref{C2:quoted}. Besides the theorems of Hecke and Weil already
mentioned, the published results that we use most are the description of the joint harmonics and
of the local theta correspondence for compact dual pairs by Kashiwara and Vergne and by
Howe~\cite{KashiwaraVergne,HoweJAMS}, the results of Rallis on theta lifts and of Harris, Kudla
and Sweet on the Weil representation of unitary
groups~\cite{RallisInjectivity,HKS,KudlaSplitting}, and Green's parametrisation of the characters
of finite general linear groups~\cite{Green1955}.

The paper is organised as follows. Section~\ref{po:sec} sets up the lift, the false lifts, the
gluing map and the comparison operator, and deduces {\cite[Propositions~14.41 and~14.42]{P46}} from statements proved in the later sections. Section~\ref{C2:g2} computes the
norm of the lift of a character and proves (IM) (\S\ref{pex:sec}), and proves the three exchanges
(\S\S\ref{pte:sec} and~\ref{ptf:sec}). Section~\ref{C2:g3} contains the facts on classes and on
the archimedean type of the descent (\S\ref{plf:sec}) and the vanishing of the periods of false
lifts (\S\ref{ppe:sec}). Section~\ref{C2:g4} proves that products span, and
Section~\ref{pfq:sec} the properties of the first factors and of the characters at the chosen
places. Section~\ref{C2:g6} proves the sharpness of the count, constructs the test form of least
valuation, and bounds the untwisted norm.


\providecommand{\auO}{\mathcal O}
\providecommand{\auI}{\mathcal I}
\providecommand{\auU}{\mathrm U}
\providecommand{\auGL}{\mathrm{GL}}
\providecommand{\auHom}{\operatorname{Hom}}
\providecommand{\auEnd}{\operatorname{End}}
\providecommand{\auAut}{\operatorname{Aut}}
\providecommand{\auNm}{\operatorname{Nm}}
\providecommand{\auA}{\mathbb A}
\providecommand{\auF}{\mathbb F}
\providecommand{\auPl}{\Phi}
\providecommand{\auG}{\overline G}
\providecommand{\audT}{d_{\mathcal T}}
\providecommand{\aubud}{\mathrm{bud}}
\providecommand{\auhdir}{h_{\mathrm{dir}}}
\providecommand{\ausinf}{s_\infty}
\providecommand{\ausfin}{s_{\mathrm{fin}}}
\providecommand{\aueu}{\varepsilon_{\mathrm u}}
\providecommand{\aurd}{\mathrm{rd}_E}
\providecommand{\aubar}[1]{\overline{#1}}
\providecommand{\auf}[1]{\texttt{#1}}
\providecommand{\psrho}{\varrho}
\providecommand{\pao}{\mathfrak o}
\providecommand{\plone}{\mathbf 1}
\providecommand{\plgb}[2]{\|#2\|_{#1}}
\providecommand{\poA}{\mathcal A}
\providecommand{\poZ}{\mathcal Z}
\providecommand{\poT}{\mathbb T}
\providecommand{\popr}{\operatorname{pr}}
\providecommand{\poglue}{\operatorname{glue}}
\providecommand{\poran}{\operatorname{im}}
\providecommand{\podel}{\delta}
\providecommand{\pop}{\mathsf p}
\providecommand{\pog}{\mathsf g}
\providecommand{\poAA}{A_\ast}

\section{One lift: the comparison operator, false lifts and Hecke values}\label{po:sec}

This section treats {\cite[Propositions~14.41 and~14.42]{P46}}. We fix $1\le j\le m$
throughout and write
\[
W=W_j,\qquad V=V_j,\qquad n=r+t=6r,\qquad r_+=k+j-1,\quad r_-=m-j,\qquad
\poAA=f\,m\,b\,(\log b)^2 .
\]

\emph{What is proved and where.} The two propositions are deduced here, with complete
arguments, from a list of statements about theta kernels and Hecke values which are stated in this
section and proved in the sections that follow. The deduction contains the linear algebra
(the determinant identity, the two eigenvalue bounds, the tight frame, the projection that does not
increase any conjugate norm), the Hecke equation of a lift, the change of lattice by which a lift
is cut into products (``gluing''), and the bookkeeping of heights. The statements and the places
where they are proved:
\begin{center}
\begin{tabular}{lll}
\hline
Statement&Content&Proved in\\
\hline
(L4)&the descent of a form on $V$ is of the type of $\poA(W)$&Prop.~\ref{plf:L4}\\
Prop.~\ref{po:exchange}&three exchanges&\S\ref{pte:sec}, \S\ref{ptf:sec}\\
Prop.~\ref{po:span}&products span&Prop.~\ref{psp:main}, Thm.~\ref{psv:main}\\
Prop.~\ref{po:periods}&periods of false lifts vanish&Thm.~\ref{ppe:main}\\
Lemma~\ref{po:ff}&first factors: integrality, size&Thm.~\ref{pfq:ff}\\
Lemma~\ref{po:classes}&numbers of classes at level $\pi_M^2$&Lemma~\ref{plf:classes}\\
Prop.~\ref{po:rallis}&the norm of a lift from $\auU(1)$&Prop.~\ref{pex:rallis}\\
Prop.~\ref{po:arch}&all conjugates of the untwisted norm&Thm.~\ref{psu:arch}, {\cite[Thm.~9.1]{P53}}\\
Prop.~\ref{po:small}&a test form of least valuation&Thm.~\ref{psu:small}, \ref{psm:attained}\\
\hline
\end{tabular}
\end{center}
Two definitions of {\cite[\S14.10]{P46}} have to be made precise for the statements to have a
meaning (false lifts; the clause of (R3) on main pairs); this is done in
Definitions~\ref{po:false}, \ref{po:denom} and~\ref{po:R3}. With these, {\cite[Propositions~14.41 and~14.42]{P46}} hold as stated, and the
proofs of {\cite[Proposition~14.43]{P46}} and of the Covolume Theorem ({\cite[Theorem~14.15]{P46}}) are unchanged.

\subsection{The lift $W\to V$: definitions and standing facts}\label{po:ss-set}

\emph{Places.} $\mathsf M$ is the set of chosen places and, for a twist $\vartheta$,
$\mathsf Q=Q_\vartheta$; $\mathsf S_\vartheta=\mathsf M\cup\mathsf Q$. A prime $w$ of $E$ is
\emph{off the level} if $w\ne\aubar w$ and $w,\aubar w\notin\mathsf S_\vartheta$. For
$w\in\mathsf S_\vartheta$, $E_w$ is identified with the completion and $\auF_w$ is the residue
field. We use (N1)--(N5) of {\cite[\S6.1]{P49}}, in particular the coordinates $T$ on
$\auHom(W_w,V_w)=M_{n\times r}(E_w)$ at the split places and the formulas {\cite[(62) and~(63)]{P49}}; the proofs of these formulas are local and apply to the pair $(W,V)$ at every
place $w\ne\aubar w$ (for the odd ranks $r$, $t$ the factor in {\cite[(62)]{P49}} is changed by an
unramified character of $E_w^\times\times E_w^\times$; we use the formula in the form
\eqref{po:eq-w5} below, where the factor on $\auU(V)$ is an unramified character $\eta_w$ that does not
involve $\mu$: on elements of $\auGL(\auO_w)$ it is trivial, and for the non-integral elements $d_1$
(Lemma~\ref{po:facts}) and $\pog_s$ (Lemma~\ref{po:slice}) only $|\eta_w(1,h)|=|\det h|_w^{-r/2}$ and its
independence of $\vartheta$ are used).

\emph{The character, the Weil representation and the datum.} $\mu=\mu_j^0\vartheta$; its tame
part at $w\in\mathsf S_\vartheta$ is a character $\psi_w$ of $\auF_w^\times$, non-trivial for
$w\in\mathsf Q$ by the definition of a twist, and for $w\in\mathsf M$ equal to the tame part of
$\mu_j^0$; we write $\psi_{\mathsf M}(\det s)=\prod_{w\in\mathsf M}\psi_w(\det s_w)$ for
$s\in\auG$. In this section $\omega_\mu$ denotes the Weil representation of
$\auU(W)(\auA)\times\auU(V)(\auA)$ on $\mathcal H(\auHom(W,V))$ in the following normalisation: the group $\auU(V)$ acts through the splitting given by
powers of $\upsilon$ alone, and the action of $\auU(W)$ is the one of that splitting multiplied by
$\mu(\det g)^{-1}$. At a place $w\ne\aubar w$ this means, in the coordinates of (N5),
\begin{equation}\label{po:eq-w5}
\omega_\mu(g,h)\Psi(T)=\chi_w(\det g)^{-1}\,\eta_w(g,h)\,\Psi(h^{-1}Tg),
\end{equation}
where $\chi$ is the base change of $\mu$ and $\eta_w$ is an unramified character of
$\auGL_r(E_w)\times\auGL_n(E_w)$ that does not depend on $\mu$, with $|\eta_w(1,h)|=|\det h|_w^{-r/2}$.
In the notation of {\cite[\S6.1.4]{P55}} this is the action $a(\det h)\,b(\det g)^{-1}\omega(x\mapsto
hxg^{-1})$ with $a=\upsilon^{r_-}$ and $b=\mu\,\upsilon^{\lfloor n/2\rfloor}$. (It differs from the
convention of {\cite[\S5.1.1]{P49}} for the upper chain, where the character twists both groups, by the
unitary character $\mu(\det h)$ of $\auU(V)(\auA)$; with the present convention the scalars of
the Hecke operators and of the change of lattice on $V$ do not depend on $\vartheta$. The same
representation is used in {\cite[\S6.1.4]{P55}}; the lift of {\cite[\S2.4]{P53}} differs from it by multiplication with a unitary character of the determinant on $\auU(V)$, which changes no norm and no determinant $Z_j$: {\cite[Remark~2.17]{P53}}.) Let
$\iota_0\in M_{n\times r}$ be the matrix of the inclusion of the first $r$ coordinates, and for
$T\in M_{n\times r}$ let $T'$ be its upper $r\times r$ block. The \emph{block datum} is
$\Phi_\vartheta=\bigotimes_{w\in\mathsf M}\plone[T\equiv\iota_0\bmod w]\otimes
\bigotimes_{w\in\mathsf Q}\plone[T\text{ integral}]\,\psi_w(\det T'\bmod w)\otimes(\text{vacuum
elsewhere})$, with $\psi_w(0)=0$.

\emph{The weight on $V$.} At each $\tau\in\auPl$ there are two Young diagrams $\mu^+_\tau$ and
$\mu^-_\tau$, with columns of heights at most $r_+$, respectively at most $r_-$; they are
determined by the parameter $\varsigma_\tau$ and by $\kappa_j(\tau)$ (explicitly in Lemma~\ref{plf:match}), and they do not depend on
$\vartheta$. We use only: $\mu^+_\tau$ has at most $4b$ columns, and $\mu^-_\tau$ has at most
$\Delta$ columns of each height. In the notation of {\cite[\S5.2.6]{P49}} put
\[
\poT^+(R^n)=\bigotimes_{\mathfrak c\text{ column of }\mu^+}\Lambda^{h_{\mathfrak c}}R^n,\qquad
\poT^-(R^n)=\bigotimes_{\mathfrak c\text{ column of }\mu^-}\Lambda^{h_{\mathfrak c}}(R^n)^\vee
\]
(tensor products over all $\tau$ as well), with the bases of pure tensors of basis vectors and the
hermitian forms for which these are orthonormal. The group $\prod_\tau\auGL_n(\C)$ acts on
$\poT^+\otimes\poT^-$, and $W_{\lambda_V}\subseteq\poT^+(\C^n)\otimes\poT^-(\C^n)$ is the span of
the vectors
\begin{equation}\label{po:eq-ev}
\mathrm{ev}(y,y')=\bigotimes_{\mathfrak c}(y_1\wedge\dots\wedge y_{h_{\mathfrak c}})\otimes
\bigotimes_{\mathfrak c}(y'_1\wedge\dots\wedge y'_{h_{\mathfrak c}}),\qquad y_i\in\C^n,\ y'_i\in
(\C^n)^\vee,\ y'_{i'}(y_i)=0\ \text{for all }i,i'
\end{equation}
(at each $\tau$). This is the irreducible representation with highest weight
$(\mu^+_\tau;0,\dots,0;-\mu^-_\tau)$; in the conventions of (N1) it has columns of heights
$h_{\mathfrak c}$ and $n-h_{\mathfrak c}$ and a power of the determinant, and
$W_{\lambda_V}(\Z)=W_{\lambda_V}\cap(\poT^+\otimes\poT^-)(\Z^n)$, with the restricted hermitian
form, by {\cite[(51)]{P49}} and because $\Lambda^h(\Z^n)^\vee\otimes\det=\Lambda^{n-h}\Z^n$
isometrically. As in {\cite[\S5.2.6]{P49}}, for $\xi\in W_{\lambda_V}\otimes\C_\ell$ the largest
absolute value of a coordinate of $\xi$ in the basis of $W_{\lambda_V}(\Z)$ is the largest
absolute value of a tensor coordinate (the two unit balls coincide, $W_{\lambda_V}(\Z)$ being
saturated in the tensors).

\emph{Forms on $V$.} $K_\vartheta\subseteq K^\circ$ is the group of the $k$ with $k_w\iota_0\equiv
\iota_0\bmod w$ for $w\in\mathsf M$ and with $k_w\bmod w$ block lower triangular, for the blocks
$r+t$, for $w\in\mathsf Q$; by \eqref{po:eq-w5}, $\omega_\mu(1,k)\Phi_\vartheta=\chi_\vartheta(k)^{-1}
\Phi_\vartheta$ for a character $\chi_\vartheta$ of $K_\vartheta$ of finite order, trivial on the
elements that are $\equiv1$ at $\mathsf S_\vartheta$. Put
$\poA(V;\vartheta)=\poA_{\lambda_V}(K_\vartheta,\chi_\vartheta^{-1})$ (so $\Xi(\gamma hk)=\chi_\vartheta(k)^{-1}\gamma\,\Xi(h)$), with the inner product and the
norms $\|\cdot\|_\sigma$ of {\cite[\S14.3]{P46}}. The level $K_\vartheta$ is not free. Let
\[
n_{\mathrm{Aut}}=\#\auGL_{2fn}(\auF_3).
\]

\begin{lemma}\label{po:naut}
For every $g\in\auU(V)(\auA_{E^+,f})$ the order of $\Gamma\cap gK^\circ g^{-1}$ divides
$n_{\mathrm{Aut}}$, and $n_{\mathrm{Aut}}\le\exp(x^{40})$. If $\Xi,\Xi'\in\poA(V;\vartheta)$ are
integral at $\aubar\iota$ and at $\iota$ respectively, then $n_{\mathrm{Aut}}\langle\Xi,\Xi'\rangle
\in\auO_\iota$.
\end{lemma}
\begin{proof}
The group is finite and is a subgroup of $\auGL_n(E)\subseteq\auGL_{2fn}(\Q)$. A finite subgroup of
$\auGL_N(\Q)$ stabilises a lattice, so it is conjugate to a subgroup of $\auGL_N(\Z)$, and the
kernel of $\auGL_N(\Z)\to\auGL_N(\auF_3)$ has no torsion (Minkowski): if $g=1+3^aX$ with $a\ge1$,
$X\not\equiv0\bmod3$ and $g^p=1$ for a prime $p$, then $pX\equiv0\bmod3^{a}$ for $p\ne3$, and
$X+3^aX^2+3^{2a-1}X^3=0$ for $p=3$; both are absurd. So the order divides $\#\auGL_N(\auF_3)\le
3^{N^2}$, $N=2fn\le12fr$, and $f\le x^{12}$, $r\le3x^5$. For the last assertion,
$\langle\Xi,\Xi'\rangle=\sum_g\#(\Gamma\cap gK_\vartheta g^{-1})^{-1}\langle\Xi(g),\Xi'(g)\rangle$,
and each $\langle\Xi(g),\Xi'(g)\rangle\in\auO_\iota$ as in {\cite[\S14.3]{P46}}.
\end{proof}

\emph{The lift and the descent.} For $x\in\auHom(W,V)\otimes_{E^+,\tau}\R=M_{n\times r}(\C)$ write
$x=(x^+\mid x^-)$ for the columns belonging to the positive and to the negative basis vectors of
$W$. For $\xi\in W_{\lambda_V}$ let $\pop_\xi$ be the vector of the Fock space of
$\auHom(W,V)\otimes\R$ given by the polynomial
\begin{equation}\label{po:eq-poly}
\pop_\xi(x)=\bigl\langle\xi,\ \mathrm{ev}\bigl(\aubar{x^+},(x^-)^\flat\bigr)\bigr\rangle ,
\end{equation}
where $\flat\colon\C^n\to(\C^n)^\vee$ is the conjugate-linear map of the standard form, the
columns of $x^+$ are taken in the fixed numbering of the positive basis vectors and those of $x^-$
in the reverse numbering of the negative ones, and $\langle\ ,\,\rangle$ is the form of
$\poT^+\otimes\poT^-$ (in the Fock model the positive variables are holomorphic and the negative
ones antiholomorphic, so that $\pop_\xi$ is a polynomial vector; it is conjugate-linear in $\xi$
and $\auU(V)(\R)$-equivariant; these are the joint harmonics of~\cite[Ch.~III, \S5]{KashiwaraVergne}). For
$\Phi\in\mathcal H_f(\auHom(W,V))$, $g\in\auU(W)(\auA)$ and $h\in\auU(V)(\auA_f)$ put
\[
k_\Phi(g,h;\xi)=\theta\bigl(\omega_\mu(g,h)(\pop_\xi\otimes\Phi)\bigr),
\]
with $\omega_\mu$ as in \eqref{po:eq-w5}. For a form
$y\in\poA(W_j)$ ({\cite[\S14.3]{P46}}; these forms are cuspidal, hence rapidly decreasing, and the kernel has moderate growth by
{\cite[Lemma~3.29]{P53}}, so the integral converges) the function $\Theta(\Phi)y\colon\auU(V)(\auA_f)\to W_{\lambda_V}$ is
defined by
\begin{equation}\label{po:eq-lift}
\bigl\langle\xi,(\Theta(\Phi)y)(h)\bigr\rangle=\int_{[\auU(W)]}y(g)\,k_\Phi(g,h;\xi)\,dg ,
\end{equation}
and $\Theta_\vartheta=\Theta(\Phi_\vartheta)$; up to a positive constant independent of
$\vartheta$ its restriction to $e_0\poA(W_j)$ is the map $u$ of {\cite[Proposition~14.37]{P46}}. As in
(N4): $(\Theta(\Phi)y)(\gamma h)=\gamma(\Theta(\Phi)y)(h)$ for $\gamma\in\Gamma$ (invariance of
$\theta$ under rational points and equivariance of $\xi\mapsto\pop_\xi$), and
\begin{equation}\label{po:eq-equiv}
R_{h'}\,\Theta(\Phi)y=\Theta\bigl(\omega_\mu(1,h')\Phi\bigr)y,\qquad
\Theta\bigl(\omega_\mu(g',1)\Phi\bigr)y=\Theta(\Phi)(R_{g'^{-1}}y)
\end{equation}
for finite adelic $g',h'$, with $R$ the right translation. Hence $\Theta_\vartheta y\in\poA(V;
\vartheta)$. The \emph{descent} of $\Xi\in\poA(V;\vartheta)$ is the function on $\auU(W)(\auA)$
\begin{equation}\label{po:eq-descent}
(D_\vartheta\Xi)(g)=\aubar{\sum_{h\in\Gamma\backslash\auU(V)(\auA_f)/K_\vartheta}\frac{k_{\Phi_\vartheta}
(g,h;\Xi(h))}{\#(\Gamma\cap hK_\vartheta h^{-1})}},\qquad\text{so that}\qquad
\langle\Xi,\Theta_\vartheta y\rangle=\int_{[\auU(W)]}y(g)\,\aubar{(D_\vartheta\Xi)(g)}\,dg
\end{equation}
by \eqref{po:eq-lift} and the definition of the inner product.

\emph{The finite group.} For $s\in\auG=\prod_{w\in\mathsf M}\auGL_r(\auF_w)$ let $\tilde s\in
\auU(W)(\auA_f)$ and $\tilde s'\in\auU(V)(\auA_f)$ be integral elements, equal to $1$ outside
$\mathsf M\cup\aubar{\mathsf M}$, that reduce to $s$ and to $\operatorname{diag}(s,1_t)$ at
$\mathsf M$. The element $\tilde s'$ normalises $K_\vartheta$, and
$\varrho_V(s)\Xi=\psi_{\mathsf M}(\det s)\,R_{\tilde s'}\Xi$ defines a representation of $\auG$ on
$\poA(V;\vartheta)$, independent of the choice of $\tilde s'$; it preserves $\|\cdot\|_\sigma$ for
every $\sigma$, the scalar being a root of unity.

\begin{lemma}[Standing facts]\label{po:facts}
Let $\vartheta$ be a twist.
\begin{enumerate}
\item[(L2)] $\Theta_\vartheta\circ\varrho(s)=\varrho_V(s)\circ\Theta_\vartheta$ for $s\in\auG$.
Hence $\Theta_\vartheta(e_0\poA(W_j))\subseteq\varrho_V(e_0)\poA(V;\vartheta)$.
\item[(L3)] Let $w$ be off the level, with residue field of $q$ elements, $P_w(X)=\sum_{i\ge0}
(-X)^iq^{\binom i2}T_{w,i}$ the operator of (N3) on $\poA(V;\vartheta)$, and $\alpha_w$ the number
with $\omega_\mu(1,d_1)\plone=\alpha_w\,\plone(d_1^{-1}T)$ at $w$ (by \eqref{po:eq-w5}, $\alpha_w=\eta_w(1,d_1)$ does not depend on $\vartheta$, and
$|\alpha_w|=q^{r/2}$). Put $X_w=(\alpha_wq^{t-1})^{-1}$. Then $P_w(X_w)\,\Theta_\vartheta y=0$ for
every $y\in\poA(W_j)$.
\end{enumerate}
\end{lemma}
\begin{proof}
(L2) By \eqref{po:eq-w5} for integral elements, $\omega_\mu(\tilde s,\tilde s')\Phi_\vartheta=
\psi_{\mathsf M}(\det s)^{-1}\Phi_\vartheta$: the unramified character is trivial, and
$\plone[\tilde s'^{-1}T\tilde s\equiv\iota_0]=\plone[T\equiv\iota_0]$ because
$\operatorname{diag}(s,1)\iota_0s^{-1}=\iota_0$. So by \eqref{po:eq-equiv}
$R_{\tilde s'}\Theta_\vartheta y=\Theta(\omega_\mu(1,\tilde s')\Phi_\vartheta)y=\psi_{\mathsf M}
(\det s)^{-1}\Theta(\omega_\mu(\tilde s^{-1},1)\Phi_\vartheta)y=\psi_{\mathsf M}(\det s)^{-1}
\Theta_\vartheta(R_{\tilde s}y)$.

(L3) The datum $\Phi_\vartheta$ is the vacuum at $w,\aubar w$, and $\mu$ is unramified there. By
\eqref{po:eq-equiv}, $T_{w,i}\Theta_\vartheta y=\Theta(T'_i\Phi_\vartheta)y$ with $T'_i=\sum_{h\in
K'd_iK'/K'}\omega_\mu(1,h)$, $K'=\auGL_n(\auO_w)$, as in the proof of {\cite[Proposition~6.9]{P49}}.
The computation in the proof of {\cite[Lemma~6.8]{P49}} (with $\alpha=\alpha_w$, and $n$, $n'$ there
equal to $r$, $n$ here) gives, as functions of $T\in M_{n\times r}(E_w)$,
\[
\sum_i(-X)^iq^{\binom i2}T'_i\,\plone=\prod_{l<n-\operatorname{rk}(T\bmod w)}(1-q^l\alpha_wX)\cdot
\plone .
\]
Since $\operatorname{rk}(T\bmod w)\le r$, the factor with $l=t-1$ occurs for every $T$, and it
vanishes at $X=X_w$.
\end{proof}

\begin{enumerate}
\item[(L4)] \emph{If $\Xi\in\varrho_V(e_0)\poA(V;\vartheta)$ is orthogonal to
$\Theta_\vartheta(e_0\poA(W_j))$, then $D_\vartheta\Xi=0$.}
\end{enumerate}
\noindent\emph{Proof of {\rm(L4)}.} This is Proposition~\ref{plf:L4}.


\subsection{Two facts of linear algebra}\label{po:ss-linalg}

\begin{lemma}\label{po:det}
Let $u\colon U\to H$ be an injective linear map of finite-dimensional complex inner product
spaces, $\popr$ the orthogonal projection of $H$ onto $u(U)$, and $T\in\auEnd(H)$. Then
$\det(u^*u)\cdot\det\bigl(\popr T|_{u(U)}\bigr)=\det(u^*Tu)$.
\end{lemma}
\begin{proof}
Let $e\colon U\to u(U)$ be the isomorphism induced by $u$. Since $u^*$ vanishes on $u(U)^\perp$,
$u^*Tu=u^*\popr Tu=(u^*u)\,e^{-1}(\popr T|_{u(U)})\,e$; take determinants.
\end{proof}

\begin{lemma}\label{po:eig}
Let $L$ be a finite-dimensional complex vector space, $S\in\auEnd(L)$, $\phi\colon\C\to F$ a ring
homomorphism to a field with an absolute value, and $B\ge0$. If $|\phi(\alpha)|\le B$ for every
eigenvalue $\alpha$ of $S$, then $|\phi(\det S)|\le B^{\dim L}$.
\end{lemma}
\begin{proof}
$\det S$ is the product of the roots of the characteristic polynomial, with multiplicities.
\end{proof}

\subsection{Conjugate pairings, Hecke operators, false lifts}\label{po:ss-false}

For $\sigma\in\auAut(\C)$ put $c_\sigma=\sigma^{-1}\circ c\circ\sigma$, an involution of the field
$\C$ that inverts roots of unity, and for $\Xi,\Xi'\in\poA(V;\vartheta)$
\[
\langle\Xi,\Xi'\rangle_\sigma=\sigma^{-1}\Bigl(\sum_{h}\#(\Gamma\cap hK_\vartheta h^{-1})^{-1}\,
\aubar{\sigma(\Xi(h))}^{\,t}\,G\,\sigma(\Xi'(h))\Bigr),
\]
where $\sigma$ is applied to the coordinates in the basis of $W_{\lambda_V}(\Z)$ and $G$ is the
integral Gram matrix. This is well defined: $\sigma(\gamma\xi)=\gamma^\sigma\sigma(\xi)$, where
$\gamma^\sigma$ acts through the embeddings $\sigma\circ\tau\circ c$, and these matrices are again
unitary, complex conjugation commuting with all embeddings of the CM field $E$; and $G$ is
invariant under unitary matrices. The form is linear in $\Xi'$, $c_\sigma$-linear in $\Xi$,
$\langle\Xi',\Xi\rangle_\sigma=c_\sigma\langle\Xi,\Xi'\rangle_\sigma$, and
$\sigma\langle\Xi,\Xi\rangle_\sigma=\|\Xi\|_\sigma^2$, which is positive for $\Xi\ne0$. For
$\sigma=1$ it is the inner product. Consequently every operator $A$ has a unique adjoint
$A^{\dagger_\sigma}$ for $\langle\ ,\,\rangle_\sigma$, and
\begin{equation}\label{po:eq-cs}
\bigl|\sigma\langle\Xi,\Xi'\rangle\bigr|\le\|\Xi\|_{c\sigma c}\,\|\Xi'\|_\sigma
\end{equation}
by the Cauchy--Schwarz inequality, since $\sigma(\aubar z)=\aubar{(c\sigma c)(z)}$.

\begin{lemma}\label{po:adj}
Let $w$ be off the level and $\sigma\in\auAut(\C)$.
\begin{enumerate}
\item[(a)] $T_{w,i}^{\dagger_\sigma}=T^*_{w,i}$, the operator of the double coset
$K_wd_i^{-1}K_w$; it does not depend on $\sigma$, and the operators $T_{w,i}$, $T^*_{w,i'}$ all
commute.
\item[(b)] For $X\in\C$ the operator $A=P_w(X)$ commutes with $A^{\dagger_\sigma}$. Hence
$\langle Av,z\rangle_\sigma=0$ for all $v$ whenever $Az=0$; and (for $\sigma=1$) every form
orthogonal to the image of $A$ is killed by $A$.
\item[(c)] For $s\in\auG$, $\varrho_V(s)^{\dagger_\sigma}=\varrho_V(s^{-1})$; for every central
idempotent $e$ of $\C[\auG]$, $\varrho_V(e)^{\dagger_\sigma}=\varrho_V(e)$. The operators
$\varrho_V(s)$ commute with the $T_{w,i}$.
\end{enumerate}
\end{lemma}
\begin{proof}
(a) $(T_{w,i}\Xi)(h)=\sum_{x\in K_wd_iK_w/K_w}\Xi(hx)$ has integer coefficients, so $\sigma\circ
T_{w,i}\Xi=T_{w,i}(\sigma\circ\Xi)$. The sum defining the pairing is a constant times the integral
over $\Gamma\backslash\auU(V)(\auA_f)$ of the function $B(h)=\aubar{\sigma\Xi(h)}^{\,t}G\,
\sigma\Xi'(h)$, and for right $K_w$-invariant functions $a,b$ on $\Gamma\backslash\auU(V)(\auA_f)$,
with the Haar measure of $\auGL_n(E_w)$ giving $K_w$ the volume $1$,
\[
\int\sum_{x\in K_wd_iK_w/K_w}a(hx)\,b(h)\,dh=\int\!\!\int_{K_wd_iK_w}a(hx)b(h)\,dx\,dh=
\int a(h)\sum_{x'\in K_wd_i^{-1}K_w/K_w}b(hx')\,dh .
\]
The Hecke algebra of $(\auGL_n(E_w),K_w)$ is commutative: transposition is an anti-automorphism
that preserves $K_w$ and every double coset (elementary divisors).
(b) By (a), $\langle P_w(X)\Xi,\Xi'\rangle_\sigma=\langle\Xi,A'\Xi'\rangle_\sigma$ with
$A'=\sum_i(-c_\sigma X)^iq^{\binom i2}T^*_{w,i}$, and $A'$ commutes with $A$. If $Az=0$ then
$\langle A'z,A'z\rangle_\sigma=\langle AA'z,z\rangle_\sigma=\langle A'Az,z\rangle_\sigma=0$, so
$A'z=0$ and $\langle Av,z\rangle_\sigma=\langle v,A'z\rangle_\sigma=0$. If $u$ is orthogonal to
the image of $A$, then $\langle Au,Au\rangle=\langle u,A'Au\rangle=\langle u,AA'u\rangle=0$.
(c) Right translation by $\tilde s'$ commutes with $\sigma$ and preserves the sum over the
classes, and $c_\sigma$ inverts the root of unity $\psi_{\mathsf M}(\det s)$; so $\langle\varrho_V(s)\Xi,\varrho_V(s)\Xi'\rangle_\sigma=\langle\Xi,\Xi'\rangle_\sigma$.
A central idempotent is a sum of idempotents $e_\chi=\chi(1)\#\auG^{-1}\sum_s\chi(s^{-1})s$, and
$c_\sigma(\chi(s^{-1}))=\chi(s)$ because $\chi(s^{-1})$ is a sum of roots of unity; hence
$\varrho_V(e_\chi)^{\dagger_\sigma}=\chi(1)\#\auG^{-1}\sum_s\chi(s)\varrho_V(s^{-1})=
\varrho_V(e_\chi)$. The last assertion holds because $\tilde s'$ is $1$ at $w$ and $\aubar w$.
\end{proof}

\begin{definition}[False lifts]\label{po:false}
Let $\poZ\subseteq\poA(V;\vartheta)$ be a subspace and $e$ a central idempotent of $\C[\auG]$. For
$w$ off the level, the \emph{false lifts at $w$ for $(e,\poZ)$} are the elements of
\[
\mathcal F_w(e,\poZ)=\varrho_V(e)\poA(V;\vartheta)\ \cap\ \ker P_w(X_w)\ \cap\ \poZ^\perp .
\]
A form is \emph{orthogonal to the false lifts for $(e,\poZ)$} if it is orthogonal to
$\mathcal F_w(e,\poZ)$ for every $w$ off the level.
\end{definition}

Thus a false lift is a form of the same type as the lifts at the chosen places, with the one
Hecke equation (L3) of a lift at one split place, and orthogonal to $\poZ$. For
$\poZ=\poZ_\vartheta:=\Theta_\vartheta(e_0\poA(W_j))$ and $e=e_0$ this is the notion used in
(R3) of {\cite[Definition~14.39]{P46}}.

\begin{proposition}[The projection does not increase any conjugate norm]\label{po:proj}
Let $e$ be a central idempotent of $\C[\auG]$, $w$ off the level, and $\poZ\subseteq\varrho_V(e)
\poA(V;\vartheta)\cap\ker P_w(X_w)$. Let $y\in\poA(V;\vartheta)$ be orthogonal to
$\mathcal F_w(e,\poZ)$, and let $\popr_\poZ$ be the orthogonal projection onto $\poZ$ (for the inner
product, $\sigma=1$). Then $\|\popr_\poZ y\|_\sigma\le\|y\|_\sigma$ for every $\sigma\in\auAut(\C)$.
\end{proposition}
\begin{proof}
Put $A=P_w(X_w)$, $\varepsilon=\varrho_V(e)$ and $R=\ker\varepsilon+\poran A$. By
Lemma~\ref{po:adj}(c) $\varepsilon$ is self-adjoint, so $(\ker\varepsilon)^\perp=\poran\varepsilon$,
and by Lemma~\ref{po:adj}(b) $(\poran A)^\perp\subseteq\ker A$. Hence
$(\poZ+R)^\perp\subseteq\poZ^\perp\cap\poran\varepsilon\cap\ker A=\mathcal F_w(e,\poZ)$, and
therefore $y\in\mathcal F_w(e,\poZ)^\perp\subseteq(\poZ+R)^{\perp\perp}=\poZ+R$. Write $y=z+x$ with
$z\in\poZ$ and $x=x_1+Av$, $\varepsilon x_1=0$. Let $z'\in\poZ$ and $\sigma\in\auAut(\C)$; then
$z'=\varepsilon z'$ and $Az'=0$, so by Lemma~\ref{po:adj}
\[
\langle x,z'\rangle_\sigma=\langle x_1,\varepsilon z'\rangle_\sigma+\langle Av,z'\rangle_\sigma=
\langle\varepsilon x_1,z'\rangle_\sigma+0=0 .
\]
For $\sigma=1$ this shows $x\perp\poZ$, that is, $z=\popr_\poZ y$. For general $\sigma$ it gives
$\langle z,x\rangle_\sigma=c_\sigma\langle x,z\rangle_\sigma=0$ and
$\langle z+x,z+x\rangle_\sigma=\langle z,z\rangle_\sigma+\langle x,x\rangle_\sigma$; applying
$\sigma$, $\|y\|_\sigma^2=\|z\|_\sigma^2+\|x\|_\sigma^2\ge\|z\|_\sigma^2$.
\end{proof}

\begin{definition}[Denominators]\label{po:denom}
Let $\iota\in\auI$ with prime $\ell$, and $\beta\ge0$. The twist $\vartheta$ \emph{has
denominators at most $\beta$ at $\iota$ for the $j$-th lift} if there is an endomorphism $\Lambda$
of $\poZ_\vartheta=\Theta_\vartheta(e_0\poA(W_j))$ with $|\det\Lambda|_\iota\ge\exp(-\dim
\poZ_\vartheta\cdot r\beta)$ such that $\Lambda(\popr_{\poZ_\vartheta}y)$ is integral at $\iota$
for every $y\in\poA(V;\vartheta)$ that is integral at $\iota$ and orthogonal to the false lifts
for $(e_0,\poZ_\vartheta)$.
\end{definition}

This is the first clause of (R3), with Definition~\ref{po:false} for the words in brackets there.

\subsection{Tight frames and products of weights}\label{po:ss-frame}

\begin{lemma}[An integral tight frame]\label{po:frame}
Let $\mu$ be a family of diagrams $(\mu_\tau)_{\tau\in\auPl}$ for $\auGL_N$, with columns of
heights at most $N$, $W_\mu=\bigotimes_\tau W_{\mu_\tau}$ with its lattice $W_\mu(\Z)$ and its
invariant form, and $H=\prod_\tau H_{\mu_\tau}$ the product of the hook numbers of
{\cite[Lemma~5.22]{P49}}. There are vectors $\eta_U\in W_\mu(\Z)$, indexed by the fillings $U$ of $\mu$
(for each $\tau$ and each column a set of rows), such that
\[
\sum_U\langle\eta_U,\xi\rangle\,\eta_U=H^2\,\xi\quad(\xi\in W_\mu),\qquad\|\eta_U\|\le H .
\]
The number of fillings is at most $N^{|\mu|}$, and $H\le|\mu|^{3|\mu|}$, where $|\mu|$ is the
total number of boxes.
\end{lemma}
\begin{proof}
In the notation of {\cite[\S5.2.6]{P49}} (for all $\tau$ at once) $W_\mu$ is identified
isometrically, through $\mathsf m$, with $\mathsf C_\mu\subseteq\poT_\mu(\C)$, the vectors $e_U$
form an orthonormal basis of $\poT_\mu(\C)$, $\pi_\mu$ is the orthogonal projection onto
$\mathsf C_\mu$, and $W_\mu(\Z)=\mathsf m(\mathsf C_\mu\cap\poT_\mu(\Z))$ by {\cite[(51)]{P49}}.
Put $\eta_U=H\,\mathsf m(\pi_\mu e_U)=H\,b_U$. By {\cite[Lemma~5.22]{P49}}, $H\Sigma_\mu\subseteq
\Sigma_\mu^\#$, so $\eta_U\in W_\mu(\Z)$. For $\xi=\mathsf m(v)$, $v\in\mathsf C_\mu$:
$\sum_U\langle\eta_U,\xi\rangle\eta_U=H^2\,\mathsf m\,\pi_\mu\bigl(\sum_U(e_U,v)e_U\bigr)=H^2\,
\mathsf m(v)$, and $\|\eta_U\|=H\|\pi_\mu e_U\|\le H$. A filling chooses a row for each box.
\end{proof}

\emph{Products.} Let $n=N''+n^\circ$ be a decomposition of the coordinates, $\C^n=\C^{N''}\oplus
\C^{n^\circ}$, and suppose $N''\ge r_+$, $n^\circ\ge r_-$. Let $W''=W_{\mu^+}(\C^{N''})\subseteq
\poT^+(\C^{N''})$ and $W^\circ=W_{\mu^-}(\C^{n^\circ})\subseteq\bigotimes_{\mathfrak c}\Lambda^{
h_{\mathfrak c}}\C^{n^\circ}$ (columns of $\mu^-$), in the tensor models of {\cite[\S5.2.6]{P49}}.
Let $\natural\colon\C^{n^\circ}\to(\C^{n^\circ})^\vee$ be the \emph{linear} map $e_i\mapsto
e_i^\vee$, and $\flat$ the conjugate-linear map of the standard hermitian form; both are applied
to each vector of each factor of a tensor. For $u''\in W''$ and $u^\circ\in W^\circ$ put
\[
u''\odot u^\circ=u''\otimes(u^\circ)^\natural\ \in\poT^+(\C^n)\otimes\poT^-(\C^n),
\]
a bilinear product.

\begin{lemma}\label{po:prod}
\begin{enumerate}
\item[(M1)] $u''\odot u^\circ\in W_{\lambda_V}$; $\|u''\odot u^\circ\|=\|u''\|\,\|u^\circ\|$; and
$\operatorname{diag}(k'',k^\circ)(u''\odot u^\circ)=k''u''\odot\aubar{k^\circ}u^\circ$ for unitary
$k'',k^\circ$.
\item[(M2)] $W''(\Z)\odot W^\circ(\Z)\subseteq W_{\lambda_V}(\Z)$, and for $\sigma\in\auAut(\C)$,
$\sigma(u''\odot u^\circ)=\sigma u''\odot\sigma u^\circ$, so that
$\|\sigma(u''\odot u^\circ)\|=\|\sigma u''\|\,\|\sigma u^\circ\|$.
\item[(M3)] For $x=(x''\,;\,x^\circ)\in M_{n\times r}(\C)$, with $x''$ the rows in $\C^{N''}$ and
$x^\circ$ those in $\C^{n^\circ}$, $\pop_{u''\odot u^\circ}(x)=\pop''_{u''}(x'')\cdot
\pop^\circ_{u^\circ}(x^\circ)$, where $\pop''_{u''}(x'')=\langle u'',\mathrm{ev}^+(\aubar{x''^+})
\rangle$ depends only on the positive columns of $x''$ and $\pop^\circ_{u^\circ}(x^\circ)=\langle
(u^\circ)^\natural,\mathrm{ev}^-(x^{\circ-})^\flat\rangle$ only on the negative columns of
$x^\circ$; both are conjugate-linear in the vector.
\end{enumerate}
\end{lemma}
\begin{proof}
(M1) $W''$ is spanned by the vectors $\bigotimes_{\mathfrak c}y_1\wedge\dots\wedge y_{h_{\mathfrak
c}}$ with $y_i\in\C^{N''}$ and $W^\circ$ by the analogous vectors with $y'_i\in\C^{n^\circ}$; for
such vectors $u''\odot u^\circ=\mathrm{ev}(y,y'^\natural)$, and $y'^\natural_{i'}(y_i)=0$ because
the two sets of coordinates are disjoint. So $u''\odot u^\circ\in W_{\lambda_V}$ by
\eqref{po:eq-ev}. The norm of a tensor product is the product of the norms, $\natural$ maps an
orthonormal basis to an orthonormal basis, and the contragredient of a unitary matrix $k^\circ$
is $\aubar{k^\circ}$, so that $(\aubar{k^\circ}u^\circ)^\natural$ is the image of
$(u^\circ)^\natural$ under $k^\circ$. (M2) The product is given by the same formula over $\Z$;
use the description of the integral structures as intersections with the integral tensors.
(M3) By \eqref{po:eq-poly}, $\pop_{u''\odot u^\circ}(x)=\langle u'',\mathrm{ev}^+(\aubar{x^+})\rangle
\cdot\langle(u^\circ)^\natural,\mathrm{ev}^-(x^-)^\flat\rangle$; the first pairing involves only the
coordinates of $\C^{N''}$ because $u''$ does, the second only those of $\C^{n^\circ}$.
\end{proof}

\subsection{Gluing}\label{po:ss-glue}

\emph{The splitting.} Put $N''=r+9m+1$ and $n^\circ=t-9m-1$, and split the coordinates of $V$ as
$V=V''\perp V^\circ$, where $V''$ has the first $N''$ coordinates (so $\iota_0$ maps into $V''$;
write $\iota_0''$ for the corresponding $N''\times r$ matrix) and $V^\circ$ the last $n^\circ$; in
$V^\circ$ we mark the first $r$ coordinates and write $\iota_z$ for the $n^\circ\times r$ matrix
of their inclusion, and $v_z$ for the block of marked rows of an $n^\circ\times r$ matrix $v$. Note
$r\le n^\circ$ and $r_+\le N''$. Let $\auO_{\mathsf M}=\prod_{w\in\mathsf M}\auO_w$; recall that
$\pi_M$ is a uniformiser at every $w\in\mathsf M$ and a unit at $\aubar w$. Let $\Gamma^1$ be the
kernel of $\auGL_r(\auO_E/\pi_M^2)\to\auGL_r(\auO_E/\pi_M)$; for $Y\in\Gamma^1$ let
$\tilde Y\in\auU(W)(\auA_f)$ be an integral lift, equal to $1$ outside $\mathsf M\cup
\aubar{\mathsf M}$; it lies in $J_M$.

For a matrix $s\in M_{3m\times N''}(\auO_E/\pi_M^2)$ let $L''_s=\{u\in\auO_{\mathsf M}^{N''}:
su\equiv0\bmod\pi_M\}$, a lattice with $\pi_M\auO_{\mathsf M}^{N''}\subseteq L''_s$ and with
$\auO_w^{N''}/L''_{s,w}$ of dimension $d_w\le3m$ over $\auF_w$. A matrix $x=(u\,;v)\in
M_{n\times r}(\auO_{\mathsf M})$, with $u$ the rows in $V''$ and $v$ those in $V^\circ$, is
\emph{normalised for $s$} if
\[
v\equiv\iota_z\bmod\pi_M^2\qquad\text{and}\qquad u\in\iota_0''+\pi_ML''_s\ \ (\text{columnwise}).
\]

\begin{lemma}[The glued lattice]\label{po:lattice}
For every $s$ there is $\pog=\pog_s\in\auU(V)(\auA_{E^+,f})$, equal to $1$ outside
$\mathsf M\cup\aubar{\mathsf M}$, such that at the places of $\mathsf M$:
\begin{enumerate}
\item[(a)] $\pog$ and $\pi_M\pog^{-1}$ are integral;
\item[(b)] for $x\in M_{n\times r}(E_{\mathsf M})$: $\pog^{-1}x$ is integral and $\equiv\iota_0
\bmod\pi_M$ if and only if $x$ is integral and $xY$ is normalised for $s$ for some
$Y\in\Gamma^1$; this $Y$ is then unique;
\item[(c)] there is $K_1\in\auGL_{N''}(\auO_{\mathsf M})$ such that, at each $w\in\mathsf M$, all
columns of $\pog^{-1}|_{V''}K_1$ except at most $3m$ are integral, and those have entries in
$\pi_M^{-1}\auO_w$;
\item[(d)] if $k\in\auGL_{N''}(\auO_{\mathsf M})$, $k\equiv1\bmod\pi_M$, $sk\equiv s\bmod\pi_M^2$,
and $k^\circ\in\auGL_{n^\circ}(\auO_{\mathsf M})$, $k^\circ\equiv1\bmod\pi_M^2$, then
$\pog^{-1}\operatorname{diag}(k,k^\circ)\pog$ is integral and $\equiv1\bmod\pi_M$.
\end{enumerate}
\end{lemma}
\begin{proof}
Let $B\in M_{N''}(\auO_{\mathsf M})$ be a matrix whose columns are a basis of $L''_s$, and
$D=\pi_MB^{-1}$; it is integral because $\pi_M\auO^{N''}\subseteq L''_s$, and $Du\equiv0\bmod
\pi_M^2$ if and only if $u\in\pi_ML''_s$. Let $\Delta_z$ be the diagonal $n^\circ\times n^\circ$
matrix with $\pi_M$ at the marked coordinates and $1$ elsewhere, $P$ the permutation matrix that
exchanges the first $r$ coordinates of $V''$ with the marked coordinates of $V^\circ$, and
\[
G=\begin{pmatrix}D&-D\,\iota_0''\,{}^t\iota_z\\0&\Delta_z\end{pmatrix},\qquad
G^{-1}=\begin{pmatrix}D^{-1}&\iota_0''\,{}^t\iota_z\Delta_z^{-1}\\0&\Delta_z^{-1}\end{pmatrix},
\qquad\pog_w=\pi_MG^{-1}P^{-1}\quad(w\in\mathsf M),
\]
with the components at $\aubar{\mathsf M}$ determined by unitarity. (a) $\pi_M\pog^{-1}=PG$ is
integral, and so is $\pi_MG^{-1}$, whose blocks are $B$, $\iota_0''\,{}^t\iota_z$ and
$\pi_M\Delta_z^{-1}$. (b) If $\pog^{-1}x$ is integral then $x=\pog(\pog^{-1}x)$ is integral. For
integral $x=(u;v)$ the condition is $Gx\equiv\pi_MP^{-1}\iota_0\bmod\pi_M^2$, where
$P^{-1}\iota_0=(0;\iota_z)$; that is,
\[
D(u-\iota_0''v_z)\equiv0,\qquad\pi_Mv_z\equiv\pi_M1_r,\qquad v-\iota_zv_z\equiv0\pmod{\pi_M^2}.
\]
If this holds, then $v_z\equiv1\bmod\pi_M$, $Y=v_z^{-1}\bmod\pi_M^2$ lies in $\Gamma^1$, and
$xY=(uY;vY)$ has $vY\equiv\iota_z$ and $D(uY-\iota_0'')\equiv D(u-\iota_0''v_z)Y\equiv0$, so it is
normalised. Conversely, if $xY$ is normalised with $Y\in\Gamma^1$, then $v_z\equiv Y^{-1}\bmod
\pi_M^2$, which determines $Y$, and the three congruences follow in the same way. (c) By the
theorem of elementary divisors $B=K_1\operatorname{diag}(\pi_M1_{d_w},1)K_2$ at $w$ with
$K_1,K_2\in\auGL_{N''}(\auO_w)$; and $\pog^{-1}|_{V''}=P\binom{B^{-1}}0$. (d) Write $k=1+\pi_MX$;
then $sX\equiv0\bmod\pi_M$, so $X=BX_1$ with $X_1$ integral, and $k^\circ=1+\pi_M^2X^\circ$. A
computation with the blocks gives
$\pog^{-1}\operatorname{diag}(k,k^\circ)\pog=1+\pi_M\,P\,G\operatorname{diag}(X,\pi_MX^\circ)
\,\pi_MG^{-1}P^{-1}\cdot\pi_M^{-1}$, and $G\operatorname{diag}(X,\pi_MX^\circ)\pi_MG^{-1}$ has the
blocks $\pi_MX_1B$, $\pi_MX_1\iota_0''\,{}^t\iota_z-\pi_M^2D\iota_0''\,{}^t\iota_zX^\circ
\Delta_z^{-1}$, $0$ and $\pi_M^2\Delta_zX^\circ\Delta_z^{-1}$, all divisible by $\pi_M$.
\end{proof}

\emph{Factorisation of the Weil representation.} $\auHom(W,V)=\auHom(W,V'')\oplus\auHom(W,
V^\circ)$, and
\begin{enumerate}
\item[(W7)] $\mathcal H(\auHom(W,V))=\mathcal H(\auHom(W,V''))\otimes\mathcal H(\auHom(W,
V^\circ))$, $\theta=\theta''\otimes\theta^\circ$, and $\omega_\mu(g,\operatorname{diag}(h,h^\circ))
=\omega''(g,h)\otimes\omega^\circ(g,h^\circ)$, where $\omega''$ and $\omega^\circ$ are Weil
representations of the two pairs: $\omega^\circ=\omega_{\upsilon^{m/2-j},\beta^\circ}$ and
$\omega''=\omega_{\nu''_\vartheta,\beta''}$ in the convention of {\cite[\S2.3]{P53}}, with
$\nu''_\vartheta=\mu\,\upsilon^{j+(t+k-1)/2}$ and $\beta^\circ,\beta''$ powers of $\upsilon$ times the
characters on $\auU(V)$, so that the characters on $\auU(W)$ multiply to $\mu\upsilon^{\lfloor
n/2\rfloor}$ (Lemma~\ref{psv:char}). In particular $\omega^\circ$ does not depend on $\vartheta$, and
its character on $\auU(W)$ is $\upsilon^{a^\circ}$ with $a^\circ=m/2-j$.
\end{enumerate}
This is the compatibility of the Heisenberg and Weil representations with orthogonal sums, and of
the splittings of~\cite{KudlaSplitting} with it (the splitting over $\auU(W)$ for $\auHom(W,V)$
is the product of those for the two summands, because the characters multiply); the theta
functionals are normalised compatibly. In the Fock model the vector $\pop''_{u''}\pop^\circ_{
u^\circ}$ of Lemma~\ref{po:prod}(M3) is the tensor product of the two polynomial vectors.

\emph{Forms on $V^\circ$, and the kernels.} Let $J^\circ\subset\auU(V^\circ)(\auA_f)$ be the
principal congruence subgroup of level $\pi_M^2$ in the stabiliser of the standard lattice; it is
free by {\cite[Lemma~14.5]{P46}}. Let $\poA^\circ$ be the space of the $\zeta\colon\auU(V^\circ)
(\auA_f)\to W^\circ$ that are right invariant under $J^\circ$ and satisfy $\zeta(\gamma h^\circ)=
\rho_{\mu^-}(\tau(\gamma))\zeta(h^\circ)$ (the rational points acting through the embeddings
$\tau$ themselves); equivalently, the coordinatewise conjugate $\aubar\zeta$ is an algebraic
modular form of weight $\mu^-$ and level $J^\circ$. The inner product is $\langle\zeta,\zeta'
\rangle=\sum_{h^\circ}\langle\zeta(h^\circ),\zeta'(h^\circ)\rangle$ over the classes,
$\|\zeta\|_\sigma^2=\sum_{h^\circ}\|\sigma\zeta(h^\circ)\|^2$, and $\zeta$ is \emph{integral at
$\iota$} if $\aubar\zeta$ is integral at $\aubar\iota$. (By Lemma~\ref{po:prod}(M1) and the
convention (N2), under which $\Gamma$ acts on $W_{\lambda_V}$ through $\tau\circ c$, the block of
$V^\circ$ acts on the second factor of $u''\odot u^\circ$ through $\tau$; this is why the forms on
$V^\circ$ are taken of this type.) Let
\[
\Psi^\circ=\plone[v\equiv\iota_z\bmod\pi_M^2]\otimes(\text{vacuum off }\mathsf M),\qquad
\Psi''_{\vartheta,s}=\plone[u\in\iota_0''+\pi_ML''_s]\otimes\bigotimes_{w\in\mathsf Q}\plone[u
\text{ integral}]\psi_w(\det u')\otimes(\text{vacuum}),
\]
\[
\theta^\circ(g,h^\circ;u^\circ)=\theta^\circ\bigl(\omega^\circ(g,h^\circ)(\pop^\circ_{u^\circ}
\otimes\Psi^\circ)\bigr),\qquad
\theta''_{\vartheta,s}(g,h;u'')=\theta''\bigl(\omega''(g,h)(\pop''_{u''}\otimes\Psi''_{\vartheta,
s})\bigr),
\]
both conjugate-linear in the vector. The \emph{kernel form} $K^\circ(g,\cdot)\in\poA^\circ$ is
defined by $\langle u^\circ,K^\circ(g,h^\circ)\rangle=\theta^\circ(g,h^\circ;u^\circ)$; it does not
depend on $\vartheta$.

\emph{Forms supported on a class.} For $x\in\auU(V)(\auA_f)$ and $v\in W_{\lambda_V}$ let
$\podel[x;v](h)=\sum\chi_\vartheta(k)^{-1}\gamma v$, the sum over the pairs $(\gamma,k)\in\Gamma\times
K_\vartheta$ with $h=\gamma xk$; there are $\#(\Gamma\cap xK_\vartheta x^{-1})$ such pairs if $h$
is in the class of $x$ and none otherwise. Then $\podel[x;v]\in\poA(V;\vartheta)$,
\begin{equation}\label{po:eq-delta}
\langle\podel[x;v],\Xi\rangle=\langle v,\Xi(x)\rangle,\qquad
D_\vartheta\podel[x;v](g)=\aubar{k_{\Phi_\vartheta}(g,x;v)},
\end{equation}
\begin{equation}\label{po:eq-deltasize}
\|\sigma\,\podel[x;v](h)\|\le n_{\mathrm{Aut}}\|\sigma v\|,\qquad\plgb h{\podel[x;v](h)}\le\plgb xv
\quad(\ell\nmid\auNm\mathsf Q),
\end{equation}
and $\podel[\gamma xk;\gamma v]=\chi_\vartheta(k)\,\podel[x;v]$ (the pairs $(\gamma',k')$ for $\gamma xk$ correspond to the pairs $(\gamma'\gamma,kk')$ for $x$). Indeed
$\Xi(x)=\chi_\vartheta(k)^{-1}\gamma\,\Xi(x)$ for each pair with $x=\gamma xk$, so each of the terms of
$\langle\podel[x;v](x),\Xi(x)\rangle$ is $\langle v,\Xi(x)\rangle$; in the same way each term of
$k_{\Phi_\vartheta}(g,x;\podel[x;v](x))$ equals $k_{\Phi_\vartheta}(g,x;v)$, by the invariance of
$\theta$ under $\Gamma$ and the definition of $\chi_\vartheta$; the bounds are (G1), (G2) of
{\cite[\S6.1]{P49}}, the unitarity of $\gamma^\sigma$, and $|\sigma\chi_\vartheta(k)|=1=|\chi_\vartheta
(k)|_\iota$.

\begin{lemma}[Slices]\label{po:slice}
Let $\pog=\pog_s$ and $c_\pog=\eta_{\mathsf M}(1,\pog)\ne0$, the scalar of \eqref{po:eq-w5} at
$\mathsf M$; it does not depend on $\vartheta$. For $h\in\auU(V'')(\auA_f)$, $h^\circ\in
\auU(V^\circ)(\auA_f)$, $u''\in W''$, $u^\circ\in W^\circ$ and $g\in\auU(W)(\auA)$,
\[
k_{\Phi_\vartheta}\bigl(g,\operatorname{diag}(h,h^\circ)\pog;u''\odot u^\circ\bigr)=c_\pog\sum_{Y\in
\Gamma^1}\theta''_{\vartheta,s}(g\tilde Y,h;u'')\ \theta^\circ(g\tilde Y,h^\circ;u^\circ).
\]
\end{lemma}
\begin{proof}
By \eqref{po:eq-w5} and Lemma~\ref{po:lattice}(b), at the places of $\mathsf M$,
\begin{align*}
\omega_\mu(1,\pog)\plone[T\equiv\iota_0]&=c_\pog\,\plone[\pog^{-1}T\equiv\iota_0]=c_\pog\sum_{Y\in
\Gamma^1}\plone[TY\text{ normalised for }s]\\
&=c_\pog\sum_{Y\in\Gamma^1}\omega_\mu(\tilde Y,1)\bigl(\plone[u\in
\iota_0''+\pi_ML''_s]\otimes\plone[v\equiv\iota_z\bmod\pi_M^2]\bigr);
\end{align*}
in the last step the scalar of \eqref{po:eq-w5} is $1$ because $\det\tilde Y\equiv1$ at
$\mathsf M$. At $\mathsf Q$ the block datum involves only rows of $V''$, and $\pog$, $\tilde Y$
are $1$ there. So $\omega_\mu(1,\pog)\Phi_\vartheta=c_\pog\sum_Y\omega_\mu(\tilde Y,1)(\Psi''_{
\vartheta,s}\otimes\Psi^\circ)$. Insert this in the definition of $k_{\Phi_\vartheta}$ and use
(W7) and Lemma~\ref{po:prod}(M3).
\end{proof}

\emph{The choices and the first factors.} There is a finite set
$I$ of \emph{choices}, independent of $\vartheta$ (Lemma~\ref{plf:I}); each $i\in I$ has a matrix $s_i\in M_{3m\times
N''}(\auO_E/\pi_M^2)$ (the matrix $\Lambda_i$ of \S\ref{pte:ss-third}), hence $\pog_i=\pog_{s_i}$, $c_i=c_{\pog_i}$, and a level
\[
\begin{aligned}J_{\vartheta,i}=\{k\in K^\circ(V''):\ &k\equiv1\bmod\pi_M,\ s_ik\equiv s_i\bmod\pi_M^2,\ \text{the last }N''-r
\text{ columns of }k\\&\text{are those of }1\bmod\pi_M^2\text{ at }\mathsf M\text{ and}\bmod w\text{ for }w\in\mathsf Q\},\end{aligned}
\]
which is free ({\cite[Lemma~14.5]{P46}}). (This is the group $J'_{\vartheta,i}$ of \S\ref{pte:ss-left}, Remark~\ref{pte:rem-final}(ii).) By Lemma~\ref{po:lattice}(d), $\pog_i^{-1}\operatorname{diag}
(k,k^\circ)\pog_i\in K_\vartheta$ for $k\in J_{\vartheta,i}$, $k^\circ\in J^\circ$; let
$\chi''_{\vartheta}(k)$ be the value of $\chi_\vartheta$ on it (it depends only on $k$ at
$\mathsf Q$). For each twist $\vartheta$ and each $i$ there is a \emph{first factor}
\[
\varphi''_{\vartheta,i}\in\poA_{\mu^+}\bigl(J_{\vartheta,i},\chi''_\vartheta{}^{-1}\bigr)\quad
\text{on }\auU(V''),\qquad\Xi_{\vartheta,i}(g)=\sum_{h\in\Gamma''\backslash\auU(V'')(\auA_f)/
J_{\vartheta,i}}\theta''_{\vartheta,s_i}\bigl(g,h;\varphi''_{\vartheta,i}(h)\bigr),
\]
the second being its theta series on $\auU(W)(\auA)$ (the summand does not depend on the
representative). The first factors $\varphi''_{\vartheta,i}$ are the functions of Definition~\ref{pfq:def-ff} (a
non-zero multiple $\kappa$, depending only on $j$, of the functions of Definition~\ref{ptf:def-ff}, through
Definition~\ref{pte:def-ff}). The character $\chi''_\vartheta$ is
$k\mapsto\chi'_\vartheta(\det\aubar k)$ (Lemma~\ref{ptf:chi}). The \emph{exponent of the denominators} is
\[
l_1=3m\sum_{\tau\in\auPl}\bigl(c^+_\tau+c^-_\tau\bigr),\qquad c^\pm_\tau=\text{the number of
columns of }\mu^\pm_\tau;\qquad l_1\le3mf(4b+m\Delta)\le13fmb .
\]

\begin{lemma}[First factors: integrality and size]\label{po:ff}
The first factors can be chosen so that $\#I\le e^b$ and, for all $i$, $\vartheta$ and $h$:
\begin{enumerate}
\item[(I1)] $\varphi''_{\vartheta,i}$ is integral at every $\iota$ with $\ell(\iota)\nmid\auNm
Q_\vartheta$;
\item[(I2)] $\|\sigma\varphi''_{1,i}(h)\|^2\le\exp(\poAA)$ for every $\sigma\in\auAut(\C)$.
\end{enumerate}
\end{lemma}
\noindent\emph{Proof.} Theorem~\ref{pfq:ff}, for the first factors of Definition~\ref{pfq:def-ff}
(a fixed non-zero multiple, depending only on $j$, of those of Definition~\ref{ptf:def-ff}); $\#I\le e^b$
is Lemma~\ref{plf:I}. Proposition~\ref{po:exchange} holds for these first factors by
Proposition~\ref{pfq:exchange}.


\begin{lemma}[Numbers of classes]\label{po:classes}
Let $X$ be a standard definite space of rank at most $2n_m$ and $J\subseteq K^\circ(X)$ a subgroup
containing the principal congruence subgroup of level $\pi_M^2$. Then $\Gamma\backslash\auU(X)
(\auA_f)/J$ has at most $e^b$ elements.
\end{lemma}
\noindent\emph{Proof.} Lemma~\ref{plf:classes}.


\begin{proposition}[Three exchanges]\label{po:exchange}
Let $H_j(\vartheta)$ be the number \eqref{po:eq-H} below. There are functions $F_i$ $(i\in I)$ on
$\auU(W)(\auA)$, independent of $\vartheta$, a constant $C_{\mathrm x}\ge0$, and for every twist
$\vartheta$ a number $c'_\vartheta\in\C$, such that
\[
c'_\vartheta\,\Xi_{\vartheta,i}=F_i\quad(i\in I),\qquad|c'_\vartheta|^2\,H_j(\vartheta)=
C_{\mathrm x}.
\]
Moreover $yF_iK^\circ(\cdot,h^\circ)$ is integrable on $[\auU(W)]$ for $y\in\poA(W_j)$, and so is
$y\,\Xi_{\vartheta,i}K^\circ(\cdot,h^\circ)$.
\end{proposition}
\begin{proof}
This is proved in \S\ref{pte:sec} and \S\ref{ptf:sec}: the functions $F_i$ are those of \eqref{pte:eq-F}, the
number is $c'_\vartheta=c_3(\vartheta)\gamma(\vartheta)/\allowbreak \#\auGL_{3m}(\auO_E/\pi_M^2)$ and $C_{\mathrm x}=C_3C_{12}/
\allowbreak (\#\auGL_{3m}(\auO_E/\pi_M^2))^2$ (\S\ref{pte:ss-proof}), with the first factors of Definition~\ref{ptf:def-ff}
and the level $J'_{\vartheta,i}$ (Remark~\ref{pte:rem-final}). For the first factors of
Definition~\ref{pfq:def-ff}, which are the ones used from now on, the number is
$c'_\vartheta=c_3(\vartheta)\gamma(\vartheta)/(\aubar\kappa\,\#\auGL_{3m}(\auO_E/\pi_M^2))$ and the constant
$C_{\mathrm x}=C_3C_{12}/(|\kappa|^2(\#\auGL_{3m}(\auO_E/\pi_M^2))^2)$ (Proposition~\ref{pfq:exchange}).
Below, $c'_\vartheta$ and $C_{\mathrm x}$ denote these numbers; the two identities of the proposition
hold for them, and they are the only properties of $c'_\vartheta$, $C_{\mathrm x}$ used.
\end{proof}


\begin{proposition}[Products span]\label{po:span}
If $y\in e_0\poA(W_j)$ and $\int_{[\auU(W)]}y(g)F_i(g)K^\circ(g,h^\circ)\,dg=0$ for all $i\in I$
and all $h^\circ$, then $y=0$.
\end{proposition}
\noindent Proposition~\ref{po:span} is Proposition~\ref{psp:main}, proved in \S\ref{psp:sec} with the inputs
of \S\ref{psq:sec}--\S\ref{psv:sec} (Theorem~\ref{psv:main}).


\begin{proposition}[Periods of false lifts]\label{po:periods}
Let $w$ be off the level and $\Xi\in\mathcal F_w(e_0,\poZ_\vartheta)$ a false lift. Then for every
$s$, every $h\in\auU(V'')(\auA_f)$, $u''\in W''$ and $g\in\auU(W)(\auA)$,
\[
\sum_{h^\circ\in\Gamma^\circ\backslash\auU(V^\circ)(\auA_f)/J^\circ}\bigl\langle u''\odot K^\circ(g,
h^\circ),\ \Xi\bigl(\operatorname{diag}(h,h^\circ)\pog_s\bigr)\bigr\rangle=0 .
\]
\end{proposition}
\begin{proof}[Proof of Proposition~\ref{po:periods}]
This is Theorem~\ref{ppe:main} of \S\ref{ppe:sec}: a false lift has vanishing descent (by
Proposition~\ref{plf:L4}), is killed by the hyperplane sums at the chosen places (dictionary between
labels and characters, Lemma~\ref{ppe:green}), and satisfies one Hecke equation at a split place;
its descents then vanish to $W$ (Proposition~\ref{ppe:bottom}), to $W\oplus\mathbb H^p$ for
$2r+2p+3<n$ (Proposition~\ref{ppe:tower}), and the periods over $\auU(V^\circ)$ vanish
(Proposition~\ref{ppe:rallis}).
\end{proof}


\emph{The gluing map.} $\poA^\oplus=\bigoplus_{i\in I}\poA^\circ$, with $\|\zeta\|_\sigma^2=\sum_i
\|\zeta_i\|_\sigma^2$. For a twist $\vartheta$,
\[
\poglue_\vartheta\colon\poA^\oplus\to\poA(V;\vartheta),\qquad
\poglue_\vartheta(\zeta)=M^{l_1}\sum_{i\in I}\ \sum_{h}\ \sum_{h^\circ}\podel\bigl[\operatorname{diag}
(h,h^\circ)\pog_i;\ \varphi''_{\vartheta,i}(h)\odot\zeta_i(h^\circ)\bigr],
\]
where $h$ runs over $\Gamma''\backslash\auU(V'')(\auA_f)/J_{\vartheta,i}$ and $h^\circ$ over
$\Gamma^\circ\backslash\auU(V^\circ)(\auA_f)/J^\circ$; the summands do not depend on the
representatives, by Lemma~\ref{po:prod}(M1), Lemma~\ref{po:lattice}(d), the rule for
$\podel[\gamma xk;\gamma v]$ and the definition of $\chi''_\vartheta$. Let
\[
\Theta^\circ\colon e_0\poA(W_j)\to\poA^\oplus,\qquad(\Theta^\circ y)_i=M^{l_1}c_i\,\#\Gamma^1
\int_{[\auU(W)]}y(g)\,F_i(g)\,K^\circ(g,\cdot)\,dg .
\]

\begin{proposition}[Properties of the gluing map]\label{po:glue}
Let $e_0\poA(W_j)\ne0$.
\begin{enumerate}
\item[(G1)] $c'_\vartheta\cdot\poglue_\vartheta^*\circ\Theta_\vartheta=\Theta^\circ$ on
$e_0\poA(W_j)$, for every twist $\vartheta$.
\item[(G2)] $\Theta^\circ$ is injective; hence $c'_\vartheta\ne0$, $C_{\mathrm x}>0$, and
$\Theta_\vartheta$ is injective on $e_0\poA(W_j)$.
\item[(G3)] $\poglue_\vartheta(\Theta^\circ y)$ is orthogonal to the false lifts for
$(e_0,\poZ_\vartheta)$, for every $y$.
\item[(G4)] If $\ell(\iota)\nmid\auNm Q_\vartheta$ and $\zeta$ is integral at $\iota$, then
$\poglue_\vartheta(\zeta)$ is integral at $\iota$.
\item[(G5)] $\|\poglue_1(\zeta)\|_\sigma^2\le\exp(2\poAA)\,\|\zeta\|_\sigma^2$ for all $\sigma\in
\auAut(\C)$.
\end{enumerate}
\end{proposition}
\begin{proof}
(G1) Let $\zeta\in\poA^\oplus$. By \eqref{po:eq-descent}, the linearity of $D_\vartheta$,
\eqref{po:eq-delta} and Lemma~\ref{po:slice},
\[
\aubar{D_\vartheta(\poglue_\vartheta\zeta)(g)}=M^{l_1}\sum_ic_i\sum_{Y\in\Gamma^1}\Xi_{\vartheta,i}
(g\tilde Y)\sum_{h^\circ}\theta^\circ\bigl(g\tilde Y,h^\circ;\zeta_i(h^\circ)\bigr)=M^{l_1}\sum_ic_i
\sum_Y\Xi_{\vartheta,i}(g\tilde Y)\,\bigl\langle\zeta_i,K^\circ(g\tilde Y,\cdot)\bigr\rangle .
\]
Since $y$ is right invariant under $\tilde Y\in J_M$ and the measure is invariant,
\[
\langle\zeta,\poglue_\vartheta^*\Theta_\vartheta y\rangle=\langle\poglue_\vartheta\zeta,
\Theta_\vartheta y\rangle=M^{l_1}\sum_ic_i\,\#\Gamma^1\int_{[\auU(W)]}y(g)\,\Xi_{\vartheta,i}(g)\,
\langle\zeta_i,K^\circ(g,\cdot)\rangle\,dg ,
\]
and multiplying by $c'_\vartheta$ gives $\langle\zeta,\Theta^\circ y\rangle$ by
Proposition~\ref{po:exchange}.
(G2) If $\Theta^\circ y=0$ then, $c_i\ne0$, the hypothesis of Proposition~\ref{po:span} holds, so
$y=0$ (Remark~\ref{psv:G2}). If $c'_\vartheta=0$ then $\Theta^\circ=0$ by (G1), against
$e_0\poA(W_j)\ne0$; so $c'_\vartheta\ne0$, and $\Theta_\vartheta$ is injective on $e_0\poA(W_j)$ by
(G1). Finally $C_{\mathrm x}=|c'_\vartheta|^2H_j(\vartheta)$ (Proposition~\ref{po:exchange}) and
$H_j(\vartheta)>0$ (Lemma~\ref{po:Hpos}), so $C_{\mathrm x}>0$. Neither the injectivity of
$\Theta^\circ$ nor Lemma~\ref{po:Hpos} uses $C_{\mathrm x}$, so the argument is not circular.
(G3) Let $\Xi\in\mathcal F_w(e_0,\poZ_\vartheta)$. By \eqref{po:eq-delta},
\[
\langle\poglue_\vartheta\Theta^\circ y,\Xi\rangle=M^{l_1}\sum_{i,h}\ \sum_{h^\circ}\bigl\langle
\varphi''_{\vartheta,i}(h)\odot(\Theta^\circ y)_i(h^\circ),\ \Xi(\operatorname{diag}(h,h^\circ)
\pog_i)\bigr\rangle ,
\]
and $(\Theta^\circ y)_i(h^\circ)$ is an integral over $g$ of multiples of $K^\circ(g,h^\circ)$; for
each $g$ the sum over $h^\circ$ vanishes by Proposition~\ref{po:periods}.
(G4) By \eqref{po:eq-deltasize} and the ultrametric inequality it suffices that
$|M|_\iota^{l_1}\,\plgb{x}{u''\odot u^\circ}\le1$ for $x=\operatorname{diag}(h,h^\circ)\pog_i$,
$u''=\varphi''_{\vartheta,i}(h)$, $u^\circ=\zeta_i(h^\circ)$. Now $\rho(x_\ell)^{-1}=\rho(\pog_{i,
\ell})^{-1}\rho(\operatorname{diag}(h,h^\circ)_\ell)^{-1}$, and by Lemma~\ref{po:prod}(M1) the second
factor maps $\iota(u''\odot u^\circ)$ to a product $a''\odot a^\circ$ in which $\|a''\|=\plgb
h{u''}\le1$ by (I1) and $\|a^\circ\|\le1$ because $\zeta$ is integral at $\iota$ (the block
$h^\circ$ acts on the factor $\poT^-$ through the contragredient, that is, through its components
at the conjugate places; this is why integrality on $V^\circ$ was defined through $\aubar\zeta$
and $\aubar\iota$); the norm of a tensor product of two vectors is the product of the norms.
If $\ell\nmid M$ then $\pog_{i,\ell}=1$ and we are done. Let $\ell\mid M$ and $\tau\in\auPl$. On the
factor $\Lambda^h$ of $\poT^+$ at $\tau$ the element $\pog_i^{-1}$ acts by
$\Lambda^h(\pog_{i,w_\tau}^{-1}|_{V''})$: if $w_\tau\in\mathsf M$ this has entries in
$\pi_M^{-3m}\auO$ by Lemma~\ref{po:lattice}(c) (after the unimodular change $K_1$ at most $3m$
columns are not integral), and otherwise $\pog_{i,w_\tau}^{-1}$ is integral, being $1$ or the
conjugate transpose of $\pog_{i,\aubar w_\tau}$. On a factor $\Lambda^h(\cdot)^\vee$ of $\poT^-$,
$h\le r_-\le m$, it acts by $\Lambda^h({}^t\pog_{i,w_\tau})$, which is integral if
$w_\tau\in\mathsf M$ and has entries in $\aubar{\pi_M}^{-h}\auO$ if $w_\tau\in\aubar{\mathsf M}$.
In all cases a column costs at most $|M|_\iota^{-3m}$, so $\|\rho(\pog_{i,\ell})^{-1}(a''\odot
a^\circ)\|\le|M|_\iota^{-l_1}$.
(G5) For $\vartheta=1$: by Lemma~\ref{po:classes}, $\poA(V;1)$ has at most $e^b$ classes, and
there are at most $\#I\cdot e^b\cdot e^b\le e^{3b}$ terms in $\poglue_1(\zeta)$; by
\eqref{po:eq-deltasize}, Lemma~\ref{po:prod}(M2) and (I2), each has values with
$\|\sigma(\cdot)\|\le n_{\mathrm{Aut}}\,e^{\poAA/2}\|\zeta\|_\sigma$. Hence
$\|\poglue_1\zeta\|_\sigma^2\le e^b\bigl(M^{l_1}e^{3b}n_{\mathrm{Aut}}\bigr)^2e^{\poAA}\|\zeta\|_\sigma^2$,
and $M^{2l_1}e^{7b}n_{\mathrm{Aut}}^2\le\exp(26\cdot22\,fmb\log b+7b+2x^{40})\le\exp(\poAA)$, a
comparison of sizes ($M\le b^{22}$).
\end{proof}

\subsection{The comparison operator; proof of \texorpdfstring{{\cite[Proposition~14.41]{P46}}}{[Row D, Proposition~14.41]}}\label{po:ss-proof}

Let $H_\circ=\prod_\tau H_{\mu^-_\tau}$ be the product of the hook numbers of the diagrams
$\mu^-_\tau$ ({\cite[Lemma~5.22]{P49}}), and
\[
U_\vartheta=n_{\mathrm{Aut}}\,H_\circ^2\ \poglue_\vartheta\circ\poglue_\vartheta^*\ \in\auEnd\bigl(
\poA(V;\vartheta)\bigr),\qquad\Delta_j(\vartheta)=\det\bigl(\popr_{\poZ_\vartheta}U_\vartheta
|_{\poZ_\vartheta}\bigr),\qquad\poZ_\vartheta=\Theta_\vartheta(e_0\poA(W_j)).
\]
We prove {\cite[Proposition~14.41]{P46}} with these $\Delta_j$, with $H_j$ given by \eqref{po:eq-H}
and with $Z_j(\vartheta)=\det(\Theta_\vartheta^*\Theta_\vartheta)$ on $e_0\poA(W_j)$ (the map $u$
of {\cite[Proposition~14.37]{P46}} is a positive multiple of $\Theta_\vartheta$, independent of
$\vartheta$; this changes only the constant $C_j$). If $d_j=0$ all determinants are $1$ and there
is nothing to prove; let $d_j>0$.

\begin{proof}[Proof of {\cite[Proposition~14.41]{P46}}, given the statements of Proposition~\ref{po:glue}]
\emph{Step 1: the identity.} By (G2), $\Theta_\vartheta$ is injective, so $\dim\poZ_\vartheta=d_j$
and Lemma~\ref{po:det} applies; by (G1),
\[
Z_j(\vartheta)\,\Delta_j(\vartheta)=\det\bigl(\Theta_\vartheta^*U_\vartheta\Theta_\vartheta\bigr)=
\det\bigl(n_{\mathrm{Aut}}H_\circ^2\,|c'_\vartheta|^{-2}\,\Theta^{\circ*}\Theta^\circ\bigr)=
\Bigl(\frac{n_{\mathrm{Aut}}H_\circ^2}{C_{\mathrm x}}\Bigr)^{d_j}\det(\Theta^{\circ*}\Theta^\circ)
\cdot H_j(\vartheta)^{d_j},
\]
by Proposition~\ref{po:exchange}. So the identity holds with $C_j=(n_{\mathrm{Aut}}H_\circ^2H_j(1)
/C_{\mathrm x})^{d_j}\det(\Theta^{\circ*}\Theta^\circ)$, which is positive because $\Theta^\circ$
is injective; and $\Delta_j(\vartheta)\ne0$.

\emph{Step 2: a sum of operators of rank one.} Let $\ell$ be a prime. By
{\cite[Corollary~6.4(a)]{P49}} choose representatives $x_c$ of the classes $c$ of
$\Gamma^\circ\backslash\auU(V^\circ)(\auA_f)/J^\circ$ that are integral above $\ell$, and let
$\eta_U\in W^\circ(\Z)$ be the frame of Lemma~\ref{po:frame} for $\mu^-$ and $N=n^\circ$. For
$a=(i,c,U)$ let $f_a\in\poA^\oplus$ be the element whose component $i$ is the form supported on
the class $c$ with value $\eta_U$ at $x_c$, the other components being $0$. Since $J^\circ$ is
free, $\poA^\circ$ is the orthogonal sum of copies of $W^\circ$ indexed by the classes; so
\[
\sum_a\langle f_a,\zeta\rangle f_a=H_\circ^2\,\zeta\quad(\zeta\in\poA^\oplus),\qquad
\|f_a\|_\sigma=\|\eta_U\|\le H_\circ\quad(\sigma\in\auAut(\C)),
\]
and $f_a$ is integral at every $\iota$ with prime $\ell$: $\aubar{f_a}$ has the value $\eta_U$, an
integral vector, at a representative that is integral above $\ell$, and (G1), (G2) of
{\cite[\S6.1]{P49}} apply as in {\cite[Corollary~6.4(b)]{P49}}. Consequently, for $\Xi\in\poA(V;
\vartheta)$,
\begin{equation}\label{po:eq-rankone}
U_\vartheta\Xi=n_{\mathrm{Aut}}\,\poglue_\vartheta\Bigl(\sum_a\langle f_a,\poglue_\vartheta^*\Xi
\rangle f_a\Bigr)=\sum_an_{\mathrm{Aut}}\,\langle\poglue_\vartheta f_a,\Xi\rangle\ \poglue_\vartheta
f_a .
\end{equation}

\emph{Step 3: $U_\vartheta$ keeps integrality.} Let $\ell=\ell(\iota)\nmid\auNm Q_\vartheta$ and
let $\Xi$ be integral at $\iota$. By (G4) the forms $\poglue_\vartheta f_a$ are integral at $\iota$
and at $\aubar\iota$; by Lemma~\ref{po:naut}, $n_{\mathrm{Aut}}\langle\poglue_\vartheta f_a,\Xi
\rangle\in\auO_\iota$. By \eqref{po:eq-rankone}, $U_\vartheta\Xi$ is integral at $\iota$.

\emph{Step 4: $U_\vartheta$ maps lifts to forms orthogonal to the false lifts.} For
$z=\Theta_\vartheta y\in\poZ_\vartheta$, (G1) gives $U_\vartheta z=n_{\mathrm{Aut}}H_\circ^2
c'^{-1}_\vartheta\,\poglue_\vartheta(\Theta^\circ y)$; apply (G3).

\emph{Step 5: proof of {\rm(b)}.} Let $\ell(\iota)\nmid\auNm Q_\vartheta$ and let $\Lambda$ be as
in Definition~\ref{po:denom}, for $\beta$. Let $\alpha$ be an eigenvalue of the endomorphism
$S=\Lambda\circ\popr_{\poZ_\vartheta}U_\vartheta|_{\poZ_\vartheta}$ of $\poZ_\vartheta$, with
eigenvector $z\ne0$. The gauges of the values of $z$ at representatives of the classes are
absolute values of elements of $\overline\Q_\ell$, not all $0$; multiplying $z$ by a suitable
constant we may assume that $z$ is integral at $\iota$ and that $az$ is integral at $\iota$ only
if $|a|_\iota\le1$. By Steps 3 and 4, $U_\vartheta z$ is integral at $\iota$ and orthogonal to the
false lifts; so $Sz=\alpha z$ is integral at $\iota$, and $|\alpha|_\iota\le1$. By
Lemma~\ref{po:eig}, $|\det\Lambda|_\iota\,|\Delta_j(\vartheta)|_\iota=|\det S|_\iota\le1$, and
$|\det\Lambda|_\iota\ge\exp(-d_jr\beta)$.

\emph{Step 6: proof of {\rm(a)}.} Let $\vartheta=1$ and $\sigma\in\auAut(\C)$. By
\eqref{po:eq-rankone}, \eqref{po:eq-cs}, (G5) and Step 2,
\[
\|U_1\Xi\|_\sigma\le n_{\mathrm{Aut}}\sum_a\|\poglue_1f_a\|_{c\sigma c}\,\|\Xi\|_\sigma\,
\|\poglue_1f_a\|_\sigma\le n_{\mathrm{Aut}}\cdot\#\{a\}\cdot e^{2\poAA}H_\circ^2\cdot\|\Xi\|_\sigma
\le e^{3\poAA}\|\Xi\|_\sigma ;
\]
indeed $\#\{a\}\le\#I\cdot e^b\cdot(n^\circ)^{|\mu^-|}$ by Lemmas~\ref{po:ff}, \ref{po:classes}
and~\ref{po:frame}, $H_\circ\le|\mu^-|^{3|\mu^-|}$, $|\mu^-|\le fm^2\Delta\le x^{60}$, and
$n_{\mathrm{Aut}}\le\exp(x^{40})$; so $n_{\mathrm{Aut}}\#\{a\}H_\circ^2\le e^{3b}\le e^{\poAA}$.
Let $\alpha$ be an eigenvalue of $\popr_{\poZ_1}U_1|_{\poZ_1}$ with eigenvector $z$. There is a
prime $w$ off the level, and $\poZ_1\subseteq\varrho_V(e_0)\poA(V;1)\cap\ker P_w(X_w)$ by
Lemma~\ref{po:facts}; by Step 4, $U_1z$ is orthogonal to $\mathcal F_w(e_0,\poZ_1)$. By
Proposition~\ref{po:proj},
\[
|\sigma\alpha|\,\|z\|_\sigma=\|\popr_{\poZ_1}U_1z\|_\sigma\le\|U_1z\|_\sigma\le e^{3\poAA}
\|z\|_\sigma ,
\]
and $\|z\|_\sigma>0$. By Lemma~\ref{po:eig}, $|\sigma\Delta_j(1)|\le\exp(3d_j\poAA)$.
\end{proof}

\subsection{Hecke values; proof of \texorpdfstring{{\cite[Proposition~14.42]{P46}}}{[Row D, Proposition~14.42]}}\label{po:ss-hecke}

\emph{The pairs.} We use the characters fixed with the admissible data in
Definition~\ref{pfq:def-chars}: $\chi_0=\chi_M$, of weight $0$, whose tame part $\omega=(\omega_w)$ is
trivial off $\mathsf M$ and has order $q_2>8m$ at every $w\in\mathsf M$; the tame part of $\mu_j^0$ is
$\omega^j$ (Proposition~\ref{pfq:tame}). For a finite set $S$ of primes of $E$ and a character $\nu$ of $\auU(1)$,
$L^S(s,\nu)$ is the $L$-function of $x\mapsto\nu(x/\aubar x)$ without the Euler factors at
$S\cup\aubar S$ (for $s=1$, the value of the continuation; the characters below are not trivial),
and $\auNm S=\prod_{w\in S}\auNm w$. The three pairs of the $j$-th lift are, with
$n'=j+(t-11m-1)/2$,
\[
(\nu_\ast,z_\ast)=\bigl(\mu_j^0\upsilon^{n'},\ \tfrac{7m}2+\tfrac k2\bigr),\qquad
(\nu_+,z_+)=\bigl(\upsilon^{-m}\chi_0^{\,j},\ m\bigr),\qquad
(\nu_-,z_-)=\bigl(\upsilon^{-m/2}\chi_0^{-j},\ \tfrac m2\bigr),
\]
with the signs $s_\ast=s_+=1$, $s_-=-1$ (here $\upsilon$ is restricted to $\auU(1)$; it has weight
$1$ and trivial tame part). They are critical pairs, of weights $-(\kappa_j(\tau)-n')$, $-m$,
$-m/2$ and tame parts $\omega^j$, $\omega^j$, $\omega^{-j}$. Put
\begin{equation}\label{po:eq-H}
H_j(\vartheta)=\prod_{i\in\{\ast,+,-\}}\Bigl((\auNm\mathsf S_\vartheta)^{z_i^2}\prod_{z=1}^{z_i}
\bigl|L^{\mathsf S_\vartheta}\bigl(z,(\nu_i\vartheta^{s_i})^{-1}\bigr)\bigr|^2\Bigr).
\end{equation}
For a critical pair $p=(\nu,z)$, a twist $\vartheta$ and a sign $s$ the \emph{ratio} is
$\rho_p(\vartheta,s)=|L^{\mathsf S_\vartheta}(z,(\nu\vartheta^{s})^{-1})|/|L^{\mathsf M}(z,
\nu^{-1})|$, and $\rho_p(\vartheta)=\rho_p(\vartheta,1)$; this is the quantity of (R2) (at the
places of $\mathsf Q$ the twisted character is ramified, so no Euler factor is removed there).
More precisely, let $q=(\nu,z)$ be one of the pairs $(\nu_i,z)$ and $1\le j\le m$. The characters
$\nu$ and $\nu\theta_\iota^{\pm j}$ are ramified at every place of $\mathsf S_{\theta_\iota}\cup
\aubar{\mathsf S_{\theta_\iota}}$: at $\mathsf M\cup\aubar{\mathsf M}$ because their tame part is
$\omega^{\pm j}\ne1$, and at $Q\cup\aubar Q$ because $\nu$ is unramified there while $\theta_\iota^{\pm j}$
has tame part $\psi_w^{\pm j}\ne1$, $\theta_\iota$ being an $m$-fold twist ({\cite[\S14.3]{P46}}); the
character $x\mapsto\nu(x/\aubar x)$ is ramified at $\aubar w$ exactly when it is ramified at $w$. Hence
$L^{\mathsf M}(z,\nu^{-1})=L(z,\nu^{-1})$ and $L^{\mathsf S_{\theta_\iota}}(z,(\nu\theta_\iota^{\pm j})^{-1})=
L(z,(\nu\theta_\iota^{\pm j})^{-1})$, and $\rho_q(\theta_\iota^{\,j},\pm1)$ is exactly the ratio bounded in (R2)
with $n=j$.

\begin{definition}[Main pairs]\label{po:main}
A critical pair $p=(\nu,z)$, with $\nu$ of weight $-c$, is a \emph{main pair for $j$} if
$m<z\le r$, the tame part of $\nu$ is $\omega^j$, and $\kappa_j(\tau)+r-t\le c(\tau)\le
\kappa_j(\tau)$ for all $\tau$. It is \emph{flat} if $c(\tau)=z$ for all $\tau$.
\end{definition}

\begin{lemma}\label{po:Hratio}
With $e=z_\ast^2+z_+^2+z_-^2$,
\[
\frac{H_j(\vartheta)}{H_j(1)}=(\auNm Q_\vartheta)^{e}\prod_{i\in\{\ast,+,-\}}\ \prod_{z=1}^{z_i}
\rho_{(\nu_i,z)}(\vartheta,s_i)^2 .
\]
All $5m+k/2$ pairs $(\nu_i,z)$ are ramified critical pairs, and the $k/2$ pairs $(\nu_\ast,z)$,
$7m/2<z\le z_\ast$, are main pairs for $j$.
\end{lemma}
\begin{proof}
The identity is the definition, with $\auNm\mathsf S_\vartheta=\auNm\mathsf M\cdot\auNm
Q_\vartheta$. The numbers of pairs are $z_\ast+z_++z_-=5m+k/2$ ($k=r+1-m$ is even). Each
$(\nu_i,z)$ is a critical pair: $1\le z\le z_i\le t$; for $i=\pm$, $c=z_i\ge z$; for $i=\ast$,
$c(\tau)=\kappa_j(\tau)-n'\ge z_\ast$ because $\kappa_j(\tau)=\kappa(\tau)+(j+\frac12)\Delta\ge\Delta\ge3t$ ({\cite[\S14.1]{P46}}), and $c(\tau)\le
\kappa_j(\tau)\le4b$. They are ramified because $\omega_w^{\pm j}\ne1$ ($1\le j\le m$, the order
of $\omega_w$ exceeds $8m$). For $7m/2<z\le z_\ast$: $m<z\le z_\ast\le r$ (indeed $z_\ast=(r+1+6m)/2\le r$, since $r=k+m-1$ and
$k\ge6m+2$ by (S7) of {\cite[Lemma~13.4]{P46}}), the tame part is
$\omega^j$, and $0\le n'\le t-r$.
\end{proof}

\begin{lemma}[Positivity of $H_j$]\label{po:Hpos}
Let $\nu\ne1$ be an automorphic character of $\auU(1)$, $S$ a finite set
of primes of $E$, and $z\ge1$ an integer. Then $L^S(z,\nu^{-1})\ne0$. In particular, for every
twist $\vartheta$ and every $i\in\{\ast,+,-\}$ the values $L^{\mathsf S_\vartheta}(z,(\nu_i
\vartheta^{s_i})^{-1})$, $1\le z\le z_i$, are non-zero, and $H_j(\vartheta)>0$.
\end{lemma}
\begin{proof}
Since $\auU(1)$ is anisotropic over $E^+$, $\auU(1)(E^+)\backslash\auU(1)(\auA_{E^+})$ is compact,
so $\nu$ is unitary, and so is the Hecke character $\chi(x)=\nu^{-1}(x/\aubar x)$ of
$\auA_E^\times/E^\times$. The map $x\mapsto x/\aubar x$ from $E_v^\times=(E\otimes E^+_v)^\times$ to
$\auU(1)(E^+_v)$ is surjective at every place $v$ of $E^+$ (Hilbert's Theorem~90), and maps
$(\auO_E\otimes\auO_{E^+,v})^\times$ onto $\auU(1)(\auO_{E^+,v})$ at every finite $v$ unramified in
$E$ (Hilbert~90 for the unramified extension of local rings); hence $x\mapsto x/\aubar x$ maps
$\auA_E^\times$ onto $\auU(1)(\auA_{E^+})$, and $\chi\ne1$ because $\nu\ne1$. Moreover $\chi$ is not of the
form $|\cdot|_E^{it}$ with $t\in\R$: $\chi$ is trivial on $\auA_{E^+}^\times$, where $x/\aubar x=1$, while
$|\cdot|_E^{it}$ restricts to $|\cdot|_{E^+}^{2it}$ there, which is trivial only for $t=0$; and $\chi\ne1$.

For real $s>1$ the Euler product $L^S(s,\chi)=\prod_{w\notin S,\ \chi\text{ unr.\ at }w}(1-\chi
(\varpi_w)q_w^{-s})^{-1}$ converges absolutely, since $|\chi(\varpi_w)q_w^{-s}|=q_w^{-s}$ and
$\sum_wq_w^{-s}\le2f\zeta(s)<\infty$; an absolutely convergent product of non-zero factors is
non-zero, so $L^S(z,\chi)\ne0$ for $z\ge2$ (Lemma~\ref{psu:euler} gives the quantitative form).
The point $z=1$ is on the edge of the region of absolute convergence. For a unitary Hecke character
$\lambda$ of $E$ write $L(s,\lambda)=\prod_{w\ \text{finite},\ \lambda\text{ unr.\ at }w}(1-\lambda
(\varpi_w)q_w^{-s})^{-1}$ (no Euler factor at the ramified primes, no archimedean factors), and, for a
finite set $R$ of primes, $L_R(s,\lambda)$ for the same product over $w\notin R$. The only analytic input
is Hecke's continuation theorem \cite{Hecke1920}: $L(s,1)=\zeta_E(s)$ extends meromorphically
to $\C$ with a single, simple pole at $s=1$ \cite[Chapter~VII, Corollary~(5.11)]{Neukirch}, and $L(s,\lambda)$
extends to an entire function for every unitary Hecke character $\lambda\ne1$ trivial on
$\auA_{E^+}^\times$ (the only ones used below). Indeed, regard $\lambda$ as a Gr\"o{\ss}encharakter as in
\cite[Chapter~VII, \S6]{Neukirch}; its component at infinity is trivial on the positive reals of
$E^+\otimes\R\subset E\otimes\R$. If $\lambda$ is ramified or its component at infinity is non-trivial, the
completed $L$-function of the associated primitive character is entire by \cite[Chapter~VII,
Theorem~(8.5)]{Neukirch}. Otherwise $\lambda$ is a non-trivial character of the ideal class group, and its
completed $L$-function is entire by \cite[Chapter~VII, the remark after Corollary~(5.10)]{Neukirch}. In both
cases dividing by the archimedean factor, which has no zeros, gives the claim (the primitive and the
original character differ at finitely many Euler factors, which are handled as below). Since the finitely many removed factors
$(1-\lambda(\varpi_w)q_w^{-s})$ are entire and non-zero on $\mathrm{Re}\,s>0$, the same holds for
$L_R$. From this we deduce $L(1,\chi)\ne0$ (this is the case $s=1$ of Hecke's non-vanishing theorem on
$\mathrm{Re}\,s=1$; for characters of finite order it is \cite[Chapter~VII, Lemma~(13.3)]{Neukirch}; see also \cite[Theorem~2.4]{FiteSatoTate}). Let $R$ be the set of
primes at which $\chi$ ramifies. First, $\chi^2$ is not of the form $|\cdot|_E^{it}$ unless $\chi^2=1$, by
the restriction argument above ($\chi^2$ is trivial on $\auA_{E^+}^\times$).

\emph{Case $\chi^2\ne1$.} For real $s>1$ put $Z(s)=L_R(s,1)^3L_R(s,\chi)^4L_R(s,\chi^2)$. Taking
logarithms of the absolutely convergent Euler products, with $\chi(\varpi_w)=e^{i\theta_w}$,
\[
\log|Z(s)|=\sum_{w\notin R}\sum_{k\ge1}\frac{q_w^{-ks}}{k}\bigl(3+4\cos k\theta_w+\cos2k\theta_w\bigr)\ge0,
\]
because $3+4\cos\phi+\cos2\phi=2(1+\cos\phi)^2$. So $|Z(s)|\ge1$ for $s>1$. Near $s=1$, $L_R(s,1)^3$
has a pole of order $3$, $L_R(s,\chi^2)$ is holomorphic, and if $L(1,\chi)=L_R(1,\chi)=0$ then
$L_R(s,\chi)^4$ has a zero of order $\ge4$; then $Z(s)\to0$ as $s\to1^+$, a contradiction.

\emph{Case $\chi^2=1$.} Then $\chi(\varpi_w)\in\{\pm1\}$ for $w\notin R$, and the Euler factor of
$D(s)=L_R(s,1)L_R(s,\chi)$ at $w$ is $(1-x)^{-2}=\sum_k(k+1)x^k$ or $(1-x^2)^{-1}=\sum_kx^{2k}$, with
$x=q_w^{-s}$; both have non-negative coefficients and coefficient $\ge1$ at every even power. Hence
$D(s)=\sum_{n\ge1}a_nn^{-s}$ ($s>1$) with $a_n\ge0$, and, with $d=[E:\Q]$, $a_{m^{2d}}\ge1$ for every
integer $m\ge1$ prime to the primes below $R$ (the ideal $(m)^2$ has norm $m^{2d}$ and the coefficient of
$\mathrm N\mathfrak a^{-s}$ is multiplicative in $\mathfrak a$). Suppose $L(1,\chi)=0$. Then the zero
cancels the simple pole of $L_R(s,1)$ and $D$ is holomorphic on $\mathrm{Re}\,s>0$. Put
$A(x)=\sum_{n\le x}a_n\ge0$; for $s>1$, $D(s)/s=\int_1^\infty A(x)x^{-s-1}dx$. By Landau's theorem
\cite[\S15.1, Lemma~15.1]{MV} (applied to $\int_1^\infty A(x)x^{-u}dx$, $u=s+1$) this integral is not
analytic at its abscissa of convergence, while $D(s)/s$ continues analytically to $\mathrm{Re}\,s>0$;
so the integral converges for every real $s>0$. Since $\int_1^\infty A(x)x^{-s-1}dx\ge\sum_na_n\int_n^\infty
x^{-s-1}dx=s^{-1}\sum_na_nn^{-s}$, the series $\sum_na_nn^{-s}$ converges for $s>0$; but at $s=1/(2d)$
it is $\ge\sum_{m}m^{-1}$ over $m$ prime to a finite set of primes, which diverges. So $L(1,\chi)\ne0$. Finally
$L^S(s,\chi)=L(s,\chi)\prod_{w\in S\cup\aubar S,\ \chi\text{ unr.\ at }w}(1-\chi(\varpi_w)q_w^{-s})$,
and at $s=1$ each removed factor satisfies $|1-\chi(\varpi_w)q_w^{-1}|\ge1-q_w^{-1}>0$; so removing
the Euler factors at $S$ creates no zero, and $L^S(1,\chi)\ne0$.

For the pairs of \eqref{po:eq-H}: $\nu_i\vartheta^{s_i}\ne1$, for two independent reasons. Its weight
is that of $\nu_i$, namely $-c$ with $c(\tau)\ge z_i\ge1$ (the twist has weight $0$; proof of
Lemma~\ref{po:Hratio}), so it is non-trivial on $\auU(1)(E^+\otimes\R)$; and its tame part at
$w\in\mathsf M$ is $\omega_w^{\pm j}\ne1$ (the twist has trivial tame part at $\mathsf M$,
\S\ref{po:ss-set}). The first part applies with $S=\mathsf S_\vartheta$. Since also $\auNm\mathsf S_\vartheta\ge1$, every
factor of \eqref{po:eq-H} is positive.
\end{proof}

\emph{The lift of a character.} Let $p=(\nu,z)$ be a critical pair with $z\ge2$ and $\vartheta$ a
twist. Apply \S\ref{po:ss-set} to the positive definite line $E$ in place of $W$ (so $r=1$, there
is no $\mu^-$, and the unitary group is $\auU(1)$), to the standard space $V^{(z)}$ of rank $2z$,
to the character $\nu\vartheta$, and to the weight with one row of length $l_\tau=c(\tau)-z$ at
$\tau$: this gives the block datum $\Phi^{(1)}_\vartheta$, the level, the space
$\poA(V^{(z)};p,\vartheta)$ with the action of $\auG_1=\prod_{w\in\mathsf M}\auF_w^\times$, the
numbers $X_w$, the lift $\Theta^{(1)}_\vartheta$ of functions on the finite set
$\auU(1)(E^+)\backslash\auU(1)(\auA_{E^+,f})/K^1_M$ ($K^1_M$ the units $\equiv1$ at $\mathsf M$)
and its adjoint, the descent $D^{(1)}_\vartheta$. The \emph{normalised lift} is
\[
\varphi_{p,\vartheta}=c_{\mathrm{vac}}^{\,2z}\,(\pi_Ms_0)^{\sum_\tau l_\tau}\cdot[K^1:K^1_M]\cdot
\Theta^{(1)}_\vartheta(\plone)\ \in\poA(V^{(z)};p,\vartheta),
\]
with the constants $c_{\mathrm{vac}}$, $s_0$ of {\cite[\S5.1.1]{P49}}: its values are sums over the
classes of $\auU(1)$ of vectors of jets at scale $\pi_Ms_0$ of theta functions of translates of
the integral vacuum, as in {\cite[(34)]{P49}}. Let $e_1\in\C[\auG_1]$ be the mean over $\auG_1$.

\begin{definition}[Test forms; the clause of {\rm(R3)} on main pairs]\label{po:R3}
A form $\Xi\in\poA(V^{(z)};p,\vartheta)$ is a \emph{test form at $\ell$} if it is integral at
every homomorphism $\C\to\C_\ell$ and, for every $w$ off the level, orthogonal to
\[
\varrho(e_1)\poA(V^{(z)};p,\vartheta)\ \cap\ \ker P_w(X_w)\ \cap\ \ker D^{(1)}_\vartheta
\]
(the false lifts for $(e_1,(\ker D^{(1)}_\vartheta)^\perp)$ of Definition~\ref{po:false}). The
second clause of {\rm(R3)} in {\cite[Definition~14.39]{P46}} is: for every $\iota$ with prime
$\ell$, every main pair $p$ for $j$, and all test forms $\Xi,\Xi'$ at $\ell$ for
$(p,\theta_\iota^{\,j})$,
\[
\bigl|\langle\popr_\varphi\Xi,\popr_\varphi\Xi'\rangle\bigr|_\iota\le e^{\beta_\ell},\qquad
\popr_\varphi\text{ the orthogonal projection onto }\C\,\varphi_{p,\theta_\iota^{\,j}} .
\]
\end{definition}

\begin{proposition}[The norm of the lift is an $L$-value]\label{po:rallis}
Let $p=(\nu,z)$ be a critical pair with $z\ge2$. There is $A_p>0$ such that for every twist
$\vartheta$
\[
\|\varphi_{p,\vartheta}\|^2=A_p\,(\auNm Q_\vartheta)^{2z-1}\,L^{\mathsf S_\vartheta}\bigl(z,(\nu
\vartheta)^{-1}\bigr);\qquad\text{hence}\quad\rho_p(\vartheta)=\frac{\|\varphi_{p,\vartheta}
\|^2}{(\auNm Q_\vartheta)^{2z-1}\,\|\varphi_{p,1}\|^2} .
\]
\end{proposition}
\noindent(The $L$-value is positive: the character $x\mapsto(\nu\vartheta)^{-1}(x/\aubar x)$ of
$E$ is carried to its complex conjugate by $c$, the Euler factors at $w\ne\aubar w$ come in
conjugate pairs, and at $w=\aubar w$, which is unramified over $E^+$, the factor is
$(1-\auNm w^{-z})^{-1}$.)
\noindent Proved in Proposition~\ref{pex:rallis}, for all tame characters of modulus $Q$.


\begin{proposition}[The untwisted norm at all conjugates]\label{po:arch}
For a main pair $p$ for $j$ the constant family $\|\varphi_{p,1}\|^2$ has height at most
$\aubud-\sum_{\tau\in\auPl}\ausinf(\kappa(\tau))$.
\end{proposition}
\noindent Proved in Theorem~\ref{psu:arch}, with the law of Galois conjugation (CL) of
\S\ref{psu:ss-arch} proved in {\cite[Theorem~9.1]{P53}}. The multiplier is $I_q(1)/\|\varphi_{p,1}\|^2$ with
$q$ flat of the same rank; its integrality, the archimedean factor $\prod_\tau l_\tau!$ from
Proposition~\ref{pex:rallis} and the Euler product bounds (Lemma~\ref{psu:euler}) are proved there.


\begin{proposition}[A test form with a pairing of least valuation]\label{po:small}
Let $p$ be a main pair for $j$ and $q_2=(\upsilon^{-2}\chi_0,2)$, a flat ramified critical pair.
There is $c_0\in\C^\times$ with $|\sigma c_0|\le\exp(fmb/k+3x^{35})$ for all $\sigma\in\auAut(\C)$
such that: for every prime $\ell$, every homomorphism $\jmath\colon\C\to\C_\ell$ and every twist
$\vartheta$ with $\ell\nmid\auNm Q_\vartheta$ there is a test form $\Xi$ at $\ell$ for
$(p,\vartheta)$ with
\[
\bigl|c_0\,\rho_{q_2}(\vartheta)\bigr|_\jmath\cdot\ell^{-F(\ell)}\le\bigl|\langle\Xi,\varphi_{p,
\vartheta}\rangle\bigr|_\jmath,\qquad F(\ell)=\sum_{\sigma\in\auPl}s_\ell(\kappa(\sigma),\sigma).
\]
\end{proposition}
\noindent\emph{Proof.} Theorem~\ref{psu:small} (with Theorem~\ref{psm:attained} and
Corollary~\ref{pst:attained} for the sharpness of the count). The test form is the triple class sum
\eqref{psu:eq-Xi} of point forms against the normalised lifts of three flat pairs; its pairing with
$\varphi_{p,\vartheta}$ is computed in Proposition~\ref{psu:pairing}, its integrality is
Lemma~\ref{psu:integral}, and its orthogonality to the false lifts (Lemma~\ref{psu:orthfalse}) is the
rank-one case of Theorem~\ref{ppe:main} (Proposition~\ref{psu:periods}, Remark~\ref{ppe:rem-rankone}).
The constant is $3x^{35}$ because one integral value $I_q(1)$ is bounded by $\exp(x^{35})$
(Lemma~\ref{psu:I}(c)).


\begin{proposition}[The height of a main pair]\label{po:mainheight}
Let $\theta$ be a regular
family ({\cite[Definition~14.39]{P46}}, with Definitions~\ref{po:denom} and~\ref{po:R3}) and $p$
a main pair for $j$. Then $\iota\mapsto\rho_p(\theta_\iota^{\,j})^{-1}$ has height at most
$6\,\aubud-\auhdir$.
\end{proposition}
\begin{proof}
Fix $\iota$, with prime $\ell$, and put $\vartheta=\theta_\iota^{\,j}$, $\varphi=\varphi_{p,
\vartheta}$, $\rho_2=\rho_{q_2}(\vartheta)>0$; by (R1), $\ell\nmid\auNm Q_\vartheta$. Let $\Xi$
and $\Xi'$ be the test forms of Proposition~\ref{po:small} for the homomorphisms $\iota$ and
$\aubar\iota$. Since $|\langle\Xi',\varphi\rangle|_{\aubar\iota}=|\langle\varphi,\Xi'\rangle|_\iota$
and $|c_0\rho_2|_{\aubar\iota}=|\aubar{c_0}\rho_2|_\iota$, multiplying the two inequalities gives
\begin{equation}\label{po:eq-two}
\bigl|c_0\aubar{c_0}\,\rho_2^2\bigr|_\iota\,\ell^{-2F(\ell)}\le\bigl|\langle\Xi,\varphi\rangle
\langle\varphi,\Xi'\rangle\bigr|_\iota ;
\end{equation}
in particular the right side is not $0$. Now $\popr_\varphi\Xi=\|\varphi\|^{-2}\langle\varphi,\Xi
\rangle\varphi$, so $\langle\popr_\varphi\Xi,\popr_\varphi\Xi'\rangle=\langle\Xi,\varphi\rangle
\langle\varphi,\Xi'\rangle/\|\varphi\|^2$, and by Proposition~\ref{po:rallis}
\[
\rho_p(\vartheta)^{-1}=(\auNm Q_\vartheta)^{2z-1}\cdot\|\varphi_{p,1}\|^2\cdot\langle\popr_\varphi
\Xi,\popr_\varphi\Xi'\rangle\cdot\rho_2^{-2}\cdot\frac{\rho_2^2}{\langle\Xi,\varphi\rangle\langle
\varphi,\Xi'\rangle} .
\]
We bound the heights of the five families of $\iota$. The first is a positive integer prime to
$\ell(\iota)$ and costs nothing. The second has height at most $\aubud-\sum_\tau\ausinf(\kappa
(\tau))$ (Proposition~\ref{po:arch}). The third has height at most $\aubud$: by the second clause
of (R3) its absolute value at $\iota$ is at most $e^{\beta_\ell}=\ell^{\beta_\ell/\log\ell}$, and
$\sum_{\ell\in T}\beta_\ell\le\aubud$. The fourth has height at most $2x_1$ by (R2), $q_2$ being
ramified. For the fifth take the multiplier $c_0\aubar{c_0}$: its conjugates are at most
$\exp(2fmb/k+6x^{35})$, and by \eqref{po:eq-two} the product has absolute value at most
$\ell^{2F(\ell)}$ at $\iota$; by {\cite[Lemma~14.9(a)]{P46}}, $\sum_{\ell\in T}2F(\ell)\log\ell\le
\sum_\sigma\ausfin(\kappa(\sigma),\sigma)$. Adding, and recalling the definition of $\auhdir$, the
height is at most
\[
2\,\aubud+2x_1+2\bigl(fmb/k+3x^{35}\bigr)-\auhdir\le6\,\aubud-\auhdir ,
\]
because $\aubud=fxb(\log b)^3/(mk)$ exceeds $x_1$ and $fmb/k+3x^{35}$ (Lemma~\ref{psu:sizes}(N6)).
\end{proof}

\begin{proof}[Proof of {\cite[Proposition~14.42]{P46}}, given Propositions~\ref{po:rallis}, \ref{po:arch},
\ref{po:small}]
By Lemma~\ref{po:Hratio}, $H_j(1)/H_j(\theta_\iota^{\,j})$ is the product of $(\auNm
Q_{\theta_\iota})^{-e}$, a rational power of a positive integer prime to $\ell(\iota)$ by (R1),
of the squares of $k/2$ inverse ratios of main pairs for $j$, and of the squares of $5m$ inverse
ratios $\rho_q(\theta_\iota^{\,j},\pm1)^{-1}$ of ramified critical pairs. By
Proposition~\ref{po:mainheight} each of the former has height at most $6\aubud-\auhdir$, and by
(R2) each of the latter has height at most $x_1$. Heights are subadditive under products, so the
height is at most $2\cdot\tfrac k2(6\aubud-\auhdir)+2\cdot5m\cdot x_1$.
\end{proof}

\subsection{What the two propositions rest on}\label{po:ss-rest}

{\cite[Proposition~14.41]{P46}} is proved here from: (L4) (Proposition~\ref{plf:L4}); Lemma~\ref{po:ff}
(Theorem~\ref{pfq:ff}, for the first factors of Definition~\ref{pfq:def-ff}); Lemma~\ref{po:classes} (Lemma~\ref{plf:classes}); Propositions~\ref{po:exchange}
(\S\ref{pte:sec}, \S\ref{ptf:sec}), \ref{po:span} (Proposition~\ref{psp:main}, Theorem~\ref{psv:main}) and
\ref{po:periods} (Theorem~\ref{ppe:main}); Lemma~\ref{po:Hpos}; the fact (W7) with the normalisation
\eqref{po:eq-w5}; and the properties of the weights $\mu^\pm$ (Lemma~\ref{plf:match}) and of the characters of
\S\ref{po:ss-set} and \S\ref{po:ss-hecke} (Definition~\ref{pfq:def-chars}, Proposition~\ref{pfq:tame}). {\cite[Proposition~14.42]{P46}} is proved from Propositions~\ref{po:arch}
(Theorem~\ref{psu:arch}) and~\ref{po:small} (Theorem~\ref{psu:small}), Proposition~\ref{po:rallis} being
Proposition~\ref{pex:rallis}; the last contains the sharpness of the count and the rank-one case of
Proposition~\ref{po:periods}. Everything else in this section is proved in full:
Lemmas~\ref{po:naut}, \ref{po:facts}, \ref{po:det}, \ref{po:eig}, \ref{po:adj}, \ref{po:frame},
\ref{po:prod}, \ref{po:lattice}, \ref{po:slice}, \ref{po:Hratio}, \ref{po:Hpos},
Propositions~\ref{po:proj}, \ref{po:glue} and~\ref{po:mainheight}, and the two deductions.
The operator is bounded directly from \eqref{po:eq-rankone}, and $n_{\mathrm{Aut}}$ is the integer of
Lemma~\ref{po:naut}.



\section{The norm of the lift of a character and the three exchanges}\label{C2:g2}
\providecommand{\auO}{\mathcal O}
\providecommand{\auU}{\mathrm U}
\providecommand{\auGL}{\mathrm{GL}}
\providecommand{\auHom}{\operatorname{Hom}}
\providecommand{\auNm}{\operatorname{Nm}}
\providecommand{\auA}{\mathbb A}
\providecommand{\auF}{\mathbb F}
\providecommand{\auPl}{\Phi}
\providecommand{\aubar}[1]{\overline{#1}}
\providecommand{\plone}{\mathbf 1}
\providecommand{\auf}[1]{\texttt{#1}}
\providecommand{\pop}{\mathsf p}
\providecommand{\pcR}{R}
\providecommand{\pcW}{\mathsf W}
\providecommand{\psHS}{\mathcal S}
\providecommand{\pao}{\mathfrak o}
\providecommand{\psvac}{{\phi^{\mathrm{vac}}}}
\providecommand{\pexK}{K^{(z)}}
\providecommand{\pexmc}{\mathfrak c}
\providecommand{\pexMat}{\mathrm M}

\subsection{The norm of the lift of a character, and integral measures}\label{pex:sec}

This section proves Proposition~\ref{po:rallis}, in the slightly stronger form needed below
(Proposition~\ref{pex:rallis}: all tame characters of modulus $Q$, not only twists, with the
constant made explicit), and deduces from it statement (IM) of {\cite[\S3.5]{P55}} for the
critical pairs with $z\ge2$ (Proposition~\ref{pex:IM}) and, with {\cite[Theorem~3.2]{P53}}, for $z=1$
(Proposition~\ref{pex:IM1}). Proposition~\ref{po:exchange} is proved in \S\ref{pte:sec} and \S\ref{ptf:sec}, with the explicit form of
{\cite[Theorem~3.2]{P53}} given in Corollary~\ref{pte:explicit}.

\subsubsection{Setting}\label{pex:ss-set}

We use the notation of \S\ref{po:ss-hecke} and {\cite[\S3.5]{P55}}. A \emph{tame character of
modulus $Q$} is as in {\cite[\S3.5]{P55}}: a character $\vartheta$ of $\auU(1)(E^+)\backslash
\auU(1)(\auA_{E^+})$, trivial at the infinite places, unramified outside $Q$ and tamely ramified
(possibly unramified) at the places of $Q$, where $Q$ is a finite set of places of $E$, prime to
$MN$, with $Q\cap\aubar Q=\emptyset$ and $Q\cap\mathsf M=\emptyset$; the twists of
\S\ref{po:ss-set} are the tame characters ramified at every place of their modulus. The trivial
character has modulus $\emptyset$. For $w\in Q$ let $\psi_w$ be the tame part of $\vartheta$ at
$w$, a character of $\auF_w^\times$ (trivial if $\vartheta$ is unramified at $w$), extended by
$\psi_w(0)=0$. Put $\mathsf S=\mathsf M\cup Q$, $\auNm Q=\prod_{w\in Q}q_w$, $q_w=\#\auF_w$.
Every place of $\mathsf S$ lies above a place of $E^+$ that splits in $E$, and $E/E^+$ is
unramified at all finite places ({\cite[\S2.3]{P53}}).

Fix a critical pair $p=(\nu,z)$ with $z\ge2$; $\nu$ has weight $-c$, $z\le c(\tau)$, and its tame
part is supported on $\mathsf M$. Put $l_\tau=c(\tau)-z\ge0$ and $\mu=\nu\vartheta$. The objects
of \S\ref{po:ss-hecke}, \emph{the lift of a character}, are defined verbatim for a tame character
$\vartheta$ of modulus $Q$ (the definition in \S\ref{po:ss-hecke} uses of the twist only $Q$ and the
tame parts $\psi_w$): the definite line $E$ in place of $W$, the standard space $V^{(z)}=E^{2z}$, the
block datum
\[
\Phi^{(1)}_{\vartheta}=\bigotimes_{w\in\mathsf M}\plone[x\equiv e_1\bmod w]\otimes
\bigotimes_{w\in Q}\plone[x\text{ integral}]\,\psi_w(x_1\bmod w)\otimes(\text{vacuum elsewhere})
\]
on $\auHom(E,V^{(z)})(\auA_f)=E^{2z}(\auA_f)$, columns written in the frame and read in the
coordinates $T=a_w^{-1}x_w$ of (N5) ({\cite[Lemma~6.32]{P53}}; so ``integral'' means $x_w\in a_w\auO_w^{2z}=\alpha_w\auO_w^{2z}$), $x_1$ the first
coordinate; the level $\pexK_\vartheta\subseteq K_{V^{(z)}}$ of the $k$ whose reduction fixes $e_1$
at the places of $\mathsf M$ and preserves the span of $e_2,\dots,e_{2z}$ at the places of $Q$; the
space $\mathcal A(V^{(z)};p,\vartheta)$ of algebraic forms of the weight with one row of length
$l_\tau$ at $\tau$, level $\pexK_\vartheta$ and the character by which $\pexK_\vartheta$ acts on
$\Phi^{(1)}_\vartheta$; the lift $\Theta^{(1)}_\vartheta$ of \eqref{po:eq-lift}, from functions on
$\auU(1)(E^+)\backslash\auU(1)(\auA_{E^+})/K^1_M$ (with $K^1_M$ the units $\equiv1$ at
$\mathsf M$, and all units at the other finite places); and the normalised lift
\begin{equation}\label{pex:eq-phi}
\varphi_{p,\vartheta}=\kappa_p\,[K^1:K^1_M]\,\Theta^{(1)}_\vartheta(\plone),\qquad\kappa_p=
c_{\mathrm{vac}}^{\,2z}\,(\pi_Ms_0)^{\sum_\tau l_\tau},
\end{equation}
where $K^1$ is the group of all integral points of $\auU(1)(\auA_{E^+,f})$ and $\kappa_p\ne0$ does not
depend on $\vartheta$. For a twist this is the $\varphi_{p,\vartheta}$ of \S\ref{po:ss-hecke}. The
norm $\|\cdot\|$ on $\mathcal A(V^{(z)};p,\vartheta)$ is that of {\cite[\S14.3]{P46}}.

\emph{The two conventions.} By \eqref{po:eq-w5} and {\cite[Remark~2.17]{P53}}, the representation
used to define $\Theta^{(1)}_\vartheta$ is $\omega'_{a,b}$ with $a=\upsilon^{r_-}=1$ (the line is
positive definite) and $b=\mu\,\upsilon^{z}$; by {\cite[(13)]{P53}} it is $\omega_{\alpha,\beta}$ of
{\cite[\S2.3]{P53}} with
\[
\alpha=\nu\vartheta\,\upsilon^{z}\ \text{ on }\auU(1),\qquad\beta=1\ \text{ on }\auU(V^{(z)}).
\]
Let $v^+$ be the highest weight vector of the weight of $\mathcal A(V^{(z)};p,\vartheta)$ and
$P=\pop_{v^+}$ the polynomial \eqref{po:eq-poly} (a polynomial of degree $l_\tau$ at $\tau$ in the
$2z$ holomorphic variables of $\auHom(E,V^{(z)})\otimes_\tau\R=\C^{2z}$). By \eqref{po:eq-lift}, for
a function $y$ on $[\auU(1)]$ and $h\in\auU(V^{(z)})(\auA)$,
\begin{equation}\label{pex:eq-f}
f_{\Theta^{(1)}_\vartheta y}(h)=\bigl\langle v^+,h_\infty^{-1}(\Theta^{(1)}_\vartheta y)(h_f)\bigr\rangle
=\int_{[\auU(1)]}y(g)\,\theta\bigl(\omega_{\alpha,\beta}(g,h)[P,\Psi_\vartheta]\bigr)\,dg
=\bigl(\Theta_{\alpha,\beta}(P,\Psi_\vartheta)y\bigr)(h),
\end{equation}
in the notation of {\cite[Lemma~2.18]{P53}} and {\cite[\S2.3]{P53}}, where $\Psi_\vartheta$ is the
function on $\pexMat_{2z\times1}(\pcR_{\mathfrak m})$, $\mathfrak m=(\pi_M)\prod_{w\in Q}w$, of which
$\Phi^{(1)}_\vartheta$ is the datum $\Psi^\natural_\vartheta\otimes\psvac^{\mathfrak m}$ of
{\cite[Definition~6.31]{P53}} (both are supported on $a_w\auO_w^{2z}$ at $w\mid\mathfrak m$ and are the same
function of $a_w^{-1}x$, the coordinate $T$ of (N5), by {\cite[Lemma~6.32(c)]{P53}}): the first
equality is the definition of $f$, the second is \eqref{po:eq-lift} with $\xi=v^+$ and the
$\auU(V^{(z)})(\R)$-equivariance of $\xi\mapsto\pop_\xi$.

\subsubsection{Local computations}\label{pex:ss-local}

\begin{lemma}[Sharpness and the finite masses]\label{pex:sharp}
Let $J\subseteq K^1$ be the group of the integral points $\equiv1$ at the places of $\mathsf M$
(no condition at the places of $Q$).
\begin{enumerate}
\item[(a)] $\Psi_\vartheta$ is sharp for $J$ ({\cite[Definition~2.15]{P53}}, for the modulus
$\mathfrak m$ and the character $\alpha$).
\item[(b)] $\|\Psi_\vartheta\|^2/\#\pexMat_{2z\times1}(\pcR_{\mathfrak m})=\prod_{w\in\mathsf M}q_w^{-2z}
\cdot\prod_{w\in Q}(q_w-1)/q_w$, and $[K^1:J]=\prod_{w\in\mathsf M}(q_w-1)$.
\item[(c)] $[K_{V^{(z)}}:\pexK_\vartheta]=c_{\mathsf M}\prod_{w\in Q}\dfrac{q_w^{2z}-1}{q_w-1}$, with
$c_{\mathsf M}=\prod_{w\in\mathsf M}(q_w^{2z}-1)$ independent of $\vartheta$.
\end{enumerate}
\end{lemma}
\begin{proof}
(a) $J$ contains $J_{E}(\mathfrak m)$ and the full components of $K^1$ outside $\mathsf S$. At
$w\in\mathsf M$ the reduction of $J$ is trivial and there is nothing to check for the first condition
of sharpness. At $w\in Q$ the tame part of $\alpha$ is that of $\vartheta$, namely $\psi_w$ ($\nu$
and $\upsilon$ are unramified at $w$), and for $\bar k\in\auF_w^\times$,
$\psi_w(\bar k)^{-1}\Psi_{\vartheta,w}(x\bar k)=\psi_w(\bar k)^{-1}\psi_w(x_1\bar k)=\Psi_{\vartheta,w}
(x)$. For the second condition take $A(x)=$ the first coordinate of $x$ at every place of
$\mathfrak m$: on the support of $\Psi_\vartheta$ it is $1$ at $w\in\mathsf M$ and a unit at $w\in Q$,
hence lies in the reduction $\bar J=\prod_{w\in\mathsf M}\{1\}\times\prod_{w\in Q}\auF_w^\times$.
(b) At $w\in\mathsf M$ the support is one vector out of $q_w^{2z}$; at $w\in Q$ it consists of the
$(q_w-1)q_w^{2z-1}$ vectors with $x_1\ne0$, where $|\Psi|=1$. The index is the order of
$\prod_{w\in\mathsf M}\auF_w^\times$. (c) At $w\in\mathsf M$ the reduction of $\pexK_\vartheta$ is
the stabiliser of the vector $e_1$ in $\auGL_{2z}(\auF_w)$, of index the number $q_w^{2z}-1$ of
non-zero vectors; at $w\in Q$ it is the stabiliser of a hyperplane, of index the number of
hyperplanes $(q_w^{2z}-1)/(q_w-1)$.
\end{proof}

\begin{lemma}[The unramified sum]\label{pex:vac}
Let $\pexmc=\pexmc_{2z,\alpha}$ be the coefficient of the vacuum of {\cite[\S2.3]{P53}} for $n=1$,
$M'=2z$ and the modulus $\mathfrak m$, on $G^{\mathfrak m}=\auU(1)(\auA^{\mathsf S}_{E^+,f})$ (the finite
points with trivial components at the places below $\mathsf S$). Then $\sum_{\gamma\in G^{\mathfrak
m}/K^{\mathfrak m}}|\pexmc(\gamma)|<\infty$ and
\[
\sum_{\gamma\in G^{\mathfrak m}/K^{\mathfrak m}}\pexmc(\gamma)=\frac{L^{\mathsf S}\bigl(z,(\nu\vartheta)^{-1}
\bigr)}{\zeta^{\mathsf S}_{E^+}(2z)},
\]
where $L^{\mathsf S}$ is as in \S\ref{po:ss-hecke} and $\zeta^{\mathsf S}_{E^+}$ is the Dedekind zeta
function of $E^+$ without the Euler factors at the places below $\mathsf S$.
\end{lemma}
\begin{proof}
$G^{\mathfrak m}/K^{\mathfrak m}$ is the restricted product of the $\auU(1)(E^+_v)/K_v$ over the
places $v$ of $E^+$ not below $\mathsf S$, and $\pexmc$ is the product of local coefficients
$\pexmc_v(\gamma_v)=\alpha_v(\gamma_v)^{-1}\langle\psvac_v,\omega_v(\gamma_v,1)\psvac_v\rangle$, equal to
$1$ on $K_v$ (by (W4); $\alpha$ is unramified at $v$). So it suffices to compute, for each $v$, the
local sum $\Sigma_v=\sum_{\gamma\in\auU(1)(E^+_v)/K_v}\pexmc_v(\gamma)$, and to check
$\sum_v(\sum|\pexmc_v|-1)<\infty$.

\emph{$v$ inert} (it is unramified). $\auU(1)(E^+_v)=\{x\in E_v^\times:x\aubar x=1\}$ consists of
units, so it equals $K_v$ and $\Sigma_v=1$. The Euler factor at the place $w=v$ of $E$ of
$L(s,(\nu\vartheta)^{-1})$ (the $L$-function of $x\mapsto(\nu\vartheta)^{-1}(x/\aubar x)$) is
$(1-q_w^{-s})^{-1}=(1-q_v^{-2s})^{-1}$, because a uniformiser $\varpi$ of $E^+_v$ is one of $E_w$
and $\varpi/\aubar\varpi=1$. So $\Sigma_v=1=L_w(z,(\nu\vartheta)^{-1})/\zeta_{E^+,v}(2z)$.

\emph{$v=w\aubar w$ split.} $\auU(1)(E^+_v)\cong E_w^\times$ by $x\mapsto x_w$ (the component at
$\aubar w$ is then $\aubar{x_w}^{-1}$), $K_v\cong\auO_w^\times$, and the classes are represented by
$\varpi_w^k$, $k\in\Z$. In the model (W5) at $v$, the space is that of Schwartz--Bruhat functions
on $E_w^{2z}$, $\psvac_v=\plone_{M_w}$ with $M_w=a_w\auO_w^{2z}$ ({\cite[Lemma~6.30]{P53}}); in the
coordinate $\tilde x=a_w^{-1}x$ of {\cite[Lemma~6.32]{P53}} this is $\plone[\tilde x\in\auO_w^{2z}]$, and
the integral below is computed in $\tilde x$, with the measure giving $\auO_w^{2z}$ volume $1$ (the
substitution $x=a_w\tilde x$ does not change its value). By \eqref{po:eq-w5} (with $n=2z$, the character
$\mu=\nu\vartheta$ on the source group, and $\alpha=\mu\upsilon^z$, whose $\upsilon$-part is the
splitting) $\pexmc_v(\varpi_w^k)=\chi_w(\varpi_w)^{-k}|\varpi_w^k|_w^{z}\int_{E_w^{2z}}\plone[\auO_w^{2z}](x)
\,\plone[\auO_w^{2z}](x\varpi_w^k)\,dx$, where $\chi_w$ is the (unramified) base change of
$\nu\vartheta$ at $w$; the modulus $|\pexmc_v(\varpi^k)|=q_w^{-|k|z}$ agrees with
{\cite[Lemma~2.11(a)]{P53}}. The integral is $1$ for $k\ge0$ and $q_w^{2zk}$ for $k<0$. With
$a=\chi_w(\varpi_w)^{-1}$, $|a|=1$, and $u=q_w^{-z}$,
\[
\Sigma_v=\sum_{k\ge0}(au)^k+\sum_{k\ge1}(a^{-1}u)^k=\frac{1-u^2}{(1-au)(1-a^{-1}u)} .
\]
Since $\chi_{\aubar w}(\varpi_{\aubar w})=\chi_w(\varpi_w)^{-1}$ (the character $x\mapsto\mu(x/\aubar x)$
is inverted by complex conjugation), the denominator is the inverse of the product of the Euler
factors at $w$ and $\aubar w$ of $L(z,(\nu\vartheta)^{-1})$, and $1-u^2=\zeta_{E^+,v}(2z)^{-1}$. The
sum of $|\pexmc_v|$ is $(1+u)/(1-u)\le1+3q_w^{-2}$ ($z\ge2$), and $\sum_wq_w^{-2}<\infty$.
Multiplying over $v$ gives the assertion; the products converge absolutely because $z\ge2$.
\end{proof}

\begin{lemma}[The archimedean factor]\label{pex:arch}
For $g\in\auU(1)(E^+\otimes\R)$, $\omega_{\alpha,\beta,\infty}(g,1)P=P$. Hence, in
{\cite[Proposition~2.16]{P53}} for the pair $(E,V^{(z)})$, the constant function on $\auU(1)_\infty$ as
the type of the source, and the polynomial $P$, the archimedean integral is
$Z_\infty=\langle P,P\rangle_{\mathrm{Fock}}>0$; it does not depend on $\vartheta$.
\end{lemma}
\begin{proof}
At $\tau$, $\auU(1)(\R)$ is the circle, acting on the Fock space of $\auHom(E,V^{(z)})\otimes_\tau\R=
\C^{2z}$ (all variables holomorphic, both spaces being definite) by the Weil representation, which in
the Fock model multiplies a homogeneous polynomial of degree $d$ in the holomorphic variables by
$\epsilon^{-d}$ times the value at $\epsilon$ of the archimedean splitting character
($\epsilon\in\auU(1)(\R)$; (N1)--(N5) of {\cite[\S6.1]{P49}}, \cite{KashiwaraVergne}). The factor
$\alpha_\tau(\epsilon)^{-1}$ in $\omega_{\alpha,\beta}$ cancels the splitting character (the
$\upsilon^z$-part of $\alpha$) and multiplies by $\epsilon^{c(\tau)-z}=\epsilon^{l_\tau}$, the
character $\nu_\tau^{-1}$ ($\nu$ has weight $-c$, $\vartheta_\tau=1$). Since $P$ is homogeneous of
degree $l_\tau$ at $\tau$, the product is trivial. Then the integrand of $Z_\infty$ in
{\cite[Lemma~2.10(A3)]{P53}} is the constant $\langle P,P\rangle\cdot1$ on a group of total mass $1$. The
vector $P$ is non-zero because $v^+\ne0$ and $\xi\mapsto\pop_\xi$ is injective on the weight
(\S\ref{po:ss-set}). (This is the condition that the constant function lifts into a space of the
weight of \S\ref{po:ss-hecke}, i.e.\ (A2) of {\cite[Lemma~2.10]{P53}} for this pair; (A1) holds because
$[\auU(1)]$ and $[\auU(V^{(z)})]$ are compact.)
\end{proof}

\subsubsection{The norm of the lift}\label{pex:ss-rallis}

\begin{proposition}[Norm of the lift of a character]\label{pex:rallis}
Let $p=(\nu,z)$ be a critical pair with $z\ge2$. There is a real number $A_p>0$, independent of
$\vartheta$, such that for every tame character $\vartheta$ of modulus $Q$
\[
\|\varphi_{p,\vartheta}\|^2=A_p\,(\auNm Q)^{2z-1}\,L^{\mathsf M\cup Q}\bigl(z,(\nu\vartheta)^{-1}\bigr),
\]
and $L^{\mathsf M\cup Q}(z,(\nu\vartheta)^{-1})>0$. Explicitly
\[
A_p=|\kappa_p|^2\,\prod_{w\in\mathsf M}(q_w-1)\cdot\dim W_{\lambda}\cdot\|v^+\|^{-2}\cdot c'\,
\langle P,P\rangle_{\mathrm{Fock}}\,\mathrm{vol}([\auU(V^{(z)})])\,\mathrm{vol}([\auU(1)])\,
\zeta_{E^+}(2z)^{-1},
\]
with $c'\ne0$ the constant of {\cite[Lemma~2.14(b)]{P53}} (all factors of $A_p$ are non-zero, and $A_p$ is real and positive, see the proof). In particular, for twists $\vartheta$, $\rho_p(\vartheta)=\|\varphi_{p,\vartheta}\|^2/
\bigl((\auNm Q_\vartheta)^{2z-1}\|\varphi_{p,1}\|^2\bigr)$.
\end{proposition}

\begin{proof}
\emph{Step 1: the inner product formula.} We apply {\cite[Proposition~2.16]{P53}} to $S=E$ (rank $n=1$,
definite), $T=V^{(z)}$ (rank $M'=2z$, definite), the characters $\alpha,\beta$ of
\S\ref{pex:ss-set}, the modulus $\mathfrak m$, the polynomial $P$, the function $\Psi_\vartheta$,
which is sharp for $J$ (Lemma~\ref{pex:sharp}(a)), and $y=y'=\plone$, which is $J$-invariant. We
check that its proof applies. Step~1 there uses (W7), (SW) and {\cite[Lemma~2.14]{P53}}: the component of
$[P,\Psi_\vartheta]$ at a place of $\mathsf M$ is supported on injective $2z\times1$ matrices
(columns $\equiv e_1$, hence non-zero), as {\cite[Lemma~2.14]{P53}} requires; (SW) is used for
$\pcW=E\oplus E^-$ of rank $2$ and $T$ of rank $2z$. Both conditions of (SW) in {\cite[Lemma~2.13]{P53}} hold: (i) the theta
integral converges absolutely because $T$ is definite at every infinite place, hence anisotropic; (ii)
$M'=2z>2=2n$. So (SW) applies directly (for $z=2$ this is the case $M'=4$, $n=1$; Weil's
identity~\cite[Th\'eor\`eme~5]{WeilSiegel} in its range of absolute convergence); the height bounds of (SW) are not needed, since $g,g'$ run over the compact space $[\auU(1)]$. Step~2 uses only that $y,y'$
are bounded on the compact space $[\auU(1)]$. In Step~3 the series is $\sum_\gamma\pexmc(\gamma)\rho
(\gamma)$, which converges absolutely by Lemma~\ref{pex:vac} ({\cite[Lemma~2.11(b)]{P53}},(c) is stated
for $n\ge2$ and is replaced by Lemma~\ref{pex:vac}); applied to $\plone$ it gives the constant
$\Sigma\cdot\plone$, $\Sigma=\sum_\gamma\pexmc(\gamma)$. The archimedean factor is $Z_\infty$ of
Lemma~\ref{pex:arch}. Hence, with \eqref{pex:eq-f},
\[
\langle f_{\Theta^{(1)}_\vartheta\plone},f_{\Theta^{(1)}_\vartheta\plone}\rangle_{[\auU(V^{(z)})]}
=c'Z_\infty\,\mathrm{vol}([\auU(V^{(z)})])\cdot\frac{\|\Psi_\vartheta\|^2}{\#\pexMat_{2z\times1}
(\pcR_{\mathfrak m})\,[K^1:J]}\cdot\mathrm{vol}([\auU(1)])\,\Sigma .
\]

\emph{Step 2: the algebraic norm.} By {\cite[Lemma~2.18]{P53}},
$\|\Theta^{(1)}_\vartheta\plone\|^2=[K_{V^{(z)}}:\pexK_\vartheta]\cdot\dim W_\lambda\cdot\|v^+\|^{-2}\cdot
\langle f,f\rangle_{[\auU(V^{(z)})]}$, $f=f_{\Theta^{(1)}_\vartheta\plone}$.

\emph{Step 3: collecting the factors.} By Lemma~\ref{pex:vac},
$\Sigma=L^{\mathsf S}(z,(\nu\vartheta)^{-1})\,\zeta_{E^+}(2z)^{-1}\prod_{w\in\mathsf S}(1-q_w^{-2z})^{-1}$
(the places of $\mathsf S$ lie above distinct split places $v$ of $E^+$ with $q_v=q_w$, because
$\mathsf S\cap\aubar{\mathsf S}=\emptyset$). Insert Lemma~\ref{pex:sharp}(b),(c). The factors
depending on a place $w$ of $\mathsf S$ are:
\[
w\in\mathsf M:\ \frac{q_w^{-2z}}{q_w-1}\cdot(q_w^{2z}-1)\cdot(1-q_w^{-2z})^{-1}=\frac1{q_w-1};\qquad
w\in Q:\ \frac{q_w-1}{q_w}\cdot\frac{q_w^{2z}-1}{q_w-1}\cdot(1-q_w^{-2z})^{-1}=q_w^{2z-1}.
\]
So $\|\Theta^{(1)}_\vartheta\plone\|^2=B\,(\auNm Q)^{2z-1}L^{\mathsf S}(z,(\nu\vartheta)^{-1})$ with
\[B=\prod_{w\in\mathsf M}(q_w-1)^{-1}\cdot\dim W_\lambda\,\|v^+\|^{-2}\,c'Z_\infty\,\mathrm{vol}([\auU(V^{(z)})])
\,\mathrm{vol}([\auU(1)])\,\zeta_{E^+}(2z)^{-1},\]
which is independent of $\vartheta$. Multiplying by
$|\kappa_p|^2[K^1:K^1_M]^2=|\kappa_p|^2\prod_{w\in\mathsf M}(q_w-1)^2$ gives the formula with the stated
$A_p$.

\emph{Positivity.} $L^{\mathsf S}(z,(\nu\vartheta)^{-1})$ is a convergent Euler product ($z\ge2$) of
non-zero factors, and it is real and positive: the factors at the places $w\ne\aubar w$ not in
$\mathsf S\cup\aubar{\mathsf S}$ come in pairs $(1-aq^{-z})^{-1}(1-\aubar aq^{-z})^{-1}=|1-aq^{-z}|^{-2}$,
$|a|=1$, and the factors at $w=\aubar w$ are $(1-q_w^{-z})^{-1}>0$ (proof of Lemma~\ref{pex:vac}).
All factors of $A_p$ are non-zero: $c'\ne0$ by {\cite[Lemma~2.14(b)]{P53}},
$Z_\infty$, hence $\langle P,P\rangle_{\mathrm{Fock}}$, is non-zero by Lemma~\ref{pex:arch}, the volumes,
$\dim W_\lambda$, $\|v^+\|^{-2}$ and $\zeta_{E^+}(2z)^{-1}$ are positive, and $\kappa_p\ne0$; so $A_p\ne0$.
Since $\|\varphi_{p,\vartheta}\|^2\ge0$ and $(\auNm Q)^{2z-1}L^{\mathsf M\cup Q}(z,(\nu\vartheta)^{-1})>0$,
the formula shows that $A_p$ is real and $\ge0$; hence $A_p>0$, and $\|\varphi_{p,\vartheta}\|^2>0$, so $\varphi_{p,\vartheta}
\ne0$. The last assertion follows by dividing the formula for $\vartheta$ by that for $\vartheta=1$
(modulus $\emptyset$) and recalling $\rho_p(\vartheta)=|L^{\mathsf M\cup Q}(z,(\nu\vartheta)^{-1})|/
|L^{\mathsf M}(z,\nu^{-1})|$.
\end{proof}

\begin{remark}\label{pex:rem-rallis}
(i) Proposition~\ref{po:rallis} is the case of twists; the extension to all
tame characters of modulus $Q$ (in particular to a twist of modulus $Q$ regarded as a tame
character of a larger modulus) is what (IM)(c) needs. (ii) The proof uses the following inputs of
other sections: {\cite[Proposition~2.16]{P53}} and {\cite[Lemmas~2.13, 2.14 and~2.18]{P53}};
(SW) is used for an anisotropic target
with $M'=2z>2n$, inside the range of {\cite[Lemma~2.13]{P53}} (absolute convergence of Weil's theorem). {\cite[Lemma~2.10]{P53}} is not used: its
content for this pair is Lemma~\ref{pex:arch}. (iii) The constant $A_p$ contains the volumes of the
two automorphic quotients and $c'$, which we do not compute; only $A_p>0$ and its independence of
$\vartheta$ are used.
\end{remark}

\subsubsection{Integral measures: statement {\rm(IM)} for $z\ge2$}\label{pex:ss-IM}

Let $\mathsf N_z=\#\auGL_{4fz}(\auF_3)$. By the proof of Lemma~\ref{po:naut}, the order of every
group $\Gamma\cap h\,\pexK h^{-1}$, $h\in\auU(V^{(z)})(\auA_{E^+,f})$, $\pexK\subseteq K_{V^{(z)}}$, divides
$\mathsf N_z$, and for forms $\Xi,\Xi'\in\mathcal A(V^{(z)};p,\vartheta)$ integral at $\aubar\iota$
and at $\iota$, $\mathsf N_z\langle\Xi,\Xi'\rangle\in\auO_\iota$ (same proof, with
$2f\cdot2z$ in place of $2fn$).

\emph{Kernel integrality.} We use the following consequence of {\cite[Theorem~7.15]{P49}}, proved exactly
as {\cite[Proposition~5.30]{P49}} and {\cite[Corollary~7.16(a)]{P49}} (the kernel is the $\auU(1)$, $\auU(2z)$
instance of {\cite[(34)]{P49}}): \emph{there is an integer $n_p\ge1$ such that for every tame
$\vartheta$ of modulus $Q$, every prime $\ell\ne d$ with $\ell\nmid\auNm Q$, every
$\jmath\colon\C\to\C_\ell$, every $x\in\auU(1)(\auA_{E^+,f})$ and $h\in\auU(V^{(z)})(\auA_{E^+,f})$, the
vector $n_p\,h_\ell^{-1}k^0_Q(x,h)$ has integral coordinates at $\jmath$}, where $k^0_Q(x,h)\in
W_\lambda$ is the kernel $k_{\Phi}(x,h;\cdot)$ of \S\ref{po:ss-set} (scaled by $\kappa_p$) for the
character $\nu$ and the \emph{point datum} $\Phi^0_Q=\bigotimes_{w\in\mathsf M}\plone[x\equiv e_1]
\otimes\bigotimes_{w\in Q}\plone[x\text{ integral},\,x_1\equiv1\bmod w]\otimes(\text{vacuum})$.
Indeed, by finite Fourier analysis at the places of $\mathsf M\cup Q$ (the indicator of a coset of a
sublattice of index $q$ of $\auO_w^{2z}$ is $q^{-1}$ times a sum of $q$ Heisenberg translates of the
local vacuum, with roots of unity as coefficients; the index is $q_w^{2z}$ at $w\in\mathsf M$ and
$q_w$ at $w\in Q$, and $\prod_{w\in\mathsf M}q_w^{2z}=M^{2fz}$), the function $M^{2fz}\auNm Q\cdot
c_{\mathrm{vac}}^{2z}\,\Phi^0_Q$ is a standard datum of level $\aubar{\pi_M}$ at every prime $\ell$;
$\auNm Q$ is a unit at $\ell$, and the values of $\nu$ and $\upsilon$ at ideles
that are units above $\ell$ are units. So $n_p$ may be taken to be $M^{2fz}$ times the constant of
{\cite[Proposition~5.30]{P49}} for this instance; it depends only on $p$
(not on $\vartheta$ or $Q$), and its value is not used (for flat pairs, $l_\tau=0$, one may take
$n_p=M^{2fz}$ at every $\ell\nmid\auNm Q$: see the proof of Lemma~\ref{psu:I}(b)).
The prime $\ell=N$ is covered by {\cite[Corollary~7.16(b)]{P49}} in the same way.

\begin{lemma}[Kernel integrality at $\ell=d$]\label{pex:dint}
Let $\ell=d$, $\iota\colon\C\to\C_d$, $\delta\in\auO_K\smallsetminus\{0\}$ and $f$ a standard datum of
level $\delta$ at $d$. Then $\Theta_f(d\,\aubar\delta s_0\,\cdot)$ is $\partial$-integral; hence for every
$k$ and every homogeneous polynomial $Q$ of degree $k$ with coefficients in $\pao$, $d^kQ(\partial)
\Theta_f(\aubar\delta s_0\,\cdot)(0)\in\pao$.
\end{lemma}
\begin{proof}
\emph{The curve at $d$.} Normalise the valuation by $v(d)=1$. In the proof of
{\cite[Proposition~5.25(C2)]{P49}}, $x$, $y$ are $\partial$-integral at every prime, $g^{(n)}(0)=x^{(n-2)}(0)$
for $n\ge3$, and $-2d\,g''(0)=\sum_{u}x(u-t)-\tfrac{d+1}{12}a_1^2$, the sum over the non-zero points of
order $d$. A point of order $d$ either has integral coordinates or lies in the group $\mathcal E_1$ of
points reducing to the origin of the integral model, which is the group of the formal group over the
integers $\mathfrak o$ of a finite extension of $\Q_d$ \cite[Chapter~VII, Proposition~2.2]{SilvermanAEC} (that proposition
assumes a minimal Weierstrass equation; the model $y^2+a_1xy=x^3+x$ of {\cite[(25)]{P49}} is
minimal at $d$, since $a_1$ is integral and its discriminant $a_1^4-64$ is a unit at the odd prime $d$,
{\cite[Lemma~4.5]{P49}}); a torsion point of
exact order $d$ there has parameter $\mathsf z$ with $v(\mathsf z)\le v(d)/(d-1)$
\cite[Chapter~IV, Theorem~6.1]{SilvermanAEC}, and $x=\mathsf z^{-2}(1+\mathsf z\,\mathfrak o[[\mathsf z]])$,
so $v(x)\ge-2/(d-1)$. Since $2\ne d$ and $12\mid d+1$ ({\cite[Proposition~5.25]{P49}}, proof of (C2)),
$v(g''(0))\ge-1-2/(d-1)>-2$ ($d\ge5$), so $d^2g''(0)\in\pao$. Hence $g(dw)=\sum_{n\ge2}g^{(n)}(0)d^nw^n/n!$
is $\partial$-integral with zero constant term, and $\theta(dw)=\exp(-g(dw))$ and $y(dw)$ are
$\partial$-integral (the rules recalled before {\cite[Proposition~5.25]{P49}}).

\emph{The datum.} {\cite[Theorem~7.15]{P49}} holds at every prime: $\Theta_f(\aubar\delta s_0z)=(\prod_j\theta
(\mu_j(z))^{m_j})^{\auNm\delta}\mathsf H(z)$ with $\mathsf H$ an integral power series in the
$y(\mu_j(z))$ and $\mu_j$ linear forms with integral coefficients ({\cite[Lemma~7.12]{P49}},
{\cite[Corollary~7.3]{P49}}, both valid at every prime). Replace $z$ by $dz$: $\mu_j(dz)=d\mu_j(z)$. A
one-variable $\partial$-integral series composed with a linear form with integral coefficients is
$\partial$-integral in several variables (expand $(\sum_ia_iz_i)^n/n!=\sum\prod_ia_i^{k_i}z_i^{k_i}/k_i!$),
so $\theta(d\mu_j(z))$ and $y(d\mu_j(z))$ are $\partial$-integral, the latter with zero constant term;
products and $\mathsf H(\dots)$ of them are $\partial$-integral. Finally a jet of degree $k$ of
$\Theta_f(\aubar\delta s_0\,\cdot)$ is $d^{-k}$ times the same jet of $\Theta_f(d\aubar\delta s_0\,\cdot)$.
\end{proof}

Consequently kernel integrality holds at $\ell=d$ (if $d\nmid N\auNm Q$) with $n_p$ replaced by
$n_pd^{\sum_\tau l_\tau}$: the reduction of {\cite[Proposition~5.30]{P49}} is independent of the prime, and the
jets are of degree $\sum_\tau l_\tau$. From now on $n_p$ includes it, so kernel integrality
holds at all $\ell\nmid\auNm Q$.

\begin{lemma}[The lift as a sum over points]\label{pex:points}
Let $Q$ be a finite set of places as above and $\mathcal P_Q=\bigl(\auU(1)(E^+)\backslash
\auU(1)(\auA_{E^+,f})/K^1_M\bigr)\times\prod_{w\in Q}\auF_w^\times$. Choose representatives $x_c$ of
the classes and, for $a=(a_w)\in\prod_{w\in Q}\auF_w^\times$, an integral $u_a\in\auU(1)(\auA_{E^+,f})$
equal to $1$ outside $Q\cup\aubar Q$ whose component at $w\in Q$ reduces to $a_w^{-1}$. Put
$x_{c,a}=x_cu_a$. Then for every tame $\vartheta$ of modulus $Q$ and every $h$,
\[
\varphi_{p,\vartheta}(h)=\sum_{(c,a)\in\mathcal P_Q}\vartheta(x_{c,a})^{-1}\,k^0_Q(x_{c,a},h).
\]
\end{lemma}
\begin{proof}
In \eqref{po:eq-w5} the character $\vartheta$ enters $\omega_\mu$, $\mu=\nu\vartheta$, only through
the scalar $\vartheta(\det g)^{-1}$ on $\auU(1)$ (it is trivial at the infinite places); so
$k_{\Phi}(g,h;\xi)$ for $\nu\vartheta$ is $\vartheta(g)^{-1}$ times the kernel for $\nu$. At $w\in Q$,
$\Phi^{(1)}_{\vartheta,w}=\sum_{a\in\auF_w^\times}\psi_w(a)\,\plone[x\text{ integral},x_1\equiv a]$, and
$\plone[x\text{ integral},x_1\equiv a_w]=\omega_\nu(u_a,1)\plone[x\text{ integral},x_1\equiv1]$ by
\eqref{po:eq-w5}, $\nu$ being unramified at $w$ (the unramified scalar is $1$ on units). Since the
component of $\vartheta$ at $w$ is $\psi_w$ on units, $\psi_w(a_w)=\vartheta_w(u_a)^{-1}$. Hence the
integrand of $\Theta^{(1)}_\vartheta(\plone)$ is $F(g)=\sum_a\vartheta(gu_a)^{-1}k^0_Q(gu_a,h)/\kappa_p$.
This function of $g\in[\auU(1)]$ is right invariant under $K^1_M$: for $k\in K^1_M$ the elements $ku_a$
run over representatives of the $u_a$ modulo the units $\equiv1$ at $Q$, under which the summand is
invariant ($\Phi^0_Q$ is invariant under them, by \eqref{po:eq-w5}). So $\int_{[\auU(1)]}F=
\mathrm{vol}(K^1_M)\sum_cF(x_c)$, the stabilisers $\auU(1)(E^+)\cap x_cK^1_Mx_c^{-1}$ being trivial
(the group $K^1_M$ is abelian, so such an element is a root of unity $\zeta\in E$ with $\zeta\equiv1$
modulo every $w\in\mathsf M$; if $\zeta$ has exact order $m>1$, then $\zeta-1$ is a unit when $m$ is not a
prime power and divides $r$ when $m$ is a power of a prime $r$, so $\zeta-1$ is a unit at every place
not above $r$; as $\mathsf M$ contains places above each of the two distinct primes $p_0,p_1$, this is
impossible, so $\zeta=1$), and
the volume of $[\auU(1)]$-measure is normalised by $\mathrm{vol}(K^1)=1$ ({\cite[\S2.3]{P53}}), so
$\mathrm{vol}(K^1_M)=[K^1:K^1_M]^{-1}$. Multiply by $\kappa_p[K^1:K^1_M]$.
\end{proof}

\begin{proposition}[Statement (IM) for $z\ge2$]\label{pex:IM}
Let $p=(\nu,z)$ be a ramified critical pair with $z\ge2$. Put $I=\{z\}$, $m=0$ and
\[
F(\vartheta)=\mathsf N_z\,n_p^2\,\frac{\|\varphi_{p,\vartheta}\|^2}{(\auNm Q)^{2z-1}}\qquad
(\vartheta\text{ tame of modulus }Q).
\]
Then: (a) $F(1)\ne0$; (b) for every $Q$ there are finitely many $x_a\in\auU(1)(\auA_{E^+,f})$ and
$c_a\in\C$, with $c_a$ integral at every prime $\ell\nmid\auNm Q$ under every embedding $\C\to\C_\ell$,
such that $F(\vartheta)=\sum_ac_a\vartheta(x_a)$ for every tame $\vartheta$ of modulus $Q$; (c) if $v\notin Q$ and $\vartheta'$ is
$\vartheta$ regarded as of modulus $Q\cup\{v\}$, then $F(\vartheta')=F(\vartheta)\prod_\pm(1-\alpha_v^{\pm1}
q_v^{-z})$, $\alpha_v$ the value at a uniformiser of the base change of $\nu\vartheta$ at $v$; (d) for
every twist $\vartheta$ of modulus $Q$, $F(\vartheta)=F(1)\,\rho_p(\vartheta)$.
\end{proposition}
\begin{proof}
By Proposition~\ref{pex:rallis}, $F(\vartheta)=\mathsf N_zn_p^2A_p\,L^{\mathsf M\cup Q}(z,(\nu\vartheta)^{-1})$
with $\mathsf N_zn_p^2A_p>0$ independent of $\vartheta$. (a) $L^{\mathsf M}(z,\nu^{-1})>0$. (d) is the
quotient of the formulas for $\vartheta$ and $1$, the $L$-values being positive. (c) The two
$L$-functions differ by the Euler factors at $v$ and $\aubar v$ (we may assume $v$, $\aubar v\notin
\mathsf M\cup Q\cup\aubar Q$, as for an auxiliary modulus; then $\vartheta$ is unramified at $v$ and the
place below $v$ splits): $L^{\mathsf S\cup\{v\}}=L^{\mathsf S}\cdot(1-aq_v^{-z})(1-a^{-1}q_v^{-z})$ with
$a$ the value at $\varpi_v$ of the base change of $(\nu\vartheta)^{-1}$, using
$\chi_{\aubar v}(\varpi_{\aubar v})=\chi_v(\varpi_v)^{-1}$ as in Lemma~\ref{pex:vac}; the product is
symmetric in $a\leftrightarrow a^{-1}=\alpha_v$.

(b) Fix $Q$ and the points $x_{c,a}$ of Lemma~\ref{pex:points}, which depend only on $Q$. The level
$\pexK_\vartheta$ and the class representatives $h$ of $\Gamma\backslash\auU(V^{(z)})(\auA_f)/\pexK_\vartheta$
depend only on $Q$; let $e_h=\#(\Gamma\cap h\pexK_\vartheta h^{-1})$. By Lemma~\ref{pex:points} and the
definition of the norm ({\cite[\S14.3]{P46}}),
\[
\|\varphi_{p,\vartheta}\|^2=\sum_h\frac1{e_h}\bigl\langle\varphi_{p,\vartheta}(h),\varphi_{p,\vartheta}(h)
\bigr\rangle=\sum_{\mathsf p,\mathsf p'\in\mathcal P_Q}\vartheta\bigl(x_{\mathsf p}x_{\mathsf p'}^{-1}\bigr)
\,G_{\mathsf p\mathsf p'},\qquad G_{\mathsf p\mathsf p'}=\sum_h\frac1{e_h}\bigl\langle k^0_Q(x_{\mathsf
p},h),k^0_Q(x_{\mathsf p'},h)\bigr\rangle,
\]
using $\aubar{\vartheta(x)^{-1}}=\vartheta(x)$ and the sesquilinearity of the form (conjugate-linear in
the first variable, the convention of {\cite[\S2.3]{P53}}); $G_{\mathsf p\mathsf p'}$ depends only on $Q$. Put $x_{(\mathsf p,\mathsf
p')}=x_{\mathsf p}x_{\mathsf p'}^{-1}$ and $c_{(\mathsf p,\mathsf p')}=\mathsf N_zn_p^2G_{\mathsf p\mathsf
p'}(\auNm Q)^{1-2z}$. Then $F(\vartheta)=\sum c_{(\mathsf p,\mathsf p')}\vartheta(x_{(\mathsf p,\mathsf
p')})$ for every tame $\vartheta$ of modulus $Q$. Let $\ell\nmid\auNm Q$ and $\jmath\colon\C\to\C_\ell$.
The form on $W_\lambda$ is invariant under $h_\ell$ and the Gram matrix of the basis of $W_\lambda(\Z)$ is integral ({\cite[\S14.3]{P46}}), so
$\langle k^0_Q(x,h),k^0_Q(x',h)\rangle=\langle h_\ell^{-1}k^0_Q(x,h),h_\ell^{-1}k^0_Q(x',h)\rangle$ is,
after multiplication by $n_p^2$, a sum of products of a coordinate integral at $\jmath\circ c$ (complex
conjugation; the conjugate of a coordinate integral at $\jmath\circ c$ is integral at $\jmath$) and one
integral at $\jmath$; so $n_p^2\langle\cdot,\cdot\rangle$ is integral at $\jmath$, by kernel
integrality applied at $\jmath$ and at $\jmath\circ c$. As $e_h\mid\mathsf N_z$, $\mathsf N_zn_p^2G_{\mathsf
p\mathsf p'}$ is integral at $\jmath$, and $\auNm Q$ is a unit at $\ell$.
\end{proof}

\begin{remark}[On the form of (IM)]\label{pex:rem-IM}
(i) In (d) the ratio is $\rho_p(\vartheta)=|L^{\mathsf M\cup Q}(z,(\nu\vartheta)^{-1})|/|L^{\mathsf M}(z,
\nu^{-1})|$, with the Euler factors at $\mathsf M$ removed in both and at $Q$ in the numerator; this is
the convention of $r(\vartheta)$ in {\cite[\S3.5]{P55}} (``conventions of
{\cite[Theorem~2.2(b)]{P55}}'') if those conventions are read as in \S\ref{po:ss-hecke}; at the places of $Q$
the twisted character is ramified, so for a twist nothing is removed there. (ii) In (c) the Euler
factor is stated for $v$ with $v,\aubar v\notin\mathsf M\cup Q\cup\aubar Q$, which is the case used in
{\cite[(6)]{P55}} ($Q\cup\{v\}$ auxiliary). (iii) In (b) the points $x_a$ are finite ideles of
$\auU(1)$ and $\vartheta(x_a)$ is the value of $\vartheta$ there; the use in {\cite[\S3.5]{P55}} of
$[x_a^{-\delta}]$ in the class group is unaffected. (iv) $\ell=d$ is covered by Lemma~\ref{pex:dint}.
\end{remark}

\subsubsection{Statement {\rm(IM)} for $z=1$}\label{pex:ss-IM1}

Let $p=(\nu,1)$ be a ramified critical pair; $\nu$ has weight $-c$, $c(\tau)\ge1$. Let $I_2=E^2$ be the
standard positive definite plane, $\vartheta$ a tame character of modulus $Q$, $\mathfrak m_\vartheta=
(\pi_M)\prod_{w\in Q}w$, $\chi_\vartheta$ the character of $\pcR_{\mathfrak m_\vartheta}^\times$ given by the
tame parts of $\nu\vartheta$ (trivial at the places of $Q$ where $\vartheta$ is unramified), extended by
$0$, and $\lambda_\vartheta=\nu\vartheta\upsilon$ (weight $-(c-1)$). Put
$P_c=\prod_\tau\det(y_\tau)^{c(\tau)-1}$, a polynomial on $\auHom(I_2,I_2)\otimes\R=\prod_\tau M_2(\C)$,
and define the \emph{period}
\[
\Lambda(\vartheta)=\sum_{h\in\Gamma\backslash\auU(I_2)(\auA_f)/K_{I_2}}\frac1{e_h}\,\theta^{I_2,I_2}_{
\lambda_\vartheta,\lambda_\vartheta}\bigl(1,h;P_c,\chi_\vartheta\circ\det\bigr),\qquad e_h=\#(\Gamma\cap
hK_{I_2}h^{-1}),
\]
in the notation of {\cite[\S2.3]{P53}} (the datum $[P_c,\chi_\vartheta\circ\det]$ modulo
$\mathfrak m_\vartheta$). The summand is right $K_{I_2}$-invariant in $h$ by {\cite[Lemma~2.9]{P53}}: for
$k\in K_{I_2}$ the factor $\lambda_\vartheta(\det k)$ is $\chi_\vartheta(\det\bar k)$ and cancels against
$\chi_\vartheta(\det(\bar k^{-1}x))$. (One may equally take the modulus $(\pi_M)^2\prod w$; as
$\chi_\vartheta$ is tame this changes nothing.)

\begin{lemma}[The isotypic vectors]\label{pex:iso}
Let $k=(k_\tau)\in\Z_{\ge0}^{\auPl}$. A polynomial $R$ on $\prod_\tau M_2(\C)$ with
$R((h_\tau^{-1}y_\tau g_\tau)_\tau)=\prod_\tau\det(g_\tau)^{k_\tau}\det(h_\tau)^{-k_\tau}R(y)$ for all
$g,h\in\prod_\tau\auU(2)$ is a multiple of $\prod_\tau\det(y_\tau)^{k_\tau}$. A polynomial with this
property for some $k\in\Z^{\auPl}$ with some $k_\tau<0$ is $0$.
\end{lemma}
\begin{proof}
Both sides are polynomial in the entries of $g_\tau,h_\tau$ and their inverses, and $\auU(2)$ is
Zariski dense in $\auGL_2(\C)$ (its Lie algebra spans $\mathfrak{gl}_2(\C)$ over $\C$), so the identity
holds for $g,h\in\prod\auGL_2(\C)$. Take $h=1$, $y=1$ and $g=y'$ invertible: $R(y')=\prod_\tau\det
(y'_\tau)^{k_\tau}R(1)$ on a dense set. If some $k_\tau<0$ the right side has a pole unless $R(1)=0$.
\end{proof}

\begin{proposition}[The equal-rank period; (E1)]\label{pex:E1}
There is $C_1>0$ such that for $\vartheta=1$ and for every twist $\vartheta$ of modulus $Q$,
\[
C_1\,|\Lambda(\vartheta)|^2=L_*\bigl(\mathfrak m_\vartheta,(\nu\vartheta)^{-1},1\bigr)=\auNm(\mathsf M\cup Q)
\,\bigl|L^{\mathsf M\cup Q}\bigl(1,(\nu\vartheta)^{-1}\bigr)\bigr|^2\ne0 .
\]
\end{proposition}
\begin{proof}
Apply {\cite[Theorem~3.2]{P53}} with $W=V=I_2$, $V'=0$: $N=a=2$, $b=0$, $e=1$, $q=q_V=q_{V'}=0$, $2q+2\le N$;
$\mu=c$, and $N\le2\mu_\tau$ since $c(\tau)\ge1$; $v$ is the (unique up to a scalar) vector $1$ of the
Fock space of the zero space. This gives $C>0$ and a Fock vector $x$ of $\auHom(I_2,I_2)\otimes\R$,
independent of the modulus and the character. The character $\nu'=(\nu\vartheta)^{-1}$ is a tame
character for $\mathfrak m_\vartheta$ of weight $c$ (its tame part is non-trivial at the places of
$\mathsf M$ because $p$ is ramified, and at those of $Q$ because $\vartheta$ is a twist and $\nu$ is
unramified there; for $\vartheta=1$, $\mathfrak m=(\pi_M)$). Then $\gamma=\nu'^{-1}\upsilon=\lambda_\vartheta$;
take $\alpha=\lambda_\vartheta$, $\alpha'=1$. Matched functions for $b=0$: $\Psi$ on $\pexMat_{2\times0}$
is the constant $1$, and $\Phi$ is matched iff $\Phi(k^{-1})=\chi(\det k)$ for $k\in\auGL_2(\pcR)$ and
$\Phi=0$ off $\auGL_2$, where $\chi$ is the character of {\cite[Definition~3.1]{P53}} attached to
$\nu'$, i.e.\ $\chi=\chi_\vartheta^{-1}$; so $\Phi=\chi_\vartheta\circ\det$. The right side of
{\cite[(12)]{P53}} is $c$ times the kernel of $\auHom(0,I_2)=0$ with character $\alpha'=1$, i.e.\ $c$. So
\begin{equation}\label{pex:eq-E1}
\int_{[\auU(I_2)]}\theta^{I_2,I_2}_{\lambda_\vartheta,\lambda_\vartheta}(g,h;x,\chi_\vartheta\circ\det)\,dh=c,
\qquad|c|^2=C\,L_*(\mathfrak m_\vartheta,\nu',1),
\end{equation}
for all $g\in\auU(I_2)(\auA)$, and $c\ne0$ ({\cite[Corollary~3.4]{P53}}).

\emph{Replacing $x$ by $P_c$.} The compact group $K_\infty=\auU(I_2)_\infty\times\auU(I_2)_\infty$ acts on
the Fock space by $\omega_{\lambda,\lambda,\infty}$, and $\theta(g k_1,hk_2;x)=\theta(g,h;\omega_\infty
(k_1,k_2)x)$. Integrate \eqref{pex:eq-E1}, with $g$ replaced by $gk_1$, against $dk_1$ over
$\auU(I_2)_\infty$ (total mass $1$), and use the invariance of $dh$ under $h\mapsto hk_2$: \eqref{pex:eq-E1}
holds with $x$ replaced by its projection $x'$ onto the $K_\infty$-invariant vectors (the integral over
$k_2$ is built into $\int dh$, the one over $k_1$ fixes the constant $c$). The $K_\infty$-invariant
vectors of the Fock space are polynomials (the $K_\infty$-types of the Fock space of a pair of definite
spaces are finite-dimensional spaces of polynomials). In the Fock model $\omega_\infty(k_1,k_2)R(y)$ is
$R(k_2^{-1}yk_1)$ times the splitting factor and $\lambda_\infty(\det k_1)^{-1}\lambda_\infty(\det k_2)$
({\cite[\S2.3]{P53}}); since $\lambda=\nu\vartheta\upsilon$ has weight $-(c-1)$ and the splitting by
$\upsilon$ is cancelled as in Lemma~\ref{pex:arch}, invariance means $R(k_2^{-1}yk_1)=\prod_\tau\det
(k_{1,\tau})^{c_\tau-1}\det(k_{2,\tau})^{-(c_\tau-1)}R(y)$ (this is the condition that characterises $P_c$). By Lemma~\ref{pex:iso}, $x'=\kappa P_c$, and $\kappa\ne0$
because $c\ne0$. $\kappa$ depends only on $x$ and the archimedean characters, not on $\vartheta$.

\emph{From the integral to the class sum.} By {\cite[Lemma~2.9]{P53}} and the invariance under
$\auU(I_2)_\infty$ just used, the integrand of \eqref{pex:eq-E1} at $g=1$ with $x'$ is a function of the
class of $h$ in $\Gamma\backslash\auU(I_2)(\auA_f)/K_{I_2}$; with $\mathrm{vol}(K_{I_2})=1$ and
$\mathrm{vol}(\auU(I_2)_\infty)=1$, $\int_{[\auU(I_2)]}=\sum_h e_h^{-1}(\cdot)(h)$. So
$\kappa\Lambda(\vartheta)=c$, and $C_1=|\kappa|^2/C$.
\end{proof}

\begin{proposition}[Statement (IM) for $z=1$]\label{pex:IM1}
Let $p=(\nu,1)$ be a ramified critical pair, $I=\{0,1\}$, $m=1$, and $F(\vartheta)=|\mathsf N_1n_1
\Lambda(\vartheta)|^2$, with $\mathsf N_1=\#\auGL_{4f}(\auF_3)$ and $n_1\ge1$ an integer, depending only on $p$, making the kernel values below integral at all
$\ell\nmid\auNm Q$ (it exists by {\cite[Proposition~5.30]{P49}}, {\cite[Corollary~7.16]{P49}} and Lemma~\ref{pex:dint}). Then {\rm(a)}--{\rm(d)}
of {\rm(IM)} hold: $F(1)\ne0$; $F$ is an integral measure at every modulus; for $v,\aubar v\notin\mathsf M
\cup Q\cup\aubar Q$, $F(\vartheta')=F(\vartheta)\prod_{e\in\{0,1\}}\prod_\pm(1-\alpha_v^{\pm1}q_v^{-e})$; and
$F(\vartheta)=F(1)\,\auNm Q\,\rho_p(\vartheta)^2$ for every twist.
\end{proposition}
\begin{proof}
(a) and (d): Proposition~\ref{pex:E1}, since $L_*(\mathfrak m_\vartheta,\nu',1)/L_*((\pi_M),\nu^{-1},1)=
\auNm Q\cdot\rho_p(\vartheta)^2$.

(b) Fix $Q$; the classes $h$, the numbers $e_h$ (which divide $\mathsf N_1$, as in Lemma~\ref{po:naut})
and the ring $\pcR=\pcR_{\mathfrak m_\vartheta}$ depend only on $Q$. Write $\chi_\vartheta\circ\det=\sum_{y\in
\auGL_2(\pcR)}\chi_\vartheta(\det y)\delta_y$ and $\delta_y(x)=\delta_1(\tilde y^{-1}x)$ for an integral lift
$\tilde y\in K_{I_2}$ of $y$; by {\cite[Lemma~2.9]{P53}}, $\theta_{\lambda,\lambda}(1,h;P_c,\delta_y)=
\lambda(\det\tilde y)^{-1}\theta_{\lambda,\lambda}(1,h\tilde y;P_c,\delta_1)$, and $\lambda(\det\tilde y)=
\chi_\vartheta(\det y)$ on the tame part. Here $\lambda=\lambda_\vartheta$ enters the kernel of $\delta_1$
only through the scalar $\lambda_\vartheta(\det h')$ on finite points $h'$ (the datum $\delta_1$ is
independent of $\vartheta$), and $\lambda_\vartheta=\vartheta\cdot\lambda_1$. Hence
$\Lambda(\vartheta)=\sum_{h,y}e_h^{-1}\vartheta(\det(h\tilde y))\,\theta^0(h\tilde y)$ with
$\theta^0=\theta_{\lambda_1,\lambda_1}(1,\cdot;P_c,\delta_1)$ independent of $\vartheta$ of modulus $Q$. So
$\mathsf N_1n_1\Lambda(\vartheta)=\sum_ac_a\vartheta(x_a)$ with $x_a=\det(h\tilde y)\in\auU(1)(\auA_f)$ and
$c_a=\mathsf N_1e_h^{-1}n_1\theta^0(h\tilde y)$, integral at every $\ell\nmid\auNm Q$ by kernel integrality
(the scalar kernel $\theta^0$ is a jet of degree $2\sum_\tau(c(\tau)-1)$ of $\Theta_f$ for a standard
datum, {\cite[Proposition~5.30]{P49}}, {\cite[Corollary~7.16]{P49}}, Lemma~\ref{pex:dint}). Then $|\sum_ac_a\vartheta
(x_a)|^2=\sum_{a,b}c_a\aubar{c_b}\vartheta(x_ax_b^{-1})$, and $\aubar{c_b}$ is integral at $\jmath$ when
$c_b$ is integral at $\jmath\circ c$.

(c) Let $\vartheta'$ be $\vartheta$ with modulus $Q\cup\{v\}$. At $v$ the datum of $\vartheta$ is the
vacuum $\plone$ and that of $\vartheta'$ is $\plone[x\text{ integral},\det x\not\equiv0]$ ($\chi_{\vartheta'}$
is trivial at $v$). By {\cite[Lemma~6.8]{P49}} (with $n=n'=2$, in which $T'_i\plone$ is computed from
{\cite[(62)]{P49}} for $h$ in $K'd_iK'$, and $\alpha$ the scalar of $d_1$ on the vacuum),
$\sum_{i}(-X)^iq_v^{\binom i2}T'_i\plone=\prod_{l<2-\mathrm{rk}(x\bmod v)}(1-q_v^l\alpha X)\plone$; at
$X=\alpha^{-1}$ this is the datum of $\vartheta'$ at $v$, by the same argument as
{\cite[(65)]{P49}}'s use there ($\prod_{l<r'}(1-q^l)$ is $1$ for $r'=0$ and $0$ otherwise). The operators
$T'_i$ are sums of translates $\omega(1,h')$, $h'\in\auU(I_2)(E^+_v)$, and the class sum is the integral over
$[\auU(I_2)]$ of a right $K_{I_2}$-invariant function, hence invariant under right translation; so
$\Lambda(\vartheta')=\sum_i(-\alpha^{-1})^iq_v^{\binom i2}\#\{K'd_iK'/K'\}\Lambda(\vartheta)=(1-\alpha^{-1})
(1-q_v\alpha^{-1})\Lambda(\vartheta)$ by {\cite[(65)]{P49}} for $V=\auF_v^2$. By {\cite[(62)]{P49}},
$\alpha^{-1}=q_v^{-1}a$ with $a$ the value of the base change of $(\nu\vartheta)^{-1}$ at $\varpi_v$ (the
splitting part cancels), $|a|=1$. Hence $|\Lambda(\vartheta')|^2=|1-aq_v^{-1}|^2|1-a|^2\,|\Lambda(\vartheta)|^2=\prod_{e\in\{0,1\}}
\prod_\pm(1-a^{\pm1}q_v^{-e})\,|\Lambda(\vartheta)|^2$, so $F(\vartheta')=F(\vartheta)\prod_{e\in\{0,1\}}\prod_\pm(1-a^{\pm1}q_v^{-e})$, symmetric in $a\leftrightarrow a^{-1}=\alpha_v$.
\end{proof}




\providecommand{\auO}{\mathcal O}
\providecommand{\auU}{\mathrm U}
\providecommand{\auGL}{\mathrm{GL}}
\providecommand{\auHom}{\operatorname{Hom}}
\providecommand{\auNm}{\operatorname{Nm}}
\providecommand{\auA}{\mathbb A}
\providecommand{\auF}{\mathbb F}
\providecommand{\auPl}{\Phi}
\providecommand{\aubar}[1]{\overline{#1}}
\providecommand{\plone}{\mathbf 1}
\providecommand{\pop}{\mathsf p}
\providecommand{\pcR}{R}
\providecommand{\pcm}{\mathfrak m}
\providecommand{\pcMat}{\mathrm M}
\providecommand{\psvac}{{\phi^{\mathrm{vac}}}}
\providecommand{\psZ}{Z}
\providecommand{\psrho}{\varrho}
\providecommand{\psHS}{\mathcal H}
\providecommand{\poA}{\mathcal A}
\providecommand{\pog}{\mathsf g}
\providecommand{\pteSD}{\mathrm{SD}}
\providecommand{\pteMP}{\mathrm{MP}}
\providecommand{\pteCS}{\mathop{\textstyle\sum^{\mathrm{cl}}}}
\providecommand{\ptetr}{\operatorname{tr}}

\subsection{Three exchanges}\label{pte:sec}

This section proves Proposition~\ref{po:exchange}. The proof has four parts: the exchange
identity with explicit determinant data at the infinite places (Corollary~\ref{pte:explicit}, a
consequence of {\cite[Theorem~3.2]{P53}}); the exchange summed over the classes of a definite group
(Lemma~\ref{pte:one}); the first two exchanges, which write the carrier as a class sum of twisted
kernels (Proposition~\ref{pte:carrier}); and the third exchange together with the identification of
its left side with the theta series $\Xi_{\vartheta,i}$ of \S\ref{po:ss-glue}
(Propositions~\ref{pte:third} and~\ref{pte:leftside}). In the course of the last step the first
factors $\varphi''_{\vartheta,i}$ of \S\ref{po:ss-glue} are defined explicitly (Definition~\ref{pte:def-ff},
Definition~\ref{ptf:def-ff}).

\subsubsection{Conventions}\label{pte:ss-conv}

We use the convention of {\cite[\S3.1.1]{P53}}: for spaces $S,T$ and characters $\alpha,\beta$
of $\auU(1)$, $\theta^{S,T}_{\alpha,\beta}(g,h;f)=\theta(\omega_{\alpha,\beta}(g,h)f)$ with
$\omega_{\alpha,\beta}(g,h)=\alpha(\det g)^{-1}\beta(\det h)\omega_{\auHom(S,T)}(g,h)$, $g\in
\auU(S)(\auA)$, $h\in\auU(T)(\auA)$; data are written $[P,\Phi]$ with $P$ a polynomial (its Fock
vector in the model of (N1)--(N5) of {\cite[\S6.1]{P49}}, holomorphic in the variables of positive
product sign and antiholomorphic in the others, as in \eqref{po:eq-poly}) and $\Phi$ a function on
matrices modulo a modulus $\pcm$ ({\cite[\S2.3]{P53}}; matrices in the frames). By
{\cite[(14)]{P53}},
\begin{equation}\label{pte:eq-twist}
\theta^{S,T}_{\alpha\lambda,\beta\lambda'}(g,h;f)=\lambda(\det g)^{-1}\lambda'(\det h)\,\theta^{S,T}_{\alpha,\beta}(g,h;f).
\end{equation}
For a positive definite standard space $I$ (frame orthonormal at the infinite places, lattice the
standard unimodular one) and a function $F$ on $\auU(I)(\auA_f)$, left invariant under
$\Gamma_I=\auU(I)(E^+)$ and right invariant under an open compact $J$, the \emph{class sum} is
\[
\pteCS_J F=\sum_{h\in\Gamma_I\backslash\auU(I)(\auA_f)/J}\frac{F(h)}{\#(\Gamma_I\cap hJh^{-1})}\cdot
\frac{1}{[K_I:J]}\qquad(J\subseteq K_I),
\]
$K_I$ the stabiliser of the standard lattice; with the Haar measure giving $K_I$ and
$\auU(I)(E^+\otimes\R)$ volume $1$ it is $\int_{\Gamma_I\backslash\auU(I)(\auA_f)}F$, so it does not
depend on $J$ (for $J=K_I$ the last factor is $1$). If $\Phi$ is a kernel datum, $h\mapsto
\theta^{S,I}_{\alpha,\beta}(g,(1,h);f)$ is written $\theta^{S,I}_{\alpha,\beta}(g,h;f)$
($h\in\auU(I)(\auA_f)$, archimedean component $1$).

\emph{Index conventions for determinant data.} Let $W$ be a space of rank $n$ whose basis is split by
sign as $m$ positive and $l$ negative vectors, and $V'$ of rank $b$ split as $p$ positive and $l$
negative vectors (the same number $l$), and let $I_a$ be positive definite of rank $a$, with
$n=b+a$; so $m=p+a$. Fix a bijection of the $m$ positive indices of $W$ with $\{1,\dots,p\}\sqcup
\{1,\dots,a\}$. At each $\tau\in\auPl$ let $\iota(\tau)$ be a finite set and $\eta_{\tau k}$
($k\in\iota(\tau)$) complex $p\times m$ matrices. On $\auHom(W,I_a)\otimes_\tau\R=\pcMat_{a\times n}
(\C)$ write $x_\tau=(x^+_\tau\mid x^-_\tau)$ (positive and negative columns), and on
$\auHom(V',W)\otimes_\tau\R=\pcMat_{n\times b}(\C)$ write $y_\tau=\bigl(\begin{smallmatrix}
y^{++}_\tau&y^{+-}_\tau\\y^{-+}_\tau&y^{--}_\tau\end{smallmatrix}\bigr)$ (rows of $W$ by sign,
columns of $V'$ by sign; $y^{++}_\tau$ is $m\times p$, $y^{--}_\tau$ is $l\times l$). Put
\begin{equation}\label{pte:eq-SDMP}
\pteSD(\eta)=\prod_\tau\prod_{k\in\iota(\tau)}\det\begin{pmatrix}\eta_{\tau k}\\x^+_\tau
\end{pmatrix},\qquad
\pteMP(\eta)=\prod_\tau\det(y^{--}_\tau)^{\#\iota(\tau)+n-2l}\prod_{k\in\iota(\tau)}\det(\eta_{\tau
k}\,y^{++}_\tau),
\end{equation}
the first a polynomial on $\auHom(W,I_a)\otimes\R$ (the $(p+a)\times m$ matrix has its columns in the
fixed order), the second on $\auHom(V',W)\otimes\R$. The \emph{standard frame} is
$\eta^0=(1_p\ 0)$, i.e.\ $\eta^0_{\tau k}$ has as rows the basis vectors of the first $p$ positive
indices; then $\pteSD(\eta^0)=\prod_\tau\det(x^{\flat}_\tau)^{\#\iota(\tau)}$, where $x^\flat_\tau$ is the
$a\times a$ block of $x^+_\tau$ in the last $a$ positive columns.

\emph{Matched functions} are those of {\cite[Definition~3.1]{P53}}, for the decomposition
$n=b+a$ of the rows of $\auHom(V',W)$ and columns of $\auHom(W,I_a)$ fixed by a bijection $\sigma$.
For a character $\alpha$ tame for $\pcm$ ({\cite[\S3.1.1]{P53}}, but with weight allowed to be any
integer vector) $\chi_\alpha$ denotes the character of $\pcR_\pcm^\times$ given by its tame parts.

\subsubsection{Three lemmas on polynomials}\label{pte:ss-poly}

\begin{lemma}[Averaging over a compact group]\label{pte:average}
Let $a\ge1$, $c=(c_\tau)\in\Z_{\ge0}^{\auPl}$ and $q$ a polynomial in the entries of square matrices
$z=(z_\tau)_\tau$, $z_\tau\in\pcMat_a(\C)$. There is $\kappa_q\in\C$ such that for every linear form
$M$ on polynomials and every $t\in\C$: if $M\bigl(q(u^{\mathsf t}z)\bigr)=t\prod_\tau\det(u_\tau)^{c_\tau}$
for all $u\in\prod_\tau\auU(a)$ (where $(u^{\mathsf t}z)_\tau=u_\tau^{\mathsf t}z_\tau$), then
$t=\kappa_q\,M\bigl(\prod_\tau\det(z_\tau)^{c_\tau}\bigr)$.
\end{lemma}
\begin{proof}
Put $\chi(u)=\prod_\tau\det(u_\tau)^{c_\tau}$ and $\bar q(z)=\int\chi(u)^{-1}q(u^{\mathsf t}z)\,du$
(probability Haar measure on the compact group $U=\prod_\tau\auU(a)$). The integrand lies in the
finite-dimensional space of polynomials of degree $\le\deg q$, so $\bar q$ is a polynomial and
$M(\bar q)=\int\chi(u)^{-1}M(q(u^{\mathsf t}z))\,du=t$. By invariance of $du$, $\bar q(v^{\mathsf t}z)=
\chi(v)\bar q(z)$ for $v\in U$. Both sides are polynomials in the entries of $v$ and $\det(v_\tau)^{-1}$, and $U$ is
Zariski dense in $\prod_\tau\auGL_a(\C)$ (its Lie algebra spans $\mathfrak{gl}_a$ over $\C$), so the
identity holds for all invertible $v$; with $z=1$, $\bar q(v^{\mathsf t})=\chi(v)\bar q(1)$, i.e.\ $\bar q=
\bar q(1)\prod_\tau\det(z_\tau)^{c_\tau}$ on a dense set, hence everywhere. So $t=M(\bar q)=\bar q(1)
M(\prod_\tau\det(z_\tau)^{c_\tau})$, and $\kappa_q=\bar q(1)$ depends only on $q$ and $c$.
\end{proof}

\begin{lemma}[Polarisation]\label{pte:polar}
Let $p<m$ and, for each $\tau$, $\eta_{\tau k}\in\pcMat_{p\times m}(\C)$ ($k\in\iota(\tau)$). There are
finitely many $u\in\prod_\tau\auU(m)$ and numbers $c_u$ such that, in the polynomial ring in the entries
of matrices $Y_\tau\in\pcMat_{m\times p}$,
\[
\prod_\tau\prod_k\det(\eta_{\tau k}Y_\tau)=\sum_uc_u\prod_\tau\prod_k\det(\eta^0u_\tau^{\mathsf t}Y_\tau).
\]
Consequently the same identity holds after substituting for the $Y_\tau$ elements of any commutative
$\C$-algebra.
\end{lemma}
\begin{proof}
For fixed $\tau$ and $c=\#\iota(\tau)$, by the Cauchy--Binet formula $\det(\eta_kY)=\sum_I
\mathrm{Pl}_I(\eta_k)\det(Y_I)$, the sum over the $p$-subsets $I$ of rows, $\mathrm{Pl}_I$ the
Pl\"ucker coordinates. So $F_\eta(Y):=\prod_k\det(\eta_kY)$ lies in the span $\mathcal S_c$ of the
products of $c$ maximal minors of $Y$, which is stable under $G=\auGL_m(\C)$ acting by $(gF)(Y)=
F(g^{\mathsf t}Y)$ (note $gF_\eta=F_{\eta g^{\mathsf t}}$). The space $\mathcal S_c$ is the degree-$c$
part of the homogeneous coordinate ring of the Grassmannian of $p$-planes in $\C^m$, which is an
irreducible $G$-module (the dual of the irreducible module of highest weight $c\varpi_p$; Borel--Weil,
see e.g.\ Fulton, \emph{Young tableaux}). For several $\tau$ the space
$\bigotimes_\tau\mathcal S_{c_\tau}$ is irreducible for $\prod_\tau G$. It contains
$F_{\eta^0}=\prod_\tau\det(Y^{\mathrm{top}}_\tau)^{c_\tau}\ne0$, so it is the span of the
$\prod_\tau G$-orbit of $F_{\eta^0}$. The map $g\mapsto gF_{\eta^0}$ is polynomial in the entries of
$g$, and $\prod_\tau\auU(m)$ is Zariski dense in $\prod_\tau G$, so every linear form vanishing on the
unitary orbit vanishes on the whole orbit: the unitary translates $F_{\eta^0u^{\mathsf t}}$ span the
same finite-dimensional space, which contains $F_\eta$.
\end{proof}

\begin{lemma}\label{pte:stackeval}
Let $Z\in\pcMat_{a\times m}(\C)$, $m=p+a$. There are $Y\in\pcMat_{m\times p}(\C)$ and $d\in\C$ with
$\det\binom{\xi}{Z}=d\det(\xi Y)$ for all $\xi\in\pcMat_{p\times m}(\C)$ (rows of $\binom\xi Z$ in the
fixed order of columns).
\end{lemma}
\begin{proof}
Choose a unitary $U=(U_1\mid U_2)\in\auU(m)$, $U_1$ of size $m\times p$, with $ZU_1=0$ (possible when
$\operatorname{rk}Z\le a$: take the columns of $U_1$ orthonormal in $\ker Z$, which has dimension
$\ge p$). Then $\binom\xi ZU=\bigl(\begin{smallmatrix}\xi U_1&\xi U_2\\0&ZU_2\end{smallmatrix}\bigr)$, so
$\det\binom\xi Z=\det(\xi U_1)\det(ZU_2)/\det U$; take $Y=U_1$, $d=\det(ZU_2)/\det U$.
\end{proof}

\subsubsection{The exchange identity with determinant data}\label{pte:ss-explicit}

\begin{corollary}[Explicit exchange]\label{pte:explicit}
Let $W$, $V'$, $I_a$, $n=b+a$, $m=p+a$, $l$, $\iota$ be as in \S\ref{pte:ss-conv}, with $a\ge1$,
$a=b+2d$ ($d\ge0$) and $2l+2\le n$. There is a real number $C_{\mathrm e}\ge0$ such that for every
modulus $\pcm\ne(1)$ and every character $\alpha$ of $\auU(1)$ tame for $\pcm$ (tame parts
non-trivial exactly at the primes of $\pcm$) of weight $(l-\#\iota(\tau))_\tau$ there is $c\in\C$
with
\[
|c|^2\,L_*\bigl(\pcm,\alpha^{-1}\upsilon^{b+d},d\bigr)=C_{\mathrm e},
\]
such that for all $\eta$, all functions $\Phi$, $\Psi$ on $\pcMat_{a\times n}(\pcR_\pcm)$,
$\pcMat_{n\times b}(\pcR_\pcm)$ matched for $\chi_\alpha^{-1}$, all characters $\beta$ and all $g\in
\auU(W)(\auA)$,
\begin{equation}\label{pte:eq-explicit}
c\,\pteCS_{K_{I_a}}\Bigl(h\mapsto\theta^{W,I_a}_{\beta,\alpha}\bigl(g,h;[\pteSD(\eta),\Phi]\bigr)\Bigr)
=\int_{[\auU(V')]}\theta^{V',W}_{\alpha,\alpha\beta^{-1}}\bigl(x,g;[\pteMP(\eta),\Psi]\bigr)\,dx .
\end{equation}
Both sides converge absolutely; $C_{\mathrm e}$ does not depend on $\pcm,\alpha,\beta,\eta,\Phi,\Psi,g$.
\end{corollary}

Unlike {\cite[Theorem~3.2]{P53}}, the corollary does not assert $c\ne0$ (that is, $C_{\mathrm e}>0$): the
explicit polynomial $\pteSD(\eta)$ replaces the vector $x$ of the theorem, and the number $\kappa$
relating them may vanish. In the application, $C_{\mathrm e}>0$ follows a posteriori
(Proposition~\ref{po:glue}(G2)).

\emph{Matching with {\cite[Theorem~3.2]{P53}}.} Put $s_0=n/2-l\ge1$, $\nu=\alpha^{-1}\upsilon^{b+d}$ and
$\gamma=\nu^{-1}\upsilon^{n/2}$. Since $n/2=b+d$, $\gamma=\alpha$. The weight of $\nu$ is
$\mu_\tau=\#\iota(\tau)-l+(b+d)=\#\iota(\tau)+s_0\ge1$, its tame parts are those of $\alpha^{-1}$, so $\nu$ is
tame for $\pcm$ of weight $\mu$ and the character $\chi$ of {\cite[Definition~3.1]{P53}} attached
to $\nu$ is $\chi_\alpha^{-1}$. The hypotheses of {\cite[Theorem~3.2]{P53}} for $(W,I_a,V')$ hold: $a=b+2d$;
$q=l=0+l=q_{I_a}+q_{V'}$; $2q+2\le n$; $n=2l+2s_0\le2\mu_\tau+2l$.

\begin{proof}
\emph{Step 0: the vector $v$.} Let $Q_0=\pteMP(\eta^0)$ (for the given $\iota$), a polynomial Fock
vector of $\auHom(V',W)\otimes\R$, which has finite moments. {\cite[Theorem~3.2]{P53}} for $(W,V,V')=(W,I_a,V')$,
$\mu$ as above and $v=Q_0$ gives $C>0$ and a Fock vector $x_0$ of $\auHom(W,I_a)\otimes\R$, and for each $\alpha$
a number $c_0$ with $|c_0|^2=C\,L_*(\pcm,\nu,d)\ne0$, such that (taking $\alpha_X=\beta$, $\alpha'=\alpha
\beta^{-1}$ in {\cite[(12)]{P53}}, so $\alpha_X\alpha'=\gamma$)
\begin{equation}\label{pte:eq-X0}
\int_{[\auU(I_a)]}\theta^{W,I_a}_{\beta,\alpha}(g,h;[x_0,\Phi])\,dh=c_0\int_{[\auU(V')]}
\theta^{V',W}_{\alpha,\alpha\beta^{-1}}(x,g;[Q_0,\Psi])\,dx
\end{equation}
for all matched $\Phi,\Psi$ and all $g$. By \eqref{pte:eq-twist} both sides of \eqref{pte:eq-X0} and of
\eqref{pte:eq-explicit} change by the same factor $\beta(\det g)^{-1}$ when $\beta$ is changed, so we
may fix $\beta$; we keep it general in the notation.

\emph{Step 1: the two functionals.} Fix $\alpha,\Phi,\Psi$ and $g$. For a Fock vector $w$ of
$\auHom(W,I_a)\otimes\R$ with finite moments and $y\in\auU(W)(E^+\otimes\R)$ put $L_y(w)=\int_{[\auU(I_a)]}
\theta^{W,I_a}_{\beta,\alpha}(gy,h;[w,\Phi])\,dh$; $[\auU(I_a)]$ is compact and the kernel depends
continuously on $w$, uniformly on compacta, so $L_y$ is a continuous linear form. For a polynomial
$Q$ on $\auHom(V',W)\otimes\R$ put $\ell'(Q)=\int_{[\auU(V')]}\theta^{V',W}_{\alpha,\alpha\beta^{-1}}
(x,g;[Q,\Psi])\,dx$, which converges absolutely ({\cite[Corollary~3.34]{P53}}: $\Psi$ is supported on
left-invertible matrices) and is linear in $Q$. Write $\omega_\infty$ for the archimedean Weil
representations (with the archimedean components of the characters). Then:
\begin{enumerate}
\item[(a)] $L_y(\omega_\infty(k,u)w)=L_{yk}(w)$ for $k\in\auU(W)_\infty$, $u\in\auU(I_a)_\infty$: the
action of $k$ moves the point $gy$, and the integral over $h$ absorbs $u$ (right invariance of $dh$);
\item[(b)] if $\omega_\infty(1,y)$, acting on $\auHom(V',W)\otimes\R$ through the second group,
carries $Q_0$ to $Q$, then $L_y(x_0)=c_0\ell'(Q)$: this is \eqref{pte:eq-X0} at the point $gy$.
\end{enumerate}

\emph{Step 2: action on the two polynomials.} By Step~0 we may take the archimedean weight of $\beta$ to be $0$.
Then, by Corollary~\ref{ptf:SDlaw}(a),(b) (the transpose conventions are fixed in \S\ref{ptf:ss-fock}), for $k$ in the
positive block $\prod_\tau\auU(m)$ of $\auU(W)_\infty$ and $u\in\auU(I_a)_\infty$,
\begin{equation}\label{pte:eq-SDlaw}
\omega_\infty(k,u)\,\pteSD(\eta k^{-\mathsf t})=\prod_\tau\det(k_\tau)^{-\#\iota(\tau)}\,\pteSD(\eta),
\qquad
\omega_\infty(1,k)\,\pteMP(\eta)=\prod_\tau\det(k_\tau)^{-\#\iota(\tau)}\,\pteMP(\eta k^{\mathsf t}),
\end{equation}
with $(\eta k^{\pm\mathsf t})_{\tau j}=\eta_{\tau j}k_\tau^{\pm\mathsf t}$, and $\omega_\infty(1,u)\pteSD(\eta)=\pteSD(\eta)$ for all
$u\in\auU(I_a)_\infty$.

\emph{Step 3: a circle mean.} For $z\in\C$, $|z|=1$, let $t_z\in\prod_\tau\auU(m)$ be diagonal with
entries $z^{-2}$ at the first $p$ positive indices and $z^{-1}$ at the last $a$, and put
$\Omega_z(w)=z^{-aN_\iota}\omega_\infty(t_z,z^{-1}1_a)w$, $N_\iota=\sum_\tau\#\iota(\tau)$. We claim:
there is $\kappa\in\C$ (depending on $x_0$) such that for every continuous linear form $\ell$ on the
Fock space with $\ell\circ\omega_\infty(1,u)=\ell$ for all $u\in\auU(I_a)_\infty$,
\begin{equation}\label{pte:eq-circle}
\kappa\,\ell\bigl(\pteSD(\eta^0)\bigr)=\frac1{2\pi}\int_0^{2\pi}\ell\bigl(\Omega_{e^{i\vartheta}}(x_0)
\bigr)\,d\vartheta .
\end{equation}
The right side is continuous in $x_0$ (uniformly over $\ell$ in bounded sets) and the set of
functions $\ell\mapsto\kappa\,\ell(\pteSD(\eta^0))$, $\kappa\in\C$, is a closed line; polynomial vectors are
dense in the Fock space; so it suffices to prove the claim for $x_0$ a polynomial $P$ in the
entries of $x$. By Corollary~\ref{ptf:SDlaw}(c), $\Omega_z$ acts on polynomials by the substitution that multiplies every
variable outside the block $x^\flat$ (the last $a$ positive columns) by $z$ and fixes the entries of
$x^\flat$ (the scalar $z^{-aN_\iota}$ cancels the vacuum character). Hence $\ell(\Omega_z P)=F(z)$ for
a polynomial $F$, and its circle mean is $F(0)=\ell(P_0)$, where $P_0$ is $P$ with all variables outside
$x^\flat$ set to $0$. Now $P_0=q(x^\flat)$ is a polynomial in square matrices, and $M(r)=\ell(r(x^\flat))$
is a linear form with $M(q(u^{\mathsf t}x^\flat))$ computed as follows: by
Corollary~\ref{ptf:SDlaw}(c), $\omega_\infty(1,u)$ acts on polynomials in $x^\flat$ by $x^\flat\mapsto u^{\mathsf t}x^\flat$
times $\prod_\tau\det(u_\tau)^{-\#\iota(\tau)}$, and $\ell$ is
invariant, so $M(q(u^{\mathsf t}x^\flat))=\ell(P_0)\prod_\tau\det(u_\tau)^{\#\iota(\tau)}$. By
Lemma~\ref{pte:average}, $\ell(P_0)=\kappa_q\,\ell(\prod_\tau\det(x^\flat_\tau)^{\#\iota(\tau)})=\kappa_q
\,\ell(\pteSD(\eta^0))$, and $\kappa=\kappa_q$ depends only on $P$.

\emph{Step 4: unitary frames.} Let $k\in\prod_\tau\auU(m)$ and $\xi=\eta^0k^{\mathsf t}$. We show
\begin{equation}\label{pte:eq-unitary}
\kappa\,L_1(\pteSD(\xi))=c_0\,\ell'(\pteMP(\xi)).
\end{equation}
By \eqref{pte:eq-SDlaw} with $\eta=\eta^0$ (and $k^{-1}$ in place of $k$),
$\kappa L_1(\pteSD(\xi))=\kappa\prod_\tau\det(k_\tau)^{\#\iota(\tau)}L_{k}(\pteSD(\eta^0))$ by (a). The
form $\ell=\prod\det(k_\tau)^{\#\iota(\tau)}L_k$ is invariant under $\auU(I_a)_\infty$ by (a), so by
\eqref{pte:eq-circle} this equals the circle mean of $\prod\det(k_\tau)^{\#\iota(\tau)}L_k(\Omega_z x_0)$.
By (a), $L_k(\Omega_zx_0)=z^{-aN_\iota}L_{kt_z}(x_0)$, and by (b) with $y=kt_z$ and \eqref{pte:eq-SDlaw}
(which gives $\omega_\infty(1,kt_z)Q_0=\prod\det(kt_z)_\tau^{-\#\iota(\tau)}\pteMP(\eta^0(kt_z)^{\mathsf t})$ and
$\eta^0(kt_z)^{\mathsf t}=z^{-2}\xi$, because $t_z$ is $z^{-2}$ on the first $p$ indices),
\[
L_{kt_z}(x_0)=c_0\prod_\tau\det(k_\tau t_{z,\tau})^{-\#\iota(\tau)}\ell'(\pteMP(z^{-2}\xi))
=c_0\prod_\tau\det(k_\tau)^{-\#\iota(\tau)}z^{(2p+a)N_\iota}z^{-2pN_\iota}\ell'(\pteMP(\xi)),
\]
using $\det t_{z,\tau}=z^{-(2p+a)}$ and the homogeneity $\pteMP(w\xi)=w^{pN_\iota}\pteMP(\xi)$. All powers
of $z$ and of $\det k_\tau$ cancel: the integrand of the circle mean is the constant $c_0\ell'(\pteMP(\xi))$.

\emph{Step 5: all $\eta$.} By Lemma~\ref{pte:polar} ($p<m$ since $a\ge1$) there are $u$'s and $c_u$
with $\prod\det(\eta_{\tau k}Y_\tau)=\sum c_u\prod\det(\eta^0u_\tau^{\mathsf t}Y_\tau)$. Substituting
$Y_\tau=y^{++}_\tau$ gives $\pteMP(\eta)=\sum_uc_u\pteMP(\eta^0u^{\mathsf t})$. Evaluating at a point
$x$, Lemma~\ref{pte:stackeval} (applied to $Z=x^+_\tau$ at each $\tau$) gives $Y_\tau$ and $d_\tau$
with $\det\binom{\zeta}{x^+_\tau}=d_\tau\det(\zeta Y_\tau)$ for all $\zeta$; so $\pteSD(\eta)(x)=
\prod_\tau d_\tau^{\#\iota(\tau)}\prod\det(\eta_{\tau k}Y_\tau)$ and the same identity with the same
$c_u$ holds for $\pteSD$ at every point: $\pteSD(\eta)=\sum_uc_u\pteSD(\eta^0u^{\mathsf t})$. By
\eqref{pte:eq-unitary} and linearity, $\kappa L_1(\pteSD(\eta))=c_0\ell'(\pteMP(\eta))$ for all $\eta$.

\emph{Step 6: class sum.} By Step 2 the integrand $h\mapsto\theta^{W,I_a}_{\beta,\alpha}(g,h;
[\pteSD(\eta),\Phi])$ is invariant under right translation by $\auU(I_a)_\infty$. It is invariant under
right translation by $K_{I_a}$: by {\cite[Lemma~2.9]{P53}}, translation by $k\in K_{I_a}$ replaces $\Phi$
by $\alpha(\det k)\Phi(\bar k^{-1}\,\cdot)$, $\bar k$ the reduction modulo $\pcm$, and $\alpha(\det k)=
\chi_\alpha(\det\bar k)$ on integral points (tameness) while $\Phi(\bar k^{-1}x)=\chi_\alpha(\det\bar k)^{-1}
\Phi(x)$ because $\Phi$ is matched for $\chi_\alpha^{-1}$ ({\cite[Definition~3.1]{P53}}: $\Phi$
lives on the last rows of inverses of invertible matrices, and left multiplication by $\bar k^{-1}$
changes $\det$ by $\det\bar k^{-1}$). With $\mathrm{vol}(K_{I_a})=\mathrm{vol}(\auU(I_a)_\infty)=1$ the
integral over $[\auU(I_a)]$ is therefore the class sum $\pteCS_{K_{I_a}}$. Hence
$\kappa\,\pteCS_{K_{I_a}}(\cdots)=c_0\ell'(\pteMP(\eta))$, which is \eqref{pte:eq-explicit} with
$c=\kappa/c_0$ ($c_0\ne0$). Then $|c|^2L_*(\pcm,\nu,d)=|\kappa|^2/C=:C_{\mathrm e}$, and $\kappa$ depends
only on $x_0$, hence not on $\pcm,\alpha$. Convergence of the right side is {\cite[Corollary~3.34]{P53}}; the
left side is a finite sum.
\end{proof}

\begin{remark}\label{pte:rem-explicit}
(i) The averaging of Step~3 takes place on the definite group $\auU(I_a)_\infty$, which is compact;
the non-compact group $\auU(V')$ is never averaged over. This is why the argument applies with
$V'\ne0$ of any signature, in contrast with the isotypic shortcut of Proposition~\ref{pex:E1}. (ii)
For $b=0$ ($V'=0$, $p=l=0$, $\pteMP=1$, $\Psi=1$ and $\Phi=\chi_\alpha\circ\det$ by
{\cite[Definition~3.1]{P53}}) the right side of \eqref{pte:eq-explicit} is $\alpha(\det g)\beta(\det
g)^{-1}$ times the volume $1$ of the point $[\auU(0)]$.
\end{remark}

\subsubsection{One exchange summed over classes}\label{pte:ss-one}

\begin{lemma}[Product data are matched]\label{pte:matched}
Let $\chi$ be a character of $\pcR^\times$ ($\pcR=\pcR_\pcm$), extended by $0$, and $\Psi'$ a function on
$\pcMat_{n\times b}(\pcR)$ with $\Psi'(y\gamma)=\Psi'(y)$ for $\gamma\in\auGL_b(\pcR)$. On
$\pcMat_{n\times(b+n)}(\pcR)\ni x=(x'\mid x'')$ and $\pcMat_{(b+n)\times b}(\pcR)\ni y=\binom{y'}{y''}$
put $\Phi(x)=\chi(\det x'')\Psi'(x''^{-1}x')$ (zero if $x''$ is not invertible) and $\Psi(y)=
\chi(\det y')\Psi'(y'')$. Then $\Phi,\Psi$ are matched for $\chi^{-1}$ ({\cite[Definition~3.1]{P53}},
with $N=b+n$, the first $b$ indices first).
\end{lemma}
\begin{proof}
Let $k\in\auGL_{b+n}(\pcR)$, $k=\bigl(\begin{smallmatrix}A&B\\C&D\end{smallmatrix}\bigr)$, $k^{-1}=
\bigl(\begin{smallmatrix}A'&B'\\C'&D'\end{smallmatrix}\bigr)$; $k_{[b]}=\binom AC$, $k^{[n]}=(C'\mid
D')$. From $k^{-1}k=1$: $C'A+D'C=0$. If $D'$ is invertible, $A$ is invertible (Schur complement:
$\det k^{-1}=\det D'\det(A'-B'D'^{-1}C')$ and $A=(A'-B'D'^{-1}C')^{-1}$), $D'^{-1}C'=-CA^{-1}$ and
$\det k=\det A/\det D'$. Then $\Phi(k^{[n]})=\chi(\det D')\Psi'(-CA^{-1})=\chi(\det D')\Psi'(C)$ (invariance
under $-A^{-1}$), and $\chi^{-1}(\det k)\Psi(k_{[b]})=\chi(\det D')\chi(\det A)^{-1}\chi(\det A)\Psi'(C)$:
equal. If $D'$ is not invertible, both vanish ($A$ is then not invertible, by the same identity read
from $k$). Supports: if $\Psi(y)\ne0$ then $y'$ is invertible and $y=k_{[b]}$ for $k=\bigl(\begin{smallmatrix}
y'&0\\y''&1\end{smallmatrix}\bigr)$; if $\Phi(x)\ne0$ then $x''$ is invertible and $x=k^{[n]}$ for
$k^{-1}=\bigl(\begin{smallmatrix}1&0\\x'&x''\end{smallmatrix}\bigr)$.
\end{proof}

\begin{lemma}[One exchange, summed]\label{pte:one}
Let $V'$, $W'$ be spaces of ranks $b$, $n$, $I_b$ the positive definite standard space, $\mathsf W=
I_b\oplus W'$, $\pcm\ne(1)$, $\chi$ a character of $\pcR_\pcm^\times$, $a_1,a_2,\beta_0$ characters, $P$ a
polynomial on $\auHom(V',W')\otimes\R$ (read on $\auHom(V',\mathsf W)$ through the rows of $W'$) and
$g\in\auU(W')(\auA)$. Suppose: (i) a function $F$ on $\auU(V')(\auA)$ satisfies $c_0\,\pteCS_{K_{I_b}}
(h'\mapsto\mathrm a(h')\theta^{V',I_b}_{a_1,\beta_0}(x,h';[1,\chi\circ\det]))=F(x)$ for all $x$, for some
number $c_0$ and function $\mathrm a$; (ii) for functions $K(\Phi,h')$ on $\auU(I_n)(\auA_f)$ and a
number $c$: for all $\Phi,\Psi$ matched for $\chi^{-1}$ as in Lemma~\ref{pte:matched} and all
$h'\in\auU(I_b)(\auA_f)$, $c\,\pteCS(K(\Phi,h'))=\int_{[\auU(V')]}\theta^{V',\mathsf W}_{a_1a_2,\beta_0}(x,
(h',g);[P,\Psi])\,dx$. Let $\Psi'$ be invariant under $\auGL_b(\pcR)$ on the right and $\Phi_{\Psi'}$ as
in Lemma~\ref{pte:matched}. Then
\[
c\,c_0\,\pteCS_{K_{I_b}}\bigl(h'\mapsto\mathrm a(h')\,\pteCS(K(\Phi_{\Psi'},h'))\bigr)=\int_{[\auU(V')]}
F(x)\,\theta^{V',W'}_{a_2,\beta_0}(x,g;[P,\Psi'])\,dx,
\]
the integrals converging absolutely.
\end{lemma}
\begin{proof}
For $\Psi=\Psi_{\Psi'}$ of Lemma~\ref{pte:matched}, $\Psi(y)=\chi(\det y')\Psi'(y'')$ is a product over
the two blocks of rows, and the polynomial involves only the rows of $W'$; by the compatibility of
kernels with orthogonal sums in the target (the Weil representation of $\auHom(V',I_b\oplus W')=
\auHom(V',I_b)\oplus\auHom(V',W')$ is the tensor product of the two Weil representations, with the
splittings multiplying, as in (W7) of \S\ref{po:ss-glue}), together with the convention
$\omega_{\alpha,\beta}=\alpha(\det x)^{-1}\beta(\det h)\,\omega$: the source group $\auU(V')$ is common to
both factors, so its character splits as $a_1a_2$, while $\det(h',g)=\det h'\det g$, so the target character $\beta_0$
appears on each block. Hence $\theta^{V',\mathsf W}_{a_1a_2,\beta_0}(x,(h',g);[P,\Psi])=
\theta^{V',I_b}_{a_1,\beta_0}(x,h';[1,\chi\circ\det])\,\theta^{V',W'}_{a_2,\beta_0}(x,g;[P,\Psi'])$. The left
factor's datum is supported on invertible matrices, so the product is integrable
({\cite[Corollary~3.34]{P53}}); by (ii) its integral is $c\,\pteCS(K(\Phi_{\Psi'},h'))$. Multiply by
$c_0\mathrm a(h')$ and sum over the finitely many classes $h'$; the finite sum commutes with the
integral, and by (i) $c_0\pteCS(\mathrm a\,\theta^{V',I_b}(x,\cdot))=F(x)$.
\end{proof}

\begin{lemma}[Translates of a datum]\label{pte:translates}
Let $F$ on $\auU(V')(\auA)$ be automorphic, rapidly decreasing and right invariant under the
stabiliser $K_{V'}$ of the lattice, $\beta'$ a character with trivial tame part, $\psi$ a function on
$\pcMat_{n\times b}(\pcR_\pcm)$, and $\pcm'\subseteq\pcm$. Then
\[
\int_{[\auU(V')]}F(x)\,\theta^{V',W'}_{\beta',a}\Bigl(x,g;\bigl[P,\textstyle\sum_{\gamma\in\auGL_b(\pcR_\pcm)}\psi(\,\cdot\,\gamma)
\bigr]\Bigr)dx=\#\auGL_b(\pcR_\pcm)\int_{[\auU(V')]}F(x)\,\theta^{V',W'}_{\beta',a}(x,g;[P,\psi])\,dx,
\]
the datum on the left read modulo $\pcm'$.
\end{lemma}
\begin{proof}
For an integral lift $\tilde\gamma\in K_{V'}$ of $\gamma$, by {\cite[Lemma~2.9]{P53}} right translation
of the kernel by $\tilde\gamma$ replaces $\psi$ by $\beta'(\det\tilde\gamma)\psi(\,\cdot\,\gamma)=
\psi(\,\cdot\,\gamma)$ ($\beta'$ has trivial tame part, so it is trivial on integral points). As $F$ is
$K_{V'}$-invariant and the measure on $[\auU(V')]$ is right invariant, each translate gives the same
integral; all integrals converge since $F$ is rapidly decreasing and kernels are of moderate growth.
Data modulo $\pcm$ and their pull-backs modulo $\pcm'$ define the same datum ({\cite[Lemma~3.24]{P53}}).
\end{proof}

\subsubsection{The carrier}\label{pte:ss-carrier}

Let $m$ be even, $I_m$ the positive definite standard space of rank $m$, $I_m^-$ its negative, and $S$
the space of signature $(2m,m)$ with the standard frame (signs $+1$ on $2m$ indices, $-1$ on $m$). The
\emph{carrier} is
\[
\varphi^\circ(x)=\int_{[\auU(I_m^-)]}\theta^{I_m^-,S}_{1,\upsilon^m}\Bigl(y,x;\Bigl[\prod_\tau\det(y^-_\tau)^{2m},\plone\Bigr]\Bigr)\,dy
\qquad(x\in\auU(S)(\auA)),
\]
the theta lift of the constant $1$, with modulus $(1)$ (vacuum at all finite places) and $y^-_\tau$ the
$m\times m$ block of variables in the negative rows of $S$. For a character $c$ of $\auU(1)$ put
$\mathrm a_c(h)=\pteCS_{K_{I_m}}\bigl(h_1\mapsto\theta^{I_m,I_{3m}}_{1,c}(h_1,h;[1,\plone])\bigr)$ for
$h\in\auU(I_{3m})(\auA_f)$.

\begin{lemma}[Negation]\label{pte:conj}
Let $S_1,T$ be spaces of ranks $n_1,n'$ and $S_1^-$ the space with the same frame and the negated form. There is
$\epsilon\in\C$, $|\epsilon|=1$, such that for all characters $a,b$, moduli $\pcm$, functions $\Psi$, polynomials $P$,
and $g,h$, $\aubar{\theta^{S_1,T}_{b,a}(g,h;[P,\Psi])}=\epsilon\,\theta^{S_1^-,T}_{b^{-1}\upsilon^{n'},a^{-1}\upsilon^{n_1}}
(g,h;[\aubar P,\aubar\Psi])$, where $\aubar P$ has conjugate coefficients and $\auU(S_1^-)=\auU(S_1)$ (the source
character is shifted by $\upsilon$ to the rank of $T$, the target character by $\upsilon$ to the rank of $S_1$).
\end{lemma}
\begin{proof}
This is Lemma~\ref{ptf:conj} (with $\alpha=b$, $\beta=a$, $n_S=n_1$, $n_T=n'$); its proof constructs the
antilinear intertwiner, computes the phase (Lemma~\ref{ptf:theta}) and the change of splitting
(Lemma~\ref{ptf:weil}).
\end{proof}

\begin{proposition}[First two exchanges]\label{pte:carrier}
Let $m$ be even. There is $C_{12}\ge0$ such that for every modulus $\pcm\ne(1)$ and every character $c$
of weight $0$, tame for $\pcm$ with residue character $\chi$, there is $\gamma\in\C$ with
\[
|\gamma|^2\,L_*(\pcm,c^{-1}\upsilon^m,m)\,L_*(\pcm,c\,\upsilon^{m/2},m/2)=C_{12},\qquad
\gamma\,\pteCS_{K_{I_{3m}}}\Bigl(h\mapsto\mathrm a_c(h)\,\theta^{S,I_{3m}}_{c,\upsilon^m}\bigl(x,h;[1,\chi\circ
\det]\bigr)\Bigr)=\varphi^\circ(x)
\]
for all $x\in\auU(S)(\auA)$.
\end{proposition}
\begin{proof}
\emph{First exchange.} Corollary~\ref{pte:explicit} with $W=I_m$, $V'=0$, $I_a=I_m$ ($n=a=m$, $b=p=l=0$,
$d=m/2$, $2l+2\le m$), $\iota(\tau)=\emptyset$ (so $\pteSD=\pteMP=1$), the character $c^{-1}$ (weight
$0=l-\#\iota$), $\beta=c^{-1}$ and the matched pair $\Phi=\chi^{-1}\circ\det$, $\Psi=1$ (Remark~\ref{pte:rem-explicit}(ii);
$\chi^{-1}$ is the residue character of $c^{-1}$) gives $\gamma_1$ with $|\gamma_1|^2L_*(\pcm,c\upsilon^{m/2},m/2)=
C_1$ and $\gamma_1\pteCS(\theta^{I_m,I_m}_{c^{-1},c^{-1}}(g,\cdot;[1,\chi^{-1}\circ\det]))=1$ for all $g$.
By Lemma~\ref{pte:conj} (with $S_1=T=I_m$, equal ranks), conjugating, $\aubar{\gamma_1}\epsilon\,\pteCS(\theta^{I_m^-,
I_m}_{c\upsilon^m,c\upsilon^m}(y,\cdot;[1,\chi\circ\det]))=1$ for all $y\in\auU(I_m^-)(\auA)$ (the datum
$\aubar{\chi^{-1}(\det x)}=\chi(\det x)$).

\emph{Second exchange.} Corollary~\ref{pte:explicit} with $W=I_m\oplus S$ (rank $4m$, $l=m$ negative
signs), $V'=I_m^-$ ($b=m$, $p=0$, $l=m$), $I_a=I_{3m}$ ($a=3m=b+2m$, $d=m$), $\iota(\tau)=\emptyset$, the
character $\alpha=c\upsilon^m$ (weight $m=l-0$, residue character $\chi$) gives $\gamma_2$ with
$|\gamma_2|^2L_*(\pcm,(c\upsilon^m)^{-1}\upsilon^{2m},m)=C_2$; note $(c\upsilon^m)^{-1}\upsilon^{2m}=c^{-1}\upsilon^m$.
Here $\pteSD=1$ and $\pteMP=\prod_\tau\det(y^{--}_\tau)^{2m}$, the negative block being the rows of $S$ of
sign $-1$ against $I_m^-$. Apply Lemma~\ref{pte:one} with $V'=I_m^-$, $W'=S$, $I_b=I_m$, the character
$\chi$, $\Psi'=1$ (so $\Phi_{\Psi'}(x)=\chi(\det x'')$), the characters $(a_1,a_2,\beta_0)=(c\upsilon^m,1,c\upsilon^m)$, $\mathrm a=$ the constant $1$, $F=1$ and $c_0=\aubar{\gamma_1}\epsilon$
(hypothesis (i) is the conjugated first exchange), and (ii) the second exchange (with $\beta=1$) at the
points $(h',x)$. It gives $\gamma_2c_0\pteCS_{K_{I_m}}(h'\mapsto\pteCS_{K_{I_{3m}}}(h\mapsto\theta^{I_m\oplus
S,I_{3m}}_{1,c\upsilon^m}((h',x),h;[1,\Phi_{\Psi'}])))=\int_{[\auU(I_m^-)]}\theta^{I_m^-,S}_{1,c\upsilon^m}(y,x;
[\prod\det(y^-)^{2m},\plone])\,dy$ (the right side of \eqref{pte:eq-explicit} carries $\theta^{I_m^-,I_m\oplus S}_{c\upsilon^m,c\upsilon^m}$, split as
$a_1a_2=c\upsilon^m\cdot1$ with target character $c\upsilon^m$; the factor $F=1$).
By the compatibility of kernels with orthogonal sums of the source space (W7), with the target character split
as $c\cdot\upsilon^m$, $\theta^{I_m\oplus S,I_{3m}}_{1,c\upsilon^m}((h',x),h;[1,\Phi_{\Psi'}])=\theta^{I_m,I_{3m}}_{1,c}
(h',h;[1,\plone])\,\theta^{S,I_{3m}}_{1,\upsilon^m}(x,h;[1,\chi\circ\det])$ (the datum is $1$ on the columns of
$I_m$ and $\chi\circ\det$ on those of $S$); exchanging the two finite class sums, the left side is
$\gamma_2c_0\,\pteCS(\mathrm a_c\,\theta^{S,I_{3m}}_{1,\upsilon^m}(x,\cdot;[1,\chi\circ\det]))$. The right side is
$\int\theta^{I_m^-,S}_{1,c\upsilon^m}(y,x;[\prod\det(y^-)^{2m},\plone])\,dy=c(\det x)\varphi^\circ(x)$ by
\eqref{pte:eq-twist}, and on the left $\theta^{S,I_{3m}}_{1,\upsilon^m}(x,h)=c(\det x)\theta^{S,I_{3m}}_{c,\upsilon^m}(x,h)$;
the factor $c(\det x)\ne0$ cancels. So the claim holds with $\gamma=\gamma_2c_0$, and $|\gamma|^2=|\gamma_2|^2|\gamma_1|^2$ gives $C_{12}=C_1C_2$.
\end{proof}

\begin{lemma}[Properties of the carrier]\label{pte:carrierprop}
Assume a character $c$ of weight $0$ tame for some $\pcm\ne(1)$ exists. Then $\varphi^\circ$ is a finite
linear combination of the functions $x\mapsto\theta^{S,I_{3m}}_{c,\upsilon^m}(x,h;[1,\chi\circ\det])$; it is
continuous, left invariant under $\auU(S)(E^+)$, right invariant under $K_S$, rapidly decreasing, and
holomorphic. Moreover $\varphi^\circ\ne0$.
\end{lemma}
\begin{proof}
If $\gamma=0$ in Proposition~\ref{pte:carrier} then $\varphi^\circ=0$; otherwise divide. The class sum is a
finite sum. Each kernel is right invariant under $K_S$ ({\cite[Lemma~2.9]{P53}}: $c(\det k)^{-1}$ cancels
$\chi(\det k)$), rapidly decreasing because its finite datum $\chi\circ\det$ is supported on invertible
matrices ({\cite[Proposition~3.26]{P53}}), and holomorphic since its polynomial is $1$ (the vacuum, which is
holomorphic for $\auU(S)_\infty$ in the order with the positive vectors first). The non-vanishing
of $\varphi^\circ$ is Lemma~\ref{pte:carriernz}, whose proof uses neither $c$ nor
Proposition~\ref{pte:carrier}.
\end{proof}

\begin{lemma}[The carrier is not zero]\label{pte:carriernz}
$\varphi^\circ\ne0$; indeed a Fourier coefficient of $\varphi^\circ$ along a unipotent subgroup of
$\auU(S)$ is not identically zero.
\end{lemma}
\begin{proof}
Write $V=I_m^-$ (frame $\varepsilon_1,\dots,\varepsilon_m$, $\langle\varepsilon_k,\varepsilon_k\rangle=-1$),
$\mathsf W=\auHom(V,S)$, and $f=[v,\plone]=v\otimes\psvac$ for the datum of the definition, with
$v=\prod_\tau\det(y^-_\tau)^{2m}$ and $\psvac$ the vacuum of the lattice $L=\alpha\auHom(L_V,L_S)$ at every
finite place ({\cite[\S2.3]{P53}}). Then $\varphi^\circ(x)=\int_{[\auU(V)]}\theta(\omega(y,x)f)\,dy$ with
$\omega(y,x)=\upsilon^m(\det x)\,\omega_{\auHom(V,S)}(y,x)$; the operator $\omega_{\auHom(V,S)}(y,x)$ is
$\omega_{\mathsf W}$ of the element $w\mapsto xwy^{-1}$ of $\auU(\mathsf W)$, and $\omega_{\mathsf W}$ is a
representation of $\auU(\mathsf W)(\auA)$. We use (W1)--(W4) of {\cite[Lemma~5.2]{P49}}, the local forms
(c), (d) of {\cite[Proposition~2.14]{P49}}, and the mixed model behind them: for an isotropic subspace $I$ with
a complement $I^*$ in duality with it, $\psHS_v(\mathsf W)$ is a space of functions on $I^*_v$ with values
in $\psHS_v(\mathsf W_I)$, $\psZ_{I,v}$ is evaluation at $0$ and $\psrho_v(u)$, $u\in I^*_v$, is translation
of the variable by $u$ (the proof of (W2), (W3) in {\cite[Lemma~5.2]{P49}}; {\cite[Proposition~2.14]{P49}}). Thus
$u\mapsto\psZ_{I,v}\psrho_v(u)F$ is the function $F$ itself.

\emph{A Witt decomposition of $S$.} Let $e_1,\dots,e_{2m}$ be the positive and $f_1,\dots,f_m$ the negative
vectors of the frame of $S$, and $\omega_0=(1+\sqrt{-d})/2$, so $\omega_0+\aubar\omega_0=1$
({\cite[\S5.1.1]{P49}}). Put $u_i=e_i+f_i$ and $u'_i=\aubar\omega_0e_i-\omega_0f_i$ ($1\le i\le m$),
$\ell=\langle u_i\rangle$, $\ell'=\langle u'_i\rangle$, $S_0=\langle e_{m+1},\dots,e_{2m}\rangle$. Then
$\langle u_i,u_j\rangle=\langle u'_i,u'_j\rangle=0$, $\langle u_i,u'_j\rangle=\delta_{ij}(\aubar\omega_0+\omega_0)
=\delta_{ij}$, and $\ell\oplus\ell'\perp S_0$, $S_0$ positive definite of rank $m$. The matrix
$\bigl(\begin{smallmatrix}1&\aubar\omega_0\\1&-\omega_0\end{smallmatrix}\bigr)$ is integral with determinant
$-1$, so $L_S=\bigoplus_i\widehat\auO_Eu_i\oplus\bigoplus_i\widehat\auO_Eu'_i\oplus\bigoplus_{k>m}
\widehat\auO_Ee_k$. Accordingly $\mathsf W=I\oplus I^*\oplus\mathsf W_0$ with $I=\auHom(V,\ell)$,
$I^*=\auHom(V,\ell')$, $\mathsf W_0=\auHom(V,S_0)$; $I$ and $I^*$ are isotropic, $I^\perp=I\oplus\mathsf W_0$,
and $\mathsf W_0\to I^\perp/I=\mathsf W_I$ is an isometry, by which we identify $\mathsf W_I$ with
$\mathsf W_0$, a definite space. Write $u\in I^*$ and $l\in I$ by their matrices $M_u,N_l\in\pcMat_m$ in
the bases $(\varepsilon_k)$, $(u'_i)$, resp.\ $(\varepsilon_k)$, $(u_i)$. Then
$L=(L\cap I)\oplus(L\cap I^*)\oplus(L\cap\mathsf W_0)$ with $L\cap I=\{l:N_l\in\alpha\pcMat_m(\widehat\auO_E)\}$,
$L\cap I^*=\{u:M_u\in\alpha\pcMat_m(\widehat\auO_E)\}$. As $\langle\varepsilon_k,\varepsilon_k\rangle=-1$,
$\langle l,u\rangle_{\mathsf W}=\ptetr(l^\dagger u)=-\ptetr(N_l^*M_u)$.

\emph{Step 1: two subgroups of $\auU(S)$.} For a skew-hermitian $B\in\pcMat_m(\auA_E)$ let $z_B\in\auU(S)(\auA)$
be $1$ on $\ell\oplus S_0$ and $u'_j\mapsto u'_j+\sum_iB_{ij}u_i$; it is unitary since $B+B^*=0$, and
$Z=\{z_B\}\cong\{B\}$ is a unipotent abelian group with $Z(E^+)\backslash Z(\auA)$ compact (Haar measure of
mass $1$). For $Q\in\auGL_m(\auA_E)$ let $m_Q$ be $u_j\mapsto\sum_iQ_{ij}u_i$ on $\ell$,
$u'_j\mapsto\sum_i(Q^{-*})_{ij}u'_i$ on $\ell'$, and $1$ on $S_0$ ($m_Q\in\auU(S)(\auA)$ when $Q$ is a unit at
almost all places). Let $N$ be the unipotent radical of the stabiliser of $\ell$: the elements fixing $\ell$
pointwise and acting trivially on $\ell^\perp/\ell$. For every $C\in\auHom(\ell',S_0)(\auA)$ there is $n\in
N(\auA)$ with $n^{-1}x\in x+C(x)+\ell$ for $x\in\ell'$: put $n^{-1}=1$ on $\ell$,
$n^{-1}x=x+C(x)+D(x)$ on $\ell'$, $n^{-1}s=s+C^\sharp(s)$ on $S_0$, where $C^\sharp\colon S_0\to\ell$ is
defined by $\langle C^\sharp s,x\rangle=-\langle s,Cx\rangle$ ($x\in\ell'$) and $D\colon\ell'\to\ell$ by
$\langle D x,y\rangle=-\tfrac12\langle Cx,Cy\rangle$ ($x,y\in\ell'$); a direct check of the three types of
pairs shows $n^{-1}$ is unitary. All these elements stabilise $I$ (their action on $\mathsf W$ is $w\mapsto
xw$); $z_B$, $n$ act trivially on $\mathsf W_I$, and so does $m_Q$.

\emph{Step 2: the laws.} Let $F\in\psHS(\mathsf W)$ and $u\in I^*(\auA)$. For $p\in\auU(S)(\auA)$ stabilising
$\ell$, (W1) gives $\psrho(u)\omega_{\mathsf W}(p)=\omega_{\mathsf W}(p)\psrho(p^{-1}u)$, and (W2) gives
$\psZ_I\omega_{\mathsf W}(p)=\chi_I(p)\omega_{\mathsf W_I}(\bar p)\psZ_I$, with $\chi_I(p)=1$ if $p$ is
unipotent. Also $\psZ_I\psrho(i)=\psZ_I$ for $i\in I$ and $\psZ_I\psrho(w_0)=\psrho(w_0)\psZ_I$ for
$w_0\in\mathsf W_0$ (W2), and $\psrho(a)\psrho(b)=\beta(a,b)\psrho(a+b)$ with $|\beta|=1$.
(a) For $p=z_B$: $z_B^{-1}u=u-bu$ with $bu\in I$ of matrix $BM_u$, and
$\psrho(u-bu)=\beta(-bu,u)^{-1}\psrho(-bu)\psrho(u)$; hence
\[
\psZ_I\psrho(u)\,\omega_{\mathsf W}(z_B)F=\psi_u(B)\,\psZ_I\psrho(u)F,\qquad
\psi_u(B)=\beta(-bu,u)^{-1}=\psi\bigl(\tfrac12\jmath\langle u,bu\rangle\bigr)=\psi\bigl(-\tfrac12\ptetr(\jmath B\,M_uM_u^*)\bigr).
\]
(b) For $p=n$ as in Step~1: $n^{-1}u=u+d+i$ with $d=C\circ u\in\mathsf W_0$, $i\in I$, so
$\psZ_I\psrho(u)\omega_{\mathsf W}(n)F=\zeta\,\psrho_{\mathsf W_I}(C\circ u)\,\psZ_I\psrho(u)F$ with
$|\zeta|=1$. (c) For $p=m_Q$: $m_Q^{-1}u$ has matrix $Q^*M_u$, so $\psZ_I\psrho(u)\omega_{\mathsf W}(m_Q)F
=\chi_I(m_Q)\,\psZ_I\psrho(m_Q^{-1}u)F$, $\chi_I(m_Q)\ne0$. (d) For $y\in\auU(V)(\auA)$, acting by $w\mapsto
wy^{-1}$ (it stabilises $I$, $I^*$, $\mathsf W_0$): $\psZ_I\psrho(u)\omega_{\mathsf W}(y)F=\chi_I(y)\,
\omega_{\mathsf W_I}(\bar y)\,\psZ_I\psrho(uy)F$, where $uy$ has matrix $M_uY$; $\chi_I=1$ on
$\auU(V)(E^+)$ and on $K_V$ ($\upsilon$ is unramified at the finite places). The same identities hold place
by place ({\cite[Proposition~2.14(c)]{P49}}, (d)).

\emph{Step 3: unfolding.} For a hermitian $H\in\pcMat_m(E)$ put
\begin{equation}\label{pte:eq-coeff}
\varphi^\circ_H(x)=\int_{Z(E^+)\backslash Z(\auA)}\varphi^\circ(z_Bx)\,\aubar{\psi_H(B)}\,dB,\qquad
\psi_H(B)=\psi\bigl(-\tfrac12\ptetr(\jmath BH)\bigr).
\end{equation}
By (W3) with this $I$, $\theta(\omega(y,z_Bx)f)=\sum_{u\in I^*(E)}\psi_u(B)\,a_u(y,x)$ with
$a_u(y,x)=\theta_{\mathsf W_I}(\psZ_I\psrho(u)\omega(y,x)f)$ (Step~2(a); $\upsilon^m(\det z_B)=1$), and by the
bound in (W3) the double series and integral over compact fundamental sets for $[\auU(V)]$ (compact, $V$
definite) and $Z(E^+)\backslash Z(\auA)$ converge absolutely. For $u\in I^*(E)$, $B\mapsto\psi_u(B)$ is a
character of $Z(E^+)\backslash Z(\auA)$, equal to $\psi_H$ if and only if $M_uM_u^*=H$: $B\mapsto\jmath B$ is a
bijection of the skew-hermitian onto the hermitian matrices, $(X,Y)\mapsto\ptetr(XY)$ is a non-degenerate
$\auA_{E^+}$-valued pairing on hermitian matrices, and $\psi$ is non-trivial on $\auA_{E^+}$. Hence
$\varphi^\circ_H(x)=\sum_{M_uM_u^*=H}\int_{[\auU(V)]}a_u(y,x)\,dy$. Take $H=1$ and $u_0\in I^*(E)$ with
$M_{u_0}=1$. The solutions of $M_uM_u^*=1$ are the $M_u=Y^{-1}$, $Y\in\auU(V)(E^+)$ ($Y^*Y=1$ is the
unitarity for $\langle\,,\rangle_V=-{}^t\bar\cdot\,\cdot$), i.e.\ the $u=u_0\gamma^{-1}$, each once. By
Step~2(d) and (W1) for $\mathsf W_I$, $a_{u_0\gamma^{-1}}(y,x)=a_{u_0}(\gamma^{-1}y,x)$ (the source character of
$\omega$ is trivial). Unfolding,
\[
\varphi^\circ_1(x)=\int_{\auU(V)(\auA)}a_{u_0}(y,x)\,dy,
\]
absolutely convergent.

\emph{Step 4: invariance in $y$.} $\omega(k,1)f=f$ for $k\in\auU(V)_\infty K_V$. At the finite places $(k,1)$
stabilises $L$ and (W4) applies. At $\tau$, Lemma~\ref{ptf:fock} with source $V$ (all signs negative,
weight $0$) and target $S$ ($n_T^+=2m$) gives $\omega_\infty(g,1)v=\prod_\tau\det(g_\tau)^{2m}\cdot
v\circ\sigma_{g,1}$, and $\sigma_{g,1}$ replaces the block $y^-=X^{--}$ by $X^{--}g^{-\mathsf t}$, which
multiplies $v$ by $\prod_\tau\det(g_\tau)^{-2m}$; so $\omega_\infty(g,1)v=v$. Since $\omega(yk,x)=\omega(y,x)
\omega(k,1)$, $a_{u_0}(yk,x)=a_{u_0}(y,x)$, and as $\auU(V)_\infty$ and $K_V$ have volume $1$,
$\varphi^\circ_1(x)=\sum_{y\in\auU(V)(\auA_f)/K_V}a_{u_0}(y,x)$.

\emph{Step 5: support at the finite places.} For $u\in I^*(\auA_f)$, $\psZ_{I,f}\psrho_f(u)\psvac=0$ unless
$u\in L\cap I^*$, and $\psZ_{I,f}\psrho_f(u)\psvac=\psZ_{I,f}\psvac$ if $u\in L\cap I^*$. Indeed, by (W4)
$\psrho_f(l)\psvac=\epsilon_L(l)\psvac$ for $l\in L$, and $\epsilon_L(l)=\psi_f(-\frac12\omega_0\jmath\langle
l,l\rangle)=1$ for $l$ in the isotropic $I$ or $I^*$. For $l\in L\cap I$, $\psrho(u)\psrho(l)=[u,l]\psrho(l)
\psrho(u)$ and $\psZ_I\psrho(l)=\psZ_I$, so $(1-[u,l])\psZ_{I,f}\psrho_f(u)\psvac=0$. Here
$[u,l]=\beta(u,l)\beta(l,u)^{-1}=\psi_f(-\jmath\,\ptetr(N_l^*M_u))$, using $\langle l,u\rangle=-\ptetr(N_l^*M_u)$
and $\psi(\jmath r)=1$ for $r\in\auA_{E^+}$. With $N_l=\alpha nE_{ik}$, $n\in\widehat\auO_E$, and
$M_u=\alpha\mu$, the condition $[u,l]=1$ for all such $l$ reads $\psi_f(-\jmath\alpha\aubar\alpha\,\bar n\mu_{ik})=1$
for all $n$, i.e.\ $\mu_{ik}\in\widehat\auO_E$ by the self-duality defining $\alpha$ ({\cite[\S5.1.1]{P49}}).
So the vector vanishes unless $M_u\in\alpha\pcMat_m(\widehat\auO_E)$; for such $u$, $\psrho_f(u)\psvac=\psvac$.
In particular $\psZ_{I,f}\psvac\ne0$: in the mixed model $\psvac$ is a non-zero function on $I^*(\auA_f)$,
whose value at $u$ is $\psZ_{I,f}\psrho_f(u)\psvac$, which is $0$ or $\psZ_{I,f}\psvac$.

\emph{Step 6: the archimedean slice.} $v\ne0$, so in the mixed model at the infinite places
$u\mapsto\psZ_{I,\infty}\psrho_\infty(u)v$ is a continuous non-zero function on $I^*_\infty=\prod_\tau
\pcMat_m(\C)$; it is non-zero on a non-empty open set, which contains a point $u_\infty$ with invertible
matrix $U_\infty$ (invertible matrices are dense). Let $Q\in\auGL_m(\auA_E)$ have archimedean component
$U_\infty^*$ and finite component $\aubar\alpha\cdot1_m$ (a unit at almost all places), and put $m=m_Q$. Then
$m^{-1}u_0$ has matrix $Q^*$: the point $u_\infty$ at the infinite places and $\alpha1_m$ at the finite ones.

\emph{Step 7: conclusion.} Let $n\in N(\auA)$ be as in Step~1 for some $C$, and $x=nm$. By Step~2(b), (c)
and $\omega(y,x)=\upsilon^m(\det m)\omega_{\mathsf W}(n)\omega_{\mathsf W}(m)\omega_{\mathsf W}(y)$ (on
$\psHS(\mathsf W)$, with the trivial source character),
\[
a_{u_0}(y,x)=\zeta\,\upsilon^m(\det m)\chi_I(m)\;\theta_{\mathsf W_I}\bigl(\psrho_{\mathsf W_I}(C\circ u_0)\,
\psZ_I\psrho(m^{-1}u_0)\,\omega_{\mathsf W}(y)f\bigr).
\]
{\emergencystretch=6em For $y\in\auU(V)(\auA_f)$, by Step~2(d) at the finite places, $\psZ_I\psrho(m^{-1}u_0)\allowbreak \omega_{\mathsf W}(y)f=
\psZ_{I,\infty}\allowbreak \psrho_\infty(u_\infty)v\otimes\chi_I(y)\omega_{\mathsf W_I,f}(\bar y)\,\psZ_{I,f}\psrho_f(u_f)
\psvac$ with $u_f$ of matrix $\alpha Y$. By Step~5 this vanishes unless $Y\in\pcMat_m(\widehat\auO_E)$, and an
integral $Y$ with $Y^*Y=1$ lies in $K_V$. So only the coset $K_V$ contributes in Step~4, and
\[
\varphi^\circ_1(nm)=\zeta\,\upsilon^m(\det m)\chi_I(m)\,\theta_{\mathsf W_I}\bigl(\psrho_{\mathsf W_I}(C\circ u_0)
\Lambda\bigr),\qquad\Lambda=\psZ_{I,\infty}\psrho_\infty(u_\infty)v\otimes\psZ_{I,f}\psvac ,
\]
with $\Lambda\ne0$ by Steps~5 and~6 and $|\zeta|=1$. As $u_0$ is an isomorphism $V\to\ell'$, $C\circ u_0$ runs
over all of $\mathsf W_0(\auA)=\mathsf W_I(\auA)$ when $C$ runs over $\auHom(\ell',S_0)(\auA)$. If
$\varphi^\circ=0$, then $\varphi^\circ_1=0$ and $\theta_{\mathsf W_I}(\psrho_{\mathsf W_I}(w)\Lambda)=0$ for
all $w\in\mathsf W_I(\auA)$, so $\Lambda$ lies in the kernel of $F\mapsto\theta_{\mathsf W_I}(\psrho(\cdot)F)$,
a proper closed invariant subspace of the irreducible representation $\psHS(\mathsf W_I)$, hence $\Lambda=0$
(as in the proof of {\cite[Lemma~5.26]{P49}}; compare Proposition~\ref{pst:I4}). This contradiction proves
$\varphi^\circ_1\ne0$, so $\varphi^\circ\ne0$.\par}
\end{proof}

\subsubsection{The third exchange}\label{pte:ss-third}

We use the notation of \S\ref{po:ss-glue} and \S\ref{po:ss-hecke}; $m$ is even, $k=r+1-m$ is even, and
$N''=r+9m+1$. Put $W'=W\oplus(E^j)^-\oplus E^{3m-j}\oplus W^\flat$, $W^\flat=E^{6m+1}$ positive definite,
of rank $N''$ with $l=(m-j)+j=m$ negative signs; $\mathsf W=I_{3m}\oplus W'$; $V''=I_{N''}$ (the space of
\S\ref{po:ss-glue}). For a twist $\vartheta$ let $\pcm^{(2)}_\vartheta:=\pi_M^2\prod_{w\in Q_\vartheta}w$ (it has the same primes as the squarefree modulus $\pcm_\vartheta$ of {\cite[\S2.3]{P53}}, so $L_*(\pcm^{(2)}_\vartheta,\cdot,\cdot)=L_*(\pcm_\vartheta,\cdot,\cdot)$ by {\cite[(11)]{P53}}), $\chi'=
\chi_0^j\vartheta$ (weight $0$, residue character $\chi'_\vartheta$ on $\pcR_{\pcm^{(2)}_\vartheta}^\times$: $\omega^j$ at
$\mathsf M$, $\psi_w$ at $Q_\vartheta$), $\nu''_\vartheta=\mu_j^0\vartheta\,\upsilon^{j+\lfloor(t+k-1)/2\rfloor}$ (the
character of the lift on $V''$) and $\nu^\flat=\nu''_\vartheta\upsilon^m/\chi'$, a character independent of
$\vartheta$ with trivial tame part. The \emph{choices} are $i=(\Lambda_i,h^\flat_i,\varrho_i)$: a matrix
$\Lambda_i\in\pcMat_{3m\times N''}(\auO_E/\pi_M^2)$, a class $h^\flat_i$ of $\auU(W^\flat)$ at level $\pi_M^2$, and
for each $\tau$ a family of $2m$-row selections of a fixed basis of short vectors of the lattice at
$h^\flat_i$; they give at each $\tau$ a family of complex $2m\times(3m+N''-m)$ matrices $\eta_i=(\eta_{i,\tau k})_{k\in
\iota(\tau)}$ (rows: basis vectors of the padded $W$ for the columns of $\mu^+_\tau$ of height $<r_+$, and the
selected short rows, read at $\tau$, for the $\Delta'$ columns of height $r_+$), with $\#\iota(\tau)=$ the number
of columns of $\mu^+_\tau$. The matrices $\eta_{i,\tau k}$ are those of Lemma~\ref{ptf:H}. Let $P_i=\pteMP(\eta_i)$; it involves only the rows of $W'$ (the rows of $\eta_i$ have
no component on $I_{3m}$), and it is the product of the minors of the padded $W$ and the determinants
$\det(\eta^\flat_{i,\tau k}y^{++}_{W^\flat,\tau})$ on $W^\flat$. Let $\psi_\Lambda(y)=\psi_E(\pi_M^{-2}\mathrm{tr}(\Lambda y))$
on $\pcMat_{N''\times3m}(\auO_E/\pi_M^2)$ and $\psi^{\mathrm{inv}}_\Lambda(y)=\sum_{\gamma\in\auGL_{3m}(\auO_E/\pi_M^2)}
\psi_\Lambda(y\gamma)$. Define
\begin{equation}\label{pte:eq-F}
F_i(g)=\int_{[\auU(S)]}\varphi^\circ(x)\,\theta^{S,W'}_{\nu^\flat\upsilon^m,\upsilon^m}\bigl(x,(g,1,1,h^\flat_i);
[P_i,\psi_{\Lambda_i}]\bigr)\,dx\qquad(g\in\auU(W)(\auA)),
\end{equation}
independent of $\vartheta$; this is the function $F_i$ of Proposition~\ref{po:exchange}
($h^\flat_i$ is the point $h^\flat_i$ there, and the twist-free datum $\psi_{s_i}$ is $\psi_{\Lambda_i}$). The
integral converges absolutely ($\varphi^\circ$ is rapidly decreasing, Lemma~\ref{pte:carrierprop}).

\begin{proposition}[Third exchange]\label{pte:third}
There are $C_3\ge0$ and for each twist $\vartheta$ a number $c_3(\vartheta)$ with $|c_3(\vartheta)|^2L_*(\pcm^{(2)}_\vartheta,
(\nu_*\vartheta)^{-1},z_*)=C_3$, such that, with $\gamma=\gamma(\vartheta)$ of Proposition~\ref{pte:carrier} for
$c=\chi'$ and $\mathsf N=\#\auGL_{3m}(\auO_E/\pi_M^2)$, for all $i$ and $g$:
\[
\frac{c_3(\vartheta)\gamma(\vartheta)}{\mathsf N}\,\mathcal L_{\vartheta,i}(g)=F_i(g),\qquad
\mathcal L_{\vartheta,i}(g)=\sum_{h'}\mathrm a_{\chi'}(h')\sum_{h}\theta^{\mathsf W,V''}_{\nu''_\vartheta,\,\chi'\nu^\flat
\upsilon^m}\bigl((h',g,1,1,h^\flat_i),h;[\pteSD(\eta_i),\Phi_{\vartheta,i}]\bigr),
\]
where $h'$, $h$ run over the classes of $\auU(I_{3m})$, $\auU(V'')$ at full level, each term weighted by the
inverse order of its stabiliser in $\Gamma$, and $\Phi_{\vartheta,i}(x'\mid x'')=\chi'_\vartheta(\det x'')\psi^{\mathrm{inv}}_{
\Lambda_i}(x''^{-1}x')$ on $\pcMat_{N''\times(3m+N'')}(\pcR_{\pcm^{(2)}_\vartheta})$.
\end{proposition}
\begin{proof}
Ranks: $n=3m+N''$, $a=N''$, $b=3m$ (the space $V'=S$, $p=2m$, $l=m$), and $a-b=N''-3m=7m+k=2z_*$, so
$d=z_*=\tfrac{7m}2+\tfrac k2$; $2l+2\le n$. The character $\alpha=\chi'\nu^\flat\upsilon^m=\nu''_\vartheta\upsilon^m$ is tame
for $\pcm^{(2)}_\vartheta$ with residue character $\chi'_\vartheta$ and weight $m-\#\iota(\tau)$: the weight of $\nu^\flat\upsilon^m$
is $m+j+\lfloor(t+k-1)/2\rfloor-\kappa_j(\tau)$ and $\#\iota(\tau)+j+\lfloor(t+k-1)/2\rfloor=\kappa_j(\tau)$ (the largest part
of $\mu^+_\tau$; see \S\ref{po:ss-set}). Corollary~\ref{pte:explicit} for $(\mathsf W,V'=S,I_a=V'')$, $\eta_i$,
$\beta=\nu''_\vartheta$ gives $c_3(\vartheta)$ with $|c_3|^2L_*(\pcm^{(2)}_\vartheta,\alpha^{-1}\upsilon^{3m+z_*},z_*)=C_3$; and
$\alpha^{-1}\upsilon^{3m+z_*}=(\nu_*\vartheta)^{-1}$ because $\nu_*=\mu_j^0\upsilon^{n'}$, $n'=j+(t-11m-1)/2$, and
$j+\lfloor(t+k-1)/2\rfloor+m-3m-z_*=n'$ (direct check with $z_*=7m/2+k/2$, $t$, $k$ odd/even as fixed in the data).
Apply Lemma~\ref{pte:one} with $V'=S$, $I_b=I_{3m}$, $W'$, the character $\chi'_\vartheta$, $\Psi'=\psi^{\mathrm{inv}}_{\Lambda_i}$
(right invariant under $\auGL_{3m}$ by construction), $F=\varphi^\circ$, $c_0=\gamma$, $\mathrm a=\mathrm a_{\chi'}$
(hypothesis (i) is Proposition~\ref{pte:carrier}), and (ii) the third exchange at the points
$(h',g,1,1,h^\flat_i)$ (the polynomial $P_i$ involves only the rows of $W'$). The right side of
\eqref{pte:eq-explicit} has the characters $\alpha=\chi'\cdot\nu^\flat\upsilon^m$ on $\auU(S)$ and
$\alpha\beta^{-1}=\upsilon^m$ on $\auU(\mathsf W)$; we take $(a_1,a_2,\beta_0)=(\chi',\nu^\flat\upsilon^m,\upsilon^m)$, so that
$\theta^{S,I_{3m}}_{a_1,\beta_0}=\theta^{S,I_{3m}}_{\chi',\upsilon^m}$ is the kernel of Proposition~\ref{pte:carrier}; the lemma gives
$c_3\gamma\,\mathcal L_{\vartheta,i}(g)=\int\varphi^\circ(x)\theta^{S,W'}_{\nu^\flat\upsilon^m,\upsilon^m}(x,(g,1,1,h^\flat_i);
[P_i,\psi^{\mathrm{inv}}_{\Lambda_i}])\,dx$. By Lemma~\ref{pte:translates} ($\nu^\flat\upsilon^m$ has trivial tame part,
$\varphi^\circ$ is right $K_S$-invariant and rapidly decreasing) the right side is $\mathsf N\,F_i(g)$.
\end{proof}

\subsubsection{The left side as the theta series of the first factors}\label{pte:ss-left}

Let $w_i\in\pcMat_{(3m+N'')\times N''}(\auO_E/\pi_M^2)$ be $\Lambda_i$ stacked over $\pi_M1_{N''}$, and $e$ the
inclusion of the $r$ indices of $W$ into those of $V''$. Let $H_{\vartheta,i}\subseteq\auGL_{N''}(\pcR_{\pcm^{(2)}_\vartheta})$ be
the group of the $A$ whose columns off $e$ are those of $1$ and whose reduction $\bar A$ modulo $\pi_M^2$
satisfies $w_i\bar A=w_i$, and $J'_{\vartheta,i}$ the group of the integral points of $\auU(V'')(\auA_f)$ whose reduction modulo
$\pcm^{(2)}_\vartheta$ lies in $H_{\vartheta,i}$. Explicitly, at $w\in\mathsf M$: $k\equiv1\bmod\pi_M$, $\Lambda_ik\equiv\Lambda_i$ and the
last $N''-r$ columns of $k$ are those of $1$ modulo $\pi_M^2$; at $w\in Q_\vartheta$: the last $N''-r$ columns of
$k$ are those of $1$ modulo $w$. It equals $J_{\vartheta,i}$ of \S\ref{po:ss-glue} (with $s_i=\Lambda_i$),
contains the principal congruence subgroup of level $\pcm^{(2)}_\vartheta$, and is free ({\cite[Lemma~14.5]{P46}}).

\begin{lemma}[The column datum is the slice datum]\label{pte:column}
For $\sigma\in\pcMat_{N''\times r}(\pcR_{\pcm^{(2)}_\vartheta})$ let $\mathrm{col}(\sigma)=\chi'_\vartheta(\det\sigma')$ if $w_i\bar\sigma=
w_i\iota''_0$ and $0$ otherwise ($\sigma'$ the rows $e$). Then $\mathrm{col}$, read as a function of integral
matrices at $\mathsf M\cup Q_\vartheta$ (integral in the coordinates $T=a_w^{-1}x$ of (N5) used in
\S\ref{po:ss-glue}, which are the coordinates of {\cite[Definition~6.31]{P53}}, {\cite[Lemma~6.32(c)]{P53}}), is the finite part at $\mathsf M\cup Q_\vartheta$ of $\Psi''_{\vartheta,\Lambda_i}$ of
\S\ref{po:ss-glue}.
\end{lemma}
\begin{proof}
$w_i\bar\sigma=w_i\iota''_0$ means $\Lambda_i(\sigma-\iota''_0)\equiv0\bmod\pi_M^2$ and $\pi_M(\sigma-\iota''_0)\equiv0\bmod\pi_M^2$,
i.e.\ $\sigma\equiv\iota''_0\bmod\pi_M$ and $\Lambda_i(\sigma-\iota''_0)\equiv0\bmod\pi_M^2$, i.e.\ $\sigma\in\iota''_0+\pi_ML''_{\Lambda_i}$
(definition of $L''_s$). On this set $\sigma'\equiv1\bmod\pi_M$, so the factor $\omega^j(\det\sigma')$ at $\mathsf M$ is
$1$. At $Q_\vartheta$ there is no condition ($w_i$ lives modulo $\pi_M^2$) and the factor is $\psi_w(\det\sigma')$, with
$\psi_w(0)=0$.
\end{proof}

\begin{lemma}[Restriction to the row stabiliser]\label{pte:rowstab}
Let $F_0$ be a function on $\pcMat_{N''\times3m}(\pcR_{\pcm^{(2)}_\vartheta})$ with $F_0(Ay)=F_0(y)$ for $A\in H_{\vartheta,i}$,
$\Phi(x'\mid x'')=\chi'_\vartheta(\det x'')F_0(x''^{-1}x')$ and $\Phi_0(x'\mid x'')=\mathrm{col}(x''e)\,F_0(x')$ if the
columns of $x''$ off $e$ are those of $1$, $0$ otherwise. Let $\theta[\cdot](h)$ be a kernel of $(\mathsf W,V'')$
whose polynomial is fixed by $\auU(V'')_\infty$ and whose character on $\auU(V'')$ is tame with residue character
$\chi'_\vartheta$. Then $\sum_{h\bmod K_{V''}}\theta[\Phi](h)=\sum_{h\bmod J'_{\vartheta,i}}\theta[\Phi_0](h)$ (sums over
classes weighted by inverse stabiliser orders).
\end{lemma}
\begin{proof}
For $A\in\auGL_{N''}(\pcR)$ put $(\mathrm{tr}_A\Psi)(x)=\chi'_\vartheta(\det A)\Psi(A^{-1}x)$; by {\cite[Lemma~2.9]{P53}},
$\theta[\Psi](hk)=\theta[\mathrm{tr}_{\bar k}\Psi](h)$ for $k\in K_{V''}$. One checks $\mathrm{tr}_A\Phi=\Phi$, and
$\Phi_0(x)=\Phi(x)$ if $x''\in H_{\vartheta,i}$, $0$ otherwise (for $x''\in H_{\vartheta,i}$: $x''e$ satisfies the condition of
$\mathrm{col}$, $\det x''=\det(x''e)'$ because the other columns are those of $1$, and $F_0(x''^{-1}x')=F_0(x')$; conversely
if $\mathrm{col}(x''e)\ne0$ and the other columns are those of $1$, then $x''\in H_{\vartheta,i}$). Hence $\Phi=\sum_{A\in
\auGL/H_{\vartheta,i}}\mathrm{tr}_A\Phi_0$ (an invertible $x''$ lies in exactly one coset $AH_{\vartheta,i}$; for non-invertible $x''$
all terms vanish). So $\theta[\Phi](h)=\sum_A\theta[\Phi_0](h\tilde A)$ for integral lifts $\tilde A$, and the function
$\theta[\Phi_0]$ is right $J'_{\vartheta,i}$-invariant. Summing over the classes modulo $K_{V''}$, the classes modulo
$J'_{\vartheta,i}$ are the $h\tilde A$, each counted with the correct weight (the tower relation of class sums:
$\sum_{h\bmod K}\sum_A f_0(h\tilde A)/\#\mathrm{Stab}_K(h)=\sum_{h\bmod J'}f_0(h)/\#\mathrm{Stab}_{J'}(h)$ because $K/J'\cong
\auGL/H_{\vartheta,i}$ and the orbit--stabiliser relation).
\end{proof}

\begin{definition}[First factors]\label{pte:def-ff}
By Lemma~\ref{pte:rowstab} with $F_0=\psi^{\mathrm{inv}}_{\Lambda_i}$ (invariant under $H_{\vartheta,i}$: for $A\in H_{\vartheta,i}$,
$\mathrm{tr}(\Lambda_iAy\gamma)=\mathrm{tr}(\Lambda_iy\gamma)$ since $\Lambda_i\bar A=\Lambda_i$), the inner sum of
$\mathcal L_{\vartheta,i}$ is a class sum modulo $J'_{\vartheta,i}$ of the kernel with datum $\Phi_0=\mathrm{col}(x''e)\cdot
\plone[x''\text{ has the columns of }1\text{ off }e]\cdot\psi^{\mathrm{inv}}_{\Lambda_i}(x')$. The first factors
$\varphi''_{\vartheta,i}$ are those of Definition~\ref{ptf:def-ff}: they come from the Laplace expansion
$\pteSD(\eta_i)=\sum_\lambda\pop''_{u_\lambda}R_\lambda$ (Lemma~\ref{ptf:laplace}), and their factorisation property and
transformation laws are Proposition~\ref{ptf:ff}.
\end{definition}

\begin{proposition}[The left side]\label{pte:leftside}
With the level $J'_{\vartheta,i}$ of \S\ref{po:ss-glue} and the first factors of Definition~\ref{pte:def-ff},
$\mathcal L_{\vartheta,i}=\Xi_{\vartheta,i}$ on $\auU(W)(\auA)$.
\end{proposition}
\begin{proof}
Proposition~\ref{ptf:leftside}; the character $\chi''_\vartheta$ there is that of \S\ref{po:ss-glue} by Lemma~\ref{ptf:chi}.
\end{proof}

\subsubsection{Proof of Proposition~\ref{po:exchange}}\label{pte:ss-proof}

Put $c'_\vartheta=c_3(\vartheta)\gamma(\vartheta)/\mathsf N$ and $C_{\mathrm x}=C_3C_{12}/\mathsf N^2\ge0$. By
Propositions~\ref{pte:third} and~\ref{pte:leftside}, $c'_\vartheta\Xi_{\vartheta,i}=F_i$ for all $i$, with $F_i$ of
\eqref{pte:eq-F} independent of $\vartheta$. By Propositions~\ref{pte:carrier} (for $c=\chi'=\chi_0^j\vartheta$, modulus
$\pcm^{(2)}_\vartheta$, whose primes are $\mathsf S_\vartheta$) and~\ref{pte:third},
\[
|c'_\vartheta|^2\,L_*(\pcm^{(2)}_\vartheta,(\nu_*\vartheta)^{-1},z_*)\,L_*(\pcm^{(2)}_\vartheta,\chi'^{-1}\upsilon^m,m)\,L_*(\pcm^{(2)}_\vartheta,\chi'\upsilon^{m/2},m/2)
=C_{\mathrm x}.
\]
Now $\chi'^{-1}\upsilon^m=(\nu_+\vartheta)^{-1}$ ($\nu_+=\upsilon^{-m}\chi_0^j$, $z_+=m$) and $\chi'\upsilon^{m/2}=(\nu_-\vartheta^{-1})^{-1}$
($\nu_-=\upsilon^{-m/2}\chi_0^{-j}$, $z_-=m/2$), and $L_*(\pcm^{(2)}_\vartheta,\nu,z)=(\auNm\mathsf S_\vartheta)^{z^2}\prod_{z'=1}^z|L^{\mathsf S_\vartheta}(z',\nu)|^2$
({\cite[(11)]{P53}}; the primes of $\pcm^{(2)}_\vartheta$ are those of $\mathsf S_\vartheta$, each counted once). So the product of the three
$L_*$ is $H_j(\vartheta)$ of \eqref{po:eq-H}, and $|c'_\vartheta|^2H_j(\vartheta)=C_{\mathrm x}$.

Integrability: $F_i=c'_1\Xi_{1,i}$ and $\Xi_{\vartheta,i}$ are finite sums of theta kernels of $(W,V'')$, which have moderate
growth on $\auU(W)(\auA)$; so is $g\mapsto K^\circ(g,h^\circ)$; and $y\in\poA(W_j)$ is rapidly decreasing (cuspidal,
{\cite[\S14.3]{P46}}). Hence $yF_iK^\circ(\cdot,h^\circ)$ and $y\Xi_{\vartheta,i}K^\circ(\cdot,h^\circ)$ are integrable on $[\auU(W)]$.
\qed

\begin{remark}\label{pte:rem-final}
(i) Corollary~\ref{pte:explicit} alone gives only $C_{\mathrm x}\ge0$; $C_{\mathrm x}>0$ is proved in Proposition~\ref{po:glue}(G2),
from $c'_\vartheta\ne0$ (injectivity of $\Theta^\circ$) and $H_j(\vartheta)>0$ (Lemma~\ref{po:Hpos}). (ii) The level $J_{\vartheta,i}$ of \S\ref{po:ss-glue} is the group $J'_{\vartheta,i}$ of \S\ref{pte:ss-left}; every use of it
in \S\ref{po:ss-glue} and Proposition~\ref{po:glue} (Lemma~\ref{po:lattice}(d), freeness, representative independence,
Lemma~\ref{po:classes}) holds for this smaller group. The character $\chi''_\vartheta$ on $J'_{\vartheta,i}$ is
$k\mapsto\chi'_\vartheta(\det\bar k)$ (Proposition~\ref{ptf:ff}(b), Lemma~\ref{ptf:chi}). (iii) The
matrices $s_i$ of \S\ref{po:ss-glue} are the $\Lambda_i$.
\end{remark}



\providecommand{\auO}{\mathcal O}
\providecommand{\auU}{\mathrm U}
\providecommand{\auGL}{\mathrm{GL}}
\providecommand{\auHom}{\operatorname{Hom}}
\providecommand{\auA}{\mathbb A}
\providecommand{\auPl}{\Phi}
\providecommand{\aubar}[1]{\overline{#1}}
\providecommand{\pcR}{R}
\providecommand{\pcm}{\mathfrak m}
\providecommand{\pcMat}{\mathrm M}
\providecommand{\psvac}{{\phi^{\mathrm{vac}}}}
\providecommand{\psrho}{\varrho}
\providecommand{\psHS}{\mathcal H}
\providecommand{\psZ}{\mathsf Z}
\providecommand{\poT}{\mathbb T}
\providecommand{\poA}{\mathcal A}
\providecommand{\pop}{\mathsf p}
\providecommand{\pog}{\mathsf g}
\providecommand{\plone}{\mathbf 1}
\providecommand{\pteSD}{\mathrm{SD}}
\providecommand{\pteMP}{\mathrm{MP}}
\providecommand{\pteCS}{\mathop{\textstyle\sum^{\mathrm{cl}}}}

\subsection{Three exchanges: conjugation, the Fock action and the first factors}\label{ptf:sec}

This section gives complete proofs of the steps of \S\ref{pte:sec} that were stated there in
abbreviated form. Each result below is self-contained; \S\ref{pte:sec} cites it at the place
indicated.

\subsubsection{Negation and complex conjugation of theta kernels}\label{ptf:ss-conj}

Let $S,T$ be spaces ({\cite[\S2.3]{P53}}) of ranks $n_S,n_T$ and $W=\auHom(S,T)$, a hermitian space of
rank $N=n_Sn_T$ with the form of {\cite[\S5.1.1]{P49}}. Let $S^-$ be the space with the same underlying
$E$-vector space and the same frame as $S$, and with the form $-\langle\,,\rangle_S$; its lattice is
the same unimodular lattice, its signature is the opposite one, and $\auU(S^-)=\auU(S)$ as groups (an
isometry of a form is an isometry of its negative). Put $W^-=\auHom(S^-,T)$. The adjoint of
$w\colon S\to T$ for $-\langle\,,\rangle_S$ is the negative of the adjoint for $\langle\,,\rangle_S$,
so \emph{$W^-$ is the $E$-vector space $W$ with the form $-\langle\,,\rangle_W$}; the frames of $W$ and
$W^-$ (the matrix units in the frames of $S$ and $T$) coincide, so the two spaces have the same
coordinates at every place, and $\auU(W)=\auU(W^-)$. The pair $(g,h)$ acts on both by $w\mapsto
hwg^{-1}$, and
\begin{equation}\label{ptf:eq-detonhom}
\det\bigl(w\mapsto hwg^{-1}\bigr)=\det(h)^{n_S}\det(g)^{-n_T}
\end{equation}
(on $\pcMat_{n_T\times n_S}$, left multiplication by $h$ acts on each of the $n_S$ columns and right
multiplication by $g^{-1}$ on each of the $n_T$ rows).

\emph{Heisenberg representations.} The bicharacter of {\cite[\S5.1.1]{P49}} of $W^-$ is
$\beta^-(x,y)=\psi(\frac12\jmath\langle y,x\rangle_{W^-})=\beta(x,y)^{-1}=\aubar{\beta(x,y)}$, and
$[x,y]^-=\aubar{[x,y]}$. If $(\psrho,\psHS)$ is a Heisenberg representation of $W$ (at one place or
adelically), then the complex conjugate space $\aubar{\psHS}$ (the same set, with $\lambda\cdot f=
\aubar\lambda f$) with the same operators $\psrho(x)$ satisfies $\psrho(x)\psrho(y)=\aubar{\beta(x,y)}
\psrho(x+y)$, so it is a Heisenberg representation of $W^-$ with the same central character
$\psi$ on the centre written as $W$ plus scalars (the centre acts through $\psi$ composed with the
real part of the form, which changes sign together with the conjugation). By the uniqueness of the
Heisenberg representation (Stone--von Neumann; at a finite place the smooth version) there is an
\emph{antilinear} bijection $A_v\colon\psHS_v(W)\to\psHS_v(W^-)$ with
\begin{equation}\label{ptf:eq-Arho}
A_v\,\psrho(x)=\psrho^-(x)\,A_v\qquad(x\in W_v),
\end{equation}
unique up to a non-zero scalar (two such maps differ by a linear self-intertwiner of an irreducible
representation, i.e.\ a scalar, by Schur's lemma). We fix the scalars as follows.

\begin{itemize}
\item[(A$_\infty$)] At $\tau\in\auPl$ the vacuum of $W_\tau$ is the unique line killed by the
annihilation operators, which are the operators $d\psrho(z)$ for $z$ in the antiholomorphic part of
$W_\tau\otimes_\R\C$ for the complex structure $J_\tau$ of the Fock model; for $W^-_\tau$ the complex
structure is $-J_\tau$ (the positive and negative coordinates are exchanged). By
\eqref{ptf:eq-Arho} $A_\tau\,d\psrho(z)=d\psrho^-(\aubar z)A_\tau$, and complex conjugation exchanges the
$\pm i$-eigenspaces of $J_\tau$, which are the $\mp i$-eigenspaces of $-J_\tau$. So $A_\tau$ carries the
vacuum line of $W_\tau$ to that of $W^-_\tau$, and we normalise $A_\tau\psvac_\tau=\psvac^-_\tau$ (the real
vacuum $1$ of the Fock model in both cases). The Fock vector of a polynomial $P$ in the holomorphic
coordinates $\zeta$ of $W_\tau$ (\eqref{po:eq-poly}; the coordinate of an entry of positive product sign is the
entry, that of an entry of negative product sign its conjugate) is $P(c_1,\dots)\psvac_\tau$, where the
$c_i$ are the creation operators attached to the coordinates (complex linear combinations of the
$d\psrho(x)$). Applying $A_\tau$ conjugates the coefficients of these combinations; the creation
operator of the coordinate $\zeta_i$ of $W_\tau$ goes to the creation operator of $\aubar{\zeta_i}$,
which is the holomorphic coordinate of the same entry in $W^-_\tau$, since all product signs have
changed. Hence
\[
A_\tau\bigl(\text{Fock vector of }P\bigr)=\text{Fock vector of }\aubar P\ \text{for }W^-_\tau,
\]
$\aubar P$ the polynomial with conjugated coefficients, in the holomorphic coordinates of $W^-_\tau$.
\item[(A$_\pcm$)] At a place $v=w\aubar w$ below a prime $w\mid\pcm$ (split, {\cite[\S2.3]{P53}}) both
$\psHS_v(W)$ and $\psHS_v(W^-)$ are realised by (W5) on Schwartz--Bruhat functions on
$\pcMat_{n_T\times n_S}(E_w)$, with $\psrho(a_w,0)$ the translation by $a_w$ and $\psrho(0,a_{\aubar w})$
the multiplication by $x\mapsto[x,a_{\aubar w}]$. Put $A_v\Psi=\aubar\Psi$. Translation commutes with
conjugation and $\aubar{[x,a]}=[x,a]^-$, so \eqref{ptf:eq-Arho} holds. In particular $A_v$ carries
$\Psi^\natural$ to $(\aubar\Psi)^\natural$ and the normalised vacuum (the normalised indicator of the
integral matrices, which is real) to itself.
\item[(A$_f$)] At the other finite places $v$: the lattice $L_v=\alpha\auHom(L_S,L_T)_v$ is the same
for $W$ and $W^-$ and self-dual for both ($[\,,]^-=\aubar{[\,,]}$). By \eqref{ptf:eq-Arho} $A_v$ carries
the line on which $\psrho(l)$, $l\in L_v$, acts by $\epsilon_L(l)$ to the line on which it acts by
$\aubar{\epsilon_L(l)}$, and $\aubar{\epsilon_L(l)}=\psi_f(+\frac12\omega\jmath\langle l,l\rangle)=\epsilon^-_L(l)$ is the
character of the vacuum of $W^-$ ({\cite[\S5.1.1]{P49}}, formula for $\epsilon_L$ with the form negated).
So $A_v\psvac_v$ is a non-zero multiple of $\psvac^-_v$; both are unit vectors for the local inner
products ($A_v$ is isometric up to a scalar, since $\langle A_vf,A_vf'\rangle$ is a conjugate of an
invariant hermitian form, hence a positive multiple of $\aubar{\langle f,f'\rangle}$), and we normalise
$A_v\psvac_v=\psvac^-_v$.
\end{itemize}
Let $A=\bigotimes_vA_v\colon\psHS(W)\to\psHS(W^-)$ (well defined on the restricted tensor product since
$A_v\psvac_v=\psvac^-_v$ for all $v\notin S(\pcm)\cup\auPl$; on the places of $\pcm$ the vacuum is also
preserved, so $A$ does not depend on $\pcm$). By (A$_\infty$), (A$_\pcm$), (A$_f$) and the definition of
the datum in {\cite[\S2.3]{P53}},
\begin{equation}\label{ptf:eq-Adatum}
A\,[P,\Psi]=[\aubar P,\aubar\Psi]\qquad\text{(datum of }W^-\text{ on the right)}.
\end{equation}

\begin{lemma}[Theta functionals]\label{ptf:theta}
There is $\kappa=\kappa(S,T)\in\C$ with $|\kappa|=1$ such that $\theta_{W^-}(Af)=\kappa\,\aubar{\theta_W(f)}$
for all $f\in\psHS(W)$ (for $f$ in the domain of $\theta$ of (W0)).
\end{lemma}
\begin{proof}
The rational points of $W$ and $W^-$ are the same set $W(E)$. The functional $f'\mapsto
\aubar{\theta_W(A^{-1}f')}$ on $\psHS(W^-)$ is linear (two antilinear maps) and, by \eqref{ptf:eq-Arho},
invariant under $\psrho^-(x)$ for $x\in W(E)$; by the uniqueness in (W1) it equals $\kappa^{-1}\theta_{W^-}$
for some $\kappa\ne0$ (non-zero since $A$ is bijective and $\theta_W\neq0$). For the modulus: the
functionals are normalised ({\cite[\S5.1.1]{P49}}, {\cite[\S2.3]{P53}}) by the condition that $x\mapsto
\theta(\psrho(x)\psvac)$, $\psvac$ the global vacuum, have mean square $1$ over $W(E)\backslash W(\auA)$.
By the normalisations above $A\psvac=\psvac^-$, so $\theta_{W^-}(\psrho^-(x)\psvac^-)=\theta_{W^-}(A\psrho(x)
\psvac)=\kappa\aubar{\theta_W(\psrho(x)\psvac)}$; taking mean squares, $1=|\kappa|^2$.
\end{proof}

\begin{lemma}[Weil representations of $W$ and $W^-$]\label{ptf:weil}
For $\gamma\in\auU(W)(\auA)=\auU(W^-)(\auA)$,
\begin{equation}\label{ptf:eq-Aomega}
A\,\omega_W(\gamma)\,A^{-1}=\upsilon(\det\gamma)\,\omega_{W^-}(\gamma).
\end{equation}
\end{lemma}
\begin{proof}
\emph{Step 1: a character.} By (W1) for $W$ and $W^-$ and \eqref{ptf:eq-Arho}, $A\omega_W(\gamma)A^{-1}$
and $\omega_{W^-}(\gamma)$ both conjugate $\psrho^-(x)$ into $\psrho^-(\gamma x)$; so
$\omega_{W^-}(\gamma)^{-1}A\omega_W(\gamma)A^{-1}$ commutes with the irreducible Heisenberg representation
and is a scalar $\lambda(\gamma)$. As both $\gamma\mapsto A\omega_W(\gamma)A^{-1}$ and $\omega_{W^-}$ are
representations (genuine, by the splitting), $\lambda$ is a character, and it is unitary ($A$ is
isometric up to a constant). It is a product $\prod_v\lambda_v$ of characters of the local groups, by
the same argument place by place.

\emph{Step 2: $\lambda=(\upsilon\circ\det)^s$ for a universal sign $s$.} The splitting of $\omega_W$ is
Kudla's splitting attached to the character $\upsilon$ ({\cite[\S5.1.1]{P49}}, \cite[Theorem~3.1]{KudlaSplitting}); we
use three of its standard properties: (K1) it is defined by scalar functions on $\auU(W)$ built from
$\psi$ (Weil indices and the character $\upsilon$ evaluated on Rao's invariants), so that replacing
$(\psi,\upsilon)$ by $(\aubar\psi,\aubar\upsilon)$ conjugates these scalars; (K2) the Heisenberg and Weil
representations of $(W,\aubar\psi)$ are those of $(W^-,\psi)$; (K3) for a unitary character $\mu$ of
$\auU(1)$, replacing $\upsilon$ by $\upsilon\cdot(\mu\circ(x\mapsto x/\aubar x))$ multiplies the splitting by
$\mu(\det\gamma)^{s}$ with a sign $s=\pm1$ that does not depend on $W$ (\cite[p.~385 and Lemma~3.5]{KudlaSplitting}, \cite[(1.7)--(1.8)]{HKS}). By (K1) and (K2), $A\omega_W(\gamma)A^{-1}$ is the Weil representation of $W^-$ defined
with $\aubar\upsilon=\upsilon^{-1}$. Since $\upsilon$ restricted to $\auA_{E^+}^\times$ is quadratic,
$\upsilon(\aubar x)=\upsilon(x)^{-1}$, hence $\upsilon(x/\aubar x)=\upsilon(x)^2$ and $\upsilon^{-1}=\upsilon\cdot
(\mu\circ(x\mapsto x/\aubar x))$ with $\mu=\upsilon|_{\auU(1)}^{-1}$. By (K3), $\lambda=(\upsilon\circ\det)^{-s}$.

\emph{Step 3: the sign, from (W2).} It remains to see $-s=1$; as $s$ is universal it suffices to check
$\lambda_v(p)=\upsilon_v(\det p)$ for one place $v$, one $W$ with an isotropic line $I\subset W_v$ and one
$p$ with $\upsilon_v(\det p)\neq\upsilon_v(\det p)^{-1}$. Let $I$ be an isotropic subspace of $W_v$, $P_I$ its
stabiliser, and let $p\in P_I$ act on $I$ by $a\in\auGL(I)$ and trivially on $I^\perp/I$; then $p$ acts on
$W_v/I^\perp\cong\aubar I^{\,*}$ by the inverse conjugate transpose of $a$, so $\det p=\det a/\aubar{\det a}$
and $\upsilon_v(\det p)=\upsilon_v(\det a)^2$. The subspace $I$ and the group $P_I$ are the same for
$W_v^-$, $I^\perp/I$ is $W_{I,v}$ with its form negated, and the zero terms $\psZ_I$, $\psZ^-_I$ of (W2)
are characterised up to a scalar as the linear maps to $\psHS(W_I)$ resp.\ $\psHS(W_I^-)$ that intertwine
the Heisenberg groups of $I^\perp$ and kill $\psrho(i)-1$, $i\in I_v$ (in the mixed model they are the
evaluation at $0$ of functions on $I^*$, and the coinvariants of $I_v$ form one copy of
$\psHS(W_I)$). Let $A_I\colon\psHS_v(W_I)\to\psHS_v(W_I^-)$ be the antilinear intertwiner as above. Then
$A_I^{-1}\psZ^-_IA$ is linear with the same characterising properties as $\psZ_I$ (by
\eqref{ptf:eq-Arho} for $W$ and for $W_I$), so $\psZ^-_IA=c\,A_I\psZ_I$ with $c\neq0$. By (W2), $\psZ_I
\omega(p)=\chi_I(p)\psZ_I$ and $\psZ^-_I\omega^-(p)=\chi_I(p)\psZ^-_I$, with the same
$\chi_I(p)=\upsilon_v(\det a)^{-1}|\det a|^{1/2}$ (the formula does not involve the form). Hence
\[
\lambda_v(p)\chi_I(p)\,\psZ^-_IA=\psZ^-_I\,\lambda_v(p)\omega^-(p)A=\psZ^-_IA\,\omega(p)=cA_I\psZ_I\omega(p)
=cA_I\chi_I(p)\psZ_I=\aubar{\chi_I(p)}\,\psZ^-_IA,
\]
so $\lambda_v(p)=\aubar{\chi_I(p)}/\chi_I(p)=\upsilon_v(\det a)^2=\upsilon_v(\det p)$. Take $v=\tau$ infinite,
$W_\tau$ indefinite and $a=\mathrm{diag}(z,1,\dots)$ with $z\in\C^\times$, $z^4\neq|z|^4$: as
$\upsilon_\tau(z)=z/|z|$, $\upsilon_\tau(\det p)^2=(z/|z|)^4\ne1$. So $-s=1$.
\end{proof}

\begin{remark}\label{ptf:rem-weil}
Step 3 alone proves \eqref{ptf:eq-Aomega} on every parabolic subgroup $P_I(E^+_v)$ (for the elements
acting trivially on $I^\perp/I$, and by induction on the Witt index for all of $P_I$); Step 2 is used
only to extend it to the whole group, in particular at places where $W$ is anisotropic. The sign check
is a consistency check of the normalisation (W2) of the splitting used in the paper with that of
Kudla.
\end{remark}

\begin{lemma}[Negation and complex conjugation]\label{ptf:conj}
Let $S,T$ be spaces of ranks $n_S,n_T$ and $S^-$ as above. There is $\epsilon=\epsilon(S,T)\in\C$,
$|\epsilon|=1$, such that for all automorphic characters $\alpha,\beta$ of $\auU(1)$, all moduli $\pcm$, all
functions $\Psi$ on $\pcMat_{n_T\times n_S}(\pcR_\pcm)$, all polynomials $P$ (more generally Fock vectors
with finite moments) and all $g\in\auU(S)(\auA)=\auU(S^-)(\auA)$, $h\in\auU(T)(\auA)$,
\begin{equation}\label{ptf:eq-conj}
\aubar{\theta^{S,T}_{\alpha,\beta}\bigl(g,h;[P,\Psi]\bigr)}=\epsilon\,\theta^{S^-,T}_{\alpha^{-1}\upsilon^{n_T},\,
\beta^{-1}\upsilon^{n_S}}\bigl(g,h;[\aubar P,\aubar\Psi]\bigr).
\end{equation}
Here the first index is the character of the source group ({\cite[\S3.1.1]{P53}}): the source
character is shifted by $\upsilon^{\mathrm{rk}\,T}$ and the target character by $\upsilon^{\mathrm{rk}\,S}$.
If instead $S^-$ is given the frame $(\varepsilon e_i)$ at the finite places, $\varepsilon\in\widehat\auO_E^\times$
with $\varepsilon\aubar\varepsilon=-1$, then $\aubar\Psi$ is replaced by $x\mapsto\aubar{\Psi(u^{-1}x)}$, $u$ the
reduction of $\varepsilon$.
\end{lemma}
\begin{proof}
Put $f=[P,\Psi]$ and $\gamma=(w\mapsto hwg^{-1})$. By Lemma~\ref{ptf:theta}, then
Lemma~\ref{ptf:weil} and \eqref{ptf:eq-detonhom}, and \eqref{ptf:eq-Adatum},
\begin{align*}
\aubar{\theta^{S,T}_{\alpha,\beta}(g,h;f)}&=\aubar{\alpha(\det g)^{-1}\beta(\det h)}\;
\aubar{\theta_W(\omega_W(\gamma)f)}=\alpha(\det g)\beta(\det h)^{-1}\kappa^{-1}\theta_{W^-}\bigl(A\omega_W(\gamma)
f\bigr)\\
&=\alpha(\det g)\beta(\det h)^{-1}\kappa^{-1}\upsilon(\det h)^{n_S}\upsilon(\det g)^{-n_T}\,
\theta_{W^-}\bigl(\omega_{W^-}(\gamma)[\aubar P,\aubar\Psi]\bigr)\\
&=\kappa^{-1}\bigl(\alpha^{-1}\upsilon^{n_T}\bigr)(\det g)^{-1}\bigl(\beta^{-1}\upsilon^{n_S}\bigr)(\det h)\,
\theta_{W^-}\bigl(\omega_{W^-}(\gamma)[\aubar P,\aubar\Psi]\bigr),
\end{align*}
which is \eqref{ptf:eq-conj} with $\epsilon=\kappa^{-1}=\aubar\kappa$; it depends only on $S,T$. (The
characters take values of modulus one, so conjugation inverts them.) For the frame $(\varepsilon e_i)$:
the coordinates of a vector of $W^-$ at a finite place are then $\varepsilon$ times those for the old
frame, and (A$_\pcm$) becomes $A_v\Psi(x)=\aubar{\Psi(\varepsilon^{-1}x)}$, which gives the stated datum;
the other steps are unchanged (the lattice $\varepsilon L_v=L_v$).
\end{proof}

\begin{remark}\label{ptf:rem-conj}
(i) Lemma~\ref{pte:conj} is the case of Lemma~\ref{ptf:conj} used in Proposition~\ref{pte:carrier}, with
$S=T=I_m$ (equal ranks). (ii) With the frames of {\cite[\S2.3]{P53}} (rational bases), $S^-$ has the same frame
as $S$, no unit $u$ appears, and this is the version used in Proposition~\ref{pte:carrier}.
\end{remark}

\subsubsection{The archimedean action on polynomial vectors}\label{ptf:ss-fock}

This section fixes the transpose conventions needed in Step~2 of the proof of
Corollary~\ref{pte:explicit} and proves the formulas used there.

\emph{Conventions.} Let $S,T$ be spaces with sign vectors $s\in\{\pm1\}^{n_S}$, $s'\in\{\pm1\}^{n_T}$ (the
same at all $\tau$), and $n_S^\pm$, $n_T^\pm$ the numbers of signs $\pm1$. At $\tau\in\auPl$ a vector of
$\auHom(S,T)\otimes_{E^+,\tau}\R$ is a matrix $w\in\pcMat_{n_T\times n_S}(\C)$ (rows indexed by $T$,
columns by $S$, in the frames), and $(g,h)$ acts by $w\mapsto hwg^{-1}$. The entry $(a,b)$ has
\emph{product sign} $s'_as_b$. The \emph{Fock variable} $X_{ab}$ of the entry $(a,b)$ is $\aubar{w_{ab}}$ if
$s'_as_b=1$ and $w_{ab}$ if $s'_as_b=-1$ ({\cite[\S5.1.1]{P49}}: the variable of the Fock space is the
conjugate of the coordinate; \eqref{po:eq-poly}). For a polynomial $P$ in the $X_{ab}$ (all $\tau$
together) its Fock vector is the function $w\mapsto P(X(w))\,G(w)$ on $\prod_\tau\pcMat_{n_T\times
n_S}(\C)$, $G$ the real Gaussian (the vacuum); we identify $P$ with it. For a block-diagonal $g\in\auU(S)_
\tau$, $g=\mathrm{diag}(g^+,g^-)$ along the signs, $g^\pm\in\auU(n_S^\pm)$, and similarly $h$. A character
$\alpha$ of $\auU(1)$ of weight $(\mu_\tau)$ has archimedean components $\alpha_\tau(z)=z^{\mu_\tau}$ on
$\auU(1)(\R)=\{|z|=1\}$. We write $\omega_\infty$ for the archimedean part of $\omega_{\alpha,\beta}$ (with the
archimedean components of the characters).

\emph{Normalisation of the compact action.} For $\gamma$ in the maximal compact subgroup of
$\auU(\auHom(S,T))_\tau$ (the elements preserving the decomposition by product signs) the Weil
representation acts on Fock vectors, read as functions of $w$, by
\begin{equation}\label{ptf:eq-compact}
(\omega(\gamma)F)(w)=\Delta(\gamma)\,F(\gamma^{-1}w),\qquad\Delta(\gamma)=\aubar{\det(\gamma^-)},
\end{equation}
$\gamma^-$ the block of $\gamma$ on the entries of negative product sign. Indeed $\gamma$ preserves
the Gaussian, so $\omega(\gamma)$ and $F\mapsto F\circ\gamma^{-1}$ both conjugate $\psrho(x)$ into
$\psrho(\gamma x)$ and differ by a scalar character $\Delta$ of the compact group; that $\Delta(\gamma)=
\aubar{\det\gamma^-}=\det(\gamma^-)^{-1}$ is the archimedean normalisation of the splitting defined
with $\upsilon$, $\upsilon_\tau(z)=z/|z|$ ({\cite[\S5.1.1]{P49}}): by {\cite[Lemma~2.20]{P49}},
$\Delta(\gamma)=\det(\gamma^+)^{e_+}\det(\gamma^-)^{e_-}$ with universal $e_\pm$ and $e_++e_-=-1$, and
$e_+=0$, $e_-=-1$ by {\cite[Proposition~2.21]{P49}} (an explicit computation with Kudla's splitting on a
doubled line, which uses the signs of $\psi_\infty$ and $\jmath$ fixed in {\cite[\S5.1.1]{P49}}). It is consistent with
Lemma~\ref{ptf:weil}: for $W^-$ the negative entries are the positive entries of $W$, so $\Delta^-(\gamma)=
\aubar{\det\gamma^+}$, and $\aubar{\Delta(\gamma)}=\det\gamma^-=\det\gamma\cdot\aubar{\det\gamma^+}=\upsilon_\tau(\det\gamma)
\Delta^-(\gamma)$, which is \eqref{ptf:eq-Aomega} at $\tau$ (with (A$_\infty$), $A_\tau F=\aubar F$ on
functions of $w$).

\begin{lemma}[Compact action on polynomials]\label{ptf:fock}
Let $\alpha,\beta$ have weights $(a_\tau)$, $(b_\tau)$, and let $g\in\auU(S)_\infty$, $h\in\auU(T)_\infty$ be
block diagonal at every $\tau$. For every polynomial $P$,
\[
\omega_\infty(g,h)P=c(g,h)\cdot P\circ\sigma_{g,h},\qquad c(g,h)=\prod_\tau\det(h^+_\tau)^{b_\tau-n_S^-}
\det(h^-_\tau)^{b_\tau-n_S^+}\det(g^+_\tau)^{n_T^--a_\tau}\det(g^-_\tau)^{n_T^+-a_\tau},
\]
where $\sigma_{g,h}$ is the linear substitution of the Fock variables which, writing $X_\tau=
\bigl(\begin{smallmatrix}X^{++}&X^{+-}\\X^{-+}&X^{--}\end{smallmatrix}\bigr)$ by the signs of the rows ($T$) and the
columns ($S$), is
\[
X^{++}\mapsto(h^+)^{\mathsf t}X^{++}(g^+)^{-\mathsf t},\quad X^{--}\mapsto(h^-)^{\mathsf t}X^{--}(g^-)^{-\mathsf t},
\quad X^{+-}\mapsto(h^+)^{-1}X^{+-}g^-,\quad X^{-+}\mapsto(h^-)^{-1}X^{-+}g^+
\]
at each $\tau$. Equivalently: on the variables of positive product sign, $X\mapsto h^{\mathsf t}Xg^{-\mathsf t}$; on
those of negative product sign, $X\mapsto h^{-1}Xg$.
\end{lemma}
\begin{proof}
Put $\gamma=(w\mapsto hwg^{-1})$, which preserves the product signs. By \eqref{ptf:eq-compact},
$(\omega_\infty(g,h)P)(w)=\alpha_\infty(\det g)^{-1}\allowbreak \beta_\infty(\det h)\allowbreak \Delta(\gamma)\allowbreak P(X(h^{-1}wg))\allowbreak G(w)$.
\emph{Substitution.} For an entry of negative product sign $X(h^{-1}wg)_{ab}=(h^{-1}wg)_{ab}$, and the
block structure gives $X^{+-}\mapsto(h^+)^{-1}X^{+-}g^-$, $X^{-+}\mapsto(h^-)^{-1}X^{-+}g^+$. For an entry of positive
product sign $X(h^{-1}wg)_{ab}=\aubar{(h^{-1}wg)_{ab}}=(\aubar{h^{-1}}\,\aubar w\,\aubar g)_{ab}$, and for unitary
blocks $\aubar{h^{-1}}=\aubar{h^*}=h^{\mathsf t}$, $\aubar g=g^{-\mathsf t}$; this gives the two diagonal blocks.
\emph{Scalar.} The negative block $\gamma^-$ consists of $X^{+-}\mapsto h^+X^{+-}(g^-)^{-1}$ on
$\pcMat_{n_T^+\times n_S^-}$ and $X^{-+}\mapsto h^-X^{-+}(g^+)^{-1}$ on $\pcMat_{n_T^-\times n_S^+}$ (in the coordinates
$w$), so by \eqref{ptf:eq-detonhom}
\[
\det\gamma^-=\det(h^+)^{n_S^-}\det(g^-)^{-n_T^+}\det(h^-)^{n_S^+}\det(g^+)^{-n_T^-},
\]
and $\Delta(\gamma)$ is its inverse (all determinants have modulus one). With $\beta_\infty(\det h)=\prod_\tau
\det(h^+_\tau)^{b_\tau}\det(h^-_\tau)^{b_\tau}$ and $\alpha_\infty(\det g)^{-1}=\prod_\tau\det(g^+_\tau)^{-a_\tau}\det(g^-_\tau)^{-a_\tau}$
the product is $c(g,h)$.
\end{proof}

\begin{corollary}[The formulas of Step~2 of Corollary~\ref{pte:explicit}]\label{ptf:SDlaw}
In the setting of \S\ref{pte:ss-conv} let the archimedean weight of $\alpha$ be $(l-\#\iota(\tau))_\tau$ and
that of $\beta$ be $0$ (Step~0 of the proof of Corollary~\ref{pte:explicit} allows this). For
$k\in\prod_\tau\auU(m)$, write also $k$ for $\mathrm{diag}(k,1_l)\in\auU(W)_\infty$ (positive block). Then:
\begin{enumerate}
\item[(a)] {\emergencystretch=6em on $\auHom(W,I_a)$ (characters $\beta$ on $\auU(W)$, $\alpha$ on $\auU(I_a)$), for $u\in\auU(I_a)_\infty$,
$\omega_\infty(k,u)\pteSD(\eta)=\allowbreak \prod_\tau\det(k_\tau)^{-\#\iota(\tau)}\pteSD(\eta k^{\mathsf t})$; equivalently
$\omega_\infty(k,u)\pteSD(\eta k^{-\mathsf t})=\prod_\tau\det(k_\tau)^{-\#\iota(\tau)}\pteSD(\eta)$. In particular
$\omega_\infty(1,u)\pteSD(\eta)=\pteSD(\eta)$.\par}
\item[(b)] on $\auHom(V',W)$ (characters $\alpha$ on $\auU(V')$, $\alpha\beta^{-1}$ on $\auU(W)$),
$\omega_\infty(1,k)\pteMP(\eta)=\prod_\tau\det(k_\tau)^{-\#\iota(\tau)}\pteMP(\eta k^{\mathsf t})$.
\item[(c)] for $|z|=1$ and $t_z,\Omega_z$ as in Step~3, $\Omega_z$ acts on polynomials in the variables
$x$ of $\auHom(W,I_a)$ by the substitution multiplying every variable outside $x^\flat$ by $z$ and fixing
$x^\flat$; and $\omega_\infty(1,u)$ acts on polynomials in $x^\flat$ alone by $q(x^\flat)\mapsto\prod_\tau
\det(u_\tau)^{-\#\iota(\tau)}q(u^{\mathsf t}x^\flat)$.
\end{enumerate}
Here $x=(x^+\mid x^-)$ are the Fock variables of $\auHom(W,I_a)$ (all of $I_a$ is positive, so $x^+$ has
positive and $x^-$ negative product sign), and $y$ those of $\auHom(V',W)$.
\end{corollary}
\begin{proof}
(a) Lemma~\ref{ptf:fock} with $(S,T)=(W,I_a)$, $g=\mathrm{diag}(k,1)$, $h=u$, source weight $0$, target
weight $l-\#\iota(\tau)$, $n_S^-=l$, $n_T^-=0$, $n_T^+=a$: $c=\prod_\tau\det(u_\tau)^{(l-\#\iota(\tau))-l}
\det(k_\tau)^{0}\det(1)^{a}=\prod_\tau\det(u_\tau)^{-\#\iota(\tau)}$, and the substitution is $x^+\mapsto u^{\mathsf t}x^+
k^{-\mathsf t}$, $x^-\mapsto u^{-1}x^-$. As $\pteSD$ involves only $x^+$, and
\[
\begin{pmatrix}\eta_{\tau j}\\u^{\mathsf t}_\tau x^+_\tau k_\tau^{-\mathsf t}\end{pmatrix}=\begin{pmatrix}1_p&0\\0&u_\tau^{\mathsf t}
\end{pmatrix}\begin{pmatrix}\eta_{\tau j}k_\tau^{\mathsf t}\\x^+_\tau\end{pmatrix}k_\tau^{-\mathsf t},
\]
each factor of $\pteSD(\eta)$ is multiplied by $\det(u_\tau)\det(k_\tau)^{-1}$ and $\eta$ becomes $\eta k^{\mathsf t}$;
the powers $\det(u_\tau)^{\#\iota(\tau)}$ cancel $c$. Replacing $\eta$ by $\eta k^{-\mathsf t}$ gives the
equivalent form, and $k=1$ the last assertion.
(b) Lemma~\ref{ptf:fock} with $(S,T)=(V',W)$, $g=1$, $h=\mathrm{diag}(k,1)$, target weight
$l-\#\iota(\tau)$ (that of $\alpha\beta^{-1}$), $n_S^-=l$: $c=\prod_\tau\det(k_\tau)^{(l-\#\iota(\tau))-l}\det(1)^{\cdots}=
\prod_\tau\det(k_\tau)^{-\#\iota(\tau)}$, and $y^{++}\mapsto k^{\mathsf t}y^{++}$, $y^{--}\mapsto y^{--}$ (the other blocks do
not enter $\pteMP$). Then $\det(\eta_{\tau j}k^{\mathsf t}_\tau y^{++}_\tau)$ is the factor of $\pteMP(\eta k^{\mathsf t})$.
(c) Lemma~\ref{ptf:fock} with $g=t_z$ (positive block; $t_z$ is $z^{-2}$ on the first $p$ and $z^{-1}$ on the
last $a$ positive indices) and $h=z^{-1}1_a$: $c=\prod_\tau(z^{-a})^{-\#\iota(\tau)}=z^{aN_\iota}$, cancelled by
the factor $z^{-aN_\iota}$ of $\Omega_z$; the substitution is $x^+\mapsto z^{-1}x^+t_z^{-1}$ (the column of a first-$p$
index is multiplied by $z^{-1}z^2=z$, that of a last-$a$ index by $z^{-1}z=1$) and $x^-\mapsto zx^-$. With
$g=1$, $h=u$ the substitution on $x^\flat$ is $x^\flat\mapsto u^{\mathsf t}x^\flat$ and $c=\prod_\tau\det(u_\tau)^{-\#\iota(\tau)}$.
\end{proof}

\begin{remark}\label{ptf:rem-fock}
In the conventions above the transpose in \eqref{pte:eq-SDlaw} is forced: the Fock variables of
positive product sign are conjugate coordinates, and conjugation turns $h^{-1}$ and $g$ into
$h^{\mathsf t}$ and $g^{-\mathsf t}$. With these formulas Steps~3--5 of the proof of
Corollary~\ref{pte:explicit} hold verbatim: Step~3 uses (c), Step~4 uses (a) with $\eta=\eta^0$ and (b),
and the cancellation of the powers of $z$ and $\det k_\tau$ there is unchanged. The weight hypothesis on
$\beta$ is harmless: by \eqref{pte:eq-twist} changing $\beta$ multiplies both sides of
\eqref{pte:eq-explicit} by the same character of $\det g$ (Step~0).
\end{remark}

\subsubsection{The first factors}\label{ptf:ss-ff}

We use the notation of \S\ref{pte:ss-third} and \S\ref{pte:ss-left}: $\mathsf W=I_{3m}\oplus W'$,
$W'=W\oplus(E^j)^-\oplus E^{3m-j}\oplus W^\flat$, $V''=I_{N''}$, $W=W_j$ with $r_+$ positive and $r_-=m-j$
negative basis vectors. Put
\[
\mathsf C=I_{3m}\oplus(E^j)^-\oplus E^{3m-j}\oplus W^\flat,\qquad \mathsf W\cong W\oplus\mathsf C
\]
(reordering the summands; the point $(h',g,1,1,h^\flat_i)$ becomes $(g,c_i(h'))$ with
$c_i(h')=(h',1,1,h^\flat_i)\in\auU(\mathsf C)(\auA)$, whose archimedean component is $1$). On
$\auHom(\mathsf W,V'')\otimes_\tau\R$ the positive columns are those of $W$ (the $N''\times r_+$ block, with
Fock variables $X_\tau$, columns in the numbering of \eqref{po:eq-poly}) and those of $\mathsf C$ (Fock
variables $Z_\tau$, $N''\times(2m+N''-r_+)$); the negative columns are those of $(E^j)^-$ and of $W$, which
$\pteSD$ does not involve. In this notation $x^+_\tau=(X_\tau\mid Z_\tau)$, and $\pteSD(\eta_i)=\prod_\tau\prod_{k\in
\iota(\tau)}\det\binom{\eta_{i,\tau k}}{X_\tau\ Z_\tau}$, with $\eta_{i,\tau k}=(\eta^X_{\tau k}\mid\eta^Z_{\tau k})$ of size
$2m\times(2m+N'')$.

\emph{The shape of the constant rows.} By the construction of the choices (\S\ref{pte:ss-third}: the rows of
$\eta_{i,\tau k}$ are basis vectors of the padded $W$ for a column of $\mu^+_\tau$ of height $<r_+$, and
short rows without component on $W$ for a column of height $r_+$), the matrices $\eta_{i,\tau k}$ have the
following property, which is all that is used:
\begin{itemize}
\item[(H)] the elements $k\in\iota(\tau)$ correspond to the columns of $\mu^+_\tau$; if $h_k$ is the height of
the column of $k$, then $\eta^X_{\tau k}=(0_{2m\times h_k}\mid E_{\tau k})$, where the $r_+-h_k$ columns of
$E_{\tau k}$ are distinct standard basis vectors of $\C^{2m}$ (so $r_+-h_k$ of the rows of $\eta_{i,\tau k}$ are the
basis vectors $e_{h_k+1},\dots,e_{r_+}$ of the positive part of $W$, and the other rows vanish on $X$).
\end{itemize}

\begin{lemma}[Laplace expansion]\label{ptf:laplace}
Under (H), for each $\tau$, $k$ there are polynomials $Q_{\tau k,R}$ in the entries of $Z_\tau$, indexed by
the $h_k$-subsets $R\subseteq\{1,\dots,N''\}$, such that
\[
\det\binom{\eta_{i,\tau k}}{X_\tau\ Z_\tau}=\sum_{R}\det\bigl(X_{\tau;R,[h_k]}\bigr)\,Q_{\tau k,R}(Z_\tau),
\]
$X_{\tau;R,[h_k]}$ the minor of $X_\tau$ on the rows $R$ and the first $h_k$ columns. Consequently there are
unique polynomials $R_\lambda$ in the entries of $Z=(Z_\tau)_\tau$ with
\begin{equation}\label{ptf:eq-SDexp}
\pteSD(\eta_i)=\sum_\lambda\pop''_{u_\lambda}(X)\,R_\lambda(Z),
\end{equation}
where $(u_\lambda)$ is a fixed basis of $W''=W_{\mu^+}$ and $\pop''_u$ is the polynomial \eqref{po:eq-poly} for the
pair $(W,V'')$.
\end{lemma}
\begin{proof}
Let $M$ be the square matrix $\binom{\eta}{X\ Z}$ (we drop $\tau,k,i$). Expand $\det M$ along its first $h=h_k$
columns (generalised Laplace expansion): $\det M=\sum_{R'}\pm\det M_{R',[h]}\det M_{R'^c,[h]^c}$, the sum over
the $h$-subsets $R'$ of the rows. By (H) the constant rows vanish in the first $h$ columns, so only $R'$
consisting of rows of $X$ contribute, and then $M_{R',[h]}=X_{R,[h]}$. In the complementary matrix
$M_{R'^c,[h]^c}$ the $r_+-h$ constant rows $e_{h+1},\dots,e_{r_+}$ each have a single entry $1$, in the columns
$h+1,\dots,r_+$; expanding along them deletes these rows and the columns $h+1,\dots,r_+$ of $X$, up to a sign.
What remains has as rows the constant rows vanishing on $X$ (restricted to the columns of $Z$) and the rows
$R^c$ of $Z$; its determinant $\pm Q_R(Z)$ depends on $Z$ alone. This is the first assertion.

Multiplying over $\tau$ and $k$: $\pteSD(\eta_i)=\sum_{(R_k)}\prod_{\tau,k}\det X_{\tau;R_k,[h_k]}\cdot\prod_{\tau,k}Q_{\tau k,R_k}
(Z_\tau)$. By \eqref{po:eq-ev} and \eqref{po:eq-poly} (all of $V''$ is positive, so only $x^+$ and the factor
$\poT^+$ occur), $\pop''_\xi(X)=\langle\xi,\mathrm{ev}(X)\rangle$ with $\mathrm{ev}(X)=\bigotimes_{\mathfrak c}(X_1\wedge\dots\wedge
X_{h_{\mathfrak c}})=\sum_{(R_{\mathfrak c})}\prod_{\mathfrak c}\det X_{R_{\mathfrak c},[h_{\mathfrak c}]}\,e_{(R_{\mathfrak c})}$ ($X_i$ the columns
of $X$, $e_{(R)}$ the tensor of the wedge basis vectors, an orthonormal basis of $\poT^+$; the Fock variables $X$
are the conjugates $\aubar{x^+}$ of the coordinates, as in \eqref{po:eq-poly}). Hence the map $\xi\mapsto\pop''_\xi$,
$\poT^+\to\C[X]$, is conjugate-linear with image the span $\Sigma_{\mu^+}$ of the products $\prod_{\mathfrak c}\det
X_{R_{\mathfrak c},[h_{\mathfrak c}]}$, and its kernel is $W''^\perp$, because the vectors $\mathrm{ev}(X)$ span $W''$
(definition of $W''$); so $u\mapsto\pop''_u$ is a bijection of $W''$ onto $\Sigma_{\mu^+}$ and $(\pop''_{u_\lambda})$ is a
basis of $\Sigma_{\mu^+}$. The expansion above shows $\pteSD(\eta_i)\in\Sigma_{\mu^+}\otimes\C[Z]$, which gives
\eqref{ptf:eq-SDexp} with unique $R_\lambda$.
\end{proof}

\begin{lemma}[Sums in the source]\label{ptf:sum}
Let $S=S_1\oplus S_2$ be an orthogonal sum of spaces, $T$ a space, $P_\nu$, $\Psi_\nu$ a polynomial and a
function for $\auHom(S_\nu,T)$, and $\alpha,\beta_1,\beta_2$ characters. With $P=P_1P_2$ (variables of the two
column blocks) and $\Psi(x_1\mid x_2)=\Psi_1(x_1)\Psi_2(x_2)$,
\[
\theta^{S,T}_{\alpha,\beta_1\beta_2}\bigl((g_1,g_2),h;[P,\Psi]\bigr)=\theta^{S_1,T}_{\alpha,\beta_1}\bigl(g_1,h;[P_1,\Psi_1]\bigr)\,
\theta^{S_2,T}_{\alpha,\beta_2}\bigl(g_2,h;[P_2,\Psi_2]\bigr).
\]
\end{lemma}
\begin{proof}
$\auHom(S,T)=\auHom(S_1,T)\oplus\auHom(S_2,T)$ orthogonally; the Heisenberg representation, the theta functional
and the datum factor as tensor products ($[P,\Psi]=[P_1,\Psi_1]\otimes[P_2,\Psi_2]$, the vacuum of the sum of
the lattices being the tensor product of the vacua, normalisations compatible), and the Weil representation of
$\auU(\auHom(S_1,T))\times\auU(\auHom(S_2,T))$ is the tensor product, the splittings defined with $\upsilon$ being
compatible with orthogonal sums \cite{KudlaSplitting}; this is the source version of (W7) of
\S\ref{po:ss-glue}. The element $(g,h)$ acts on the two summands by $((g_1,h),(g_2,h))$, and the scalar
$\alpha(\det g)^{-1}\beta_1\beta_2(\det h)$ is $\alpha(\det g_1)^{-1}\beta_1(\det h)\cdot\alpha(\det g_2)^{-1}\beta_2(\det h)$.
\end{proof}

\begin{definition}[First factors]\label{ptf:def-ff}
Let $\nu''=\nu''_\vartheta$ and $\Phi^{\mathsf C}_i$ the function on the columns of $\mathsf C$:
$\Phi^{\mathsf C}_i(x_{\mathsf C})=\psi^{\mathrm{inv}}_{\Lambda_i}(x')\cdot\plone[\text{the columns of }x_{\mathsf C}\text{ on }(E^j)^-\oplus E^{3m-j}
\oplus W^\flat\text{ are those of }1]$, $x'$ the columns on $I_{3m}$ (read modulo $\pi_M^2$). For $h\in
\auU(V'')(\auA_f)$ put
\[
c_{\lambda}(h)=\pteCS_{K_{I_{3m}}}\Bigl(h'\mapsto\mathrm a_{\chi'}(h')\,\theta^{\mathsf C,V''}_{\nu'',\nu''\upsilon^{j}}\bigl(c_i(h'),h;
[R_\lambda,\Phi^{\mathsf C}_i]\bigr)\Bigr),\qquad\varphi''_{\vartheta,i}(h)=\sum_\lambda\aubar{c_\lambda(h)}\,u_\lambda\in W''.
\]
\end{definition}

The definition does not depend on the basis $(u_\lambda)$: by \eqref{ptf:eq-SDexp}, $\varphi''_{\vartheta,i}(h)$ is the
unique $\varphi\in W''$ with $\pop''_{\varphi}(X)=\sum_\lambda c_\lambda(h)\pop''_{u_\lambda}(X)$, and $\sum_\lambda R_\lambda\otimes
\pop''_{u_\lambda}$ is basis-free.

\begin{proposition}[Factorisation and transformation law]\label{ptf:ff}
Let $\theta''_{\vartheta,\Lambda}(g,h;u)=\theta^{W,V''}_{\nu'',\upsilon^{m-j}}(g,h;[\pop''_u,\Psi''_{\vartheta,\Lambda}])$, the kernel of
\S\ref{po:ss-glue} (conjugate-linear in $u$), and let $\Phi_0$ be the datum of Definition~\ref{pte:def-ff}. Then for
$g\in\auU(W)(\auA)$, $h\in\auU(V'')(\auA_f)$:
\begin{enumerate}
\item[(a)] $\displaystyle\pteCS_{K_{I_{3m}}}\Bigl(h'\mapsto\mathrm a_{\chi'}(h')\,\theta^{\mathsf W,V''}_{\nu'',\nu''\upsilon^m}\bigl((h',g,1,1,h^\flat_i),h;
[\pteSD(\eta_i),\Phi_0]\bigr)\Bigr)=\theta''_{\vartheta,\Lambda_i}\bigl(g,h;\varphi''_{\vartheta,i}(h)\bigr)$;
\item[(b)] $\varphi''_{\vartheta,i}(hk)=\chi'_\vartheta(\det\bar k)^{-1}\varphi''_{\vartheta,i}(h)$ for $k\in J'_{\vartheta,i}$, $\bar k$ its
reduction modulo $\pcm_\vartheta$; and $\theta''_{\vartheta,\Lambda_i}(g,hk;u)=\chi'_\vartheta(\det\bar k)^{-1}\theta''_{\vartheta,\Lambda_i}(g,h;u)$;
\item[(c)] $\varphi''_{\vartheta,i}(\gamma h)=\rho''(\gamma_\infty)\varphi''_{\vartheta,i}(h)$ for $\gamma\in\auU(V'')(E^+)$, where
$\rho''$ is the representation of $\auU(V'')_\infty$ on $W''$ defined by $\omega''_\infty(1,u)\pop''_\xi=\pop''_{\rho''(u)\xi}$
(the equivariance of \eqref{po:eq-poly}; this is $\rho_{\mu^+}$ in the conventions of \S\ref{po:ss-glue}).
\end{enumerate}
In particular the summand $h\mapsto\theta''_{\vartheta,\Lambda_i}(g,h;\varphi''_{\vartheta,i}(h))$ is right invariant under
$J'_{\vartheta,i}$ and left invariant under $\auU(V'')(E^+)$.
\end{proposition}
\begin{proof}
(a) The datum $\Phi_0(x_W\mid x_{\mathsf C})=\mathrm{col}(x_W)\cdot\Phi^{\mathsf C}_i(x_{\mathsf C})$ is a product over the
column blocks $W$, $\mathsf C$ (the columns $x''e$ are the columns of $W$; the condition on the other columns
of $x''$ and the factor $\psi^{\mathrm{inv}}_{\Lambda_i}(x')$ concern $\mathsf C$), and $\mathrm{col}$ is the finite part at
$\pcm_\vartheta$ of $\Psi''_{\vartheta,\Lambda_i}$ (Lemma~\ref{pte:column}). By \eqref{ptf:eq-SDexp} and linearity of the
kernel in the polynomial, then Lemma~\ref{ptf:sum} with $\alpha=\nu''$, $\beta_1=\upsilon^{m-j}$, $\beta_2=\nu''\upsilon^j$
($\beta_1\beta_2=\nu''\upsilon^m$),
\[
\theta^{\mathsf W,V''}_{\nu'',\nu''\upsilon^m}\bigl((g,c_i(h')),h;[\pteSD(\eta_i),\Phi_0]\bigr)=\sum_\lambda\theta^{W,V''}_{\nu'',\upsilon^{m-j}}
\bigl(g,h;[\pop''_{u_\lambda},\Psi''_{\vartheta,\Lambda_i}]\bigr)\,\theta^{\mathsf C,V''}_{\nu'',\nu''\upsilon^j}\bigl(c_i(h'),h;[R_\lambda,
\Phi^{\mathsf C}_i]\bigr).
\]
The first factor is $\theta''_{\vartheta,\Lambda_i}(g,h;u_\lambda)$ and does not involve $h'$; the finite class sum over
$h'$ with the weights $\mathrm a_{\chi'}(h')$ therefore produces $\sum_\lambda c_\lambda(h)\theta''_{\vartheta,\Lambda_i}(g,h;u_\lambda)=
\theta''_{\vartheta,\Lambda_i}(g,h;\sum_\lambda\aubar{c_\lambda(h)}u_\lambda)$ by conjugate-linearity.

(b) Let $k\in J'_{\vartheta,i}$ and $A=\bar k$; $A\in H_{\vartheta,i}$, so the columns of $A$ off $e$ are those of
$1$ and $w_iA\equiv w_i\bmod\pi_M^2$; the same holds for $A^{-1}$. By {\cite[Lemma~2.9]{P53}}, right translation
of a kernel of $(S,V'')$ with target character $\beta$ by $k$ replaces the finite datum $\Psi$ by $\beta(\det
k)\Psi(A^{-1}\,\cdot)$. \emph{The $\mathsf C$-factor:} $A^{-1}$ maps the columns of $1$ off $e$ to the columns of
$A^{-1}$ off $e$, which are those of $1$, and $\psi^{\mathrm{inv}}_{\Lambda_i}(A^{-1}x')=\psi^{\mathrm{inv}}_{\Lambda_i}(x')$ because
$\Lambda_iA^{-1}\equiv\Lambda_i\bmod\pi_M^2$; so $\Phi^{\mathsf C}_i(A^{-1}\cdot)=\Phi^{\mathsf C}_i$, the polynomial is unchanged ($k$ is
finite), and $c_\lambda(hk)=(\nu''\upsilon^j)(\det k)\,c_\lambda(h)$. Now $\upsilon$ is unramified at the finite places, so
$\upsilon(\det k)=1$, and $\nu''$ is tame for $\pcm_\vartheta$ with residue character $\chi'_\vartheta$ (the character
$\nu''\upsilon^m$ is, \S\ref{pte:ss-third}), so $\nu''(\det k)=\chi'_\vartheta(\det A)$. Hence $\varphi''_{\vartheta,i}(hk)=
\aubar{\chi'_\vartheta(\det A)}\varphi''_{\vartheta,i}(h)=\chi'_\vartheta(\det A)^{-1}\varphi''_{\vartheta,i}(h)$. \emph{The $W$-factor:}
write $A^{-1}=\bigl(\begin{smallmatrix}B&0\\C&1\end{smallmatrix}\bigr)$ with the indices $e$ first. For $\sigma\in\pcMat_{N''\times
r}$, $(A^{-1}\sigma)'=B\sigma'$ and $\det B=\det A^{-1}$, while $w_iA^{-1}\sigma=w_i\sigma$; so $\mathrm{col}(A^{-1}\sigma)=\chi'_\vartheta
(\det A)^{-1}\mathrm{col}(\sigma)$. With $\upsilon^{m-j}(\det k)=1$ this gives the second assertion. 

(c) Let $\gamma\in\auU(V'')(E^+)$, $c\in\auU(\mathsf C)(\auA_f)$. By (W1) the kernel is invariant under the diagonal
rational point: $\theta^{\mathsf C,V''}(c,(1,\gamma h);f)=\theta^{\mathsf C,V''}(c,(\gamma_\infty^{-1},h);f)$ (the source point $c$ is
unchanged, $\gamma$ acting on $V''$ only), and the archimedean component $\gamma_\infty^{-1}$ acts on the polynomial
through $\omega_\infty(1,\gamma_\infty^{-1})$. By Corollary~\ref{ptf:SDlaw}(a) (applied to $(\mathsf W,V'')$, with $k=1$),
$\omega_\infty(1,u)\pteSD(\eta_i)=\pteSD(\eta_i)$ for $u\in\auU(V'')_\infty$; by Lemma~\ref{ptf:sum} at infinity the left
side is $\sum_\lambda\omega''_\infty(1,u)\pop''_{u_\lambda}\cdot\omega^{\mathsf C}_\infty(1,u)R_\lambda=\sum_\lambda\pop''_{\rho''(u)u_\lambda}\cdot
\omega^{\mathsf C}_\infty(1,u)R_\lambda$. Write $\rho''(u)u_\lambda=\sum_\mu a_{\mu\lambda}(u)u_\mu$; then $\pop''_{\rho''(u)u_\lambda}=
\sum_\mu\aubar{a_{\mu\lambda}(u)}\pop''_{u_\mu}$, and comparing coefficients of the basis $\pop''_{u_\mu}$ (unique by
Lemma~\ref{ptf:laplace}) gives $\sum_\lambda\aubar{a_{\mu\lambda}(u)}\,\omega^{\mathsf C}_\infty(1,u)R_\lambda=R_\mu$, i.e.\
$\omega^{\mathsf C}_\infty(1,u^{-1})R_\mu=\sum_\lambda\aubar{a_{\mu\lambda}(u)}R_\lambda$. With $u=\gamma_\infty$, $c_\mu(\gamma h)=\sum_\lambda\aubar{a_{\mu\lambda}
(\gamma_\infty)}c_\lambda(h)$, so $\varphi''_{\vartheta,i}(\gamma h)=\sum_{\mu,\lambda}a_{\mu\lambda}(\gamma_\infty)\aubar{c_\lambda(h)}u_\mu=\rho''(\gamma_\infty)
\varphi''_{\vartheta,i}(h)$. The last sentence follows from (a), (b), (c) and the invariance of the left side of (a):
under $J'_{\vartheta,i}$ by (b) (the two characters cancel by conjugate-linearity), under rational points by (W1).
\end{proof}

\subsubsection{The character bookkeeping on $V''$}\label{ptf:ss-left}

\begin{proposition}[The left side]\label{ptf:leftside}
With $J'_{\vartheta,i}$ in place of $J_{\vartheta,i}$ in \S\ref{po:ss-glue}, the character $\chi''_\vartheta(k)=\chi'_\vartheta(\det\bar k)$ on
$J'_{\vartheta,i}$, and the first factors of Definition~\ref{ptf:def-ff},
\[
\mathcal L_{\vartheta,i}(g)=\sum_{h\in\Gamma''\backslash\auU(V'')(\auA_f)/J'_{\vartheta,i}}\theta''_{\vartheta,\Lambda_i}\bigl(g,h;\varphi''_{\vartheta,i}(h)\bigr)
=\Xi_{\vartheta,i}(g)\qquad(g\in\auU(W)(\auA)),
\]
and $\varphi''_{\vartheta,i}\in\poA_{\mu^+}(J'_{\vartheta,i},\chi''_\vartheta{}^{-1})$.
\end{proposition}
\begin{proof}
\emph{Characters.} The kernel in $\mathcal L_{\vartheta,i}$ is $\theta^{\mathsf W,V''}_{\nu'',\nu''\upsilon^m}$: source character $\nu''$
on $\auU(\mathsf W)$, target character $\chi'\nu^\flat\upsilon^m=\nu''\upsilon^m$ on $\auU(V'')$. In Lemma~\ref{ptf:sum} (sums in
the source) the source character is common to the summands and the target character is distributed; we
distributed it as $\upsilon^{m-j}$ on the summand $W$ and $\nu''\upsilon^j$ on $\mathsf C$. The kernel $\theta''$ of
\S\ref{po:ss-glue} is the $V''$-factor of $\omega_\mu$ in (W7): by \eqref{po:eq-w5}, $\omega_\mu=\omega'_{a,b}=\omega_{b,a}$
({\cite[(13)]{P53}}) with $a=\upsilon^{r_-}$ on $\auU(V)$ and $b=\mu\upsilon^{\lfloor n/2\rfloor}$ on $\auU(W)$. In (W7) the sum is in
the target $V=V''\oplus V^\circ$, so (by the target analogue of Lemma~\ref{ptf:sum}) the target character
$\upsilon^{r_-}$ is common to the two factors and the source character $b$ is distributed as $\nu''\cdot\upsilon^{m/2-j}$
($\nu''$ on the $V''$-factor, $\upsilon^{m/2-j}$, independent of $\vartheta$, on the $V^\circ$-factor); the identity
$\nu''\upsilon^{m/2-j}=\mu\upsilon^{\lfloor n/2\rfloor}$ follows from the definition of $\nu''_\vartheta$ in \S\ref{pte:ss-third}: with $\mu=\mu^0_j\vartheta$, $n=N''+
\mathrm{rk}\,V^\circ=6r$, $t+k-1=5r+k-1$ even and $r=k+m-1$, one has $j+\frac{t+k-1}2+\frac m2-j=\frac{5r+k-1+m}2=3r=\frac n2$
(compare Lemma~\ref{psv:char}). Hence $\omega''=\omega_{\nu'',\upsilon^{r_-}}$, and $r_-=m-j$. This is exactly the pair of characters of the
$W$-factor in Proposition~\ref{ptf:ff}(a), so the kernel there is the $\theta''_{\vartheta,\Lambda_i}$ of \S\ref{po:ss-glue},
with the datum $\Psi''_{\vartheta,\Lambda_i}$ (Lemma~\ref{pte:column}) and the polynomial $\pop''$ of \eqref{po:eq-poly}.

\emph{Sums.} By Lemma~\ref{pte:rowstab} with $F_0=\psi^{\mathrm{inv}}_{\Lambda_i}$ (invariant under $H_{\vartheta,i}$, as in
Definition~\ref{pte:def-ff}), the inner sum of $\mathcal L_{\vartheta,i}$ is the class sum modulo $J'_{\vartheta,i}$ of the kernel with
datum $\Phi_0$. Exchange it with the finite class sum over $h'$ and apply Proposition~\ref{ptf:ff}(a): $\mathcal L_{\vartheta,i}(g)
=\sum_{h\bmod J'_{\vartheta,i}}\#(\Gamma''\cap hJ'_{\vartheta,i}h^{-1})^{-1}\theta''_{\vartheta,\Lambda_i}(g,h;\varphi''_{\vartheta,i}(h))$, and the stabilisers are
trivial because $J'_{\vartheta,i}$ is free ({\cite[Lemma~14.5]{P46}}). By Proposition~\ref{ptf:ff}(b), (c) the summand is a
function on the double cosets and $\varphi''_{\vartheta,i}$ has the stated type ($\rho''=\rho_{\mu^+}$). This is the definition of
$\Xi_{\vartheta,i}$ with the level $J'_{\vartheta,i}$ and the character $\chi''_\vartheta$.
\end{proof}


\subsubsection{Weights, the constant rows and the character $\chi''_\vartheta$}\label{ptf:ss-close}

\emph{Weights.} A character of $\auU(1)$ of weight $(c_\tau)$ is one with $\mu_\tau(u)=\tau(u)^{c_\tau}$ at the
infinite places ({\cite[Lemma~8.6]{P49}}; this is the meaning of ``weight'' in {\cite[\S14.3]{P46}} and
{\cite[\S3.1.1]{P53}}). This is the convention used in \S\ref{ptf:ss-fock}, so Lemma~\ref{ptf:fock} and
Corollary~\ref{ptf:SDlaw} hold as stated. As a check, $\upsilon$ has weight $1$ ($\upsilon_\tau(z)=z/|z|$, which is
$z$ on $\auU(1)(\R)$).

\begin{lemma}[The constant rows satisfy (H)]\label{ptf:H}
Let the matrices $\eta_{i,\tau k}$ be defined as follows (this is the precise form of the description in
\S\ref{pte:ss-third}). Their columns are the $r_+$ positive indices of $W$, numbered $1,\dots,r_+$ as in
\eqref{po:eq-poly}, followed by the indices of $E^{3m-j}\oplus W^\flat\oplus I_{3m}$. The rows are numbered
$a=0,\dots,2m-1$.
\begin{itemize}
\item For $k$ belonging to a column of $\mu^+_\tau$ of height $h<r_+$, row $a$ is the basis vector $e_{h+a+1}$ of
$W$ if $h+a<r_+$, and otherwise the basis vector number $a$ of $E^{3m-j}$.
\item For $k$ belonging to a column of height $r_+$ (the $\Delta'$ columns), the rows are the $2m$ short rows
selected by $\varrho_i$, read at $\tau$ and placed in the columns of $W^\flat$, with $0$ elsewhere.
\end{itemize}
Then (H) holds.
\end{lemma}
\begin{proof}
For a column of height $r_+$, $\eta^X_{\tau k}=0$, which is (H) with $h_k=r_+$. For height $h<r_+$, the rows with
$h+a<r_+$, that is $a=0,\dots,r_+-h-1$, are $e_{h+1},\dots,e_{r_+}$, each exactly once. The remaining rows lie in
$E^{3m-j}$ and vanish on the columns of $X$. So $\eta^X_{\tau k}=(0_{2m\times h}\mid E_{\tau k})$, where $E_{\tau k}$
has as columns the distinct basis vectors $e_0,\dots,e_{r_+-h-1}$ of $\C^{2m}$. This uses $r_+-h\le2m$, i.e.\ every column of $\mu^+_\tau$ of height $h<r_+$ has $h\ge r_+-2m$, which we now
derive from the diagram. By Lemma~\ref{plf:match} and its proof, for $h<r_+$ the number of columns of height $h$ is
$c^+_\tau(h)=\varsigma_\tau(h-1)-\varsigma_\tau(h)-1$, where $\varsigma_\tau(0)>\varsigma_\tau(1)>\cdots$ are the entries of
$\varsigma_\tau$ on the positive coordinates of $W_j$ in decreasing order. By {\cite[\S14.3]{P46}} these are the
$k_0-1$ entries $\frac{k_0}2-1,\frac{k_0}2-2,\dots,1-\frac{k_0}2$ followed by $-\kappa(\tau)-a\Delta$, $1\le a\le j$ (so $r_+=k_0-1+j$), where $k_0$ denotes the integer weight of {\cite[\S14.3]{P46}} (written $k_0$ here to keep $k$ for the elements of $\iota(\tau)$).
Hence $\varsigma_\tau(h-1)-\varsigma_\tau(h)=1$ and $c^+_\tau(h)=0$ for $1\le h\le k_0-2$. So a column of height $h<r_+$ has
$h\ge k_0-1$, and $r_+-h\le r_+-(k_0-1)=j\le m\le2m$.
\end{proof}

\begin{lemma}[The character $\chi''_\vartheta$]\label{ptf:chi}
Let $k\in J'_{\vartheta,i}$. Then the character $\chi''_\vartheta(k)$ of \S\ref{po:ss-glue} (the value of $\chi_\vartheta$ on
$\pog_i^{-1}\mathrm{diag}(k,k^\circ)\pog_i$) equals $\chi'_\vartheta(\det\bar k)$.
\end{lemma}
\begin{proof}
\emph{$\chi_\vartheta$.} Let $k\in K_\vartheta$. By \eqref{po:eq-w5} with $g=1$, $\omega_\mu(1,k)\Phi_\vartheta(T)=\eta_w(1,k)\Phi_\vartheta(k^{-1}T)$,
and $\eta_w(1,k)=1$ because $\eta_w$ is unramified and $k$ is integral. At $w\in\mathsf M$ the factor
$\plone[T\equiv\iota_0]$ is invariant, since $k\iota_0\equiv\iota_0$. At $w\in\mathsf Q$, $k\bmod w$ is block lower
triangular for the blocks $r+t$, so $k^{-1}=\bigl(\begin{smallmatrix}A&0\\ *&*\end{smallmatrix}\bigr)$ modulo $w$ with
$A=k_{11}^{-1}$, where $k_{11}$ is the upper left $r\times r$ block. Hence $(k^{-1}T)'=AT'$ and
$\psi_w(\det(k^{-1}T)')=\psi_w(\det k_{11})^{-1}\psi_w(\det T')$. Therefore
$\chi_\vartheta(k)=\prod_{w\in\mathsf Q}\psi_w(\det k_{11}\bmod w)$.

\emph{$\chi''_\vartheta$.} At $\mathsf Q$ the element $\pog_i$ is $1$ (proof of Lemma~\ref{po:slice}). So the $\mathsf Q$-component of
$\pog_i^{-1}\mathrm{diag}(k,k^\circ)\pog_i$ is $\mathrm{diag}(k,k^\circ)$, whose upper left $r\times r$ block is the block $k_{11}$
of $k$ on the first $r$ indices of $V''$, which are the indices $e$. Hence $\chi''_\vartheta(k)=\prod_{w\in\mathsf Q}\psi_w(\det
k_{11})$.

\emph{Comparison.} For $k\in J'_{\vartheta,i}$ and $w\in\mathsf Q$, the columns of $\bar k_w$ off $e$ are those of $1$, so
$\bar k_w=\bigl(\begin{smallmatrix}k_{11}&0\\ *&1\end{smallmatrix}\bigr)$ and $\det\bar k_w=\det k_{11}$. At $w\in\mathsf M$,
$k\equiv1\bmod\pi_M$, so the factor $\omega^j(\det\bar k_w)$ of $\chi'_\vartheta$ is $1$. The residue character of $\chi'$ at
$w\in\mathsf Q$ is $\psi_w$: here $\psi_w$ is the tame part at $w$ of $\mu=\mu^0_j\vartheta$ (\S\ref{po:ss-set}), which is
that of $\vartheta$ because $\mu^0_j$ is unramified outside the chosen places ({\cite[\S14.3]{P46}}); and
$\chi'=\chi_0^{\,j}\vartheta$ with $\chi_0$ of trivial tame part off $\mathsf M$ (\S\ref{po:ss-hecke}), so the residue
character of $\chi'$ at $w\in\mathsf Q$ (\S\ref{pte:ss-third}) is also the tame part of $\vartheta$ at $w$. Hence $\chi'_\vartheta(\det\bar k)=\prod_{w\in\mathsf Q}\psi_w(\det k_{11})=
\chi''_\vartheta(k)$.
\end{proof}

By Lemma~\ref{ptf:chi}, the character $\chi''_\vartheta$ of \S\ref{po:ss-glue} restricted
to $J'_{\vartheta,i}$ is the character of Proposition~\ref{ptf:ff}(b), and Proposition~\ref{ptf:leftside} holds with the
definition of \S\ref{po:ss-glue} unchanged.



\section{Classes, first factors and periods of false lifts}\label{C2:g3}
\providecommand{\auO}{\mathcal O}
\providecommand{\auU}{\mathrm U}
\providecommand{\auGL}{\mathrm{GL}}
\providecommand{\auHom}{\operatorname{Hom}}
\providecommand{\auA}{\mathbb A}
\providecommand{\auF}{\mathbb F}
\providecommand{\auPl}{\Phi}
\providecommand{\auNm}{\operatorname{Nm}}
\providecommand{\auAut}{\operatorname{Aut}}
\providecommand{\aubar}[1]{\overline{#1}}
\providecommand{\pao}{\mathfrak o}
\providecommand{\psvac}{\mathbf v}
\providecommand{\plone}{\mathbf 1}
\providecommand{\plgb}[2]{\|#2\|_{#1}}
\providecommand{\psHS}{\mathcal H}
\providecommand{\psrho}{\varrho}
\providecommand{\pzF}{\mathfrak F}
\providecommand{\poA}{\mathcal A}
\providecommand{\poAA}{A_\ast}
\providecommand{\pop}{\mathsf p}
\providecommand{\pog}{\mathsf g}
\providecommand{\plfK}{\mathsf K}
\providecommand{\pkm}{\mathsf m}

\subsection{One lift: classes, the type of the descent, first factors}\label{plf:sec}

This section proves two of the statements that \S\ref{po:sec} isolates: Lemma~\ref{po:classes}
(numbers of classes) and the statement (L4) (the descent of a form on $V_j$ lies in $\poA(W_j)$). It
also proves the kernel scales (Lemma~\ref{plf:KS}) and the count of choices (Lemma~\ref{plf:I}), which are
used in \S\ref{pfq:sec}, where
Lemma~\ref{po:ff} is proved (Theorem~\ref{pfq:ff}). We keep the notation of \S\ref{po:sec}: $1\le j\le m$
is fixed, $W=W_j$, $V=V_j$, $n=r+t=6r$, $r_+=k+j-1$, $r_-=m-j$, $\poAA=fmb(\log b)^2$.
Propositions~\ref{po:periods} and~\ref{po:span} are proved in \S\ref{ppe:sec} (Theorem~\ref{ppe:main}) and
\S\ref{psp:sec} (Proposition~\ref{psp:main}, with Theorem~\ref{psv:main}).

\subsubsection{Numbers of classes}\label{plf:ss-classes}

\begin{lemma}[Lemma~\ref{po:classes}]\label{plf:classes}
Let $X=E^{n_X}$ be a standard definite space with $n_X\le2n_m$, $G=\auU(X)$, $\Gamma=G(E^+)$, and
let $J$ be a subgroup of $G_f$ containing the principal congruence subgroup $K(\pi_M^2)$ of level
$\pi_M^2$ in the stabiliser $K^\circ$ of $\widehat\auO_E^{\,n_X}$. Then $\#(\Gamma\backslash G_f/J)\le
\exp(2x^{26})\le e^b$.
\end{lemma}
\begin{proof}
The proof of {\cite[Proposition~7.20]{P49}} does not use the hypothesis $n\le n_m$ except in its last
paragraph (sizes): {\cite[Lemma~14.21]{P46}} holds for every rank, and the counting and the index
of the level are valid for every $n$ and every integer $\mathsf G\ge1$. Since
$\pi_M\aubar{\pi_M}=M$, the integer $M^2$ lies in $\pi_M^2\auO_E$, so $K(M^2)\subseteq K(\pi_M^2)
\subseteq J$, and {\cite[Proposition~7.20]{P49}} with $\mathsf G=M^2$ gives
\[
\#(\Gamma\backslash G_f/J)\le(2\mathsf B+1)^{4fn_X^2}M^{4fn_X^2},\qquad\mathsf B=n_X^2C_{\mathrm{sb}}^2
\cdot2f\sqrt d,\quad C_{\mathrm{sb}}=6q_1d(fn_X)^2 .
\]
Sizes: $f\le x^{12}$, $n_X\le2n_m\le80x^6$, and $k=r+1-m\ge x^5/2$, so $n_X/k\le160x$ and, as in
the proof of {\cite[Proposition~7.20]{P49}}, $C_{\mathrm{sb}}=6q_1d(fk)^2(n_X/k)^2\le x^{100}\cdot
(160x)^2\le x^{103}$, $\mathsf B_0=2f\sqrt d\le C_{\mathrm{sb}}$, $n_X\le C_{\mathrm{sb}}$ and $2\mathsf
B+1\le C_{\mathrm{sb}}^6$. Hence $4fn_X^2\log(2\mathsf B+1)\le4x^{12}\cdot6400x^{12}\cdot618\log x\le
x^{26}$, and, with $\log M\le23\log b\le x$ by {\cite[(36) and~(38)]{P46}},
$4fn_X^2\log M\le25600x^{24}\cdot x^{1/100}\cdot23\le x^{26}$. Finally $2x^{26}\le b$ by
{\cite[(36)]{P46}}.
\end{proof}

\begin{remark}\label{plf:rem-classes}
The spaces to which Lemma~\ref{po:classes} is applied in \S\ref{po:sec} and below are $V_j$
(rank $6r\le2n_m$), $V''$ and $V^\circ$ (ranks $N''$, $n^\circ<6r$), and the auxiliary spaces
$E^m$, $E^{3m}$, $W^\flat=E^{6m+1}$ of \S\ref{plf:ss-ff}; their levels ($K_\vartheta$, $J_{\vartheta,i}$,
$J^\circ$, $K^\circ$, $K(\pi_M^2)$) all contain $K(\pi_M^2)$, except $K_\vartheta$ and $J_{\vartheta,i}$
for $\vartheta\ne1$, which are defined also by congruences at $\mathsf Q$; for those
{\cite[Proposition~7.20]{P49}} with $\mathsf G=M^2\auNm Q_\vartheta$ applies, but only $\vartheta=1$ is
used with a bound on the number of classes (proof of (G5)). In \S\ref{pfq:sec} we also use the
order of the stabilisers: put
\[
n'_{\mathrm{Aut}}=\#\auGL_{4fn_m}(\auF_3)\le3^{16f^2n_m^2}\le\exp(x^{40}).
\]
By the proof of Lemma~\ref{po:naut}, for every standard definite $X$ of rank at most $2n_m$ and
every $g\in\auU(X)(\auA_f)$ the order of $\Gamma\cap gK^\circ g^{-1}$ divides $n'_{\mathrm{Aut}}$
($\auGL_{2fn_X}(\auF_3)$ is a subgroup of $\auGL_{4fn_m}(\auF_3)$).
\end{remark}

\subsubsection{The type of the descent: statement (L4)}\label{plf:ss-L4}

\emph{The type of $\poA(W_j)$.} At $\tau\in\auPl$ the group $\auU(W)(E^+_\tau)$ is $\auU(r_+,r_-)$
for the signs of $W_j$, and $\poA(W_j)$ is the space of forms of type $e_W=\bigotimes_\tau e_{W,\tau}$
and level $J_M$ in the sense of {\cite[Definition~8.1]{P49}} (with $G_j$ replaced by $\auU(W)$):
$e_{W,\tau}=d_\tau\aubar{c_{W,\tau}}$, $c_{W,\tau}(x)=\langle v,\pi(x)v\rangle$ for the discrete
series $\pi=\pi_{\varsigma_\tau}$ with Harish-Chandra parameter $\varsigma_\tau$, the signs being those
of {\cite[\S14.3]{P46}}, and $v$ a unit highest weight vector of its lowest $K$-type. We number the
$r$ entries of $\varsigma_\tau$ decreasingly, $\varsigma_\tau(0)>\dots>\varsigma_\tau(r-1)$; they are
\[
\tfrac k2-1,\dots,1-\tfrac k2\ \ (k-1\text{ entries}),\qquad-\kappa(\tau)-a\Delta\ \ (1\le a\le m),
\]
all integers ($k$ is even, $\Delta$ is even); the first $r_+=k-1+j$ of them lie on the positive
coordinates of $W_j$ and the last $r_-=m-j$ on the negative ones ({\cite[\S14.3]{P46}}). Two
entries on coordinates of different signs differ by at least $\Delta=2r^8$.

\begin{lemma}\label{plf:A12}
The conclusions of {\cite[Theorem~3.2]{P49}} (with $B=x^{25}$, $N_j$ replaced by $r$) and of
{\cite[Corollary~3.3]{P49}} hold for $c_W=\bigotimes_\tau c_{W,\tau}$. Consequently the conclusions of
{\cite[Proposition~8.4{\rm(a)}]{P49}} hold for forms of type $e_W$ or $\aubar{e_W}$ on $\auU(W)$: they
are cusp forms, rapidly decreasing on Siegel sets, and those of a given level form a
finite-dimensional space.
\end{lemma}
\begin{proof}
The proofs of {\cite[Theorem~3.2]{P49}} and {\cite[Corollary~3.3]{P49}} use only the decay bound
{\cite[(10)]{P49}} with $c=B+N^2+1$, respectively $c=5N^3$, for the parameter at hand ({\cite[Remark~3.1]{P49}},
proof of {\cite[Corollary~3.3]{P49}}). By {\cite[Theorem~3.17]{P49}} and {\cite[Theorem~5.6]{P51}}, the bound holds
for every regular parameter in $(\frac{N-1}2+\Z)^N$ whose gap $h$ satisfies $h\ge N^2+1$, for every
integer $c\le h$. Here $N=r$, the gap of $\varsigma_\tau$ is at least $\Delta=2r^8$, and
$2r^8\ge x^{25}+r^2+1$, $2r^8\ge5r^3$ ($r\ge x^5$). The last assertion is the proof of
{\cite[Proposition~8.4(a)]{P49}}, which uses only (A1) and (A2) for the group at hand.
\end{proof}

\emph{The Fock model.} At $\tau$ we use the Fock model $\pzF_\tau$ of {\cite[\S3.1.4]{P49}} for the
pair $(\auU(W),\auU(V))$: rows $R$ = the basis of $W_j$ (signs $\varepsilon_a$), columns $C$ = the
basis of $V_j$ (all signs $\varepsilon'_c=+1$), $n_R=r$, $n_C=n=r+t$, variables $z_{ac}$, operators
$E_{ab}$, $E'_{cd}$ as in {\cite[(12)]{P49}}, with constants $\gamma,\gamma'$.

\begin{lemma}[Constants]\label{plf:const}
For the representation $\omega_\mu$ of \eqref{po:eq-w5} the derived actions at $\tau$ are
{\cite[(12)]{P49}} after conjugation by an element of the diagonal tori, with
$\gamma=\kappa_j(\tau)+\frac n2$; the constant $\gamma'$ does not enter below.
\end{lemma}
\begin{proof}
This is the computation of the proof of {\cite[Lemma~3.6]{P49}}. By (V$'$) of {\cite[\S3.2.7]{P49}}, the
Weil representation defined with $\upsilon$ contributes $\det(g)^{n/2}$ at $\tau$; its proof
({\cite[Theorem~3.29]{P49}}) uses only the signature of the hermitian space $\auHom(W,V)\otimes_{E,\tau}\C$:
for $j<m$ this space is indefinite ($W_j$ is), and Step~2 of that proof applies; for $j=m$ it is
definite and Step~3 applies, with (S). In the convention of \eqref{po:eq-w5} the group $\auU(W)$ is
further twisted by $b(\det g)^{-1}$ with $b=\mu\,\upsilon^{\lfloor n/2\rfloor}$ (\S\ref{po:ss-set}),
and at $\tau$, $\mu=\mu^0_j\vartheta$ has weight $-\kappa_j(\tau)$ ($\vartheta$ has weight $0$) and
$\upsilon_\tau(z)=z$ on $\auU(1)$. So $b(\det g)^{-1}$ contributes $\det(g)^{\kappa_j(\tau)-\lfloor n/2\rfloor}$, and, as in
{\cite[Lemma~3.6]{P49}}, $\gamma=\tau_0+\frac n2$ with $\tau_0=\kappa_j(\tau)+\frac n2-\lfloor\frac
n2\rfloor=\kappa_j(\tau)$, since $n=6r$ is even. (The twist $a(\det h)$ of $\auU(V)$ changes only
$\gamma'$.)
\end{proof}

\emph{Matching.} Put $\lambda=-\varsigma_\tau$ (the parameter of $\aubar{\pi}$) and
$\xi_a=\lambda_a-\gamma+\frac n2=-\varsigma_\tau(a)-\kappa_j(\tau)$ for $a\in R$. Recall
$\kappa_j=\kappa+(j+\frac12)\Delta$.

\begin{lemma}[Layers]\label{plf:match}
Let $p(a)$ be the number of the row $a$: its position among the rows of its sign, the positive ones
numbered by decreasing $\varsigma_\tau$ and the negative ones by increasing $\varsigma_\tau$, from $0$.
Define $c\colon R\to C$ by: $c(a)$ is the coordinate number $p(a)$ of $V$ if $\varepsilon_a=+1$,
and the coordinate number $n-1-p(a)$ if $\varepsilon_a=-1$; let $B\subset C$ be the $t$ remaining
coordinates (the coordinates $r_+,\dots,r_++t-1$). Let $\gamma'$ be given, $s_0=\gamma+\gamma'-
\frac{r+n}2$, and define $\lambda'\in(\gamma'-\tfrac r2+\Z)^{C}$ by $\lambda'_{c(a)}=-\lambda_a+s_0$ and, on $B$ in
increasing order of coordinates, $\lambda'=\gamma'-\frac r2+\frac{t-1}2-b$ $(0\le b<t)$. Then:
\begin{enumerate}
\item[(a)] $(\lambda,\lambda')$ match in the sense of {\cite[(13)]{P49}}, with this $c$ and $B$;
\item[(b)] with $w_a$ as in {\cite[\S3.1.5]{P49}}, $|w_a|$ is the length of the row number $p(a)$ of
$\mu^+_\tau$ if $\varepsilon_a=+1$, and of the row number $p(a)$ of $\mu^-_\tau$ if
$\varepsilon_a=-1$ (rows of a diagram numbered from $0$, longest first).
\end{enumerate}
Here $\mu^\pm_\tau$ are the diagrams of \S\ref{po:ss-set}, defined by their numbers of columns of
each height $h\ge1$:
\[
c^+_\tau(h)=\bigl(\varsigma_\tau(h-1)-\max(\varsigma_\tau(h),\tfrac{t-1}2-\kappa_j(\tau))-1\bigr)^+,
\quad
c^-_\tau(h)=\bigl(\min(\varsigma_\tau(r-h-1),-\kappa_j(\tau)-\tfrac{t-1}2)-\varsigma_\tau(r-h)-1\bigr)^+ ,
\]
for $1\le h\le r$, with the convention $\varsigma_\tau(r)=-\infty$ (so that the maximum is
$\frac{t-1}2-\kappa_j(\tau)$ for $h=r$), and $c^-_\tau(r)=0$; the proof shows $c^\pm_\tau(h)=0$ for $h>r_\pm$.
\end{lemma}
\begin{proof}
\emph{Signs and sizes.} The entries on positive rows are $\ge-\kappa(\tau)-j\Delta=-\kappa_j(\tau)+
\frac\Delta2$, so $\xi_a\le-\frac\Delta2<-\frac t2$; those on negative rows are $-\kappa-a'\Delta$ with
$a'>j$, so $\xi_a=(a'-j-\frac12)\Delta\ge\frac\Delta2>\frac t2$. Hence $|\xi_a|>\frac t2$ and
$\operatorname{sgn}\xi_a=-\varepsilon_a$; since $\varepsilon'\equiv1$, $\eta_{a\,c(a)}=-\varepsilon_a
\varepsilon'_{c(a)}=\operatorname{sgn}\xi_a$. With $\xi'_c=\lambda'_c-\gamma'+\frac r2$ we get
$\xi'_{c(a)}=-\lambda_a+s_0-\gamma'+\frac r2=-\xi_a$ (because $s_0-\gamma'+\frac r2=\gamma-\frac n2$),
and $\{\xi'_c:c\in B\}=\{\frac{t-1}2-b\}$; $\varepsilon'$ is constant on $B$. All entries are
integers ($t$ odd). This is (a).

\emph{Row lengths.} By the computation in the proof of {\cite[Lemma~3.10]{P49}}, for $\xi_a<0$
\[
2w_a=2\xi_a+1+2\!\!\sum_{\xi_a<\xi_k<0}\!\!\varepsilon_a\varepsilon_k+\varepsilon_a\sum_{c\in B}
\varepsilon'_c,
\]
and for $\xi_a>0$ the formula displayed there. (For $\xi_a<0$: in $2w_a=2\xi_a+A_a-\beta_a$ one has
$\beta_a=\varepsilon_a\sum_k\varepsilon_k\operatorname{sgn}\xi_k-\varepsilon_a\sum_B\varepsilon'_c$,
and $\operatorname{sgn}(\xi_k-\xi_a)-\operatorname{sgn}\xi_k$ equals $2$ if $\xi_a<\xi_k<0$ and $0$
otherwise.) Let $\varepsilon_a=+1$, $p=p(a)$, so $\varsigma_\tau(a)=\varsigma_\tau(p)$. The $k$ with
$\xi_a<\xi_k<0$ are the positive rows with smaller $\varsigma$, $r_+-1-p$ of them, and
$\sum_B\varepsilon'=t$; so $|w_a|=-w_a=\varsigma_\tau(p)+\kappa_j(\tau)-\frac{t-1}2-(r_+-p)$. On the
other hand the row number $p$ of $\mu^+_\tau$ has length $\sum_{h>p}c^+_\tau(h)$. For $h<r_+$,
$\varsigma_\tau(h)>-\kappa_j+\frac\Delta2>\frac{t-1}2-\kappa_j$, so $c^+_\tau(h)=\varsigma_\tau(h-1)-
\varsigma_\tau(h)-1$; for $h=r_+$, $\varsigma_\tau(r_+)\le-\kappa_j-\frac\Delta2$ and $c^+_\tau(r_+)=
\varsigma_\tau(r_+-1)-\frac{t-1}2+\kappa_j-1$; for $h>r_+$, $\varsigma_\tau(h-1)<\frac{t-1}2-\kappa_j$
and $c^+_\tau(h)=0$. Summing (telescoping) gives $\varsigma_\tau(p)-\frac{t-1}2+\kappa_j-(r_+-p)=|w_a|$.
Let $\varepsilon_a=-1$, $p=p(a)$; then $\varsigma_\tau(a)=\varsigma_\tau(r-1-p)=-\kappa-a'\Delta$ with
$a'=m-p$. The $k$ with $0<\xi_k<\xi_a$ are the negative rows $-\kappa-a''\Delta$ with $j<a''<a'$, that is
$a'-j-1$ of them, with $\varepsilon_a\varepsilon_k
=1$; so $w_a=\xi_a-\frac12-(a'-j-1)-\frac t2$, i.e.\ $|w_a|=(a'-j-\frac12)\Delta-\frac{t+1}2-(a'-j-1)$.
The row number $p$ of $\mu^-_\tau$ has length $\sum_{h>p}c^-_\tau(h)$; for $h<r_-$,
$c^-_\tau(h)=\varsigma_\tau(r-h-1)-\varsigma_\tau(r-h)-1$; for $h=r_-$, $\varsigma_\tau(r-r_--1)=
-\kappa-j\Delta>-\kappa_j-\frac{t-1}2$ and $c^-_\tau(r_-)=-\kappa_j-\frac{t-1}2-\varsigma_\tau(r-r_-)-1$; for $h>r_-$ the minimum is
$-\kappa_j-\frac{t-1}2<\varsigma_\tau(r-h)$ and $c^-_\tau(h)=0$. Telescoping, the length is
$-\kappa_j-\frac{t-1}2-\varsigma_\tau(r-1-p)-(r_--p)=(a'-j-\frac12)\Delta-\frac{t-1}2-(a'-j)$, which is
$|w_a|$.
\end{proof}

\begin{remark}\label{plf:rem-mu}
Lemma~\ref{plf:match} makes precise the diagrams $\mu^\pm_\tau$ of \S\ref{po:ss-set}. The two properties used in
\S\ref{po:sec} follow: $\mu^+_\tau$ has $\sum_hc^+_\tau(h)=\varsigma_\tau(0)+\kappa_j-\frac{t-1}2-r_+
\le\frac k2+\kappa_j\le4b$ columns, by {\cite[(38)]{P46}} and the definition of $\kappa$; and
$c^-_\tau(h)\le\Delta$ for each $h$. The column heights are at most $r_\pm$, the last displayed
identities of the proof.
\end{remark}

\emph{The polynomial of a highest weight vector.} Let $\xi_0\in W_{\lambda_V}$ be the vector dual
to the pure tensor of the first basis vectors: $\langle\xi_0,\mathrm{ev}(y,y')\rangle=\prod_{\mathfrak
c}\det\bigl(y_{i}(e_{l})\bigr)_{i,l<h_{\mathfrak c}}\cdot\prod_{\mathfrak c'}\det\bigl(y'_{i}(e_{n-1-l})
\bigr)_{i,l<h_{\mathfrak c'}}$ up to complex conjugation of all entries, the first product over the
columns of $\mu^+$, the second over those of $\mu^-$, in \eqref{po:eq-ev}. By \eqref{po:eq-poly},
written in the Fock variables (for a positive row $z_{ac}$ is the conjugate of $x_{ac}$, for a
negative row the conjugate-linear map $\flat$ produces the variable $z_{ac}$), $\pop_{\xi_0}$ is
\[
\pop_{\xi_0}=\prod_{\mathfrak c\in\mu^+}\det\bigl(z_{a\,c(a')}\bigr)_{a,a'\in R^+_{<h_{\mathfrak c}}}
\cdot\prod_{\mathfrak c'\in\mu^-}\det\bigl(z_{a\,c(a')}\bigr)_{a,a'\in R^-_{<h_{\mathfrak c'}}},
\]
where $R^\pm_{<h}$ is the set of rows of sign $\pm$ with number $<h$: the columns of the positive
rows are taken in the fixed numbering and those of the negative rows in the reverse one
(\S\ref{po:ss-set}), and the dual coordinates of the highest weight vector of $\Lambda^h(\C^n)^\vee$
are the last ones.

\begin{lemma}\label{plf:harm}
$\pop_{\xi_0}$ is, up to a non-zero scalar, the joint harmonic $\Phi$ of
{\cite[Proposition~3.11]{P49}} for the matching pair $(\lambda,\lambda')$ of Lemma~\ref{plf:match}.
\end{lemma}
\begin{proof}
By Lemma~\ref{plf:match}(b) a row $a$ has $|w_a|$ layers, and the rows having a pair of level
$\ell$ are those with $|w_a|>\ell$, that is, since the row lengths of a diagram decrease, the rows
of sign $\varepsilon$ with number less than the height of the column number $\ell$ of
$\mu^{\varepsilon}_\tau$ (columns numbered from $0$, highest first). In the notation of
{\cite[Proposition~3.11]{P49}}, $\varepsilon'\equiv1$, so $L_{\ell,\epsilon,1}=R^{\epsilon}_{<h}$ with
$h$ that height, and $M_{\ell,\epsilon,1}=\det(z_{a\,c(a')})_{a,a'\in R^\epsilon_{<h}}$. The product
over $\ell$ and $\epsilon$ is the product over the columns of $\mu^+$ and of $\mu^-$ displayed
above.
\end{proof}

\begin{proposition}[(L4)]\label{plf:L4}
Let $\vartheta$ be a twist and $\Xi\in\poA(V;\vartheta)$. Then $D_\vartheta\Xi\in\poA(W_j)$ (forms of
type $e_W$ and level $J_M$). Consequently statement {\rm(L4)} of \S\ref{po:ss-set} holds: if
$\Xi\in\varrho_V(e_0)\poA(V;\vartheta)$ is orthogonal to $\Theta_\vartheta(e_0\poA(W_j))$, then
$D_\vartheta\Xi=0$.
\end{proposition}
\begin{proof}
\emph{Step 1: the archimedean identity.} Let $\sigma=\omega_\infty(\cdot,1)$ on $\pzF=\bigotimes_\tau
\pzF_\tau$, and $\mathsf P=\int e_W(x)\sigma(x^{-1})\,dx$, which converges absolutely in $\pzF_s$ by
Lemma~\ref{plf:A12} and~{\cite[Lemma~8.2(a)]{P49}} ($s=D+1$, $D=frn$, $2s\le x^{25}$). We show
$\mathsf P\pop_\xi=\pop_\xi$ for all $\xi\in W_{\lambda_V}$. As in the proof of
{\cite[Theorem~3.15]{P49}}, $\mathsf P=\bigotimes_\tau\sigma_\tau(e_{\lambda})$ with $\lambda=-\varsigma_\tau$
at $\tau$ (because $e_W(x^{-1})=\aubar{e_W(x)}$ and $\aubar{c_{W,\tau}}=c_{-\varsigma_\tau}$), and
$\pop_{\xi_0}=\bigotimes_\tau\pop_{\xi_0,\tau}$; so it suffices to show $\sigma_\tau(e_\lambda)\Phi=
\Phi$ at each $\tau$ for $\Phi=\pop_{\xi_0,\tau}$. We apply {\cite[Lemma~3.5]{P49}}. By
Lemma~\ref{plf:harm} and {\cite[Definition~3.9]{P49}}, $\Phi$ has $T$-weight $\Lambda^K(\lambda)$
and $E_{ab}\Phi=0$ for $\varepsilon_a=\varepsilon_b$, $\lambda_a>\lambda_b$; and $\Phi$ has
$T'$-weight $\Lambda'=\Lambda^K(\lambda')=\lambda'-\rho'$ (the group $\auU(n)$ is compact, all roots
are compact) and is killed by $E'_{cd}$ for $\lambda'_c>\lambda'_d$, so it is a highest weight vector
for $\mathfrak{gl}_n$ and $\Omega'\Phi=(|\Lambda'+\rho'|^2-|\rho'|^2)\Phi=\kappa_{\lambda'}\Phi$. Since
$\sum_aE_{aa}\Phi=(\sum_a\Lambda_a)\Phi=(\sum_a\lambda_a)\Phi$ and $(\lambda,\lambda')$ match
(Lemma~\ref{plf:match}(a)), {\cite[Lemma~3.14]{P49}} gives
$\Omega\Phi=\bigl(\Omega'+2s_0\mathcal N-c\bigr)\Phi=\kappa_\lambda\Phi$, with $c$ the constant of
that lemma. The function $e_\lambda$ is integrable (Lemma~\ref{plf:A12}). By {\cite[Lemma~3.5]{P49}},
$\sigma_\tau(e_\lambda)\Phi=\Phi$. Finally the operators $\mathsf P$ and $\omega_\infty(1,y)$,
$y\in\auU(V)(E^+\otimes\R)$, commute, and $\xi\mapsto\pop_\xi$ is conjugate-linear and
$\auU(V)(\R)$-equivariant (\S\ref{po:ss-set}); so $\{\xi:\mathsf P\pop_\xi=\pop_\xi\}$ is an invariant
subspace of the irreducible representation $W_{\lambda_V}$ containing $\xi_0\ne0$, hence everything.

\emph{Step 2: the type of the descent.} By \eqref{po:eq-descent}, $D_\vartheta\Xi=\aubar F$ with
$F(g)=\sum_hc_h\,\theta\bigl(\omega_\mu(g,h)(\pop_{\Xi(h)}\otimes\Phi_\vartheta)\bigr)$, a finite sum
($c_h>0$). By {\cite[Lemma~8.2]{P49}}, $|F(gx)|\le C(g)\|x\|^{2s}$ for $x\in\auU(W)(E^+\otimes\R)$, and
$\Psi\mapsto\theta(\omega_\mu(g,h)(\Psi\otimes\Phi_\vartheta))$ is continuous on $\pzF_s$; by
Lemma~\ref{plf:A12} and Fubini,
\[
\int\aubar{e_W(x)}F(gx)\,dx=\sum_hc_h\,\theta\Bigl(\omega_\mu(g,h)\Bigl(\int\aubar{e_W(x)}\,
\omega_\infty(x,1)\pop_{\Xi(h)}\,dx\otimes\Phi_\vartheta\Bigr)\Bigr)=F(g),
\]
since $\int\aubar{e_W(x)}\omega_\infty(x,1)\,dx=\int e_W(x^{-1})\omega_\infty(x,1)\,dx=\mathsf P$ and
$\mathsf P\pop_\xi=\pop_\xi$ by Step~1. Taking complex conjugates, $R(e_W)D_\vartheta\Xi=D_\vartheta\Xi$.
The function $D_\vartheta\Xi$ is continuous and left invariant under $\auU(W)(E^+)$ (invariance of
$\theta$), and right invariant under $J_M$ (\S\ref{po:ss-set}). It is bounded: the truncation
argument at the end of the proof of {\cite[Proposition~8.4(b)]{P49}} applies verbatim, with the
finite-dimensional space of forms of type $e_W$ and level $J_M$ (Lemma~\ref{plf:A12}). So
$D_\vartheta\Xi\in\poA(W_j)$.

\emph{Step 3: (L4).} If $\Xi$ is orthogonal to $\Theta_\vartheta(e_0\poA(W_j))$, then by (L2) and
Lemma~\ref{po:adj}(c) it is orthogonal to $\Theta_\vartheta y$ for all $y\in\poA(W_j)$ (decompose
$y=\varrho(e_0)y+(1-\varrho(e_0))y$; $\Theta_\vartheta\varrho(1-e_0)y=\varrho_V(1-e_0)\Theta_\vartheta y$
is orthogonal to $\varrho_V(e_0)\poA(V;\vartheta)\ni\Xi$). By \eqref{po:eq-descent},
$\int_{[\auU(W)]}y\,\aubar{D_\vartheta\Xi}=0$ for all $y\in\poA(W_j)$; take $y=D_\vartheta\Xi$, which is
allowed by Step~2. So $D_\vartheta\Xi=0$.
\end{proof}

\begin{remark}\label{plf:rem-L4}
Step~1 is the analogue of (A3) ({\cite[Theorem~3.15]{P49}}) for the pair $(W_j,V_j)$, with the compact
group $\auU(V_j)(\R)$ on the column side; on that side rigidity is automatic, and only the row side
needs {\cite[Lemma~3.5]{P49}}, that is, the published results (D1)--(D3) of {\cite[\S3.1.3]{P49}}. So
Proposition~\ref{plf:L4} is proved modulo exactly the inputs of {\cite[Theorem~3.15]{P49}}: (D1)--(D3),
and (V$'$), proved in {\cite[Theorem~3.29]{P49}} modulo (W1$_\infty$), (W2$_\infty$), (M), and (S) for
$j=m$.
\end{remark}

\subsubsection{Kernel scales for definite pairs}\label{plf:ss-KS}

The first factors are built from theta kernels of pairs of standard definite spaces with data at
the chosen places and at $\mathsf Q$. We state the integrality and size of such kernels in the
generality needed, and prove them from {\cite[Theorem~7.15]{P49}}, {\cite[Lemma~5.7]{P49}} and the argument of
{\cite[\S5.1.5]{P49}}.

\emph{Setting.} Let $L=E^{l}$ and $N=E^{n_N}$ be standard definite spaces, $\mathsf W=\auHom(L,N)=
M_{n_N\times l}$, and let $a,b$ be unitary Hecke characters of $\auU(1)(\auA_{E^+})$ with
algebraic weights at the infinite places and tame ramification, at split places only. Let
$\omega_{a,b}$ be the Weil representation of $\auU(L)(\auA)\times\auU(N)(\auA)$ on $\psHS(\mathsf W)$
twisted by $a(\det g)^{-1}$ and $b(\det h)$, as in \eqref{po:eq-w5}. A \emph{modulus} is an ideal
$\mathfrak m$ of $\auO_E$ divisible only by split primes $w$ (with $\aubar w\nmid\mathfrak m$); for a
function $\Psi$ on $M_{n_N\times l}(\auO_E/\mathfrak m)$ let $f_\Psi\in\psHS_f(\mathsf W)$ be the
datum $x\mapsto\plone[x\in L_w]\,\Psi(a_w^{-1}x\bmod\mathfrak m)$ at the places $w\mid\mathfrak m$,
$x$ the matrix in the frames (model (W5)), and the vacuum elsewhere; here $a_w$ is the generator of
$\alpha_w\auO_w$ of {\cite[Definition~6.31]{P53}} and $L_w=a_wM_{n_N\times l}(\auO_w)$, which is the
$w$-component $M_w$ of the self-dual lattice of {\cite[Lemma~6.30]{P53}}. In the coordinates $T=a_w^{-1}x$ of
(N5) ({\cite[Lemma~6.32(c)]{P53}}) this is $T\mapsto\plone[T\ \text{integral}]\,\Psi(T\bmod\mathfrak m)$. By {\cite[Proposition~6.9]{P53}}, $a_w=1$ and $L_w$ is the
lattice of integral matrices except at primes $w$ above $q_1$ dividing $\omega$, which can divide
$\mathfrak m$ only through $Q_\vartheta$; this is the datum $\Psi^\natural$ of {\cite[\S2.3]{P53}}. For a polynomial $Q$ in the variables
$z_{\tau,c,a}$ ($\tau\in\auPl$, $c\le n_N$, $a\le l$) let $\pop_Q\in\pzF$ be its polynomial vector,
and
\[
\plfK_{a,b}[\Psi,Q](g,h)=\theta\bigl(\omega_{a,b}(g,h)(\pop_Q\otimes f_\Psi)\bigr)\qquad(g\in\auU(L)(\auA_f),\
h\in\auU(N)(\auA_f)),
\]
the archimedean components being $1$. For $Q^\sigma$ we apply $\sigma$ to the coefficients of $Q$.

\begin{lemma}[Kernel scales]\label{plf:KS}
For all $l,n_N\ge1$ and $e\ge0$ there is $c=c_{l,n_N,e}\in\C^\times$ such that, for all $a,b$,
$\mathfrak m$, $\Psi$ and all $Q$ homogeneous of degree $e$:
\begin{enumerate}
\item[(KS1)] Let $q\ge1$ be an integer with $qM^2\in\mathfrak m$, let $\iota\colon\C\to\C_\ell$ with
$\ell\nmid q$, and suppose that the values of $\Psi$ and the coefficients of $Q$ are integral at
$\iota$ and that $g,h$ are integral with integral inverses at the places above $\ell$. Then
$|c\,\plfK_{a,b}[\Psi,Q](g,h)|_\iota\le1$.
\item[(KS2)] Suppose $l,n_N\le\Delta$, $e\le M$, $M^2\in\mathfrak m$, $|\sigma'(\Psi(x))|\le A$ for all
$x$ and all $\sigma'\in\auAut(\C)$, and $|Q^\sigma(z)|\le B$ for all $z$ with $\max_i|z_i|\le M^4$.
Then $|\sigma(c\,\plfK_{a,b}[\Psi,Q](g,h))|\le e^{x_1}AB$ for all $g,h$.
\end{enumerate}
\end{lemma}
\begin{proof}
\emph{Jets.} Write $\Theta_f$ for the function of {\cite[(42)]{P49}} attached to $f\in\psHS_f(\mathsf
W)$ ({\cite[\S5.2.8]{P49}}, with $\mathsf C=[n_N]\times[l]$). The polynomial vectors are the
derivatives at $0$ of the coherent states $\psrho_\infty(z)1$ times $e^{|z|^2/2}$; hence, as in
{\cite[(34)]{P49}}, there is a constant $\kappa_{\mathrm F}\ne0$, depending only on the
normalisation of the Fock model, with $\theta(\pop_Q\otimes f)=\kappa_{\mathrm F}^{\,e}\,
Q(\partial)[\Theta_f](0)$ for $Q$ homogeneous of degree $e$. For $s\in\C^\times$, $Q(\partial)[\Theta_f
(s\,\cdot)](0)=s^eQ(\partial)[\Theta_f](0)$.

\emph{The datum as a combination of translates of the vacuum.} At a place $w\mid\mathfrak m$, the
indicator of $x+\mathfrak mL_w$ ($L_w=a_wM_{n_N\times l}(\auO_w)$ as above, $x\in L_w$) is a root of unity times
$\psrho^\circ(x)\plone_{\mathfrak mL_w}$, and by Fourier inversion on the finite group
$\mathfrak m^{-1}L^\vee_w/L^\vee_w$ in the variables of $\aubar w$,
$\plone_{\mathfrak mL_w}=[L_w:\mathfrak mL_w]^{-1}\sum_y\zeta_y\,\psrho^\circ(y)\psvac_{L,w}$, a sum of
$[L_w:\mathfrak mL_w]$ terms with roots of unity $\zeta_y$ and translations $y$ with $\mathfrak my$
integral. Taking the product over $w\mid\mathfrak m$ and writing $\psvac=c_{\mathrm{vac}}^{-ln_N}
\psvac^\circ$ ({\cite[(58)]{P49}}), we get
\begin{equation}\label{plf:eq-vacexp}
c_{\mathrm{vac}}^{\,ln_N}\,\auNm(\mathfrak m)^{ln_N}f_\Psi=\sum_{x,y}\zeta_{x,y}\,\Psi(x)\,
\psrho^\circ(x+y)\psvac^\circ ,
\end{equation}
a sum of $\auNm(\mathfrak m)^{2ln_N}$ terms, $\zeta_{x,y}$ roots of unity, $\mathfrak m(x+y)$ integral
at the places of $\mathfrak m$ and $x+y=0$ elsewhere.
Here $\psvac_{L,w}$ is the vacuum, which is $\plone_{L_w}$ up to a scalar by {\cite[Lemma~6.30]{P53}}
\emph{because} $L_w=a_wM_{n_N\times l}(\auO_w)$ is the $w$-component of the self-dual lattice (with
the lattice of integral matrices this would fail at a prime where $\alpha_w$ is not a unit),
$L_w^\vee$ is its dual $M_{\aubar w}$, $\Psi(x)$ stands for $\Psi(a_w^{-1}x\bmod\mathfrak m)$,
``$\mathfrak m(x+y)$ integral'' means $\mathfrak mx\in\mathfrak mL_w$, $\mathfrak my\in L^\vee_w$, and
$[L_w:\mathfrak mL_w]=\auNm(\mathfrak m_w)^{ln_N}$ does not depend on $a_w$. A prime $w\mid\mathfrak m$
with $a_w\ne1$ lies above some prime $p\mid q_1$, and $p\nmid M$ ($p_0,p_1\equiv1\pmod{q_1}$); in
(KS1) below $qM^2\in\mathfrak m\subseteq w$, so $p\mid q$ and $\ell\ne p$. Hence such $w$ do not lie above
$\ell$, and at the places above $\ell$ the translations are exactly as in the case $a_w=1$. (Such
primes can occur: they are split over $E^+$, {\cite[Proposition~6.9(b)]{P53}}, of norm prime to $MN$, so
the definition of a twist allows them in $Q_\vartheta$.)

\emph{Definition of $c$.} Let $s_0$ be the scale of {\cite[\S5.2.8]{P49}} and put
\[
c_{l,n_N,e}=c_{\mathrm{vac}}^{\,ln_N}\,M^{4fln_N}\,\bigl(dM^2s_0/\kappa_{\mathrm F}\bigr)^e .
\]

\emph{(KS1).} Since $qM^2\in\mathfrak m$, $\auNm(\mathfrak m)$ divides $(qM^2)^{2f}$, so
$M^{4fln_N}\auNm(\mathfrak m)^{-ln_N}$ is integral at $\ell\nmid q$; and at the places above $\ell$
(which do not divide $q$) $M^2(x+y)$ is integral. So by \eqref{plf:eq-vacexp}
$c_{\mathrm{vac}}^{ln_N}M^{4fln_N}f_\Psi$ is a standard datum of level $M^2$ at $\ell$ in the sense of
{\cite[\S7.4]{P49}}; the values of $\Psi$ are integral and the coefficients are roots of unity.
By (S1) at the level $M^2$ ({\cite[\S7.4]{P49}}), $\omega_{a,b}(g,h)$ preserves standard data of
level $M^2$ at $\ell$ up to the scalars $a(\det g)^{-1}b(\det h)$ and the unramified factors of
\eqref{po:eq-w5}, which are units at $\ell$ for $g,h$ integral with integral inverse above $\ell$
(values of algebraic Hecke characters at ideles that are units above $\ell$, and powers of residue
characteristics different from $\ell$). Let $f'$ be the resulting standard datum. Then
\[
c\,\plfK_{a,b}[\Psi,Q](g,h)=(\text{unit})\cdot Q(\partial)\bigl[\Theta_{f'}(d\cdot M^2s_0\,\cdot)\bigr](0).
\]
By {\cite[Theorem~7.15]{P49}} with $\delta=M^2$ ($\aubar\delta=\delta$), $\Theta_{f'}(M^2s_0z)=\bigl(\prod_j
\theta(\mu_j(z))^{m_j}\bigr)^{\auNm\delta}\mathsf H(z)$ with $\mu_j$ linear forms with integral
coefficients and $\mathsf H$ an integral power series in the $y(\mu_j(z))$. We claim that
$\Theta_{f'}(dM^2s_0\,\cdot)$ is $\partial$-integral at every prime $\ell$, including $\ell=d$. For
$\ell\ne d$ this is the last assertion of {\cite[Theorem~7.15]{P49}} with $c=d$. For $\ell=d$: $y$ is
$\partial$-integral at every prime (proof of (C2) in {\cite[Proposition~5.25]{P49}}), so $y(d\mu_j(z))$
and $\mathsf H(dz)$ are; and $\theta=\exp(-g)$ with $g(0)=g'(0)=0$, $g^{(n)}(0)=x^{(n-2)}(0)\in\pao$ for
$n\ge3$ (same proof) and $d^2g''(0)\in\pao$ (at $\ell=d$ the numbers $x(u-t)$ in the proof of (C2) need
not be integral, but $v(g''(0))>-2$ for the valuation with $v(d)=1$: proof of Lemma~\ref{pex:dint}); hence
$z\mapsto g(dz)$ has derivatives $d^ng^{(n)}(0)\in
\pao$ and vanishes at $0$, and $\theta(d\,\cdot)=\exp(-g(d\,\cdot))$ is $\partial$-integral, as are
products and compositions with linear forms with integral coefficients ({\cite[\S4.3]{P49}}). This
proves the claim, and the jet $Q(\partial)[\cdot](0)$ of a $\partial$-integral series along a polynomial
with integral coefficients is integral.

\emph{(KS2).} We may assume $\sigma|_K=1$ (proof of {\cite[Proposition~14.23]{P46}}). By
\eqref{plf:eq-vacexp} and Step~1 of that proof, $c_{\mathrm{vac}}^{ln_N}\auNm(\mathfrak m)^{ln_N}
\omega_{a,b}(g,h)f_\Psi$ is a sum of $\auNm(\mathfrak m)^{2ln_N}\le M^{8fln_N}$ vectors
$\epsilon_{x,y}\Psi(x)\,\psrho^\circ(x')\omega(k)\psvac^\circ$ with $|\sigma\epsilon_{x,y}|=1$. For each
of them, Step~2 of that proof, with $\pi_M$ replaced by the rational integer $dM^2$, gives
\[
\sigma\bigl(Q(\partial)[\Theta_{f}(dM^2s_0\,\cdot)](0)\bigr)=\int Q^\sigma\bigl(dM^2t_\sigma\bar
y\bigr)e^{-|y|^2/2}\,\theta\bigl(\psrho_\infty(y)(1\otimes\Lambda_\sigma f)\bigr)\,d\mu(y),
\]
with $|t_\sigma|^2=\Omega^{-1}$ and $|\theta(\psrho_\infty(y)(1\otimes\Lambda_\sigma f))|\le
\exp(C_0(q_1ln_N+1)^2)$ by {\cite[Lemma~5.7(ii)]{P49}}. Since $Q^\sigma$ is homogeneous of degree $e$,
$|Q^\sigma(z)|\le B\max(1,|z|/M^4)^e$. Here $|dM^2t_\sigma\bar y|/M^4\le\rho|y|$ with
$\rho=d\,\Omega^{-1/2}M^{-2}\le b^2\cdot x^{7}\cdot b^{-40}\le b^{-37}$ ($d\le b^2$ since $\sqrt d\le b$, $\Omega^{-1}\le
x^{13}$ by the definition of $\Omega$ and {\cite[(38)]{P46}}, $M\ge b^{20}$; {\cite[Lemma~13.4]{P46}}).
With $D=fln_N$ and the normalisation of $d\mu$ of {\cite[\S5.1.5]{P49}},
$\int(1+(\rho|y|)^e)e^{-|y|^2/2}d\mu\le2^D\bigl(1+(\rho^2(2D+e))^{e/2}\bigr)\le2^{D+1}$, because
$\rho^2(2D+e)\le b^{-74}(4\Delta^2f+M)\le1$ (as $M\le b^{22}$). Collecting: by the definition of $c$ and
\eqref{plf:eq-vacexp}, $c\,\plfK_{a,b}[\Psi,Q](g,h)=M^{4fln_N}\auNm(\mathfrak m)^{-ln_N}\sum_{x,y}
\epsilon_{x,y}\Psi(x)\,Q(\partial)[\Theta_{f_{x,y}}(dM^2s_0\,\cdot)](0)$ with $f_{x,y}=\psrho^\circ(x')
\omega(k)\psvac^\circ$ (the factor $c_{\mathrm{vac}}^{ln_N}$ is absorbed by writing the datum through
$\psvac^\circ$), and $\auNm(\mathfrak m)$ is a positive integer, so
\[
|\sigma(c\,\plfK_{a,b}[\Psi,Q](g,h))|\le AB\cdot2^{D+1}\,M^{12fln_N}\,
\exp\bigl(C_0(q_1ln_N+1)^2\bigr).
\]
With $l,n_N\le\Delta=2r^8\le2\cdot3^8x^{40}$, $f,q_1\le x^{12}$, $\log M\le23\log b\le x$ and
$C_0\le\log\log b$ ({\cite[Lemma~8.7]{P49}}), the logarithm of the factor after $AB$ is
at most $x^{200}\le x_1$ by {\cite[(36)]{P46}}.
\end{proof}

\begin{remark}\label{plf:rem-KS}
(a) (KS1) is the integrality statement (K1) of {\cite[Lemma~4.13]{P49}} for definite pairs with
arbitrary data at split places, at every prime $\ell\nmid q$, \emph{including} $\ell\mid2q_1M$
({\cite[Theorem~7.15]{P49}}), $\ell=N$ (the data are unramified at $N$), and $\ell=d$ (the extra factor
$d^e$ in the scale). The factor $d^e$ is harmless for the size: it is part of $\rho$ above.
(b) The size bound (KS2) is far from optimal; only its form $e^{x_1}AB$ is used.
\end{remark}

\subsubsection{The set of choices}\label{plf:ss-ff}

\emph{Spaces.} Besides $V''$ (rank $N''=r+9m+1$) we use the standard definite
spaces $E^m$, $E^{3m}$ and $W^\flat=E^{6m+1}$; their levels are $K^\circ$ (the stabiliser of the
standard lattice) or $K(\pi_M^2)$.

\emph{The choices.} A \emph{choice} $i=(s_i,h^\flat_i,\mathsf r_i)$ consists of a matrix $s_i\in
M_{3m\times N''}(\auO_E/\pi_M^2)$, a class $h^\flat_i\in\Gamma^\flat\backslash\auU(W^\flat)(\auA_f)/
K^\flat(\pi_M^2)$ (with a fixed representative), and, for each $\tau\in\auPl$ and each of the
$\Delta'=r^8-\frac{t+1}2$ columns of height $r_+$ of $\mu^+_\tau$, a list of $2m$ indices in
$\{1,\dots,6m+1\}$. For each class $h^\flat$ fix $Y=Y_{h^\flat}\in\auGL_{6m+1}(E)$ as in
{\cite[Lemma~14.21]{P46}} for $g=h^\flat$ ($h^{\flat,-1}Y$ integral, all conjugates of all entries at
most $C_{\mathrm{sb}}$).

\begin{lemma}\label{plf:I}
The set $I$ of choices has at most $\exp(x^{54})\le e^b$ elements.
\end{lemma}
\begin{proof}
$\#(\auO_E/\pi_M^2)=M^{2f}$ ($\pi_M\in\auO_K$ has norm $M$ to $\Q$, and $[E:K]=f$), so there are
$M^{6fmN''}$ matrices $s_i$; by Lemma~\ref{plf:classes} there are at most $\exp(2x^{26})$ classes
$h^\flat$; and at most $(6m+1)^{2mf\Delta'}$ lists. By {\cite[(36) and~(38)]{P46}} and
$m\le2\sqrt x$, $r\le3x^5$: $6fmN''\log M\le6x^{12}\cdot2\sqrt x\cdot4x^5\cdot x\le x^{19}$ and
$2mf\Delta'\log(6m+1)\le4\sqrt x\,x^{12}\,3^8x^{40}\log(13\sqrt x)\le x^{53}$.
\end{proof}

The first factors attached to the choices are those of Definition~\ref{ptf:def-ff}; their
integrality and size (Lemma~\ref{po:ff}) are proved in \S\ref{pfq:sec} (Theorem~\ref{pfq:ff}).





\providecommand{\auO}{\mathcal O}
\providecommand{\auU}{\mathrm U}
\providecommand{\auGL}{\mathrm{GL}}
\providecommand{\auHom}{\operatorname{Hom}}
\providecommand{\auA}{\mathbb A}
\providecommand{\auF}{\mathbb F}
\providecommand{\auPl}{\Phi}
\providecommand{\auG}{\overline G}
\providecommand{\aubar}[1]{\overline{#1}}
\providecommand{\plone}{\mathbf 1}
\providecommand{\psrho}{\varrho}
\providecommand{\poA}{\mathcal A}
\providecommand{\poZ}{\mathcal Z}
\providecommand{\pop}{\mathsf p}
\providecommand{\pog}{\mathsf g}
\providecommand{\ppeD}{\mathcal D}
\providecommand{\ppeVan}{\mathrm{Van}}
\providecommand{\ppeS}{\mathcal S}
\providecommand{\ppeH}{\mathbb H}
\providecommand{\ppeN}{\mathsf N}
\providecommand{\pexK}{K^{(z)}}

\subsection{Periods of false lifts}\label{ppe:sec}

This section proves Proposition~\ref{po:periods}. We keep the notation of \S\ref{po:sec}:
$1\le j\le m$ is fixed, $W=W_j$ (rank $r$, signature $(r_++1,r_-)$-type as in
{\cite[\S14.3]{P46}}), $V=V_j$ the positive definite standard space of rank $n=r+t=6r$,
$\vartheta$ a twist, $\mathsf S_\vartheta=\mathsf M\cup\mathsf Q$ the places of the level, $\mu=\mu_j^0
\vartheta$ with tame parts $\psi_w$ ($w\in\mathsf S_\vartheta$), $\omega_\mu$ the Weil
representation \eqref{po:eq-w5}, $\Phi_\vartheta$ the block datum, $K_\vartheta$, $\chi_\vartheta$,
$\poA(V;\vartheta)$, the descent $D_\vartheta$ \eqref{po:eq-descent}, $\varrho_V$, $e_0$, and the
splitting $V=V''\perp V^\circ$ of \S\ref{po:ss-glue} with $N''=r+9m+1$, $n^\circ=t-9m-1$. The
argument is the classical one of the theta tower (Rallis~\cite[Chapter~I, \S1]{RallisHoweDuality}; see also~\cite{RallisInjectivity}): the descents of a
false lift vanish to $W$, then to every $W\oplus\ppeH^p$ with $2r+2p+3<n$, and a Fourier
coefficient of the descent to the split space $\ppeH^{N''}$ is the period over $\auU(V^\circ)$.
The new ingredient, which replaces the use of poles of $L$-functions in the classical argument,
is that a false lift satisfies \emph{one} Hecke equation at one split place, at the smallest root
of a lift; this forces bounded descents to vanish.

\subsubsection{Set-up}\label{ppe:ss-setup}

\emph{Hermitian spaces and data.} For a non-degenerate hermitian space $X$ over $E$ of rank
$m_X$ (signatures at the infinite places arbitrary) we use the Weil representation
$\omega^X_a$ of $\auU(X)(\auA)\times\auU(V)(\auA)$ on $\ppeS(X)=\mathcal H(\auHom(X,V))$, in the
normalisation of \eqref{po:eq-w5}: $\auU(V)$ acts through the splitting given by powers of
$\upsilon$ multiplied by the character $a(\det h)$, where $a$ is a character of
$\auU(1)(E^+)\backslash\auU(1)(\auA_{E^+})$, and $\auU(X)$ acts through the splitting attached to
$X$ (an extra character $b(\det g)^{-1}$ on $\auU(X)$ only multiplies descents by $b(\det g)^{-1}$
and is never needed). For $X=W$ and $a=1$, $\omega^W_1$ is $\omega_\mu$ up to the character
$\mu(\det g)^{-1}$ of $\auU(W)$. Let $\theta$ be the theta functional. Since $V$ is definite,
$[\auU(V)]=\auU(V)(E^+)\backslash\auU(V)(\auA)$ is compact. For a continuous function $Y$ on
$[\auU(V)]$ and $f\in\ppeS(X)$ put
\[
\ppeD^X_a(f)Y(g)=\int_{[\auU(V)]}\theta\bigl(\omega^X_a(g,h)f\bigr)\,Y(h)\,dh\qquad(g\in
\auU(X)(\auA)),
\]
the \emph{descent} of $Y$ to $X$ by $f$; the integral converges absolutely because $[\auU(V)]$
is compact and $h\mapsto\theta(\omega(g,h)f)$ is continuous. Let $\Sigma=\mathsf S_\vartheta\cup
\aubar{\mathsf S_\vartheta}$ and let $\ppeS_\Sigma(X)\subseteq\ppeS(X)$ be the closed subspace of
the data that are the vacuum at every finite place outside $\Sigma$ (arbitrary at the infinite
places and at $\Sigma$). We say that \emph{the descents of $Y$ to $X$ vanish}, and write
$\ppeVan(X,a,Y)$, if $\ppeD^X_a(f)Y=0$ for every $f\in\ppeS_\Sigma(X)$.

For $p\ge0$ let $\ppeH^p$ be the sum of $p$ hyperbolic planes, with an isotropic frame
$e_1,\dots,e_p,f_1,\dots,f_p$, $\langle e_i,f_k\rangle=\delta_{ik}$, adapted (orthonormal up to the
standard scaling at the infinite places and spanning a self-dual lattice at the finite places),
and $W_p=W\oplus\ppeH^p$. Put $\upsilon_1=\upsilon|_{\auU(1)}$ (the restriction of the splitting
character to the norm-one ideles).

\emph{Matrix coefficients.} For $\Xi\in\poA(V;\vartheta)$ and $\zeta\in W_{\lambda_V}$ let
\[
c_\zeta\Xi(h)=\bigl\langle\zeta,\ \rho_{\lambda_V}(h_\infty)^{-1}\Xi(h_f)\bigr\rangle,\qquad
Y_\zeta=\aubar{c_\zeta\Xi}\qquad(h=h_\infty h_f\in\auU(V)(\auA)),
\]
with the form of $W_{\lambda_V}$ conjugate-linear in the first variable (as in
\eqref{po:eq-delta}). Since $\Xi(\gamma h_f)=\gamma\,\Xi(h_f)$, with $\gamma$ acting through its
archimedean components, $c_\zeta\Xi$ is left invariant under $\auU(V)(E^+)$; it is right invariant
under the elements of $K_\vartheta$ on which $\chi_\vartheta$ is trivial, in particular under
$K'_\vartheta=\{k\in K^\circ:\ k_w=1\ (w\in\Sigma)\}$, and it is continuous. If $T$ is a Hecke
operator at a place $w\notin\Sigma$ (sum of right translations by elements of $\auU(V)(\auA_f)$
supported at $w$), then $c_\zeta(T\Xi)=T(c_\zeta\Xi)$ and $Y_\zeta$ transforms in the same way
with the complex conjugate coefficients. For $x\in\auU(V)(E^+\otimes\R)$,
$c_\zeta\Xi(hx)=c_{\rho(x)\zeta}\Xi(h)$.

\begin{theorem}[Proposition~\ref{po:periods}]\label{ppe:main}
Let $w$ be off the level and $\Xi\in\mathcal F_w(e_0,\poZ_\vartheta)$ a false lift. Then for every
$s$, every $h\in\auU(V'')(\auA_f)$, $u''\in W''$ and $g\in\auU(W)(\auA)$,
\[
\sum_{h^\circ\in\Gamma^\circ\backslash\auU(V^\circ)(\auA_f)/J^\circ}\bigl\langle u''\odot
K^\circ(g,h^\circ),\ \Xi\bigl(\operatorname{diag}(h,h^\circ)\pog_s\bigr)\bigr\rangle=0 .
\]
The same holds with $J^\circ$ replaced by any open compact subgroup of $J^\circ$, and for every
$\Xi\in\poA(V;\vartheta)$ that satisfies (F1)--(F3) of Proposition~\ref{ppe:tower} below.
\end{theorem}

The pairing is that of \eqref{po:eq-delta}, conjugate-linear in the first variable; this is the
convention under which $\langle u^\circ,K^\circ(g,h^\circ)\rangle=\theta^\circ(g,h^\circ;u^\circ)$
is conjugate-linear in $u^\circ$, as stated in \S\ref{po:ss-glue}. (If the other convention is
used, both sides of the identity are complex conjugated and nothing changes.)

\emph{Plan.} A false lift has three properties (\S\ref{ppe:ss-three}): (F1) $D_\vartheta\Xi=0$;
(F2) $\Xi$ is killed by the hyperplane sums at the chosen places; (F3) $P_w(X_w)\Xi=0$.
From (F1), (F2) and $\psi_w\ne1$ on $\mathsf Q$ the descents of every $Y_\zeta$ to $W$ vanish
(\S\ref{ppe:ss-bottom}); by (F3) they vanish to $W_p$ for $2r+2p+3<n$
(\S\ref{ppe:ss-planes}); and the vanishing of descents to $\ppeH^{N''}$ kills the periods over
$\auU(V^\circ)$ (\S\ref{ppe:ss-periods}).

\subsubsection{The three properties of a false lift}\label{ppe:ss-three}

For a chosen place $w\in\mathsf M$ and a row $\lambda\in\auF_w^r\smallsetminus\{0\}$ let
$\ppeN_\lambda=\{s\in\auGL_r(\auF_w):\ s=1+u\lambda\text{ for some column }u\}$; it is a subgroup
($(1+u\lambda)(1+u'\lambda)=1+(u+u'+u(\lambda u'))\lambda$, and the inverse of $1+u\lambda$ is
$1-(1+\lambda u)^{-1}u\lambda$). For a character $\psi$ of $\auF_w^\times$ the \emph{hyperplane
sum} is
\[
\eta_{\psi,\lambda}=\sum_{s\in\ppeN_\lambda}\psi(\det s)^{-1}\,s\ \in\C[\auGL_r(\auF_w)],
\]
viewed in $\C[\auG]$ through the factor at $w$. Note $\det(1+u\lambda)=1+\lambda u$, so
$\det\colon\ppeN_\lambda\to\auF_w^\times$ is surjective.

\begin{lemma}[Characters of label idempotents on hyperplane groups]\label{ppe:green}
Let $\auF$ be a finite field, $c$ a label of degree $r$ over $\auF$, $e_c$ the central primitive
idempotent of $\C[\auGL_r(\auF)]$ of the character $\Xi_c$, $\psi$ a character of $\auF^\times$
and $x_\psi\in\auF^\times$ its name in Green's parametrisation (the element with
$\psi=\theta_{x_\psi}|_{\auF^\times}$ for the fixed injective character $\theta$ of
$\aubar{\auF}^\times$ used to define the $\Xi_c$). If $x_\psi$ is not a root of $c$, then
$\eta_{\psi,\lambda}\,e_c=0$ for every $\lambda\ne0$.
\end{lemma}
\begin{proof}
Let $\pi_c$ be an irreducible representation with character $\Xi_c$. Since $e_c$ is central and
$\C[\auGL_r(\auF)]e_c\cong\operatorname{End}(\pi_c)$, it suffices to show $\pi_c(\eta_{\psi,
\lambda})=0$. Now $\#\ppeN_\lambda^{-1}\pi_c(\eta_{\psi,\lambda})$ is the projection of $\pi_c$
onto the subspace where $\ppeN_\lambda$ acts by $\psi\circ\det$, of dimension
$\langle\Xi_c|_{\ppeN_\lambda},\psi\circ\det\rangle_{\ppeN_\lambda}$. Choose a basis
$v_1,\dots,v_r$ with $\lambda v_k=0$ for $k\ge2$ and $\lambda v_1=1$. In it $\ppeN_\lambda$ is the
group of matrices $\begin{pmatrix}a&0\\ y&1_{r-1}\end{pmatrix}$ ($a\in\auF^\times$,
$y\in\auF^{r-1}$), the semidirect product of $\auGL_1$ (embedded as $\operatorname{diag}(a,1)$) and
the unipotent radical $U$ of the parabolic subgroup $P$ of type $(1,r-1)$ containing it. Hence
\[
\langle\Xi_c|_{\ppeN_\lambda},\psi\circ\det\rangle=\frac1{q-1}\sum_{a}\psi(a)^{-1}\,
\frac1{\#U}\sum_{y\in U}\Xi_c\bigl(\operatorname{diag}(a,1)y\bigr)
=\bigl\langle\operatorname{r}_P\Xi_c\,,\ \psi\otimes\plone\bigr\rangle_{\auGL_1\times\{1\}},
\]
where $\operatorname{r}_P\Xi_c$ is the Harish-Chandra restriction (the character of the
$U$-invariants) to the Levi subgroup $\auGL_1\times\auGL_{r-1}$, evaluated on $\auGL_1\times\{1\}$.
By the restriction formula for Green's characters (the comultiplication of the Hopf algebra of
characters of the groups $\auGL_k(\auF)$, which is the Harish-Chandra restriction and is a ring homomorphism for which the cuspidal characters are primitive, \cite[\S\S9.1--9.2, pp.~110--111]{ZelevinskyLNM}; equivalently
the Mackey formula applied to Green's description of $\Xi_c$ as a product of the cuspidal
characters of the irreducible factors of $c$, \cite[p.~403 and Thm.~13]{Green1955}), $\operatorname{r}_P\Xi_c=
\sum_{\alpha}\theta_\alpha\otimes\Xi_{c/(X-\alpha)}$, the sum over the roots $\alpha\in\auF^\times$
of $c$ (a factor of degree $\ge2$ of $c$ is cuspidal, by property (G) of {\cite[\S4.1.8]{P51}}, and contributes nothing to a restriction
with a $\auGL_1$-block). Pairing with $\psi\otimes\plone$ on $\auGL_1\times\{1\}$ gives
$\sum_\alpha\langle\theta_\alpha,\psi\rangle\,\Xi_{c/(X-\alpha)}(1)$, and
$\langle\theta_\alpha,\psi\rangle=0$ unless $\alpha=x_\psi$. So the dimension is $0$.
\end{proof}

\begin{lemma}[The three properties]\label{ppe:three}
Let $w$ be off the level and $\Xi\in\mathcal F_w(e_0,\poZ_\vartheta)$. Then
\begin{enumerate}
\item[(F1)] $D_\vartheta\Xi=0$;
\item[(F2)] for every chosen place $v\in\mathsf M$ and every $\lambda\in\auF_v^r\smallsetminus
\{0\}$, $\varrho_V(\eta_{\psi_v,\lambda})\Xi=0$;
\item[(F3)] $P_w(X_w)\Xi=0$, with $X_w$ as in Lemma~\textup{\ref{po:facts}(L3)}; $|X_w|=
q^{-r/2-(t-1)}$, $q=\#\auF_w$.
\end{enumerate}
\end{lemma}
\begin{proof}
(F1) By definition $\Xi\in\varrho_V(e_0)\poA(V;\vartheta)$ is orthogonal to $\poZ_\vartheta=
\Theta_\vartheta(e_0\poA(W_j))$; Proposition~\ref{plf:L4} gives $D_\vartheta\Xi=0$. (F3) is part of
the definition of $\mathcal F_w$. (F2) The tame part of $\mu=\mu^0_j\vartheta$ at $v\in\mathsf M$ is
that of $\mu^0_j$ ($\vartheta$ is unramified at $\mathsf M$), namely $\omega_v^{\,j}$
(Definition~\ref{pfq:def-chars}, Proposition~\ref{pfq:tame}(T3)), a power of the character $\omega_v$
of order $q_2$; so
$\psi_v^{q_2}=1$ and its name $x_{\psi_v}$ satisfies $x_{\psi_v}^{q_2}=1$ (the naming map is a
group isomorphism). Write $\Xi=\varrho_V(e_0)\Xi$ and $e_0=\sum_{(c_u)\text{ good}}\bigotimes_u
e_{c_u}$. For each good family, $c_v$ has no root in $\auF_v$ that is a $q_2$-th root of unity, so
$\eta_{\psi_v,\lambda}e_{c_v}=0$ by Lemma~\ref{ppe:green}; since $\eta_{\psi_v,\lambda}$ lives in
the factor at $v$, $\eta_{\psi_v,\lambda}\bigotimes_ue_{c_u}=0$. Hence $\eta_{\psi_v,\lambda}e_0=0$
and $\varrho_V(\eta_{\psi_v,\lambda})\Xi=\varrho_V(\eta_{\psi_v,\lambda}e_0)\Xi=0$.
\end{proof}

\begin{remark}\label{ppe:rem-T}
The only property of the tame parts used in (F2) is $\psi_v^{q_2}=1$ at the chosen places, together
with the definition of good labels, which excludes the $q_2$-th roots of unity for this reason. That
the characters $\mu^0_j$ have tame part $\omega_v^j$, with $\omega_v^{q_2}=1$, at every chosen place is
Proposition~\ref{pfq:tame}(T3); their existence is Lemma~\ref{pfq:mu} (the compatibility condition on
the roots of unity of $E$ is checked in Definition~\ref{pfq:def-chars} and Proposition~\ref{pfq:tame}(T1)).
\end{remark}

\begin{lemma}[The descent and the matrix coefficients]\label{ppe:normal}
Let $\xi_0\in W_{\lambda_V}$ be a highest weight vector and $\Xi\in\poA(V;\vartheta)$. Then
$D_\vartheta\Xi=0$ if and only if $\ppeD^W_1(\pop_{\xi_0}\otimes\Phi_\vartheta)Y_{\xi_0}=0$, where
$\pop_{\xi_0}$ is the polynomial \eqref{po:eq-poly}.
\end{lemma}
\begin{proof}
Let $F(g)=\ppeD^W_1(\pop_{\xi_0}\otimes\Phi_\vartheta)Y_{\xi_0}(g)$. Integrate first over
$\auU(V)(E^+\otimes\R)$, which is compact; with $dh_\infty$ of total mass one and $c>0$ the
volume constant,
\[
F(g)=c\sum_{h_f}\frac1{\#(\Gamma\cap h_fK_\vartheta h_f^{-1})}\int_{\auU(V)(\R)}\theta\bigl(
\omega(g,h_f)\omega(1,h_\infty)(\pop_{\xi_0}\otimes\Phi_\vartheta)\bigr)\,\aubar{\bigl\langle
\rho(h_\infty)\xi_0,\Xi(h_f)\bigr\rangle}\,dh_\infty
\]
(the integrand is invariant under $K_\vartheta$ by the definition of $\chi_\vartheta$, and the sum
runs over $\Gamma\backslash\auU(V)(\auA_f)/K_\vartheta$). By the equivariance of
$\xi\mapsto\pop_\xi$ (\S\ref{po:ss-set}), $\omega(1,h_\infty)\pop_{\xi_0}=\pop_{\rho(h_\infty)
\xi_0}$, and $\xi\mapsto\theta(\omega(g,h_f)(\pop_\xi\otimes\Phi_\vartheta))=k_{\Phi_\vartheta}
(g,h_f;\xi)$ is conjugate-linear, so $k_{\Phi_\vartheta}(g,h_f;\xi)=\langle\xi,A\rangle$ for a
unique $A=A(g,h_f)\in W_{\lambda_V}$. Also $c_{\xi_0}\Xi(h_\infty h_f)=\langle\rho(h_\infty)\xi_0,
B\rangle$ with $B=\Xi(h_f)$, $\rho$ being unitary. Schur orthogonality for the irreducible unitary
representation $\rho=\rho_{\lambda_V}$ of the compact group $\auU(V)(\R)$ gives
\[
\int_{\auU(V)(\R)}\langle\rho(h)\xi_0,A\rangle\,\aubar{\langle\rho(h)\xi_0,B\rangle}\,dh=
\frac{\|\xi_0\|^2}{\dim W_{\lambda_V}}\,\langle B,A\rangle .
\]
Hence $F(g)=c\,\|\xi_0\|^2(\dim W_{\lambda_V})^{-1}\sum_{h_f}\#(\Gamma\cap h_fK_\vartheta
h_f^{-1})^{-1}k_{\Phi_\vartheta}(g,h_f;\Xi(h_f))$; by \eqref{po:eq-descent} the sum is the complex
conjugate of $D_\vartheta\Xi(g)$, up to the factor $\mu(\det g)$ that relates $\omega^W_1$ and
$\omega_\mu$. So $F=0$ if and only if $D_\vartheta\Xi=0$.
\end{proof}

\subsubsection{The bottom of the tower}\label{ppe:ss-bottom}

At a place $u\in\mathsf S_\vartheta$ we use the coordinates $T\in M_{n\times r}(E_u)$ of (N5) on
$\auHom(W,V)\otimes_{E^+}E^+_{u}$; the local datum at $u$ is a Schwartz--Bruhat function of $T$,
and by \eqref{po:eq-w5} an integral element $k$ of $\auU(V)(E^+_u)=\auGL_n(E_u)$ acts by
$\Psi\mapsto\Psi(k^{-1}T)$ and an integral element $g$ of $\auU(W)(E^+_u)=\auGL_r(E_u)$ by
$\Psi\mapsto\Psi(Tg)$ times $\chi_u(\det g)^{-1}$ (the unramified factors are trivial on integral
elements). For $T$ we write $T'$ for the $r\times r$ block of the rows of $\iota_0$ and $T''$ for
the other rows. For a Schwartz--Bruhat function $\Psi$ on $\prod_{u\in\mathsf S_\vartheta}
M_{n\times r}(E_u)$ and a polynomial Fock vector $v$ at the infinite places let
$f(v,\Psi)=v\otimes\Psi\otimes(\text{vacuum elsewhere})\in\ppeS_\Sigma(W)$.

\begin{lemma}[Special data suffice]\label{ppe:special}
Let $X$ be a hermitian space and $Y$ continuous on $[\auU(V)]$. If $\ppeD^X_a(f(v,\Psi))Y=0$ for
all polynomial Fock vectors $v$ and all $\Psi$ that are indicator functions of cosets
$z+\varpi^i M_{n\times m_X}(\auO_{\mathsf S_\vartheta})$ ($z$ arbitrary, $i\ge0$, $\varpi$ a
fixed element with $\varpi\auO_u$ the maximal ideal at each $u\in\mathsf S_\vartheta$), then
$\ppeVan(X,a,Y)$.
\end{lemma}
\begin{proof}
For fixed $g$ the map $f\mapsto\ppeD^X_a(f)Y(g)$ is a continuous linear functional on
$\ppeS(X)$ for its Schwartz (Fr\'echet) topology: $|\theta(\omega(g,h)f)|\le C(g)\,p(f)$, $p$ a
continuous seminorm of the Schwartz topology, uniformly for $h$ in a compact fundamental
domain (moderate growth of theta functions, {\cite[Lemma~5.2]{P49}}). The space $\ppeS_\Sigma(X)$ is
the completed tensor product of the archimedean Schwartz space and of the spaces of Schwartz--Bruhat
functions below; in the archimedean factor the polynomial Fock vectors correspond, in the
Schr\"odinger model, to polynomials times the Gaussian, which are dense in the Schwartz space for its
Fr\'echet topology (Hermite expansion). The finite factors are the spaces of Schwartz--Bruhat functions at the places of $\mathsf S_\vartheta$
(by (N5) the local Weil representation at the split place below $u$ is realised on functions of
$T\in M_{n\times m_X}(E_u)$, the place $\aubar u$ carrying no further variable), and every
Schwartz--Bruhat function is a finite linear combination of indicators of such cosets. Hence the
functional vanishes on a dense subspace, so everywhere.
\end{proof}

\begin{lemma}[Harmonic generation]\label{ppe:harm}
Let $\Xi\in\poA(V;\vartheta)$ satisfy (F1) and let $\xi_0$ be a highest weight vector. Then
$\ppeD^W_1(f(v,\Phi_\vartheta))Y_{\xi_0}=0$ for every polynomial Fock vector $v$ (here
$\Phi_\vartheta$ denotes the finite part of the block datum).
\end{lemma}
\begin{proof}
Write $L(v)(g)=\ppeD^W_1(f(v,\Phi_\vartheta))Y_{\xi_0}(g)$. (a) For $\xi\in W_{\lambda_V}$,
$L(\pop_\xi)=0$: the computation of the proof of Lemma~\ref{ppe:normal} with $\pop_\xi$ in place of
$\pop_{\xi_0}$ gives $L(\pop_\xi)(g)=c\,\langle\xi_0,\xi\rangle^{\ast}(\dim W_{\lambda_V})^{-1}
\aubar{D_\vartheta\Xi(g)}\mu(\det g)$ (Schur orthogonality with $\langle\xi_0,\xi\rangle$ in place
of $\|\xi_0\|^2$, the star denoting the appropriate complex conjugate), which is $0$ by (F1).
(b) For $x$ in the Lie algebra $\mathfrak g_W$ of $\auU(W)(E^+\otimes\R)$ and $v$ smooth,
$L(\omega(x,1)v)=R_xL(v)$, the derivative of the right translate; this is \eqref{po:eq-equiv}
at the infinite places (differentiate $\theta(\omega(gk,h)f)$ in $k$). (c) By (b) and
Schur orthogonality over the compact group $\auU(V)(\R)$, $L(v)=L(e_{\lambda}v)$, where $e_\lambda$
is the projection of the Fock space onto its $\rho_{\lambda_V}^\vee$-isotypic subspace for
$\omega(1,\cdot)$ (the integral over $h_\infty$ of $\omega(1,h_\infty)v$ against the coefficient
$\aubar{c_{\xi_0}\Xi}$ only sees this component). (d) By the theory of joint harmonics of
Kashiwara--Vergne and Howe (\cite[Ch.~III, Lemma~5.3 and Corollary~5.4]{KashiwaraVergne}; \cite[Proposition~3.1, (3.10)(b)]{HoweJAMS}), the polynomial vectors of that isotypic subspace are
$U(\mathfrak g_{W,\C})\cdot\{\pop_\xi:\xi\in W_{\lambda_V}\}$: the space of joint harmonics of
type $\lambda_V$ is $\{\pop_\xi\}$ (this is the choice of $\pop_\xi$ in \S\ref{po:ss-set}, see
Lemma~\ref{plf:harm}), and the isotypic component of the polynomial Fock space is generated by
its harmonics under $U(\mathfrak g_{W,\C})$. By (a), (b), (c), $L(v)=0$ for every polynomial $v$.
\end{proof}

\begin{definition}\label{ppe:def-ratio}
Let $u\in\mathsf S_\vartheta$ and $T\in M_{n\times r}(E_u)$. $T$ has an \emph{integral block
ratio} at $u$ if $T'$ is invertible, $T''T'^{-1}$ is integral, and moreover $T''T'^{-1}\equiv0
\bmod u$ if $u\in\mathsf M$. An \emph{annihilating row} for $T$ at $u$ is an integral row
$l=(\lambda,\lambda')$ ($\lambda$ indexed by the rows of $\iota_0$) with $lT=0$, $\lambda\not
\equiv0\bmod u$, and $\lambda'\equiv0\bmod u$ if $u\in\mathsf Q$.
\end{definition}

\begin{lemma}\label{ppe:row}
If $T$ has no integral block ratio at $u$, it has an annihilating row at $u$.
\end{lemma}
\begin{proof}
Let $|\cdot|$ be the absolute value at $u$. It suffices to find a row $\mu\ne0$ with $\mu T=0$ and
an index $b$ of $\iota_0$ with $|\mu_b|=\max_a|\mu_a|$, and $|\mu_a|<|\mu_b|$ for $a$ outside
$\iota_0$ when $u\in\mathsf Q$: then $l=\mu_b^{-1}\mu$ is an annihilating row. If $T'$ is
singular take $\mu=(\lambda,0)$ with $\lambda T'=0$, $\lambda\ne0$. Otherwise let $y=T''T'^{-1}$;
by hypothesis some entry $y_{ab}$ has $|y_{ab}|>1$, or $u\in\mathsf M$ and $|y_{ab}|\ge1$. Take
$\mu=(y_a,-e_a)$ (the row $a$ of $y$ on $\iota_0$, and $-1$ at $a$); then $\mu T=y_aT'-T''_a=0$,
and an index $b$ maximising $|y_{ab}|$ works, since $|\mu_a|=1$ and the other entries outside
$\iota_0$ vanish.
\end{proof}

\begin{proposition}[Bottom of the tower]\label{ppe:bottom}
Let $\Xi\in\poA(V;\vartheta)$ satisfy (F1) and (F2), and let $\xi_0$ be a highest weight vector.
Then $\ppeVan(W,1,Y_{\xi_0})$.
\end{proposition}
\begin{proof}
For $f\in\ppeS_\Sigma(W)$ put $L_\Xi(f)=\ppeD^W_1(f)Y_{\xi_0}$ (a function of $g$; we suppress $g$
and $v$). Two rules: if $k\in\auU(V)(\auA_f)$ then, substituting $h\mapsto hk$,
\begin{equation}\label{ppe:eq-move}
L_{R_k\Xi}(f)=L_\Xi\bigl(\omega(1,k^{-1})f\bigr),
\end{equation}
because $Y_{\xi_0}$ for $R_k\Xi$ is the right translate of $Y_{\xi_0}$ for $\Xi$; and for $k\in
K_\vartheta$, $R_k\Xi=\chi_\vartheta(k)^{-1}\Xi$ (definition of $\poA(V;\vartheta)=
\poA_{\lambda_V}(K_\vartheta,\chi_\vartheta^{-1})$, \S\ref{po:ss-set}) and $Y_{\xi_0}$ is
conjugate-linear in $\Xi$, so $L_{R_k\Xi}=\chi_\vartheta(k)L_\Xi$ ($\chi_\vartheta$ is unitary) and
\eqref{ppe:eq-move} gives $L_\Xi(\omega(1,k)f)=\chi_\vartheta(k)^{-1}L_\Xi(f)$.
By Lemma~\ref{ppe:special} it suffices to show $L_\Xi(f(v,\plone_C))=0$ for $C=z+\varpi^i
M_{n\times r}(\auO_{\mathsf S_\vartheta})$ and $i$ large (a coset for small $i$ is a finite union
of cosets for large $i$). Fix $z$.

\emph{Case A: $z_u$ has an integral block ratio at every $u\in\mathsf S_\vartheta$.} Let $A\in
\auU(W)(\auA_f)$ be $1$ outside $\mathsf S_\vartheta\cup\aubar{\mathsf S_\vartheta}$ with
$A_u=(z_u')^{-1}$ ($u\in\mathsf S_\vartheta$), and $\Psi_0(T)=\Phi_\vartheta(TA)$, so that
$f(v,\Psi_0)$ is a non-zero multiple of $\omega(A^{-1},1)f(v,\Phi_\vartheta)$ (by \eqref{po:eq-w5}
at $\mathsf S_\vartheta$; at $\aubar{\mathsf S_\vartheta}$ the component of $A$ is determined by
unitarity and acts on the vacuum of the conjugate place by a scalar). Hence $L_\Xi(f(v,\Psi_0))(g)$
is a multiple of $L_\Xi(f(v,\Phi_\vartheta))(gA^{-1})=0$ by Lemma~\ref{ppe:harm} and
\eqref{po:eq-equiv}. Now $z_uA_u=\binom{1}{z''_u z_u'^{-1}}$ is integral, $\equiv\iota_0\bmod u$ at
$\mathsf M$, and has upper block $1$ at $\mathsf Q$; so $\Psi_0(z)=1$, and since $\Psi_0$ is
locally constant, $\Psi_0\equiv1$ on $C$ for $i$ large. The support of $\Psi_0$ is one orbit of
$K_\vartheta$: if $\Psi_0(x)\ne0$ then $y=xA$ is integral at $\mathsf S_\vartheta$, $y\equiv\iota_0$
at $\mathsf M$ and $\det y'$ is a unit at $\mathsf Q$, and the element $k$ with components
$\begin{pmatrix}y'&0\\y''&1\end{pmatrix}$ at $\mathsf S_\vartheta$ (and $1$ elsewhere in
$K^\circ$) lies in $K_\vartheta$ and satisfies $k\,\iota_0A^{-1}=x$; and
$\Psi_0(kT)=\chi'(k)\Psi_0(T)$ for $k\in K_\vartheta$, with $\chi'(k)=\prod_{u\in\mathsf Q}
\psi_u(\det k'_u)$, because left multiplication by $k$ multiplies the upper block of $TA$ by
$k'$ at $\mathsf Q$ ($k$ is block lower triangular there) and preserves the condition at
$\mathsf M$. Let $\Psi_1=\int_{K_\vartheta}\chi_\vartheta(k)\,\omega(1,k)\plone_C\,dk$ (probability Haar
measure; the integrand is locally constant on the compact group $K_\vartheta$, so this is a finite
sum). By the rule above, $L_\Xi(f(v,\Psi_1))=\int\chi_\vartheta(k)\chi_\vartheta(k)^{-1}
L_\Xi(f(v,\plone_C))\,dk=L_\Xi(f(v,\plone_C))$, and $\omega(1,k')\Psi_1=\chi_\vartheta(k')^{-1}
\Psi_1$ for $k'\in K_\vartheta$ (substitute $k\mapsto k'^{-1}k$). The function $\Psi_1$ depends only on the components
at $\mathsf S_\vartheta$, is supported in $K_\vartheta\cdot C\subseteq\operatorname{supp}\Psi_0$, and
transforms under $K_\vartheta$ by the same character as $\Psi_0$ (the character by which
$\omega(1,k)$ acts on $f(v,\Phi_\vartheta)$ is $\chi_\vartheta(k)^{-1}$, \S\ref{po:ss-set}, and the
same holds for $\Psi_0$). A function supported on one orbit and transforming by a character is
determined by its value at one point, so $\Psi_1=\Psi_1(z)\Psi_0$, and
$L_\Xi(f(v,\plone_C))=\Psi_1(z)L_\Xi(f(v,\Psi_0))=0$.

\emph{Case B: $z_u$ has no integral block ratio at some $u\in\mathsf S_\vartheta$.} Let
$l=(\lambda,\lambda')$ be an annihilating row (Lemma~\ref{ppe:row}) and $\bar\lambda\ne0$ the
reduction of $\lambda$. For $s\in\ppeN_{\bar\lambda}$ write $s=1+\bar v\bar\lambda$, lift $\bar v$
to an integral column $v$, and let $k_s\in K^\circ$ have component $1+\iota_0vl$ at $u$ and $1$ at
the other places of $\mathsf S_\vartheta$ (the determinant $1+lv\equiv\det s$ is a unit). Since
$lz_u=0$, $k_sz_u=z_u$, and $k_s$ is integral, so $\omega(1,k_s^{-1})\plone_C=\plone_C$ and by
\eqref{ppe:eq-move} $L_{R_{k_s}\Xi}(f(v,\plone_C))=L_\Xi(f(v,\plone_C))$. Modulo $u$,
$k_s\equiv\begin{pmatrix}s&\bar v\bar\lambda'\\0&1\end{pmatrix}$. If $u\in\mathsf M$, then
$k_s=\tilde s'\kappa$ with $\kappa$ fixing $\iota_0$ modulo $u$, so $\kappa\in K_\vartheta$ with
$\chi_\vartheta(\kappa)=1$ (the datum at $\mathsf M$ is $\plone[T\equiv\iota_0]$, invariant under
such $\kappa$ without character), and $R_{k_s}\Xi=R_{\tilde s'}\Xi=\psi_u(\det s)^{-1}
\varrho_V(s)\Xi$; summing over $s$,
\[
\#\ppeN_{\bar\lambda}\cdot L_\Xi(f(v,\plone_C))=L_{\varrho_V(\eta_{\psi_u,\bar\lambda})\Xi}
(f(v,\plone_C))=0
\]
by (F2) and linearity of $\Xi\mapsto L_\Xi$. If $u\in\mathsf Q$, then $\bar\lambda'=0$, $k_s\equiv
\operatorname{diag}(s,1)$ lies in $K_\vartheta$, and $R_{k_s}\Xi=\chi_\vartheta(k_s)^{-1}\Xi$ with
$\chi_\vartheta(k_s)=\psi_u(\det s)^{\pm1}$ (by \eqref{po:eq-w5} and the factor $\psi_u(\det T')$
of the datum); summing over $s$ gives $\#\ppeN_{\bar\lambda}L_\Xi=(\sum_s\psi_u(\det s)^{\pm1})
L_\Xi=0$, because $\det\colon\ppeN_{\bar\lambda}\to\auF_u^\times$ is a surjective homomorphism and
$\psi_u\ne1$ ($\vartheta$ is a twist).
\end{proof}

\subsubsection{Hyperbolic planes}\label{ppe:ss-planes}

Let $w$ be a split place off $\Sigma$, $q=\#\auF_w$, $K^X_w$ the stabiliser of the standard
lattice in $\auU(X)(E^+_w)\cong\auGL_{m_X}(E_w)$, and for a $K^X_w$-invariant function $F$ on
$\auU(X)(\auA)$ let $P^X_w(X')F=\sum_i(-X')^iq^{\binom i2}T^X_{w,i}F$, with $T^X_{w,i}F(g)=
\sum_{x\in K_wd_iK_w/K_w}F(gx)$, $d_i=\operatorname{diag}(\varpi_w1_i,1)$, as in (N3).

\begin{lemma}[Small roots]\label{ppe:small}
Let $F$ be a bounded $K^X_w$-invariant function on $\auU(X)(\auA)$ with $P^X_w(X')F=0$, where
$q^{m_X+1}|X'|\le1$. Then $F=0$.
\end{lemma}
\begin{proof}
The cosets $K_wd_iK_w/K_w$ correspond to the subspaces of $\auF_w^{m_X}$ of dimension $i$ (proof
of {\cite[Lemma~6.8]{P49}}), so $P^X_w(X')F(g)=\sum_{U}c_UF(gx_U)$ over all subspaces $U$, with
$c_U=(-X')^{\dim U}q^{\binom{\dim U}2}$ and $x_0=1$. By {\cite[(65)]{P49}} at $-|X'|$,
$\sum_U|c_U|=\prod_{l<m_X}(1+q^l|X'|)\le\exp(\sum_{l<m_X}q^l|X'|)\le\exp(q^{m_X}|X'|)\le
\exp(1/q)\le e^{1/2}<2$. Let $S=\sup|F|$. From $F(g)=-\sum_{U\ne0}c_UF(gx_U)$,
$S\le(\sum_U|c_U|-1)S$ with $\sum_U|c_U|-1<1$; so $S=0$.
\end{proof}

\begin{lemma}[Hecke transfer]\label{ppe:transfer}
Let $Y$ be continuous on $[\auU(V)]$ and right invariant under $K^V_w$, $f\in\ppeS_\Sigma(X)$ with
$m_X\le n$, $k=n-m_X$, $a$ unramified at $w$, and $F=\ppeD^X_a(f)Y$. Then $F$ is right invariant
under $K^X_w$, and there are $\alpha,\beta\in\C$ with $|\alpha|=q^{m_X/2}$, $|\beta|=q^{n/2}$
(independent of $Y$ and $f$) such that for every $X_0\in\C$
\[
\ppeD^X_a(f)\bigl(P^V_w(X_0)Y\bigr)=\prod_{l<k}(1-q^l\alpha X_0)\cdot P^X_{w^\ast}\bigl(q^k\alpha
\beta^{-1}X_0\bigr)F,
\]
where $w^\ast$ is $w$ or $\aubar w$ (the same for all $X_0$), $P^X_{\aubar w}$ being defined with the
identification $\auU(X)(E^+_w)\cong\auGL_{m_X}(E_{\aubar w})$.
\end{lemma}
\begin{proof}
The datum is the vacuum at $w,\aubar w$, fixed up to the unramified characters (trivial on
integral elements) by $K^X_w\times K^V_w$ (\eqref{po:eq-w5}); this gives the invariance. Since
$f$ is $K^V_w$-fixed, substituting $h\mapsto hx^{-1}$ in each term,
$\ppeD(f)(T^V_{w,i}Y)=\ppeD(T'_if)Y$ with $T'_i=\sum_{y\in K_wd_i^{-1}K_w/K_w}\omega(1,y)$; under the
identification $\auGL_n(E_w)\cong\auGL_n(E_{\aubar w})$, $g\mapsto{}^t\aubar g^{-1}$, the double coset
of $d_i^{-1}$ becomes that of $d_i$, and the vacuum in the coordinates at $w$ is the vacuum in those
at $\aubar w$. So $\ppeD(f)(P^V_w(X_0)Y)=\ppeD(P'(X_0)f)Y$ with $P'$ the operator of
{\cite[Lemma~6.8]{P49}} for the pair $(X,V)$ at $\aubar w$ (bigger space $V$ of rank $n$, smaller $X$
of rank $m_X$, $k_0=k$; that computation is local and uses only \eqref{po:eq-w5}). By
{\cite[Lemma~6.8]{P49}}, $P'(X_0)f=\prod_{l<k}(1-q^l\alpha X_0)P^-(q^k\alpha\beta^{-1}X_0)f$, where
$P^-$ is built from $\sum_{g\in K d_i^{-1}K/K}\omega(g,1)$; and $\ppeD(\omega(g,1)f)Y=R_gF$. Translating
the double cosets back gives $P^X_{w^\ast}$.
\end{proof}

\begin{lemma}[The zero term]\label{ppe:zero}
Let $p\ge1$, $W_p=W_{p-1}\oplus\ppeH$ with isotropic vectors $e=e_p$, $f_p$, and for
$f\in\ppeS(W_p)$ let $f[x]\in\ppeS(W_{p-1})$ ($x\in V(\auA)$) be its mixed model along the
polarisation $\auHom(Ee,V)\oplus\auHom(Ef_p,V)$, and $\theta_0(f)=\theta(f[0])$ its zero term. If
$\ppeVan(W_{p-1},a,Y)$, then for every $f\in\ppeS_\Sigma(W_p)$ and $g\in\auU(W_p)(\auA)$,
$\int_{[\auU(V)]}\theta_0(\omega^{W_p}_{a\upsilon_1}(g,h)f)Y(h)\,dh=0$.
\end{lemma}
\begin{proof}
Write $g=n\,m(q,g_0)\kappa$ (Iwasawa decomposition for the parabolic subgroup stabilising $Ee$,
with $\kappa$ in the standard maximal compact subgroup). Put $f'=\omega(\kappa,1)f$; it lies in
$\ppeS_\Sigma(W_p)$ (at the places outside $\Sigma$, $\kappa$ fixes the vacuum up to a scalar). By
{\cite[Lemma~3.62]{P53}} with $x=0$, the radical acts trivially on the zero slice and
$(\omega(m(q,g_0),h)f')[0]=\nu(q)\,\omega^{W_{p-1}}_a(g_0,h)(f'[0])$ with $\nu(q)=\chi_V(\det q)^{-1}
|\det q|^{n/2}$ independent of $h$ (the shift of the character by $\upsilon_1$ is the one in that
lemma). The slice $f'[0]$ lies in $\ppeS_\Sigma(W_{p-1})$ (the mixed model of the vacuum at a place
outside $\Sigma$ is the vacuum times the characteristic function of the integral points). Hence the
integral is $\nu(q)\,\ppeD^{W_{p-1}}_a(f'[0])Y(g_0)=0$.
\end{proof}

\begin{lemma}[Bounded descents]\label{ppe:bdd}
In the situation of Lemma~\textup{\ref{ppe:zero}}, let $f=v\otimes f_0$ with $v$ a polynomial Fock
vector and $f_0$ a finite datum in $\ppeS_\Sigma$ (a finite combination of finite Heisenberg
translates of the vacuum). Then $\ppeD^{W_p}_{a\upsilon_1}(f)Y$ is bounded on $\auU(W_p)(\auA)$.
\end{lemma}
\begin{proof}
Let $F=\ppeD^{W_p}_{a\upsilon_1}(f)Y$, $\Lambda=\theta$ on $\ppeS(W_p)$ and $Z=\{x\in\auHom(W_p,V)(E):
x\text{ factors through }w'\mapsto\langle w',e\rangle\}$, an isotropic subspace; the mean of
$\Lambda\circ\rho$ over $Z(\auA)/Z(E)$ is $\Lambda_0=\theta_0$. Since $F$ is left invariant under
$\auU(W_p)(E)$, $F(g)=F(\gamma g)$ for $\gamma$ rational, and by Lemma~\ref{ppe:zero}
$F(\gamma g)=\int_{[\auU(V)]}(\Lambda-\Lambda_0)(\omega(\gamma g,h)f)Y(h)\,dh$. Let $A$ be as in (R2)
of {\cite[\S3.1.6]{P53}}; $|(\Lambda-\Lambda_0)(\omega(g',h)f)|\le2Y_0\|f\|\mathrm{ht}(g')^A$ by (R2)
($\mathrm{ht}(h)=1$, $V$ being definite). If $\mathrm{ht}(g)\le2$ take $\gamma=1$. If
$\mathrm{ht}(g)>2$, there is $0\ne d\in\Lambda_g$ isotropic with $|d|_g\le2/\mathrm{ht}(g)$ (by (R3)
and the remark after it), and by Witt's theorem a $\gamma\in\auU(W_p)(E)$ with $\gamma d$ a
multiple of $e$. Since $\Lambda_{\gamma g}=\gamma\Lambda_g$ isometrically, $\gamma d\in
\Lambda_{\gamma g}$ with $|\gamma d|_{\gamma g}=|d|_g\le2/\mathrm{ht}(g)$, and
$\mathrm{ht}(\gamma g)=\mathrm{ht}(g)$, since $\Lambda_{\gamma g}=\gamma\Lambda_g$ isometrically and the
height is defined from the lattice (so (R2) at $\gamma g$ gives the factor $\mathrm{ht}(g)^A$). As
$F$ is an integral over $[\auU(V)]$, we may let $h$ run over a compact fundamental set
$\mathcal F$ ($V$ definite). Let $d_1,\dots,d_n$ be a basis of the standard lattice of $V$ and $N$
a positive integer with $Nh^{-1}d_i$ integral at the finite places for all $h\in\mathcal F$; the
elements $z_i=Nd_i\langle\cdot,\gamma d\rangle$ span $Z$ (the same line as $\langle\cdot,e\rangle$),
lie in the lattice at $k=(\gamma g,h)$, and have length $\le c\,\mathrm{ht}(g)^{-1}$ there,
uniformly in $h\in\mathcal F$. {\cite[Lemma~3.32]{P53}} with
$\delta=c\,\mathrm{ht}(g)^{-1}$ and $l=A$ gives $|(\Lambda-\Lambda_0)(\omega(\gamma g,h)f)|\le
C\,\mathrm{ht}(g)^A\delta^A\|v\|_A\le C'$. Hence $|F(g)|\le C'\sup|Y|\operatorname{vol}[\auU(V)]$.
\end{proof}

\begin{proposition}[The tower]\label{ppe:tower}
Let $\Xi\in\poA(V;\vartheta)$ satisfy
\textup{(F1)} $D_\vartheta\Xi=0$, \textup{(F2)} of Lemma~\textup{\ref{ppe:three}}, and
\textup{(F3)} $P_w(X_w)\Xi=0$ for a place $w$ off the level. Let $p\ge0$ with $2r+2p+3<n$. Then
$\ppeVan(W_p,\upsilon_1^p,Y_\zeta)$ for every $\zeta\in W_{\lambda_V}$.
\end{proposition}
\begin{proof}
First $\zeta=\xi_0$. For $p=0$ this is Proposition~\ref{ppe:bottom}. Let $p\ge1$ and assume the
claim for $p-1$. By Lemma~\ref{ppe:special} it suffices to treat $f=v\otimes f_0$ as in
Lemma~\ref{ppe:bdd}; let $F=\ppeD^{W_p}_{\upsilon_1^p}(f)Y_{\xi_0}$, bounded by Lemma~\ref{ppe:bdd}.
The function $Y_{\xi_0}$ is right invariant under $K^V_w$ ($w\notin\Sigma$) and, the Hecke operators
having integer coefficients, $P_w(\aubar{X_w})Y_{\xi_0}=\aubar{c_{\xi_0}(P_w(X_w)\Xi)}=0$. By
Lemma~\ref{ppe:transfer} with $m_X=r+2p$, $k=n-m_X$, $X_0=\aubar{X_w}$:
$\prod_{l<k}(1-q^l\alpha X_0)\,P^{W_p}_{w^\ast}(X')F=0$, $X'=q^k\alpha\beta^{-1}X_0$. Here
$|X_0|=q^{-(r/2+t-1)}$, so $|q^l\alpha X_0|\le q^{k-1+m_X/2-r/2-t+1}=q^{-(m_X-r)/2}<1$ for $l<k$ (as
$k+m_X/2=n-m_X/2$ and $m_X\ge r+2$): the product is non-zero. And $|X'|=q^{k/2}|X_0|$, so
$q^{m_X+1}|X'|=q^{m_X+1+(n-m_X)/2-r/2-t+1}=q^{(m_X+4+r-n)/2}$, and $m_X+4+r-n=2r+2p+4-n\le0$.
By Lemma~\ref{ppe:small} (at $w^\ast$), $F=0$. For general $\zeta$: for $x\in\auU(V)(E^+\otimes\R)$,
$Y_{\rho(x)\xi_0}$ is the right translate of $Y_{\xi_0}$ by $x$, and
$\ppeD(f)(R_xY)=\ppeD(\omega(1,x^{-1})f)Y$ with $\omega(1,x^{-1})f\in\ppeS_\Sigma(W_p)$; so
$\ppeVan$ holds for $Y_{\rho(x)\xi_0}$. The vectors $\rho(x)\xi_0$ span $W_{\lambda_V}$
(irreducibility) and $\zeta\mapsto Y_\zeta$ is linear (conjugate of a conjugate-linear map), so it
holds for all $Y_\zeta$.
\end{proof}

\subsubsection{From descents to periods; proof of Theorem~\ref{ppe:main}}\label{ppe:ss-periods}

\begin{proposition}[Descents and periods]\label{ppe:rallis}
Let $V=V_1\perp V_2$ be a splitting of the standard coordinates, of ranks $n_1,n_2$, with
$o_1\in M_{n\times n_1}$ the inclusion of $V_1$. Let $Y$ be continuous on $[\auU(V)]$ and right
invariant under $K''=\{k\in K^\circ:k_u=1\ (u\in\Sigma)\}$. If $\ppeVan(\ppeH^{n_1},\upsilon_1^{n_1},Y)$,
then $I(h_1)=\int_{[\auU(V_2)]}Y(\operatorname{diag}(h_1,h_2))\,dh_2=0$ for every $h_1\in\auU(V_1)(\auA)$.
\end{proposition}
\begin{proof}
Let $X=\ppeH^{n_1}$ with the Lagrangian $L=\langle e_1,\dots,e_{n_1}\rangle$, Siegel parabolic
$P=MN$, $m(q)$ ($q\in\auGL_{n_1}(\auA_E)$) and $n(z)$ ($z$ hermitian; $n(z)$ is the element
$n(0,-\jmath^{-1}z)$ of {\cite[\S3.2.3]{P53}}, with skew-hermitian argument $-\jmath^{-1}z$,
and $J_V=1$ since the standard $V$ is used, so that the character
$\psi(-\frac12\jmath\operatorname{tr}(z'x^*J_Vx))$ of {\cite[Lemma~3.62]{P53}} with $z'=-\jmath^{-1}z$ is the
one written below). Since $X=L\oplus L'$, the
mixed model of {\cite[Lemma~3.62]{P53}} (with $W_0=0$) is the Schr\"odinger model: $f\mapsto f[\,\cdot\,]$,
a function of $x\in M_{n\times n_1}(\auA_E)=\auHom(L',V)(\auA)$, and
\[
(\omega(n(z)m(q),h)f)[x]=\psi\bigl(\tfrac12\operatorname{tr}(z\,x^*x)\bigr)\,\nu(q)\,f[h^{-1}xq],
\qquad\theta(f)=\sum_{x\in M_{n\times n_1}(E)}f[x],
\]
with $\nu(q)\ne0$. The image of $\ppeS_\Sigma(X)$ consists of the functions $\phi_\infty\otimes
\phi_\Sigma\otimes\plone_{\mathrm{int}}$, $\phi_\infty$ in the Schr\"odinger image of the Fock space
(which contains all products of polynomials and the Gaussian, dense in the continuous functions on
compact sets), $\phi_\Sigma$ any Schwartz--Bruhat function at $\Sigma$, and $\plone_{\mathrm{int}}$
the characteristic function of the integral matrices outside $\Sigma$ (the vacuum of the adapted
frame). Let $F=\ppeD^X_{\upsilon_1^{n_1}}(f)Y=0$. Its Fourier coefficient along $N$ at $\beta=1$,
$\int_{N(E)\backslash N(\auA)}\psi(-\frac12\operatorname{tr}z)F(n(z)m(q))\,dz$, equals (inserting the
formula, exchanging sum and integrals by absolute convergence: $Y$ is bounded, $[\auU(V)]$ compact,
$\sum_x\sup_h|f[h^{-1}xq]|<\infty$) $\nu(q)\sum_{x^*x=1}\int_{[\auU(V)]}f[h^{-1}xq]Y(h)\,dh$. By
Witt's theorem the rational frames with $x^*x=1$ form the orbit $\auU(V)(E)\,o_1$, with stabiliser
$\auU(V_2)(E)$; unfolding (and $Y(\gamma h)=Y(h)$),
\[
0=\int_{\auU(V_2)(E)\backslash\auU(V)(\auA)}f[h^{-1}o_1q]\,Y(h)\,dh .
\]
Take $q=h_1\in\auU(V_1)(\auA)\subseteq\auGL_{n_1}(\auA_E)$; then $o_1h_1=\operatorname{diag}(h_1,1)o_1$,
and the substitution $h\mapsto\operatorname{diag}(h_1,1)h$ (which normalises $\auU(V_2)$) gives
$0=\int_{\auU(V_2)(\auA)\backslash\auU(V)(\auA)}f[h^{-1}o_1]\,J(h)\,dh$ with
$J(h)=\int_{[\auU(V_2)]}Y(\operatorname{diag}(h_1,h_2)h)\,dh_2$, a continuous function with $J(1)=
I(h_1)$, right invariant under $K''$. The map $h\mapsto h^{-1}o_1$ identifies
$\auU(V_2)(\auA)\backslash\auU(V)(\auA)$ with the closed set $\Omega$ of adelic frames with $x^*x=1$
(Witt at each place), and $J$ with a continuous function $\tilde J$ on $\Omega$. Outside $\Sigma$
the factor $\plone_{\mathrm{int}}$ restricts to the integral frames, on which $\prod_{u\notin\Sigma}
K^\circ_u$ acts transitively (if $x$ is integral with $x^*x=1$, the lattice $L_1=x\auO_u^{n_1}$ is
unimodular, so $\auO_u^n=L_1\perp L_1^\perp$ with $L_1^\perp$ unimodular of rank $n_2$; since
$E/E^+$ is unramified at all finite places (choice of $E$), unimodular hermitian lattices over
$\auO_u$ of a given rank are all isometric, as each has an orthonormal basis (proof: the reduction of the lattice is a non-degenerate hermitian space over the quadratic extension of residue fields, so it has a vector $\bar x$ with $\bar h(\bar x,\bar x)\ne0$ (otherwise polarisation and the surjectivity of the residue trace give $\bar h=0$); a lift $x$ has $h(x,x)$ a unit of the base ring, hence the norm of a unit since the extension is unramified, so after scaling $h(x,x)=1$; then the lattice is $\mathcal O x\perp(x^\perp\cap\text{lattice})$ with unimodular second summand, and one inducts on the rank; see also \cite{Jacobowitz}), and at a split place this is the
statement that $\auGL_n(\auO)$ acts transitively on bases; hence $L_1^\perp\cong\auO_u^{n_2}$, and the
isometry $o_1\auO_u^{n_1}\perp(o_1\auO_u^{n_1})^\perp\to L_1\perp L_1^\perp$ sending $o_1$ to $x$ lies
in $K^\circ_u$), so $\tilde J(x)=\tilde J_0(x_\infty,x_\Sigma)$
there. Thus $\int_{\Omega_\infty\times\Omega_\Sigma}\phi_\infty\phi_\Sigma\tilde J_0\,d\mu=0$ for all
admissible $\phi_\infty,\phi_\Sigma$, where $\Omega_\infty$ is compact ($V$ definite) and the
restrictions of the $\phi_\Sigma$ to $\Omega_\Sigma$ are all locally constant compactly supported
functions. By density, $\tilde J_0=0$ on $\Omega_\infty\times\Omega_\Sigma$; at $h=1$ the frame is
$o_1$, which is integral, so $I(h_1)=J(1)=0$.
\end{proof}

\begin{proof}[Proof of Theorem~\ref{ppe:main}]
By Lemma~\ref{ppe:three}, $\Xi$ satisfies (F1)--(F3). Let $V_1=V''$ (rank $n_1=N''=r+9m+1$),
$V_2=V^\circ$. Since $r\ge x^5$ and $m\le2\sqrt x$ ({\cite[Lemma~13.4]{P46}}),
$2r+2N''+3=4r+18m+5<6r=n$. Fix $s$, $h$, $u''$, $g$, and write
$\langle u''\odot K^\circ(g,h^\circ),\Xi(\cdot)\rangle=\sum_j\aubar{\theta^\circ(g,h^\circ;e_j)}\,
\langle u''\odot e_j,\Xi(\cdot)\rangle$ for an orthonormal basis $e_j$ of $W^\circ$ (the pairing is
conjugate-linear in the first slot and $\langle e_j,K^\circ\rangle=\theta^\circ(\cdot;e_j)$). By
Lemma~\ref{po:prod}(M1) and the equivariance of $\pop^\circ$, the function
$h^\circ\mapsto\sum_j\aubar{\theta^\circ(g,h^\circ;e_j)}c_{u''\odot e_j}\Xi(\operatorname{diag}(h,
h^\circ)\pog_s)$ on $\auU(V^\circ)(\auA)$ is invariant under $\auU(V^\circ)(\R)$ and equals the
summand of the theorem on finite points; it is left invariant under $\auU(V^\circ)(E)$ and right
invariant under $J^\circ$. Since $J^\circ$ is free, the sum over the classes is
$\operatorname{vol}^{-1}$ times its integral over $[\auU(V^\circ)]$. So it suffices that, for each
$j$ and with $\zeta=u''\odot e_j$,
\begin{equation}\label{ppe:eq-goal}
\int_{[\auU(V^\circ)]}\aubar{\theta^\circ(g,h_2;e_j)}\,c_\zeta\Xi(\operatorname{diag}(h,h_2)\pog_s)\,
dh_2=0 .
\end{equation}
Let $Y_1(h')=Y_\zeta(h'\pog_s)=\aubar{c_\zeta\Xi(h'\pog_s)}$. Since $\pog_s$ is $1$ outside
$\mathsf M\cup\aubar{\mathsf M}$, the argument at the end of the proof of
Proposition~\ref{ppe:tower} (with $\pog_s$ in place of a real element; $\omega(1,\pog_s^{-1})$
preserves $\ppeS_\Sigma$) and Proposition~\ref{ppe:tower} with $p=N''$ give
$\ppeVan(W_{N''},\upsilon_1^{N''},Y_1)$. Choose (by (W7) of \S\ref{po:ss-glue}, applied to
$W_{N''}=\ppeH^{N''}\oplus W$ and to $V=V''\perp V^\circ$) a datum $\phi''\in\ppeS_\Sigma$ for the
pair $(W,V'')$ with $\theta''(\omega''(g,h)\phi'')\ne0$. Such a datum exists. Work in a Schr\"odinger model
on $\auHom(W,V'')(\auA)$ in which the vacuum outside $\Sigma$ is the characteristic function
$\plone_{\mathrm{int}}$ of the integral points, and let $P$ be the finite set of finite places
outside $\Sigma$ at which $g$ or $h$ is not integral. For $\phi=\phi_\infty\otimes\phi_\Sigma\otimes
(\text{vacuum outside }\Sigma)$ we have $\omega''(g,h)\phi=\phi'_\infty\otimes\phi'_\Sigma\otimes
\phi_P\otimes\plone_{\mathrm{int}}$, where $\phi_P=\omega''(g,h)_P\plone_{\mathrm{int}}\ne0$ is
fixed, and, the local operators $\omega''(g,h)_\infty$, $\omega''(g,h)_\Sigma$ being invertible on
the local Schwartz spaces and $\ppeS_\Sigma$ being arbitrary at $\infty$ and $\Sigma$,
$\phi'_\Sigma$ is an arbitrary Schwartz--Bruhat function and $\phi'_\infty$ an arbitrary Schwartz
function. By strong approximation for the vector group $\auHom(W,V'')$ choose a rational $x_0$,
integral outside $\Sigma\cup P$, with $\phi_P(x_0)\ne0$, and an open compact neighbourhood
$U_\Sigma$ of $x_0$ at $\Sigma$; take $\phi'_\Sigma=\plone_{U_\Sigma}$. The datum at $P$ is the fixed
Schwartz--Bruhat function $\phi_P$, compactly supported and locally constant. Let $D$ be the set of
rational $y$ with $y_\Sigma\in U_\Sigma$, $y_P\in\operatorname{supp}\phi_P$ and $y$ integral outside
$\Sigma\cup P$. The conditions at the finite places bound the denominators and confine $y$ to a
compact open set of finite adeles, so $D$ is a finite union of cosets of a lattice in the real space
$\auHom(W,V'')(E\otimes\R)$; in particular $D$ is discrete, $x_0\in D$, and $\phi_P$ takes finitely many
values on $D$. Hence
$\theta''(\omega''(g,h)\phi)=\sum_{y\in D}\phi_P(y)\,\phi'_\infty(y)$. With
$\phi'_\infty(y)=\exp(-T|y-x_0|^2)$ the term $y=x_0$ equals $\phi_P(x_0)\ne0$, and the other terms
are bounded in total by $\max|\phi_P|\sum_{y\in D\setminus\{x_0\}}e^{-T|y-x_0|^2}$, which tends to $0$
as $T\to\infty$ ($D$ discrete and a finite union of lattice cosets). So the theta value is non-zero
for $T$ large; $\phi''=\phi$ for the corresponding $\phi_\infty$ works. Let $\phi^\circ=
\pop^\circ_{e_j}\otimes\Psi^\circ$ and $\Phi=\phi''\otimes\phi^\circ\in\ppeS_\Sigma(W)$, and put
$Y_2(h')=\theta(\omega^W(g,h')\Phi)\,Y_1(h')$. Then $Y_2$ is continuous, left invariant under
$\auU(V)(E)$, right invariant under $K''$ (the datum, $\Xi$ and $\pog_s$ are), and its descents to
$\ppeH^{N''}$ vanish: for $f'\in\ppeS_\Sigma(\ppeH^{N''})$ the kernel of $f'\otimes\Phi$ on
$W_{N''}$ at $((g',g),h')$ is the product of the kernels of $f'$ and $\Phi$ (compatibility of theta
functionals with orthogonal sums of the source space), so $\ppeD^{\ppeH^{N''}}(f')Y_2(g')=
\ppeD^{W_{N''}}(f'\otimes\Phi)Y_1((g',g))=0$. By Proposition~\ref{ppe:rallis},
$\int_{[\auU(V^\circ)]}Y_2(\operatorname{diag}(h,h_2))\,dh_2=0$. By (W7) the kernel factors,
$\theta(\omega^W(g,\operatorname{diag}(h,h_2))\Phi)=\theta''(\omega''(g,h)\phi'')\cdot\theta^\circ(g,
h_2;e_j)$ (the characters of the two factors being those of (W7)), so dividing by the non-zero first
factor, $\int\theta^\circ(g,h_2;e_j)\aubar{c_\zeta\Xi(\operatorname{diag}(h,h_2)\pog_s)}\,dh_2=0$, the
complex conjugate of \eqref{ppe:eq-goal}. The same argument applies to any open compact subgroup of
$J^\circ$ and to any $\Xi$ with (F1)--(F3).
\end{proof}

\begin{remark}[What is used]\label{ppe:rem-used}
Besides results proved in this series of papers (Proposition~\ref{plf:L4}, {\cite[Lemma~6.8]{P49}} and~{\cite[Lemmas~3.32 and~3.62]{P53}} and the facts (R1)--(R3) of {\cite[\S3.1.6]{P53}}, (W7) of
\S\ref{po:ss-glue}), the proof uses the following published results: the restriction formula for
Green's characters (Lemma~\ref{ppe:green}; \cite[\S9]{ZelevinskyLNM}, \cite[Thm.~13]{Green1955}); the generation of
the isotypic components of the polynomial Fock space by joint harmonics (Lemma~\ref{ppe:harm};
\cite[Ch.~III, Corollary~5.4]{KashiwaraVergne}, \cite[Proposition~3.1]{HoweJAMS}); Schur orthogonality; Witt's theorem; and the property
$\psi_v^{q_2}=1$ of the tame parts at the chosen places (Remark~\ref{ppe:rem-T}) (proved in Proposition~\ref{pfq:tame}).
\end{remark}

\begin{remark}[Rank one]\label{ppe:rem-rankone}
Theorem~\ref{ppe:main} and its proof are used in rank one in Proposition~\ref{psu:periods}: $W=E$
(rank $r=1$, definite), $V=V^{(z)}$ of rank $n=2z$, so that $t=n-r=2z-1$; the character $\nu\vartheta$ of a
critical pair $(\nu,z)$; the idempotent $e_1$ (the mean over $\auG_1$) in place of $e_0$; a splitting
$V=V'\perp V_1$ with $V'$ of rank $n_1=20$ in place of $V''\perp V^\circ$; an integral element $x'$
equal to $1$ outside $\Sigma$ in place of $\pog_s$; and the level of the flat pair on $V_1$ in place of
$J^\circ$. We list the steps of the proof and what each needs for $r=1$.
\begin{enumerate}
\item[(1)] \emph{Set-up} (\S\ref{ppe:ss-setup}). Nothing depends on $r$; the descents $\ppeD^X_a$ are
defined for any $X$, and an extra character of $\det g$ on $\auU(1)$ only multiplies descents.
\item[(2)] \emph{(F1).} For $r=1$ this is the hypothesis $D^{(1)}_\vartheta y=0$ of
Proposition~\ref{psu:periods}; Lemma~\ref{ppe:normal} (Schur orthogonality over the compact group
$\auU(V)(\R)$) does not use $r$.
\item[(3)] \emph{(F2).} Lemma~\ref{ppe:three} uses the label idempotents of $e_0$; for $r=1$ it is
replaced by the direct argument of Proposition~\ref{psu:periods}: $\auF_v^1=\auF_v$, the hyperplane
sum is $\eta=\sum_{s\in\auF_v^\times}\psi_v(s)^{-1}s$, and $\eta e_1=0$ as soon as $\psi_v\ne1$; here
$\psi_v=\omega_v^{\,j}\ne1$. No condition of order $q_2$ is needed in rank one.
\item[(4)] \emph{(F3).} $P_w(X_w)y=0$ is the hypothesis, with $|X_w|=q^{-r/2-(t-1)}=q^{-1/2-(2z-2)}$,
the value used in Lemma~\ref{ppe:three}(F3) and in the tower.
\item[(5)] \emph{Bottom of the tower} (Lemmas~\ref{ppe:special}, \ref{ppe:harm}, \ref{ppe:row},
Proposition~\ref{ppe:bottom}). The coordinates $T\in M_{n\times r}$ become columns $T\in M_{n\times1}$,
$T'$ is the entry of the row of $\iota_0$; the arguments (annihilating rows, the support of the block
datum, the non-triviality of $\psi_w$ at $\mathsf Q$) are written for arbitrary $r\ge1$. At
$w\in\mathsf Q$ the tame part of $\nu\vartheta$ is that of $\vartheta$, which is non-trivial ($\nu$ is
unramified at $\mathsf Q$). In detail: Lemmas~\ref{ppe:normal} and~\ref{ppe:special} do not involve
$r$ (the first is Schur orthogonality on $\auU(V)(\R)$, the second is stated for every space $X$).
Lemma~\ref{ppe:row} and Cases A and B of Proposition~\ref{ppe:bottom} are literally the same with
$r=1$: $T'$ is the first entry of the column $T$, $A_u=(z'_u)^{-1}\in E_u^\times$, the element
$k=\bigl(\begin{smallmatrix}y'&0\\y''&1\end{smallmatrix}\bigr)$ lies in $\pexK_\vartheta$, the group
$\ppeN_{\bar\lambda}$ is $\auF_u^\times$, and at $u\in\mathsf M$ the sum over $s$ is
$\varrho_V(\eta)$ with $\eta$ as in (3). In Lemma~\ref{ppe:harm} steps (a) and (c) are unchanged, and
step (d) is elementary, so that the theorem on joint harmonics is not needed: the line is positive
definite, so the Fock space of $\auHom(E,V)\otimes\R$ consists of holomorphic functions of
$(x_\tau)_\tau\in\prod_\tau\C^{2z}$, the group $\auU(V)(\R)$ acts on polynomials by linear substitution
(the splitting on $\auU(V)$ is $\upsilon^{r_-}=1$), and the polynomials that are homogeneous of degree
$d_\tau$ at each $\tau$ form an irreducible representation, different degrees giving inequivalent
ones. By \eqref{po:eq-poly}, $\pop_\xi(x)=\langle\xi,\bigotimes_\tau\aubar{x_\tau}^{\otimes l_\tau}\rangle$,
and the vectors $\bigotimes_\tau y_\tau^{\otimes l_\tau}$ span $W_\lambda=\bigotimes_\tau
\operatorname{Sym}^{l_\tau}\C^{2z}$; so $\{\pop_\xi\}$ is the whole space of polynomials of degree
$l_\tau$ at each $\tau$, which is the isotypic subspace of step (c). Hence $L(v)=L(e_\lambda v)$ with
$e_\lambda v\in\{\pop_\xi\}$, and $L(v)=0$ by (a).
\item[(6)] \emph{The tower} (Lemmas~\ref{ppe:small}--\ref{ppe:bdd}, Proposition~\ref{ppe:tower}). Stated
for arbitrary $r$; needs $2r+2p+3<n$, here with $p=n_1=20$: $45<2z$. This holds because the pairs
to which Proposition~\ref{psu:periods} is applied have $z>m$ and $m\ge\sqrt x>23$
({\cite[Lemma~13.4]{P46}}). The inequality $|q^l\alpha X_0|<1$ uses only $m_X\ge r+2$, true for $p\ge1$.
The four lemmas are applied to the spaces $W_p=E\oplus\ppeH^p$ and their statements and proofs do not
involve $W$ otherwise: Lemma~\ref{ppe:small} concerns one function on $\auGL_{m_X}(E_w)$;
Lemma~\ref{ppe:transfer} is the local computation of {\cite[Lemma~6.8]{P49}} at a place where all data
are the vacuum; Lemma~\ref{ppe:zero} uses {\cite[Lemma~3.62]{P53}} for the isotropic line $Ee_p$ of $W_p$;
and Lemma~\ref{ppe:bdd} uses (R2), (R3) of {\cite[\S3.1.6]{P53}} and {\cite[Lemma~3.32]{P53}}, which are
stated for arbitrary spaces, and Witt's theorem in $W_p$, $p\ge1$. In the proof of
Proposition~\ref{ppe:tower} the exponents are $q^{-(m_X-r)/2}$ and $q^{(2r+2p+4-n)/2}$ with $r=1$,
$n=2z$, and the passage from $\xi_0$ to all $\zeta$ is unchanged.
\item[(7)] \emph{Descents and periods} (Proposition~\ref{ppe:rallis}). Stated for an arbitrary splitting
of the standard coordinates; used with $V_1=V'$, $n_1=20$.
\item[(8)] \emph{Final step.} Let $y$, $u$, $h'$, $g$, $x'$ be as in Proposition~\ref{psu:periods},
$Y_u=\aubar{c_uy}$ and $Y_1(h)=Y_u(hx')$. By (6), $\ppeVan(W_{20},\upsilon_1^{20},Y_u)$; as
$\ppeD(f)(R_{x'}Y_u)=\ppeD(\omega(1,x'^{-1})f)Y_u$ and $\omega(1,x'^{-1})$ preserves $\ppeS_\Sigma$
($x'$ is $1$ outside $\Sigma$), also $\ppeVan(W_{20},\upsilon_1^{20},Y_1)$. The datum
$\phi^\circ=1\otimes\Psi_{q_1,1}$ of $\auHom(E,V_1)$ (constant polynomial, block datum of the flat pair)
lies in $\ppeS_\Sigma$, and so does $\Phi=\phi''\otimes\phi^\circ$ for $\phi''\in\ppeS_\Sigma$ on
$\auHom(E,V')$. By (W7) for $V=V'\perp V_1$, $\theta(\omega^E(g,[h',h_1])\Phi)=\theta''(\omega''(g,h')
\phi'')\,k_1(g,h_1)$, where $\omega^E$ is a Weil representation of $(E,V)$ whose character on $\auU(1)$
is the product of the character of $k_1$ and of that of $\omega''$; the latter is at our disposal,
since a character of $\det g$ on $\auU(1)$ only multiplies descents (step (1)). The datum $\phi''$ with
$\theta''(\omega''(g,h')\phi'')\ne0$ is constructed as in the proof of Theorem~\ref{ppe:main}; that
construction uses only that $\auHom(E,V')$ is a vector group. Then $Y_2(h)=\theta(\omega^E(g,h)\Phi)
Y_1(h)$ is continuous, left invariant under $\auU(V)(E^+)$ and right invariant under $K''$ (an element
of $K''$ fixes $\Phi$, commutes with $x'$ and lies in $\pexK_\vartheta$ with trivial character), and
$\ppeD^{\ppeH^{20}}_{\upsilon_1^{20}}(f')Y_2(g')=\ppeD^{W_{20}}_{\upsilon_1^{20}}(f'\otimes\Phi)Y_1((g',g))=0$.
Proposition~\ref{ppe:rallis} with $V_1,V_2$ there equal to $V',V_1$ here gives $\int_{[\auU(V_1)]}Y_2([h',
h_1])\,dh_1=0$, and dividing by $\theta''(\omega''(g,h')\phi'')$ we get $\int_{[\auU(V_1)]}k_1(g,h_1)\,
\aubar{c_uy([h',h_1]x')}\,dh_1=0$. The hypothesis that $\pop_u$ involves only the variables of $V'$
means that $u\in\bigotimes_\tau\operatorname{Sym}^{l_\tau}\C^{20}$, so $u$ is fixed by the block of $V_1$;
this gives the invariance of the integrand under $\auU(V_1)(E^+\otimes\R)$ and, on the left, under
$\auU(V_1)(E^+)$ (the kernel $k_1$ has trivial weight). The level on $V_1$ need not be
free, and then the sum over the classes weighted by $e_{h_1}^{-1}=\#(\text{stabiliser})^{-1}$ is the
integral over $[\auU(V_1)]$ divided by the volume of the level, which is the form in which
Proposition~\ref{psu:periods} is stated. For the summand to be a function on the classes (right
invariance under the level on $V_1$, which in \S\ref{ppe:ss-periods} comes from
Lemma~\ref{po:lattice}(d)) one needs the condition on $x'$ stated in Proposition~\ref{psu:periods}:
for $k$ in the level on $V_1$ the element $x'^{-1}\operatorname{diag}(1,k)\,x'$ lies in the level of
$V$ and acts on the block datum by the scalar by which $k$ acts on the datum on $V_1$.
\end{enumerate}
So Theorem~\ref{ppe:main} holds in rank one in the form of Proposition~\ref{psu:periods}, with (F2)
replaced by step (3).
\end{remark}



\section{Products span}\label{C2:g4}
\providecommand{\auO}{\mathcal O}
\providecommand{\auU}{\mathrm U}
\providecommand{\auGL}{\mathrm{GL}}
\providecommand{\auHom}{\operatorname{Hom}}
\providecommand{\auA}{\mathbb A}
\providecommand{\auPl}{\Phi}
\providecommand{\aubar}[1]{\overline{#1}}
\providecommand{\poA}{\mathcal A}
\providecommand{\pcMat}{\mathrm M}
\providecommand{\psvac}{{\phi^{\mathrm{vac}}}}
\providecommand{\pspF}{\mathfrak F}
\providecommand{\pspl}{\lambda}
\providecommand{\pspn}{n_*}
\providecommand{\pspm}{\mathsf m}
\providecommand{\pspip}[2]{\langle #1,#2\rangle}
\providecommand{\pspG}{G_\infty}

\subsection{Products span}\label{psp:sec}

This section proves Proposition~\ref{po:span}. We use the notation of \S\ref{po:ss-glue} and of
\S\ref{pte:ss-carrier}--\S\ref{pte:ss-third}: $W=W_j$ of rank $r$, $1\le j\le m$; the space $S$ of
signature $(2m,m)$ and the carrier $\varphi^\circ$; the padded space $\widetilde W=W\oplus(E^j)^-\oplus
E^{3m-j}$, the space $W^\flat=E^{6m+1}$, $W'=\widetilde W\oplus W^\flat$; the choices $i$ and the
functions $F_i$ of \eqref{pte:eq-F}; the definite space $V^\circ$ and the kernel $K^\circ(g,h^\circ)$ of
$(W,V^\circ)$, with values in the Weyl module of the diagrams $\mu^-_\tau$; the modulus $\pi_M^2$.

\emph{What is proved and what is not.} The argument has four steps (\S\ref{psp:ss-outline}). Steps~3
and~4 rest on three estimates which are proved here in full: the decay of Heisenberg coefficients
of polynomial Fock vectors (Lemma~\ref{psp:decay}), the comparison of a theta kernel with a Fock
coefficient on a lattice without short vectors (Proposition~\ref{psp:sparse}), and the isolation
estimate (Proposition~\ref{psp:isol}). The proof of Proposition~\ref{po:span} from these is complete
(\S\ref{psp:ss-proof}). It also uses three \emph{inputs}, collected with exact statements in
\S\ref{psp:ss-inputs}: a reproducing formula for pairings with holomorphic discrete series
coefficients in the stable range (R), elementary properties of holomorphic functions of a given
weight on $\auU(p,q)$ (H), and the fact that peaks of bounded holomorphic forms lie away from the
cusps (P). They are proved in \S\ref{psq:sec}--\S\ref{psv:sec} (Theorem~\ref{psv:main}); (P) is
Siegel's argument, cf.\ Poor--Yuen~\cite{PoorYuen} for Siegel modular forms. Everything else is proved
here.

\subsubsection{Statement}\label{psp:ss-statement}

\begin{proposition}[Products span; Proposition~\ref{po:span}]\label{psp:main}
Let $y\in e_0\poA(W_j)$. Suppose that for every choice $i$ and every class $h^\circ$ of
$\auU(V^\circ)$ at level $\pi_M^2$,
\[
\int_{[\auU(W)]}y(g)\,F_i(g)\,K^\circ(g,h^\circ)\,dg=0
\]
(a vector of the Weyl module). Then $y=0$.
\end{proposition}

This is the printed statement; the integrals converge by Proposition~\ref{po:exchange}. Since the
forms on $V^\circ$ of \S\ref{po:ss-glue} are arbitrary functions on the classes, the hypothesis is
equivalent to $\int y\,F_i\,\pspip{\zeta}{K^\circ(\cdot,\cdot)}=0$ for all $i$ and all forms $\zeta$ on
$V^\circ$, which is the hypothesis $\Theta^\circ y=0$ of (G2) in Proposition~\ref{po:glue} (there
$c_i\ne0$). So (G2) follows as written there.

\subsubsection{Notation}\label{psp:ss-notation}

$f=\#\auPl$, $t=5r$, $\pspn=2ft^3$, and $b$ is the scale of the data. For a space $X$ (with its frame)
and $g\in\auU(X)(\auA)$, $\Lambda_g$, $|\gamma|_g$ and $\lambda_1(g)$ are as in
{\cite[Definition~3.27]{P53}}, and $\pspl(X,g)=\min(1,\lambda_1(g))$. For $h\in\auU(X)_\infty=\prod_\tau
\auU(X_\tau)$ we write $|h|=\max_\tau\|h_\tau\|\ge1$ (operator norms in the frame, in which the form
is $\mathrm{diag}(\pm1)$). The Fock coordinates of $\auHom(X,X')\otimes\R$ are $x_\tau=c_\psi\tau(w)$
with $c_\psi=\sqrt{2\pi/\sqrt d}$ ({\cite[\S5.2.1]{P49}}); $|x|$ is the euclidean norm of
$(x_\tau)_\tau$.

\emph{Weights.} On $\auU(p,q)$ with the frame ordered by a total order $o$ that puts the $p$ positive
indices first, a weight is $u\in\Z^{p+q}$; it has \emph{gap} $\gamma$ if $u_a+\gamma\le u_c$ whenever
$a$ is positive and $c$ negative, and it is \emph{dominant} if it is non-increasing in the order $o$
within each sign block. For a dominant weight $u$ with gap $\ge N^2+N+1$ ($N=p+q$) let $\pspm_u$ be the
coefficient $\langle v,\pi(\cdot)v\rangle$ of a unit highest weight vector $v$ of the lowest $K$-type of the
discrete series of parameter $\lambda=u+\rho_{\mathrm c}-\rho_{\mathrm n}$ ({\cite[\S5.1]{P51}}), for which
$\Lambda^K(\lambda)=u$ (Lemma~\ref{psq:param}); so $\pspm_u(1)=1$; and for a family $u=(u_\tau)_\tau$ put $\pspm_u(h)=\prod_\tau\pspm_{u_\tau}(h_\tau)$
on $\pspG=\auU(X)_\infty$; the Haar measure on $\pspG$ is that of {\cite[\S5.4]{P51}}. On $W=W_j$ the
order $o$ is that of Definition~\ref{psv:order} (positive indices first, and within each sign block by
increasing $\varsigma_\tau$); all weights on $W$ are taken for this order.

\emph{Peaks.} A \emph{peak} of a bounded function $\Phi$ is a point $x$ with $|\Phi(x)|>\frac12\sup|
\Phi|$. Then $\Phi(x)\ne0$, $|\Phi|\le2|\Phi(x)|$ everywhere, and a point $x'$ with $|\Phi(x')|\ge|
\Phi(x)|$ is again a peak. A non-zero bounded function has a peak.

\begin{lemma}[Sizes]\label{psp:sizes}
\textup{(i)} $\pspn\le b$, and $2ft^2\le\pspn/t$. \textup{(ii)} $300\pspn+t\le\Delta/2$, $\Delta=2r^8$.
\textup{(iii)} $q_1\le x^{12}$, $d\le x^{24}$, $f\le x^{12}$, $t\le x^6$; hence
$24\,q_1^2\,d\,f\,t\le x^{80}\le b^{1/2}$. \textup{(iv)} $2^{25}\le b$, $t\ge14$, and
$100(b^{3/2}+4\pspn)\le b^2$. \textup{(v)} $5\,b^{12.75}\le b^{13}$.
\end{lemma}
\begin{proof}
By (S2), (S3) and (S1): $r\le3x^5$, $t\le15x^5\le x^6$, $f\le x^{12}$, so $\pspn\le2x^{12}x^{18}\le
b$; and $2ft^2=\pspn/t$. (ii): $300\pspn+t\le600\cdot125fr^3+5r\le r^8$ because $3r^3f\le r^7$ (S7)(b)
gives $75000fr^3\le25000r^7\le r^8/2$ ($r\ge x^5\ge10^5$). (iii): (S3); $\sqrt d\le\mathrm{rd}_E\le x^{12}$
by (S7)(e) and (S4); then $24q_1^2dft\le24x^{24+24+12+6}\le x^{80}$, and $x^{160}\le b$ by (S1).
(iv): $\log b\ge128$ (S1); $t\ge r\ge100$; $b^{3/2}+4\pspn\le5b^{3/2}$ by (i), and $500\le b^{1/2}$.
(v): $5\le b^{1/4}$.
\end{proof}

\subsubsection{The inputs}\label{psp:ss-inputs}

\emph{Holomorphic functions.} Let $X$ have signature $(p,q)$ at every place, with the order $o$ of
the frame putting the positive indices first, and let $u=(u_\tau)$ be a family of weights. A
continuous function $\phi$ on $\pspG$ is \emph{holomorphic of weight $u$} if for each $\tau$ it
extends to a holomorphic function $\tilde\phi$ on a neighbourhood of $\pspG\cdot B$ in
$\prod_\tau\auGL_N(\C)$, $B$ the group of invertible matrices upper triangular for $o$, with
$\tilde\phi(hb)=\prod_\tau\prod_ab_{\tau,aa}^{u_{\tau,a}}\phi(h)$. A function $\Phi$ on
$\auU(X)(\auA)$ is \emph{holomorphic of weight $u$} if $h\mapsto\Phi(gh)$ is, for every $g$. Let
$\mathrm{rot}(\theta)$, $\theta\in\R/\Z$, be the element $\mathrm{diag}(e^{2\pi i\theta o_a})_a$ at
every place.

\begin{proposition}[Inputs]\label{psp:inputs}
\textup{(H)} Let $X$ be of rank $\le t$.
\textup{(H1)} If $\phi,\phi'$ are holomorphic of any weights $u,u'$, then $\phi\phi'$ is holomorphic of
weight $u+u'$ \textup{(}Lemma~\ref{psq:H12}\textup{(i))}; and for every dominant weight $u$ of gap $\ge N^2+N+1$
\textup{(}$N$ the rank at each place\textup{)}, for which $\pspm_u$ is defined,
$\pspm_u$ is holomorphic of weight $u$ \textup{(}Proposition~\ref{psq:mhol}\textup{)}. In \textup{(H2)}, \textup{(H3)} let $u$ be dominant of gap $\ge300\pspn$. \textup{(H2)} If $\phi$ is holomorphic of
weight $u$ then $\phi(h\,\mathrm{rot}(\theta))=e^{2\pi i\theta\langle o,u\rangle}\phi(h)$ with
$\langle o,u\rangle=\sum_{\tau,a}o_au_{\tau,a}$, and
$\int_{\R/\Z}\phi(\mathrm{rot}(\theta)h)\,e^{-2\pi i\theta\langle o,u\rangle}d\theta=\phi(1)\,\pspm_u(h)$.
\textup{(H3)} If $\phi$ is holomorphic of weight $u$ and $\phi\,\aubar{\pspm_u}$ is integrable, then
$\int_{\pspG}\phi\,\aubar{\pspm_u}=\phi(1)\int_{\pspG}|\pspm_u|^2$.
\textup{(H4)} For $y\in e_0\poA(W_j)$ and every combination $F$ of the $F_i$, $\aubar y$ is
holomorphic of a weight $u$ and $F$ of the weight $u-u^\circ$, where $u^\circ$ is the weight of
$g\mapsto K^\circ(g,h^\circ)$; $u,u^\circ$ are dominant with gap $\ge300\pspn$, $|u_{\tau,a}|\le5b$.
The carrier $\varphi^\circ$ is holomorphic of weight $(0;3m)$ (entry $0$ on the positive, $3m$ on
the negative indices of $S$). \textup{(H5)} Let $\widetilde P$ be the fixed product of minors of
\S\ref{pte:ss-third} on $\auHom(S,\widetilde W)$ and $c_{\widetilde P}(h',h)=\pspip{v}{\Omega(h',h)v}$
its Fock coefficient \textup{(}$v$ the Fock vector of $\widetilde P$, $\Omega$ the real Weil
representation with the determinant twists of \eqref{pte:eq-F}\textup{)}. Then
$c_{\widetilde P}(h',1)=c_{\widetilde P}(1,1)\,\pspm_{u_2}(h')\,\aubar{\pspm_{u_3}(h')}$ for the
weights $u_2=(\Delta';\mu^+_{1}+r+2m)$, $u_3=(0;\mu^+_1+N''+2m)$, where $\Delta'=\Delta/2-(t+1)/2$ and
$\mu^+_1=\mu^+_{1,\tau}$ is the number of columns of $\mu^+_\tau$; and $u_3=(0;3m)+u_1+u_2$ with
$u_1=(-\Delta';6m+1)$.

\textup{(R)} Let $X$ be a space of rank $n_X\le t$, $u$ a dominant weight on $X$ of gap $\ge300\pspn$,
and $Y$ a positive definite space of rank $>2n_X$. There are finitely many classes $h_p$ of $\auU(Y)$
at level $\pi_M^2$, polynomials $Q_p$ on $\auHom(X,Y)\otimes\R$ in the variables of the positive
rows of $X$ of the weight $u$, and constants $a_p$, such that for every $g\in\auU(X)(\auA)$ and every
rapidly decreasing $\psi$ on $[\auU(X)]$ right invariant under the principal congruence subgroup of
level $\pi_M^2$,
\[
\int_{\pspG}\psi(gh)\,\pspm_u(h)\,dh=\sum_pa_p(g)\int_{[\auU(X)]}\psi(x)\,\theta^{X,Y}\bigl(x,h_p;
[Q_p,\delta_\iota]\bigr)\,dx ,
\]
with the point datum $\delta_\iota$ modulo $\pi_M^2$ at the inclusion. Here the characters of the kernel are
unramified outside $(\pi_M^2)$ and tamely ramified at the primes of $\pi_M$ (in the applications they are unramified at every finite place, Theorem~\ref{pss:a}), and the index $p$ runs over pairs (class, element of a basis of the space
of the $Q$'s), as in Theorem~\ref{psr:R}. Moreover the $Q_p$ may be taken
in any space of polynomials stable under $\auU(Y)_\infty$ containing a non-zero holomorphic vector of
weight $u$ (Definition~\ref{psr:def-hol}): for $(X,Y)=(S,W^\flat)$, $u=u_1$, the products of $2m$-minors with the short rows of the
classes of $\auU(W^\flat)$ (the polynomials $Q_r$ of \S\ref{pte:ss-third}); for $(X,Y)=(W,V^\circ)$,
$u=u^\circ$, the coordinates of the kernel $K^\circ$ against a basis of the Weyl module.

\textup{(P)} Let $\Phi$ be a bounded function on $\auU(X)(\auA)$, left invariant under $\auU(X)(E^+)$,
right invariant under the principal congruence subgroup of a level $\mathfrak p$ generated by
$p\in\auO_E$ with $\aubar p=-p$, $0<\Iml\tau(p)\le A'$ for all $\tau$, $A'\ge1$, and holomorphic of a
weight with gap $\ge1$ and entries in $[-A,A]$, $A\ge2$. Then every peak $g$ of $\Phi$ has
$\pspl(X,g)\ge(AA')^{-1/2}$.
\end{proposition}
\noindent Proposition~\ref{psp:inputs} is proved: (P), (H1)--(H3) in \S\ref{psq:sec}, (R) in \S\ref{psr:sec} with its
hypotheses verified in \S\ref{pss:sec} and \S\ref{psv:sec}, (H5) in \S\ref{pss:sec}, (H4) in \S\ref{pss:sec} and
\S\ref{psv:sec}; see Theorem~\ref{psv:main} for the list of published results used.

\emph{Archimedean bounds.} By {\cite[\S\S5.2--5.6]{P51}}, whose proofs use only that the
gap $h$ of the parameter satisfies $h\ge N^2+1$ and the bound $L$ on its entries, for a dominant
weight $u$ on $\auU(p,q)$, $N=p+q\le t$, of gap $\ge300\pspn$, with $l=\lambda$ the parameter of
Lemma~\ref{psq:param}, of gap $h\ge300\pspn-t$ and entries $\le5b+t$:
\begin{equation}\label{psp:eq-ds}
|\pspm_u(g)|\le\bigl(2/\|g\|\bigr)^c\quad(0\le c\le h),\qquad
\int|\pspm_u|^2=d_l^{-1}>0,\qquad d_l\int\bigl(2/\|g\|\bigr)^{10N^3}dg\le2^{13t^3}b^{14t^2}
\end{equation}
when moreover $|u_a|\le5b$; the last bound is {\cite[Theorem~5.9]{P51}} with $\mathsf s=10N^3>4N^2$,
$L\le10b+t$, $\delta^{-1}=12N(L+2N)$: $4(5/\delta)^{2N^2}\mathsf s/(\mathsf s-4N^2)\cdot2^{10N^3}\le
8(60t(10b+3t))^{2t^2}2^{10t^3}\le2^{13t^3}b^{14t^2}$ (as $60t(10b+3t)\le b^7$, Lemma~\ref{psp:sizes}).
Here $\pspm_u$ is the coefficient $c_\tau$ of {\cite[\S5]{P51}} for the parameter $l$, and
$\int|\pspm_u|^2=d_l^{-1}$ is (DS1)(ii); $h\ge300\pspn-t\ge10N^3+6\pspn$, and the entries of $l$ are at most
$5b+t\le10b+t$, so the bounds above, and those used below ($c\le10N^3+6\pspn\le h$, $L\le10b+t$), hold.

\begin{lemma}[Integrals against $\pspm_{u'}$]\label{psp:int}
Let $X$ have rank $N\le t$, and let $u,u'$ be dominant of gap $\ge300\pspn$ with $|u_{\tau,a}|\le5b$.
Put $I_u=\int_{\pspG}|\pspm_u|^2>0$. Then $h\mapsto|\pspm_{u'}(h)|\,|h|^{6\pspn}$ is integrable and
$\int_{\pspG}|\pspm_{u'}(h)||h|^{6\pspn}dh\le b^{\pspn}I_u$. Consequently, if $A$ is continuous with
$|A(h)|\le C|h|^{6\pspn}$, then $A\pspm_{u'}$ is integrable and $|\int A\pspm_{u'}|\le Cb^{\pspn}I_u$.
\end{lemma}
\begin{proof}
Let $\tau_0$ be a place with $\|h_{\tau_0}\|=|h|$. By \eqref{psp:eq-ds} with $c=10N^3$ at $\tau\ne\tau_0$
and $c=10N^3+6\pspn\le300\pspn$ at $\tau_0$,
$|\pspm_{u'}(h)||h|^{6\pspn}\le2^{6\pspn}\prod_\tau(2/\|h_\tau\|)^{10N^3}$. Write
$2^{6\pspn}=2^{6\pspn}I_u\prod_\tau d_{l_\tau}$ ($l_\tau$ the parameter of Lemma~\ref{psq:param} for $u_\tau$,$I_u=\prod_\tau d_{l_\tau}^{-1}$).
By Fubini and \eqref{psp:eq-ds}, the integral is at most
$2^{6\pspn}I_u(2^{13t^3}b^{14t^2})^f=I_u\,2^{25ft^3}b^{14ft^2}\le I_u\,b^{ft^3}b^{ft^3}=b^{\pspn}I_u$,
using $2^{25}\le b$ and $14\le t$ (Lemma~\ref{psp:sizes}(iv)). The consequence is immediate.
\end{proof}

\subsubsection{The Fock model; decay of Heisenberg coefficients}\label{psp:ss-fock}

Let $\mathcal V=\C^D$ with a hermitian form $\pspip{x}{y}_\varepsilon=x^*\varepsilon y$,
$\varepsilon=\mathrm{diag}(\pm1)$, the symplectic form $\Iml\pspip xy_\varepsilon$, and the euclidean
norm $|x|$. Let $\mu$ be the Lebesgue measure normalised by $\int e^{-|x|^2}d\mu=1$, so that
$\int e^{-s|x|^2}d\mu=s^{-D}$ for $s>0$. As in {\cite[Lemma~3.30]{P53}}, the Fock space $\pspF$ is the
space of $c\in L^1(\mu)\cap L^2(\mu)$ with $c(x)=\int c(y)\beta(y,x)e(x-y)\,d\mu(y)$, where
$e(x)=e^{-|x|^2/2}$ is the vacuum and $\beta(y,x)=\exp(i\,\Iml\pspip yx_\varepsilon)$ (this exact
normalisation, matched with $e$ and $\mu$, is the one for which the reproducing identity, the
projection property and the computation of $\pspip{e}{\rho(z)e}$ in Lemma~\ref{psp:geom} hold; it is
fixed throughout); the Heisenberg translation is
$(\rho(y)c)(x)=\beta(y,x)c(x-y)$, and $\pspip{c}{c'}=\int\aubar cc'\,d\mu$. The polynomial vector of a
polynomial $P$ is $v_P(x)=P(\zeta(x))e(x)$, $\zeta(x)_a=\aubar{x_a}$ at the positive and $x_a$ at the
negative coordinates. The defining identity says $c(x)=\pspip{\rho(x)e}{c}$ (as $\aubar{\beta(x,y)}=
\beta(y,x)$ and $e$ is even), and $c=\int c(x)\rho(x)e\,d\mu(x)$ ({\cite[Lemma~3.30]{P53}}). The group
$\auU(\varepsilon)$ of complex matrices preserving $\pspip{}{}_\varepsilon$ preserves $\mu$ (as
$|\det\sigma|=1$) and the symplectic form; its Weil representation $\omega$ on $\pspF$ is unitary
and satisfies $\omega(\sigma)\rho(y)\omega(\sigma)^{-1}=\rho(\sigma y)$ ((W1)); it is determined up to
a character of modulus one, which does not affect the moduli below. $\|\sigma\|$ is the operator
norm; $\|\sigma\|\ge1$ since $\sigma^{-1}=\varepsilon\sigma^*\varepsilon$.

\begin{lemma}[Geometric form of the Weil action]\label{psp:geom}
For $\sigma\in\auU(\varepsilon)$ put $A_\sigma c=\int c(x)\,\rho(\sigma x)e\,d\mu(x)$. Then
$A_\sigma=\kappa(\sigma)\,\omega(\sigma)$ with $|\kappa(\sigma)|\ge\|\sigma\|^{-2D}$.
\end{lemma}
\begin{proof}
$A_\sigma$ is the composite of $c\mapsto c\circ\sigma^{-1}$, unitary on $L^2(\mu)$, with the operator
$c\mapsto\int c(x)\rho(x)e\,d\mu(x)$, which is the orthogonal projection onto $\pspF$ (its kernel is
$\aubar{(\rho(x')e)(x)}$-reproducing); so it is bounded. Since $\sigma$ preserves the symplectic form,
$\rho(\sigma y)\rho(\sigma x)=\beta'(y,x)\rho(\sigma(x+y))$ with the same cocycle $\beta'$ as for
$\rho(y)\rho(x)$; hence $\rho(\sigma y)A_\sigma\rho(y)^{-1}=A_\sigma$ (change of variables $x\mapsto x+y$
in $\int(\rho(y)^{-1}c)(x)\rho(\sigma x)e\,d\mu$). With (W1), $A_\sigma\omega(\sigma)^{-1}$ commutes with
all $\rho(y)$, so it is a scalar $\kappa(\sigma)$ by the irreducibility of $\rho$ on $\pspF$ (Stone--von
Neumann). Next $\pspip{e}{\rho(z)e}=\int e(x)\beta(z,x)e(x-z)\,d\mu(x)=e^{-|z|^2/2}$: completing the
square, $x=z/2+w$, the phase is $e^{i\Iml\pspip zw_\varepsilon}$ (as $\pspip zz_\varepsilon\in\R$) and
$\int e^{-|w|^2+i\Iml\pspip zw_\varepsilon}d\mu(w)=e^{-|z|^2/4}$. Hence
\[
\kappa(\sigma)\pspip{e}{\omega(\sigma)e}=\pspip{e}{A_\sigma e}=\int
e(x)\pspip{e}{\rho(\sigma x)e}\,d\mu(x)=\int e^{-(|x|^2+|\sigma x|^2)/2}d\mu(x)\ge
\Bigl(\frac{2}{1+\|\sigma\|^2}\Bigr)^D ,
\]
since $e$ is real. As $|\pspip e{\omega(\sigma)e}|\le1$ and $(1+\|\sigma\|^2)/2\le
\|\sigma\|^2$, $|\kappa(\sigma)|\ge\|\sigma\|^{-2D}$.
\end{proof}

\begin{lemma}[Decay]\label{psp:decay}
Let $P$ be a polynomial with $P(sx)=s^dP(x)$ for real $s$, $v=v_P$, $\sigma\in\auU(\varepsilon)$,
$y\in\mathcal V$, $y\ne0$, and let $A\ge d+D$ and $n\ge0$ be integers. Then
\[
\bigl|\pspip{v}{\rho(y)\omega(\sigma)v}\bigr|\le\|\sigma\|^{2D+2n}\,\bigl(4(A+n)\bigr)^n\,|y|^{-2n}\,
\pspip vv .
\]
\end{lemma}
\begin{proof}
By Lemma~\ref{psp:geom} and $\omega(\sigma)v=\kappa(\sigma)^{-1}\int v(x)\rho(\sigma x)e\,d\mu(x)$,
\[
\pspip{v}{\rho(y)\omega(\sigma)v}=\kappa(\sigma)^{-1}\int v(x)\pspip{v}{\rho(y)\rho(\sigma x)e}\,
d\mu(x),
\]
and $\rho(y)\rho(\sigma x)$ is a unimodular multiple of $\rho(\sigma x+y)$, while $|\pspip
v{\rho(z)e}|=|v(z)|$. So the left side is at most $\|\sigma\|^{2D}\int|v(x)||v(\sigma x+y)|\,d\mu(x)$.
Put $c=\sqrt{1+n/A}$ and $G(x)=c^{2d}|P(\zeta(x/c))|^2e^{-|x|^2/c^2}$. By homogeneity
$|v(x)|=G(x)^{1/2}e^{-a|x|^2/2}$ with $a=1-c^{-2}=n/(n+A)$, and by the change of variables $x\mapsto
cx$, $\int G\,d\mu=c^{2d+2D}\int|v|^2d\mu\le e^{n(d+D)/A}\pspip vv\le e^n\pspip vv$. Since $|y|\le
\|\sigma\||\sigma x+y|+\|\sigma\||x|$ (as $|y|\le|\sigma x+y|+|\sigma x|$ and $\|\sigma\|\ge1$),
$|y|^2\le2\|\sigma\|^2(|x|^2+|\sigma x+y|^2)$, so $\frac a2(|x|^2+|\sigma x+y|^2)\ge
T:=\frac{n|y|^2}{4(n+A)\|\sigma\|^2}$. With $2st\le s^2+t^2$ and $e^{-T}\le(n/eT)^n$,
\[
|v(x)||v(\sigma x+y)|\le\tfrac12\bigl(G(x)+G(\sigma x+y)\bigr)e^{-T}\le\tfrac12\bigl(G(x)+G(\sigma
x+y)\bigr)\Bigl(\frac{4(n+A)\|\sigma\|^2}{e|y|^2}\Bigr)^n .
\]
The map $x\mapsto\sigma x+y$ preserves $\mu$, so integrating gives at most $e^n\pspip vv\cdot
(4(n+A)\|\sigma\|^2/e|y|^2)^n$, which is the claim. (For $n=0$ read $e^{-T}\le1$.)
\end{proof}

\begin{lemma}[Separated sets]\label{psp:pack}
Let $\mathcal Y\subset\R^k$ be a set with $|y|\ge m$ and $|y-y'|\ge m$ for distinct $y,y'\in\mathcal Y$,
$m>0$, and let $s\ge k+2$ be an integer. Then $\sum_{y\in\mathcal Y}|y|^{-s}\le(5/m)^s$.
\end{lemma}
\begin{proof}
The balls $B(y,m/2)$ are disjoint and contained in $\{|z|\ge|y|/2\}$ (as $|y|\ge m$) and in
$\{|z|\le\frac32|y|\}$. For $z\in B(y,m/2)$, $|z|\le\frac32|y|$, so $|y|^{-s}\le(3/2)^s|z|^{-s}$, and
\[
\sum_y|y|^{-s}\le\frac{(3/2)^s}{\mathrm{vol}\,B(m/2)}\int_{|z|\ge m/2}|z|^{-s}dz=
\frac{(3/2)^s\,k\,\mathrm{vol}\,B(1)}{\mathrm{vol}\,B(1)(m/2)^k}\cdot\frac{(m/2)^{k-s}}{s-k}
=\frac{k}{s-k}\Bigl(\frac{3}{m}\Bigr)^s\le\Bigl(\frac5m\Bigr)^s ,
\]
since $k/(s-k)\le k/2\le(5/3)^s$ for $s\ge k+2$.
\end{proof}

\subsubsection{Theta kernels on lattices without short vectors}\label{psp:ss-sparse}

{\emergencystretch=6em Let $X,Y$ be spaces of ranks $n,n'$ and $\mathsf H=\auHom(X,Y)$. For $(g,h)\in\auU(X)(\auA)\times
\auU(Y)(\auA)$ the lattice of $\mathsf H$ at $(g,h)$ is $\Lambda_{g,h}=\{z\in\mathsf H(E):\allowbreak h^{-1}zg
\text{ is integral at all finite places}\}$, with $|z|_{g,h}$ the euclidean norm of the archimedean
coordinates of $h^{-1}zg$, and $\lambda_1(\mathsf H;g,h)$ its first minimum. Put $N_d=q_1d$; then
$N_d\in\mathfrak d_E$, because $\mathfrak d_E=\mathfrak d_{E/K}\mathfrak d_K$, $\mathfrak d_K=
(\sqrt{-d})$ divides $d$, and $\mathfrak d_{E/K}=(\Phi'_{q_1}(\zeta))$ (proof of
{\cite[Lemma~5.10]{P51}}) divides $q_1$, as $\Phi_{q_1}$ divides $X^{q_1}-1$.\par}

\begin{lemma}[First minimum of $\auHom$]\label{psp:hom}
$\lambda_1(\mathsf H;g,h)\ge\pspl(X,g)\,\pspl(Y,h)/(fnN_d)$.
\end{lemma}
\begin{proof}
Identify $X(E)\otimes_\Q\R$ with $\prod_\tau\C^n$ by $\gamma\mapsto((g^{-1}\gamma)_\tau)_\tau$, with the
euclidean product $\Rel\sum_\tau a_\tau^*b_\tau$; $\Lambda_g$ becomes a lattice $L$ of rank $2fn$. At
the infinite places $g$ is unitary and the form is $\varepsilon=\mathrm{diag}(\pm1)$, so
$\mathrm{Tr}_{E/\Q}\pspip{\gamma}{\gamma'}=\Rel\sum_\tau a_\tau^*(2\varepsilon b_\tau)$. The hermitian
lattice $\Lambda_g$ is unimodular, so its dual for $\mathrm{Tr}_{E/\Q}\pspip{}{}$ is $\mathfrak
d_E^{-1}\Lambda_g$, and the euclidean dual lattice is $L^\vee=2\varepsilon(\mathfrak d_E^{-1}
\Lambda_g)$. As $\varepsilon$ is an isometry and $N_d\mathfrak d_E^{-1}\Lambda_g\subseteq\Lambda_g$,
$\lambda_1(L^\vee)\ge2\lambda_1(g)/N_d$. By Banaszczyk's transference theorem~\cite[Theorem~2.1]{Banaszczyk}
($\lambda_k(L)\lambda_1(L^\vee)\le k$ in rank $k$), the $2fn$-th successive
minimum satisfies $\lambda_{2fn}(L)\le fnN_d/\lambda_1(g)$. Choose $\Z$-linearly independent
$Y_1,\dots,Y_{2fn}\in\Lambda_g$ with $|Y_k|_g\le fnN_d/\lambda_1(g)$; they span $X(E)$ over $\Q$. Let
$0\ne z\in\Lambda_{g,h}$; then $zY_k\ne0$ for some $k$, and $zY_k\in\Lambda_h$ since $h^{-1}zY_k=
(h^{-1}zg)(g^{-1}Y_k)$. At each $\tau$ the operator norm of $(h^{-1}zg)_\tau$ is at most its euclidean
norm, hence at most $|z|_{g,h}$, so $\pspl(Y,h)\le\lambda_1(h)\le|zY_k|_h\le|z|_{g,h}|Y_k|_g\le
|z|_{g,h}fnN_d/\lambda_1(g)$. Since $\lambda_1(g)\ge\pspl(X,g)$, the claim follows.
\end{proof}

\begin{lemma}[Products of theta functionals]\label{psp:prodtheta}
Let $\mathcal W$ be a symplectic space over $E^+$ and $F,F'$ data (smooth at the infinite places).
Then $\theta(F)\,\aubar{\theta(F')}=\sum_{y\in\mathcal W(E^+)}\pspip{F'}{\rho(y)F}$, the sum being
absolutely convergent.
\end{lemma}
\begin{proof}
Both sides are unchanged under a unitary isomorphism of models commuting with $\rho$ (which carries
$\theta$ to $\theta$, by the uniqueness in (W1) and the normalisation of $\theta$), so we may use the
Schrödinger model of a polarisation $\mathcal W=\mathcal X\oplus\mathcal Y$: $F\in\mathcal S(\mathcal
X(\auA))$, $\theta(F)=\sum_{\xi\in\mathcal X(E^+)}F(\xi)$,
$(\rho(a,b)F)(x)=\psi(\pspip xb+\frac12\pspip ab)F(x+a)$, with the self-dual measure $dx$, for which
$\mathcal X(E^+)\backslash\mathcal X(\auA)$ has volume $1$ and Poisson summation
$\sum_{b\in\mathcal Y(E^+)}\hat\phi(b)\psi(\pspip cb)=\sum_\xi\phi(\xi-c)$ holds,
$\hat\phi(b)=\int\phi(x)\psi(\pspip xb)dx$. For fixed $a\in\mathcal X(E^+)$ and
$\phi(x)=\aubar{F'(x)}F(x+a)$,
\[
\sum_{b}\pspip{F'}{\rho(a,b)F}=\sum_b\hat\phi(b)\psi(\tfrac12\pspip ab)=\sum_{\xi}\aubar{F'(\xi-\tfrac a2)}
F(\xi+\tfrac a2).
\]
Summing over $a$ and putting $\xi'=\xi-\frac a2$, $\xi''=\xi+\frac a2$, which is a bijection of
$\mathcal X(E^+)^2$, gives $\sum_{\xi''}F(\xi'')\sum_{\xi'}\aubar{F'(\xi')}$. Absolute convergence: the
matrix coefficient $y\mapsto\pspip{F'}{\rho(y)F}$ of two smooth (Schwartz) data is a Schwartz--Bruhat
function on $\mathcal W(\auA)$, so it is summable over the lattice $\mathcal W(E^+)$; and smooth
Fock vectors correspond to Schwartz functions under the Bargmann transform.
\end{proof}

For the modulus $\pi_M^2$ and $T\in\pcMat_{n\times n'}(\auO_E/\pi_M^2)$ let
$\Psi_T(x)=\psi_{M}(\mathrm{tr}(xT))$, where $\psi_M$ is a primitive additive character of
$\auO_E/\pi_M^2$; then $\frac1{\#}\sum_T\Psi_T(x)=[x=0]$ ($\#$ the number of $T$).

\begin{lemma}[Averaged coefficients]\label{psp:avg}
Let $v,w$ be smooth Fock vectors, $p=(g,h)$, and $z\in\mathsf H(E)$; put $z'=h^{-1}zg$. Then
\[
\frac1\#\sum_T\pspip{\omega(p)[v,\Psi_T]}{\rho(z)\,\omega(p)[w,\Psi_T]}=\chi(z)\,\pspip{v}{\rho(z'_\infty)
w}\cdot\mathbf 1\bigl[z'\in\alpha\auHom(L_X,L_Y)\text{ and }z'_{\mathfrak w}\equiv0\bmod\pi_M^2\
(\mathfrak w\mid\pi_M)\bigr],
\]
with $|\chi(z)|=1$ and $\chi(0)=1$. Every $z\ne0$ for which the right side is non-zero satisfies
$(q_1/\pi_M^2)z\in\Lambda_{g,h}$.
\end{lemma}
\begin{proof}
$\omega(p)$ is unitary and $\omega(p)^{-1}\rho(z)\omega(p)=\rho(p^{-1}z)$ with $p^{-1}z=z'$ ((W1)), so
we may take $p=1$, $z=z'$. The pairing factors over the places. At the finite places off $\pi_M$ the
datum is the vacuum $\psvac$ of $L=\alpha\auHom(L_X,L_Y)$; by (W4), $\rho(z)\psvac$ is a unimodular
multiple of $\psvac$ if $z\in L$ and lies on the line of a different character of $L$ otherwise,
hence is orthogonal to $\psvac$. At a place $v=\mathfrak w\aubar{\mathfrak w}$ with $\mathfrak w\mid
\pi_M$ we use (W5): functions on $\pcMat(E_{\mathfrak w})$, $\rho(a,0)$ the translation and $\rho(0,a')$
multiplication by $x\mapsto[x,a']$, with the datum $\Psi_T^\natural$ supported on integral matrices.
Up to a unimodular factor depending on $z$ only,
\[
\pspip{\Psi_T^\natural}{\rho(z)\Psi_T^\natural}=\int_{x,\,x+z_{\mathfrak w}\text{ integral}}
\aubar{\Psi_T(x)}\Psi_T(x+z_{\mathfrak w})[x,z_{\aubar{\mathfrak w}}]\,dx=
\mathbf 1[z_{\mathfrak w},z_{\aubar{\mathfrak w}}\text{ integral}]\;\Psi_T(z_{\mathfrak w}),
\]
because $\Psi_T$ is additive, the integral matrices have volume $1$ and are their own dual for
$[\,,\,]$ ($\psi$ is unramified at $p_0,p_1$, which do not divide the different, and $\alpha=1$
there, since the primes of $\alpha$ divide $(1+d)/4<b^{10}\le p_0,p_1$ by Lemma~\ref{psp:sizes}). The
average over $T$ of $\Psi_T(z_{\mathfrak w})$ is $[z_{\mathfrak w}\equiv0\bmod\pi_M^2]$. At $z=0$ all
factors are $1$ (unit vectors; Lemma~\ref{psq:unitvac}). For the last assertion: $q_1\alpha$ is integral ($\alpha=1/P'(\zeta)$
and $P'(\zeta)$ divides $q_1$), $\pi_M$ is a unit at the places $\aubar{\mathfrak w}$ and off $M$,
and $z_{\mathfrak w}\in\pi_M^2\pcMat(\auO_{\mathfrak w})$; so $(q_1/\pi_M^2)z$ is integral everywhere.
\end{proof}

\begin{proposition}[Sparse kernels]\label{psp:sparse}
Let $\alpha',\beta'$ be characters as in the definition of the kernel $\theta_{\alpha',\beta'}$
({\cite[\S2.3]{P53}}), $Q$ a polynomial on $\mathsf H\otimes\R$, homogeneous of degree $d_Q$ for real
scalars, $v$ its Fock vector, and $c_Q(u,u')=\pspip{v}{\Omega(u,u')v}$ for $(u,u')\in\auU(X)_\infty
\times\auU(Y)_\infty$, $\Omega$ the real Weil representation with the archimedean components of the
characters. Let $D=fnn'$, let $q>D$ be an integer and $m=c_\psi M\lambda_1(\mathsf H;g,h)/q_1$. There
is a function $\Psi$ modulo $\pi_M^2$ such that for all $u,u'$
\[
\bigl|\theta_{\alpha',\beta'}\bigl(gu,hu';[Q,\Psi]\bigr)-c_Q(u,u')\bigr|\le(|u||u'|)^{2D+2q}
\Bigl(\frac{100(d_Q+2q)}{m^2}\Bigr)^q c_Q(1,1).
\]
\end{proposition}
\begin{proof}
Let $k_T=\theta_{\alpha',\beta'}(g,h;[Q,\Psi_T])$ and $\Psi=\frac1\#\sum_T\aubar{k_T}\Psi_T$. The kernel is
linear in the finite datum, and $\omega(u,u')[v,\Psi_T]=[\omega_\infty(u,u')v,\Psi_T]$; the characters
contribute $\chi_1=\alpha'(\det g)^{-1}\beta'(\det h)$, of modulus $1$, at $(g,h)$, and their
archimedean components at $(u,u')$ are absorbed in $\Omega$. With $w=\Omega(u,u')v$,
$F_T=\omega(p)[v,\Psi_T]$, $F'_T=\omega(p)[w,\Psi_T]$, $p=(g,h)$:
\[
\theta_{\alpha',\beta'}(gu,hu';[Q,\Psi])=\frac1\#\sum_T\aubar{\chi_1\theta(F_T)}\,\chi_1\theta(F'_T)
=\sum_{z\in\mathsf H(E)}\frac1\#\sum_T\pspip{F_T}{\rho(z)F'_T}
\]
by Lemma~\ref{psp:prodtheta}. By Lemma~\ref{psp:avg} the term of $z=0$ is $\pspip vw=c_Q(u,u')$, and
the other terms vanish unless $z\in\mathcal Z=\{z\ne0:(q_1/\pi_M^2)z\in\Lambda_{g,h}\}$. For
$z\in\mathcal Z$ let $x_z$ be the Fock coordinates of $z'_\infty$; as $|\tau(q_1/\pi_M^2)|=q_1/M$,
$|x_z|\ge c_\psi(M/q_1)\lambda_1(\mathsf H;g,h)=m$, and likewise $|x_z-x_{z'}|\ge m$ for $z\ne z'$ in
$\mathcal Z$ (their difference is in $\mathcal Z$). The term of $z$ has modulus at most
$|\pspip{v}{\rho(x_z)\omega_\infty(u,u')v}|$, and $\omega_\infty(u,u')$ is the Weil operator of
$\sigma\colon x\mapsto u'xu^{-1}$, with $\|\sigma\|\le|u'||u|$ (the euclidean norm of $u'xu^{-1}$ is at
most $\|u'\|\|u^{-1}\|$ times that of $x$, and $\|u^{-1}\|=\|u\|$). The Fock model of
{\cite[\S5.2.1]{P49}} is identified with that of \S\ref{psp:ss-fock} (complex dimension $D$) by
$F\mapsto(x\mapsto\pspip{\rho(x)e}{F})$, a unitary map commuting with $\rho$ and $\omega$, carrying
$Q$ to $v_Q$ up to the conjugation of coordinates fixed there. By Lemma~\ref{psp:decay} with
$A=d_Q+D$ and $n=q$, and Lemma~\ref{psp:pack} with $k=2D$, $s=2q\ge2D+2$,
\[
|\theta-c_Q(u,u')|\le\|\sigma\|^{2D+2q}\bigl(4(d_Q+D+q)\bigr)^q\pspip vv\sum_{z\in\mathcal Z}|x_z|^{-2q}
\le(|u||u'|)^{2D+2q}\bigl(4(d_Q+2q)\bigr)^q(5/m)^{2q}c_Q(1,1),
\]
which is the claim ($4\cdot25=100$, $c_Q(1,1)=\pspip vv$).
\end{proof}

\begin{corollary}\label{psp:sparse2}
If moreover $n,n'\le t$, $d_Q\le b^{3/2}$ and $\pspl(X,g)\pspl(Y,h)\ge b^{-13}$, then $\Psi$ can be
chosen with
\[
\bigl|\theta_{\alpha',\beta'}(gu,hu';[Q,\Psi])-c_Q(u,u')\bigr|\le b^{-20\pspn}(|u||u'|)^{6\pspn}
c_Q(1,1)\qquad\text{for all }u,u'.
\]
\end{corollary}
\begin{proof}
Take $q=2\pspn$. Then $D\le ft^2\le\pspn<q$ and $2D+2q\le6\pspn$ (Lemma~\ref{psp:sizes}(i)); as
$|u|,|u'|\ge1$ we may raise the exponent to $6\pspn$. By Lemma~\ref{psp:hom},
$\lambda_1(\mathsf H;g,h)\ge b^{-13}/(ftN_d)$, so with $c_\psi=\sqrt{2\pi/\sqrt d}\ge x^{-6}\ge
b^{-1/2}$, $M\ge b^{20}$ ({\cite[Lemma~13.4(S5)]{P46}}) and Lemma~\ref{psp:sizes}(iii),
$m\ge b^{-1/2}b^{20}b^{-13}/(q_1^2dft)\ge b^{6}$. Then $100(d_Q+4\pspn)/m^2\le b^{-10}$ by
Lemma~\ref{psp:sizes}(iv), and $(b^{-10})^{2\pspn}=b^{-20\pspn}$.
\end{proof}

\subsubsection{The isolation estimate}\label{psp:ss-isol}

\begin{proposition}[Isolation]\label{psp:isol}
Let $X$ have rank $\le t$, let $u,u'$ be dominant weights of gap $\ge300\pspn$ with $|u_{\tau,a}|\le5b$,
and $I_u=\int_{\pspG}|\pspm_u|^2$. Let $\Phi$ be a continuous bounded function on $\auU(X)(\auA)$,
holomorphic of weight $u$, with a peak $z$; let $K_0$ be a function on $\pspG$ with $|K_0|\le C$, and
$\delta\ge0$. Suppose that for every peak $x$ of $\Phi$ there is a function $K_x$ on $\pspG$,
holomorphic of weight $u-u'$, with $|K_x(h)-K_0(h)|\le\delta|h|^{6\pspn}$ for all $h$ and
\[
\int_{\pspG}\aubar{\Phi(xh)}\,K_x(h)\,\pspm_{u'}(h)\,dh=0 .
\]
Then $|K_0(1)|\le(4b^{\pspn}+1)\delta$.
\end{proposition}
\begin{proof}
Put $e(\theta)=e^{2\pi i\theta\langle o,u\rangle}$. By (H2), $|\Phi(z\,\mathrm{rot}(\theta))|=|\Phi(z)|$,
so every $z\,\mathrm{rot}(\theta)$ is a peak; let $K_\theta=K_{z\mathrm{rot}(\theta)}$ and
$W_\theta(h)=\aubar{\Phi(z\,\mathrm{rot}(\theta)h)}\,e(\theta)$. Then $|W_\theta|\le2|\Phi(z)|$, $W$ is
jointly continuous, $\int W_\theta K_\theta\pspm_{u'}=0$, and by (H2) applied to $\phi=\Phi(z\,\cdot)$,
$\int_{\R/\Z}W_\theta(h)\,d\theta=\aubar{\Phi(z)\pspm_u(h)}$. Let $K_*=K_{\theta=0}$. Since
$|K_\theta(h)|\le(C+\delta)|h|^{6\pspn}$ for all $\theta$ and the $K_\theta$ are continuous,
Lemma~\ref{psp:int} shows that $K_*\pspm_{u'}$ and $K_\theta\pspm_{u'}W_\theta$ are integrable, and
\[
\Bigl|\int K_*\pspm_{u'}W_\theta\Bigr|=\Bigl|\int(K_*-K_\theta)\pspm_{u'}W_\theta\Bigr|\le
4\delta\,|\Phi(z)|\,b^{\pspn}I_u ,
\]
as $|(K_*-K_\theta)W_\theta|\le4\delta|\Phi(z)||h|^{6\pspn}$. The function $(\theta,h)\mapsto
K_*(h)\pspm_{u'}(h)W_\theta(h)$ is integrable on $\R/\Z\times\pspG$ ($W$ bounded), so by Fubini
\[
\int_{\R/\Z}\int K_*\pspm_{u'}W_\theta\,dh\,d\theta=\aubar{\Phi(z)}\int K_*\pspm_{u'}\,\aubar{\pspm_u}
=\aubar{\Phi(z)}\,K_*(1)\,I_u ,
\]
by (H3): $K_*\pspm_{u'}$ is holomorphic of weight $u$ by (H1), equals $K_*(1)$ at $1$
($\pspm_{u'}(1)=1$), and $K_*\pspm_{u'}\aubar{\pspm_u}$ is integrable as $|\pspm_u|\le1$. Hence
$I_u|\Phi(z)||K_*(1)|\le4\delta|\Phi(z)|b^{\pspn}I_u$; as $I_u>0$ and $\Phi(z)\ne0$,
$|K_*(1)|\le4\delta b^{\pspn}$, and $|K_0(1)|\le|K_*(1)|+\delta$ since $|1|=1$.
\end{proof}

\subsubsection{The four steps}\label{psp:ss-outline}

Let $\widetilde P$ be the product of minors on $\auHom(S,\widetilde W)$ of \S\ref{pte:ss-third} (the
factor of $P_i$ in the rows of $\widetilde W$), of degree $d_{\widetilde P}$, and for a choice $i$ write
$P_i=\widetilde P\cdot Q_{r(i)}$, where $Q_{r(i)}$ is the product of the determinants on $W^\flat$
attached to the class $h^\flat_i$ and the row selection of $i$; the choice $i$ is the triple
$(\Lambda_i,h^\flat_i,\text{rows})$ (\S\ref{pte:ss-third}). Let $\theta^{\widetilde W}[\Psi]$ be the kernel
of $(S,\widetilde W)$ with datum $[\widetilde P,\Psi]$, $\Psi$ a function modulo $\pi_M^2$, and with the
characters of \eqref{pte:eq-F} restricted to $\widetilde W$; write $(g,1)$ for the point $(g,1,1)$ of
$\auU(\widetilde W)$. Let $u_1=(-\Delta';6m+1)$ and $u_3$ be as in (H5).

\begin{lemma}[Weights]\label{psp:weights}
$u_1$ and $u_3$ are dominant of gap $\ge300\pspn$ on $S$, and $|u_{1,a}|,|u_{3,a}|\le5b$; and
$d_{\widetilde P}\le b^{3/2}$.
\end{lemma}
\begin{proof}
Both are constant on each sign block, so dominant. Gap of $u_1$: $\Delta'+6m+1\ge\Delta/2-(t+1)/2\ge
300\pspn$ by Lemma~\ref{psp:sizes}(ii). Gap of $u_3$: $\mu^+_1\ge\Delta'$ (the diagram $\mu^+_\tau$ has
$\Delta'$ columns of height $r_+$, \S\ref{pte:ss-third}). Entries: $\mu^+_1\le\kappa_j(\tau)\le2b$
(the largest part of $\mu^+_\tau$ is $\kappa_j(\tau)$, proof of Proposition~\ref{pte:third}, and
$\kappa_j(\tau)\le\kappa(\tau)+m\Delta\le2b$ by (S5), (S7)(g)), and $N''+2m=r+11m+1\le b$; for $u_1$, $\Delta'\le\Delta/2=r^8\le3^8x^{40}\le x^{160}\le b$ (by (S2) and (S1)) and
$6m+1\le N''+2m\le b$. The degree: $d_{\widetilde P}\le7fmb\le b^{3/2}$ by Lemma~\ref{psq:degree}.
\end{proof}

\begin{lemma}[Step 1]\label{psp:step1}
For $g_1\in\auU(S)(\auA)$ and every $\Psi$ there are constants $c_i$ with
\[
\sum_ic_iF_i(g)=\int_{\auU(S)_\infty}\varphi^\circ(g_1h)\,\theta^{\widetilde W}[\Psi]\bigl(g_1h,(g,1)\bigr)
\,\pspm_{u_1}(h)\,dh\qquad(g\in\auU(W)(\auA)).
\]
\end{lemma}
\begin{proof}
Fix $g$ and put $\psi(x)=\varphi^\circ(x)\theta^{\widetilde W}[\Psi](x,(g,1))$, rapidly decreasing
($\varphi^\circ$ is, Lemma~\ref{pte:carrierprop}, and kernels have moderate growth, {\cite[Lemma~3.29]{P53}}) and right
invariant under the principal congruence subgroup of level $\pi_M^2$ ({\cite[Lemma~2.9]{P53}}). By (R)
for $(X,Y)=(S,W^\flat)$, $\operatorname{rk}W^\flat=6m+1>6m$, $u=u_1$ (Lemma~\ref{psp:weights}), the
right side is $\sum_ra_r(g_1)\int_{[\auU(S)]}\psi(x)\theta^{W^\flat}(x,h^\flat_r;[Q_r,\delta_\iota])dx$,
$r$ running over the pairs (class, rows); the constants do not depend on $g$. The product of the
kernels of $(S,\widetilde W)$ and $(S,W^\flat)$ is the kernel of $(S,\widetilde W\oplus W^\flat)=(S,W')$
with the product datum $[\widetilde PQ_r,\Psi\otimes\delta_\iota]$ ((W7) of \S\ref{po:ss-glue}, the
compatibility with orthogonal sums of the target; the characters are those of \eqref{pte:eq-F}; the factor
$(\nu^\flat\upsilon^m)(\det x)^{-1}$ is attached to the kernel of $(S,\widetilde W)$, so that the kernel of
$(S,W^\flat)$ has the characters $(1,\upsilon^m)$, Lemma~\ref{psv:split}). By
finite Fourier analysis on $\pcMat_{3m\times N''}(\auO_E/\pi_M^2)$, $\Psi\otimes\delta_\iota=\sum_\Lambda
f(\Lambda)\psi_\Lambda$ for some $f$, and the kernel is linear in the datum. So the right side is
$\sum_{r,\Lambda}a_r(g_1)f(\Lambda)\int\varphi^\circ(x)\theta^{S,W'}(x,(g,1,1,h^\flat_r);[P_{(\Lambda,r)},
\psi_\Lambda])dx=\sum_{r,\Lambda}a_r(g_1)f(\Lambda)F_{(\Lambda,r)}(g)$ by \eqref{pte:eq-F}.
\end{proof}

\begin{lemma}[Step 2]\label{psp:step2}
Under the hypothesis of Proposition~\ref{psp:main}, for every combination $F=\sum_ic_iF_i$ and every
$x\in\auU(W)(\auA)$, $\int_{\pspG}y(xh)F(xh)\pspm_{u^\circ}(h)\,dh=0$.
\end{lemma}
\begin{proof}
$\psi=yF$ is rapidly decreasing ($y$ is cuspidal, $F$ of moderate growth) and right invariant under
the principal congruence subgroup of level $\pi_M^2$. By (R) for $(X,Y)=(W,V^\circ)$,
$\operatorname{rk}V^\circ>2r$, $u=u^\circ$, the integral is $\sum_pa_p(x)\int_{[\auU(W)]}y\,F\,
\pspip{\zeta_p}{K^\circ(\cdot,h^\circ_p)}$ for vectors $\zeta_p$ of the Weyl module, and each term is
$\sum_ic_i\pspip{\zeta_p}{\int yF_iK^\circ(\cdot,h^\circ_p)}=0$ by hypothesis (the integrals converge
absolutely, Proposition~\ref{po:exchange}).
\end{proof}

\begin{proposition}[Step 3]\label{psp:approx}
There are a bounded function $F^0$ on $\auU(W)_\infty$ with $F^0(1)\ne0$ such that for every
$g\in\auU(W)(\auA)$ with $\pspl(W,g)\ge(5b^{24})^{-1/2}$ some combination $F=\sum_ic_iF_i$ satisfies
\[
|F(gh)-F^0(h)|\le2b^{-19\pspn}|F^0(1)|\,|h|^{6\pspn}\qquad(h\in\auU(W)_\infty).
\]
\end{proposition}
\begin{proof}
\emph{A peak of the carrier.} $\varphi^\circ$ is bounded, non-zero, left invariant under the rational
points, right invariant under $K_S$ (Lemmas~\ref{pte:carrierprop} and~\ref{pte:carriernz}), hence under the principal
congruence subgroup of level $p=\pm\sqrt{-d}$, the sign chosen as in Remark~\ref{psq:Papp} ($\aubar p=-p$, $\Iml\tau(p)=\sqrt d\le x^{12}\le
b^{1/2}$), and holomorphic of weight $(0;3m)$ (H4), of gap $3m\ge1$ and entries in $[-5b,5b]$. By
(P) with $A=5b$, $A'=b^{1/2}$ it has a peak $g_1$ with $\pspl(S,g_1)\ge(5b^{3/2})^{-1/2}$.

\emph{The function $F^0$.} Let $c(h',h)=c_{\widetilde P}(h',\iota h)$, $\iota h=(h,1,1)$, the Fock
coefficient of (H5), and $J=I_{u_3}>0$. By unitarity and Cauchy--Schwarz $|c|\le c(1,1)=\pspip vv>0$
($\widetilde P\ne0$). Put
\[
F^0(h)=\int_{\auU(S)_\infty}\varphi^\circ(g_1h')\,c(h',h)\,\pspm_{u_1}(h')\,dh' .
\]
By Lemma~\ref{psp:int} (with $u=u_3$, $u'=u_1$, Lemma~\ref{psp:weights}) and $|\varphi^\circ(g_1h')|\le2
|\varphi^\circ(g_1)|$, the integral converges and $|F^0(h)|\le2|\varphi^\circ(g_1)|c(1,1)b^{\pspn}J$. By
(H5) and $u_3=(0;3m)+u_1+u_2$, the function $h'\mapsto\varphi^\circ(g_1h')\pspm_{u_1}(h')c(1,1)
\pspm_{u_2}(h')$ is holomorphic of weight $u_3$ ((H1), (H4); $\pspm_{u_2}$ is defined and holomorphic of weight $u_2$
by (H1), $u_2$ being constant on each sign block, hence dominant, and its gap being $\ge9m^2+3m+1=N^2+N+1$ with $N=3m$, as checked in the proof of
Theorem~\ref{pss:H5}), so by (H3)
$F^0(1)=J\,\varphi^\circ(g_1)\,c(1,1)\ne0$.

\emph{Approximation.} Let $g$ be as in the statement. The lattice of $\widetilde W$ at $(g,1)$ is the
orthogonal sum of $\Lambda_g$ and the standard lattice $\auO_E^{3m}$ of the padding space $(E^j)^-\oplus E^{3m-j}$, whose non-zero vectors
have length $\ge1$ (a non-zero integral vector $\gamma'$ has $\sum_\tau|\tau\gamma'|^2\ge f\ge1$ by the
arithmetic--geometric mean inequality and $|\mathrm N\gamma'_a|\ge1$); so $\pspl(\widetilde W,(g,1))\ge
\pspl(W,g)$ and $\pspl(S,g_1)\pspl(\widetilde W,(g,1))\ge(5b^{3/2})^{-1/2}(5b^{24})^{-1/2}\ge b^{-13}$
(Lemma~\ref{psp:sizes}(v)). By Corollary~\ref{psp:sparse2} for $(X,Y)=(S,\widetilde W)$ (ranks $3m$ and
$r+3m\le t$, $d_{\widetilde P}\le b^{3/2}$) there is $\Psi$ with
$|\theta^{\widetilde W}[\Psi](g_1h',(g,1)\iota h)-c(h',h)|\le b^{-20\pspn}(|h'||h|)^{6\pspn}c(1,1)$, using
$|\iota h|\le|h|$. By Lemma~\ref{psp:step1} some $F=\sum c_iF_i$ has
$F(gh)=\int\varphi^\circ(g_1h')\theta^{\widetilde W}[\Psi](g_1h',(gh,1))\pspm_{u_1}(h')dh'$, and
$(gh,1)=(g,1)\iota h$. Hence $F(gh)-F^0(h)=\int\varphi^\circ(g_1h')A(h')\pspm_{u_1}(h')dh'$ with
$|\varphi^\circ(g_1h')A(h')|\le2|\varphi^\circ(g_1)|\,b^{-20\pspn}|h|^{6\pspn}c(1,1)\,|h'|^{6\pspn}$, and
Lemma~\ref{psp:int} gives
\[
|F(gh)-F^0(h)|\le2|\varphi^\circ(g_1)|c(1,1)J\,b^{-20\pspn}b^{\pspn}|h|^{6\pspn}=2b^{-19\pspn}|F^0(1)|
\,|h|^{6\pspn}. \qedhere
\]
\end{proof}

\subsubsection{Proof of Proposition~\ref{psp:main}}\label{psp:ss-proof}

Suppose $y\ne0$ and let $\Phi=\aubar y$, a bounded continuous non-zero function on $\auU(W)(\auA)$
(cuspidal, hence bounded), holomorphic of the weight $u$ of (H4); here and below (P), (H1)--(H5) and (R) are the inputs listed in Proposition~\ref{psp:inputs}. It is right invariant under the
principal congruence subgroup of level $\pi_M$ (the level of $\poA(W_j)$, \S\ref{po:ss-set}), which
contains that of level $p=\pm M\sqrt{-d}$ (sign as in Remark~\ref{psq:Papp}); $\aubar p=-p$ and $0<\Iml\tau(p)=M\sqrt d\le b^{22}x^{12}\le
b^{23}$ ((S5)). By (P) with $A=5b$, $A'=b^{23}$, every peak $x$ of $\Phi$ has $\pspl(W,x)\ge
(5b^{24})^{-1/2}$. Let $F^0$ be as in Proposition~\ref{psp:approx}, $|F^0|\le C$, and for a peak $x$
let $K_x(h)=F(xh)$ with $F$ the combination given there for $g=x$. Then $K_x$ is holomorphic of weight
$u-u^\circ$ (H4), $|K_x-F^0|\le\delta|h|^{6\pspn}$ with $\delta=2b^{-19\pspn}|F^0(1)|$, and
$\int\aubar{\Phi(xh)}K_x(h)\pspm_{u^\circ}(h)dh=\int y(xh)F(xh)\pspm_{u^\circ}(h)dh=0$ by
Lemma~\ref{psp:step2}. $\Phi$ has a peak. Proposition~\ref{psp:isol} (with $u'=u^\circ$; the bounds on
$u,u^\circ$ are part of (H4)) gives
\[
|F^0(1)|\le(4b^{\pspn}+1)\cdot2b^{-19\pspn}|F^0(1)| .
\]
With $B=b^{\pspn}>10$, $(4B+1)\cdot2<B^2\le B^{19}$, so $(4B+1)2B^{-19}<1$, and $F^0(1)=0$, a
contradiction. Hence $y=0$. \qed

\subsubsection{What this proof rests on}\label{psp:ss-rest}

Proposition~\ref{psp:main} is proved from Lemmas~\ref{psp:sizes}--\ref{psp:step2},
Propositions~\ref{psp:sparse}, \ref{psp:isol}, \ref{psp:approx}, the cited results of the paper
({\cite[Lemma~13.4]{P46}}, {\cite[\S5]{P51}}, {\cite[Lemma~3.30]{P53}}, (W1)--(W5), (W7),
{\cite[Lemma~2.9]{P53}}, Lemmas~\ref{pte:carrierprop} and~\ref{pte:carriernz}, \eqref{pte:eq-F}, Proposition~\ref{po:exchange}),
Banaszczyk's transference theorem and the Stone--von Neumann theorem, and the inputs (R), (H), (P) of
Proposition~\ref{psp:inputs}, proved in \S\ref{psq:sec}--\S\ref{psv:sec} (Theorem~\ref{psv:main}). With (G1) and Proposition~\ref{po:span},
the injectivity of $\Theta^\circ$ in (G2) of Proposition~\ref{po:glue} holds (Remark~\ref{psv:G2}).



\providecommand{\auO}{\mathcal O}
\providecommand{\auU}{\mathrm U}
\providecommand{\auGL}{\mathrm{GL}}
\providecommand{\auA}{\mathbb A}
\providecommand{\auPl}{\Phi}
\providecommand{\aubar}[1]{\overline{#1}}
\providecommand{\pspl}{\lambda}
\providecommand{\pspG}{G_\infty}
\providecommand{\pspip}[2]{\langle #1,#2\rangle}
\providecommand{\psqH}{\mathfrak H}

\subsection{Inputs for products span}\label{psq:sec}

This section proves the inputs (P), (H1), (H2), (H3) of Proposition~\ref{psp:inputs}
(\S\ref{psq:ss-P}, \S\ref{psq:ss-H}, \S\ref{psq:ss-H3}), and four facts used in \S\ref{psp:sec}
without written proof: the normalisation of $\theta$ in Lemma~\ref{psp:prodtheta}
(Lemma~\ref{psq:thetanorm}) and the unit vacuum (Lemma~\ref{psq:unitvac}), the level of
$e_0\poA(W_j)$ used in \S\ref{psp:ss-proof} (Remark~\ref{psq:level}), and the degree of
$\widetilde P$ used in Lemma~\ref{psp:weights} (Lemma~\ref{psq:degree}). The rest of
Proposition~\ref{psp:inputs}, namely (H4), (H5), (R), is proved in \S\ref{psr:sec}--\S\ref{psv:sec}.

\emph{Conventions.} $X$ is a hermitian space of rank $n$ over $E$ with a frame as in
{\cite[Definition~3.27]{P53}}: in the frame the form has the matrix $\varepsilon=\mathrm{diag}(s_1,\dots,
s_n)$, $s_a=\pm1$, at every infinite place, and at the finite places a matrix $S_f\in\auGL_n(\widehat
\auO_E)$ (the lattice of integral vectors is self-dual; this is the unimodularity used in fact (R3) stated after
{\cite[Definition~3.27]{P53}}). Forms are linear in the first variable ((Cv1)): $\langle w,w'\rangle=w'^*Sw$ in
coordinates. The order $o$ is a numbering of the indices putting the positive ones first; we number
the indices $1,\dots,n$ in this order, so that $s_a=+1$ for $a\le n_+$ and $s_a=-1$ for $a>n_+$. $B$
is the group of invertible matrices upper triangular in this order (at every place), and for a
diagonal or upper triangular $b$ and $u\in\Z^n$, $b^u=\prod_ab_{aa}^{u_a}$. Holomorphy of weight $u$
is that of \S\ref{psp:ss-inputs}. For a non-zero $p\in\auO_E$ the \emph{principal congruence subgroup
of level $p$} is the group $K(p)$ of the $k\in\auU(X)(\auA_f)$ whose matrix in the frame is integral
and $\equiv1\bmod p$.

\subsubsection{Holomorphic functions: the formal part}\label{psq:ss-H}

\begin{lemma}[(H1), first part, and (H2), first part]\label{psq:H12}
\textup{(i)} If $\phi,\phi'$ are holomorphic of weights $u,u'$ on $\pspG$, then $\phi\phi'$ is
holomorphic of weight $u+u'$. \textup{(ii)} If $\phi$ is holomorphic of weight $u$, then
$\phi(h\,\mathrm{rot}(\theta))=e^{2\pi i\theta\langle o,u\rangle}\phi(h)$ for $h\in\pspG$, $\theta\in\R/\Z$.
\end{lemma}
\begin{proof}
(i) If $\tilde\phi,\tilde\phi'$ are holomorphic extensions on neighbourhoods $U,U'$ of $\pspG B$, then
$\tilde\phi\tilde\phi'$ is holomorphic on $U\cap U'$, a neighbourhood of $\pspG B$, and
$\tilde\phi\tilde\phi'(hb)=b^ub^{u'}\phi(h)\phi'(h)=b^{u+u'}(\phi\phi')(h)$. (ii) $\mathrm{rot}(\theta)$ is
diagonal and unitary at every place, so it lies in $\pspG\cap B$; hence $h\,\mathrm{rot}(\theta)\in\pspG$
and, reading $h\,\mathrm{rot}(\theta)$ as $h\cdot b$ with $b=\mathrm{rot}(\theta)$,
$\phi(h\,\mathrm{rot}(\theta))=\tilde\phi(h\,\mathrm{rot}(\theta))=\mathrm{rot}(\theta)^u\phi(h)
=\prod_{\tau,a}e^{2\pi i\theta o_au_{\tau,a}}\phi(h)$.
\end{proof}

\subsubsection{Peaks lie away from the cusp}\label{psq:ss-P}

This is Siegel's argument; we write it out for unitary groups, for vector-valued weights and with
peaks in place of maxima. The analytic core is Lemma~\ref{psq:analytic}.

\begin{lemma}[Transvections]\label{psq:transv}
Let $v\in\C^n$ with $v^*\varepsilon v=0$ and $\alpha,\beta\in\C$. Put $T_\alpha=1-\alpha vv^*\varepsilon$.
Then $T_\alpha T_\beta=T_{\alpha+\beta}$, and $T_\alpha^*\varepsilon T_\alpha=\varepsilon-2\Rel(\alpha)\,
\varepsilon vv^*\varepsilon$. In particular $T_\alpha\in\auU(\varepsilon)$ if $\Rel\alpha=0$.
\end{lemma}
\begin{proof}
$T_\alpha T_\beta=1-(\alpha+\beta)vv^*\varepsilon+\alpha\beta v(v^*\varepsilon v)v^*\varepsilon$, and the
last term is $0$. Likewise $T_\alpha^*\varepsilon T_\alpha=(1-\aubar\alpha\varepsilon vv^*)\varepsilon(1-
\alpha vv^*\varepsilon)=\varepsilon-(\alpha+\aubar\alpha)\varepsilon vv^*\varepsilon+|\alpha|^2\varepsilon
v(v^*\varepsilon v)v^*\varepsilon$, using $\varepsilon^2=1$.
\end{proof}

\begin{lemma}[Completing squares]\label{psq:factor}
Let $0\ne v\in\C^n$ with $v^*\varepsilon v=0$, $t\in\C$, $\tau=2\Rel t$, and put $S_k=\sum_{l\le k}s_l|v_l|^2$
and $p_k=1-\tau S_k$ for $0\le k\le n$ \textup($S_0=0$, $p_0=p_n=1$\textup). If $p_k>0$ for all $k$,
then $T_t=1-t\,vv^*\varepsilon=hb$ with $h\in\auU(\varepsilon)$, $b\in B$, and $b_{kk}=(p_k/p_{k-1})^{1/2}>0$.
\end{lemma}
\begin{proof}
By Lemma~\ref{psq:transv}, $G:=T_t^*\varepsilon T_t=\varepsilon-\tau ww^*$ with $w=\varepsilon v$, so
$w_l=s_lv_l$. Let $G_k$, $\varepsilon_k$, $w_{(k)}$ be the leading $k\times k$ blocks and the first $k$
coordinates. By the matrix determinant lemma, $\det G_k=\det\varepsilon_k\,(1-\tau w_{(k)}^*\varepsilon_k
w_{(k)})=\bigl(\prod_{l\le k}s_l\bigr)\bigl(1-\tau\sum_{l\le k}s_l|v_l|^2\bigr)=\bigl(\prod_{l\le k}s_l
\bigr)p_k$, which is non-zero. A hermitian matrix all of whose leading principal minors $\delta_k$ are
non-zero is uniquely $c^*Dc$ with $c$ upper unitriangular and $D$ real diagonal, $D_{kk}=\delta_k/
\delta_{k-1}$ ($\delta_0=1$): by induction on $n$, writing $G=\bigl(\begin{smallmatrix}G'&g\\g^*&
\gamma\end{smallmatrix}\bigr)$ with $G'=c'^*D'c'$, take the last column of $c$ to be
$(c'G'^{-1}g,1)$ (so that $c'^*D'(c'G'^{-1}g)=c'^*D'c'G'^{-1}g=g$) and $D_{nn}=\gamma-g^*G'^{-1}g=\det G/\det G'$ (Schur complement). Here
$D_{kk}=s_kp_k/p_{k-1}$. Put $d=\mathrm{diag}((p_k/p_{k-1})^{1/2})$ and $b=dc\in B$; as $d$ commutes
with $\varepsilon$, $b^*\varepsilon b=c^*d\varepsilon dc=c^*Dc=G$. Hence $h=T_tb^{-1}$ satisfies
$h^*\varepsilon h=b^{-*}Gb^{-1}=\varepsilon$, and $b_{kk}=d_{kk}$.
\end{proof}

\begin{lemma}[The weight factor]\label{psq:bound}
Let $u\in\Z^n$ have gap $\ge1$ \textup($u_a+1\le u_c$ for $a\le n_+<c$\textup) and $|u_a|\le A$,
$A\ge1$. Let $0\ne v\in\C^n$ with $v^*\varepsilon v=0$, $a_0=i\alpha$ with $\alpha>0$, and suppose
$\eta=\alpha|v|^2\le\frac12$. For $z$ in the upper half plane $\psqH$ put $t(z)=a_0(z-i)$. Then
$T_{t(z)}=h(z)b(z)$ with $h(z)\in\auU(\varepsilon)$, $b(z)\in B$ with positive diagonal, and
\[
0<b(z)^u\le e^{2A\eta}\,(1+\eta\,\Iml z)^{-1/2}.
\]
\end{lemma}
\begin{proof}
Write $y=\Iml z>0$. Then $\Rel t(z)=\Rel(i\alpha(z-i))=\alpha(1-y)$, so $\tau=2\alpha(1-y)$. Since $v$ is
isotropic, $\sum_{a\le n_+}|v_a|^2=\sum_{a>n_+}|v_a|^2=|v|^2/2$; in particular both signs occur. The
partial sums $S_k$ increase from $0$ to $|v|^2/2$ for $k\le n_+$ and decrease back to $0$ for
$k\ge n_+$, so $S_k\in[0,|v|^2/2]$ and $p_k=1-\tau S_k$ lies between $1$ and
\[
Y:=p_{n_+}=1-\alpha(1-y)|v|^2=1+\eta(y-1)\ \ge 1-\eta\ge\tfrac12 .
\]
So all $p_k>0$, and Lemma~\ref{psq:factor} gives $h(z),b(z)$ with $b(z)^u=\exp(\frac12\sum_kL_ku_k)$,
$L_k=\log p_k-\log p_{k-1}$. We have $\sum_kL_k=0$ and $\sum_{k\le n_+}L_k=\log Y$. Moreover
$p_k-p_{k-1}=-\tau s_k|v_k|^2$, so all $L_k$ with $k\le n_+$ have the sign of $y-1$ (or vanish) and all
$L_k$ with $k>n_+$ the opposite sign.

\emph{Case $y\ge1$.} Let $A_0=\min_{c>n_+}u_c$; by the gap, $u_a-A_0\le-1$ for $a\le n_+$, and
$u_c-A_0\ge0$ for $c>n_+$. As $L_k\ge0$ for $k\le n_+$ and $L_k\le0$ for $k>n_+$,
\[
\sum_kL_ku_k=\sum_kL_k(u_k-A_0)\le-\sum_{k\le n_+}L_k=-\log Y ,
\]
so $b^u\le Y^{-1/2}$. Now $(1+2\eta)Y-(1+\eta y)=\eta(1-2\eta)+2\eta^2y\ge0$ as $\eta\le\frac12$. Hence $Y^{-1/2}\le
(1+2\eta)^{1/2}(1+\eta y)^{-1/2}\le e^{\eta}(1+\eta y)^{-1/2}\le e^{2A\eta}(1+\eta y)^{-1/2}$.

\emph{Case $y<1$.} Then $\sum_kL_ku_k\le A\sum_k|L_k|=2A|\log Y|=-2A\log Y$ (the $L_k$ with $k\le n_+$
sum to $\log Y$ and have one sign, those with $k>n_+$ sum to $-\log Y$), so $b^u\le Y^{-A}\le
(1-\eta)^{-A}$, as $Y\ge1-\eta$ here. By convexity $-\log(1-\eta)\le2\log2\cdot\eta\le1.39\,\eta$ for
$0\le\eta\le\frac12$, and $(1+\eta y)^{1/2}\le(1+\eta)^{1/2}\le e^{\eta/2}$; so
$b^u\le e^{(1.39A+0.5)\eta}(1+\eta y)^{-1/2}\le e^{2A\eta}(1+\eta y)^{-1/2}$ since $A\ge1$.
\end{proof}

\begin{lemma}[Schwarz at the cusp]\label{psq:schwarz}
Let $F$ be holomorphic on $\psqH$ with $F(z+1)=F(z)$, $|F|\le C'$, and $F(z)\to0$ as $\Iml z\to\infty$
uniformly in $\Rel z$. Then $|F(i)|\le C'e^{-2\pi}$.
\end{lemma}
\begin{proof}
$F(z)=g(e^{2\pi iz})$ with $g$ holomorphic on the punctured unit disc and bounded by $C'$; by Riemann's
theorem $g$ extends holomorphically to the disc, with $g(0)=0$ by the hypothesis at $i\infty$. By the
Schwarz lemma applied to $g/C'$, $|g(q)|\le C'|q|$, and $z=i$ corresponds to $|q|=e^{-2\pi}$.
\end{proof}

\begin{lemma}[Analytic part]\label{psq:analytic}
Let $u=(u_\tau)_\tau$ be a weight with gap $\ge1$ at every $\tau$ and entries in $[-A,A]$, $A\ge2$. Let
$\phi$ be holomorphic of weight $u$ on $\pspG$ with $|\phi|\le C$. Let $t\in\pspG$ with
$t_\tau=1-a_\tau v_\tau v_\tau^*\varepsilon$, where $v_\tau\in\C^n$ is isotropic for $\varepsilon$,
$a_\tau\in i\R$, $\eta_\tau=\Iml(a_\tau)|v_\tau|^2>0$, and $A\sum_\tau\eta_\tau\le1$. If $\phi(th)=\phi(h)$
for all $h\in\pspG$, then $|\phi(1)|\le Ce^{2-2\pi}$.
\end{lemma}
\begin{proof}
For $\zeta\in\C$ let $x(\zeta)\in\prod_\tau\auGL_n(\C)$ have components $x(\zeta)_\tau=1-a_\tau(\zeta-i)
v_\tau v_\tau^*\varepsilon$ (invertible, with inverse given by Lemma~\ref{psq:transv}); it depends
holomorphically (affinely) on $\zeta$, $x(i)=1$, and $t\,x(\zeta)=x(\zeta+1)$ by Lemma~\ref{psq:transv}.
Each $\eta_\tau\le1/A\le\frac12$. For $z\in\psqH$, Lemma~\ref{psq:bound} at each $\tau$ (with
$a_0=a_\tau$, $v=v_\tau$) gives $x(z)=h(z)b(z)$ with $h(z)\in\pspG$, $b(z)\in B$ and
$0<b(z)_\tau^{u_\tau}\le e^{2A\eta_\tau}(1+\eta_\tau\Iml z)^{-1/2}$. In particular $x(z)\in\pspG B$, and
$f(z)=\tilde\phi(x(z))$ is holomorphic on $\psqH$, with
\[
f(z)=\Bigl(\prod_\tau b(z)_\tau^{u_\tau}\Bigr)\phi(h(z)),\qquad |f(z)|\le C\prod_\tau e^{2A\eta_\tau}
(1+\eta_\tau\Iml z)^{-1/2}\le Ce^{2}\prod_\tau(1+\eta_\tau\Iml z)^{-1/2}.
\]
So $|f|\le Ce^2$ and $f(z)\to0$ as $\Iml z\to\infty$, uniformly in $\Rel z$. Periodicity:
$x(z+1)=t\,h(z)\,b(z)$ with $t\,h(z)\in\pspG$, so $f(z+1)=b(z)^u\phi(t\,h(z))=b(z)^u\phi(h(z))=f(z)$.
Finally $f(i)=\tilde\phi(1)=\phi(1)$. Lemma~\ref{psq:schwarz} with $C'=Ce^2$ gives the claim.
\end{proof}

\begin{theorem}[Input (P)]\label{psq:peaks}
Let $\Phi$ be a bounded function on $\auU(X)(\auA)$, left invariant under $\auU(X)(E^+)$, right
invariant under $K(p)$ for some $p\in\auO_E$ with $\aubar p=-p$ and $0<\Iml\tau(p)\le A'$ for all
$\tau\in\auPl$, $A'\ge1$, and holomorphic of a weight $u$ with gap $\ge1$ and entries in $[-A,A]$,
$A\ge2$. Then every peak $g$ of $\Phi$ satisfies $\pspl(X,g)\ge(AA')^{-1/2}$. This is the input (P)
of Proposition~\ref{psp:inputs}.
\end{theorem}
\begin{proof}
Let $C=\sup|\Phi|$, positive since $\Phi(g)\ne0$. As $(AA')^{-1/2}\le1$ it suffices to show that no
$0\ne\gamma\in\Lambda_g$ has $|\gamma|_g<(AA')^{-1/2}$. Suppose $\gamma$ is one, and let
$Y=g^{-1}\gamma\in\auA_E^n$ be its coordinates at $g$; so $Y_f$ is integral and $\sum_\tau|Y_\tau|^2=
|\gamma|_g^2<(AA')^{-1}\le1$.

\emph{Isotropy.} By fact (R3) stated after {\cite[Definition~3.27]{P53}} with $\gamma'=\gamma$, $\langle\gamma,\gamma\rangle=0$;
hence $Y^*SY=0$ at every place ($g$ is unitary), in particular $Y_\tau^*\varepsilon Y_\tau=0$.

\emph{The transvection.} Let $T_\gamma(w)=w-p\langle w,\gamma\rangle\gamma$ on $X$. Since
$\aubar p=-p$ and $\langle\gamma,\gamma\rangle=0$, the computation of Lemma~\ref{psq:transv} (in any
basis) shows $T_\gamma\in\auU(X)(E^+)$. Put $t=g^{-1}T_\gamma g\in\auU(X)(\auA)$. As $g$ is unitary,
$t(w)=w-p\langle gw,\gamma\rangle g^{-1}\gamma=w-p\langle w,Y\rangle Y$, so the matrix of $t$ in the
frame is $1-p\,YY^*S$. At the finite places this is $1-p\,Y_fY_f^*S_f$ with $Y_f$ and $S_f$ integral,
so $t_f\in K(p)$; at $\tau$ it is $t_\tau=1-\tau(p)Y_\tau Y_\tau^*\varepsilon$, and $\tau(p)\in i\R$
because $\aubar p=-p$.

\emph{Invariance.} Let $\phi(h)=\Phi(gh)$ on $\pspG$; it is holomorphic of weight $u$ and $|\phi|\le C$.
For $h\in\pspG$, as $t_f$ commutes with $\pspG$, $g\,t_\infty h=g\,t\,h\,t_f^{-1}$, so
\[
\phi(t_\infty h)=\Phi(g\,t\,h\,t_f^{-1})=\Phi(g\,t\,h)=\Phi(T_\gamma\,g\,h)=\Phi(gh)=\phi(h),
\]
using $t_f^{-1}\in K(p)$ and $T_\gamma\in\auU(X)(E^+)$.

\emph{Sizes.} $\eta_\tau=\Iml\tau(p)|Y_\tau|^2>0$, since $\gamma\ne0$ forces $Y_\tau=g_\tau^{-1}\tau(\gamma)\ne0$;
and $A\sum_\tau\eta_\tau\le AA'\sum_\tau|Y_\tau|^2<1$.

Lemma~\ref{psq:analytic} for $\phi$, $t_\infty$, $v_\tau=Y_\tau$, $a_\tau=\tau(p)$ gives $|\Phi(g)|=
|\phi(1)|\le Ce^{2-2\pi}<C/2$, as $2\pi-2>\log2$. This contradicts $|\Phi(g)|>C/2$.
\end{proof}

\begin{remark}[The applications in \S\ref{psp:sec}]\label{psq:Papp}
In Proposition~\ref{psp:approx} (P) is applied with $p=\sqrt{-d}$ to the carrier, and in
\S\ref{psp:ss-proof} with $p=M\sqrt{-d}$ to $\aubar y$. In both cases $\aubar p=-p$, and
$\Iml\tau(p)$ has the same sign for all $\tau\in\auPl$: $\auPl$ is identified with
$\operatorname{Gal}(E/K)$, so all $\tau$ restrict to one embedding of $K=\Q(\sqrt{-d})$ and
$\tau(\sqrt{-d})=\pm i\sqrt d$ with a sign independent of $\tau$. If the sign is $-$, replace $p$ by
$-p$, which generates the same ideal and so gives the same group $K(p)$. Thus $0<\Iml\tau(p)\le A'$
holds with $A'=\sqrt d$, resp.\ $A'=M\sqrt d$, as used there. The holomorphy hypotheses are the
inputs (H4), proved in \S\S\ref{pss:sec}, \ref{psv:sec}.
\end{remark}

\subsubsection{Three checks}\label{psq:ss-checks}

\begin{lemma}[Normalisation of $\theta$]\label{psq:thetanorm}
Let $\mathcal W$ be a symplectic space over $E^+$ of the form $\auHom(X,X')$, with the inner product on
$\psHS(\mathcal W)=\mathfrak F\otimes\psHS_f(\mathcal W)$ for which the vacuum of {\cite[\S5.2.1]{P49}}
is a unit vector \textup(the inner product used in Lemma~\ref{psp:avg}\textup), and with $\theta$
normalised as in the proof of {\cite[Proposition~5.27(ii)]{P49}}, where $\int|\theta(\psrho(w)F)|^2dw=\|F\|^2$ for the self-dual measure, which gives $\mathcal W(E^+)\backslash\mathcal W(\auA)$ volume $1$; in particular $w\mapsto|\theta(\psrho(w)(1\otimes\psvac))|^2$ has mean $1$ over
$\mathcal W(E^+)\backslash\mathcal W(\auA)$ \textup(total mass $1$\textup). Then for every unitary
isomorphism $U$ of $\psHS(\mathcal W)$ onto a Schr\"odinger model $\mathcal S(\mathcal X(\auA))$
commuting with the Heisenberg operators there is $\zeta$ with $|\zeta|=1$ and $\theta=\zeta\,\theta_S\circ U$,
$\theta_S(F)=\sum_{\xi\in\mathcal X(E^+)}F(\xi)$. Consequently Lemma~\ref{psp:prodtheta} holds for the
paper's $\theta$ exactly as stated, without a constant.
\end{lemma}
\begin{proof}
$\theta_S$ is invariant under $\psrho(w)$, $w\in\mathcal W(E^+)$ (Poisson summation, as in the proof of
Lemma~\ref{psp:prodtheta}), so $\theta\circ U^{-1}$ and $\theta_S$ are both invariant functionals on the
model, and by the uniqueness in (W1) of {\cite[Lemma~5.2]{P49}}, $\theta=\zeta\theta_S\circ U$ for some
$\zeta\in\C$. For $F\in\mathcal S(\mathcal X(\auA))$ and $w=(a,b)$ with $a\in\mathcal X(\auA)$,
$b\in\mathcal Y(\auA)$, with the notation of the proof of Lemma~\ref{psp:prodtheta},
$\theta_S(\psrho(a,b)F)=\psi(\frac12\langle a,b\rangle)\sum_\xi\psi(\langle\xi,b\rangle)F(\xi+a)$. The
modulus is invariant under $\mathcal W(E^+)$ (since $\psrho(\gamma)\psrho(w)$ is a unimodular multiple
of $\psrho(\gamma+w)$ and $\theta_S\circ\psrho(\gamma)=\theta_S$). The characters $b\mapsto\psi(\langle\xi,b
\rangle)$, $\xi\in\mathcal X(E^+)$, of the compact group $\mathcal Y(E^+)\backslash\mathcal Y(\auA)$ (of mass
$1$) are distinct, so by Parseval and the absolute summability of $\xi\mapsto F(\xi+a)$,
\[
\int_{\mathcal Y(E^+)\backslash\mathcal Y(\auA)}|\theta_S(\psrho(a,b)F)|^2db=\sum_\xi|F(\xi+a)|^2,\qquad
\int_{\mathcal X(E^+)\backslash\mathcal X(\auA)}\sum_\xi|F(\xi+a)|^2da=\int_{\mathcal X(\auA)}|F|^2=\|F\|^2 .
\]
Hence the mean of $|\theta(\psrho(w)F)|^2$ over $\mathcal W(E^+)\backslash\mathcal W(\auA)$ (which is
$\mathcal X(E^+)\backslash\mathcal X(\auA)\times\mathcal Y(E^+)\backslash\mathcal Y(\auA)$ as a measure space) is
$|\zeta|^2\|F\|^2$ for every $F$. For $F=1\otimes\psvac$, a unit vector, the normalisation gives
$|\zeta|=1$; conversely the mean is then $\|F\|^2$ for every $F$, so the normalisation by the unit vacuum agrees with the one of {\cite[Proposition~5.27(ii)]{P49}}. In Lemma~\ref{psp:prodtheta} the proof computes $\theta_S(UF)\aubar{\theta_S(UF')}$; the
factor $|\zeta|^2=1$ makes this equal to $\theta(F)\aubar{\theta(F')}$, and $U$ preserves the inner
products $\langle F',\psrho(y)F\rangle$.
\end{proof}

\noindent So no constant enters
Proposition~\ref{psp:sparse}; the only requirement is that the vacuum be a unit vector for the inner
product used, which is how Lemma~\ref{psp:avg} uses it.

\begin{remark}[Level of $e_0\poA(W_j)$]\label{psq:level}
By the definitions preceding {\cite[Definition~14.6]{P46}} and the sentence after it, $\poA(W_j)$
is defined on $W_j$ ``with the parameter $\varsigma$ and the level $J_M$'', and $J_M$ is the principal
congruence subgroup of level $\pi_M$ in the stabiliser of the lattice; $e_0\poA(W_j)\subseteq
\poA(W_j)$, since $\varrho(e_0)$ is an operator on $\poA(W_j)$. So every $y\in e_0\poA(W_j)$ is right
invariant under $J_M=K(\pi_M)$, where $K(\cdot)$ is taken in the frame whose integral vectors form that
lattice, as in {\cite[Definition~3.27]{P53}}. If $\pi_M\mid p$ in $\auO_E$, then $K(p)\subseteq
K(\pi_M)$ (a matrix $\equiv1\bmod p$ is $\equiv1\bmod\pi_M$). Since $\pi_M\mid\pi_M\aubar{\pi_M}=M$, this
applies to $p=M\sqrt{-d}$ and to $p=\pi_M^2$. This is the invariance used in \S\ref{psp:ss-proof} (level $M\sqrt{-d}$) and in Lemma~\ref{psp:step2}
(level $\pi_M^2$, as $\pi_M\mid\pi_M^2$).
\end{remark}

\begin{lemma}[Degree of $\widetilde P$]\label{psq:degree}
In the notation of \S\ref{pte:ss-third} and Lemma~\ref{psp:weights}, $\widetilde P$ is homogeneous for
real scalars of degree $d_{\widetilde P}=\sum_\tau d_\tau$,
\[
d_\tau=2m\bigl(\mu^+_{1,\tau}-\Delta'\bigr)+m\bigl(\mu^+_{1,\tau}+N''+m\bigr),
\]
where $\mu^+_{1,\tau}=\#\iota(\tau)$ is the number of columns of $\mu^+_\tau$; and $d_{\widetilde P}\le
7fmb\le b^{3/2}$.
\end{lemma}
\begin{proof}
By \eqref{pte:eq-SDMP} with $\mathsf W$ of rank $n=3m+N''$, $l=m$ negative indices, and $V'=S$ ($p=2m$,
$l=m$), $P_i=\prod_\tau\det(y^{--}_\tau)^{\#\iota(\tau)+n-2l}\prod_{k\in\iota(\tau)}\det(\eta_{i,\tau k}
y^{++}_\tau)$. The negative indices of $\mathsf W$ lie in $\widetilde W$ ($W^\flat$ is positive definite),
so the first factor lives on $\auHom(S,\widetilde W)$; it is homogeneous of degree $m(\#\iota(\tau)+
n-2m)=m(\mu^+_{1,\tau}+N''+m)$, as $n-2m=N''+m$. Each $\det(\eta_{i,\tau k}y^{++}_\tau)$ is a $2m\times2m$
determinant whose entries are linear in $y^{++}_\tau$, so it is homogeneous of degree $2m$; by
\S\ref{pte:ss-third} the rows of $\eta_{i,\tau k}$ are basis vectors of $\widetilde W$ for the
$\mu^+_{1,\tau}-\Delta'$ columns of height $<r_+$, independent of $i$, and short rows of $W^\flat$ for the
$\Delta'$ columns of height $r_+$. Hence $P_i=\widetilde P\cdot Q_{r(i)}$ with $\widetilde P$ the product
of the first factor and of the $\mu^+_{1,\tau}-\Delta'$ factors with rows in $\widetilde W$, and its
degree at $\tau$ is $d_\tau$. Each factor is homogeneous under $y\mapsto sy$, $s$ real, at all places at
once, so the total degree is $\sum_\tau d_\tau$.

Bound: by the proof of Proposition~\ref{pte:third}, $\mu^+_{1,\tau}=\#\iota(\tau)\le\kappa_j(\tau)=
\kappa(\tau)+(j+\frac12)\Delta\le b+(m+1)\Delta\le2b$ (as $\kappa(\tau)\le b$ and $(m+1)\Delta\le
3\sqrt x\cdot2(3x^5)^8\le b$ by {\cite[(36)]{P46}}); and $N''+m=r+10m+1\le b$. So $d_\tau\le4mb+3mb=7mb$
and $d_{\widetilde P}\le7fmb$. Finally $f\le x^{12}$ (Lemma~\ref{psp:sizes}(iii)), $m\le2\sqrt x$, so
$7fm\le14x^{12.5}\le x^{13}\le b^{1/2}$, as $x^{26}\le b$ by {\cite[(36)]{P46}}; hence
$d_{\widetilde P}\le b^{3/2}$.
\end{proof}

\noindent This is the formula used in Lemma~\ref{psp:weights}.

\subsubsection{The vacuum is a unit vector}\label{psq:ss-vac}

\begin{lemma}[Inner product and vacuum]\label{psq:unitvac}
Let $\mathcal W=\auHom(X,X')$ and, for smooth $F,F'\in\psHS(\mathcal W)$,
\[
\langle F,F'\rangle_\theta=\int_{\mathcal W(E^+)\backslash\mathcal W(\auA)}\aubar{\theta(\psrho(w)F)}\,
\theta(\psrho(w)F')\,dw\qquad(\text{total mass }1).
\]
\textup{(a)} This is well defined, positive definite and invariant under $\psrho$ and $\omega$.
\textup{(b)} $\|1\otimes\psvac\|_\theta=1$. \textup{(c)} Fix at each place $v$ an invariant inner
product $\langle\,,\rangle_v$ on $\psHS_v$ for which the local vacuum $\psvac_v$ is a unit vector, and
write $\psvac=c\bigotimes_v\psvac_v$, $c\in\C^\times$. Then for vectors $F=c\bigotimes_vF_v$,
$F'=c\bigotimes_vF'_v$ with $F_v=F'_v=\psvac_v$ for almost all $v$,
$\langle F,F'\rangle_\theta=\prod_v\langle F_v,F'_v\rangle_v$.
\end{lemma}
\begin{proof}
(a) The integrand is $\mathcal W(E^+)$-invariant, as in the proof of Lemma~\ref{psq:thetanorm}, and by
that proof, with $\theta=\zeta\theta_S\circ U$ for an isometry $U$ onto a Schr\"odinger model (for any
fixed invariant inner product $\langle\,,\rangle_0$), polarisation gives $\langle F,F'\rangle_\theta=
|\zeta|^2\langle F,F'\rangle_0$, with $\zeta\ne0$ as $\theta\ne0$. Invariance under $\psrho$ is the
invariance of $dw$ under translation, and under $\omega$ follows from $\langle\,,\rangle_0$ being
$\omega$-invariant ($\omega$ is unitary, {\cite[Lemma~5.2]{P49}}). (b) is the normalisation of $\theta$ in Lemma~\ref{psq:thetanorm} (equivalently, of the proof of
{\cite[Proposition~5.27(ii)]{P49}}). (c) At each place the Stone--von Neumann theorem gives a unitary
isomorphism $U_v$ of $(\psHS_v,\langle\,,\rangle_v)$ onto a local Schr\"odinger model commuting with
the Heisenberg operators; $U=\bigotimes_vU_v$ (restricted tensor product with respect to $\psvac_v$ and
$U_v\psvac_v$) is then an isometry of $(\psHS,\bigotimes_v\langle\,,\rangle_v)$ onto a Schr\"odinger
model $\mathcal S(\mathcal X(\auA))$ commuting with the Heisenberg operators. By (a), applied with
$\langle\,,\rangle_0=\bigotimes_v\langle\,,\rangle_v$, we get $\langle\,,\rangle_\theta=\kappa\bigotimes_v
\langle\,,\rangle_v$ with $\kappa=|\zeta|^2>0$. Evaluating at $1\otimes\psvac$ and using (b), $1=\kappa|c|^2$, which gives (c).
\end{proof}

\noindent Thus, with $\langle\,,\rangle_\theta$ as the inner product, the vacuum is a unit vector (by
the paper's normalisation, not by an extra convention), Lemma~\ref{psq:thetanorm} applies with
$|\zeta|=1$, and the pairings of Lemma~\ref{psp:avg} factor over the places with unit local vacua,
provided the data $[v,\Psi_T]$ are read as $\psvac$ with its components at the places above $\pi_M$
(and at infinity) replaced, which is how kernel data are built from $\psvac$.

\subsubsection{Holomorphic discrete series: (H1)--(H3)}\label{psq:ss-H3}

Fix one infinite place and $G=\auU(\varepsilon)$, $\varepsilon=\mathrm{diag}(s_a)$, $N=n$, indices
numbered by $o$ (positive first), $p=n_+$, $q=n-p$, $T$ the diagonal unitary torus, $B$ as above. The
statements for $\pspG=\prod_\tau G_\tau$ follow by taking products. Let $u\in\Z^N$ be dominant with
gap $\gamma\ge N^2+N+1$.

\begin{lemma}[The parameter]\label{psq:param}
Let $\rho_{\mathrm c},\rho_{\mathrm n}$ be the half sums of {\cite[\S5.1]{P51}} for the order
in which the indices are decreasing as follows: the negative indices in the order $o$, then the
positive ones in the order $o$. Then $\lambda=u+\rho_{\mathrm c}-\rho_{\mathrm n}$ is a regular
parameter whose entries decrease in exactly this order, $\Lambda^K(\lambda)=u$, the gap of $\lambda$ is
$\ge\gamma-N\ge N^2+1$, and $\max|\lambda_a|\le\max|u_a|+N$.
\end{lemma}
\begin{proof}
For $a\le p$: $\rho_{\mathrm n,a}=-q/2$ and $\rho_{\mathrm c,a}=\frac{p+1}2-a$; for $c>p$:
$\rho_{\mathrm n,c}=p/2$ and $\rho_{\mathrm c,c}=\frac{q+1}2-(c-p)$. Within a block $u$ is
non-increasing and $\rho_{\mathrm c}$ strictly decreasing, $\rho_{\mathrm n}$ constant, so $\lambda$ is
strictly decreasing in $o$ there. For $a\le p<c$, $\lambda_c-\lambda_a\ge\gamma-\frac{p-1}2-\frac{q-1}2-
\frac N2\ge\gamma-N>0$. So the order of $\lambda$ is the stated one, the half sums computed for it are
those used, and $\Lambda^K(\lambda)=\lambda+\rho_{\mathrm n}-\rho_{\mathrm c}=u$. The entry bound is
immediate.
\end{proof}

\noindent Let $\pi$ be the discrete series of parameter $\lambda$, $v$ a unit highest weight vector
of its lowest $K$-type (highest for the positive compact roots $e_a-e_b$, $\lambda_a>\lambda_b$,
i.e.\ $a$ before $b$ in one block), and $\mathsf m_u(g)=\langle v,\pi(g)v\rangle$, the coefficient
$c_\tau$ of {\cite[\S5]{P51}} for $\lambda$ (inner products conjugate-linear in the first variable).
This is the precise meaning of $\mathsf m_u$ in \S\ref{psp:ss-notation}; the bounds
\eqref{psp:eq-ds} hold for it with the gap and entry bound of Lemma~\ref{psq:param}. Let
$\Delta_k$ be the leading principal $k\times k$ minor and
\[
\phi_u(x)=\prod_{k=1}^{N-1}\Delta_k(x)^{u_k-u_{k+1}}\cdot\det(x)^{u_N}.
\]

\begin{lemma}[$G B$ and $\phi_u$]\label{psq:GB}
\textup{(i)} $G\cap B=T$, and $GB$ is open and connected in $\auGL_N(\C)$. \textup{(ii)} $\phi_u$ is
holomorphic on $GB$ and holomorphic of weight $u$; $\phi_u(1)=1$; and $\phi_u(tx)=t^u\phi_u(x)$ for
diagonal $t$.
\end{lemma}
\begin{proof}
(i) If $g\in G\cap B$, $g^{-1}=\varepsilon g^*\varepsilon$ is upper and lower triangular, so $g$ is diagonal
unitary. The map $G\times B\to\auGL_N(\C)$, $(h,b)\mapsto hb$, has differential at $(h,b)$ with image
$h(\mathfrak g+\mathfrak b)b$; $\mathfrak g\cap\mathfrak b=\mathfrak t$, so $\dim_\R(\mathfrak g+\mathfrak b)
=N^2+N(N+1)-N=2N^2$ and the map is a submersion; its image is open, and connected since $G$ and $B$
are. (ii) For $b\in B$ the leading $k\times k$ block of $xb$ is that of $x$ times that of $b$, so
$\Delta_k(xb)=\Delta_k(x)\prod_{l\le k}b_{ll}$; similarly $\Delta_k(tx)=\prod_{l\le k}t_{ll}\Delta_k(x)$.
Hence $\phi_u(xb)=b^u\phi_u(x)$ and $\phi_u(tx)=t^u\phi_u(x)$ (the exponent of $b_{ll}$ is
$\sum_{k\ge l}(u_k-u_{k+1})+u_N=u_l$). The exponents $u_k-u_{k+1}$ are $\ge0$ for $k\ne p$ (dominance),
so $\phi_u$ is holomorphic wherever $\Delta_p\det\ne0$. For $h\in G$, $h_{++}^*h_{++}=1+h_{-+}^*h_{-+}$
(from $h^*\varepsilon h=\varepsilon$), so $\Delta_p(h)=\det h_{++}\ne0$, and $\Delta_p(hb)=\Delta_p(h)\prod_{l\le p}
b_{ll}\ne0$.
\end{proof}

\begin{proposition}[Uniqueness]\label{psq:unique}
Let $\psi$ be holomorphic of weight $u$ on $G$ with $\psi(\mathrm{rot}(\theta)h)=e^{2\pi i\theta\langle
o,u\rangle}\psi(h)$ for all $\theta,h$. Then $\psi=\psi(1)\phi_u$ on $G$.
\end{proposition}
\begin{proof}
Let $\tilde\psi$ be the extension; on the open set $GB$ it satisfies $\tilde\psi(hb)=b^u\psi(h)$, hence
$\tilde\psi(xb)=b^u\tilde\psi(x)$ and $\tilde\psi(\mathrm{rot}(\theta)x)=e^{2\pi i\theta\langle o,u\rangle}
\tilde\psi(x)$ for $x\in GB$. Near $1$ all $\Delta_k\ne0$; let $f=\tilde\psi/\phi_u$ there, so
$f(xb)=f(x)$ and $f(\mathrm{rot}(\theta)x)=f(x)$ (Lemma~\ref{psq:GB}(ii)). For $X$ strictly lower
triangular and small, $1+X\in GB$; put $F(X)=f(1+X)$, holomorphic near $0$. Since $\mathrm{rot}(\theta)$ is
diagonal, $F(\mathrm{rot}\,X\,\mathrm{rot}^{-1})=f(\mathrm{rot}(1+X)\mathrm{rot}^{-1})=f(1+X)=F(X)$, and
conjugation multiplies $X_{ab}$ ($a>b$) by $e^{2\pi i\theta(o_a-o_b)}$ with $o_a-o_b>0$. Comparing
Taylor coefficients, every non-constant monomial has a non-zero total weight, so its coefficient
vanishes: $F$ is constant near $0$. Since $(X,b)\mapsto(1+X)b$ covers a neighbourhood of $1$ (big
cell), $\tilde\psi=f(1)\phi_u$ near $1$, hence on the connected open set $GB$ by the identity theorem;
$f(1)=\psi(1)$.
\end{proof}

\begin{proposition}[(H1)]\label{psq:mhol}
$\mathsf m_u$ is holomorphic of weight $u$, $\mathsf m_u(\mathrm{rot}(\theta)h)=e^{2\pi i\theta\langle
o,u\rangle}\mathsf m_u(h)$, and $\mathsf m_u=\phi_u$ on $G$.
\end{proposition}
\begin{proof}
$v$ is $K$-finite, hence smooth, and for real $X\in\mathfrak g$ the right derivative is $R_X\mathsf m_u(h)
=\langle v,\pi(h)d\pi(X)v\rangle$; extend $d\pi$ and $R$ complex-linearly to $\mathfrak g_\C=\mathfrak
{gl}_N(\C)$. \emph{Claim: $d\pi(Y)v=u(Y)v$ for $Y\in\mathfrak b$, $u(Y)=\sum_au_aY_{aa}$.} For diagonal
$Y$ this says $v$ has $T$-weight $\Lambda^K(\lambda)=u$. For $Y=E_{ab}$, $a$ before $b$ in one block,
$e_a-e_b$ is a positive compact root and $d\pi(E_{ab})v=0$ as $v$ is a highest weight vector. For
$Y=E_{ac}$, $a\le p<c$: $d\pi(E_{ac})v$ has $T$-weight $u+e_a-e_c$; if it were non-zero, some $K$-type
$\mu$ of $\pi$ would contain this weight, so $u+e_a-e_c=\mu-\sum m_\alpha\alpha$ with $m_\alpha\ge0$ over
positive compact roots, and by (D4) of {\cite[\S4.2.3]{P53}} $\mu=u+\sum n_\beta\beta$ over positive
noncompact roots, $n_\beta\ge0$. Apply the linear form $H(x)=\sum_{c'>p}x_{c'}$: $H$ vanishes on compact
roots, is $+1$ on positive noncompact roots (these are $e_{c'}-e_{a'}$, $a'\le p<c'$, since
$\lambda_{c'}>\lambda_{a'}$ by Lemma~\ref{psq:param}), and is $-1$ on $e_a-e_c$; so $-1=\sum n_\beta\ge0$,
a contradiction.

\emph{Extension.} As $\pi(t)v=t^uv$ for $t\in T$, $\mathsf m_u(ht)=t^u\mathsf m_u(h)$; by
Lemma~\ref{psq:GB}(i), $\tilde{\mathsf m}(hb)=b^u\mathsf m_u(h)$ is well defined on $GB$, and smooth because
$(h,b)\mapsto b^u\mathsf m_u(h)$ is smooth and constant on the fibres of the surjective submersion
$(h,b)\mapsto hb$. At $x=hb$ and $Z\in\mathfrak{gl}_N(\C)$, $\tilde{\mathsf m}(xe^{sZ})=b^u\tilde{\mathsf
m}(he^{sZ'})$ with $Z'=bZb^{-1}$. Write $Z'=X+Y$, $X\in\mathfrak g$, $Y\in\mathfrak b$; to first order
$he^{sZ'}=he^{sX}e^{sY}$, and $\tilde{\mathsf m}(he^{sX}e^{sY})=e^{su(Y)}\mathsf m_u(he^{sX})$, so the
derivative is $R_X\mathsf m_u(h)+u(Y)\mathsf m_u(h)=\langle v,\pi(h)d\pi(X+Y)v\rangle$ by the claim. This
is complex-linear in $Z'$, hence in $Z$: $\tilde{\mathsf m}$ has complex-linear differential, so it is
holomorphic on $GB$. Left equivariance: $\mathsf m_u(\mathrm{rot}\,h)=\langle\pi(\mathrm{rot})^{-1}v,
\pi(h)v\rangle=\aubar{\mathrm{rot}^{-u}}\mathsf m_u(h)=e^{2\pi i\theta\langle o,u\rangle}\mathsf m_u(h)$.
Proposition~\ref{psq:unique} and $\mathsf m_u(1)=1$ give $\mathsf m_u=\phi_u$.
\end{proof}

\begin{theorem}[(H1)--(H3)]\label{psq:H123}
For dominant weights of gap $\ge N^2+N+1$ at every place: \textup{(H1)} and \textup{(H2)} of
Proposition~\ref{psp:inputs} hold, and \textup{(H3)} holds: if $\phi$ is holomorphic of weight $u$ and
$\phi\,\aubar{\mathsf m_u}$ is integrable on $\pspG$, then $\int\phi\aubar{\mathsf m_u}=\phi(1)\int|\mathsf
m_u|^2$.
\end{theorem}
\begin{proof}
(H1): Lemma~\ref{psq:H12}(i) and Proposition~\ref{psq:mhol}. (H2): the first part is
Lemma~\ref{psq:H12}(ii). Let $\psi(h)=\int_{\R/\Z}\phi(\mathrm{rot}(\theta)h)e^{-2\pi i\theta\langle o,u
\rangle}d\theta$. Then $\tilde\psi(x)=\int\tilde\phi(\mathrm{rot}(\theta)x)e^{-2\pi i\theta\langle o,u\rangle}
d\theta$ is holomorphic on $\pspG B$ (the integrand is holomorphic in $x$, continuous in $(\theta,x)$, and
$\mathrm{rot}(\theta)\pspG B=\pspG B$), $\tilde\psi(hb)=b^u\psi(h)$, and $\psi$ is left equivariant by the
substitution $\theta\mapsto\theta-\theta'$. By Proposition~\ref{psq:unique} for $\pspG=\prod_\tau G_\tau$ (the same proof applies:
conjugation by the diagonal $\mathrm{rot}(\theta)$ multiplies each coordinate $X_{\tau,ab}$, $a>b$, by
$e^{2\pi i\theta(o_a-o_b)}$, which has positive weight, $\pspG B$ is connected and open, and
$\phi_u=\prod_\tau\phi_{u_\tau}$; it does not apply place by place, since $\psi$ is equivariant only under
the diagonal $\mathrm{rot}(\theta)$) $\psi=\psi(1)\phi_u=\psi(1)\mathsf m_u$, and $\psi(1)=\int\phi(1\cdot\mathrm{rot}(\theta))
e^{-2\pi i\theta\langle o,u\rangle}d\theta=\phi(1)$ by Lemma~\ref{psq:H12}(ii). (H3): Haar measure on $\pspG$
is invariant under left translation by $\mathrm{rot}(\theta)$, and
$(\theta,h)\mapsto\phi(\mathrm{rot}(\theta)h)\aubar{\mathsf m_u(\mathrm{rot}(\theta)h)}$ is integrable on
$\R/\Z\times\pspG$ (its $h$-integral of the modulus is $\int|\phi\mathsf m_u|$ for each $\theta$). Using
$\aubar{\mathsf m_u(\mathrm{rot}(\theta)h)}=e^{-2\pi i\theta\langle o,u\rangle}\aubar{\mathsf m_u(h)}$ and
Fubini, $\int\phi\aubar{\mathsf m_u}=\int\psi\,\aubar{\mathsf m_u}=\phi(1)\int|\mathsf m_u|^2$.
\end{proof}




\providecommand{\auO}{\mathcal O}
\providecommand{\auU}{\mathrm U}
\providecommand{\auGL}{\mathrm{GL}}
\providecommand{\auA}{\mathbb A}
\providecommand{\aubar}[1]{\overline{#1}}
\providecommand{\pspG}{G_\infty}
\providecommand{\pspm}{\mathsf m}
\providecommand{\pcMat}{\mathrm M}
\providecommand{\psvac}{{\phi^{\mathrm{vac}}}}
\providecommand{\psZ}{Z}
\providecommand{\psrR}{\mathcal R}
\providecommand{\psrmc}{\mathsf{mc}}
\providecommand{\psrm}{\mathfrak m}
\providecommand{\psrF}{\mathbb F}

\subsection{The reproducing formula}\label{psr:sec}

This section proves the input (R) of Proposition~\ref{psp:inputs}, in the general form of
Theorem~\ref{psr:R}: for a holomorphic discrete series coefficient $\pspm_u$ on a unitary group
$\auU(X)$ and a positive definite space $Y$ of rank $>2\operatorname{rk}X$, convolution with $\pspm_u$
over the real points is, at each point, a fixed finite combination of theta lifts to $\auU(Y)$, the
same for all rapidly decreasing functions of a given level. The argument is linear duality
(\S\ref{psr:ss-R}) on top of three facts: an inner product formula of Rallis' type for rapidly
decreasing functions in Weil's range (\S\ref{psr:ss-rallis}), whose unramified part is a series of
translates which we show to be injective in the whole range $M'>2n$ (\S\ref{psr:ss-series}); and the
identification of the matrix coefficient of a holomorphic Fock vector with $\pspm_u$
(\S\ref{psr:ss-hol}). The applications of Theorem~\ref{psr:R} in \S\ref{psp:sec} need, in addition,
that its hypotheses hold for the data used there; this is proved in \S\ref{pss:sec} and \S\ref{psv:sec}
(the datum of $K^\circ$ is $\delta_\iota$, Proposition~\ref{psv:datum}).

\emph{Conventions.} Spaces, frames, moduli, the groups $K_S$, $J_S(\psrm)$, the representations
$\omega_{\alpha,\beta}$, data $[P,\Psi]$, kernels $\theta_{\alpha,\beta}(g,h;P,\Psi)$ and lifts
$\Theta_{\alpha,\beta}(P,\Psi)$ are those of {\cite[\S2.3]{P53}}; Haar measures are normalised as
there ($K_S$ of volume $1$, $\auU(Y)_\infty$ of mass $1$ for definite $Y$). Inner products of Hilbert
spaces and the Petersson product $\langle\varphi',\varphi\rangle=\int_{[\auU(S)]}\aubar{\varphi'}
\varphi$ are conjugate-linear in the first variable. A function $\psi$ on $[\auU(S)]$ is
\emph{rapidly decreasing} if it is continuous and $|\psi(g)|\le C_k\,\mathrm{ht}(g)^{-k}$ for every
$k$, $\mathrm{ht}$ the height of (SW) in {\cite[Lemma~2.13]{P53}}; right translates of rapidly decreasing
functions are rapidly decreasing ($\mathrm{ht}(gx)\le c(x)\mathrm{ht}(g)$), and the product of a
rapidly decreasing function with a function bounded by a power of the height is integrable on
$[\auU(S)]$ (reduction theory). For an integrable function $\kappa$ on $\auU(S)_\infty$ and a bounded
measurable $\psi$ on $\auU(S)(\auA)$ put
\[
(\kappa*\psi)(g)=\int_{\auU(S)_\infty}\kappa(x)\,\psi(gx)\,dx .
\]
If $\psi$ is continuous and bounded, so is $\kappa*\psi$ (dominated convergence), and
$\kappa*\psi$ is invariant on the left under whatever $\psi$ is, and on the right under every finite
adelic element that leaves $\psi$ invariant on the right (finite and infinite components commute).

\subsubsection{Holomorphic vectors and their coefficients}\label{psr:ss-hol}

Let $X$ be a space of signature $(p,q)$ at every infinite place, $N=p+q$, with the order $o$ of its
frame putting the positive indices first; $\pspG=\auU(X)_\infty=\prod_\tau G_\tau$, $B$ and $\mathfrak b$
(the upper triangular matrices for $o$ at every place) are as in \S\ref{psq:sec}, and for
$Y=(Y_\tau)\in\mathfrak b$ and a family of weights $u=(u_\tau)$ put $u(Y)=\sum_{\tau,a}u_{\tau,a}
Y_{\tau,aa}$.

\begin{definition}\label{psr:def-hol}
Let $\sigma$ be a unitary representation of $\pspG$ on a Hilbert space $\mathcal H$. A
\emph{holomorphic vector of weight $u$} for $\sigma$ is a smooth vector $v\in\mathcal H$ with
$d\sigma(Y)v=u(Y)v$ for all $Y\in\mathfrak b$, where $d\sigma$ is extended complex-linearly to the
complexified Lie algebra $\prod_\tau\mathfrak{gl}_N(\C)\supseteq\mathfrak b$.
\end{definition}

\noindent For the real Weil representation $\Omega$ on the Fock space of $\auHom(X,Y)\otimes\R$
(with the archimedean components of characters), the polynomial vectors are smooth, and the
complexified Lie algebra acts on them by polynomial differential operators preserving the space of
polynomials (Howe~\cite[(2.2)--(2.3)]{HoweJAMS}: in the Fock model the $\widetilde{\auU}$-finite vectors are the polynomials, and the complexified Lie algebra acts by the operators $z_iz_j$, $\partial^2/\partial z_i\partial z_j$, $\frac12(z_i\partial/\partial z_j+\partial/\partial z_j\, z_i)$; $K$-finite vectors are smooth). So for a polynomial vector
the condition of Definition~\ref{psr:def-hol} is a finite algebraic condition.

\begin{lemma}[Coefficients of holomorphic vectors]\label{psr:coef}
Let $u$ be dominant of gap $\ge N^2+N+1$ at every place, $\sigma$ a unitary representation of
$\pspG$, and $v$ a holomorphic vector of weight $u$ for $\sigma$. Then for $h\in\pspG$
\[
\langle v,\sigma(h)v\rangle=\langle v,v\rangle\,\pspm_u(h),
\]
with $\pspm_u$ the coefficient of \S\ref{psp:ss-notation}, as made precise before
Lemma~\ref{psq:GB}. In particular $h\mapsto\langle v,\sigma(h)v\rangle$ is integrable on $\pspG$.
\end{lemma}
\begin{proof}
Put $c(h)=\langle v,\sigma(h)v\rangle$. For $Y$ in the Lie algebra $\mathfrak t$ of the diagonal
unitary torus $T$, $\frac d{ds}\sigma(e^{sY})v=\sigma(e^{sY})d\sigma(Y)v=u(Y)\sigma(e^{sY})v$, so
$\sigma(e^Y)v=e^{u(Y)}v$; as $T$ is connected, $\sigma(t)v=t^uv$ for $t\in T$. The part
\emph{Extension} of the proof of Proposition~\ref{psq:mhol} uses of the vector only that it is smooth,
that $\sigma(t)v=t^uv$ on $T$, and the identity $d\sigma(Y)v=u(Y)v$ on $\mathfrak b$ (the
\emph{Claim} there); it therefore applies verbatim to $c$ and shows that $c$ is holomorphic of weight
$u$ (the \emph{Extension} step is local, so for the product group it is done at each place). Next, for
$\theta\in\R/\Z$, $\mathrm{rot}(\theta)\in T$ and
\[
c(\mathrm{rot}(\theta)h)=\langle\sigma(\mathrm{rot}(\theta))^{-1}v,\sigma(h)v\rangle=
\aubar{\mathrm{rot}(\theta)^{-u}}\,c(h)=e^{2\pi i\theta\langle o,u\rangle}c(h).
\]
By Proposition~\ref{psq:unique} for $\pspG=\prod_\tau G_\tau$ (its proof applies to the product group with the diagonal $\mathrm{rot}(\theta)$, as in the proof of Theorem~\ref{psq:H123}), $c=c(1)\phi_u=\langle v,v\rangle\phi_u$, and $\phi_u=\pspm_u$ by
Proposition~\ref{psq:mhol}. Finally $\pspm_u$ is the coefficient $c_\tau$ of {\cite[\S5]{P51}} for the
parameter of Lemma~\ref{psq:param}, whose gap is $\ge N^2+1$, at each place; it is integrable by
{\cite[Theorem~5.9]{P51}}, and so is the product over the places.
\end{proof}

\subsubsection{The series of translates in Weil's range}\label{psr:ss-series}

{\cite[Lemma~2.11(b)]{P53}},(c) bounds the series of Rallis' formula only for $M'\ge6n$ and $n=r$.
Here we need every $M'>2n$. We use the notation of {\cite[\S2.3]{P53}} before {\cite[Lemma~2.11]{P53}}:
$S$ a space of rank $n$, $M'\ge1$, $\alpha$ unramified outside the modulus $\psrm$, $G^\psrm$, $K^\psrm$
and $\psrmc=\psrmc_{M',\alpha}$ on $G^\psrm$. As in the proof of Lemma~\ref{pex:vac} (which does not
use $n=1$ at this point), $G^\psrm/K^\psrm$ is the restricted product of the $G_v/K_v$ over the places
$v\notin S(\psrm)$ of $E^+$, and $\psrmc(\gamma)=\prod_v\psrmc_v(\gamma_v)$ with
$\psrmc_v(\gamma_v)=\alpha_v(\det\gamma_v)^{-1}\langle\psvac_v,\omega_v(\gamma_v,1)\psvac_v\rangle$, equal to
$1$ on $K_v$; each $\psrmc_v$ is bi-invariant under $K_v$, and by {\cite[Lemma~2.11(a)]{P53}}
$|\psrmc_v(\gamma_v)|=[L_v:L_v\cap\gamma_vL_v]^{-M'/2}$, $L_v$ the localisation of $L_S$. Put
\[
\varepsilon_v=\sum_{\gamma\in G_v/K_v,\ \gamma\notin K_v}|\psrmc_v(\gamma)| .
\]
Recall that every residue field of $E$ has at least $5$ elements (proof of {\cite[Lemma~2.11(b)]{P53}}),
and that $E/E^+$ is unramified at the finite places, so each finite place $v$ of $E^+$ is split or
inert.

\begin{lemma}[Neighbours]\label{psr:neigh}
Let $F$ be a non-archimedean local field, $K/F$ an unramified quadratic extension with ring of
integers $\auO$, $\varpi\in F$ a uniformiser, $q=\#\auO/\varpi\auO$ and $q_0=q^{1/2}$ the order of the
residue field of $F$. Let $V$ be a non-degenerate hermitian space of dimension $n$ over $K$. For a
lattice $L$ put $L^\sharp=\{x\in V:\langle x,L\rangle\subseteq\auO\}$; $L$ is \emph{self-dual} if
$L^\sharp=L$. For self-dual $L,L'$ define $d(L,L')$ by $q^{d(L,L')}=[L:L\cap L']$, and call $L'$ a
\emph{neighbour} of $L$ if $d(L,L')=1$.
\begin{enumerate}
\item[(i)] If $L\ne L'$ are self-dual, $L$ has a neighbour $L_1$ with $d(L_1,L')\le d(L,L')-1$.
\item[(ii)] Every self-dual $L$ has at most $\mathsf N=q_0(q^n-1)/(q-1)$ neighbours.
\item[(iii)] For every self-dual $L$ and $k\ge0$, $\#\{L'\text{ self-dual}:d(L,L')\le k\}\le
(\mathsf N+1)^k$.
\end{enumerate}
\end{lemma}
\begin{proof}
Lattices are $\auO$-modules; all indices are powers of $q$. Fix a Haar measure on $V$. The product
$\mathrm{vol}(M^\sharp)\mathrm{vol}(M)$ does not depend on the lattice $M$: writing $M=aM_0$ with
$a\in\auGL(V)$, one has $M^\sharp=(a^*)^{-1}M_0^\sharp$ with $a^*$ the adjoint for the form, and
$|\det a^*|_K=|\aubar{\det a}|_K=|\det a|_K$. Hence $\mathrm{vol}(M^\sharp)=\mathrm{vol}(M)$ for every
self-dual $M$, $[M^\sharp:M'^\sharp]=[M':M]$ for
$M\subseteq M'$, and $\mathrm{vol}(H^\sharp)=\mathrm{vol}(L)^2/\mathrm{vol}(H)$ for every lattice $H$, $L$
self-dual. In particular all self-dual lattices have the same volume, so $[L:L\cap
L']=[L':L\cap L']$ and $d$ is symmetric; and an integral lattice ($M\subseteq M^\sharp$) with the volume
of a self-dual lattice is self-dual. Two self-dual lattices one of which contains the other are equal.
Since $\varpi\in F$, $\aubar\varpi=\varpi$, and the residue field automorphism induced by the conjugation
of $K/F$ is the non-trivial one ($K/F$ unramified).

(i) $L'\not\subseteq L$, as $L'\ne L$. Take $z\in L'\smallsetminus L$ and $j\ge1$ minimal with
$\varpi^jz\in L$, and put $x=\varpi^{j-1}z$: then $x\in L'\smallsetminus L$ and $\varpi x\in L$. For
$l\in L$, $\langle l,\varpi x\rangle\in\auO$, so $\langle l,x\rangle\in\varpi^{-1}\auO$, and the map
$l\mapsto\langle l,x\rangle\bmod\auO$ from $L$ to $\varpi^{-1}\auO/\auO$ (a group of order $q$) is
$\auO$-semilinear and non-zero (otherwise $x\in L^\sharp=L$), hence surjective. Its kernel
$L_x=\{l\in L:\langle l,x\rangle\in\auO\}$ has index $q$ in $L$. Put $L_1=L_x+\auO x$. It is integral:
$L_x$ is, $\langle L_x,x\rangle\subseteq\auO$, and $\langle x,x\rangle\in\auO$ as $x\in L'$. Moreover
$\varpi x\in L_x$ ($\langle\varpi x,x\rangle=\varpi\langle x,x\rangle$) and $x\notin L_x$, so
$L_1/L_x\cong\auO/\varpi$ has order $q$ and $\mathrm{vol}(L_1)=\mathrm{vol}(L)$: $L_1$ is self-dual.
$L\cap L_1\supseteq L_x$, and $L\cap L_1\ne L$ (else $L\subseteq L_1$, so $L=L_1\ni x$); as $L/L_x$ is
simple, $L\cap L_1=L_x$ and $d(L,L_1)=1$. Finally let $\Lambda=L\cap L'$. For $l\in\Lambda$,
$\langle l,x\rangle\in\auO$ since $l,x\in L'$; so $\Lambda\subseteq L_x\subseteq L_1$, and
$L_1\cap L'\supseteq\Lambda+\auO x$. Since $x\notin\Lambda$ and $\varpi x\in L\cap L'=\Lambda$,
$[\Lambda+\auO x:\Lambda]=q$, so $[L':L_1\cap L']\le[L':\Lambda]/q$, i.e.\ $d(L_1,L')\le d(L,L')-1$.

(ii) Let $L_1$ be a neighbour and $H=L\cap L_1$, of index $q$ in $L$. $L/H$ is an $\auO$-module of
length one, so $\varpi L\subseteq H$, and $H$ corresponds to a hyperplane of the $n$-dimensional
$\auO/\varpi$-space $L/\varpi L$: there are at most $(q^n-1)/(q-1)$ choices of $H$. Fix $H$. The
integral lattices containing $H$ lie in $H^\sharp$, and $[H^\sharp:H]=q^2$. If $H^\sharp/H$ were
cyclic, its only submodule of order $q$ would be $L/H$, and $L_1=L$; so $H^\sharp/H\cong(\auO/\varpi)^2$
and $\varpi H^\sharp\subseteq H$. For $x,y\in H^\sharp$, $\langle\varpi x,y\rangle\in\auO$, so
$\langle x,y\rangle\in\varpi^{-1}\auO$, and $b(\bar x,\bar y)=\varpi\langle x,y\rangle\bmod\varpi$ is a
well defined hermitian form on the plane $\mathsf V=H^\sharp/H$ over $\psrF_q=\auO/\varpi$ (relative to
$\psrF_q/\psrF_{q_0}$); it is non-degenerate, because $b(\bar x,\cdot)=0$ means $\langle x,
H^\sharp\rangle\subseteq\auO$, i.e.\ $x\in H^{\sharp\sharp}=H$. The lattice $H+\auO x$ ($x\in H^\sharp$)
is integral iff $\langle x,x\rangle\in\auO$, iff $b(\bar x,\bar x)=0$ (as $\langle x,x\rangle\in
\varpi^{-1}\auO\cap F$). So $L_1/H$ is an isotropic line of $\mathsf V$ other than the isotropic line
$e_1=L/H$. Choose $e_2\in\mathsf V$ with $b(e_1,e_2)=1$ (non-degeneracy). The lines other than $\psrF_q e_1$ are
the $\psrF_q(e_2+se_1)$, $s\in\psrF_q$, and $b(e_2+se_1,e_2+se_1)=b(e_2,e_2)+s+\aubar s$ (or with $s$
and $\aubar s$ exchanged, according to the side on which $b$ is linear). As $b(e_2,e_2)\in\psrF_{q_0}$ and
the trace $\psrF_q\to\psrF_{q_0}$ is surjective and $\psrF_{q_0}$-linear, exactly $q/q_0=q_0$ values of $s$
give an isotropic line. Hence at most $q_0$ neighbours contain $H$, and at most $\mathsf N$ in all.

(iii) Let $a_k$ be the supremum over self-dual $L$ of the left side; $a_0=1$. By (i), every $L'\ne L$
with $d(L,L')\le k$ has $d(L_1,L')\le k-1$ for some neighbour $L_1$ of $L$, so by (ii)
$a_k\le1+\mathsf Na_{k-1}$, and $a_k\le(\mathsf N+1)^k$ by induction.
\end{proof}

\begin{lemma}[Local sums]\label{psr:local}
Let $M'\ge2n+1$ and let $v\notin S(\psrm)$ be a finite place of $E^+$, $w$ a place of $E$ above it and
$q_w$ the order of its residue field. If $v$ splits, $\varepsilon_v\le3.6\,q_w^{-3/2}$ and
$\varepsilon_v\le0.28$. If $v$ is inert, $q_w\ge9$, $\varepsilon_v\le1.8\,q_w^{-1}$ and $\varepsilon_v\le
0.2$. Consequently $\sum_v\varepsilon_v<\infty$.
\end{lemma}
\begin{proof}
\emph{Split $v$.} The proof of {\cite[Lemma~2.11(b)]{P53}}, up to and including the identity
$\bigl(\sum_ks_kq^{-kM'/2}\bigr)^2=\prod_{i=0}^{n-1}(1-q^{i-M'/2})^{-2}$, does not use $M'\ge6n$; it gives
$1+\varepsilon_v\le\prod_{i=0}^{n-1}(1-q^{i-M'/2})^{-2}$, $q=q_w\ge5$. As $i-M'/2\le-\frac32-(n-1-i)$,
\[
\log(1+\varepsilon_v)\le2\sum_{j\ge0}-\log\bigl(1-q^{-3/2-j}\bigr)\le\frac{2\sum_jq^{-3/2-j}}{1-q^{-3/2}}
=\frac{2q^{-3/2}}{(1-q^{-1})(1-q^{-3/2})}\le2.75\,q^{-3/2},
\]
using $-\log(1-t)\le t/(1-t)$ and $(1-\frac15)(1-5^{-3/2})\ge0.728$. This is $\le2.75\cdot5^{-3/2}\le
0.246$, so $\varepsilon_v\le e^{0.246}-1\le0.28$, and $\varepsilon_v\le e^{0.246}\cdot2.75\,q^{-3/2}\le
3.6\,q^{-3/2}$.

\emph{Inert $v$.} Here $E_w/E^+_v$ is unramified quadratic and $q_w=q_v^2$; as $q_w\ge5$, $q_v\ge3$ and
$q_w\ge9$. The map $\gamma\mapsto\gamma L_v$ is injective from $G_v/K_v$ to the set of self-dual
lattices of $S_v$ ($L_v$ is self-dual because $L_S$ is unimodular), and $|\psrmc_v(\gamma)|=
q_w^{-d(L_v,\gamma L_v)M'/2}$ in the notation of Lemma~\ref{psr:neigh}. By Lemma~\ref{psr:neigh}(iii),
with $X=q_w^{-M'/2}\le q_w^{-n-1/2}$,
\[
\varepsilon_v\le\sum_{k\ge1}\#\{L':d(L_v,L')=k\}X^k\le\sum_{k\ge1}\bigl((\mathsf N+1)X\bigr)^k .
\]
Now $\mathsf N\le q_w^{1/2}q_w^{n-1}\cdot q_w/(q_w-1)\le\frac98q_w^{n-1/2}$, so $(\mathsf N+1)X\le
q_w^{-1}(\frac98+q_w^{1/2-n})\le\frac32q_w^{-1}\le\frac16$, and $\varepsilon_v\le\frac{3/2}{1-1/6}q_w^{-1}
=1.8\,q_w^{-1}\le0.2$.

\emph{Summation.} Above a rational prime $p$ there are at most $f=[E^+:\Q]$ places $v$ of $E^+$, and
$q_w\ge p$ if $v$ splits, $q_w\ge p^2$ if $v$ is inert. So $\sum_v\varepsilon_v\le f\sum_p(3.6\,
p^{-3/2}+1.8\,p^{-2})<\infty$.
\end{proof}

\begin{proposition}[The series for $M'>2n$]\label{psr:series}
Let $M'>2n$. \textup{(a)} $\sum_{\gamma\in G^\psrm/K^\psrm}|\psrmc(\gamma)|<\infty$. \textup{(b)} On the space
$\mathcal B$ of bounded functions on $\auU(S)(\auA)$ that are right invariant under $K^\psrm$, the
series $\psrR=\sum_{\gamma\in G^\psrm/K^\psrm}\psrmc(\gamma)\rho(\gamma)$ converges absolutely in the
supremum norm and maps $\mathcal B$ to itself; it commutes with left translations, with right
translations by $\auU(S)_\infty$ and by the components at $S(\psrm)$, and with $\kappa*$ for every
integrable $\kappa$ on $\auU(S)_\infty$; it maps continuous functions to continuous functions.
\textup{(c)} $\psrR$ is injective on $\mathcal B$.
\end{proposition}
\begin{proof}
(a) The sum is $\prod_v(1+\varepsilon_v)\le\exp\sum_v\varepsilon_v$, finite by Lemma~\ref{psr:local} (the
product formula holds for the sums of the non-negative multiplicative function $|\psrmc|$ over the
restricted product). (b) $\psrmc$ is right $K^\psrm$-invariant, so the terms are well defined on $\mathcal
B$; it is left $K^\psrm$-invariant, so (reindexing the cosets $\gamma\mapsto k\gamma$) the sum is again
right $K^\psrm$-invariant. Absolute convergence in the supremum norm follows from (a), and gives the
continuity statement (uniform limit of continuous functions). The commutation with left translations
is clear; right translations by $\auU(S)_\infty$ and by the components at $S(\psrm)$ commute with the
$\rho(\gamma)$; and $\kappa*(\psrR\psi)=\psrR(\kappa*\psi)$ by Fubini, the double series--integral
converging absolutely ($\int|\kappa|\cdot\sum|\psrmc|\cdot\sup|\psi|<\infty$).

(c) By Lemma~\ref{psr:local} choose a finite set $\mathsf T$ of places $v\notin S(\psrm)$ with
$\sum_{v\notin\mathsf T\cup S(\psrm)}\varepsilon_v\le\frac13$. For $v\in\mathsf T$ let
$\psrR_v=\sum_{\gamma\in G_v/K_v}\psrmc_v(\gamma)\rho(\gamma)$, and let $\psrR^{\mathsf T}$ be the analogous
series over the restricted product of the $G_v/K_v$, $v\notin\mathsf T\cup S(\psrm)$, with the
coefficient $\prod_{v\notin\mathsf T}\psrmc_v$. By (a) and the product structure of $\psrmc$ and of
$G^\psrm/K^\psrm$, the absolutely convergent series $\psrR$ may be summed iteratively:
$\psrR=\psrR_{v_1}\circ\dots\circ\psrR_{v_k}\circ\psrR^{\mathsf T}$ on $\mathcal B$, $\mathsf T=\{v_1,\dots,v_k\}$,
each factor mapping $\mathcal B$ to itself by the argument of (b). The coset of $1$ contributes the
identity (as $\psrmc_v=1$ on $K_v$), so
\[
\|\psrR_v\psi-\psi\|_\infty\le\varepsilon_v\|\psi\|_\infty,\qquad
\|\psrR^{\mathsf T}\psi-\psi\|_\infty\le\Bigl(\prod_{v\notin\mathsf T\cup S(\psrm)}(1+\varepsilon_v)-1\Bigr)
\|\psi\|_\infty\le(e^{1/3}-1)\|\psi\|_\infty .
\]
Since $\varepsilon_v\le0.28<1$ (Lemma~\ref{psr:local}) and $e^{1/3}-1<\frac12$, each factor satisfies
$\|\psrR_\bullet\psi\|_\infty\ge c\|\psi\|_\infty$ with $c>0$, hence is injective on $\mathcal B$, and so is
the composite.
\end{proof}

\noindent For $M'\ge6n$ and $n=r$ this is consistent with {\cite[Lemma~2.11(c)]{P53}}, which gives
the stronger bound $\|\psrR\psi\|_\infty\ge\frac12\|\psi\|_\infty$ directly.

\subsubsection{The inner product formula for rapidly decreasing functions}\label{psr:ss-rallis}

\begin{lemma}[Weil's formula for a definite target]\label{psr:SW}
The statement \textup{(SW)} of {\cite[Lemma~2.13]{P53}}, including the absolute convergence and the bounds by
a power of the height of $g\in\auU(\mathcal W)(\auA)$, holds when $T$ is positive definite at every infinite place and $M'>2n$: these imply the
hypotheses (i), (ii) of {\cite[Lemma~2.13]{P53}}.
\end{lemma}
\begin{proof}[References and comments]
This is the argument of Step~1 of the proof of Proposition~\ref{pex:rallis}, where it is given for
$n=1$, $M'=4$; it does not use $n=1$. Weil's identity~\cite[Th\'eor\`eme~5]{WeilSiegel} holds under (i) absolute
convergence of the theta integral for all $\Phi$ and (ii) absolute convergence of the Siegel Eisenstein
series of $\auU(\mathcal W)\cong\auU(n,n)$ attached to $\Phi$. (i): $T$ is anisotropic over $E^+$, being
definite at the infinite places, so $[\auU(T)]$ is compact, and $h\mapsto\theta(\omega(g,h)\Phi)$ is
continuous; $\int_{[\auU(T)]}|\theta(\omega(g,h)\Phi)|\,dh$ is bounded by a power of the height of $g$
because $|\theta(\omega(g,h)\Phi)|$ is, uniformly for $h$ in a compact set of representatives (moderate
growth of theta functions, {\cite[Lemma~3.29]{P53}}). (ii): the Eisenstein series at the point attached to $M'$ converges
absolutely for $M'>2n$~\cite[Th\'eor\`eme~1]{WeilSiegel}, and $\sum_L|\psZ_L(\omega(g,1)\Phi)|$ is bounded by a
power of the height of $g$ by Lemma~\ref{psr:Egrowth} below. The identification of the terms
with the Lagrangians is as in the comments on {\cite[Lemma~2.13]{P53}}; that part of the comments rests on~\cite{WeilSiegel} alone and comes before the bounds by powers of the height, which are deduced there from Lemma~\ref{psr:Egrowth}.

What the applications use is only this: in Step~1 of Proposition~\ref{psr:rallis} (and of
{\cite[Proposition~2.16]{P53}}), the bound of $\sum_\delta|k(g'^{-1}\delta g)|$ by a power of the heights of
$g,g'\in\auU(S)(\auA)$, which by {\cite[Lemma~2.14]{P53}} is $|c'|^{-1}\sum_L|\psZ_L(\omega((g,g'),1)\Phi)|$ for
$\Phi=F\otimes\mathsf CF$; it is a bound for the series of absolute values, so it does not follow from the
identity $I(\Phi)=E(\Phi)$ and the growth of the (finite-sum, definite-target) theta side alone, and
Lemma~\ref{psr:Egrowth} is needed. The heights of $\auU(\pcW)$ and of $\auU(S)$ compare as follows: for the frame
of $\pcW=S\oplus S^-$ that is the sum of the frames, $\Lambda_{(g,g')}=\Lambda_g\oplus\Lambda_{g'}$ with
$|(\xi,\xi')|^2=|\xi|_g^2+|\xi'|_{g'}^2$, so $\mathrm{ht}(g,g')=\max(\mathrm{ht}(g),\mathrm{ht}(g'))$; another frame
of $\pcW$ changes $\lambda_1$ by at most a constant factor (a fixed rational change of basis changes lengths
by bounded factors and integrality by a bounded denominator).
\end{proof}

\begin{lemma}[Growth of the Siegel Eisenstein series]\label{psr:Egrowth}
Let $\pcW$ be a space of rank $2n$ with a Lagrangian, $T$ a space of rank $M'>2n$, and $\Phi$ a datum of
$\auHom(\pcW,T)$. There are $C,A$ with $\sum_L|\psZ_L(\omega(x,1)\Phi)|\le C\,\mathrm{ht}(x)^A$ for all
$x\in\auU(\pcW)(\auA)$, the sum over all Lagrangians of $\pcW$.
\end{lemma}
\begin{proof}
Put $G=\auU(\pcW)$, $P$ the stabiliser of a Lagrangian $L_0$, so that the sum is $\sum_{\gamma\in P(E^+)
\backslash G(E^+)}f_\Phi(\gamma x)$ with $f_\Phi(y)=|\psZ_{L_0}(\omega(y,1)\Phi)|$ (comments on
{\cite[Lemma~2.13]{P53}}); in Weil's notation $f_\Phi$ is the function $f_\Phi(\pi(\mathbf S))=|\mathbf S\Phi(0)|$ of
\cite[no.~39]{WeilSiegel}, on the group $Ps_\auA=G(\auA)$, and $E(\Phi)$ is the series $\sum_{s\in R}$ over
representatives $R$ of $P(E^+)\backslash G(E^+)$. Weil's Lemme~18~\cite[no.~39]{WeilSiegel} says: if $f_0$ is
continuous and $>0$ on $G(\auA)$ with $f_0(py)=\psi(p)f_0(y)$ for $p\in P(\auA)$ ($\psi(p)=|\lambda|_X^{1/2}$, $\lambda$
the Levi component of $p$), then $E(\Phi)$ converges absolutely for all $\Phi$ if and only if $\sum_Rf_0(s)<\infty$;
such $f_0$ exist (p.~57). We need, for our fixed $\Phi$, a constant $c$ with $f_\Phi\le cf_0$ on $G(\auA)$, and
we prove it directly. Let $C_0$ be compact with $G(\auA)=P(\auA)C_0$ ($P(\auA)\backslash G(\auA)$ is compact). The Siegel parabolic acts on the zero term by a character: for $p\in P(\auA)$ with Levi
component $\lambda$, the Weil representation gives $(\omega(p,1)\Phi')(0)=\chi(p)\,|\lambda|_X^{1/2}\Phi'(0)$ for
every datum $\Phi'$, with $|\chi(p)|=1$ (the unipotent radical acts at $0$ by the trivial character and $\lambda$
acts by $\Phi'\mapsto\chi(\lambda)|\lambda|_X^{1/2}\Phi'(\cdot\,\lambda)$, with $|\lambda|_X$ the module of $\lambda$ on
$X_\auA$, exactly as in \cite[no.~39]{WeilSiegel}). Hence $f_\Phi(py)=\psi(p)f_\Phi(y)$, the same rule as for
$f_0$. The function $f_\Phi$ is continuous ($y\mapsto\omega(y,1)\Phi$ is continuous into the Schwartz space and
evaluation of the zero term is a continuous functional), so $f_\Phi/f_0$ is continuous and left
$P(\auA)$-invariant on $G(\auA)$, hence bounded by its maximum $c$ on $C_0$. By Th\'eor\`eme~1~\cite[no.~40]{WeilSiegel} ($M'>2n$ is condition (B) in the hermitian case,
see the comments on {\cite[Lemma~2.13]{P53}}) $E(\Phi)$ converges absolutely for all $\Phi$, so
$E_0(y):=\sum_{\gamma}f_0(\gamma y)$ is finite at $y=1$; also $E_0(\gamma y)=E_0(y)$ for rational $\gamma$.

\emph{Translation.} For
$y=pc\in G(\auA)$ and $u\in G(\auA)$ write $cu=p'c'$ with $c'\in C_0$; then $f_0(yu)=\psi(p)\psi(p')f_0(c')
\le\psi(p)\,\psi(p')\max_{C_0}f_0$ and $f_0(y)\ge\psi(p)\min_{C_0}f_0$, so
$f_0(yu)\le\kappa(u)f_0(y)$ with $\kappa(u)=\sup_{c\in C_0}\psi(p'(c,u))\cdot\max_{C_0}f_0/\min_{C_0}f_0$. Summing,
$E_0(u)\le\kappa(u)E_0(1)$, and $\sum_L|\psZ_L(\omega(u,1)\Phi)|=\sum_\gamma f_\Phi(\gamma u)\le cE_0(u)$.

\emph{Bound for $\kappa$.} $\psi(p')$ is $|\det(p'|_{L_0})|^{\kappa_0}$ (idelic absolute value) for a real
$\kappa_0$, since $|\lambda|_X$ is the module of $\lambda$ acting on $X_\auA=\auHom(L_0,T)(\auA)$-type
coordinates, a power of $|\det\lambda|$. With $p'=cuc'^{-1}$ and $p'^{-1}=c'u^{-1}c^{-1}$ we get
$|\det(p'|_{L_0})|^{\pm1}\le\prod_v\|\wedge^n p'^{\pm1}_v\|_v\le C_5(\|u\|_\auA\|u^{-1}\|_\auA)^n$, where
$\|y\|_\auA=\prod_v\|y_v\|_v$ (operator norms in the frame; almost all factors are $1$) and $c,c'$ run over a
compact set. Hence $\kappa(u)\le C_6(\|u\|_\auA\|u^{-1}\|_\auA)^{n|\kappa_0|}$.

\emph{Reduction.} Given $x$, {\cite[Lemma~3.28]{P53}} gives $\gamma$ rational with $u=\gamma x$, $u_f,u_f^{-1}$ in a
fixed compact set (so their finite norms are bounded) and $\|u_\infty\|,\|u_\infty^{-1}\|\le C\,\mathrm{ht}(x)$.
So $\sum_L|\psZ_L(\omega(x,1)\Phi)|=\sum_\gamma f_\Phi(\gamma u)\le cE_0(1)\kappa(u)\le C\,\mathrm{ht}(x)^{2n|\kappa_0|}$.
\end{proof}

\noindent The comments on Lemma~\ref{psr:SW} are complete with Lemma~\ref{psr:Egrowth} and, for the theta
integral, {\cite[Lemma~3.29]{P53}} (with $\mathrm{ht}(h)$ bounded on a compact set of representatives of
$[\auU(T)]$).

\begin{proposition}[Inner product formula]\label{psr:rallis}
Let $S$ be a space of rank $n$ (any signature), $T$ a positive definite space of rank $M'>2n$,
$\psrm$ a modulus with $S(\psrm)\ne\emptyset$, $\alpha,\beta$ characters of $\auU(1)$ unramified outside
$\psrm$ and tamely ramified at $\psrm$, $\Psi$ a function on $\pcMat_{M'\times n}(\auO_E/\psrm)$ sharp for
a group $J$ ({\cite[Definition~2.15]{P53}}), and $P$ a smooth vector of the Fock space of
$\auHom(S,T)\otimes\R$ such that
\[
\kappa(x)=\langle P,\omega_{\alpha,\beta,\infty}(x,1)P\rangle\qquad(x\in\auU(S)_\infty)
\]
is integrable. Put $\Theta=\Theta_{\alpha,\beta}(P,\Psi)$, $\psrR=\psrR_{M',\alpha}$
(Proposition~\ref{psr:series}) and
\[
A=c'\,\mathrm{vol}([\auU(T)])\,\frac{\|\Psi\|^2}{\#\pcMat_{M'\times n}(\auO_E/\psrm)\,[K_S:J]}\ne0 ,
\]
$c'$ as in {\cite[Lemma~2.14(b)]{P53}}. Then for all rapidly decreasing $\psi,\psi'$ on $[\auU(S)]$ with $\psi$
right invariant under $J$,
\[
\langle\Theta\psi',\Theta\psi\rangle_{[\auU(T)]}=A\int_{\auU(S)_\infty}\kappa(x)\,\langle\psi',
\rho(x)\psrR\psi\rangle\,dx=A\,\langle\psi',\psrR(\kappa*\psi)\rangle ,
\]
all integrals converging absolutely.
\end{proposition}
\begin{proof}
We follow the proof of {\cite[Proposition~2.16]{P53}}, with $y=\psi$, $y'=\psi'$; $F=[P,\Psi]$,
$\omega_\alpha(g)=\omega_{\alpha,\beta}(g,1)$, $k(g)=\langle F,\omega_\alpha(g)F\rangle$.

\emph{Step 1} there (doubling) uses (SW), {\cite[Lemma~2.14]{P53}} and (W7). {\cite[Lemma~2.14]{P53}} does not
use $M'\ge6n$, and its hypothesis holds: at a place of $S(\psrm)$, which splits, the component of $F$ is
$\Psi^\natural$, supported on matrices whose reduction $\bar x$ has $A_0\bar x\in\bar J$ invertible for
the matrix $A_0$ of the definition of sharpness, hence on injective matrices. By
Lemma~\ref{psr:SW} in place of (SW), {\cite[(7)]{P53}} holds, with
$\sum_\delta|k(g'^{-1}\delta g)|$ bounded by a power of the heights of $g,g'$.

\emph{Step 2} there (unfolding) uses only that $y,y'$ are rapidly decreasing, which is our hypothesis:
the function $(g,g')\mapsto|\psi'(g')||\psi(g)|\sum_\delta|k(g'^{-1}\delta g)|$ is integrable on
$[\auU(S)]^2$, and
\[
\langle\Theta\psi',\Theta\psi\rangle=c'\,\mathrm{vol}([\auU(T)])\int_{\auU(S)(\auA)}k(g)\,\langle\psi',
\rho(g)\psi\rangle\,dg ,
\]
the integral converging absolutely.

\emph{Step 3.} As in Step~3 there, $k(g)=\kappa(g_\infty)\,\psrmc(g^\psrm)\,k_\psrm(g_\psrm)$ with
$k_\psrm=\|\Psi^\natural\|^2\,\mathbf 1_J$ on the components at $S(\psrm)$ (this uses only the sharpness
of $\Psi$ and {\cite[Lemma~2.9]{P53}}) and $\|\Psi^\natural\|^2=\|\Psi\|^2/\#\pcMat_{M'\times n}(\auO_E/\psrm)$.
Since $|\langle\psi',\rho(g)\psi\rangle|\le\|\psi'\|_{L^1}\|\psi\|_\infty$, and $\kappa$, $\psrmc$ (summed
over $G^\psrm/K^\psrm$, Proposition~\ref{psr:series}(a); $K^\psrm$ has volume $1$) and $k_\psrm$ are
integrable, Fubini's theorem allows us to integrate first over the finite adelic points. As $\psi$
is right invariant under $J\supseteq K^\psrm$, $\int_{G^\psrm}\psrmc(g)\rho(g)\psi\,dg=\psrR\psi$, and the
integral over the components at $S(\psrm)$ multiplies by $\|\Psi^\natural\|^2\mathrm{vol}(J_\psrm)=
\|\Psi^\natural\|^2/[K_S:J]$. This gives the first formula; the second is
$\int\kappa(x)\langle\psi',\rho(x)\psrR\psi\rangle dx=\langle\psi',\kappa*\psrR\psi\rangle$ (Fubini, as
$\kappa\in L^1$, $\psrR\psi$ bounded, $\psi'\in L^1$) and $\kappa*\psrR\psi=\psrR(\kappa*\psi)$
(Proposition~\ref{psr:series}(b)). Finally $A\ne0$: $c'\ne0$, the volume is positive and $\Psi\ne0$.
\end{proof}

\begin{corollary}[Vanishing lifts]\label{psr:vanish}
In the situation of Proposition~\ref{psr:rallis}, if $\psi$ is rapidly
decreasing, right invariant under $J$, and $\Theta\psi=0$, then $\kappa*\psi=0$.
\end{corollary}
\begin{proof}
Put $\psi_1=\kappa*\psi$ and $z=\psrR\psi_1$. As $\psi$ is continuous and bounded, left invariant under
$\auU(S)(E^+)$ and right invariant under $J\supseteq K^\psrm$, $\psi_1$ is continuous, bounded, left
$\auU(S)(E^+)$-invariant and right $J$-invariant, and so is $z$ (Proposition~\ref{psr:series}(b); $J$
lies in $K_S$, and $\psrR$ commutes with the components at $S(\psrm)$ and preserves $K^\psrm$-invariance).
By Proposition~\ref{psr:rallis} with $\Theta\psi=0$, $\langle\psi',z\rangle=0$ for every rapidly
decreasing $\psi'$. For $\gamma\in G^\psrm$ and $x\in\auU(S)_\infty$ the translate
$\psi_{\gamma x}=\psi(\cdot\,\gamma x)$ is rapidly decreasing, so $\langle z,\psi_{\gamma x}\rangle=
\aubar{\langle\psi_{\gamma x},z\rangle}=0$, and by Fubini ($z$ bounded, $\psi\in L^1([\auU(S)])$, $\kappa\in
L^1$)
\[
\langle z,\rho(\gamma)\psi_1\rangle=\int_{\auU(S)_\infty}\kappa(x)\,\langle z,\psi_{\gamma x}\rangle\,dx=0 ,
\]
using $\psi(g\gamma x)=\psi(gx\gamma)$. Hence, $[\auU(S)]$ having finite volume and the series converging
absolutely and uniformly, $\langle z,z\rangle=\sum_\gamma\psrmc(\gamma)\langle z,\rho(\gamma)\psi_1\rangle=0$.
So $z=0$ almost everywhere, hence everywhere by continuity, and $\psi_1=0$ by
Proposition~\ref{psr:series}(c).
\end{proof}

\subsubsection{The reproducing formula}\label{psr:ss-R}

\begin{theorem}[Input (R)]\label{psr:R}
Let $X$ be a space of rank $n_X$ and signature $(p,q)$ at every infinite place, with the order $o$
putting the positive indices first; $Y$ a positive definite space of rank $M'>2n_X$; $\psrm$ a modulus
with $S(\psrm)\ne\emptyset$; $\alpha,\beta$ characters of $\auU(1)$ unramified outside $\psrm$ and tamely
ramified at $\psrm$; $\Omega=\omega_{\alpha,\beta,\infty}(\cdot,1)$ and $\Omega'=\omega_{\alpha,\beta,\infty}(1,
\cdot)$ the actions of $\pspG=\auU(X)_\infty$ and $\auU(Y)_\infty$ on the Fock space of $\auHom(X,Y)
\otimes\R$. Let $u$ be a dominant weight on $X$ of gap $\ge n_X^2+n_X+1$ at every place. Let $\mathcal Q$
be a finite-dimensional space of smooth Fock vectors \textup{(}for instance polynomials\textup{)}, stable
under $\Omega'(y)$ for all $y\in\auU(Y)_\infty$, that contains a non-zero holomorphic vector $v_0$ of
weight $u$ for $\Omega$ \textup{(}Definition~\ref{psr:def-hol}\textup{)}, and let $Q_1,\dots,Q_s$ be a
basis of $\mathcal Q$. Let $\iota\in\pcMat_{M'\times n_X}(\auO_E/\psrm)$ be the matrix of the inclusion of
frames and $h_1,\dots,h_{P_Y}\in\auU(Y)(\auA_f)$ representatives of the finite set
$\auU(Y)(E^+)\backslash\auU(Y)(\auA_f)/J_Y(\psrm)$. Then for every $g\in\auU(X)(\auA)$ there are
constants $a_{p,i}(g)$ such that for every rapidly decreasing $\psi$ on $[\auU(X)]$ right invariant
under $J_X(\psrm)$
\[
\int_{\pspG}\psi(gh)\,\pspm_u(h)\,dh=\sum_{p=1}^{P_Y}\sum_{i=1}^sa_{p,i}(g)\int_{[\auU(X)]}\psi(x)\,
\theta_{\alpha,\beta}\bigl(x,h_p;Q_i,\delta_\iota\bigr)\,dx ,
\]
all integrals converging absolutely.
\end{theorem}

The finiteness of the set of classes is {\cite[Proposition~7.20]{P49}} for a standard space $Y=E^{M'}$
(take $\mathsf G\in\psrm\cap\Z$, so that $J_Y(\psrm)\supseteq K(\mathsf G)$), and Borel's finiteness
theorem~\cite[Theorem~5.1]{BorelFin} in general. This is (R) of Proposition~\ref{psp:inputs},
with $\psrm=(\pi_M^2)$, the classes $h_p$ of $\auU(Y)$ at level $\pi_M^2$, the polynomials
$Q_{(p,i)}=Q_i$ and the constants $a_{(p,i)}$, under the hypotheses on $\mathcal Q$, $\alpha$, $\beta$ stated
here; the gap $\ge300\pspn$ of Proposition~\ref{psp:inputs} exceeds $t^2+t+1\ge n_X^2+n_X+1$.

\begin{proof}
\emph{Convergence.} $\pspm_u$ is integrable (Lemma~\ref{psr:coef}) and $\psi$ bounded; the kernel has
moderate growth in $x$ ({\cite[Lemma~3.29]{P53}}) and $\psi$ decreases rapidly.

\emph{Linear forms.} Let $\mathcal V$ be the complex vector space of rapidly decreasing functions on
$[\auU(X)]$ right invariant under $J=J_X(\psrm)$. Fix $g$ and define linear forms on $\mathcal V$:
\[
\ell(\psi)=\int_{\pspG}\psi(gh)\pspm_u(h)\,dh,\qquad L_{p,i}(\psi)=\bigl(\Theta_{\alpha,\beta}(Q_i,
\delta_\iota)\psi\bigr)(h_p) .
\]
If $\bigcap_{p,i}\ker L_{p,i}\subseteq\ker\ell$, then $\ell$ is a linear combination of the finitely many
$L_{p,i}$ (the map $\psi\mapsto(L_{p,i}(\psi))_{p,i}$ factors $\ell$ through its image in $\C^{P_Ys}$, and a
linear form on a subspace of $\C^{P_Ys}$ extends to $\C^{P_Ys}$), which is the assertion. So let
$\psi\in\mathcal V$ with $L_{p,i}(\psi)=0$ for all $p,i$; we show $\ell(\psi)=0$.

\emph{Step 1: the lift by $v_0$ vanishes.} For a smooth Fock vector $v$ write $F_v=[v,\delta_\iota]$ and
$\Lambda(v)(h)=\int_{[\auU(X)]}\psi(x)\theta(\omega_{\alpha,\beta}(x,h)F_v)\,dx$, linear in $v$. Let
$h\in\auU(Y)(\auA)$ and write $h_f=\gamma h_pk$ with $\gamma\in\auU(Y)(E^+)$, $k\in J_Y(\psrm)$, and put
$y=\gamma^{-1}h_\infty\in\auU(Y)_\infty$ ($\gamma$ taken at the infinite places); then $h=\gamma\cdot
h_p\,k\,y$, the finite and infinite components commuting. The
element $(1,k)$ fixes $F_v$: by {\cite[Lemma~2.9]{P53}}, $\omega_{\alpha,\beta}(1,k)[v,\delta_\iota]=[v,\Psi']$
with $\Psi'(x)=\beta_\psrm(\det\bar k)\delta_\iota(\bar k^{-1}x)=\delta_\iota(x)$, as $\bar k=1$. The theta
functional is invariant under the rational points and $\beta$ is trivial on them ((W1), (W3) of
{\cite[Lemma~5.2]{P49}}), so
\[
\Lambda(v)(h)=\int\psi(x)\theta\bigl(\omega_{\alpha,\beta}(x,h_p)\omega_{\alpha,\beta}(1,k)
\omega_{\alpha,\beta}(1,y)F_v\bigr)dx=\Lambda(\Omega'(y)v)(h_p),
\]
since $\omega_{\alpha,\beta}(1,y)F_v=F_{\Omega'(y)v}$. By hypothesis $\Omega'(y)v_0=\sum_ic_iQ_i$ for some
$c_i\in\C$, and $\Lambda(Q_i)(h_p)=L_{p,i}(\psi)=0$. Hence $\Lambda(v_0)=0$, i.e.\
$\Theta_{\alpha,\beta}(v_0,\delta_\iota)\psi=0$ on $\auU(Y)(\auA)$.

\emph{Step 2: convolution.} The point datum $\delta_\iota$ is sharp for $J_X(\psrm)$
({\cite[Definition~2.15]{P53}}): the reduction of $J_X(\psrm)$ is trivial, so the first condition is empty,
and the matrix $A_0={}^t\iota$ satisfies $A_0\iota=1$. By Lemma~\ref{psr:coef},
$\kappa(h)=\langle v_0,\Omega(h)v_0\rangle=\langle v_0,v_0\rangle\pspm_u(h)$ is integrable. By
Corollary~\ref{psr:vanish} (with $S=X$, $T=Y$, $P=v_0$, $\Psi=\delta_\iota$, $J=J_X(\psrm)$),
$\kappa*\psi=0$; at $g$ this reads $\langle v_0,v_0\rangle\,\ell(\psi)=0$, and $\langle v_0,v_0\rangle\ne0$.
\end{proof}

\begin{remark}[On the hypotheses of Theorem~\ref{psr:R}]\label{psr:rem}
(i) Theorem~\ref{psr:R} is (R) of Proposition~\ref{psp:inputs} for $Y$ of rank $>2n_X$, the gap
$\ge300\pspn$, and characters of the kernel unramified outside $\psrm=(\pi_M^2)$ and tamely ramified at
$\psrm$ (in the applications they are unramified at every finite place, Theorem~\ref{pss:a}); it makes the ``Moreover'' clause of (R) precise: the $Q_p$ are a basis of a
finite-dimensional space of Fock vectors stable under $\auU(Y)_\infty$ and containing a non-zero
holomorphic vector of weight $u$ in the sense of Definition~\ref{psr:def-hol}, and the index runs over
pairs (class, basis vector). (ii) The proof does not use the gap beyond $n_X^2+n_X+1$ (needed to identify
the coefficient with $\pspm_u$, Lemma~\ref{psr:coef}), and it uses of the level only that $\delta_\iota$
is sharp for $J_X(\psrm)$ and invariant under $J_Y(\psrm)$. (iii) The inner product formula is used only
through Corollary~\ref{psr:vanish}; the constant $A$ is not computed.
\end{remark}




\providecommand{\auU}{\mathrm U}
\providecommand{\auGL}{\mathrm{GL}}
\providecommand{\auHom}{\operatorname{Hom}}
\providecommand{\auA}{\mathbb A}
\providecommand{\aubar}[1]{\overline{#1}}
\providecommand{\pspG}{G_\infty}
\providecommand{\pspm}{\mathsf m}
\providecommand{\pcMat}{\mathrm M}
\providecommand{\pssg}{\mathfrak g}
\providecommand{\pop}{\mathsf p}

\subsection{Holomorphic polynomial Fock vectors}\label{pss:sec}

This section supplies the tool needed for (H4), (H5) and the hypotheses of Theorem~\ref{psr:R}: a
criterion for a polynomial Fock vector to be holomorphic in the sense of
Definition~\ref{psr:def-hol}, with its weight (\S\ref{pss:ss-tool}). With it we prove the input
(H5) of Proposition~\ref{psp:inputs} (\S\ref{pss:ss-H5}), the hypothesis on the polynomials $Q_r$ in
the application of Theorem~\ref{psr:R} to $(S,W^\flat)$ (\S\ref{pss:ss-b}), the part of (H4) concerning
the carrier (\S\ref{pss:ss-carrier}), and the hypothesis on the characters in both applications of
Theorem~\ref{psr:R} (\S\ref{pss:ss-a}). Proposition~\ref{pss:status} collects these inputs.

\subsubsection{The tool}\label{pss:ss-tool}

\emph{Setting.} Fix an infinite place $\tau$ (statements for $\prod_\tau$ follow by taking products,
as in \S\ref{psq:ss-H3}). Let $X,T$ be hermitian spaces over $\C$ with orthonormal frames of signs
$s\in\{\pm1\}^{N}$, $s'\in\{\pm1\}^{n'}$, the indices of $X$ numbered by its order $o$ (positive
first, $p$ positive, $q=N-p$ negative); $n_T^\pm$, $n_X^\pm$ count the signs. We use the conventions of
\S\ref{ptf:ss-fock}: a vector of $\auHom(X,T)$ is a matrix $w\in\pcMat_{n'\times N}(\C)$, $(g,h)$ acts
by $w\mapsto hwg^{-1}$, the entry $(r,c)$ has product sign $\varepsilon_{rc}=s'_rs_c$ and Fock variable
$z_{rc}$ ($=\aubar{w_{rc}}$ if $\varepsilon_{rc}=1$, $w_{rc}$ otherwise); a polynomial $P$ in the $z_{rc}$
is identified with its Fock vector. Let $\omega_\infty$ be the archimedean Weil representation with
characters of weights $a$ on $\auU(X)$ and $b$ on $\auU(T)$ (Lemma~\ref{ptf:fock}),
$\Omega(g)=\omega_\infty(g,1)$, $\Omega'(h)=\omega_\infty(1,h)$; $G=\auU(X)$, $\pssg$ its Lie algebra,
$\pssg_\C=\mathfrak{gl}_N(\C)$, $\mathfrak b$ (resp.\ $\mathfrak b^-$) the upper (resp.\ lower)
triangular matrices, $\mathfrak k_\C$ the block diagonal ones, $T_X\subset G$ the diagonal torus, and
similarly for $\auU(T)$ with its order (positive first). For $u\in\Z^N$ and $Z\in\pssg_\C$ put
$u(Z)=\sum_cu_cZ_{cc}$.

We use three standard facts. (KV) The polynomial vectors are smooth vectors for the action of the
semidirect product of the Heisenberg group with $\auU(X)\times\auU(T)$, and $d\Omega$, $d\Omega'$ and
$d\varrho$ map polynomial vectors to polynomial vectors (Howe~\cite[(2.2)--(2.3)]{HoweJAMS}, for the symplectic part: the polynomials are the $K$-finite, hence smooth, vectors, and the derived action is by the operators of (2.2); for the Heisenberg factor, a polynomial vector corresponds in the Schr\"odinger model to a polynomial times a Gaussian, a Schwartz function, and Schwartz functions are smooth vectors of the Schr\"odinger representation: the Heisenberg group acts on $L^2$ by translations, multiplications by characters $e^{2\pi i\langle y,x\rangle}$ and central scalars, so for a Schwartz $f$ the difference quotients of the orbit map at the identity converge in $L^2$, by dominated convergence, to constant multiples of $\partial_jf$, $x_jf$ and $f$, which are again Schwartz; by induction and equivariance the orbit map is $C^\infty$; see also Folland~\cite{FollandPS}). (Hs) On polynomial vectors the derived Heisenberg action is
{\cite[(32)]{P53}} (or its variant in {\cite[Remark~4.23]{P53}}); in either normalisation, since
$d\varrho(w)-i\,d\varrho(iw)$ and $d\varrho(w)+i\,d\varrho(iw)$ are the complex-linear and antilinear parts,
the complex span of the operators $d\varrho(w)$ is the span $V_{\mathrm{lin}}$ of the operators $z_i$
(multiplication) and $\partial_i=\partial/\partial z_i$, $i=(r,c)$. (Cv) $\Omega(g)\varrho(w)\Omega(g)^{-1}=
\varrho(\iota(g)w)$ with $\iota(g)w=wg^{-1}$ ((W1$_\infty$)); likewise for $\Omega'$.

\begin{lemma}[Quadratic form of the derived action]\label{pss:weyl}
For $Z\in\pssg_\C$ there are an element $D_0$ of the span $\mathcal Q_2$ of the operators $z_iz_j$,
$z_i\partial_j$, $\partial_i\partial_j$ and a constant $c$ such that $d\Omega(Z)=D_0+c$ on polynomials.
The same holds for $d\Omega'$.
\end{lemma}
\begin{proof}
By complex linearity we may take $Z\in\pssg$. Differentiating (Cv) along $e^{sZ}$ on polynomial vectors
(smooth by (KV)) gives $[d\Omega(Z),d\varrho(w)]=d\varrho(\iota(Z)w)$, so by (Hs) $D=d\Omega(Z)$ satisfies
$[D,V_{\mathrm{lin}}]\subseteq V_{\mathrm{lin}}$, and $A_Z=\mathrm{ad}(D)|_{V_{\mathrm{lin}}}$ preserves the
alternating form $\langle\ell,m\rangle=[\ell,m]\in\C$ infinitesimally ($[D,[\ell,m]]=0$ by Jacobi), which is
non-degenerate ($[\partial_i,z_j]=\delta_{ij}$): $A_Z\in\mathfrak{sp}(V_{\mathrm{lin}})$. Let $Q$ be the span of
the symmetrised products $\ell m+m\ell$, $\ell,m\in V_{\mathrm{lin}}$; $Q\subseteq\mathcal Q_2+\C$. We have $\dim Q=k(2k+1)$
($k=\dim V_{\mathrm{lin}}/2$): $Q$ is spanned by the $k(2k+1)$ elements $z_iz_j$ ($i\le j$), $\partial_i\partial_j$
($i\le j$), $\frac12(z_i\partial_j+\partial_jz_i)=z_i\partial_j+\frac12\delta_{ij}$ (all $i,j$), and these are
linearly independent, since a combination $\sum_{i\le j}a_{ij}z_iz_j+\sum_{i\le j}b_{ij}\partial_i\partial_j+
\sum_{i,j}c_{ij}(z_i\partial_j+\frac12\delta_{ij})$ that vanishes on polynomials has, applied to $1$, constant
term $\frac12\sum_ic_{ii}$ and $z_iz_j$-coefficients $a_{ij}$, so $a=0$; applied to $z_l$ it gives
$\sum_ic_{il}z_i+\frac12(\sum_ic_{ii})z_l=0$, whence $c_{il}=0$ for $i\ne l$ and then $c_{ll}+\frac12\sum_ic_{ii}=0$
for all $l$, which forces all $c_{ll}$ equal to some $\gamma$ with $\gamma(1+k/2)=0$ (as $\sum_ic_{ii}=k\gamma$), so $c=0$; applied to
$z_lz_{l'}$ it gives $b_{ll'}(1+\delta_{ll'})=0$, so $b=0$. Moreover $\mathrm{ad}(q)$ maps $V_{\mathrm{lin}}$ into itself
($[\ell m,x]=\ell[m,x]+[\ell,x]m$) and lies in $\mathfrak{sp}(V_{\mathrm{lin}})$. If $\mathrm{ad}(q)=0$, then $q$
commutes with all $z_i$ and $\partial_i$, so $q(P)=Pq(1)$ and $\partial_iq(1)=q(\partial_i1)=0$: $q$ is a scalar;
writing $q=\sum a_{ij}z_iz_j+\sum b_{ij}\partial_i\partial_j+\sum c_{ij}z_i\partial_j+\frac12\sum c_{ii}$ and
applying it to $1$, $z_l$, $z_lz_{l'}$ gives $a=0$, $c=0$, $b=0$, so $q=0$. Hence $\mathrm{ad}\colon Q\to
\mathfrak{sp}(V_{\mathrm{lin}})$ is injective and, by dimensions ($\dim\mathfrak{sp}_{2k}=k(2k+1)$), bijective.
Choose $q$ with $\mathrm{ad}(q)=A_Z$; then $D-q$ commutes with all $z_i$, $\partial_i$ on polynomials,
hence is a scalar (same argument). Take $D_0=q$ minus its constant part.
\end{proof}

\begin{lemma}[Torus weights]\label{pss:torus}
\textup{(i)} For $t\in T_X$, $\Omega(t)P=c(t)\,P(\lambda(t)z)$ with $\lambda_{rc}(t)=t_c^{-\varepsilon_{rc}}$ and a
character $c$; hence $\Omega(t)z_{rc}\Omega(t)^{-1}=t_c^{-\varepsilon_{rc}}z_{rc}$ and
$\Omega(t)\partial_{rc}\Omega(t)^{-1}=t_c^{\varepsilon_{rc}}\partial_{rc}$. We say $z_{rc}$ has \emph{$X$-weight}
$-\varepsilon_{rc}e_c$ and $\partial_{rc}$ has $X$-weight $\varepsilon_{rc}e_c$.
\textup{(ii)} For $a\ne b$, $d\Omega(E_{ab})$ is a linear combination of the products $xy$ of two
operators $x,y\in\{z_i,\partial_i\}$ whose $X$-weights add up to $e_a-e_b$.
\textup{(iii)} The same holds for $\Omega'$ and $t\in T_T$ with \emph{$T$-weights} $\varepsilon_{rc}e_r$ for
$z_{rc}$ and $-\varepsilon_{rc}e_r$ for $\partial_{rc}$.
\end{lemma}
\begin{proof}
(i) is Lemma~\ref{ptf:fock} with $g=t$, $h=1$: variables of positive product sign go to $z\,t^{-\mathsf t}$, i.e.\
$z_{rc}\mapsto t_c^{-1}z_{rc}$, those of negative sign to $zt$; the conjugation formulas follow by the
chain rule. (ii) $\Omega(t)d\Omega(Z)\Omega(t)^{-1}=d\Omega(\mathrm{Ad}(t)Z)$ and $\mathrm{Ad}(t)E_{ab}=t_at_b^{-1}E_{ab}$.
Write $d\Omega(E_{ab})=D_0+c$ as in Lemma~\ref{pss:weyl} in the basis $1$, $z_iz_j$ ($i\le j$), $z_i\partial_j$,
$\partial_i\partial_j$ ($i\le j$) (linearly independent, by the argument at the end of the proof of
Lemma~\ref{pss:weyl}); each basis element is an eigenvector of conjugation by $\Omega(t)$ with the
character given by its weight, by (i). Characters with distinct exponents are distinct, so only basis
elements of weight $e_a-e_b$ occur; the constant has weight $0$. (iii) Lemma~\ref{ptf:fock} with $g=1$,
$h=t$: positive product sign $z\mapsto t^{\mathsf t}z$, negative $z\mapsto t^{-1}z$.
\end{proof}

\begin{lemma}[The compact part]\label{pss:cpt}
For $Y\in\mathfrak k_\C$ and every polynomial $P$,
\[
d\Omega(Y)P=-\sum_{\varepsilon_{rc}=1}\sum_{c'}Y_{cc'}z_{rc'}\partial_{rc}P+\sum_{\varepsilon_{rc}=-1}\sum_{c'}
Y_{c'c}z_{rc'}\partial_{rc}P+\Bigl(\sum_{s_c=1}(n_T^--a)Y_{cc}+\sum_{s_c=-1}(n_T^+-a)Y_{cc}\Bigr)P,
\]
and for $Y\in\mathfrak k'_\C$ (block diagonal for $T$)
\[
d\Omega'(Y)P=\sum_{\varepsilon_{rc}=1}\sum_{r'}Y_{r'r}z_{r'c}\partial_{rc}P-\sum_{\varepsilon_{rc}=-1}\sum_{r'}
Y_{rr'}z_{r'c}\partial_{rc}P+\Bigl(\sum_{s'_r=1}(b-n_X^-)Y_{rr}+\sum_{s'_r=-1}(b-n_X^+)Y_{rr}\Bigr)P .
\]
\end{lemma}
\begin{proof}
For real $Y$ differentiate the formula of Lemma~\ref{ptf:fock} at $g=e^{sY}$ (resp.\ $h=e^{sY}$):
$(z e^{-sY^{\mathsf t}})_{rc}$, $(ze^{sY})_{rc}$, $(e^{sY^{\mathsf t}}z)_{rc}$, $(e^{-sY}z)_{rc}$ have derivatives
$-\sum_{c'}Y_{cc'}z_{rc'}$, $\sum_{c'}Y_{c'c}z_{rc'}$, $\sum_{r'}Y_{r'r}z_{r'c}$, $-\sum_{r'}Y_{rr'}z_{r'c}$, and the
derivatives of the characters $c(g,1)$, $c(1,h)$ of Lemma~\ref{ptf:fock} are the traces shown. Both sides
are complex-linear in $Y$.
\end{proof}

\begin{lemma}[Holomorphy criterion]\label{pss:ext}
Let $\phi$ be a smooth function on $G$ and $u\in\Z^N$. For real $X\in\pssg$ let $R_X\phi(h)=\frac d{ds}\phi(he^{sX})|_{s=0}$,
extended complex-linearly to $\pssg_\C=\pssg\oplus i\pssg$. If $R_Y\phi=u(Y)\phi$ for all $Y\in\mathfrak b$, then
$\phi$ is holomorphic of weight $u$ \textup(\S\ref{psp:ss-inputs}\textup).
\end{lemma}
\begin{proof}
This is the part \emph{Extension} of the proof of Proposition~\ref{psq:mhol}, which uses of the
coefficient only the identity $R_Y\phi=u(Y)\phi$ on $\mathfrak b$. In detail: for $Y\in i\R^N$ diagonal,
$Y\in\pssg\cap\mathfrak b$ and $\phi(he^{Y})=e^{u(Y)}\phi(h)$, so $\phi(ht)=t^u\phi(h)$ on $T_X$; since $G\cap B=T_X$
(Lemma~\ref{psq:GB}), $\tilde\phi(hb)=b^u\phi(h)$ is well defined on the open set $GB$; it is smooth because $(h,b)\mapsto b^u\phi(h)$ is smooth on $G\times B$ and
constant on the fibres of the surjective submersion $G\times B\to GB$, $(h,b)\mapsto hb$ (Lemma~\ref{psq:GB}(i)), which has
smooth local sections. At $x=hb$,
for $Z\in\pssg_\C$ put $Z'=bZb^{-1}=X+Y$ with $X\in\pssg$, $Y\in\mathfrak b$; then
$\frac d{ds}\tilde\phi(xe^{sZ})|_0=b^u(R_X\phi(h)+u(Y)\phi(h))=b^uR_{Z'}\phi(h)$, complex-linear in $Z$.
So $\tilde\phi$ is holomorphic, and $\tilde\phi(hb)=b^u\phi(h)$.
\end{proof}

\begin{corollary}[Coefficients]\label{pss:coefs}
Let $\sigma$ be a unitary representation of $G$ (or of $\pspG$), $\mathcal H^\infty$ its Fr\'echet space of
smooth vectors.
\textup{(a)} If $v$ is holomorphic of weight $u$ and $\Lambda$ is a continuous linear form on
$\mathcal H^\infty$, then $h\mapsto\Lambda(\sigma(h)v)$ is holomorphic of weight $u$. In particular so is
$h\mapsto\langle w,\sigma(h)v\rangle$ for every $w\in\mathcal H$. More generally, the same holds if
$\mathcal E\subset\mathcal H$ is a Fr\'echet space, continuously embedded, such that $h\mapsto\sigma(h)v$ is smooth
into $\mathcal E$ with derivative $h\mapsto\sigma(h)d\sigma(X)v$, and $\Lambda$ is a continuous linear form on $\mathcal E$.
\textup{(b)} If $v\ne0$ is smooth with $d\sigma(Z)v=\lambda(Z)v$ for all $Z\in\mathfrak b^-$, $\lambda\in\Z^N$, then
$h\mapsto\langle\sigma(h)v,v\rangle$ is holomorphic of weight $-\lambda$; if moreover $-\lambda$ is dominant of gap
$\ge N^2+N+1$, then $\langle v,\sigma(h)v\rangle=\langle v,v\rangle\,\aubar{\pspm_{-\lambda}(h)}$.
\end{corollary}
\begin{proof}
$h\mapsto\sigma(h)v\in\mathcal H^\infty$ is smooth (G\aa rding), so $\phi(h)=\Lambda(\sigma(h)v)$ is smooth with
$R_X\phi(h)=\Lambda(\sigma(h)d\sigma(X)v)$, and by complex linearity $R_Y\phi=u(Y)\phi$ on $\mathfrak b$;
Lemma~\ref{pss:ext} gives (a); in the general form the same computation applies verbatim with $\mathcal E$
in place of $\mathcal H^\infty$. (b) Let $g(h)=\langle\sigma(h)v,v\rangle$. For $Y=Y_1+iY_2$, $Y_k\in\pssg$,
$R_Yg=\langle\sigma(h)d\sigma(Y_1-iY_2)v,v\rangle$ (conjugate-linearity in the first variable), and $Y_1-iY_2=
Y^\natural:=-\varepsilon Y^*\varepsilon$, the conjugation of $\pssg_\C$ fixing $\pssg$. For $Y\in\mathfrak b$,
$Y^\natural\in\mathfrak b^-$, so $d\sigma(Y^\natural)v=\lambda(Y^\natural)v=-\sum_c\lambda_c\aubar{Y_{cc}}v$ and
$R_Yg=-\sum_c\lambda_cY_{cc}\,g$. By Lemma~\ref{pss:ext} $g$ is holomorphic of weight $-\lambda$. As $\sigma(t)v=t^\lambda v$
on $T_X$ (diagonal elements, connectedness), $g(\mathrm{rot}(\theta)h)=\langle\sigma(h)v,\sigma(\mathrm{rot}(\theta))^{-1}v\rangle
=e^{2\pi i\theta\langle o,-\lambda\rangle}g(h)$. By Propositions~\ref{psq:unique} and~\ref{psq:mhol} (at each place),
$g=g(1)\pspm_{-\lambda}$, and $\langle v,\sigma(h)v\rangle=\aubar{g(h)}$.
\end{proof}

\begin{proposition}[Holomorphic polynomial vectors]\label{pss:hol}
Let $P$ be a polynomial and $d\in\Z_{\ge0}^N$.
\textup{(i)} Let $T$ be positive definite, $P$ involve only the variables $z_{rc}$ with $s_c=1$, and
satisfy $\sum_rz_{rc'}\partial_{rc}P=0$ for positive $c\ne c'$ with $c$ before $c'$, and $\sum_rz_{rc}\partial_{rc}P=d_cP$.
Then $P$ is holomorphic for $\Omega$ of weight $u$, $u_c=-d_c-a$ for $s_c=1$ and $u_c=n'-a$ for $s_c=-1$.
\textup{(ii)} Let $T$ be negative definite, $P$ involve only the $z_{rc}$ with $s_c=-1$, with
$\sum_rz_{rc'}\partial_{rc}P=0$ for negative $c\ne c'$ with $c$ after $c'$, and $\sum_rz_{rc}\partial_{rc}P=d_cP$. Then
$d\Omega(Z)P=\lambda(Z)P$ for $Z\in\mathfrak b^-$, $\lambda_c=n'-a$ for $s_c=1$, $\lambda_c=-d_c-a$ for $s_c=-1$.
\textup{(iii)} Let $X$ be negative definite of rank $N$, $P$ involve only the $z_{rc}$ with $s'_r=-1$, with
$\sum_cz_{rc}\partial_{r'c}P=0$ for negative rows $r\ne r'$ with $r$ before $r'$, and $\sum_cz_{rc}\partial_{rc}P=d_rP$.
Then $P$ is holomorphic for $\Omega'$ \textup(for the order of $T$\textup) of weight $u'$, $u'_r=b-N$ for
$s'_r=1$, $u'_r=d_r+b$ for $s'_r=-1$.
\end{proposition}
\begin{proof}
$\mathfrak b$ is spanned by the diagonal, the $E_{cc'}$ ($c$ before $c'$, same sign) and the $E_{ac}$
($s_a=1$, $s_c=-1$); $\mathfrak b^-$ likewise with the transposes. (i) Here $\varepsilon_{rc}=s_c$. By
Lemma~\ref{pss:cpt}, for $c,c'$ positive, $d\Omega(E_{cc'})P=-\sum_rz_{rc'}\partial_{rc}P$ (the second sum
involves $\partial$ of variables absent from $P$), which is $0$ for $c$ before $c'$ and $(-d_c-a)P$ after
adding the trace term for $c=c'$ ($n_T^-=0$); for $c,c'$ negative only the trace term survives:
$(n'-a)\delta_{cc'}P$. For $E_{ac}$: by Lemma~\ref{pss:torus}(ii) it is a combination of products of two
operators with $X$-weights adding to $e_a-e_c$; the $X$-weights are $\mp e_a$ for $z_{ra},\partial_{ra}$ and
$\pm e_c$ for $z_{rc},\partial_{rc}$, $\pm e_{c''}$ for other columns; $e_a-e_c$ arises only as $\partial_{ra}\partial_{r'c}$,
and $\partial_{r'c}P=0$. (ii) Here $\varepsilon_{rc}=-s_c$; the variables of $P$ have $\varepsilon=1$ and the
same computation gives the compact part ($n_T^+=0$, $n_T^-=n'$). For $E_{ca}$ ($s_c=-1$, $s_a=1$): $z_{rc}$,
$\partial_{rc}$ have weights $-e_c,e_c$ and $z_{ra},\partial_{ra}$ have $e_a,-e_a$; $e_c-e_a$ arises only as
$\partial_{rc}\partial_{r'a}$, and $\partial_{r'a}P=0$. (iii) Here $\varepsilon_{rc}=-s'_r$; by Lemma~\ref{pss:cpt},
for negative rows $d\Omega'(E_{rr'})P=\sum_cz_{rc}\partial_{r'c}P$, zero for $r$ before $r'$ and with trace term
$b-n_X^+=b$ on the diagonal; on positive rows only the trace $b-n_X^-=b-N$. For $E_{ac}$ ($s'_a=1$,
$s'_c=-1$), Lemma~\ref{pss:torus}(iii): $z_{a\cdot},\partial_{a\cdot}$ have $T$-weights $-e_a,e_a$ and
$z_{c\cdot},\partial_{c\cdot}$ have $e_c,-e_c$; $e_a-e_c$ arises only as $\partial_{a\cdot}\partial_{c\cdot}$, and
$\partial_{a\cdot}P=0$.
\end{proof}

\begin{corollary}[Products of minors]\label{pss:minors}
\textup{(i)} With $T$ positive definite, for matrices
$\eta_k\in\pcMat_{p\times n'}(\C)$, $1\le k\le d$, the polynomial $\prod_k\det(\eta_kz^+)$ \textup($z^+$ the
$n'\times p$ block of positive columns\textup) satisfies (i) with $d_c=d$; it is holomorphic of weight
$(-d-a;\,n'-a)$. \textup{(ii)} With $T$ negative definite of rank $q$, $\det(z^-)^e$ satisfies (ii) with
$d_c=e$. \textup{(iii)} With $X$ negative definite of rank $N$ and $T$ with $N$ negative rows,
$\det(z^{-\cdot})^e$ \textup(the block of negative rows\textup) satisfies (iii) with $d_r=e$.
\end{corollary}
\begin{proof}
$\sum_rz_{rc'}\partial_{rc}$ is a derivation and on $\det(\eta_kz^+)$ replaces the column $c$ of $\eta_kz^+$ by the
column $c'$, giving $0$ for $c\ne c'$ and the determinant for $c=c'$; so it acts on the product by $0$, resp.\
by $d$. (ii), (iii) are the same with columns, resp.\ rows.
\end{proof}

\subsubsection{The factorisation (H5)}\label{pss:ss-H5}

\begin{theorem}[(H5)]\label{pss:H5}
\textup{(H5)} of Proposition~\ref{psp:inputs} holds: with $\Omega$ the real Weil representation of
$\auU(S)_\infty$ on the Fock space of $\auHom(S,\widetilde W)\otimes\R$ with the character
$\nu^\flat\upsilon^m$ of \eqref{pte:eq-F} on $\auU(S)$ \textup(and its restriction $\upsilon^m$ on
$\auU(\widetilde W)$\textup), $c_{\widetilde P}(h',1)=c_{\widetilde P}(1,1)\,\pspm_{u_2}(h')\,\aubar{\pspm_{u_3}(h')}$ for
$u_2=(\Delta';\mu^+_1+r+2m)$, $u_3=(0;\mu^+_1+N''+2m)$.
\end{theorem}
\begin{proof}
Fix $\tau$; write $\mu_1=\mu^+_{1,\tau}$. By the proof of Proposition~\ref{pte:third} the weight of
$\nu^\flat\upsilon^m$ at $\tau$ is $a=m-\mu_1$. $\widetilde W$ has $r+2m$ positive and $m$ negative indices
(\S\ref{pte:ss-third}). By Lemma~\ref{psq:degree} and its proof,
$\widetilde P=P^+P^-$ with $P^+=\prod_{k}\det(\eta_kz^{++})$, a product of $d^+=\mu_1-\Delta'$ minors of the block
$z^{++}$ (positive rows, positive columns), and $P^-=\det(z^{--})^e$, $e=\mu_1+N''+m$.

\emph{Total weight.} By Lemma~\ref{ptf:fock} with $g=t\in T_S$, $h=1$: $\Omega(t)\widetilde P=t^{\mathrm{wt}}\widetilde P$ with
$\mathrm{wt}_c=-d^++(m-a)=\Delta'$ for $c$ positive and $\mathrm{wt}_c=-e+(r+2m-a)=r-N''$ for $c$ negative. Hence
$c(h't)=t^{\mathrm{wt}}c(h')$ for $c(h')=c_{\widetilde P}(h',1)$, and $\mathrm{wt}=u_2-u_3$.

\emph{Splitting.} Let $\widetilde W_\tau=\widetilde W^+\oplus\widetilde W^-$ be the orthogonal decomposition by signs.
The natural map $F_1\otimes F_2\mapsto F_1(w_1)F_2(w_2)$ identifies the Fock space $\mathfrak F$ of
$\auHom(S,\widetilde W_\tau)$ with $\mathfrak F^+\otimes\mathfrak F^-$ unitarily, intertwining the Heisenberg
representation of the sum with the tensor product (the cocycle $\beta$ of \S\ref{psp:ss-fock} and the
Gaussian are multiplicative over orthogonal sums), and carries $\widetilde P$ to $P^+\otimes P^-$. Let
$\Omega^\pm$ be the Weil representations of $\auU(S_\tau)$ on $\mathfrak F^\pm$ of Lemma~\ref{ptf:fock} with
trivial characters. $\Omega(h)$ and $\Omega^+(h)\otimes\Omega^-(h)$ satisfy (Cv) for the same linear map,
so by irreducibility (Stone--von Neumann, Schur) $\Omega=\chi\cdot(\Omega^+\otimes\Omega^-)$ with
$\chi(h)\in\C$, $|\chi|=1$; $\chi$ is a continuous character of the connected group $\auU(S_\tau)$, so its
differential vanishes on the semisimple $\mathfrak{su}(2m,m)$ and $\chi=\det^k$, with $k\in\Z$ by restriction
to $\mathrm{diag}(z,1,\dots,1)$. Hence $c=\chi\,c^+c^-$, $c^\pm(h)=\langle P^\pm,\Omega^\pm(h)P^\pm\rangle$.

\emph{The factors.} By Corollary~\ref{pss:minors}(i) ($T=\widetilde W^+$, $n'=r+2m$, $a=0$) $P^+$ is holomorphic
of weight $w^+=(-d^+;r+2m)$, dominant of gap $r+2m+d^+\ge9m^2+3m+1$ ($r\ge x^5$, $m\le2\sqrt x$); by
Lemma~\ref{psr:coef}, $c^+=\langle P^+,P^+\rangle\pspm_{w^+}$. By Corollary~\ref{pss:minors}(ii)
($T=\widetilde W^-$, $n'=m$, $a=0$) and Proposition~\ref{pss:hol}(ii), $P^-$ satisfies the hypothesis of
Corollary~\ref{pss:coefs}(b) with $\lambda=(m;-e)$, and $-\lambda=(-m;e)$ is dominant of gap $e+m$; so
$c^-=\langle P^-,P^-\rangle\aubar{\pspm_{(-m;e)}}$. Since $\phi_{u+k\mathbf 1}=\det^k\phi_u$ (definition of $\phi_u$
before Lemma~\ref{psq:GB}) and $\pspm=\phi$ (Proposition~\ref{psq:mhol}), with $|\det h|=1$,
$\aubar{\pspm_{(-m;e)}}=\det^m\aubar{\pspm_{u_3}}$ ($u_3=(0;e+m)$) and $\chi\pspm_{w^+}\det^m=\pspm_{w'}$,
$w'=w^++(k+m)\mathbf 1$. Thus $c=c(1)\,\pspm_{w'}\,\aubar{\pspm_{u_3}}$ ($c(1)=\|P^+\|^2\|P^-\|^2$, as all
three functions are $1$ at $1$). Comparing right $T_S$-weights, $t^{\mathrm{wt}}=t^{w'-u_3}$, so $w'=u_3+
\mathrm{wt}=u_2$. Taking the product over $\tau$ gives the claim.
\end{proof}

\noindent In particular $u_3=(0;3m)+u_1+u_2$ (checked in Lemma~\ref{psp:weights}'s notation:
$N''+2m=r+11m+1$), as used in Proposition~\ref{psp:approx}.

\subsubsection{The hypothesis (b) for $(S,W^\flat)$}\label{pss:ss-b}

In Lemma~\ref{psp:step1} the kernel of $(S,W')$ with the characters $(\nu^\flat\upsilon^m,\upsilon^m)$ of
\eqref{pte:eq-F} is written by (W7) as the product of kernels of $(S,\widetilde W)$ and $(S,W^\flat)$. The scalar
$(\nu^\flat\upsilon^m)(\det g)^{-1}$ may be attached to either factor; we attach it to $(S,\widetilde W)$, so that the
kernel of $(S,W^\flat)$ has the characters $(1,\upsilon^m)$, and the kernel of $(S,\widetilde W)$ those used in
Theorem~\ref{pss:H5}.

\begin{theorem}[Hypothesis (b)]\label{pss:b}
Let $I$ be a set of $2m$ indices of $W^\flat$, $v_0=\prod_\tau\det(z^{I,+}_\tau)^{\Delta'}$ on $\auHom(S,W^\flat)\otimes\R$,
and $\mathcal Q$ the span of the $\Omega'(y)v_0$, $y\in\auU(W^\flat)_\infty$. Then $\mathcal Q$ is finite
dimensional, stable under $\auU(W^\flat)_\infty$, and $v_0\in\mathcal Q$ is non-zero and holomorphic of weight
$u_1=(-\Delta';6m+1)$ for $\Omega$ \textup(characters $(1,\upsilon^m)$\textup). Moreover for every invertible
$\mathsf Y\in\pcMat_{6m+1}(E)$, $\mathcal Q$ is contained in the span $\mathcal Q_{\mathsf Y}$ of the polynomials
$\prod_\tau\prod_{k=1}^{\Delta'}\det\bigl((\tau(\mathsf Y)z^+_\tau)_{I_{\tau k}}\bigr)$, $I_{\tau k}$ sets of $2m$ rows.
\end{theorem}
\begin{proof}
Holomorphy and weight: Corollary~\ref{pss:minors}(i) at each $\tau$ ($n'=6m+1$, $a=0$, $d=\Delta'$, $\eta_k$ the
selection of the rows $I$). By Lemma~\ref{ptf:fock}, $\Omega'(y)$ acts on polynomials in $z^+$ by a scalar
times $z^+\mapsto y^{\mathsf t}z^+$, preserving the finite-dimensional space of polynomials in $z^+$ of degree
$2m\Delta'$ at each $\tau$; so $\mathcal Q$ is finite dimensional and stable. Row $i$ of $y_\tau^{\mathsf t}z^+_\tau$ is
$\sum_jM_{ij}(\tau(\mathsf Y)z^+_\tau)_j$ for some $M$ ($\tau(\mathsf Y)$ invertible), and by multilinearity of the determinant in
its rows $\det((y^{\mathsf t}_\tau z^+_\tau)_I)$ is a combination of minors $\det((\tau(\mathsf Y)z^+_\tau)_J)$; taking powers and
the product over $\tau$, $\Omega'(y)v_0\in\mathcal Q_{\mathsf Y}$.
\end{proof}

\begin{corollary}\label{pss:bcor}
Theorem~\ref{psr:R} applies to $(X,Y)=(S,W^\flat)$, $u=u_1$, the characters $(1,\upsilon^m)$ and $\mathcal Q$
\textup(the hypothesis on the characters is Theorem~\ref{pss:a}\textup). Expanding each basis vector of $\mathcal Q$,
at the class $h_p$, in the products of minors of $\mathcal Q_{\mathsf Y_p}$, $\mathsf Y_p$ the matrix of the basis of short vectors
at $h_p$ fixed in \S\ref{pte:ss-third}, the formula takes the form used in Lemma~\ref{psp:step1}, with
polynomials $Q_r$ indexed by pairs \textup(class, row selections\textup).
\end{corollary}
\begin{proof}
The kernel is linear in the polynomial; the gap of $u_1$ is $\ge300n_*$ (Lemma~\ref{psp:weights}).
\end{proof}

\subsubsection{The carrier}\label{pss:ss-carrier}

\begin{proposition}[Part of (H4)]\label{pss:carrier}
The carrier $\varphi^\circ$ is holomorphic of weight $(0;3m)$.
\end{proposition}
\begin{proof}
Fix $x\in\auU(S)(\auA)$. By the definition in \S\ref{pte:ss-carrier},
$\varphi^\circ(xh)=\Lambda_x(\Omega'(h)v)$ for $h\in\auU(S)_\infty$, with $v=\prod_\tau\det(y^-_\tau)^{2m}$ on
$\auHom(I_m^-,S)\otimes\R$ and $\Lambda_x(w)=\int_{[\auU(I_m^-)]}\theta(\omega(y,x)[w\otimes\mathbf 1])\,dy$. We apply Corollary~\ref{pss:coefs}(a) in its general form with
$\mathcal E=\mathcal S$, the Schwartz space of a Schr\"odinger model (as in the proof of Lemma~\ref{psp:prodtheta}).
The map $h\mapsto\Omega'(h)v$ is $C^\infty$ into $\mathcal S$ (with its Fr\'echet topology), with derivative
$\Omega'(h)d\Omega'(X)v$: $\auU(S)_\infty$ acts through a smooth homomorphism into the metaplectic group, the
metaplectic operators act continuously on $\mathcal S$, locally uniformly in the group element, and the derived
action of the symplectic Lie algebra is by differential operators of order $\le2$ with polynomial coefficients of
degree $\le2$, which map $\mathcal S$ continuously to itself (Folland~\cite[Ch.~4, Prop.~4.27, Thm.~4.45 and (4.44)]{FollandPS}; local uniformity: with $\mathsf N$ the number operator (harmonic oscillator), the seminorms $\|(\mathsf N+1)^k\cdot\|_2$ define the topology of $\mathcal S$ and dominate every monomial of degree $\le2k$ in the $x_j$, $\partial_j$ \cite[Appendix to \S V.3, Lemmas~1--2]{ReedSimonI}, and $\Omega'(h)^{-1}(\mathsf N+1)\Omega'(h)$ is a quadratic polynomial in the $x_j,\partial_j$ with coefficients continuous in $h$, by \cite[(4.23)]{FollandPS}); hence for
$w\in\mathcal S$, $\frac d{ds}\Omega'(he^{sX})w=\Omega'(he^{sX})d\Omega'(X)w$ in $\mathcal S$, and induction on the
order of the derivative gives smoothness. The theta functional
$F\mapsto\sum_\xi F(\xi)$ is continuous on $\mathcal S$ (the sum is bounded by a Schwartz seminorm), and
$\omega(y,x)$ acts continuously on $\mathcal S$, uniformly for $(y,x)$ in compact sets; as $[\auU(I_m^-)]$ is compact,
$\Lambda_x$ is continuous on $\mathcal S$. By Corollary~\ref{pss:minors}(iii) ($X=I_m^-$, $N=m$, $e=2m$,
$b=m$ the weight of $\upsilon^m$, as $\upsilon_\tau|_{\auU(1)}(z)=z$) and Proposition~\ref{pss:hol}(iii), $v$ is
holomorphic of weight $(m-m;2m+m)=(0;3m)$; Corollary~\ref{pss:coefs}(a) concludes.
\end{proof}

\subsubsection{The characters (a)}\label{pss:ss-a}

\begin{theorem}[Hypothesis (a)]\label{pss:a}
The characters of the kernels used in the two applications of Theorem~\ref{psr:R} are unramified at every
finite place; in particular unramified outside $(\pi_M^2)$ and tamely ramified there.
\end{theorem}
\begin{proof}
For $(S,W^\flat)$ they are $1$ and $\upsilon^m$ (\S\ref{pss:ss-b}); $\upsilon$ is unramified at every finite place
(proof of {\cite[Lemma~7.12]{P49}}). For $(W,V^\circ)$, $K^\circ$ is defined from $\omega^\circ$ (\S\ref{po:ss-glue}),
whose character on $\auU(W)$ is $\alpha^\circ=\upsilon^{m/2-j}$, an integral power of $\upsilon$ (Lemma~\ref{psv:char}),
hence unramified at every finite place.
\end{proof}

\subsubsection{The inputs of Proposition~\ref{psp:inputs}}\label{pss:ss-open}

\begin{proposition}[Inputs of Proposition~\ref{psp:inputs}]\label{pss:status}
\textup{(P)}, \textup{(H1)}--\textup{(H3)} (\S\ref{psq:sec}), \textup{(H5)} (Theorem~\ref{pss:H5}), the carrier part of
\textup{(H4)} (Proposition~\ref{pss:carrier}), and \textup{(R)} in its application to $(S,W^\flat)$
(Theorem~\ref{psr:R} with Theorems~\ref{pss:a}, \ref{pss:b}) are proved, modulo the published results listed
in \S\ref{psr:sec}, \S\ref{psq:sec} and here: Howe~\cite[(2.2)--(2.3)]{HoweJAMS} (smoothness of polynomial vectors and
polynomiality of the derived action), the Stone--von Neumann theorem, G\aa rding's theorem on smooth
vectors. The proofs here do not use the vacuum coefficient of {\cite[\S4.2.2]{P53}} ({\cite[Lemma~4.24]{P53}}).
\end{proposition}

\noindent The rest of (H4), concerning $\aubar y$ and the combinations of the $F_i$, with the bounds on $u,u^\circ$,
is Theorem~\ref{psv:H4}; hypothesis (c) for $(W,V^\circ)$ is Theorem~\ref{psv:c}; the datum of $K^\circ$ is the
point datum of Theorem~\ref{psr:R} (Proposition~\ref{psv:datum}). All inputs of Proposition~\ref{psp:inputs} are
thereby proved (Theorem~\ref{psv:main}).



\providecommand{\auU}{\mathrm U}
\providecommand{\auGL}{\mathrm{GL}}
\providecommand{\auHom}{\operatorname{Hom}}
\providecommand{\auA}{\mathbb A}
\providecommand{\aubar}[1]{\overline{#1}}
\providecommand{\poA}{\mathcal A}
\providecommand{\pspG}{G_\infty}
\providecommand{\pspm}{\mathsf m}
\providecommand{\pspn}{n_*}
\providecommand{\pcMat}{\mathrm M}
\providecommand{\pop}{\mathsf p}
\providecommand{\pssg}{\mathfrak g}
\providecommand{\psvmu}{\mu}
\providecommand{\pcm}{\mathfrak m}
\providecommand{\auO}{\mathcal O}

\subsection{The last inputs of products span}\label{psv:sec}

This section proves the inputs of Proposition~\ref{psp:inputs} not covered by
Proposition~\ref{pss:status}. We fix the order of the frame of
$W=W_j$ (\S\ref{psv:ss-order}), prove a fourth case of Proposition~\ref{pss:hol}
(\S\ref{psv:ss-hol}), and with it the hypothesis (c) of Theorem~\ref{psr:R} for $(W,V^\circ)$
(\S\ref{psv:ss-c}); we show that the datum of $K^\circ$ is the point datum of Theorem~\ref{psr:R}
(\S\ref{psv:ss-datum}); we prove (H4) for $\aubar y$ and for the combinations of the $F_i$, with the
bounds on $u$ and $u^\circ$ (\S\ref{psv:ss-H4}); we settle the character split of
Lemma~\ref{psp:step1} (\S\ref{psv:ss-split}); and we conclude (\S\ref{psv:ss-main}).

Notation is that of \S\ref{psp:sec}, \S\ref{pss:ss-tool} and \S\ref{po:ss-glue}. In particular
$n=r+t$ with $t=5r$, $r=k+m-1$ ($k$ even, $m$ even, so $r$ is odd), $r_+=k-1+j$, $r_-=m-j$,
$N''=r+9m+1$, $n^\circ=t-9m-1$, and $\Delta=2r^8$. At $\tau\in\auPl$ the parameter $\varsigma_\tau$ of
$\poA(W_j)$ and the diagrams $\mu^\pm_\tau$ are those of \S\ref{plf:ss-L4} and
Lemma~\ref{plf:match}. Inner products of Hilbert spaces are conjugate-linear in the first
variable.

\subsubsection{The order of the frame of $W$}\label{psv:ss-order}

The definition of holomorphy in \S\ref{psp:ss-inputs} depends on the order $o$ of the frame, which
\S\ref{psp:ss-notation} only requires to put the positive indices first. For $W=W_j$ we fix it as
follows. Number the frame of $W$ by $0,\dots,r-1$ (the \emph{fixed numbering} of
{\cite[\S14.3]{P46}}: $0,\dots,r_+-1$ positive, $r_+,\dots,r-1$ negative); recall that $\varsigma_\tau(a)$
sits on the coordinate $a$ and is strictly decreasing in $a$ (\S\ref{plf:ss-L4}).

\begin{definition}\label{psv:order}
The order $o$ of the frame of $W$ is: first the positive indices in the order $r_+-1,r_+-2,\dots,0$,
then the negative indices in the order $r-1,r-2,\dots,r_+$. Equivalently, positive indices first,
and within each sign block the order of increasing $\varsigma_\tau$ (the same at every $\tau$).
\end{definition}

\noindent Thus, within each sign block, ``$a$ before $b$ in $o$'' means $\varsigma_\tau(a)<\varsigma_\tau(b)$.
The position of a negative index $c$ in $o$ among the negative ones, counted from $0$, is the number
$p(c)$ of Lemma~\ref{plf:match} (negative rows numbered by increasing $\varsigma_\tau$), and the
position of a positive index $a$ in the fixed numbering is the number $p(a)$ there (positive rows
numbered by decreasing $\varsigma_\tau$). The ``reverse numbering of the negative basis vectors'' of
\eqref{po:eq-poly} is the order $o$ on the negative block, and the ``fixed numbering of the positive
ones'' is the reverse of $o$ on the positive block. All weights on $W$ below are written in the fixed
numbering of the indices, and dominance and gaps are taken for the order $o$. Nothing in
\S\ref{psp:sec}, \S\ref{psq:sec}, \S\ref{psr:sec}, \S\ref{pss:sec} depends on the choice of $o$ among
the orders that put the positive indices first, other than through the weights, which we compute
for this $o$.

\begin{lemma}[Two row lengths]\label{psv:rows}
At $\tau\in\auPl$, for $0\le p<r_+$ the row number $p$ \textup(from $0$, longest first\textup) of
$\mu^+_\tau$ has length $\ell^+_p=\varsigma_\tau(p)+\kappa_j(\tau)-\frac{t-1}2-(r_+-p)$; for $0\le p<r_-$ the
row number $p$ of $\mu^-_\tau$ has length $\ell^-_p=-\kappa_j(\tau)-\frac{t-1}2-\varsigma_\tau(r-1-p)-(r_--p)$,
and $\ell^-_p\ge\frac\Delta2-\frac{t+1}2$. Moreover $\ell^-_p\le m\Delta\le b$ and $\max_a|\varsigma_\tau(a)|\le2b$.
\end{lemma}
\begin{proof}
The two formulas are the telescoped sums in the proof of Lemma~\ref{plf:match} (end of the parts
$\varepsilon_a=+1$ and $\varepsilon_a=-1$). In the second, $\varsigma_\tau(r-1-p)=-\kappa(\tau)-a'\Delta$ with
$a'=m-p>j$ and $\kappa_j=\kappa+(j+\frac12)\Delta$, so $\ell^-_p=(a'-j-\frac12)\Delta-\frac{t-1}2-(a'-j)$; with
$1\le a'-j\le m$ this lies between $\frac\Delta2-\frac{t+1}2$ and $m\Delta$. Finally $m\Delta\le(m+1)\Delta\le b$
and $\kappa(\tau)\le b$ (proof of Lemma~\ref{psq:degree}), and the entries of $\varsigma_\tau$ lie between
$-\kappa(\tau)-m\Delta$ and $\frac k2-1\le b$ (\S\ref{plf:ss-L4}).
\end{proof}

\subsubsection{Polynomials in the negative columns}\label{psv:ss-hol}

We add to Proposition~\ref{pss:hol} the case of a positive definite target and a polynomial in the
negative columns of the source. The setting and notation are those of \S\ref{pss:ss-tool} (one place
$\tau$; source $X$ with signs $s$ and order $o$, positive first; target $T$ of rank $n'$; matrices
$w\in\pcMat_{n'\times N}(\C)$; Fock variables $z_{rc}$; $\Omega$ with a character of weight $a$ on
$\auU(X)$).

\begin{proposition}[Holomorphic polynomial vectors, case (iv)]\label{psv:hol4}
Let $T$ be positive definite, $d\in\Z_{\ge0}^N$, and let $P$ be a polynomial involving only the
variables $z_{rc}$ with $s_c=-1$, such that $\sum_rz_{rc}\partial_{rc'}P=0$ for negative $c\ne c'$ with
$c$ before $c'$ in $o$, and $\sum_rz_{rc}\partial_{rc}P=d_cP$ for negative $c$. Then $P$ is holomorphic for
$\Omega$ \textup(Definition~\ref{psr:def-hol}\textup) of weight $u$, $u_c=-a$ for $s_c=1$ and
$u_c=n'-a+d_c$ for $s_c=-1$.
\end{proposition}
\begin{proof}
Here $s'\equiv1$, so $\varepsilon_{rc}=s_c$, $n_T^+=n'$, $n_T^-=0$, and the variables of $P$ have $\varepsilon=-1$.
As in the proof of Proposition~\ref{pss:hol}, $\mathfrak b$ is spanned by the diagonal matrices, the
$E_{cc'}$ with $c$ before $c'$ of the same sign, and the $E_{ac}$ with $s_a=1$, $s_c=-1$.

\emph{Compact part.} Let $Y=E_{cc'}$, $c,c'$ of the same sign. In the formula of
Lemma~\ref{pss:cpt} for $d\Omega(Y)P$ the first sum is $-\sum_{\varepsilon_{rc}=1}z_{rc'}\partial_{rc}P$ (it is
non-zero only when $c$, hence $c'$, is positive), which vanishes because $P$ does not involve the
variables of positive columns; the second sum is $\sum_{\varepsilon_{rc'}=-1}z_{rc}\partial_{rc'}P$ (non-zero only
when $c'$, hence $c$, is negative). The trace term is $(n_T^--a)\delta_{cc'}=-a\,\delta_{cc'}$ for positive
$c$, and $(n_T^+-a)\delta_{cc'}=(n'-a)\delta_{cc'}$ for negative $c$. Hence for $c,c'$ positive
$d\Omega(E_{cc'})P=-a\,\delta_{cc'}P$; for $c,c'$ negative $d\Omega(E_{cc'})P=\sum_rz_{rc}\partial_{rc'}P+(n'-a)
\delta_{cc'}P$, which is $0$ for $c$ before $c'$ by hypothesis and $(d_c+n'-a)P$ for $c=c'$.

\emph{Non-compact part.} Let $Y=E_{ac}$, $s_a=1$, $s_c=-1$. By Lemma~\ref{pss:torus}(ii),
$d\Omega(E_{ac})$ is a linear combination of products $xy$, $x,y\in\{z_i,\partial_i\}$, whose $X$-weights add
up to $e_a-e_c$. By Lemma~\ref{pss:torus}(i) the $X$-weight of $z_{rc''}$ is $-s_{c''}e_{c''}$ and that of
$\partial_{rc''}$ is $s_{c''}e_{c''}$; so the operators of the column $a$ have weights $\mp e_a$ ($z_{ra}$,
$\partial_{ra}$), those of the column $c$ have weights $\pm e_c$ ($z_{rc}$, $\partial_{rc}$), and those of any other
column $c''$ have weights $\pm e_{c''}$ with $c''\ne a,c$. A sum of two such weights equals $e_a-e_c$ only
for $\partial_{ra}\partial_{r'c}$. Since $\partial_{ra}P=0$ ($a$ is positive) and the $\partial$'s commute,
$d\Omega(E_{ac})P=0$. Hence $d\Omega(Y)P=u(Y)P$ for all $Y\in\mathfrak b$.
\end{proof}

\begin{corollary}\label{psv:minors}
Let $T$ be positive definite, and let $c_0,c_1,\dots,c_{q-1}$ be the negative indices of $X$ in the
order $o$. For a finite family $(h_{\mathfrak c})_{\mathfrak c}$ of integers $1\le h_{\mathfrak c}\le q$ and row sets
$R_{\mathfrak c}$ of cardinality $h_{\mathfrak c}$, put $J_{\mathfrak c}=\{c_0,\dots,c_{h_{\mathfrak c}-1}\}$ and
$D=\prod_{\mathfrak c}\det\bigl(z_{R_{\mathfrak c},J_{\mathfrak c}}\bigr)$ \textup(the minor on the rows $R_{\mathfrak c}$ and the
columns $J_{\mathfrak c}$, in the order $o$\textup). Then every linear combination of such products, for a
fixed family $(h_{\mathfrak c})$ and arbitrary row sets, satisfies the hypothesis of
Proposition~\ref{psv:hol4} with $d_{c_p}=\#\{\mathfrak c:h_{\mathfrak c}>p\}$; it is holomorphic of weight
$u_c=-a$ \textup($s_c=1$\textup), $u_{c_p}=n'-a+\#\{\mathfrak c:h_{\mathfrak c}>p\}$.
\end{corollary}
\begin{proof}
For $c\ne c'$ the operator $L=\sum_rz_{rc}\partial_{rc'}$ is a derivation, and on $\det(z_{R,J})$ it replaces
the column $c'$ by the column $c$ if $c'\in J$, and gives $0$ if $c'\notin J$. If $c=c_p$ is before
$c'=c_{p'}$ ($p<p'$) and $c'\in J_{\mathfrak c}$, then $c\in J_{\mathfrak c}$ ($J_{\mathfrak c}$ is an initial segment), and the
resulting determinant has two equal columns. So $L$ kills each factor, hence $D$. For $c=c'=c_p$, $L$
multiplies $\det(z_{R,J_{\mathfrak c}})$ by $1$ if $c_p\in J_{\mathfrak c}$, i.e.\ $h_{\mathfrak c}>p$, and by $0$
otherwise; so $LD=d_{c_p}D$, with $d_{c_p}$ independent of the row sets. Both conditions are linear in
$D$.
\end{proof}

\subsubsection{The hypothesis (c) for $(W,V^\circ)$}\label{psv:ss-c}

By (W7) of \S\ref{po:ss-glue}, $\omega^\circ$ is a Weil representation of the pair $(W,V^\circ)$ twisted by
powers of $\upsilon$ only; in the convention of {\cite[\S2.3]{P53}} (and \S\ref{pte:ss-conv}),
$\omega^\circ=\omega_{\alpha^\circ,\beta^\circ}$ with $\alpha^\circ,\beta^\circ$ powers of $\upsilon$. We need the
exponent of $\alpha^\circ$.

\begin{lemma}[The character of $\omega^\circ$ on $\auU(W)$]\label{psv:char}
$\alpha^\circ=\upsilon^{m/2-j}$; its archimedean weight is $a^\circ=m/2-j$ at every $\tau$. The character of
$\omega''$ on $\auU(W)$ is $\nu''_\vartheta=\mu^0_j\vartheta\,\upsilon^{j+(t+k-1)/2}$, of weight
$a''_\tau=j+\frac{t+k-1}2-\kappa_j(\tau)$, and $a''_\tau+a^\circ=\frac n2-\kappa_j(\tau)$.
\end{lemma}
\begin{proof}
By \S\ref{po:ss-set}, $\auU(W)$ acts in $\omega_\mu$ through $b(\det g)^{-1}$ times the representation
defined with $\upsilon$, with $b=\mu\,\upsilon^{\lfloor n/2\rfloor}$; that is, $\omega_\mu$ restricted to $\auU(W)$ is
$\omega_{b,\cdot}$ in the convention of {\cite[\S2.3]{P53}}. By (W7) the characters of the two factors
multiply to $b$ (``the splitting over $\auU(W)$ for $\auHom(W,V)$ is the product of those for the two
summands''), and by the proof of Proposition~\ref{pte:leftside} the character of $\omega''$ on $\auU(W)$ is
that of the kernel $\theta^{\mathsf W,V''}_{\nu''_\vartheta,\,\cdot}$ restricted to $W$, namely $\nu''_\vartheta=\mu^0_j\vartheta
\upsilon^{j+\lfloor(t+k-1)/2\rfloor}$ (\S\ref{pte:ss-third}). As $\mu=\mu^0_j\vartheta$, $\alpha^\circ=b/\nu''_\vartheta=
\upsilon^{\lfloor n/2\rfloor-j-\lfloor(t+k-1)/2\rfloor}$. Here $n=6r$ and $t+k-1=5r+k-1$ are even ($r$ odd, $k$
even), and $3r-j-\frac{5r+k-1}2=\frac{r-k+1}2-j=\frac m2-j$ since $r=k+m-1$. The weights: $\upsilon_\tau(z)=z$ on
$\auU(1)$, $\vartheta$ has weight $0$ and $\mu^0_j$ has weight $-\kappa_j(\tau)$ (proof of Lemma~\ref{plf:const}).
Finally $a''_\tau+a^\circ=\frac{t+k-1+m}2-\kappa_j(\tau)=3r-\kappa_j(\tau)$.
\end{proof}

Let $n^\circ=\operatorname{rk}V^\circ$ (positive definite). For $u^\circ\in W^\circ=W_{\mu^-}(\C^{n^\circ})$ the
polynomial $\pop^\circ_{u^\circ}$ of Lemma~\ref{po:prod}(M3) is
$\pop^\circ_{u^\circ}(x^\circ)=\langle(u^\circ)^\natural,\mathrm{ev}^-(x^{\circ-})^\flat\rangle$, where
$\mathrm{ev}^-(y')=\bigotimes_{\mathfrak c}y'_1\wedge\dots\wedge y'_{h_{\mathfrak c}}$, the product over the columns
$\mathfrak c$ of $\mu^-_\tau$ (all $\tau$), with $y'_i$ the image under $\flat$ of the column number $i$ of
$x^{\circ-}$ in the reverse numbering of the negative basis vectors of $W$, i.e.\ (Definition~\ref{psv:order})
of the negative column $c_{i-1}$ of $W$ in the order $o$.

\begin{lemma}[The coordinates of $K^\circ$]\label{psv:coords}
For every $u^\circ\in W^\circ$, $\pop^\circ_{u^\circ}$ is a linear combination of the products
$\prod_\tau\prod_{\mathfrak c}\det(z_{R_{\mathfrak c},J_{\mathfrak c}})$ of Corollary~\ref{psv:minors} \textup(at each
$\tau$, with $X=W$, $T=V^\circ$, the columns $\mathfrak c$ of $\mu^-_\tau$ of heights $h_{\mathfrak c}$, and row sets
$R_{\mathfrak c}$ of $V^\circ$\textup). The span $\mathcal Q^\circ$ of the $\pop^\circ_{u^\circ}$, $u^\circ\in W^\circ$, is
finite dimensional and stable under $\Omega'(y)$, $y\in\auU(V^\circ)_\infty$, and $\pop^\circ_{u^\circ}\ne0$ for
$u^\circ\ne0$.
\end{lemma}
\begin{proof}
Expanding the wedge products in the basis of $\bigotimes_{\mathfrak c}\Lambda^{h_{\mathfrak c}}(\C^{n^\circ})^\vee$
dual to the basis of pure tensors of basis vectors, the coordinate of $\mathrm{ev}^-(y')$ on the basis vector
of index $(R_{\mathfrak c})$ is $\prod_{\mathfrak c}\det\bigl(y'_i(e_l)\bigr)_{l\in R_{\mathfrak c},\,i\le h_{\mathfrak c}}$, and
$y'_i(e_l)$ is the entry of $x^\circ$ in the row $l$ and the column $c_{i-1}$, or its conjugate, according to
the convention for $\flat$. By \eqref{po:eq-poly} and the remark following it, $\pop^\circ_{u^\circ}$ is a
polynomial in the Fock variables, which for the entries of negative product sign (rows of $V^\circ$,
negative columns of $W$) are the entries themselves (\S\ref{ptf:ss-fock}); this fixes the conjugations,
as in the computation of $\pop_{\xi_0}$ in \S\ref{plf:ss-L4}, where the negative factor is displayed as a
product of minors $\det(z_{a\,c(a')})$ in the Fock variables. Pairing with $(u^\circ)^\natural$ gives
$\pop^\circ_{u^\circ}=\sum_{(R_{\mathfrak c})}\gamma_{(R_{\mathfrak c})}(u^\circ)\prod_{\mathfrak c}\det(z_{R_{\mathfrak c},J_{\mathfrak c}})$ with
$J_{\mathfrak c}=\{c_0,\dots,c_{h_{\mathfrak c}-1}\}$ and scalars $\gamma$. This is the first assertion. $\mathcal Q^\circ$ lies
in the finite-dimensional space of polynomials of bounded degree in finitely many variables. The map
$u^\circ\mapsto\pop^\circ_{u^\circ}$ is conjugate-linear and $\auU(V^\circ)(\R)$-equivariant (\S\ref{po:ss-set}:
$\xi\mapsto\pop_\xi$ is $\auU(V)(\R)$-equivariant, and by (W7) and Lemma~\ref{po:prod}(M1),(M3) the block
$V^\circ$ acts on the factor $\pop^\circ$), so $\mathcal Q^\circ$ is stable under $\Omega'$. Finally, the vectors
$\mathrm{ev}^-(y')$ span $W^\circ$ (definition of $W^\circ$ as in \eqref{po:eq-ev}), so for $u^\circ\ne0$ some
$\langle(u^\circ)^\natural,\mathrm{ev}^-(y')\rangle\ne0$; as $x^{\circ-}\mapsto(y'_i)$ is onto, $\pop^\circ_{u^\circ}\ne0$.
\end{proof}

\begin{theorem}[Hypothesis (c)]\label{psv:c}
Let $u^\circ$ be the weight on $W$ given at $\tau$, in the fixed numbering, by
\[
u^\circ_{\tau,a}=-a^\circ=j-\tfrac m2\quad(0\le a<r_+),\qquad u^\circ_{\tau,c}=n^\circ-a^\circ+\ell^-_{p(c)}
\quad(r_+\le c<r),
\]
with $p(c)=r-1-c$ the position of $c$ in $o$ among the negative indices and $\ell^-_p$ the row lengths of
$\mu^-_\tau$ \textup(Lemma~\ref{psv:rows}\textup). Then every non-zero element of $\mathcal Q^\circ$ is a non-zero
holomorphic vector of weight $u^\circ$ for $\Omega=\omega^\circ_\infty(\cdot,1)$
\textup(Definition~\ref{psr:def-hol}\textup). Consequently the hypothesis on $\mathcal Q$ in
Theorem~\ref{psr:R} holds for $(X,Y)=(W,V^\circ)$, the weight $u^\circ$, the characters of $\omega^\circ$, and
$\mathcal Q=\mathcal Q^\circ$, the span of the coordinates of $K^\circ$ against a basis of $W^\circ$. Moreover, for every
$h^\circ\in\auU(V^\circ)(\auA_f)$, every coordinate of $g\mapsto K^\circ(g,h^\circ)$ is holomorphic of weight
$u^\circ$ on $\auU(W)(\auA)$.
\end{theorem}
\begin{proof}
At $\tau$, by Lemma~\ref{psv:coords} every element of $\mathcal Q^\circ$ is a linear combination of products of
minors as in Corollary~\ref{psv:minors} with $X=W$ (order $o$), $T=V^\circ$, $n'=n^\circ$, $a=a^\circ$
(Lemma~\ref{psv:char}) and the column heights of $\mu^-_\tau$; there $d_{c_p}=\#\{\mathfrak c:h_{\mathfrak c}>p\}$ is the
length $\ell^-_p$ of the row number $p$ of $\mu^-_\tau$, and $c_p$ is the negative index $c$ with $p(c)=p$,
i.e.\ $c=r-1-p$. So Corollary~\ref{psv:minors} gives the $\mathfrak b$-condition at $\tau$ with the weight
$u^\circ_\tau$; at the other places the polynomial is a product of factors, and $d\Omega$ of an element of the
Lie algebra at $\tau$ acts on the factor at $\tau$ only (\S\ref{pss:ss-tool}, statements for $\prod_\tau$ by
taking products), so the condition holds for $\prod_\tau\mathfrak b$. Polynomial vectors are smooth (Howe~\cite[(2.2)--(2.3)]{HoweJAMS},
\S\ref{pss:ss-tool}). Non-vanishing and $\Omega'$-stability are Lemma~\ref{psv:coords}.

For the last assertion fix $g\in\auU(W)(\auA)$ and $h^\circ$. By the definition in \S\ref{po:ss-glue}, for
$x\in\auU(W)_\infty$, $\langle u^\circ,K^\circ(gx,h^\circ)\rangle=\theta^\circ\bigl(\omega^\circ(g,h^\circ)(\Omega(x)
\pop^\circ_{u^\circ}\otimes\Psi^\circ)\bigr)=\Lambda(\Omega(x)\pop^\circ_{u^\circ})$ with $\Lambda(w)=\theta^\circ(\omega^\circ(g,h^\circ)
(w\otimes\Psi^\circ))$. As in the proof of Proposition~\ref{pss:carrier}, $\Lambda$ is a continuous linear form on
the smooth vectors (smooth vectors are Schwartz functions in a Schr\"odinger model,
Lemma~\ref{psp:prodtheta}, on which $\omega^\circ(g,h^\circ)$ and the theta functional are continuous).
Corollary~\ref{pss:coefs}(a) shows that $x\mapsto\langle u^\circ,K^\circ(gx,h^\circ)\rangle$ is holomorphic of weight
$u^\circ$. Taking $u^\circ$ in a basis of $W^\circ$ gives the coordinates of $K^\circ(\cdot,h^\circ)$ (up to complex
conjugation of the coefficients of the basis change, which does not affect holomorphy).
\end{proof}

\noindent The weights $u^\circ_\tau$ are determined by the shift $j-\frac m2$, the rank $n^\circ$ of $V^\circ$, and the
row lengths of $\mu^-$ on the negative indices.

\subsubsection{The datum of $K^\circ$}\label{psv:ss-datum}

\begin{proposition}[The datum of $K^\circ$ is the point datum]\label{psv:datum}
Let $\pcm=(\pi_M^2)$ and let $\iota\in\pcMat_{n^\circ\times r}(\auO_E/\pcm)$ be the matrix of the inclusion of the
frame of $W$ into that of $V^\circ$ as the marked coordinates \textup(the first $r$ coordinates of $V^\circ$,
\S\ref{po:ss-glue}\textup), i.e.\ $\iota=\iota_z\bmod\pcm$. Then the finite part $\Psi^\circ=\mathbf 1[v\equiv\iota_z
\bmod\pi_M^2]\otimes(\text{vacuum off }\mathsf M)$ of the datum of $K^\circ$ is the finite part
$\delta_\iota^\natural\otimes\phi^{\mathrm{vac},\pcm}$ of the datum $[\,\cdot\,,\delta_\iota]$ of {\cite[\S2.3]{P53}}.
Hence for $u^\circ\in W^\circ$, $g\in\auU(W)(\auA)$ and $h^\circ\in\auU(V^\circ)(\auA_f)$,
\[
\langle u^\circ,K^\circ(g,h^\circ)\rangle=\theta_{\alpha^\circ,\beta^\circ}\bigl(g,h^\circ;\pop^\circ_{u^\circ},\delta_\iota\bigr)
\]
in the notation of {\cite[\S2.3]{P53}}, and Theorem~\ref{psr:R} applies to $(X,Y)=(W,V^\circ)$ with the datum of
$K^\circ$ as it stands: no variant of Theorem~\ref{psr:R} is needed.
\end{proposition}
\begin{proof}
\emph{The modulus.} $\pi_M$ is a uniformiser at every $w\in\mathsf M$ and a unit at $\aubar w$ and at all other
primes (\S\ref{po:ss-glue}), so $\pcm=\prod_{w\in\mathsf M}w^2$ is a modulus in the sense of {\cite[\S2.3]{P53}}
(prime to $\aubar\pcm$), $S(\pcm)$ is the set of places of $E^+$ below $\mathsf M$ (non-empty), and
$\auO_E/\pcm=\prod_{w\in\mathsf M}\auO_w/w^2$.

\emph{The point datum.} By {\cite[\S2.3]{P53}}, $\delta_\iota$ is the indicator of the single point $\iota$ of
$\pcMat_{n^\circ\times r}(\auO_E/\pcm)$, and $\delta_\iota^\natural$ is the function on $\prod_{w\in\mathsf M}
\pcMat_{n^\circ\times r}(E_w)$ equal to $\delta_\iota(x\bmod\pcm)$ at integral $x$ and $0$ elsewhere, that is,
$\delta_\iota^\natural(x)=\mathbf 1[x\text{ integral and }x\equiv\iota_z\bmod\pi_M^2]$, which is the factor
$\mathbf 1[v\equiv\iota_z\bmod\pi_M^2]$ of $\Psi^\circ$ (both are functions of the $E_w$-coordinates, $w\in\mathsf M$,
in the model (W5) at the split places of $S(\pcm)$, with matrices written in the frames; the normalised
datum of \S\ref{po:ss-glue} is a function of integral matrices $v\in\pcMat_{n^\circ\times r}(\auO_{\mathsf M})$). Off
$S(\pcm)$ both data are the normalised vacuum of the same lattice. With the convention of
{\cite[Definition~6.31]{P53}}, $\delta_\iota^\natural(x)=\delta_\iota(a^{-1}x\bmod\pcm)$ on $a\,\pcMat(\widehat\auO)$; at the
primes of $\mathsf M$ the element $\alpha_w$ is a unit (the primes where it is not lie above $q_1$,
{\cite[Proposition~6.9(a)]{P53}}, and the chosen places do not), so $a_w=1$ and the two conventions agree there.

\emph{The kernel.} By Lemma~\ref{psv:char} and (W7), $\omega^\circ=\omega_{\alpha^\circ,\beta^\circ}$ and $\theta^\circ$
is the theta functional; so $\theta^\circ(g,h^\circ;u^\circ)=\theta(\omega_{\alpha^\circ,\beta^\circ}(g,h^\circ)[\pop^\circ_{u^\circ},
\delta_\iota])$, which is the displayed identity.

\emph{Where the datum enters the proof of Theorem~\ref{psr:R}.} It enters in exactly two places
(Remark~\ref{psr:rem}(ii)): Step~1 uses that $\omega_{\alpha,\beta}(1,k)$ fixes $[v,\delta_\iota]$ for $k\in J_Y(\pcm)$
({\cite[Lemma~2.9]{P53}}: $\bar k=1$); Step~2 uses that $\delta_\iota$ is sharp for $J_X(\pcm)$ (with $A_0={}^t\iota$,
$A_0\iota=1$), which is also what makes {\cite[Lemma~2.14]{P53}} applicable in Proposition~\ref{psr:rallis}. Both
are properties of $\delta_\iota$ as a function on $\pcMat_{n^\circ\times r}(\auO_E/\pcm)$, and we have just seen that
the datum of $K^\circ$ is $\delta_\iota$ itself. The characters $\alpha^\circ,\beta^\circ$ are powers of $\upsilon$, unramified at
every finite place (Theorem~\ref{pss:a}), and $V^\circ$ is positive definite of rank $n^\circ=5r-9m-1>2r$ (as
$r\ge x^5$ and $m\le2\sqrt x$).
\end{proof}

\subsubsection{The input (H4)}\label{psv:ss-H4}

\emph{The conjugate of $y$.} Recall (\S\ref{plf:ss-L4}, {\cite[Definition~8.1]{P49}}) that $\poA(W_j)$ consists
of the bounded continuous functions $y$ on $\auU(W)(\auA)$, left invariant under $\auU(W)(E^+)$, of level $J_M$,
with $R(e_W)y=y$, where $e_W=\bigotimes_\tau d_\tau\aubar{c_{W,\tau}}$, $c_{W,\tau}(x)=\langle v,\pi(x)v\rangle$, $\pi=\pi_{\varsigma_\tau}$
the discrete series of Harish-Chandra parameter $\varsigma_\tau$ and $v$ a unit highest weight vector of its lowest
$K$-type. Let $\Lambda^K(\varsigma_\tau)=\varsigma_\tau+\rho_{\mathrm n}-\rho_{\mathrm c}$ be the highest weight of that $K$-type,
$\rho_{\mathrm c},\rho_{\mathrm n}$ the half sums of the compact and non-compact roots positive for $\varsigma_\tau$.

\begin{lemma}[The vector $v$]\label{psv:v}
At $\tau$, let $\lambda=\Lambda^K(\varsigma_\tau)$ and let $\mathfrak b^-$ be the lower triangular matrices for the order
$o$ of Definition~\ref{psv:order}. Then $d\pi(Z)v=\lambda(Z)v$ for $Z\in\mathfrak b^-$, and in the fixed numbering
\[
\lambda_a=\varsigma_\tau(a)+\tfrac{r_-}2-\tfrac{r_+-1}2+a\quad(0\le a<r_+),\qquad
\lambda_c=\varsigma_\tau(c)-\tfrac{r_+}2-\tfrac{r_--1}2+(c-r_+)\quad(r_+\le c<r);
\]
these are integers.
\end{lemma}
\begin{proof}
The positive roots for $\varsigma_\tau$ are the $e_a-e_b$ with $\varsigma_\tau(a)>\varsigma_\tau(b)$: the compact ones join
indices of one sign block with $a<b$ in the fixed numbering, and every non-compact one is $e_a-e_c$
with $a$ positive and $c$ negative, since every entry of $\varsigma_\tau$ on a positive coordinate exceeds every
entry on a negative one (\S\ref{plf:ss-L4}). So, in the fixed numbering, $\rho_{\mathrm c,a}=\frac{r_+-1}2-a$,
$\rho_{\mathrm n,a}=\frac{r_-}2$ for positive $a$, and $\rho_{\mathrm c,c}=\frac{r_--1}2-(c-r_+)$, $\rho_{\mathrm n,c}=-\frac{r_+}2$
for negative $c$, which gives the formula; $\frac{r_-}2-\frac{r_+-1}2=\frac{r-2r_++1}2$ and $-\frac{r_+}2-\frac{r_--1}2=
-\frac{r-1}2$ are integers as $r$ is odd, and the entries of $\varsigma_\tau$ are integers.

The diagonal part of $\mathfrak b^-$ acts on $v$ by its weight $\lambda$. A compact element $E_{ba}$ of
$\mathfrak b^-$ has $a$ before $b$ in $o$ in one sign block, i.e.\ $\varsigma_\tau(b)>\varsigma_\tau(a)$
(Definition~\ref{psv:order}); so $e_b-e_a$ is a positive compact root and $d\pi(E_{ba})v=0$, $v$ being a highest
weight vector. A non-compact element of $\mathfrak b^-$ is $E_{ca}$ with $c$ negative and $a$ positive. If
$d\pi(E_{ca})v\ne0$, it is a vector of $T$-weight $\lambda+e_c-e_a$ in some $K$-type of $\pi$, of highest weight
$\mu'$, so $\lambda+e_c-e_a=\mu'-\sum_\alpha m_\alpha\alpha$ with $m_\alpha\ge0$ over positive compact roots, and by
(D4) of {\cite[\S4.2.3]{P53}} $\mu'=\lambda+\sum_\beta n_\beta\beta$, $n_\beta\ge0$, over the positive non-compact roots
$\beta=e_{a'}-e_{c'}$. The linear form $H(x)=\sum_{c'\text{ negative}}x_{c'}$ vanishes on compact roots, is $-1$ on
each $\beta$ and $+1$ on $e_c-e_a$; so $1=-\sum n_\beta\le0$, a contradiction.
\end{proof}

\begin{theorem}[(H4) for $\aubar y$]\label{psv:ybar}
Let $u$ be the weight on $W$ given at $\tau$ by $u_\tau=-\Lambda^K(\varsigma_\tau)$, i.e.\ in the fixed numbering
\[
u_{\tau,a}=-\varsigma_\tau(a)-\tfrac{r-2r_++1}2-a\quad(0\le a<r_+),\qquad
u_{\tau,c}=-\varsigma_\tau(c)+\tfrac{r-1}2-(c-r_+)\quad(r_+\le c<r).
\]
For every form $y$ of type $e_W$ \textup(in particular every $y\in e_0\poA(W_j)$\textup), $\aubar y$ is holomorphic of
weight $u$ on $\auU(W)(\auA)$ for the order $o$.
\end{theorem}
\begin{proof}
Fix $g\in\auU(W)(\auA)$ and put $f(h)=\aubar y(gh)$ for $h\in\pspG=\auU(W)_\infty$. Write $c=c_W=\bigotimes_\tau c_{W,\tau}$,
$d=\prod_\tau d_\tau$, and $\kappa(z)=c(z^{-1})=\langle\pi(z)v,v\rangle$ ($\pi=\bigotimes_\tau\pi_{\varsigma_\tau}$, $v=\bigotimes_\tau v_\tau$).
Conjugating $y=R(e_W)y$ and substituting $x\mapsto h^{-1}x$ (Haar measure is invariant),
\[
f(h)=d\int_{\pspG}c(x)\,\aubar y(ghx)\,dx=d\int_{\pspG}c(h^{-1}x)\,\aubar y(gx)\,dx=d\int_{\pspG}\kappa(x^{-1}h)\,\aubar y(gx)\,dx .
\]
\emph{Smoothness.} By Lemma~\ref{plf:A12} the conclusion of {\cite[Theorem~3.2]{P49}} holds for $c$: it is smooth and
$\mathsf Xc$ is integrable for every left- or right-invariant differential operator $\mathsf X$; as inversion
preserves Haar measure and exchanges left and right invariant operators, the same holds for $\kappa$. For real
$X\in\pssg$ and $R_X$ the right derivative, $\kappa(x^{-1}he^{sX})-\kappa(x^{-1}h)=\int_0^s(R_X\kappa)(x^{-1}he^{\sigma X})\,
d\sigma$; since $\aubar y$ is bounded and $R_X\kappa$ is integrable, Fubini's theorem and the continuity of right
translation in $L^1(\pspG)$ give that $f$ is differentiable along $X$ with
$R_Xf(h)=d\int(R_X\kappa)(x^{-1}h)\,\aubar y(gx)\,dx$, a continuous function of $h$. Iterating, $f$ is smooth, and by
complex linearity the same formula holds for $R_Y$, $Y\in\pssg_\C$.

\emph{The $\mathfrak b$-condition.} For real $X$, $(R_X\kappa)(z)=\frac d{ds}\langle\pi(ze^{sX})v,v\rangle|_0=\langle\pi(z)d\pi(X)v,v\rangle$.
For $Y=Y_1+iY_2$, $Y_1,Y_2\in\pssg$, since the inner product is conjugate-linear in the first variable,
$(R_Y\kappa)(z)=\langle\pi(z)d\pi(Y_1)v,v\rangle+i\langle\pi(z)d\pi(Y_2)v,v\rangle=\langle\pi(z)d\pi(Y^\natural)v,v\rangle$ with
$Y^\natural=Y_1-iY_2=-\varepsilon Y^*\varepsilon$, the conjugation of $\pssg_\C$ fixing $\pssg$ (proof of
Corollary~\ref{pss:coefs}). For $Y\in\prod_\tau\mathfrak b$, $Y^\natural$ is lower triangular, so by Lemma~\ref{psv:v}
$d\pi(Y^\natural)v=\lambda(Y^\natural)v$ with $\lambda(Y^\natural)=\sum\lambda_{\tau,a}(Y^\natural)_{\tau,aa}=-\sum\lambda_{\tau,a}\aubar{Y_{\tau,aa}}$.
Hence $(R_Y\kappa)(z)=\aubar{\lambda(Y^\natural)}\kappa(z)=(-\lambda)(Y)\kappa(z)=u(Y)\kappa(z)$, and by the formula for $R_Yf$,
$R_Yf=u(Y)f$ for all $Y\in\prod_\tau\mathfrak b$. Lemma~\ref{pss:ext} shows that $f$ is holomorphic of weight $u$. As $g$ is
arbitrary, $\aubar y$ is holomorphic of weight $u$.
\end{proof}

\emph{The functions $F_i$.} By Proposition~\ref{pte:third} and Proposition~\ref{pte:leftside} for the trivial twist
$\vartheta=1$, for every choice $i$ there is $c_i^{\mathrm F}\in\C$ (namely $c_3(1)\gamma(1)/\mathsf N$) with
\begin{equation}\label{psv:eq-F}
F_i=c_i^{\mathrm F}\,\Xi_{1,i},\qquad\Xi_{1,i}(g)=\sum_{h}\theta''_{1,\Lambda_i}\bigl(g,h;\varphi''_{1,i}(h)\bigr),
\end{equation}
a finite sum over the classes $h$ of $\auU(V'')$ modulo $J'_{1,i}$ (Definition~\ref{pte:def-ff}), where
$\theta''_{1,\Lambda_i}(g,h;u'')=\theta''(\omega''(g,h)(\pop''_{u''}\otimes\Psi''_{1,\Lambda_i}))$ for $u''\in W''=W_{\mu^+}$
(\S\ref{po:ss-glue}).

\begin{theorem}[(H4) for the $F_i$]\label{psv:F}
Let $u''$ be the weight on $W$ given at $\tau$ by
\[
u''_{\tau,a}=-\ell^+_a-a''_\tau\quad(0\le a<r_+),\qquad u''_{\tau,c}=N''-a''_\tau\quad(r_+\le c<r),
\]
with $\ell^+_a$ the row lengths of $\mu^+_\tau$ \textup(Lemma~\ref{psv:rows}\textup) and $a''_\tau$ as in
Lemma~\ref{psv:char}. Then every linear combination $F=\sum_ic_iF_i$ is holomorphic of weight $u''$ on
$\auU(W)(\auA)$ for the order $o$, and $u''=u-u^\circ$, with $u$ of Theorem~\ref{psv:ybar} and $u^\circ$ of
Theorem~\ref{psv:c}.
\end{theorem}
\begin{proof}
\emph{Holomorphy.} By \eqref{psv:eq-F} it suffices that $g\mapsto\theta''_{1,\Lambda_i}(g,h;u'')$ is holomorphic of weight
$u''$ for each $h$ and $u''\in W''$ (sums of holomorphic functions of one weight are holomorphic of that weight:
take the sum of the extensions on the intersection of their domains). By the proof of Lemma~\ref{ptf:laplace},
$\pop''_{u''}$ is a linear combination of the products $\prod_\tau\prod_{\mathfrak c}\det X_{\tau;R_{\mathfrak c},[h_{\mathfrak c}]}$ over
the columns $\mathfrak c$ of $\mu^+_\tau$, where $X_\tau$ are the Fock variables of the positive columns of $W$ (rows of
$V''$, product sign $+1$) and $[h]$ is the set of the first $h$ positive columns in the fixed numbering, i.e.\ the
\emph{last} $h$ positive columns in the order $o$. At $\tau$ apply Proposition~\ref{pss:hol}(i) with $X=W$, $T=V''$
(positive definite, $n'=N''$) and $a=a''_\tau$ (Lemma~\ref{psv:char}): $P$ involves only variables of positive columns;
for positive $c\ne c'$ with $c$ before $c'$ in $o$, the derivation $\sum_rz_{rc'}\partial_{rc}$ replaces the column $c$ by
the column $c'$ in each minor containing $c$, and then $c'\in[h_{\mathfrak c}]$ as well ($[h_{\mathfrak c}]$ is a final segment of
$o$), so it kills each product; and $\sum_rz_{rc}\partial_{rc}$ multiplies each product by $d_c=\#\{\mathfrak c:c<h_{\mathfrak c}\}$
($c$ in the fixed numbering), which is the length $\ell^+_c$ of the row number $c$ of $\mu^+_\tau$. So $\pop''_{u''}$ is a
holomorphic vector of weight $u''$ for $\Omega''=\omega''_\infty(\cdot,1)$ (at each place, hence for the product, as
in the proof of Theorem~\ref{psv:c}). For $x\in\auU(W)_\infty$, $\theta''_{1,\Lambda_i}(gx,h;u'')=\Lambda(\Omega''(x)\pop''_{u''})$
with $\Lambda(w)=\theta''(\omega''(g,h)(w\otimes\Psi''_{1,\Lambda_i}))$ continuous on smooth vectors (as in the proof of
Theorem~\ref{psv:c}), and Corollary~\ref{pss:coefs}(a) gives the holomorphy.

\emph{The identity $u''+u^\circ=u$.} Positive $a$: by Lemma~\ref{psv:rows} and Lemma~\ref{psv:char},
$u''_a+u^\circ_a=-\ell^+_a-(a''+a^\circ)=-\varsigma_\tau(a)-\kappa_j+\frac{t-1}2+(r_+-a)-\frac n2+\kappa_j=-\varsigma_\tau(a)-\frac{r+1}2+r_+-a$,
which is $u_a$. Negative $c$, with $p=p(c)=r-1-c$: $u''_c+u^\circ_c=N''+n^\circ-(a''+a^\circ)+\ell^-_p=n-\frac n2+\kappa_j+\ell^-_p$,
and $\ell^-_p=-\kappa_j-\frac{t-1}2-\varsigma_\tau(c)-(r_--p)$ with $r_--p=c-r_++1$; so the sum is
$\frac n2-\frac{t-1}2-\varsigma_\tau(c)-(c-r_+)-1=-\varsigma_\tau(c)+\frac{r-1}2-(c-r_+)=u_c$ (using $N''+n^\circ=n=r+t$).
\end{proof}

\begin{proposition}[Bounds on $u$ and $u^\circ$]\label{psv:bounds}
$u$ and $u^\circ$ are dominant for $o$ with gap $\ge300\pspn$, and $|u_{\tau,a}|,|u^\circ_{\tau,a}|\le5b$.
\end{proposition}
\begin{proof}
\emph{$u$.} Along $o$ in the positive block the fixed index $a$ decreases; $u_a=-\varsigma_\tau(a)-a+\mathrm{const}$ and
$\varsigma_\tau(a)-\varsigma_\tau(a+1)\ge1$ (integers, strictly decreasing) give $u_{a+1}\ge u_a$, so $u$ is non-increasing along
$o$; the negative block is the same. Gap: for $a$ positive and $c$ negative, $u_c-u_a=\varsigma_\tau(a)-\varsigma_\tau(c)+
\frac{r-1}2+\frac{r-2r_++1}2+a-(c-r_+)=\varsigma_\tau(a)-\varsigma_\tau(c)+r_--(c-r_+)+a\ge\Delta+1$, since entries on coordinates of
different signs differ by at least $\Delta$ (\S\ref{plf:ss-L4}) and $c-r_+\le r_--1$; and $\Delta\ge300\pspn$ by
Lemma~\ref{psp:sizes}(ii). Entries: $|u_{\tau,a}|\le\max|\varsigma_\tau|+r\le2b+b$ by Lemma~\ref{psv:rows} and $r\le t\le\pspn\le b$
(Lemma~\ref{psp:sizes}(i)).

\emph{$u^\circ$.} It is constant on the positive block. Along $o$ in the negative block $p(c)$ increases and
$\ell^-_p$ is non-increasing in $p$ (rows of a diagram, longest first), so $u^\circ$ is non-increasing. Gap:
$u^\circ_c-u^\circ_a=n^\circ+\ell^-_{p(c)}\ge\ell^-_{p(c)}\ge\frac\Delta2-\frac{t+1}2\ge300\pspn$ by Lemma~\ref{psv:rows} and
Lemma~\ref{psp:sizes}(ii) ($300\pspn+t\le\Delta/2$). Entries: $|a^\circ|\le m/2$, $n^\circ\le t\le b$ and $\ell^-_p\le b$
(Lemma~\ref{psv:rows}), so $|u^\circ_{\tau,a}|\le m+2b\le3b$.
\end{proof}

\begin{theorem}[(H4)]\label{psv:H4}
The input \textup{(H4)} of Proposition~\ref{psp:inputs} holds, for the order $o$ of Definition~\ref{psv:order} on
$W$ and the standard order on $S$: for $y\in e_0\poA(W_j)$, $\aubar y$ is holomorphic of the weight $u$ of
Theorem~\ref{psv:ybar}; every combination $F$ of the $F_i$ is holomorphic of weight $u-u^\circ$, where $u^\circ$
\textup(Theorem~\ref{psv:c}\textup) is the weight of $g\mapsto K^\circ(g,h^\circ)$; $u,u^\circ$ are dominant with gap
$\ge300\pspn$ and entries of modulus $\le5b$; and the carrier is holomorphic of weight $(0;3m)$.
\end{theorem}
\begin{proof}
Theorems~\ref{psv:ybar}, \ref{psv:F}, \ref{psv:c}, Proposition~\ref{psv:bounds} and Proposition~\ref{pss:carrier}.
\end{proof}

\subsubsection{The character split in Lemma~\ref{psp:step1}}\label{psv:ss-split}

\begin{lemma}[Splitting the kernel of $(S,W')$]\label{psv:split}
Let $\alpha=\nu^\flat\upsilon^m$ and $\beta=\upsilon^m$ be the characters of \eqref{pte:eq-F}, on $\auU(S)$ and on $\auU(W')$. For
$x\in\auU(S)(\auA)$, $\tilde h\in\auU(\widetilde W)(\auA)$, $h^\flat\in\auU(W^\flat)(\auA)$, polynomials $\widetilde P$ on
$\auHom(S,\widetilde W)\otimes\R$ and $Q$ on $\auHom(S,W^\flat)\otimes\R$, and functions $\Psi$, $\Psi^\flat$ modulo $\pi_M^2$,
\[
\theta^{S,W'}_{\alpha,\beta}\bigl(x,(\tilde h,h^\flat);[\widetilde PQ,\Psi\otimes\Psi^\flat]\bigr)=
\theta^{S,\widetilde W}_{\alpha,\beta}\bigl(x,\tilde h;[\widetilde P,\Psi]\bigr)\cdot\theta^{S,W^\flat}_{1,\beta}\bigl(x,h^\flat;[Q,\Psi^\flat]\bigr).
\]
In particular the kernel of $(S,W^\flat)$ in Lemma~\ref{psp:step1} has the characters $(1,\upsilon^m)$ and that of
$(S,\widetilde W)$ the characters $(\nu^\flat\upsilon^m,\ \upsilon^m|_{\auU(\widetilde W)})$, as used in
Theorems~\ref{pss:H5}, \ref{pss:b}, \ref{pss:a} and Lemma~\ref{psp:weights}. The other attachment (the factor
$\alpha(\det x)^{-1}$ on the $(S,W^\flat)$ kernel) gives the same product but is incompatible with the weights used there.
\end{lemma}
\begin{proof}
By the convention of \S\ref{pte:ss-conv}, $\omega_{\alpha,\beta}(x,h)=\alpha(\det x)^{-1}\beta(\det h)\,\omega_{\auHom(S,W')}(x,h)$. For
$h=(\tilde h,h^\flat)$ block diagonal, $\beta(\det h)=\beta(\det\tilde h)\beta(\det h^\flat)$. The compatibility of the
Heisenberg and Weil representations, of the splittings defined with $\upsilon$, and of the theta functionals with orthogonal
sums, which is the content of (W7) in \S\ref{po:ss-glue} (there for the summands of the target $V=V''\perp V^\circ$ of
$\auHom(W,V)$; the argument is the same for the summands $W'=\widetilde W\perp W^\flat$ of the target of $\auHom(S,W')$:
$\auHom(S,W')=\auHom(S,\widetilde W)\oplus\auHom(S,W^\flat)$ orthogonally, and the splitting of $\auU(S)$ for the sum is the
product of the splittings of the summands, the characters multiplying; at the infinite places this is visible on
\eqref{ptf:eq-compact}, $\Delta(\gamma)=\aubar{\det\gamma^-}$ being multiplicative on block-diagonal $\gamma$), gives
$\omega_{\auHom(S,W')}(x,h)=\omega_{\auHom(S,\widetilde W)}(x,\tilde h)\otimes\omega_{\auHom(S,W^\flat)}(x,h^\flat)$,
$\theta=\theta\otimes\theta$, and the datum $[\widetilde PQ,\Psi\otimes\Psi^\flat]$ is the tensor product of the two data (the
Fock vector of a product of polynomials in disjoint sets of variables is the tensor product of their Fock vectors, and the
finite data are products along the rows). The scalar $\alpha(\det x)^{-1}$ appears once; attaching it to the first
factor gives the identity. For the last sentence: with the characters $(\alpha,\beta)$ on $(S,W^\flat)$ the vector
$v_0$ of Theorem~\ref{pss:b} would have weight $(-\Delta'-a_\tau;6m+1-a_\tau)$, $a_\tau=m-\mu^+_{1,\tau}$, by
Corollary~\ref{pss:minors}(i), and with $(1,\beta)$ on $(S,\widetilde W)$ the weight $u_2$ of Theorem~\ref{pss:H5} would
change by $+a_\tau$ on all entries; the relation $u_3=(0;3m)+u_1+u_2$ is unchanged, but $u_1$ would no longer be
the weight $(-\Delta';6m+1)$ fixed in \S\ref{psp:ss-outline}, whose gap is checked in Lemma~\ref{psp:weights}. With the
attachment of the statement, the relation $u_3=(0;3m)+u_1+u_2$ holds: $(-\Delta';6m+1)+(\Delta';\mu^+_1+r+2m)+(0;3m)=
(0;\mu^+_1+r+11m+1)=(0;\mu^+_1+N''+2m)$.
\end{proof}

\noindent Lemma~\ref{psv:split} is the character split used in the proof of Lemma~\ref{psp:step1}.

\subsubsection{Conclusion}\label{psv:ss-main}

\begin{theorem}[Products span]\label{psv:main}
All inputs of Proposition~\ref{psp:inputs} hold in the form in which they are used in \S\ref{psp:sec}:
\begin{enumerate}
\item[(H1)--(H3)] Theorem~\ref{psq:H123} \textup(the gap $300\pspn$ exceeds $N^2+N+1$ for $N\le t$\textup);
\item[(H4)] Theorem~\ref{psv:H4}, for the order $o$ of Definition~\ref{psv:order} on $W$;
\item[(H5)] Theorem~\ref{pss:H5};
\item[(R)] Theorem~\ref{psr:R}, applied in Lemma~\ref{psp:step1} to $(S,W^\flat)$, $u=u_1$, the characters
$(1,\upsilon^m)$ \textup(Lemma~\ref{psv:split}\textup) and $\mathcal Q$ of Theorem~\ref{pss:b} \textup(Corollary~\ref{pss:bcor};
characters: Theorem~\ref{pss:a}\textup), and in Lemma~\ref{psp:step2} to $(W,V^\circ)$, $u=u^\circ$, the characters of
$\omega^\circ$ \textup(Theorem~\ref{pss:a}\textup), $\mathcal Q=\mathcal Q^\circ$ \textup(Theorem~\ref{psv:c}\textup) and the datum of
$K^\circ$ \textup(Proposition~\ref{psv:datum}\textup); in both cases the target is positive definite of rank more than
twice that of the source ($6m+1>6m$; $n^\circ>2r$);
\item[(P)] Theorem~\ref{psq:peaks}.
\end{enumerate}
Consequently Proposition~\ref{psp:main}, that is Proposition~\ref{po:span}, holds, modulo the published results
listed below.
\end{theorem}
\begin{proof}
The list is the content of the cited statements. In Lemma~\ref{psp:step2} the kernel $\theta_{\alpha^\circ,\beta^\circ}(x,h_p;Q_i,
\delta_\iota)$ of Theorem~\ref{psr:R}, for $Q_i=\pop^\circ_{u_i}$ with $(u_i)$ a basis of $W^\circ$, is
$\langle u_i,K^\circ(x,h_p)\rangle$ (Proposition~\ref{psv:datum}), so the right side of Theorem~\ref{psr:R} is
$\sum_p\int y\,F\,\langle\zeta_p,K^\circ(\cdot,h_p)\rangle$ for vectors $\zeta_p\in W^\circ$, as written there; the classes
$h_p$ of $\auU(V^\circ)$ modulo $J_{V^\circ}(\pi_M^2)$ are the classes modulo $J^\circ$ of \S\ref{po:ss-glue}. The proof
of Proposition~\ref{psp:main} in \S\ref{psp:ss-proof} uses exactly these inputs, the weights $u$ and $u^\circ$
entering through (H4) and Proposition~\ref{psp:isol}; with the order $o$ fixed on $W$, $\pspm_{u^\circ}$ and the
element $\mathrm{rot}(\theta)$ are taken for this order, which changes nothing in \S\ref{psp:ss-isol}--\S\ref{psp:ss-proof}.
\end{proof}

\emph{Published results used} (with the form used):
\begin{itemize}
\item Howe~\cite[(2.2)--(2.3)]{HoweJAMS}: polynomial Fock vectors are smooth vectors of the archimedean Weil
representation, and the derived action of the complexified Lie algebras maps polynomials to polynomials
(\S\ref{pss:ss-tool}, Theorems~\ref{psv:c}, \ref{psv:F}).
\item The Stone--von Neumann theorem (irreducibility of the Heisenberg representation and uniqueness of the
intertwiner up to a scalar) and G\aa rding's theorem (orbit maps of smooth vectors are smooth)
(\S\ref{psp:ss-fock}, \S\ref{pss:sec}).
\item Banaszczyk's transference theorem~\cite[Theorem~2.1]{Banaszczyk}: $\lambda_k(L)\lambda_1(L^\vee)\le k$ in rank $k$
(Lemma~\ref{psp:hom}).
\item Weil's identity~\cite{WeilSiegel} for absolutely convergent theta integrals with a definite target and
$M'>2n$, and the moderate growth of absolutely convergent Siegel Eisenstein series (Lemma~\ref{psr:SW}).
\item Schmid's theorem (D4) of {\cite[\S4.2.3]{P53}}: the $K$-types of $\pi_\lambda$ have highest weights
$\Lambda^K(\lambda)+\sum n_\beta\beta$, $n_\beta\ge0$, over positive non-compact roots (Lemma~\ref{psv:v}).
\item The inputs of {\cite[Theorem~3.2]{P49}} through Lemma~\ref{plf:A12} (the decay bounds proved in the paper and
{\cite[Lemma~4.7]{P51}}), for the integrability of the invariant derivatives of $c_W$ (Theorem~\ref{psv:ybar}).
\item The continuity of the Weil representation on the Schwartz space and of the theta functional on it (used as in
Proposition~\ref{pss:carrier}, Theorems~\ref{psv:c}, \ref{psv:F}).
\end{itemize}
Borel's finiteness theorem is not needed here: both targets $W^\flat$ and $V^\circ$ are standard spaces, for which
the finiteness of the classes is {\cite[Proposition~7.20]{P49}}.

\begin{remark}[(G2) of Proposition~\ref{po:glue}]\label{psv:G2}
The proof of (G2) in \S\ref{po:ss-glue} reads: if $\Theta^\circ y=0$ then, $c_i\ne0$, the hypothesis of
Proposition~\ref{po:span} holds, so $y=0$. Indeed $(\Theta^\circ y)_i=M^{l_1}c_i\,\#\Gamma^1\int y\,F_i\,K^\circ(\cdot,\cdot)$
with $M^{l_1}\#\Gamma^1\ne0$ and $c_i=c_{\mathsf g_i}\ne0$ (Lemma~\ref{po:slice}), so $\Theta^\circ y=0$ is the vanishing of
$\int y\,F_iK^\circ(\cdot,h^\circ)$ for all $i$ and all classes $h^\circ$ modulo $J^\circ$, the hypothesis of
Proposition~\ref{psp:main}; Theorem~\ref{psv:main} gives $y=0$. So $\Theta^\circ$ is injective. If $c'_\vartheta=0$ then
$\Theta^\circ=0$ by (G1), against $e_0\poA(W_j)\ne0$; so $c'_\vartheta\ne0$, and $\Theta_\vartheta$ is injective on $e_0\poA(W_j)$
by (G1). The remaining clause $C_{\mathrm x}>0$ follows from $C_{\mathrm x}=|c'_\vartheta|^2H_j(\vartheta)$
(Proposition~\ref{po:exchange}) and the non-vanishing of the partial $L$-values in \eqref{po:eq-H}, which is
Lemma~\ref{po:Hpos} (absolute convergence for $z\ge2$, Hecke's theorem on the line $\mathrm{Re}\,s=1$ for $z=1$).
\end{remark}



\providecommand{\auO}{\mathcal O}
\providecommand{\auU}{\mathrm U}
\providecommand{\auGL}{\mathrm{GL}}
\providecommand{\auHom}{\operatorname{Hom}}
\providecommand{\auA}{\mathbb A}
\providecommand{\auF}{\mathbb F}
\providecommand{\auPl}{\Phi}
\providecommand{\auNm}{\operatorname{Nm}}
\providecommand{\auAut}{\operatorname{Aut}}
\providecommand{\aubar}[1]{\overline{#1}}
\providecommand{\pao}{\mathfrak o}
\providecommand{\plone}{\mathbf 1}
\providecommand{\plgb}[2]{\|#2\|_{#1}}
\providecommand{\poA}{\mathcal A}
\providecommand{\poT}{\mathbb T}
\providecommand{\poAA}{A_\ast}
\providecommand{\pop}{\mathsf p}
\providecommand{\plfK}{\mathsf K}
\providecommand{\pteSD}{\mathrm{SD}}
\providecommand{\pteCS}{\mathop{\textstyle\sum^{\mathrm{cl}}}}
\providecommand{\pfqC}{\mathsf C}
\providecommand{\pfqW}{W^{\mathrm{add}}}
\providecommand{\pfqm}{\mathfrak m_\vartheta}

\section{The first factors and the tame parts at the chosen places}\label{pfq:sec}

This section treats two points of the proof in \S\ref{po:sec}: the first factors and the tame parts
of the characters at the chosen places.

\emph{First factors} (\S\S\ref{pfq:ss-hook}--\ref{pfq:ss-ff}). The exchange identity,
Proposition~\ref{po:exchange}, is proved in \S\S\ref{pte:sec}--\ref{ptf:sec} for the first factors of
Definition~\ref{ptf:def-ff}. We prove (I1) and (I2) for an
explicit non-zero constant multiple of these first factors
(Theorem~\ref{pfq:ff}), and check that Proposition~\ref{po:exchange} holds for this multiple
(Proposition~\ref{pfq:exchange}). This proves Lemma~\ref{po:ff}.

\emph{Tame parts} (\S\ref{pfq:ss-tame}). We define the residue characters $\omega_w$ at the chosen
places, the characters $\chi_M$ ($=\chi_0$ of \S\ref{po:ss-hecke}), $\mu^\flat$ and $\mu^0_j$, prove the
existence lemma for characters of $\auU(1)$ with prescribed weight and prescribed tame parts
(Lemma~\ref{pfq:mu}; compare {\cite[Lemma~8.6]{P49}}), and check that the tame parts have the
properties used in \S\ref{po:sec} and in Lemma~\ref{ppe:three}.

Throughout, $j$ ($1\le j\le m$) is fixed, and we use the notation of \S\ref{po:ss-glue},
\S\ref{pte:ss-third}, \S\ref{pte:ss-left} and \S\ref{ptf:ss-ff}: $V''=I_{N''}$, $W=W_j$ with $r_+$
positive basis vectors, $\pfqm=\pi_M^2\prod_{w\in Q_\vartheta}w$, the characters
$\chi'=\chi_M^{\,j}\vartheta$ and $\nu''=\nu''_\vartheta=\mu^0_j\vartheta\,\upsilon^{j+\lfloor(t+k-1)/2\rfloor}$, the
choices $i=(\Lambda_i,h^\flat_i,\varrho_i)$, and the space
\[
\pfqC=I_{3m}\oplus(E^j)^-\oplus E^{3m-j}\oplus W^\flat
\]
of \S\ref{ptf:ss-ff}. We write $\#\iota(\tau)$ for the number of columns of $\mu^+_\tau$ and
$\mathsf q=\sum_\tau\#\iota(\tau)\le4fb$. By the proof of Proposition~\ref{pte:third},
\begin{equation}\label{pfq:eq-weight}
\nu''\ \text{has weight}\ -\#\iota(\tau)\ \text{at}\ \tau ,
\end{equation}
since $\mu^0_j$ has weight $-\kappa_j$, $\upsilon$ weight $1$, $\vartheta$ weight $0$, and
$\#\iota(\tau)+j+\lfloor\tfrac{t+k-1}2\rfloor=\kappa_j(\tau)$. Kernels are those of {\cite[\S3.1.1]{P53}}:
$\theta^{S,T}_{\alpha,\beta}(g,h;[P,\Psi])$, the first index being the character of the source group.
The kernels $\plfK^{(L,N)}_{a,b}[\Psi,Q]$ of \S\ref{plf:ss-KS} are the kernels
$\theta^{L,N}_{a,b}(g,h;[Q,\Psi])$ of {\cite[\S2.3]{P53}} at finite points $g,h$ (archimedean components
$1$), with the same convention for the indices (the characters twist the Weil representation by
$a(\det g)^{-1}b(\det h)$); this is how Lemma~\ref{plf:KS} is used below. The identification includes
the primes $w$ of $Q_\vartheta$ at which $\alpha_w$ is not a unit (primes above $q_1$ dividing $\omega$,
{\cite[Proposition~6.9]{P53}}; the definition of a twist does not exclude them): there the datum of
{\cite[\S2.3]{P53}} is $\Psi^\natural$ of {\cite[Definition~6.31]{P53}}, supported on $a_wM(\auO_w)$, and
the datum $f_\Psi$ of \S\ref{plf:ss-KS} is the same function of the frame coordinate $x$ (it is defined
with the same lattice $L_w=a_wM(\auO_w)$ and the same generator $a_w$, that of {\cite[\S5.1.1]{P49}}),
equivalently the same function $\plone[T\ \text{integral}]\Psi(T\bmod\pcm)$ of the coordinate $T$ of (N5)
({\cite[Lemma~6.32(c)]{P53}}); so $f_\Psi=\Psi^\natural$ literally. The proof of
(KS1) and (KS2) covers this case (the paragraph after \eqref{plf:eq-vacexp}). So (I1) and (I2) below
need no restriction on $Q_\vartheta$ beyond the definition of a twist and $\ell\nmid\auNm Q_\vartheta$.

\subsection{A complementary hook number}\label{pfq:ss-hook}

We use the tensor model of {\cite[\S5.2.6]{P49}}: for a family $\mu=(\mu_\tau)$ of diagrams for $\auGL_N$,
$\poT_\mu(R)=\bigotimes_{\tau}\bigotimes_{\mathfrak c}\Lambda^{h_{\mathfrak c}}R^N$ with the orthonormal basis
$(e_U)$ indexed by the fillings, $\mathsf C_\mu\subseteq\poT_\mu(\C)$ the span of the vectors
$\mathrm{ev}_y$, $\pi_\mu$ the orthogonal projection onto $\mathsf C_\mu$, and $W_\mu(\Z)=\mathsf m(\mathsf
C_\mu\cap\poT_\mu(\Z))$ ({\cite[(51)]{P49}}); $\mathsf m$ identifies $W_\mu$ isometrically with
$\mathsf C_\mu$ (proof of Lemma~\ref{po:frame}). Since $\mathsf C_\mu\cap\poT_\mu(\Z)$ is the intersection of
$\poT_\mu(\Z)$ with a subspace, it is a direct summand of $\poT_\mu(\Z)$; hence for every valuation
ring $\pao\supseteq\Z$ of $\C$, $\mathsf C_\mu\cap\poT_\mu(\pao)=(\mathsf C_\mu\cap\poT_\mu(\Z))\otimes\pao$,
and a vector of $W_\mu$ has all coordinates in the basis of $W_\mu(\Z)$ in $\pao$ if and only if all its
tensor coordinates are in $\pao$. (This is the remark ``the two unit balls coincide'' of
\S\ref{po:ss-set}.)

\begin{lemma}[Complementary hook number]\label{pfq:hook}
Let $\mu=(\mu_\tau)$ be diagrams for $\auGL_N$ all of whose columns have heights $h_{\mathfrak c}<N$, let
$\mu^{\mathrm c}_\tau$ be the diagram whose columns have the heights $N-h_{\mathfrak c}$, and
$H^{\mathrm c}=\prod_\tau H_{\mu^{\mathrm c}_\tau}$ ({\cite[Lemma~5.22]{P49}}). Then
$H^{\mathrm c}\pi_\mu(e_U)\in\mathsf C_\mu\cap\poT_\mu(\Z)$ for every filling $U$. Consequently
$H^{\mathrm c}\pi_\mu(\poT_\mu(R))\subseteq\mathsf C_\mu\cap\poT_\mu(R)$ for every subring $R\subseteq\C$.
Moreover $H^{\mathrm c}\le |\mu^{\mathrm c}|^{3|\mu^{\mathrm c}|}=e^{3e}$, where the base is the integer $e=|\mu^{\mathrm c}|=\sum_\tau\sum_{\mathfrak c}(N-h_{\mathfrak c})$.
\end{lemma}
\begin{proof}
All objects are tensor products over $\tau$ and over the columns; it suffices to treat one $\tau$. For
$0\le h\le N$ let $\star\colon\Lambda^h\Z^N\to\Lambda^{N-h}\Z^N$ be the map with $v\wedge w=(\star v,w)\,
e_1\wedge\dots\wedge e_N$ for all $w\in\Lambda^{N-h}$, $(\,,)$ the symmetric bilinear form with orthonormal
basis $(e_R)$; thus $\star e_R=\pm e_{R^c}$, and $\star$ is an isometric bijection of lattices. For
$g\in\auGL_N(\C)$, $(gv)\wedge(gw)=\det g\,(v\wedge w)$ and $(g\,x,w)=(x,g^{\mathsf t}w)$ on exterior powers, so
$(g^{\mathsf t}\star(gv),w)=\det g\,(\star v,w)$ for all $w$, that is, $\star(gv)=\det g\cdot g^{-\mathsf t}\star v$.
Let $\mathrm{St}=\bigotimes_{\mathfrak c}\star\colon\poT_\mu\to\poT_{\mu^{\mathrm c}}$; it maps the orthonormal
basis to $\pm$ the orthonormal basis, is real, and satisfies $\mathrm{St}(gv)=(\det g)^{c}g^{-\mathsf t}\mathrm{St}(v)$,
$c$ the number of columns.

For $y\in M_N(\C)$, $\mathrm{ev}_y=y\cdot e_{U_0}$ with $e_{U_0}=\bigotimes_{\mathfrak c}e_1\wedge\dots\wedge e_{h_{\mathfrak c}}$
(functoriality of $\Lambda^h$); since $\auGL_N(\C)$ is dense in $M_N(\C)$ and $\mathsf C_\mu$ is finite-dimensional,
$\mathsf C_\mu$ is the span of the $g\,e_{U_0}$, $g\in\auGL_N(\C)$. Let $w_0$ be the permutation matrix of
$i\mapsto N+1-i$. Then $\star(e_1\wedge\dots\wedge e_h)=\pm e_{h+1}\wedge\dots\wedge e_N=\pm w_0(e_1\wedge\dots\wedge e_{N-h})$,
so $\mathrm{St}(e_{U_0})=\pm w_0e_{U_0^{\mathrm c}}$, $U_0^{\mathrm c}$ the leading filling of $\mu^{\mathrm c}$. Hence
$\mathrm{St}(\mathsf C_\mu)$ is the span of the $(\det g)^cg^{-\mathsf t}w_0e_{U^{\mathrm c}_0}$, which is
$\mathsf C_{\mu^{\mathrm c}}$. As $\mathrm{St}$ is unitary, $\mathrm{St}\circ\pi_\mu=\pi_{\mu^{\mathrm c}}\circ\mathrm{St}$.
By {\cite[Lemma~5.22]{P49}} for $\mu^{\mathrm c}$ and {\cite[(51)]{P49}}, $H^{\mathrm c}\pi_{\mu^{\mathrm c}}(e_{U'})\in
\mathsf C_{\mu^{\mathrm c}}\cap\poT_{\mu^{\mathrm c}}(\Z)$ for every filling $U'$ (this is the argument of the
proof of Lemma~\ref{po:frame}). With $\mathrm{St}(e_U)=\pm e_{U'}$,
$H^{\mathrm c}\pi_\mu(e_U)=\pm\mathrm{St}^{-1}(H^{\mathrm c}\pi_{\mu^{\mathrm c}}(e_{U'}))\in\mathsf C_\mu\cap\poT_\mu(\Z)$.
The second assertion follows by linearity, the bound from {\cite[Lemma~5.22]{P49}} ($H_\mu\le|\mu|^{3|\mu|}$,
and $\prod_\tau|\mu^{\mathrm c}_\tau|^{3|\mu^{\mathrm c}_\tau|}\le |\mu^{\mathrm c}|^{3|\mu^{\mathrm c}|}$).
\end{proof}

\subsection{The first factors as a projection}\label{pfq:ss-proj}

\emph{Fillings and Laplace coefficients.} For $Y\in\auGL_{6m+1}(E)$ let $\eta^Y_i=(\eta^Y_{i,\tau k})$ be the
constant rows of Lemma~\ref{ptf:H}, built with the short vectors given by the columns of $Y$ (the rows
of $\eta^Y_{i,\tau k}$ for a column of height $r_+$ are the transposes of the columns of $\tau(Y)$ selected by
$\varrho_i$, placed in the columns of $W^\flat$); for $Y=Y_i:=Y_{h^\flat_i}$ of {\cite[Lemma~14.21]{P46}}
these are the rows $\eta_i$ of \S\ref{pte:ss-third}. A \emph{filling} $U=(U_{\tau k})$ chooses, for each $\tau$
and each $k\in\iota(\tau)$, a subset $U_{\tau k}\subseteq\{1,\dots,N''\}$ of $h_k$ rows; put
$b_U(X)=\prod_{\tau,k}\det X_{\tau;U_{\tau k},[h_k]}$ (the product of minors of {\cite[\S5.2.6]{P49}}) and
$Q^Y_U(Z)=\prod_{\tau,k}Q_{\tau k,U_{\tau k}}(Z_\tau)$, with $Q_{\tau k,R}$ the polynomials of
Lemma~\ref{ptf:laplace} for the rows $\eta^Y_i$. Let $\pfqW=I_{3m}\oplus E^{3m-j}\oplus W^\flat$, the
positive definite summand of $\pfqC$ (so $\pfqC=(E^j)^-\oplus\pfqW$ after reordering the frame); the
variables $Z$ are the Fock variables of $\auHom(\pfqW,V'')$.

\begin{lemma}[Laplace coefficients]\label{pfq:Q}
\begin{enumerate}
\item[(Q1)] $\pteSD(\eta^Y_i)=\sum_Ub_U(X)\,Q^Y_U(Z)$.
\item[(Q2)] Every $Q^Y_U$ is homogeneous in $Z$ of degree $e=\sum_\tau\sum_{k\in\iota(\tau)}(N''-h_k)$, and
$e\le(10m+1)\mathsf q\le M$.
\item[(Q3)] $Q^Y_U$ is a polynomial whose coefficients are polynomials with integer coefficients, independent
of $Y$, in the numbers $\tau(Y_{ab})$ ($\tau\in\auPl$, $1\le a,b\le6m+1$).
\item[(Q4)] If $|\tau'(Y_{ab})|\le M^4$ for all embeddings $\tau'\colon E\to\C$ and all $a,b$, and all entries of
$z$ have absolute value at most $M^4$, then $\sum_U|Q^Y_U(z)|^2\le M^{9(12m+1)\mathsf q}$; the same holds for the
polynomials obtained from $Q^Y_U$ by applying any $\sigma\in\auAut(\C)$ to the coefficients.
\item[(Q5)] Let $\gamma\in\auU(W^\flat)(E^+)$ and $B=\mathrm{diag}(1_{3m},1_{3m-j},\gamma)\in\auU(\pfqW)(E^+)$. Then
$Q^Y_U(Z_\tau\,\tau(B)^{-\mathsf t})_\tau=\prod_\tau\tau(\det\gamma)^{-\#\iota(\tau)}\,Q^{\gamma Y}_U(Z)$. For
$\gamma'\in\auU(I_{3m})(E^+)$ and $B'=\mathrm{diag}(\gamma',1,1)$,
$Q^Y_U(Z_\tau\,\tau(B')^{-\mathsf t})_\tau=\prod_\tau\tau(\det\gamma')^{-\#\iota(\tau)}\,Q^Y_U(Z)$.
\end{enumerate}
\end{lemma}
\begin{proof}
(Q1) is the expansion in the proof of Lemma~\ref{ptf:laplace} (multiplied over $\tau$ and $k$). By that proof,
$Q_{\tau k,R}(Z_\tau)=\pm\det G_{\tau k,R}$, where $G_{\tau k,R}$ is the square matrix of size $12m+1-j$ whose
rows are the $2m-(r_+-h_k)$ constant rows of $\eta^Y_{i,\tau k}$ vanishing on the columns of $X$ (restricted to the
columns of $\pfqW$) and the $N''-h_k$ rows of $Z_\tau$ outside $R$; the sign does not depend on $Y$ or $Z$.
(Q2) Each $\det G_{\tau k,R}$ is homogeneous of degree $N''-h_k$ in $Z$. By Lemma~\ref{ptf:H} and its proof,
every column of $\iota(\tau)$ has height $h_k\ge r-m$ (one less than the parameter $r+1-m$ of that lemma; the subscript $k$ here is a column index), and $N''-(r-m)=10m+1$; so $e\le(10m+1)\mathsf q$, and $(10m+1)\mathsf q\le44fmb\le M$ by
{\cite[(36) and~(38)]{P46}} and $M\ge b^{20}$ ({\cite[Lemma~13.4]{P46}}). (Q3) The constant
rows have entries $0$, $1$ and $\tau(Y_{ab})$. (Q4) For fixed $\tau,k$, the numbers $\det G_{\tau k,R}$, $R$ running
over the $h_k$-subsets, are some of the maximal minors of the matrix $G_{\tau k}$ formed by the constant rows
vanishing on $X$ and all rows of $Z_\tau$; by the Cauchy--Binet formula and Hadamard's inequality,
$\sum_R|\det G_{\tau k,R}|^2\le\det(G^*_{\tau k}G_{\tau k})\le\prod(\text{squared lengths of the columns of }G_{\tau k})$.
There are $12m+1-j\le12m+1$ columns, each with at most $2m+N''\le M$ entries of absolute value at most $M^4$, so
the product is at most $M^{9(12m+1)}$. Since $\sum_U|Q^Y_U(z)|^2=\prod_{\tau,k}\sum_R|\det G_{\tau k,R}|^2$, the
bound follows. Applying $\sigma$ to the coefficients replaces $\tau(Y_{ab})$ by $\sigma\tau(Y_{ab})$ (by (Q3)), which
satisfies the same bound. (Q5) Multiplying on the right by $\tau(B)^{-\mathsf t}$ the rows of $Z_\tau$ in $G_{\tau k,R}$
gives the matrix $\bigl(\begin{smallmatrix}c\,\tau(B)^{\mathsf t}\\ Z_{R^c}\end{smallmatrix}\bigr)\tau(B)^{-\mathsf t}$, $c$ the
constant rows; its determinant is $\tau(\det\gamma)^{-1}\det\bigl(\begin{smallmatrix}c\,\tau(B)^{\mathsf t}\\ Z_{R^c}
\end{smallmatrix}\bigr)$. The constant rows vanish on $I_{3m}$, are basis vectors on $E^{3m-j}$ (where $B=1$), and on
$W^\flat$ are rows ${}^ty$, $y$ a column of $\tau(Y)$, which become ${}^t(\tau(\gamma)y)$, the corresponding column of
$\tau(\gamma Y)$. So $c\,\tau(B)^{\mathsf t}$ are the constant rows for $\gamma Y$. For $B'$ the constant rows vanish on
$I_{3m}$, so $c\,\tau(B')^{\mathsf t}=c$. Multiply over $(\tau,k)$.
\end{proof}

\begin{lemma}[Factorisation]\label{pfq:factor}
Let $\Psi^{(j)}$ be the indicator of the matrix of the inclusion $E^j\to V''$ (on $N''\times j$ matrices modulo
$\pfqm$), and $\Psi^{\mathrm{add}}_i$ the function on $N''\times(12m+1-j)$ matrices modulo $\pfqm$ which is
$\psi^{\mathrm{inv}}_{\Lambda_i}(x')$ if the columns on $E^{3m-j}\oplus W^\flat$ are those of the inclusion and $0$ otherwise
($x'$ the columns on $I_{3m}$). Then $\Phi^{\mathsf C}_i(x_{(E^j)^-}\mid x_{\pfqW})=\Psi^{(j)}(x_{(E^j)^-})
\Psi^{\mathrm{add}}_i(x_{\pfqW})$, and for every polynomial $R$ in $Z$ and $h'\in\auU(I_{3m})(\auA_f)$,
\[
\theta^{\pfqC,V''}_{\nu'',\nu''\upsilon^j}\bigl(c_i(h'),h;[R,\Phi^{\pfqC}_i]\bigr)=\epsilon_j^{-1}\,\aubar{k_1(h)}\cdot
\theta^{\pfqW,V''}_{\nu'',\nu''}\bigl((h',1,h^\flat_i),h;[R,\Psi^{\mathrm{add}}_i]\bigr),
\]
where $k_1(h)=\theta^{E^j,V''}_{1,1}(1,h;[1,\Psi^{(j)}])$ and $\epsilon_j=\epsilon(E^j,V'')$ is the number of modulus one of
Lemma~\ref{ptf:conj}; it depends only on $j$.
\end{lemma}
\begin{proof}
The product decomposition of $\Phi^{\mathsf C}_i$ is its definition (Definition~\ref{ptf:def-ff}). Apply
Lemma~\ref{ptf:sum} with $S_1=(E^j)^-$, $S_2=\pfqW$ (the order of the summands of a frame is immaterial: a
permutation of the frame of $\pfqC$ is an isometry of standard lattices and changes kernels, data and polynomials
compatibly), $\alpha=\nu''$, $\beta_1=\upsilon^j$, $\beta_2=\nu''$, $P_1=1$, $P_2=R$ (the polynomial $R$ involves only
the positive columns of $\pfqC$, which are those of $\pfqW$), at the source point $c_i(h')=(1,(h',1,h^\flat_i))$:
\[
\theta^{\pfqC,V''}_{\nu'',\nu''\upsilon^j}(c_i(h'),h;[R,\Phi^{\mathsf C}_i])=\theta^{(E^j)^-,V''}_{\nu'',\upsilon^j}(1,h;[1,
\Psi^{(j)}])\,\theta^{\pfqW,V''}_{\nu'',\nu''}((h',1,h^\flat_i),h;[R,\Psi^{\mathrm{add}}_i]).
\]
By Lemma~\ref{ptf:conj} with $S=E^j$, $T=V''$, $\alpha=\nu''^{-1}\upsilon^{N''}$, $\beta=1$ (so that
$\alpha^{-1}\upsilon^{N''}=\nu''$ and $\beta^{-1}\upsilon^{j}=\upsilon^j$), and as $\Psi^{(j)}$ is real,
$\aubar{\theta^{E^j,V''}_{\alpha,1}(1,h;[1,\Psi^{(j)}])}=\epsilon_j\,\theta^{(E^j)^-,V''}_{\nu'',\upsilon^j}(1,h;[1,\Psi^{(j)}])$;
and $\theta^{E^j,V''}_{\alpha,1}(1,h;\cdot)=\theta^{E^j,V''}_{1,1}(1,h;\cdot)$ by \eqref{pte:eq-twist}, the source point being $1$.
\end{proof}

\begin{lemma}[Moving the class of $W^\flat$]\label{pfq:move}
Let $x\in\auU(\pfqW)(\auA_f)$, $\gamma\in\auU(W^\flat)(E^+)$, $B=\mathrm{diag}(1,1,\gamma)$, and $\Psi$ any datum. Then
\[
\theta^{\pfqW,V''}_{\nu'',\nu''}(x,h;[Q^Y_U,\Psi])=\theta^{\pfqW,V''}_{\nu'',\nu''}(B_fx,h;[Q^{\gamma Y}_U,\Psi]),
\]
$B_f$ the finite part of $B$. The same holds with $B'=\mathrm{diag}(\gamma',1,1)$, $\gamma'\in\auU(I_{3m})(E^+)$, and
$Q^Y_U$ on both sides.
\end{lemma}
\begin{proof}
Let $\delta$ be a rational point of $\auU(\pfqW)$. The operator $\omega_{\nu'',\nu''}(\delta,1)$ leaves $\theta$ invariant
((W1) and $\nu''(\det\delta)=1$, $\nu''$ being automorphic), and $\delta=\delta_\infty\delta_f$, with $\delta_\infty$ commuting with
finite points; hence $\theta^{\pfqW,V''}_{\nu'',\nu''}(\delta_fx,h;[P,\Psi])=\theta^{\pfqW,V''}_{\nu'',\nu''}(x,h;[\omega_\infty
(\delta_\infty^{-1},1)P,\Psi])$. By Lemma~\ref{ptf:fock} (all product signs positive, $g=\delta_\infty^{-1}$, $h=1$, source weight
$-\#\iota(\tau)$ by \eqref{pfq:eq-weight}), $\omega_\infty(\delta_\infty^{-1},1)P=\prod_\tau\tau(\det\delta)^{-\#\iota(\tau)}\,
P(Z_\tau\tau(\delta)^{\mathsf t})_\tau$. Take $\delta=B$ and $P=Q^Y_U$. Since $\tau(B)^{\mathsf t}=\tau(B^{-1})^{-\mathsf t}$, (Q5) for the
element $\gamma^{-1}$ gives $Q^Y_U(Z_\tau\tau(B)^{\mathsf t})=\prod_\tau\tau(\det\gamma)^{\#\iota(\tau)}Q^{\gamma^{-1}Y}_U(Z)$, and
$\det B=\det\gamma$; the two scalars cancel, so $\theta(B_fx,h;[Q^Y_U,\Psi])=\theta(x,h;[Q^{\gamma^{-1}Y}_U,\Psi])$ for all $x$ and
$Y$. Replacing $Y$ by $\gamma Y$ gives the claim. The case of $B'$ is the same, with the second part of (Q5).
\end{proof}

\begin{lemma}[The first factors as a projection]\label{pfq:proj}
For $h\in\auU(V'')(\auA_f)$ and a filling $U$ put
\[
a_U(h)=\pteCS_{K_{I_{3m}}}\Bigl(h'\mapsto\mathrm a_{\chi'}(h')\,\theta^{\pfqC,V''}_{\nu'',\nu''\upsilon^j}\bigl(c_i(h'),h;[Q^{Y_i}_U,
\Phi^{\mathsf C}_i]\bigr)\Bigr),\qquad a(h)=\sum_Ua_U(h)\,e_U\in\poT_{\mu^+}(\C^{N''}).
\]
Then the first factor of Definition~\ref{ptf:def-ff} is $\varphi''_{\vartheta,i}(h)=\pi_{\mu^+}\bigl(\aubar{a(h)}\bigr)$, where the bar
conjugates the tensor coordinates and $W''=W_{\mu^+}(\C^{N''})$ is identified with $\mathsf C_{\mu^+}$ as in
\S\ref{pfq:ss-hook}.
\end{lemma}
\begin{proof}
By the proof of Lemma~\ref{ptf:laplace}, $u\mapsto\pop''_u$ is a conjugate-linear bijection of $W''=\mathsf C_{\mu^+}$ onto the span
$\Sigma_{\mu^+}\otimes\C$ of the $b_U$, given by $\pop''_u=\mathsf m(\aubar u)$, where $\mathsf m(v)=\sum_Uv_Ub_U$ is the multiplication
map of {\cite[\S5.2.6]{P49}} (tensor coordinates $v_U$); its kernel on $\poT_{\mu^+}(\C)$ is $\mathsf C_{\mu^+}^\perp$, so
$\mathsf m(v)=\mathsf m(\pi_{\mu^+}v)=\pop''_{\pi_{\mu^+}\aubar v}$ for every $v$ ($\pi_{\mu^+}$ is defined over $\Q$ and commutes with
the conjugation of coordinates). Let $L_h$ be the linear form $R\mapsto\pteCS(h'\mapsto\mathrm a_{\chi'}(h')\theta^{\pfqC,V''}(c_i(h'),h;
[R,\Phi^{\mathsf C}_i]))$ on polynomials in $Z$. By definition $c_\lambda(h)=L_h(R_\lambda)$, and $\pop''_{\varphi''_{\vartheta,i}(h)}=\sum_\lambda
c_\lambda(h)\pop''_{u_\lambda}=(\mathrm{id}\otimes L_h)\bigl(\sum_\lambda\pop''_{u_\lambda}\otimes R_\lambda\bigr)$. By \eqref{ptf:eq-SDexp} and (Q1),
$\sum_\lambda\pop''_{u_\lambda}\otimes R_\lambda=\pteSD(\eta_i)=\sum_Ub_U\otimes Q^{Y_i}_U$, so $\pop''_{\varphi''_{\vartheta,i}(h)}=\sum_Ua_U(h)b_U=
\mathsf m(a(h))=\pop''_{\pi_{\mu^+}\aubar{a(h)}}$. Injectivity of $u\mapsto\pop''_u$ on $W''$ gives the claim.
\end{proof}

\begin{remark}[Convention]\label{pfq:rem-linear}
The map $u\mapsto\pop''_u$ is conjugate-linear: by \eqref{po:eq-poly}, $\pop_\xi(x)=\langle\xi,\mathrm{ev}(\cdot)\rangle$
with $\langle\ ,\,\rangle$ the inner product of the representation space $\poT^+\otimes\poT^-$, which, like all
inner products of representation spaces and of functions, is conjugate-linear in the first variable;
(Cv1) of {\cite[\S2.1]{P51}} concerns hermitian forms on the spaces $W_j$, $Y_j$, and states that inner
products of functions and of representation spaces follow the opposite convention; the form of
$\poT^+\otimes\poT^-$ is of this second kind, so there is no conflict with (Cv1). This is the convention of \S\ref{po:ss-glue}, of the
proof of Lemma~\ref{ptf:laplace} and of Proposition~\ref{ptf:ff}, and it is used throughout. (Had the map been
linear, Lemma~\ref{pfq:proj} would read $\varphi''_{\vartheta,i}=\pi_{\mu^+}(a(h))$ and the arguments below would
go through with $\aubar\kappa$ replaced by $\kappa$ and $\iota$, $\aubar\iota$ interchanged in the proof of (I1),
(KS1) holding at both; we do not use this.)
\end{remark}

\begin{remark}[The summands are functions on the classes]\label{pfq:rem-classes}
The function $h'\mapsto\mathrm a_{\chi'}(h')\theta^{\pfqC,V''}(c_i(h'),h;[Q^{Y_i}_U,\Phi^{\mathsf C}_i])$ is left invariant under
$\auU(I_{3m})(E^+)$: by Lemma~\ref{pfq:factor} and the second part of Lemma~\ref{pfq:move} for the kernel, and because
$\theta^{I_m,I_{3m}}_{1,\chi'}$ has the vacuum at infinity and characters of weight $0$, so that rational points act trivially
(Lemma~\ref{ptf:fock}). It is right invariant under $K_{I_{3m}}$: by {\cite[Lemma~2.9]{P53}}, $k\in K_{I_{3m}}$ multiplies the kernel by
$\nu''_{\pfqm}(\det\bar k)^{-1}$ ($\psi^{\mathrm{inv}}_{\Lambda_i}$ is invariant under right multiplication by $\auGL_{3m}$) and
$\mathrm a_{\chi'}$ by $\chi'_{\pfqm}(\det\bar k)$; the tame parts of $\nu''$ and $\chi'$ at every prime of $\pfqm$ coincide
($\omega_w^j$ at $w\in\mathsf M$ and $\psi_w$ at $w\in Q_\vartheta$, Proposition~\ref{pfq:tame}(T4)). So the class sums of
Lemma~\ref{pfq:proj} (and of Definition~\ref{ptf:def-ff}) do not depend on the representatives.
\end{remark}

\subsection{The first factors used in \S\ref{po:sec}; integrality and size}\label{pfq:ss-ff}

Let $c_{l,n_N,e}$ be the kernel scales of Lemma~\ref{plf:KS}, $n'_{\mathrm{Aut}}$ the integer of Remark~\ref{plf:rem-classes},
$e$ the degree of (Q2), and $H^{\mathrm c}$ the number of Lemma~\ref{pfq:hook} for $\mu^+$ and $N=N''$ (the columns of
$\mu^+$ have heights at most $r_+<N''$). Put
\[
\kappa_1=n'^{\,2}_{\mathrm{Aut}}\;c_{m,3m,0}\;\aubar{c_{j,N'',0}}\;\epsilon_j\;c_{12m+1-j,N'',e},\qquad\kappa=H^{\mathrm c}\,\aubar{\kappa_1}.
\]

\begin{definition}[First factors]\label{pfq:def-ff}
The first factors of \S\ref{po:ss-glue} are $\varphi^\natural_{\vartheta,i}=\kappa\,\varphi''_{\vartheta,i}$, with $\varphi''_{\vartheta,i}$ as in
Definition~\ref{ptf:def-ff}. The number $\kappa\ne0$ depends only on $j$ (not on $\vartheta$, $i$, $h$).
\end{definition}

\begin{proposition}[The exchange identity for $\varphi^\natural$]\label{pfq:exchange}
The functions $\varphi^\natural_{\vartheta,i}$ lie in $\poA_{\mu^+}(J'_{\vartheta,i},\chi''^{-1}_\vartheta)$, and Proposition~\ref{po:exchange} holds for
them, with the same functions $F_i$, with $c'_\vartheta$ replaced by $c'_\vartheta/\aubar\kappa$ and $C_{\mathrm x}$ by $C_{\mathrm x}/|\kappa|^2$.
\end{proposition}
\begin{proof}
The type and level are those of $\varphi''_{\vartheta,i}$ (Proposition~\ref{ptf:leftside}), multiplied by a constant. The theta series
$\Xi_{\vartheta,i}$ of \S\ref{po:ss-glue} is conjugate-linear in the first factor ($\theta''_{\vartheta,s}$ is conjugate-linear in the
vector), so $\Xi^\natural_{\vartheta,i}=\aubar\kappa\,\Xi_{\vartheta,i}$, and $c'_\vartheta\Xi_{\vartheta,i}=F_i$ (\S\ref{pte:ss-proof}) gives
$(c'_\vartheta/\aubar\kappa)\Xi^\natural_{\vartheta,i}=F_i$, while $|c'_\vartheta/\aubar\kappa|^2H_j(\vartheta)=C_{\mathrm x}/|\kappa|^2$. The integrability
statements are unchanged.
\end{proof}

\begin{theorem}[Lemma~\ref{po:ff}]\label{pfq:ff}
With the first factors of Definition~\ref{pfq:def-ff} and the set $I$ of choices of \S\ref{pte:ss-third} ($\#I\le e^b$ by
Lemma~\ref{plf:I}: the choices there and in Lemma~\ref{ptf:H} are the same data), for every $i\in I$, every twist $\vartheta$
and every $h\in\auU(V'')(\auA_f)$:
\begin{enumerate}
\item[(I1)] $\plgb h{\varphi^\natural_{\vartheta,i}(h)}\le1$ at every $\iota$ with $\ell(\iota)\nmid\auNm Q_\vartheta$ (the gauge of
{\cite[\S6.1]{P49}}); that is, $\varphi^\natural_{\vartheta,i}$ is integral at $\iota$;
\item[(I2)] $\|\sigma\varphi^\natural_{1,i}(h)\|^2\le\exp(\poAA)$ for every $\sigma\in\auAut(\C)$.
\end{enumerate}
\end{theorem}
\begin{proof}
\emph{Expansion.} By Lemmas~\ref{pfq:proj} and~\ref{pfq:factor}, and expanding the class sum $\mathrm a_{\chi'}$ (\S\ref{pte:ss-carrier},
a class sum over $K_{I_m}$, which for the level $K_{I_m}$ carries no index factor),
\begin{equation}\label{pfq:eq-expand}
\varphi^\natural_{\vartheta,i}(h)=H^{\mathrm c}\,\pi_{\mu^+}\Bigl(\aubar{\kappa_1a(h)}\Bigr),\qquad
\kappa_1a_U(h)=\sum_{h'}\frac{n'_{\mathrm{Aut}}}{s(h')}\,A(h')\,B(h)\,C_U(h',h),
\end{equation}
the sum over representatives $h'$ of $\Gamma\backslash\auU(I_{3m})(\auA_f)/K_{I_{3m}}$, $s(h')$ the order of the stabiliser, and
\[
A(h')=\sum_{h_1}\frac{n'_{\mathrm{Aut}}}{s(h_1)}\,c_{m,3m,0}\,\theta^{I_m,I_{3m}}_{1,\chi'}(h_1,h';[1,1]),\qquad
B(h)=\aubar{c_{j,N'',0}\,k_1(h)},
\]
\[
C_U(h',h)=c_{12m+1-j,N'',e}\;\theta^{\pfqW,V''}_{\nu'',\nu''}\bigl((h',1,h^\flat_i),h;[Q^{Y_i}_U,\Psi^{\mathrm{add}}_i]\bigr),
\]
$h_1$ running over representatives of $\Gamma\backslash\auU(I_m)(\auA_f)/K_{I_m}$. The factor $\epsilon_j$ of $\kappa_1$ cancels the
$\epsilon_j^{-1}$ of Lemma~\ref{pfq:factor}. The stabiliser orders divide $n'_{\mathrm{Aut}}$ (Remark~\ref{plf:rem-classes}), so the
coefficients $n'_{\mathrm{Aut}}/s$ are positive integers at most $n'_{\mathrm{Aut}}$. The kernels are those of Lemma~\ref{plf:KS}:
$\theta^{I_m,I_{3m}}_{1,\chi'}(h_1,h';[1,1])=\plfK^{(E^m,E^{3m})}_{1,\chi'}[1,1](h_1,h')$ (trivial modulus), $k_1(h)=\plfK^{(E^j,V'')}_{1,1}
[\Psi^{(j)},1](1,h)$ and $\theta^{\pfqW,V''}_{\nu'',\nu''}(x,h;[Q,\Psi^{\mathrm{add}}_i])=\plfK^{(\pfqW,V'')}_{\nu'',\nu''}[\Psi^{\mathrm{add}}_i,Q](x,h)$
(modulus $\pfqm$), all for standard definite spaces; the characters $1,\chi',\nu''$ have algebraic weights and tame
ramification at split places only.

\emph{(I1).} Let $\ell=\ell(\iota)\nmid\auNm Q_\vartheta$, and write $\aubar\iota=\iota\circ c$.
\emph{Reduction to $h$ integral above $\ell$.} By {\cite[Corollary~6.4(a)]{P49}} there is $\gamma\in\auU(V'')(E^+)$ with $\gamma h$
integral above $\ell$. By Proposition~\ref{pfq:exchange}, $\varphi^\natural_{\vartheta,i}(\gamma h)=\gamma\,\varphi^\natural_{\vartheta,i}(h)$, so by (G1) of
{\cite[\S6.1]{P49}}, $\plgb h{\varphi^\natural(h)}=\plgb{\gamma h}{\varphi^\natural(\gamma h)}$. So we may assume $h$ integral above $\ell$;
then $\plgb h\xi=\|\iota(\xi)\|$ by (G2), the largest absolute value of a coordinate of $\iota(\xi)$ in the basis of $W''(\Z)$.
\emph{Reduction to the coefficients.} By \S\ref{pfq:ss-hook} it suffices that all tensor coordinates of $\iota(\varphi^\natural(h))$ lie
in $\pao=\{|\cdot|_\ell\le1\}$. By Lemma~\ref{pfq:hook} with $R=\iota^{-1}(\pao)$ and \eqref{pfq:eq-expand}, it suffices that
$\aubar{\kappa_1a_U(h)}\in\iota^{-1}(\pao)$, i.e.\ $|\kappa_1a_U(h)|_{\aubar\iota}\le1$, for every filling $U$.
\emph{The terms.} By Remark~\ref{pfq:rem-classes} and {\cite[Corollary~6.4(a)]{P49}} we may take the representatives $h'$ and $h_1$
integral above $\ell$. (a) $|c_{m,3m,0}\theta^{I_m,I_{3m}}_{1,\chi'}(h_1,h';[1,1])|_{\aubar\iota}\le1$ by (KS1) with the trivial modulus
and $q=1$; so $|A(h')|_{\aubar\iota}\le1$. (b) $|c_{j,N'',0}k_1(h)|_\iota\le1$ by (KS1) with $q=\auNm Q_\vartheta$ ($qM^2\in\pfqm$, as
$\pi_M\mid M$; the datum takes the values $0,1$); hence $|B(h)|_{\aubar\iota}\le1$. (c) Choose $\gamma^\flat\in\auU(W^\flat)(E^+)$ with
$\gamma^\flat h^\flat_i$ integral above $\ell$ ({\cite[Corollary~6.4(a)]{P49}}), and put $B=\mathrm{diag}(1,1,\gamma^\flat)$. By
Lemma~\ref{pfq:move}, $C_U(h',h)=c_{12m+1-j,N'',e}\theta^{\pfqW,V''}_{\nu'',\nu''}((h',1,\gamma^\flat h^\flat_i),h;[Q^{\gamma^\flat Y_i}_U,
\Psi^{\mathrm{add}}_i])$. By {\cite[Lemma~14.21]{P46}}, $h^{\flat,-1}_iY_i$ is integral, so $\gamma^\flat Y_i=(\gamma^\flat h^\flat_i)
(h^{\flat,-1}_iY_i)$ is integral at the places above $\ell$; its entries are elements of $E$ integral at every prime above $\ell$,
so all $\tau(\gamma^\flat Y_i)_{ab}$ are integral at $\aubar\iota$, and by (Q3) the coefficients of $Q^{\gamma^\flat Y_i}_U$ are
integral at $\aubar\iota$. The values of $\Psi^{\mathrm{add}}_i$ are sums of roots of unity, $Q^{\gamma^\flat Y_i}_U$ is homogeneous of
degree $e$ (Q2), and the points $(h',1,\gamma^\flat h^\flat_i)$ and $h$ are integral above $\ell$. So (KS1) with $q=\auNm Q_\vartheta$ at
$\aubar\iota$ gives $|C_U(h',h)|_{\aubar\iota}\le1$. By \eqref{pfq:eq-expand} and the ultrametric inequality,
$|\kappa_1a_U(h)|_{\aubar\iota}\le1$.

\emph{(I2).} Let $\vartheta=1$ (so $\pfqm=\pi_M^2$), $\sigma\in\auAut(\C)$ and $\sigma'=\sigma\circ c$. The norm on $W''$ is the
restriction of the norm of $\poT_{\mu^+}$ (\S\ref{pfq:ss-hook}); $\pi_{\mu^+}$ has rational coefficients, so it commutes with
$\sigma$ applied to the coordinates, and has norm $1$. By \eqref{pfq:eq-expand},
\[
\|\sigma\varphi^\natural_{1,i}(h)\|\le H^{\mathrm c}\,\|\sigma'(\kappa_1a(h))\|\le H^{\mathrm c}\sum_{h'}n'_{\mathrm{Aut}}\,|\sigma'A(h')|\,
|\sigma'B(h)|\,\bigl\|\bigl(\sigma'C_U(h',h)\bigr)_U\bigr\| .
\]
By (KS2) with $A=B=1$ (ranks $m,3m,j,N''\le\Delta$, degree $0$, trivial modulus resp.\ $\pi_M^2$): $|\sigma'A(h')|\le n'_{\mathrm{Aut}}
e^{2x^{26}}e^{x_1}$ (Lemma~\ref{plf:classes} for the number of classes $h_1$), and $|\sigma'B(h)|=|(\sigma'c)(c_{j,N'',0}k_1(h))|\le
e^{x_1}$. For the vector $C(h',h)=(C_U(h',h))_U$ we use duality: $\|\sigma'C\|=\sup|\sum_U\phi_U\,\sigma'C_U|$ over the $\phi$ with
$\sum|\phi_U|^2\le1$, and $\sum_U\phi_U\,\sigma'C_U=\sigma'\bigl(c_{12m+1-j,N'',e}\theta^{\pfqW,V''}_{\nu'',\nu''}((h',1,h^\flat_i),h;
[Q_\phi,\Psi^{\mathrm{add}}_i])\bigr)$ with $Q_\phi=\sum_U\sigma'^{-1}(\phi_U)Q^{Y_i}_U$, homogeneous of degree $e\le M$ (Q2). Its
$\sigma'$-conjugate $Q_\phi^{\sigma'}=\sum_U\phi_U(Q^{Y_i}_U)^{\sigma'}$ satisfies, by the Cauchy--Schwarz inequality and (Q4),
$|Q^{\sigma'}_\phi(z)|\le M^{9(12m+1)\mathsf q/2}$ whenever all $|z_c|\le M^4$; the hypothesis of (Q4) holds because all conjugates of
the entries of $Y_i$ are at most $C_{\mathrm{sb}}\le x^{103}\le M^4$ ({\cite[Lemma~14.21]{P46}}; the bound on $C_{\mathrm{sb}}$ is that of
the proof of Lemma~\ref{plf:classes}, rank $6m+1\le2n_m$). The values of $\Psi^{\mathrm{add}}_i$ and their conjugates are at most
$\#\auGL_{3m}(\auO_E/\pi_M^2)\le M^{18fm^2}$. Since $12m+1-j,N''\le\Delta$ and $M^2\in\pi_M^2\auO_E$, (KS2) gives
$\|\sigma'C(h',h)\|\le e^{x_1}M^{18fm^2}M^{9(12m+1)\mathsf q/2}$. With at most $e^{2x^{26}}$ classes $h'$ (Lemma~\ref{plf:classes}),
\[
\|\sigma\varphi^\natural_{1,i}(h)\|\le H^{\mathrm c}\,n'^{\,2}_{\mathrm{Aut}}\,e^{4x^{26}}\,e^{3x_1}\,M^{18fm^2}\,M^{9(12m+1)\mathsf q/2}.
\]
By Lemma~\ref{pfq:hook}, $H^{\mathrm c}\le e^{3e}$, the base being the integer $e$ of (Q2), so $\log H^{\mathrm c}\le3e\log e$. Taking logarithms and doubling: $2\log H^{\mathrm c}\le6e\log e\le6\cdot11m\cdot4fb\cdot23\log b$ (as
$e\le(10m+1)\mathsf q\le M\le b^{22}$), $4\log n'_{\mathrm{Aut}}\le4x^{40}$, $8x^{26}+6x_1\le7x_1$, $36fm^2\log M\le x^{14}$, and
$9(12m+1)\mathsf q\log M\le9\cdot13m\cdot4fb\cdot23\log b$. The total is at most $1.7\cdot10^4fmb\log b+8x_1\le fmb(\log b)^2=\poAA$, since
$\log b\ge2\cdot10^4+1$ and $8x_1\le fmb\log b$ ({\cite[Lemma~13.4 and~(36)]{P46}}).
\end{proof}

\begin{remark}[Uses in \S\ref{po:sec}]\label{pfq:rem-uses}
In \S\ref{po:ss-glue} and the proof of {\cite[Proposition~14.41]{P46}}, the first factors enter only through: their type and level
(the summands of $\poglue_\vartheta$ do not depend on the representatives) --- Proposition~\ref{pfq:exchange}; (G1) --- the identity
$c'_\vartheta\Xi_{\vartheta,i}=F_i$, Proposition~\ref{pfq:exchange}; (G3) --- Proposition~\ref{po:periods}, valid for every vector $u''$; (G4) ---
(I1), Theorem~\ref{pfq:ff}; (G5) and Step~6 --- (I2) and $\#I\le e^b$. Lemmas~\ref{po:lattice} and~\ref{po:slice} do not involve the
first factors. The constant $C_{\mathrm x}$ is replaced by $C_{\mathrm x}/|\kappa|^2$, which changes only the constant $C_j$ of Step~1
and keeps $C_{\mathrm x}>0$ in (G2).
\end{remark}

\subsection{Tame parts at the chosen places}\label{pfq:ss-tame}

Recall from {\cite[Definition~13.3]{P46}} and {\cite[\S13.1]{P46}}: $E=K(\zeta)$, $\zeta$ a primitive $q_1$-th root of unity,
$4\mid q_1$, $q_1d\mid q_2=(q_0d)^{s+6}$ with $2^s=m^{22}$; $\kappa(\tau)\in\{b,xb/x_1\}$ and $q_1x_1\mid b$ (D4); the level primes
$p_0,p_1\equiv1\pmod{q_2^2}$; $\Delta/2=r^8$; and $\auPl$ is identified with $\operatorname{Gal}(E/K)\cong(\Z/q_1)^\times$, the
embeddings extending a fixed $K\subset\C$. For a finite place $v$ of $E^+$ split as $w\aubar w$ we identify
$\auU(1)(E^+_v)$ with $E_w^\times$ by the component at $w$.

\begin{lemma}[Roots of unity]\label{pfq:roots}
\begin{enumerate}
\item[(a)] Every root of unity of $E$ is a $q_1$-th root of unity: $\mu_E=\mu_{q_1}$.
\item[(b)] $\prod_{\tau\in\auPl}\tau(\xi)=1$ for every $\xi\in\mu_E$.
\item[(c)] $q_1$ divides $\kappa(\tau)$ for every $\tau$; hence $\prod_\tau\tau(\xi)^{c\kappa(\tau)+c'r^8}=1$ for $\xi\in\mu_E$ and all
integers $c,c'$.
\item[(d)] At every chosen place $w$, $\auF_w=\auF_{p_i}$, $q_2^2\mid\#\auF_w^\times$, and the reduction of every $\xi\in\mu_E$ is a
$q_2$-th power in $\auF_w^\times$.
\end{enumerate}
\end{lemma}
\begin{proof}
(a) $K=\Q(\sqrt{-d})\subseteq\Q(\zeta_d)$ ($d\equiv3\pmod4$ is prime), so $E\subseteq\Q(\zeta_{q_1d})$, whose roots of unity are
$\mu_{q_1d}=\mu_{q_1}\times\mu_d$ ($q_1d$ is even). If $E$ contained a root of unity of order divisible by $d$ it would contain
$\zeta_d$, hence the compositum of $\Q(\zeta_d)$ and $\Q(\zeta_{q_1})$, of degree $(d-1)\varphi(q_1)>2\varphi(q_1)=[E:\Q]$
($d\ge7$, coprime conductors); impossible. So $\mu_E\subseteq\mu_{q_1}\subseteq E$. (b) Write $\xi=\zeta^n$; the embeddings in
$\auPl$ send $\zeta$ to $\zeta_0^a$, $a\in(\Z/q_1)^\times$, for a fixed primitive $\zeta_0$, so $\prod_\tau\tau(\xi)=\zeta_0^{n\sum a}$, and
$\sum_{a\in(\Z/q_1)^\times}a\equiv0\pmod{q_1}$ (pair $a$ with $q_1-a$; $a\ne q_1-a$ because $q_1/2$ is even). (c) $q_1\mid b$ by (D4),
and $xb/x_1=x\,q_1(b/(q_1x_1))$; then $\tau(\xi)^{c\kappa(\tau)}=1$ by (a), and $\prod_\tau\tau(\xi)^{c'r^8}=1$ by (b). (d) $p_i\equiv1$
modulo $q_2^2$, hence modulo $q_1d$, so $p_i$ splits completely in $\Q(\zeta_{q_1d})\supseteq E$ and $\auF_w=\auF_{p_i}$. The group
$\auF_w^\times$ is cyclic of order $p_i-1$, divisible by $q_2^2$; an element $y$ with $y^{q_1}=1$ is a $q_2$-th power because
$y^{(p_i-1)/q_2}=1$, as $q_1\mid q_2\mid(p_i-1)/q_2$.
\end{proof}

\begin{lemma}[Characters of $\auU(1)$ with prescribed weight and tame parts]\label{pfq:mu}
Let $S$ be a finite set of primes of $E$, no two above the same place of $E^+$ and each above a place of $E^+$ split in $E$; let
$(c_\tau)_{\tau\in\auPl}$ be integers and, for $w\in S$, $\psi_w$ a character of $\auF_w^\times$, such that
\begin{equation}\label{pfq:eq-compat}
\prod_{\tau\in\auPl}\tau(\xi)^{c_\tau}\cdot\prod_{w\in S}\psi_w(\xi\bmod w)=1\qquad\text{for every }\xi\in\mu_E .
\end{equation}
Then there is a unitary character $\mu$ of $\auU(1)(E^+)\backslash\auU(1)(\auA_{E^+})$ with
\[
\mu(u)=\prod_\tau\tau(u_\tau)^{c_\tau}\cdot\prod_{w\in S}\psi_w(u_w\bmod w)\qquad\text{for all }u\in\mathsf U=\auU(1)(E^+\otimes\R)\times\prod_v
\auU(1)(\auO_v),
\]
that is: $\mu_\tau(u)=\tau(u)^{c_\tau}$ at the infinite places, $\mu$ has tame part $\psi_w$ at $w\in S$, and is unramified at all
other finite places. Conversely \eqref{pfq:eq-compat} is necessary for the existence of such $\mu$. If the reduction
$\mu_E\to\auF_{w_0}^\times$ is injective for some $w_0\in S$, then for every $(c_\tau)$ and every choice of $\psi_w$, $w\ne w_0$, there is $\psi_{w_0}$ satisfying \eqref{pfq:eq-compat}.
\end{lemma}
\begin{proof}
$\mathsf U$ is an open subgroup of $\auU(1)(\auA_{E^+})$ and $\auU(1)(E^+)\cap\mathsf U$ consists of the $\xi\in E$ with $\xi\aubar\xi=1$
that are integral at every finite place; all their conjugates have absolute value $1$, so they are roots of unity (Kronecker), and
$\auU(1)(E^+)\cap\mathsf U=\mu_E$. The formula defines a continuous unitary character $\mu_{\mathsf U}$ of $\mathsf U$ (at $w\in S$,
$\auU(1)(\auO_v)=\auO_w^\times\to\auF_w^\times$), and \eqref{pfq:eq-compat} says $\mu_{\mathsf U}(\xi)=1$ for $\xi\in\mu_E$ (a rational point
$\xi$ has components $\tau(\xi)$ at the infinite places and $\xi$ at the finite ones). Hence $\gamma u\mapsto\mu_{\mathsf U}(u)$ is a
well defined character of the open subgroup $\auU(1)(E^+)\mathsf U$, trivial on $\auU(1)(E^+)$. A unitary character of an open
subgroup of a locally compact abelian group extends to a unitary character of the group (Pontryagin duality: the restriction map
of dual groups is onto for closed subgroups); the extension is trivial on $\auU(1)(E^+)$ and is the required $\mu$. Necessity: if
$\mu$ exists, $1=\mu(\xi)=\mu_{\mathsf U}(\xi)$ for $\xi\in\mu_E$. The last assertion: $\xi\mapsto\prod_\tau\tau(\xi)^{-c_\tau}\prod_{w\ne w_0}
\psi_w(\xi)^{-1}$ is a character of the subgroup $\mu_E$ of the cyclic group $\auF_{w_0}^\times$, and extends.
\end{proof}

\begin{definition}[The characters attached to the data]\label{pfq:def-chars}
\begin{enumerate}
\item[(i)] For each chosen place $w$ fix a character $\omega_w$ of $\auF_w^\times$ of order exactly $q_2$ (it exists by
Lemma~\ref{pfq:roots}(d)); for the other finite places put $\omega_w=1$. Write $\omega=(\omega_w)_w$.
\item[(ii)] $\chi_M$ is a character of $\auU(1)$ of weight $0$ with tame part $\omega$: Lemma~\ref{pfq:mu} with $S$ the chosen places,
$c_\tau=0$, $\psi_w=\omega_w$. Condition \eqref{pfq:eq-compat} holds factor by factor: $\omega_w(\xi\bmod w)=1$ because $\xi\bmod w$ is a
$q_2$-th power (Lemma~\ref{pfq:roots}(d)) and $\omega_w^{q_2}=1$. We write $\chi_{M,w}=\omega_w$ for its tame parts.
\item[(iii)] $\mu^\flat$ is a character of $\auU(1)$ of weight $-\kappa_0=-(\kappa+\frac\Delta2)$, unramified at all finite places:
Lemma~\ref{pfq:mu} with $S=\emptyset$, $c_\tau=-\kappa(\tau)-r^8$; \eqref{pfq:eq-compat} is Lemma~\ref{pfq:roots}(c).
\item[(iv)] $\mu^0_j=\mu^\flat\,(\upsilon^{-\Delta}\chi_M)^j$ for $1\le j\le m$, $\upsilon$ restricted to $\auU(1)$.
\end{enumerate}
\end{definition}

\begin{proposition}[Properties of the tame parts]\label{pfq:tame}
\begin{enumerate}
\item[(T1)] $\mu^0_j$ has weight $-\kappa_j$, tame part $\omega_w^j$ at every chosen place $w$, and is unramified at all other finite
places; so it is a character as required in {\cite[\S14.3]{P46}}. The same characters are obtained, up to a character of weight $0$
unramified everywhere, from Lemma~\ref{pfq:mu} with $c_\tau=-\kappa_j(\tau)$ and $\psi_w=\omega_w^j$: condition \eqref{pfq:eq-compat}
holds for every $j$.
\item[(T2)] $\omega_w$ has order $q_2>8m$ at every chosen place; hence $\omega_w^{\,n}\ne1$ for $1\le|n|\le8m$, and in particular
$\omega_w^{\pm j}\ne1$ for $1\le j\le m$.
\item[(T3)] For every twist $\vartheta$, the tame part of $\mu=\mu^0_j\vartheta$ at a chosen place $v$ is $\psi_v=\omega_v^{\,j}$, and
$\psi_v^{q_2}=1$. This is the property used in (F2) of Lemma~\ref{ppe:three} (Remark~\ref{ppe:rem-T}).
\item[(T4)] The character $\chi_0$ of \S\ref{po:ss-hecke} is $\chi_M$, with tame part $\omega$ trivial off $\mathsf M$ and of order
$>8m$ at $\mathsf M$; the tame part of $\mu^0_j$ is $\omega^j$. The characters $\nu_\ast=\mu^0_j\upsilon^{n'}$, $\nu_\pm$,
$\chi'=\chi_M^j\vartheta$, $\nu''_\vartheta$ of \S\ref{po:ss-hecke}, \S\ref{pte:ss-third} have the weights and tame parts stated there,
and $\nu^\flat=\nu''_\vartheta\upsilon^m/\chi'=\mu^\flat\upsilon^{-j\Delta+j+\lfloor(t+k-1)/2\rfloor+m}$ is independent of $\vartheta$ with trivial tame part. At every
prime of $\pfqm$ the tame parts of $\nu''_\vartheta$ and $\chi'$ coincide ($\omega^j_w$ at $\mathsf M$, $\psi_w$ at $Q_\vartheta$).
\end{enumerate}
\end{proposition}
\begin{proof}
Weights add and tame parts multiply under products of characters (the defining formula on $\mathsf U$ is multiplicative), and
$\upsilon|_{\auU(1)}$ has weight $1$ and is unramified at the finite places: $\upsilon$ is unramified at every
finite place of $E$ ({\cite[\S5.2.1]{P49}}: $\upsilon(x)=\prod_\tau x_\tau/|x_\tau|$ when the finite part of $x$ is a unit), and
$\auU(1)(\auO_v)$ consists of units.
(T1) The weight of $\mu^0_j$ is $-\kappa-\frac\Delta2-j\Delta=-\kappa_j$, and its tame part is $\omega^j$. For the second statement,
\eqref{pfq:eq-compat} reads $\prod_\tau\tau(\xi)^{-\kappa(\tau)-(2j+1)r^8}\prod_w\omega_w(\xi)^j=1$, which holds by Lemma~\ref{pfq:roots}(c)
and (ii) of Definition~\ref{pfq:def-chars}; two characters with the same weight and tame parts differ by a character trivial on
$\mathsf U$. (T2) $q_2=(q_0d)^{s+6}\ge2^s=m^{22}>8m$. (T3) $\vartheta$ is unramified at the chosen places (its set $Q$ is prime to
$M$), so the tame part is that of $\mu^0_j$, and $(\omega_v^j)^{q_2}=1$. (T4) The definition of $\chi_0$ in \S\ref{po:ss-hecke} asks
exactly for weight $0$ and a tame part trivial off $\mathsf M$ of order $>8m$ on $\mathsf M$, with the tame part of $\mu^0_j$ equal to
$\omega^j$; (ii), (T1), (T2) give this with $\chi_0=\chi_M$. The other characters are products of $\mu^\flat$, $\chi_M$, $\upsilon$ and
$\vartheta$, and their weights and tame parts follow by addition and multiplication; for $\nu^\flat$ the factors $\chi_M^j\vartheta$
cancel. Finally the tame part of $\nu''_\vartheta=\mu^0_j\vartheta\upsilon^{\cdots}$ is $\omega^j\psi$, and so is that of $\chi'$.
\end{proof}

\begin{remark}[Every use of {\cite[Lemma~8.6]{P49}}]\label{pfq:rem-pzmu}
{\cite[Lemma~8.6]{P49}}, as stated in {\cite[\S8]{P49}}, is Lemma~\ref{pfq:mu} (same statement and proof; the weaker form with a
tame part prescribed at one chosen place $w_0$ only is the case $S=\{w_0\}$, $\psi_{w_0}$ given by the last assertion; $\mu_E\to\auF_{w_0}^\times$ is injective because $p_0>b^{10}$ does not divide $\#\mu_E=q_1$). {\cite[Lemma~8.6]{P49}} is used: (1) in {\cite[\S14.3]{P46}}, for the
existence of $\mu^0_j$ --- now Definition~\ref{pfq:def-chars}(iv) and (T1), which give in addition the tame parts $\omega^j$ needed
in \S\ref{po:ss-hecke} and in Lemma~\ref{ppe:three}; (2) in the summary at the end of {\cite[\S8]{P49}} (types, characters, levels) ---
unchanged; (3) in Remark~\ref{ppe:rem-T} --- now (T3); (4) in \S\ref{ptf:ss-close}, for the meaning of ``weight''
($\mu_\tau(u)=\tau(u)^{c_\tau}$) --- the same convention in Lemma~\ref{pfq:mu}. The existence of twists with prescribed tame parts
trivial on $\mu_E$ ({\cite[\S14.3]{P46}}) is Lemma~\ref{pfq:mu} with $c_\tau=0$ and $S=Q$.
\end{remark}



\section{The test form of least valuation}\label{C2:g6}
\providecommand{\auO}{\mathcal O}
\providecommand{\auU}{\mathrm U}
\providecommand{\auAut}{\operatorname{Aut}}
\providecommand{\auNm}{\operatorname{Nm}}
\providecommand{\auA}{\mathbb A}
\providecommand{\auPl}{\Phi}
\providecommand{\auPa}{\Phi_a}
\providecommand{\auPb}{\Phi_b}
\providecommand{\aubud}{\mathrm{bud}}
\providecommand{\auhdir}{h_{\mathrm{dir}}}
\providecommand{\ausinf}{s_\infty}
\providecommand{\ausfin}{s_{\mathrm{fin}}}
\providecommand{\aubar}[1]{\overline{#1}}
\providecommand{\pao}{\mathfrak o}
\providecommand{\pam}{\mathfrak m}
\providecommand{\paCp}{\C_p}
\providecommand{\popr}{\operatorname{pr}}
\providecommand{\psvac}{\mathbf v}
\providecommand{\psrho}{\varrho}
\providecommand{\ancl}{\operatorname{cl}}
\providecommand{\psmval}{\operatorname{val}}
\providecommand{\psmR}{R}
\providecommand{\psmns}{n_\sharp}
\providecommand{\psmpr}{\mathrm{pr}}
\providecommand{\psmnot}[1]{\par\smallskip{\raggedright\noindent\textbf{[Not proved here:} #1\textbf{]}\par}\smallskip}

\subsection{The test form of least valuation: sharpness of the count}\label{psm:sec}

This section concerns Proposition~\ref{po:small}. In \S\ref{psm:ss-statement} we fix its exact
statement and check that Proposition~\ref{po:mainheight} and {\cite[Proposition~14.42]{P46}} follow from
it in the direction stated. The rest of the section proves the core of the proposition, the
\emph{sharpness of the count} (part (i) of the note to Proposition~\ref{po:small}): at every prime
some Taylor coefficient of the right degrees of a theta function with complex multiplication has
valuation at most the prime sum $F(p)$, up to a factor whose conjugates are controlled uniformly
in $p$. The $p$-adic part of this (Theorem~\ref{psm:twelve}) is proved in full in
\S\S\ref{psm:ss-least}--\ref{psm:ss-twelve}; its transfer to the indicator data
(Theorem~\ref{psm:attained}) is proved in \S\ref{psm:ss-attained} from four inputs on the
Heisenberg model, proved in \S\ref{pst:sec}. Parts (ii)--(iv) of the note (the test form itself)
are proved in \S\ref{psu:sec}, which also treats Proposition~\ref{po:arch}.

\subsubsection{The exact statement}\label{psm:ss-statement}

Recall from {\cite[Definition~14.8]{P46}} the local prime sum
$s_p(y,\sigma)=v(\pi_p)\allowbreak \sum_{e\in I_p}\allowbreak \bigl(y/p^e-\kappa(u(-p^e)\sigma)\bigr)^+$, and put
\begin{equation}\label{psm:eq-F}
F(p)=\sum_{\sigma\in\auPl}s_p\bigl(\kappa(\sigma),\sigma\bigr)\qquad(p\text{ prime}),
\end{equation}
the prime sum evaluated \emph{at the weights $\kappa$}, not at the degrees of differentiation.
Thus $F(p)=0$ if $p$ splits in $K$ or $p\ge x_1$, and $\sum_p2\log p\cdot F(p)=\sum_\sigma
\ausfin(\kappa(\sigma),\sigma)$ by {\cite[Lemma~14.9(a)]{P46}}.

\begin{proposition}[Proposition~\ref{po:small}, exact form]\label{psm:small}
Let $p$ be a main pair for $j$ and $q_2$ a critical pair with empty row on a block of size $4$
\textup{(}in the application $q_2=(\upsilon^{-2}\chi_0,2)$, flat and ramified\textup{)}. There is
$c_0\in\C^\times$ with
\[
|\sigma c_0|\le\exp\bigl(fmb/k+3x^{35}\bigr)\qquad(\sigma\in\auAut(\C))
\]
such that for every prime $\ell$, every homomorphism $\jmath\colon\C\to\C_\ell$ and every twist
$\vartheta$ with $\ell\nmid\auNm Q_\vartheta$ there is a test form $\Xi$ at $\ell$ for
$(p,\vartheta)$ with
\begin{equation}\label{psm:eq-small}
\bigl|c_0\,\rho_{q_2}(\vartheta)\bigr|_\jmath\cdot\ell^{-F(\ell)}\ \le\ \bigl|\langle\Xi,
\varphi_{p,\vartheta}\rangle\bigr|_\jmath ,
\end{equation}
where $\varphi_{p,\vartheta}$ is the normalised lift and $F$ is \eqref{psm:eq-F}.
\end{proposition}

\emph{Comparison with the printed statement.} The statement printed as Proposition~\ref{po:small}
is correct as it stands: the function there, $F(\ell)=\sum_{\sigma\in\auPl}s_\ell(\kappa(\sigma),
\sigma)$, is \eqref{psm:eq-F}; the inequality is a \emph{lower} bound for the pairing (an upper
bound for its valuation); the constant $c_0$ does not depend on $\ell$, $\jmath$ or $\vartheta$; and
the twist enters only through the ratio $\rho_{q_2}(\vartheta)$ of the auxiliary pair. The only
hypothesis on $q_2$ that is used in the proof is that it has an empty row on a block of size $4$;
being ramified is used only afterwards, through (R2), in Proposition~\ref{po:mainheight}.

\emph{The deductions.} (1) Proposition~\ref{po:mainheight}: its proof uses
Proposition~\ref{psm:small} exactly in the form \eqref{psm:eq-small}, for $\jmath=\iota$ and
$\jmath=\aubar\iota$, multiplies the two inequalities into \eqref{po:eq-two}, and needs
$\sum_{\ell\in T}2F(\ell)\log\ell\le\sum_\sigma\ausfin(\kappa(\sigma),\sigma)$; this holds (with
equality when $T$ contains all primes $<x_1$ that do not split) because every term of
$\sum_p2\log p\cdot F(p)$ is $\ge0$ and the sum is $\sum_\sigma\ausfin(\kappa(\sigma),\sigma)$.
The fifth factor in that proof has multiplier $c_0\aubar{c_0}$, whose conjugates are at most
$\exp(2fmb/k+6x^{35})$, and its absolute value at $\iota$ is at most $\ell^{2F(\ell)}$; adding the
heights, the total is
\[
\bigl(\aubud-\textstyle\sum_\tau\ausinf(\kappa(\tau))\bigr)+\aubud+2x_1+2\bigl(fmb/k+3x^{35}
\bigr)+\textstyle\sum_\sigma\ausfin(\kappa(\sigma),\sigma)\ \le\ 6\,\aubud-\auhdir ,
\]
as stated. The sign is the right one: the prime sums $\ausfin$ enter with the sign $+$ in the
height of the inverse ratio, so that they cancel the $-\ausfin$ inside $-\auhdir=-\sum_\sigma
(\ausinf(\kappa(\sigma))-\ausfin(\kappa(\sigma),\sigma))$ and leave $-\sum_\sigma\ausinf
(\kappa(\sigma))$, which is cancelled by Proposition~\ref{po:arch}. A version of
Proposition~\ref{psm:small} with $F(\ell)$ replaced by a larger quantity would give a weaker bound
($6\aubud-\auhdir$ plus the excess), and with an upper bound for the pairing instead of a lower
bound the argument would not apply at all. (2) {\cite[Proposition~14.42]{P46}} is deduced from
Proposition~\ref{po:mainheight} and (R2) by subadditivity of heights; it uses nothing further from
Proposition~\ref{psm:small}. The constant $fmb/k+3x^{35}$ is bounded by
$\aubud$ before the heights are added (Lemma~\ref{psu:sizes}(N6)), which gives the same final bound.

\emph{Where the exponent $F(\ell)$ comes from.} In the proof of Proposition~\ref{psm:small} the
number $c_0$ is a product $t\cdot I_{q_2}(1)\cdot L$ (note to Proposition~\ref{po:small}), the
pairing $\langle\Xi,\varphi_{p,\vartheta}\rangle$ is, up to $\rho_{q_2}(\vartheta)$, a class number
and the two other factors, an average over characters of Taylor coefficients of a theta function in
$12$ variables, and the whole of the exponent $F(\ell)$ comes from one statement on these Taylor
coefficients, Theorem~\ref{psm:attained} below, applied with the degrees $l_\tau=c(\tau)-z$ of the
main pair (which lie in $[\kappa(\tau),\kappa(\tau)+x^{42}]$). That theorem is the converse of the
upper bound side of the count ({\cite[Proposition~14.10]{P46}}, proved in {\cite[\S4.2]{P49}}): there,
every jet along some $Q_{T,y}$ or along an entry of the Schur matrix has valuation \emph{at least} the predicted divisibility; here,
some jet has valuation \emph{at most} $F(p)$ plus a loss $2\sum_\tau v_p(B_\tau)$ that is absorbed in
the number $t$ uniformly in $p$.

\subsubsection{Least valuations of coefficients of theta products}\label{psm:ss-least}

Throughout \S\S\ref{psm:ss-least}--\ref{psm:ss-twelve}, $p$ is any prime (split, inert or $d$),
$\C\to\C_p$ is fixed, $v$ is the valuation of $\C_p$ with $v(p)=1$, $\pao$ its valuation ring and
$\pam$ its maximal ideal. Let $(\Lambda,t,a_1)$ be a normalised lattice ({\cite[\S7.1]{P49}}) and
$\theta$, $x$, $y$, $g=-\log\theta$ the series of {\cite[\S4.3]{P49}}, read in $\C_p[[z]]$. By
{\cite[Lemma~7.2]{P49}} the curve $\mathcal E\colon y^2+a_1xy=x^3+x$ has good reduction at $p$ and $P_0$
does not reduce to the origin; so the construction of the ring
\[
\psmR=\pao[[y]]\subset\C_p[[z]]
\]
in {\cite[\S4.3]{P49}} applies at every prime: $\psmR$ is closed under coefficientwise limits,
contains $x$, is stable under $d/dz$, and is the completion of the smooth model of $\mathcal E$
along $P_0$. Since $y=z+O(z^2)$, a series in $\psmR$ has integral coefficient in degree $0$ and
$1+y\psmR$ is a multiplicative group.

Let $X$ be a finite set of variables $z=(z_x)_{x\in X}$. For a series $F\in\C_p[[z]]$ and
$\nu\in\N^X$ write $[\nu]F$ for the coefficient of $z^\nu$, $|\nu|=\sum_x\nu(x)$ and $\nu!=
\prod_x\nu(x)!$. Let $\mu=(\mu_i)_{i\in J}$ be a finite family of linear forms in $z$ with
coefficients in $\pao$ and $m=(m_i)$ multiplicities. Put
\[
\Theta=\Theta_{\mu,m}=\prod_{i\in J}\theta(\mu_i(z))^{m_i},\qquad
\psmR_\mu=\Bigl\{\textstyle\sum_{e\in\N^J}c_e\prod_iy(\mu_i(z))^{e_i}:\ c_e\in\pao\Bigr\}.
\]
$\psmR_\mu$ is a ring containing $G(\mu_i(z))$ for $G\in\psmR$, and $\Theta(0)=1$.

\begin{definition}[Least valuation on $n$ copies]\label{psm:def-val}
For $n\ge1$, $\nu\in\N^X$ and $r\in\R$ we write $\psmval_n(\nu)\le r$ if there are
$H_1,\dots,H_n\in\psmR_\mu$ and $\alpha_1,\dots,\alpha_n\in\N^X$ with $\sum_c\alpha_c=\nu$ and
\[
\sum_{c=1}^nv\bigl([\alpha_c](H_c\Theta)\bigr)\le r .
\]
Equivalently: some coefficient of degree $\nu$ (after identifying the $n$ copies of the variables)
of the series $\prod_c(H_c\Theta)(z_c)$ in $n$ copies $z_1,\dots,z_n$ of the variables has
valuation at most $r$, the coefficient being a product of one coefficient of each factor.
\end{definition}

\begin{lemma}[Rules]\label{psm:rules}
\begin{enumerate}
\item[(a)] If $\psmval_n(\nu)\le r\le r'$ then $\psmval_n(\nu)\le r'$; and $\psmval_n(0)\le0$.
\item[(b)] If $\psmval_n(\nu)\le r$ and $\psmval_{n'}(\nu')\le r'$ then $\psmval_{n+n'}(\nu+\nu')
\le r+r'$.
\item[(c)] \emph{(Integral substitution.)} Let $H_1,\dots,H_n\in\psmR_\mu$, let $K$ be a finite set,
$A_{c,k}\in\pao$, and $w=(w_{x,k})$ variables. If the coefficient of $w^\gamma$, $\gamma\in
\N^{X\times K}$, in $\prod_{c}(H_c\Theta)\bigl(\sum_kA_{c,k}w_{\cdot,k}\bigr)$ has valuation at
most $r$, then $\psmval_n(\bar\gamma)\le r$, where $\bar\gamma(x)=\sum_k\gamma(x,k)$.
\item[(d)] \emph{(Monomials.)} If $\psmval_n(\nu)\le r$, the $H_c$ in
Definition~\ref{psm:def-val} may be taken to be monomials $\prod_iy(\mu_i(z))^{e_{c,i}}$ with
$\sum_ie_{c,i}\le|\nu|$.
\end{enumerate}
\end{lemma}
\begin{proof}
(a) Clear; for the second claim take $H_c=1$, $\alpha_c=0$, as $\Theta(0)=1$. (b) Concatenate the
two families. (c) The coefficient of $w^\gamma$ in the product is $\sum\prod_c[\beta_c]\bigl((H_c
\Theta)(\sum_kA_{c,k}w_{\cdot,k})\bigr)$ over the decompositions $\gamma=\sum_c\beta_c$; by the
ultrametric inequality one term has valuation at most $r$. For a series $G$ in $z$,
\[
[\beta]\,G\Bigl(\sum_kA_kw_{\cdot,k}\Bigr)=[\bar\beta]G\cdot\prod_x\frac{\bar\beta(x)!}{\prod_k
\beta(x,k)!}\cdot\prod_{x,k}A_k^{\beta(x,k)} ,
\]
an integral multinomial coefficient times integral numbers times $[\bar\beta]G$; so
$v([\bar\beta_c](H_c\Theta))\le v([\beta_c](\cdots))$, and $\alpha_c=\bar\beta_c$ satisfy
$\sum_c\alpha_c=\bar\gamma$. (d) Let $H=\sum_ec_e\prod_iy(\mu_i)^{e_i}\in\psmR_\mu$ and $\alpha\in
\N^X$. Since the $y(\mu_i(z))$ have no constant term, only the $e$ with $|e|\le|\alpha|$ contribute
to $[\alpha](H\Theta)$, and $[\alpha](H\Theta)=\sum_{|e|\le|\alpha|}c_e\,[\alpha](\prod_iy(\mu_i)^{
e_i}\Theta)$ with $c_e\in\pao$. By the ultrametric inequality one of the finitely many monomials
gives a coefficient of valuation at most $v([\alpha](H\Theta))$; and $|\alpha_c|\le|\nu|$.
\end{proof}

\subsubsection{One curve: units, the rate $\gamma_p$, and the leading coefficient}\label{psm:ss-curve}

\begin{lemma}[Isogeny units]\label{psm:units}
\begin{enumerate}
\item[(a)] For every $\alpha\in\auO_K$, $\alpha\ne0$, the series $G(\alpha z)$ lies in $\psmR$ for all
$G\in\psmR$, and $f_\alpha(z)=\theta(\alpha z)\,\theta(z)^{-\auNm\alpha}$ lies in $1+y\psmR$.
\item[(b)] Let $(\delta,s)$ lie in the closure in $\C_p^2$ of $\{(\alpha,\auNm\alpha):\alpha\in
\auO_K\setminus0\}$. Then $G(\delta z)\in\psmR$ for $G\in\psmR$, and $\exp\bigl(s\,g(z)-g(\delta
z)\bigr)\in1+y\psmR$. This applies to $(\delta,\delta^2)$ for every $\delta\in\Z_p$.
\end{enumerate}
\end{lemma}
\begin{proof}
(a) In the uniformisation $z\mapsto Q(z)$ of {\cite[\S7.2]{P49}}, $z\mapsto Q(z)-P_0$ is a group
isomorphism, so $z\mapsto\alpha z$ is the morphism $Q\mapsto\alpha(Q-P_0)+P_0$ of $\mathcal E$,
which fixes $P_0=Q(0)$. It extends to the smooth model, which is the N\'eron model of its generic
fibre~\cite[\S1.2, Definition~1, p.~12, and Proposition~8, p.~15]{BLR} (\cite[IV, Corollary~6.3, p.~329, and the N\'eron mapping property, p.~319]{SilvermanAT}), and fixes the section $P_0$; hence it preserves $\psmR$, as in
{\cite[(26)]{P49}}. No congruence condition on $\alpha$ is needed for this. By
{\cite[Proposition~7.6]{P49}} with $u=0$ and $\aubar\delta=\alpha$, $\theta(\alpha z)=\theta(z)^{
\auNm\alpha}\Gamma(y(z))$ with $\Gamma\in\pao[[Y]]$, at every prime; comparing constant terms,
$\Gamma(0)=1$, so $f_\alpha=\Gamma(y)\in1+y\psmR$.

(b) Let $(\alpha_n,\auNm\alpha_n)\to(\delta,s)$. The coefficients of $G(\alpha_nz)$ and of
$\exp(\auNm\alpha_n\,g(z)-g(\alpha_nz))=f_{\alpha_n}^{-1}$ are polynomials in $(\alpha_n,
\auNm\alpha_n)$ (with fixed coefficients), and they lie in the closed sets $\psmR$ and $1+y\psmR$
by (a) ($1+y\psmR$ is a closed group). Pass to the limit. For $\delta\in\Z_p$ take
$\alpha_n\in\Z$, $\alpha_n\to\delta$; then $\auNm\alpha_n=\alpha_n^2\to\delta^2$.
\end{proof}

\begin{lemma}[Six copies]\label{psm:six}
There are $\delta_0=1,\delta_1,\dots,\delta_5\in\Z_p$ such that, for every family $(\mu,m)$ as above,
$U=\prod_{c=0}^5\Theta(\delta_cz)$ lies in $\psmR_\mu$, has constant term $1$, and has an inverse in
$\psmR_\mu$. If moreover $v(\varpi)\ge\rho_p$ \textup{(}with $\rho_p$ as in
Lemma~\ref{psm:param}\textup{)}, then $H(\varpi z)$ is integral for every $H\in\psmR_\mu$.
\end{lemma}
\begin{proof}
Every element of $\Z_p$ is a sum of four squares: the set of sums of four squares of elements of
$\Z_p$ is compact, hence closed, and contains $\N$ (Lagrange), which is dense. Choose
$a_1,\dots,a_4\in\Z_p$ with $\sum a_i^2=-1$ and put $(\delta_c)=(1,a_1,a_2,a_3,a_4,0)$, so that
$\sum_c\delta_c^2=0$. Then
\[
u(z):=\prod_c\theta(\delta_cz)=\exp\Bigl(-\sum_cg(\delta_cz)\Bigr)=\prod_c\exp\bigl(\delta_c^2g(z)
-g(\delta_cz)\bigr)\in1+y\psmR
\]
by Lemma~\ref{psm:units}(b), and $u^{-1}\in1+y\psmR$. Since $U=\prod_iu(\mu_i(z))^{m_i}$, the first
assertion follows. For the second, $y(\varpi w)$ is an integral series in $w$ when
$v(\varpi)\ge\rho_p$ (Lemma~\ref{psm:param}(c)), and the $\mu_i$ have integral coefficients.
\end{proof}

For the next lemma let $\pi$, $q$ be as follows: if $p$ does not split, $\pi=\pi_p$ and $q=q_p$ as
in {\cite[\S4.1]{P49}}; if $p$ splits, $\pi=p$ and $q=p$. Put $\rho_p=v(\pi)/(q-1)$ and
$\gamma_p=\frac1{p-1}-\rho_p$; thus $\gamma_p=0$ if $p$ splits, $\gamma_p=p/(p^2-1)$ if $p$ is
inert and $\gamma_p=1/(2p-2)$ if $p=d$, in agreement with {\cite[\S4.1]{P49}}.

\begin{lemma}[The formal parameter at every prime]\label{psm:param}
Let $T=\frac1{q-1}\sum_{\zeta\in\mu_{q-1}}\zeta^{-1}y(\zeta z)$, where $\mu_{q-1}\subset\C_p$
\textup{(}for split $p$, $\mu_{p-1}\subset\Z_p$\textup{)}.
\begin{enumerate}
\item[(a)] $T\in y+y^2\psmR$, $\psmR=\pao[[T]]$, and $T(\pi z)=\pi T+\sum_{k\ge1}b_kT^{1+k(q-1)}$
with $b_k\in\pao$, $b_1\in\pao^\times$.
\item[(b)] With $\varpi^{q-1}=-\pi/b_1$ \textup{(}so $v(\varpi)=\rho_p$\textup{)}, $Y(z)=T(\varpi z)/
\varpi$ is integral and $Y\equiv\sum_{j\ge0}z^{q^j}\pmod{\pi}$.
\item[(c)] For $G\in\psmR$ and $v(\varpi')\ge\rho_p$ the series $G(\varpi'z)$ is integral.
\item[(d)] For every $j\ge0$, $v\bigl((q^j)!\,[q^j]T\bigr)=\gamma_p\,(q^j-1)$.
\end{enumerate}
\end{lemma}
\begin{proof}
For non-split $p$, (a)--(c) are {\cite[Lemma~4.6]{P49}}, whose proof uses {\cite[(26)]{P49}} for
$\alpha\in\auO_{K_p}$; this is now Lemma~\ref{psm:units}(b) (for $\zeta\in\mu_{q-1}$, the pair
$(\zeta,\zeta\cdot\zeta^{p})$ or $(\zeta,\zeta^2)$ is a limit of pairs $(\alpha,\auNm\alpha)$, as
in the proof of {\cite[Lemma~4.7(b)]{P49}}). For split $p$ the same proof applies with $\pi=p$, $q=p$,
once we know that the reduction of $[p]$, written in the parameter $T$, has order exactly $p$: the
reduction of $\mathcal E$ is ordinary~\cite[Chapter~13, Theorem~12]{LangEF}, so its formal group has
height one and $[p]$ is congruent modulo $\pam$ to a unit times $T^p$ plus terms of higher degree;
$\mu_{p-1}\subset\Z_p$, and $G(\zeta z)\in\psmR$ by Lemma~\ref{psm:units}(b). (For $p=2$, $\mu_1=\{1\}$
and $T=y$.) The computation of {\cite[Lemma~4.6(b)]{P49}} uses only $v(\pi^{n-1})>0$ for $n\ge2$ and
(E3), valid also for $p=2$. (d) By (b), $\varpi^{q^j}[q^j]T=\varpi\,[q^j]Y$ and $[q^j]Y\equiv1$, so
$v([q^j]T)=-(q^j-1)\rho_p$; and $v((q^j)!)=(q^j-1)/(p-1)$ by Legendre's formula, $q^j$ being a power
of $p$.
\end{proof}

\begin{lemma}[The rate $\gamma_p$ is attained]\label{psm:katz}
For every $n\ge0$ there is $G\in\psmR$ with $v\bigl(n!\,[n]G\bigr)\le\gamma_p\,n$.
\end{lemma}
\begin{proof}
For $G\in\psmR$, $n![n]G=(d/dz)^nG(0)$ is integral, because $\psmR$ is stable under $d/dz$ and the
constant term of an element of $\psmR$ is integral. Let $c(n)=\max_{G\in\psmR}|n![n]G|$; the
maximum exists and is attained at a power $y^i$, $i\le n$, by the argument of
Lemma~\ref{psm:rules}(d). Then $c(0)=1$ and $0\le c(n)\le1$.

\emph{Submultiplicativity.} Let $c(m+n)=|(m+n)![m+n]y^i|$. The formal group law of the smooth model
with origin $P_0$ is integral (for the origin $O$ and the parameter $-x/y$ of a Weierstrass equation
with integral coefficients this is \cite[Chapter~IV, \S1]{SilvermanAEC}; for $P_0$ we argue directly): the morphism $(Q,Q')\mapsto Q+Q'-P_0$
extends to the smooth model and fixes $(P_0,P_0)$, hence induces a map of the completions along
$(P_0,P_0)$ and $P_0$; since the completion of the smooth model along $P_0$ is $\pao[[y]]$
({\cite[\S4.3]{P49}}), that of the product along $(P_0,P_0)$ is $\pao[[y(z_0),y(z_1)]]$, and so $y(z_0+z_1)^i=\sum_{a,b}c_{ab}y(z_0)^ay(z_1)^b$
with $c_{ab}\in\pao$. Since $(m+n)![m+n]y^i=m!\,n!\,[z_0^mz_1^n]\,y(z_0+z_1)^i$, the ultrametric
inequality gives $c(m+n)\le\max_{a,b}|m![m]y^a|\cdot|n![n]y^b|\le c(m)c(n)$.

\emph{Powers of $q$.} By Lemma~\ref{psm:param}(d), $c(q^j)\ge p^{-\gamma_p(q^j-1)}\ge
\beta^{q^j}$ with $\beta=p^{-\gamma_p}\in(0,1]$.

\emph{Fekete.} Suppose $c(n)<\beta^n$ for some $n\ge1$, and let $\varepsilon=c(n)/\beta^n<1$. For
$Q=q^j\ge n$ write $Q=nk+r$, $0\le r<n$. Then $\beta^Q\le c(Q)\le c(n)^kc(r)\le c(n)^k=
\varepsilon^k\beta^{nk}$, so $\beta^{r}\le\varepsilon^k$, that is $\beta^n\le\varepsilon^k$ for all
large $j$, which is absurd as $k\to\infty$. So $c(n)\ge\beta^n$, which is the claim.
\end{proof}

\begin{lemma}[Unit forms]\label{psm:unitform}
\begin{enumerate}
\item[(a)] If one form $\mu_{i_0}$ has all its coefficients units, then for every $\nu$,
$\psmval_6(\nu)\le\gamma_p|\nu|-v(\nu!)$.
\item[(b)] If there are $a_{x,i}\in\pao$ and a unit $s$ with $\sum_ia_{x,i}\mu_i(z)=s\,z_x$ for
all $x\in X$, then $\psmval_6(\nu)\le0$ for every $\nu$.
\end{enumerate}
\end{lemma}
\begin{proof}
(a) Let $n=|\nu|$, $G$ as in Lemma~\ref{psm:katz}, and $(\delta_c)$, $U$ as in
Lemma~\ref{psm:six}. Put $H_0=G(\mu_{i_0}(z))\,U^{-1}$ and $H_c=1$ for $c\ge1$; all lie in
$\psmR_\mu$. As $\delta_0=1$, $\prod_c(H_c\Theta)(\delta_cz)=G(\mu_{i_0}(z))U(z)^{-1}U(z)=
G(\mu_{i_0}(z))$, whose coefficient of $z^\nu$ is $[n]G\cdot\frac{n!}{\nu!}\prod_x\mu_{i_0,x}^{\nu(x)}$,
of valuation $v(n![n]G)-v(\nu!)\le\gamma_pn-v(\nu!)$. Apply Lemma~\ref{psm:rules}(c) with one new
copy of the variables and $A_c=\delta_c$.

(b) Put $G_x=\sum_ia_{x,i}\,y(\mu_i(z))\in\psmR_\mu$. Since $y=z+O(z^2)$, $G_x=s\,z_x+(\text{terms of
degree}\ge2)$, so $\bigl(\prod_xG_x^{\nu(x)}\bigr)\Theta=s^{|\nu|}z^\nu+(\text{terms of degree}>
|\nu|)$, and its coefficient of $z^\nu$ is the unit $s^{|\nu|}$. This gives $\psmval_1(\nu)\le0$;
add five copies in degree $0$ (Lemma~\ref{psm:rules}(a),(b)).
\end{proof}

\begin{lemma}[The leading coefficient is a unit]\label{psm:lambda}
Let $p$ be inert. The number $\lambda$ of {\cite[Theorem~4.9]{P49}} is a unit; hence in the normal
form of {\cite[Proposition~4.12]{P49}}, if the scalar $c$ there is a unit, all $\lambda_{e,j}$ $(e$
odd$)$ are units; they are $c^{(p^e+1)p^j}\lambda$.
\end{lemma}
\begin{proof}
Let $h_n=[z^{p-1}]x^n$. From {\cite[(25)]{P49}}, $x'^2=(2y+a_1x)^2=4x^3+a_1^2x^2+4x$ and
$x''=6x^2+a_1^2x+2$, so
\[
\bigl(x^nx'\bigr)'=nx^{n-1}x'^2+x^nx''=(4n+6)x^{n+2}+(n+1)a_1^2x^{n+1}+(4n+2)x^n .
\]
The coefficient of $z^{p-1}$ of the left side is $p\,[z^p](x^nx')$, and $x^nx'\in\psmR$ has integral
coefficients in degrees $<q=p^2$ ({\cite[Lemma~4.6(c)]{P49}}); so
$(4n+6)h_{n+2}+(n+1)a_1^2h_{n+1}+(4n+2)h_n\in\pam$ for all $n\ge0$, with $h_n$ integral and $a_1$
integral ({\cite[Lemma~4.5]{P49}}). Write $p=2m'+1$. If $h_1\in\pam$, then, as $h_0=0$ and
$4n+6=2(2n+3)$ is a unit for $n+2\le m'$ ($2n+3\le p-2$), induction gives $h_n\in\pam$ for all
$n\le m'$. But $x=z^2+O(z^3)$, so $h_{m'}=[z^{2m'}]x^{m'}=1$. Hence $x_{p-1}=h_1$ is a unit.

By {\cite[(28)]{P49}}, $(p+1)p\,g_{p+1}=x_{p-1}$. In the proof of {\cite[Theorem~4.9]{P49}},
$\Xi_1(z)=-g^\flat(\varpi z)/c_\varpi$ with $c_\varpi=\varpi^{p+1}/p$, and $p+1$ is one of the degrees
kept by $\flat$; so $[p+1]\Xi_1=-p\,g_{p+1}=-x_{p-1}/(p+1)$, a unit. On the other hand
$\Xi_1\equiv\lambda Y^{p+1}\pmod p$ and $[p+1]Y^{p+1}\equiv1$ (Step~3 of that proof); so
$\lambda\equiv[p+1]\Xi_1$ is a unit. The statement on $\lambda_{e,j}$ is read off from
{\cite[Theorem~4.9]{P49}} and the proof of {\cite[Proposition~4.12]{P49}}, where the substitution
$z\mapsto cz$ multiplies $\lambda_{e,j}$ by $c^{(p^e+1)p^j}$.
\end{proof}

\subsubsection{Couplings}\label{psm:ss-couplings}

In this section $p$ is inert in $K$ and $p>2$. Let $\auPl$ be a finite set, $X=\auPl$ (one
variable $z_\tau$ for each $\tau$), and let $\psmpr_e\colon\auPl\to\auPl$ ($e$ odd) be maps. Let
$(\mu,m)$ be a family as in \S\ref{psm:ss-least} such that $\Theta=\Theta_{\mu,m}$ has the normal
form {\cite[(20)]{P49}} for the partner maps $\psmpr_e$:
\[
\Theta(\varpi z)=W(z)\exp\bigl(\kappa_p\,\Xi(z)\bigr),\qquad\kappa_p=\frac{\varpi^{p+1}}p,\qquad
\Xi=\sum_{e\text{ odd}}\sum_{j\ge0}\lambda_{e,j}\Gamma_e^{\,p^j},\qquad\Gamma_e=\sum_\tau
z_\tau^{\,p^e}z_{\psmpr_e\tau},
\]
with $v(\varpi)=\rho_p=1/(p^2-1)$, $W$ integral, $\lambda_{e,j}\in\pao$, and assume that
$\lambda_{e,0}$ is a unit for every odd $e$. Then $v(\kappa_p)=\frac1{p-1}-1<0$ and $W(0)=1$, so
$W$ is a unit of $\pao[[z]]$.

\begin{definition}\label{psm:def-family}
A \emph{coupling family} is a subset $\auPl_0\subseteq\auPl$ with odd exponents $e(\tau)$,
$\tau\in\auPl_0$, such that: for $\tau,\tau_1\in\auPl_0$ and odd $e'\le e(\tau)$,
$\psmpr_{e'}(\tau_1)=\psmpr_{e(\tau)}(\tau)$ implies $\tau_1=\tau$ and $e'=e(\tau)$. Write
$\tau^+=\psmpr_{e(\tau)}(\tau)$. For $j\colon\auPl_0\to\N$ put
\[
\nu'=\sum_{\tau\in\auPl_0}p^{e(\tau)}j(\tau)\,\mathbf 1_\tau,\qquad\nu''=\sum_{\tau\in\auPl_0}j(\tau)
\,\mathbf 1_{\tau^+},\qquad|j|=\sum_\tau j(\tau),
\]
the degrees of the two factors of $\prod_\tau(z_\tau^{p^{e(\tau)}}z_{\tau^+})^{j(\tau)}$.
\end{definition}

\begin{proposition}[Couplings]\label{psm:coupling}
For a coupling family and every $j$,
\[
\psmval_6(\nu'+\nu'')\ \le\ \frac{|j|}{p-1}-|j|-\rho_p|\nu'+\nu''|-v(\nu''!) .
\]
Consequently $\psmval_6(\nu'+\nu'')\le-|j|$ and $v(\nu'!)+v(\nu''!)+\psmval_6(\nu'+\nu'')\le
\gamma_p|\nu'+\nu''|-|j|$.
\end{proposition}
\begin{proof}
Let $(\delta_c)$ and $U$ be as in Lemma~\ref{psm:six}, $F_\varpi(z)=\Theta(\varpi z)$ and
$U'(z)=U(\varpi z)=\prod_cF_\varpi(\delta_cz)$; $U'$ is integral with constant term $1$
(Lemma~\ref{psm:six}). In two copies $z',z''$ of the variables put
\[
G(z',z'')=F_\varpi(z'+z'')\prod_{c=1}^5F_\varpi(\delta_cz'),\qquad G(z',0)=U'(z') .
\]
\emph{Step 1: derivatives in $z''$.} Only the first factor involves $z''$, and
$\partial_{z_x}F_\varpi=(\kappa_p\,\partial_x\Xi+\partial_xW/W)F_\varpi$; so $\partial_{z''_x}G=(\kappa_pA_x+B_x)G$
with $A_x=(\partial_x\Xi)(z'+z'')$, $B_x=(\partial_xW/W)(z'+z'')$ integral. We claim: for variables
$t_1,\dots,t_N$ among the $z''_x$ there is an integral $g$ with $\partial_{t_1}\cdots\partial_{t_N}G=
\kappa_p^Ng\,G$ and $g\equiv\prod_iA_{t_i}\pmod\pam$. For $N=0$ take $g=1$. If $\partial(g G)$ is
to be computed for one more derivative $\partial=\partial_t$, then
\[
\partial(gG)=\bigl(\partial g+g(\kappa_pA_t+B_t)\bigr)G=\kappa_p\bigl(gA_t+\kappa_p^{-1}(\partial g+
gB_t)\bigr)G,
\]
and $\kappa_p^{-1}\in\pam$. Apply this to the multiset in which each $z''_{\tau^+}$ occurs $j(\tau)$
times ($N=|j|$): $\partial_{z''}^{\,\nu''}G=\kappa_p^{|j|}g\,G$ with
\[
g(z',0)\equiv\prod_{\tau\in\auPl_0}\bigl(\partial_{\tau^+}\Xi\bigr)(z')^{j(\tau)}\pmod\pam .
\]
\emph{Step 2: a unit coefficient.} We show that $[z'^{\nu'}]\bigl(g(z',0)U'(z')\bigr)$ is a unit.
Since $\nu'$ is supported on $\auPl_0$ we may put $z'_\sigma=0$ for $\sigma\notin\auPl_0$ and reduce
modulo $\pam$. Modulo $p$, $\partial(\Gamma_e^{p^j})=p^j\Gamma_e^{p^j-1}\partial\Gamma_e\equiv0$ for
$j\ge1$, and $\partial_{\tau^+}\Gamma_e\equiv\sum_{\sigma:\,\psmpr_e\sigma=\tau^+}z_\sigma^{p^e}$
(the other term carries the factor $p^e$). So, after putting $z_\sigma=0$ off $\auPl_0$,
\[
\partial_{\tau^+}\Xi\equiv\sum_{e\text{ odd}}\lambda_{e,0}\sum_{\sigma\in\auPl_0,\ \psmpr_e\sigma=
\tau^+}z_\sigma^{\,p^e}\pmod\pam ,
\]
a sum of monomials of degree $p^e$; by the definition of a coupling family the only one of degree
$\le p^{e(\tau)}$ is $\lambda_{e(\tau),0}z_\tau^{p^{e(\tau)}}$. Hence the product over $\tau$, times
$U'\equiv1+(\text{higher degree})$, is $\prod_\tau\lambda_{e(\tau),0}^{j(\tau)}\,z^{\nu'}$ plus terms of
higher total degree, and $\prod_\tau\lambda_{e(\tau),0}^{j(\tau)}$ is a unit.

\emph{Step 3: the coefficient.} The jet $\partial_{z'}^{\nu'}\partial_{z''}^{\nu''}G(0)$ equals
$\nu'!\,\nu''!\,[z'^{\nu'}z''^{\nu''}]G$ and also, by Step~1 at $z''=0$, $\kappa_p^{|j|}\nu'!\,
[z'^{\nu'}](g(z',0)U'(z'))$. By Step~2,
\[
v\bigl([z'^{\nu'}z''^{\nu''}]G\bigr)=|j|\,v(\kappa_p)-v(\nu''!)=|j|\bigl(\tfrac1{p-1}-1\bigr)-
v(\nu''!) .
\]
\emph{Step 4: unscaling.} $G(z',z'')=G_0(\varpi z',\varpi z'')$ with $G_0=\Theta(z'+z'')\prod_{c\ge1}
\Theta(\delta_cz')$, so $[z'^{\nu'}z''^{\nu''}]G_0$ has valuation $|j|(\frac1{p-1}-1)-v(\nu''!)-
\rho_p|\nu'+\nu''|$. Now $G_0=\prod_{c=0}^5\Theta(A_cz'+B_cz'')$ with $(A_0,B_0)=(1,1)$ and
$(A_c,B_c)=(\delta_c,0)$, integral; Lemma~\ref{psm:rules}(c) with $H_c=1$ gives the first claim.

For the consequences: $|\nu'+\nu''|=|\nu'|+|j|\ge(p+1)|j|$ (as $e(\tau)\ge1$), and
$\rho_p(p+1)=\frac1{p-1}$; so the bound is $\le-|j|-v(\nu''!)\le-|j|$. And by Legendre's formula
$v(\nu'!)\le|\nu'|/(p-1)$, so $v(\nu'!)+v(\nu''!)+\psmval_6\le\frac{|\nu'|+|j|}{p-1}-\rho_p|\nu'+
\nu''|-|j|=\gamma_p|\nu'+\nu''|-|j|$.
\end{proof}

\subsubsection{The bound at every prime}\label{psm:ss-twelve}

Now let $D$ be admissible data ({\cite[Definition~13.3]{P46}}), $\auPl$ its set of infinite places,
$\kappa$ its weight, $\auPb$ its $b$-places, $l_0=l_0(s)$, $k=r+1-m$, $a=xb$, and for a prime $p$
let $u(-p^e)\tau$ be the partner of $\tau$ ({\cite[\S14.4]{P46}}). Put
\[
\psmns=\lceil6b/k\rceil,\qquad l^-=(l-\psmns)^+,\qquad B(l)=(l^-+1)(l^-+2)\cdots(l^-+\psmns)
\qquad(l\in\N).
\]
Let $\mathsf w=(\mathsf w_i)_{i\in J}$ be a family in $\mathfrak a_0$ with multiplicities $c_i$
forming a \emph{tight frame of length $M$}: $\sum_ic_i\,\tau(\mathsf w_i)\aubar{\tau'(\mathsf w_i)}=
M\delta_{\tau\tau'}$ (cf. {\cite[\S7.4]{P49}}). Let $\mathsf s\in\C$ be the number
$\varpi_{\mathrm{vac}}/\mathsf u$ of {\cite[\S7.4]{P49}} and, at the prime $p$, consider the linear
forms in $z=(z_\tau)_{\tau\in\auPl}$
\[
\mu_i(z)=\mathsf s\sum_\tau\aubar{\tau(\mathsf w_i)}\,z_\tau=\mathfrak s\,u_i^\dagger(z),\qquad
u_i=\mathsf w_i/\aubar{\pi_M},\quad\mathfrak s=\mathsf s\,\pi_M .
\]
Their coefficients are integral at every prime ({\cite[Lemma~7.12]{P49}} and {\cite[Corollary~7.3]{P49}}).

\begin{lemma}[Facts at a non-split prime]\label{psm:nonsplit}
Let $p$ be a prime that does not split in $K$. Then $p\ge13$, $p\nmid q_1M$, the vectors $u_i$
are integral above $p$, they form with the multiplicities $c_i$ a resolution of the identity
{\cite[(30)]{P49}}, and $\mathfrak s$ is a $p$-adic unit. If $p$ is inert, $\Theta_{\mu,c}$ has the
normal form of \S\ref{psm:ss-couplings} with all $\lambda_{e,0}$ units, and there are
$a_{\tau,i}\in\pao$ with $\sum_ia_{\tau,i}\mu_i(z)=\mathfrak s\,z_\tau$.
\end{lemma}
\begin{proof}
$p\ge13$ by {\cite[(34)]{P46}}; the primes of $q_1$ split, and $p_0,p_1\equiv1\pmod{q_2^2}$ split
(their residue classes are $1$ modulo $d\mid q_2$, hence squares modulo $d$). The resolution property
is $\sum_ic_i\tau(u_i)\aubar{\tau'(u_i)}=\delta_{\tau\tau'}$, which is the frame identity divided by
$\pi_M\aubar{\pi_M}=M$. Since $q_1\mathfrak a_0\subseteq\auO_E$, $q_1M\,u_i=q_1\mathsf w_i\cdot
\pi_M$ is integral, and $q_1M$ is prime to $p$. $\pi_M$ is a unit at $p$ (its norm $M$ is), and
$\mathsf s$ is a unit at $p\nmid q_1$ ({\cite[Lemma~7.12]{P49}}: $|\varpi_{\mathrm{vac}}|=
|\alpha|_w^{-1}=1$ as $\alpha$ is a unit off $q_1$; and $\mathsf u$ is a unit by
{\cite[Corollary~7.3]{P49}}). For inert $p$, {\cite[Proposition~4.12]{P49}} (with $G=1$ and the unit scalar
$\mathfrak s$) gives the normal form and Lemma~\ref{psm:lambda} the units. Finally $a_{\tau,i}=c_i
\tau(u_i)$ works by {\cite[(30)]{P49}}.
\end{proof}

Write $I_p$, $v(\pi_p)$ as in {\cite[\S14.4]{P46}} and recall $s_p$ and $F(p)$
\eqref{psm:eq-F}; for $p=d$, $v(\pi_d)=\frac12$ and $I_d=\{1,2,\dots\}$. We use the following size
inequalities.

\begin{lemma}[Sizes]\label{psm:sizes}
\begin{enumerate}
\item[(N1)] $x^{42}\le6b/k\le\psmns$, and $\kappa(\tau)+x^{42}\le2b$ for all $\tau$.
\item[(N2)] $3(\log b+2)+x^{42}+1\le a/x_1$, and $a/x_1\ge1$.
\item[(N3)] $(81s+1)\,a/x_1\le b/k^2$.
\item[(N4)] $d\le k\le2^s$, $27s<l_0\le81s$, and $k\le5r(r+1-m)$.
\end{enumerate}
\end{lemma}
\begin{proof}
By {\cite[Lemma~13.4]{P46}}: $k\le r\le3x^5$, $x^{10^6}\le x_1$, $x_1^{10^6}\le b$, $s\le32\log x$,
$x^{11}\le2^s$, $\log x\ge10^{14}d$, $\kappa(\tau)\le b$. (N1): $x^{42}k\le3x^{47}\le6b$ and
$x^{42}\le b$. (N2): $a/x_1=xb/x_1\ge b^{1-10^{-6}}\ge x^{43}$, which exceeds
$3(\log b+2)+x^{42}+1$ as $\log b\le x^{1/100}$. (N3): equivalent to $(81s+1)k^2x\le x_1$, and
$(81s+1)k^2x\le(2592\log x+1)\,9x^{11}\le x^{12}\le x_1$. (N4): $d\le\log x\le x^5/2\le k$ (S2), and
$k\le3x^5\le x^{11}\le2^s$; the bounds for $l_0$ are in {\cite[\S13.1]{P46}}; $5r(r+1-m)=5rk\ge k$.
\end{proof}

\emph{Near pairs.} For a prime $p\le k$ that does not split, a \emph{near pair} is a pair $(\tau,e)$
with $e\in I_p$, $2e<l_0$, $\tau\in\auPb$ and $u(-p^e)\tau\in\auPb$.

\begin{lemma}[Near pairs]\label{psm:near}
Let $p\le k$ be a prime that does not split in $K$.
\begin{enumerate}
\item[(a)] If $p=d$ there are no near pairs. If $(\tau,e)$ is a near pair, then $p$ is inert, $e$ is
odd and $u(-p^e)\tau=u(-p)^e\tau$.
\item[(b)] Two near pairs with the same start have the same exponent; two near pairs with the same
end are equal; the end of a near pair is not the start of a near pair.
\end{enumerate}
\end{lemma}
\begin{proof}
(a) All $b$-places have class $-1$ modulo $q_0$, so a near pair needs $\ancl_{q_0}(u(-p^e))=1$,
i.e.\ $p^e\equiv-1\pmod{q_0}$. For $p=d$: with $h=q_0/2$ we have $d\equiv h-1\pmod{q_0}$ and
$h^2\equiv0\pmod{q_0}$ ($4\mid q_0$), so $d^e\equiv(-1)^e+e(-1)^{e-1}h$. If $e$ is odd this is
$-1+eh\equiv-1+h\not\equiv-1$; if $e$ is even it is $1-eh\equiv1\not\equiv-1$. So $p\ne d$, $p$ is
inert and $I_p$ consists of odd numbers; then $-p^e=(-p)^e$.

(b) By (D5) the set $\auPb$ is sparse for the primes up to $5r(r+1-m)\ge k\ge p$
({\cite[Definition~13.2]{P46}}); in the classes modulo the largest divisor $q'$ of $q_1$ prime to
$p$, $u(-p)^i\sigma/\sigma$ has class $(-1)^ip^i$. Hence for $0<i<i'<l_0$ the places
$\sigma,u(-p)^i\sigma,u(-p)^{i'}\sigma$ are not all in $\auPb$. Same start $\tau$, exponents
$e<e'$: take $\sigma=\tau$, $(i,i')=(e,e')$, $e'<l_0/2$. Same end, $e<e'$: then $\tau=u(-p)^{e'-e}
\tau'$, and take $\sigma=\tau'$, $(i,i')=(e'-e,e')$; so $e=e'$ and then $\tau=\tau'$. End of
$(\tau,e)$ equal to the start of $(\tau',e')$: take $\sigma=\tau$, $(i,i')=(e,e+e')$, with
$e+e'<l_0$.
\end{proof}

\begin{theorem}[The bound on twelve copies]\label{psm:twelve}
Let $(\mathsf w,c)$ be a tight frame of length $M$ such that, at the prime $p$, some form
$\mu_{i_0}$ has all coefficients units. Let $l\in\N^{\auPl}$ with $\kappa(\tau)\le l_\tau\le
\kappa(\tau)+x^{42}$ for all $\tau$. Then, for $\Theta=\Theta_{\mu,c}$,
\[
\psmval_{12}(l)\ \le\ F(p)+2\sum_\tau v\bigl(B(l_\tau)\bigr)-\sum_\tau v(l_\tau!) .
\]
\end{theorem}
\begin{proof}
\emph{Case 1: $p$ splits.} Then $F(p)=0$ and $\gamma_p=0$; Lemma~\ref{psm:unitform}(a) gives
$\psmval_6(l)\le-v(l!)$, and six copies in degree $0$ give the claim.

\emph{Case 2: $p$ does not split and $p>k$.} Then $p>d$ (N4), so $p$ is inert, $v(\pi_p)=1$, and
$\psmpr_1\colon\tau\mapsto\tau^+=u(-p)\tau$ is a bijection; so $\auPl_0=\auPl$ with $e\equiv1$ is a
coupling family. Put $i(\tau)=\lfloor l_\tau^-/p\rfloor$ and $j(\tau)=\min(i(\tau),l_{\tau^+})$.
At a place $t=\tau^+$ we have $\nu'(t)=p\,j(t)\le l_t^-$ and $\nu''(t)=j(\tau)\le\min(\psmns,l_t)$,
because $i(\tau)\le l_\tau/p\le2b/k\le\psmns$ by (N1); hence $\nu'+\nu''\le l$, as
$l^-_t+\min(\psmns,l_t)=l_t$. Let $\eta=l-\nu'-\nu''$. By Proposition~\ref{psm:coupling},
Lemma~\ref{psm:nonsplit} and Lemma~\ref{psm:unitform}(b), $\psmval_{12}(l)\le-|j|+0$. It remains
to show, for each $\tau$,
\[
v(l_\tau!)-j(\tau)\le s_p(\kappa(\tau),\tau)+2v(B(l_\tau)).
\]
Since $(l^-+\psmns)!=l^-!\,B(l)$ and $l\le l^-+\psmns$, $v(l!)\le v(l^-!)+v(B(l))$; by Legendre,
$v(l^-!)=i+v(i!)$ with $i=\lfloor l^-/p\rfloor$; and $v(i!)\le v(\psmns!)\le v(B(l))$, as $i\le
\psmns$ and a product of $\psmns$ consecutive integers is divisible by $\psmns!$. So
$v(l_\tau!)\le i(\tau)+2v(B(l_\tau))$. Further $l_\tau^-\le\kappa(\tau)$ by (N1), so
$i(\tau)\le\kappa(\tau)/p$, and $l_{\tau^+}\ge\kappa(\tau^+)$; hence $i(\tau)-j(\tau)=(i(\tau)-
l_{\tau^+})^+\le(\kappa(\tau)/p-\kappa(u(-p)\tau))^+$, the term $e=1$ of $s_p(\kappa(\tau),\tau)$.

\emph{Case 3: $p$ does not split and $p\le k$.} Let $\auPl_0$ be the set of starts of near pairs and,
for $\tau\in\auPl_0$, $e(\tau)$ the exponent (Lemma~\ref{psm:near}(b)) and $j(\tau)=\lfloor l_\tau/
p^{e(\tau)}\rfloor$. If $\auPl_0=\emptyset$ (in particular if $p=d$) put $\nu'=\nu''=0$, $j=0$.
Otherwise $p$ is inert and $\auPl_0$ is a coupling family: if $\tau,\tau_1\in\auPl_0$, $e'\le
e(\tau)$ is odd and $u(-p^{e'})\tau_1=\tau^+$, then $(\tau_1,e')$ is a near pair ($e'\in I_p$,
$2e'<l_0$, $\tau_1\in\auPb$, and its end $\tau^+\in\auPb$) with the same end as $(\tau,e(\tau))$,
so $\tau_1=\tau$, $e'=e(\tau)$. We have $\nu'+\nu''\le l$: at $t\in\auPl_0$, $\nu'(t)=p^{e(t)}j(t)\le
l_t$ and $\nu''(t)=0$ (an end is not a start); at $t\notin\auPl_0$, $\nu'(t)=0$ and $\nu''(t)$ is
$j(\tau)$ for the unique near pair $(\tau,e(\tau))$ ending at $t$ (or $0$), and $j(\tau)\le l_\tau/2
\le b\le\kappa(t)\le l_t$ by (N1), as $t\in\auPb$. Let $\eta=l-\nu'-\nu''$. By
Proposition~\ref{psm:coupling} (Lemma~\ref{psm:nonsplit}) and Lemma~\ref{psm:unitform}(a),
\[
\psmval_{12}(l)\le\gamma_p|l|-|j|-v(\nu'!)-v(\nu''!)-v(\eta!) .
\]
For $n_1+n_2+n_3=n\le2b$ and $p\ge3$, Legendre's formula $\frac n{p-1}-\log_p(n+1)-1\le v(n!)\le
\frac n{p-1}$ gives $v(n!)\le v(n_1!)+v(n_2!)+v(n_3!)+3(\log b+2)$. Since $l_\tau\le2b$ (N1),
\[
\psmval_{12}(l)+\sum_\tau v(l_\tau!)\le\gamma_p|l|-|j|+3f(\log b+2).
\]
Now $\gamma_p\le1$ and $l_\tau-\kappa(\tau)\le x^{42}$, so $\gamma_p|l|\le\gamma_p\sum_\tau
\kappa(\tau)+fx^{42}$; and by $\gamma_p=v(\pi_p)\sum_{e\in I_p}p^{-e}$ {\cite[(18)]{P49}} and
$y-(y-y')^+=\min(y,y')$,
\[
\gamma_p\sum_\tau\kappa(\tau)-F(p)=v(\pi_p)\sum_\tau R_\tau\le\sum_\tau R_\tau,\qquad
R_\tau=\sum_{e\in I_p}\min\bigl(\kappa(\tau)/p^e,\kappa(u(-p^e)\tau)\bigr).
\]
So it suffices to show, for each $\tau$,
\begin{equation}\label{psm:eq-place}
3(\log b+2)+x^{42}+R_\tau-j(\tau)\le v\bigl(B(l_\tau)\bigr)
\end{equation}
(with $j(\tau)=0$ for $\tau\notin\auPl_0$). Let $n=\lceil l_0/2\rceil$. For $e<n$ (so $2e<l_0$) the
term of $R_\tau$ is: at most $\kappa(\tau)/p^{e}\le l_\tau/p^e<j(\tau)+1$ if $(\tau,e)$ is a near pair
(then $e=e(\tau)$); at most $a/x_1$ otherwise (if $\tau\in\auPb$ and $u(-p^e)\tau\notin\auPb$ the
minimum is at most $\kappa(u(-p^e)\tau)=a/x_1$; if $\tau\notin\auPb$, it is at most
$\kappa(\tau)=a/x_1$). For $e\ge n$ the terms are at most $b/p^e$, with sum at most $2b/p^n\le
2b/(kp)$, because $n>s$ and $p^{n-1}\ge2^s\ge k$ (N4). As $n\le l_0$,
\[
R_\tau\le j(\tau)+1+l_0\,a/x_1+2b/(kp).
\]
By (N2) and (N3) (with $p\le k$ and $l_0\le81s$), the left side of \eqref{psm:eq-place} is at most
$(l_0+1)a/x_1+2b/(kp)\le3b/(kp)$, and $b/(kp)\ge(l_0+1)a/x_1\ge1$. On the other hand
$v(B(l))\ge v(\psmns!)\ge\psmns/p-1\ge6b/(kp)-1\ge3b/(kp)$. This proves \eqref{psm:eq-place}.
\end{proof}

\begin{remark}
Theorem~\ref{psm:twelve} is the converse of {\cite[Proposition~14.10]{P46}} at the weights: the
couplings that produce the divisibility $s_p$ in the upper bound are exactly those whose leading
coefficients $\lambda_{e,0}$ are units, and choosing the exponent $j(\tau)$ of each coupling as
large as the degrees allow realises the prime sum up to the loss $2\sum_\tau v(B(l_\tau))$. For
primes $p>k$ all places are coupled at $e=1$; for $p\le k$ only near pairs of $b$-places can
contribute more than the loss, and sparseness makes them disjoint.
\end{remark}

\subsubsection{Transfer to the indicator data}\label{psm:ss-attained}

Let $W_{12}=\operatorname{Hom}(E,E^{12})$ with the standard hermitian structure, $L$ its standard
self-dual lattice, $\psvac$ the vacuum and $\psvac^\circ=c_{\mathrm{vac}}^{12}\psvac$ the integral
vacuum {\cite[(58)]{P49}}. For $\mathfrak m=(\pi_M)$ and a class $y_0\in(\auO_E/\mathfrak m)^{12}$
let $\phi_{y_0}$ be the \emph{indicator datum}: the vector of the Heisenberg model whose component
at the places dividing $\mathfrak m$ is the indicator of $y_0$ and which agrees with $\psvac$ at
all other places (including the Gaussian at infinity). For a vector $\Phi$ let
$\Theta_\Phi(\pi_Ms_0\,z)$, $z=(z_{\tau,b})_{\tau\in\auPl,\,1\le b\le12}$, be its theta function
at the scale of {\cite[(45)]{P49}}, and $[\gamma]$ its Taylor coefficient at $\gamma\in
\N^{\auPl\times12}$.

The proof uses the following four statements on the Heisenberg model of $W_{12}$; they are
proved in \S\ref{pst:sec} (Propositions~\ref{pst:I1}--\ref{pst:I4}). Here, as in \S\ref{pst:sec},
the value of a vector $\Phi$ at the finite adelic point $(0,x)$, $x\in W_{12}(\auA_f)$, is
$\Phi(0,x):=\theta(1\otimes\psrho_f(x)\Phi)$; it is linear in $\Phi$.
\begin{enumerate}
\item[(I1)] \emph{A frame.} There is a tight frame $(\mathsf w,c)$ of length $M$ with $c_i\ge6fb$
for all $i$ such that at every prime $p$ (and every $\C\to\C_p$) some form $\mu_{i_0}$ has unit
coefficients.
\item[(I2)] \emph{Products of monomial sections.} For exponents $e_b(i)\in\N$ $(1\le b\le12$,
$i\in J)$ with $3e_b(i)\le c_i$ there is a vector $\Phi$ with $\Theta_\Phi(\pi_Ms_0z)=\prod_{b}
\prod_iy(\mu_i(z_b))^{e_b(i)}\theta(\mu_i(z_b))^{c_i}$, whose values $\Phi(0,x)$,
$x\in W_{12}(\auA_f)$, are integral at every prime, and with $\psrho^\circ(l)\Phi=\Phi$ for $l\in\pi_ML$.
\item[(I3)] \emph{Fourier expansion.} For $\Phi$ as in (I2), $M^{12f}\Phi=\sum_ia_i\,\psrho(u_i)
\psvac$, a finite sum with $\aubar{\pi_M}u_i\in\mathfrak a_0^{12}$, such that $a_i\psvac(0,x)$
is integral for every $x\in W_{12}(\auA_f)$; and each $\psrho(u)\psvac$ with $\aubar{\pi_M}u$ integral
is a sum of the $\phi_{y_0}$ with roots of unity as coefficients.
\item[(I4)] \emph{The vacuum constant.} There is a finite adelic $y$ with $c_{\mathrm{vac},y}:=\psvac^\circ(0,
y)=\theta(1\otimes\psrho_f(y)\psvac^\circ)\ne0$, and $|\sigma c_{\mathrm{vac},y}|\le
\exp(12C_0q_1^2)\le\exp(b)$ for all $\sigma\in\auAut(\C)$.
\end{enumerate}

\begin{theorem}[Sharpness of the count]\label{psm:attained}
Let $l\in\N^{\auPl}$ with $\kappa(\tau)\le l_\tau\le\kappa(\tau)+x^{42}$.
There is $t\in\C^\times$ with $|\sigma t|\le\exp(fmb/k)$ for all $\sigma\in\auAut(\C)$ such that
for every prime $p$ and every $\C\to\C_p$ there are $\gamma\in\N^{\auPl\times12}$ with
$\sum_b\gamma(\tau,b)=l_\tau$ for all $\tau$, and a class $y_0$ modulo $\pi_M$, with
\[
p^{-F(p)}\ \le\ \Bigl|\,\frac{c_{\mathrm{vac}}^{12}}t\,\prod_\tau l_\tau!\cdot[\gamma]\,
\Theta_{\phi_{y_0}}(\pi_Ms_0\,\cdot)\Bigr|_p .
\]
\end{theorem}
\begin{proof}
Let $(\mathsf w,c)$ be the frame of (I1), $y$ and $c_{\mathrm{vac},y}$ as in (I4), and put
\[
t=M^{12f}c_{\mathrm{vac},y}\prod_\tau B(l_\tau)^2 .
\]
\emph{Size of $t$.} $B(l)\le(l^-+\psmns)^{\psmns}\le(2b+\psmns)^{\psmns}$ by (N1), $\psmns\le
7b/k$, $\log(2b+\psmns)\le2\log b$ and $\log M\le22\log b$ ({\cite[Lemma~13.4(S5)]{P46}}); so
$\log|\sigma t|\le264f\log b+28fb\log b/k+b$. Since $k\le3x^5\le f$, $b\le fb/k$; since
$b/k\ge b/(3x^5)\ge264\log b$, $264f\log b\le fb/k$; and $m\ge\sqrt x\ge(\log b)^{50}\ge
28\log b+2$. Hence $\log|\sigma t|\le fmb/k$.

\emph{The coefficient.} Fix $p$ and $\C\to\C_p$, and let $r=F(p)+2\sum_\tau v(B(l_\tau))-
\sum_\tau v(l_\tau!)$. By Theorem~\ref{psm:twelve} (with $i_0$ from (I1)) and
Lemma~\ref{psm:rules}(d) there are monomials $H_b=\prod_iy(\mu_i)^{e_b(i)}$ with $\sum_ie_b(i)\le
|l|\le2fb$, hence $3e_b(i)\le6fb\le c_i$, and $\alpha_1+\dots+\alpha_{12}=l$, with $\sum_bv([\alpha_b]
(H_b\Theta))\le r$. Let $\Phi$ be the vector of (I2) for these exponents and
$\gamma(\tau,b)=\alpha_b(\tau)$; then $\sum_b\gamma(\tau,b)=l_\tau$, and the coefficient
$A=[\gamma]\Theta_\Phi(\pi_Ms_0\,\cdot)=\prod_b[\alpha_b](H_b\Theta)$ has $v(A)\le r$. Let
$\Lambda(\Phi')=c_{\mathrm{vac}}^{12}[\gamma]\Theta_{\Phi'}(\pi_Ms_0\,\cdot)$, a linear form. By (I3),
\[
M^{12f}\Lambda(\Phi)\,\psvac(0,y)=\sum_i\sum_{y_0}\zeta_{i,y_0}\,\bigl(a_i\psvac(0,y)
\bigr)\,\Lambda(\phi_{y_0})
\]
with roots of unity $\zeta_{i,y_0}$ and integral $a_i\psvac(0,y)$. Since $\psvac^\circ=c_{\mathrm{vac}}^{12}\psvac$ and the value at $(0,y)$ is linear,
$c_{\mathrm{vac}}^{12}\psvac(0,y)=\psvac^\circ(0,y)=c_{\mathrm{vac},y}$. Choose $y_0$ with
$v(\Lambda(\phi_{y_0}))$ least. The ultrametric inequality gives
\[
v\bigl(\Lambda(\phi_{y_0})\bigr)\le v\bigl(M^{12f}c_{\mathrm{vac}}^{12}A\,\psvac(0,y)\bigr)=
v(A)+v(M^{12f}c_{\mathrm{vac},y})\le r+v(M^{12f}c_{\mathrm{vac},y}).
\]
Therefore
\[
v\Bigl(\frac{c_{\mathrm{vac}}^{12}}t\prod_\tau l_\tau!\,[\gamma]\Theta_{\phi_{y_0}}(\pi_Ms_0\,\cdot)
\Bigr)=v\bigl(\Lambda(\phi_{y_0})\bigr)-v(M^{12f}c_{\mathrm{vac},y})-2\sum_\tau v(B(l_\tau))+\sum_\tau v(l_\tau!)
\le F(p),
\]
which is the claim.
\end{proof}

\emph{Application in Proposition~\ref{psm:small}.} For a main pair $p=(\nu,z)$ the degrees
$l_\tau=c(\tau)-z$ satisfy $\kappa(\tau)\le l_\tau\le\kappa(\tau)+x^{42}$, so Theorem~\ref{psm:attained} applies; this is part (i) of
the note to Proposition~\ref{po:small}, with the exponent $F(\ell)$ of \eqref{psm:eq-F}. The
number $t$ is the first factor of $c_0=t\,I_{q_2}(1)\,L$, and $|\sigma t|\le\exp(fmb/k)$ is the
first term of the bound $\exp(fmb/k+3x^{35})$.




\providecommand{\auO}{\mathcal O}
\providecommand{\auA}{\mathbb A}
\providecommand{\auPl}{\Phi}
\providecommand{\auNm}{\operatorname{Nm}}
\providecommand{\auAut}{\operatorname{Aut}}
\providecommand{\auHom}{\operatorname{Hom}}
\providecommand{\aubar}[1]{\overline{#1}}
\providecommand{\pao}{\mathfrak o}
\providecommand{\pam}{\mathfrak m}
\providecommand{\psvac}{\mathbf v}
\providecommand{\psrho}{\varrho}
\providecommand{\psHS}{\mathcal H}
\providecommand{\pfT}{\operatorname{Tr}}
\providecommand{\pfbr}[1]{\theta[#1]}
\providecommand{\pstS}{\mathcal S}

\subsection{The four inputs of the sharpness of the count}\label{pst:sec}

This section proves the statements (I1)--(I4) that were isolated before
Theorem~\ref{psm:attained}; with them, Theorem~\ref{psm:attained} is proved in full. The notation is
that of \S\ref{psm:ss-twelve} and \S\ref{psm:ss-attained}: $W_{12}=\auHom(E,E^{12})=E^{\mathsf C}$
with $\mathsf C=[12]\times[1]$, $L=\alpha\widehat\auO_E^{\,\mathsf C}$, $\Lambda_L=\mathfrak
a_0^{\mathsf C}$, the coordinates $\gamma^\natural=(c_\psi\tau(\gamma_b))_{\tau,b}$ of
{\cite[\S5.2.8]{P49}}, $\Lambda=\Omega_0\auO_K$ the normalised lattice of {\cite[(45)]{P49}}, with theta
function $\theta$, hermitian form $H$ and multiplier $\chi$ {\cite[(27)]{P49}}, and $\mathsf s=
\varpi_{\mathrm{vac}}/\mathsf u$. Recall $\pfbr u=e^{-H(u,u)/2}\theta(u)$ and the spaces $\pstS_m$
of {\cite[(71)]{P49}}. For $G\in\pstS_m$ and $u\in\Q\Lambda$ we call $e^{-mH(u,u)/2}G(u)$ the
\emph{torsion value} of $G$ at $u$; for $G=\theta^m$ it is $\pfbr u^m$.

\subsubsection{(I1): a frame of length $M$ with a unit form at every prime}\label{pst:ss-frame}

\begin{lemma}[Size of the frame of length prime to $q_1$]\label{pst:hsize}
The length $\mathsf h$ of the tight frame of {\cite[Lemma~7.11]{P49}}, built with the least admissible
exponent $J_0$ (the exponent $J$ of that proof), satisfies $\mathsf h\le b$.
\end{lemma}
\begin{proof}
Let $\zeta=\zeta_{q_1}$, $T=q_1/q_0$. Every prime dividing $q_1$ divides $q_0$
({\cite[Lemma~13.4(S7)]{P46}}), so the minimal polynomial of $\zeta$ over $K$ is
$P(X)=\Phi_{q_1}(X)=\Phi_{q_0}(X^{T})$ ($\Phi_n$ the cyclotomic polynomials; $K\cap\Q(\zeta)=\Q$
since $d\nmid q_1$). The roots of $\Phi_{q_0}$ have absolute value $1$, so its coefficients are
bounded by binomial coefficients $\binom{\varphi(q_0)}{k}\le2^{q_0}$, and so are those of $P$.
Hence, for every $\tau\in\auPl$: (i) the coefficients $b_i=\sum_{j>i}P_j\zeta^{j-i-1}$ of
$P(X)/(X-\zeta)$ satisfy $|\tau(b_i)|\le q_1 2^{q_0}$; (ii) $\tau(P'(\zeta))=P'(\eta')$ with
$\eta'=\tau(\zeta)$, and $P'(\eta')=T\eta'^{\,T-1}\Phi_{q_0}'(\eta'^{\,T})$, where $\eta'^{\,T}$ is
a primitive $q_0$-th root of unity. The number $\Phi_{q_0}'(\eta)$, $\eta$ primitive of order $q_0$, is
a non-zero algebraic integer of $\Q(\zeta_{q_0})$ whose conjugates $\Phi_{q_0}'(\eta'')$ have absolute
value at most $q_0^22^{q_0}$; its norm is a non-zero integer, so $|\Phi'_{q_0}(\eta)|\ge(q_0^22^{q_0}
)^{-\varphi(q_0)}\ge e^{-3q_0^2}$. Hence $|\tau(P'(\zeta))|\ge e^{-3q_0^2}$.

Every prime $p\mid q_1$ divides $q_0$ ({\cite[Lemma~13.4(S7)]{P46}}), hence $4p\mid d+1$ ({\cite[Definition~13.1]{P46}} and the remark after it), so $p$ divides $\omega\aubar\omega=(1+d)/4$; hence the least $J_0$ with $q_1\mid(\omega
\aubar\omega)^{J_0}$ is at most $\max_pv_p(q_1)\le\log_2q_1$; and $|\omega|^2=(1+d)/4\le d$. With
$e_i=\aubar\omega^{\,J_0}b_i/P'(\zeta)$ as in the proof of {\cite[Lemma~7.11]{P49}},
\[
|\tau(e_i)|^2\le d^{J_0}q_1^24^{q_0}e^{6q_0^2}=:B,\qquad\nu(e_i)=\sum_\tau|\tau(e_i)|^2\le fB,
\]
and $\mathsf h=q_1(1+\sum_{i<f}\nu(e_i))+\omega^{J_0}+\aubar\omega^{\,J_0}\le q_1(1+f^2B)+2d^{J_0/2}\le
4q_1^3B$, using $f\le q_1$. By {\cite[Lemma~13.4(S1)]{P46}},(S3), $q_1\le x^{12}$ and $\log x\ge
10^{14}(q_0+2d)$; so $J_0\le18\log x$, $\log d\le d\le10^{-14}\log x$, $q_0\le10^{-14}\log x$, and
\[
\log\mathsf h\le\log4+36\log x+18\log d\log x+24\log x+2q_0+6q_0^2\le(\log x)^2 .
\]
By (S1), $\log b\ge(\log x_1)^6\ge(\log x)^6$. So $\mathsf h\le b$.
\end{proof}

\begin{proposition}[(I1)]\label{pst:I1}
There is a tight frame $(\mathsf w_i,c_i)_{i\in J}$ of length $M$ with $c_i\ge6fb$ for all $i$ such
that, for every prime $p$ and every $\C\to\C_p$, some $i_0$ has the property that all coefficients
$\mathsf s\,\aubar{\tau(\mathsf w_{i_0})}$ $(\tau\in\auPl)$ of $\mu_{i_0}$ are units. The members
$\mathsf w_i$ do not depend on $M$; in particular this holds at $p_0$ and $p_1$.
\end{proposition}
\begin{proof}
\emph{The frame.} Let $(v_j)_{j\in\mathsf I}$ be the frame of {\cite[Lemma~7.11]{P49}}, of length
$\mathsf h$ prime to $q_1$ and $\le b$ (Lemma~\ref{pst:hsize}), and $(\zeta^k)_{k<q_1}$ the frame of
length $q_1$. Let $R=\lceil6fb\rceil$. Since $\gcd(q_1,\mathsf h)=1$, every integer
$\ge(q_1-1)(\mathsf h-1)$ is of the form $\alpha q_1+\beta\mathsf h$ with $\alpha,\beta\in\N$.
By {\cite[Lemma~13.4(S3)]{P46}},(S5), $q_1\le x^{12}\le b$, $R\le7b^2$ and $M\ge b^{20}$, so
$M-R(q_1+\mathsf h)\ge b^{20}-14b^3\ge q_1\mathsf h$; write $M-R(q_1+\mathsf h)=\alpha q_1+\beta
\mathsf h$. Take the members $\zeta^k$ $(k<q_1)$ with multiplicity $\alpha+R$ and the members $v_j$
with multiplicity $\beta+R$. A disjoint union of tight frames, with multiplicities, is a tight frame
whose length is the sum of the lengths; so this is a tight frame of length $(\alpha+R)q_1+(\beta+R)
\mathsf h=M$, all multiplicities are $\ge R\ge6fb$, and all members lie in $\mathfrak a_0$.

\emph{The unit form.} Fix $p$ and $\iota\colon\C\to\C_p$, and for $\tau\in\auPl$ let $w_\tau$ be
the place of $E$ induced by $\iota\circ\tau\circ c$, so that $|\iota(\aubar{\tau(\mathsf v)})|=|\mathsf
v|_{w_\tau}$. By {\cite[Lemma~7.12]{P49}} and {\cite[Corollary~7.3]{P49}} ($\mathsf u$ a unit),
\[
|\iota(\mathsf s\,\aubar{\tau(\mathsf v)})|=|\mathsf v|_{w_\tau}/|\alpha|_{w_\tau}\qquad(\mathsf v
\in E^\times).
\]
The places $w_\tau$ lie above one place $\mathfrak q$ of $K$. If $\omega\notin\mathfrak q$, then
$\alpha_{w_\tau}=1$ for all $\tau$ ({\cite[\S5.2.1]{P49}}), and the member $\mathsf w_{i_0}=\zeta^0=1$
works. If $\omega\in\mathfrak q$, then $\alpha_{w_\tau}=1/P'(\zeta)$ for all $\tau$, and
$\aubar\omega\notin\mathfrak q$ ($\omega+\aubar\omega=1$). Take $\mathsf w_{i_0}=\zeta e_{f-1}$; it is
a member of the frame of {\cite[Lemma~7.11]{P49}} (the member $\zeta^k\rho(e_{f-1})$ with $k=1$, $\rho=1$,
which is not one of the replaced pairs, as $q_1>1$). Since $P$ is monic of degree $f$, $b_{f-1}=1$
and $\zeta e_{f-1}=\zeta\aubar\omega^{\,J_0}/P'(\zeta)$; as $\zeta$ and $\aubar\omega$ are units at
$w_\tau$, $|\zeta e_{f-1}|_{w_\tau}=|1/P'(\zeta)|_{w_\tau}=|\alpha|_{w_\tau}$. In both cases all
coefficients of $\mu_{i_0}$ are units. Nothing here depends on whether $p$ divides $M$.
\end{proof}

\subsubsection{(I2): products of monomial sections}\label{pst:ss-product}

\begin{lemma}[Torsion values of monomial sections]\label{pst:torsion}
Let $e,c\in\N$ with $3e\le c$ and $G=y^e\theta^c$. Then $G$ is entire, $G\in\pstS_c$, and for every
$u\in\Q\Lambda$, every prime $\ell$ and every $\iota$ the torsion value $e^{-cH(u,u)/2}G(u)$ is an
algebraic integer.
\end{lemma}
\begin{proof}
The function $y$ has poles of order $3$ exactly at $t+\Lambda$, where $\theta$ has simple zeros; so
$G$ is entire, and $G\in\pstS_c$ by {\cite[(71)]{P49}} (the polynomial $y^e$ has a pole of
order $3e\le c$ at $O$). Let $u\in\Q\Lambda$. The torsion value is algebraic: $Q(u)$ is a torsion point of the curve $y^2+a_1xy=x^3+x$, defined over $\Q(a_1)\subset\overline\Q$, so $x(u),y(u)\in\overline\Q$, and $\pfbr u$ is algebraic by {\cite[Proposition~7.5]{P49}}; it remains to show it is integral at every $\ell$ and $\iota$.

If $u\notin t+\Lambda$, the torsion value is $y(u)^e\pfbr u^{\,c}$. If $u-t\notin E_1$, then
$\pfbr u$ is a unit ({\cite[Proposition~7.5]{P49}}) and $Q(u)$ does not reduce to the origin of the
smooth model ({\cite[Lemma~7.2]{P49}}), so $y(u)$ is integral. If $u-t\in E_1$, then $|x(u)|>1$ and
$|\pfbr u|=|x(u)|^{-1/2}$ ({\cite[Proposition~7.5]{P49}}); from $y^2+a_1xy=x^3+x$ with $a_1$
integral, $|y(u)|>|x(u)|$ (otherwise the left side would have absolute value $\le|x(u)|^2<
|x(u)|^3$), hence $|y(u)|^2=|x(u)|^3$; so $|y(u)^e\pfbr u^{\,c}|=|x(u)|^{(3e-c)/2}\le1$.

If $u\in t+\Lambda$ and $c>3e$ the value is $0$. Let $c=3e$, so $G=(y\theta^3)^e$. For $u=t+\lambda$,
$\lambda\in\Lambda$, the functional equation of $\pstS_c$ gives $e^{-cH(u,u)/2}G(u)=\chi(\lambda)^c
e^{ic\,\Iml H(\lambda,t+\lambda/2)}\cdot e^{-cH(t,t)/2}G(t)$, and the middle factor is a root of
unity, since $\Iml H$ is $\boldsymbol\pi$ times an integral alternating form on $\Lambda$ and
$2t,\lambda\in\Lambda$. So it suffices to treat $u=t$. As $w\to0$: $x(t+w)+\wp(t)=\wp(w)$, so
$x(t+w)\sim w^{-2}$ and $y(t+w)=\tfrac12\wp'(w)-\tfrac12a_1x(t+w)\sim-w^{-3}$; and $e^{-H(t+w,t+w)/2}
\theta(t+w)=\pfbr{t+w}=e^{-\boldsymbol\pi iE(w,t)}e^{-H(w,w)/2}\theta_t(w)\sim\theta_t'(0)w$, because
$\theta_t(0)=\pfbr t=0$. Hence $e^{-3H(t,t)/2}(y\theta^3)(t)=-\theta_t'(0)^3$, and $\theta_t'(0)$ is a
unit at every prime by the proof of {\cite[Corollary~7.16(c)]{P49}}.
\end{proof}

Let $(\mathsf w_i,c_i)_{i\in J}$ be a tight frame of length $M$, and put, for $z=(z_{\tau,b})\in
\C^{\auPl\times\mathsf C}$ and $b\in[12]$,
\[
a_{i}(z_b)=\frac{\Omega_0}{\pi_Mc_\psi}\sum_{\tau\in\auPl}\aubar{\tau(\mathsf w_i)}\,z_{\tau,b} ,
\qquad\text{so that}\qquad a_i\bigl((\pi_Ms_0z)_b\bigr)=\mathsf s\sum_\tau\aubar{\tau(\mathsf w_i)}
z_{\tau,b}=\mu_i(z_b)
\]
by {\cite[(45)]{P49}} ($s_0\Omega_0/c_\psi=\mathsf s$). For $\gamma\in W_{12}(E)$,
$a_i(\gamma^\natural_b)=\frac{\Omega_0}{\pi_M}\pfT_{E/K}(\aubar{\mathsf w_i}\gamma_b)$, since the
$\tau\in\auPl$ fix $K$ and commute with $c$.

\begin{lemma}[Isometry]\label{pst:iso}
$\sum_{b}\sum_ic_i\,H\bigl(a_i(w_b),a_i(z_b)\bigr)=\langle z,w\rangle=\sum_{\tau,b}\aubar{w_{\tau,b}}
z_{\tau,b}$ for all $z,w$.
\end{lemma}
\begin{proof}
$H(v,v')=\boldsymbol\pi\aubar vv'/A$ with $A=|\Omega_0|^2\sqrt d/2$, and $\boldsymbol\pi|\Omega_0|^2
/(Ac_\psi^2)=1$ ({\cite[\S5.2.8]{P49}}). So the left side is $M^{-1}\sum_b\sum_{\tau,\tau'}\aubar{w_{
\tau,b}}z_{\tau',b}\sum_ic_i\tau(\mathsf w_i)\aubar{\tau'(\mathsf w_i)}$, and the frame identity gives
$\sum_ic_i\tau(\mathsf w_i)\aubar{\tau'(\mathsf w_i)}=M\delta_{\tau\tau'}$.
\end{proof}

\begin{proposition}[(I2)]\label{pst:I2}
Let $e_b(i)\in\N$ with $3e_b(i)\le c_i$, and $G_{i,b}=y^{e_b(i)}\theta^{c_i}\in\pstS_{c_i}$. There is a
unique $\Phi\in\psHS_f(W_{12})$ with
\[
\Theta_\Phi(z)=\prod_{b\in[12]}\prod_{i\in J}G_{i,b}\bigl(a_i(z_b)\bigr),\qquad\text{hence}\qquad
\Theta_\Phi(\pi_Ms_0z)=\prod_b\prod_iy(\mu_i(z_b))^{e_b(i)}\theta(\mu_i(z_b))^{c_i}.
\]
Moreover {\rm(a)} $\psrho_f(l)\Phi=\epsilon_L(l)\Phi$ for all $l\in\pi_ML$, that is,
$\psrho^\circ(l)\Phi=\Phi$ (notation of {\cite[Lemma~7.14]{P49}}); {\rm(b)} for every $x\in
W_{12}(\auA_f)$ the value $\theta(1\otimes\psrho_f(x)\Phi)$ is an algebraic integer.
\end{proposition}
\begin{proof}
Write $S(z)$ for the product. Let $\lambda=\pi_M\lambda'$ with $\lambda'\in\Lambda_L$. Then
$\omega_{i,b}:=a_i(\lambda^\natural_b)=\Omega_0\nu_{i,b}$ with $\nu_{i,b}=\pfT_{E/K}(\aubar{\mathsf
w_i}\lambda'_b)\in\pfT_{E/K}(\aubar{\mathfrak a_0}\mathfrak a_0)=\pfT_{E/K}(\mathfrak d_{E/K}^{-1})
\subseteq\auO_K$, so $\omega_{i,b}\in\Lambda$. The functional equation of $\pstS_{c_i}$ and
Lemma~\ref{pst:iso} give
\[
S(z+\lambda^\natural)=\epsilon(\lambda)\,e^{\langle z,\lambda^\natural\rangle+|\lambda^\natural|^2/2}
S(z),\qquad\epsilon(\lambda)=\prod_{i,b}\chi(\Omega_0\nu_{i,b})^{c_i}=(-1)^{\sum_{i,b}c_i\auNm
\nu_{i,b}}
\]
by {\cite[Lemma~7.1(b)]{P49}}. Now $\sum_ic_i\auNm\nu_{i,b}=\sum_ic_i|\sum_\tau\aubar{\tau(\mathsf w_i)}
\tau(\lambda'_b)|^2=M\sum_\tau|\tau(\lambda'_b)|^2$ by the frame identity, and summing over $b$,
$\sum_{i,b}c_i\auNm\nu_{i,b}=M\pfT_{E/K}\langle\lambda',\lambda'\rangle=\pfT_{E/K}\langle\lambda,
\lambda\rangle$, as $\pi_M\aubar{\pi_M}=M$. By {\cite[Lemma~7.4]{P49}}, $\epsilon(\lambda)=\epsilon_L(
\lambda)$ on $\pi_M\Lambda_L$. In particular $S$ satisfies the functional equation of
{\cite[Lemma~5.26]{P49}} with $\mathsf N=M$ ($M\Lambda_L\subseteq\pi_M\Lambda_L$), so $S=\Theta_\Phi$ for a
unique $\Phi$; the second formula follows from $a_i((\pi_Ms_0z)_b)=\mu_i(z_b)$.

(a) As at the end of the proof of {\cite[Lemma~7.13]{P49}}: the computation in the proof of
{\cite[Lemma~5.26]{P49}} shows that $\Theta_{\psrho_f(\lambda)\Phi}(z)=e^{\langle z,\lambda^\natural
\rangle-|\lambda^\natural|^2/2}\Theta_\Phi(z-\lambda^\natural)$ for $\lambda\in\Lambda_L$ (the root
of unity there is $1$ for lattice vectors, as in the proof of {\cite[Lemma~7.13]{P49}}), which
by the functional equation just proved (applied at $z-\lambda^\natural$) is $\epsilon_L(\lambda)^{-1}\Theta_\Phi(z)=\epsilon_L(\lambda)\Theta_\Phi(z)$ for $\lambda\in\pi_M
\Lambda_L$; so $\psrho_f(\lambda)\Phi=\epsilon_L(\lambda)\Phi$ by injectivity
($\epsilon_L=\epsilon_L^{-1}$); it extends to $\pi_ML$ by continuity, $\pi_M\Lambda_L$ being dense in
$\pi_ML$ and $\epsilon_L$ locally constant.

(b) Since $W_{12}(E)+\pi_ML=W_{12}(\auA_f)$, by (a) $\psrho_f(x)\Phi$ is a root of unity times
$\psrho_f(\gamma)\Phi$ for some $\gamma\in W_{12}(E)$. By {\cite[Lemma~5.26]{P49}} at $z=0$,
$\theta(1\otimes\psrho_f(\gamma)\Phi)=\Theta_{\psrho_f(\gamma)\Phi}(0)$ is a root of unity times
$e^{-|\gamma^\natural|^2/2}S(-\gamma^\natural)$. By Lemma~\ref{pst:iso}, $|\gamma^\natural|^2=
\sum_{i,b}c_iH(u_{i,b},u_{i,b})$ with $u_{i,b}=-a_i(\gamma^\natural_b)=-\frac{\Omega_0}{\pi_M}
\pfT_{E/K}(\aubar{\mathsf w_i}\gamma_b)\in K\Omega_0=\Q\Lambda$. Hence
\[
e^{-|\gamma^\natural|^2/2}S(-\gamma^\natural)=\prod_{i,b}e^{-c_iH(u_{i,b},u_{i,b})/2}G_{i,b}(u_{i,b}),
\]
a product of torsion values of monomial sections, which are algebraic integers by
Lemma~\ref{pst:torsion}.
\end{proof}

\emph{The scale.} The Taylor series are taken at the scale $\pi_Ms_0$, and the scale of the linear
forms is $\mathsf s=\varpi_{\mathrm{vac}}/\mathsf u$. This confirms the identification of
\S\ref{psm:ss-twelve}: the forms $\mu_i$ there are those of Proposition~\ref{pst:I2}.

\subsubsection{(I3): Fourier expansion in translates of the vacuum}\label{pst:ss-fourier}

\begin{proposition}[(I3)]\label{pst:I3}
Let $\Phi$ be as in Proposition~\ref{pst:I2}, and let $Y\subset\aubar{\pi_M}^{\,-1}\Lambda_L$ be a
set of representatives of $\aubar{\pi_M}^{\,-1}L/L$. Then $\#Y=M^{12f}$ and there are $a_y\in\C$
$(y\in Y)$ with
\begin{enumerate}
\item[(a)] $M^{12f}\Phi=\sum_{y\in Y}a_y\,\psrho_f(y)\psvac$;
\item[(b)] $a_y\,\theta(1\otimes\psrho_f(x)\psvac)$ is an algebraic integer for every $y\in Y$ and
every $x\in W_{12}(\auA_f)$.
\end{enumerate}
Further, for every $u\in W_{12}(\auA_f)$ with $\aubar{\pi_M}u\in L$, $\psrho_f(u)\psvac=\sum_{y_0}
\zeta_{u,y_0}\phi_{y_0}$ with roots of unity $\zeta_{u,y_0}$, the sum over the classes $y_0\in
(\auO_E/\pi_M\auO_E)^{12}$.
\end{proposition}
\begin{proof}
\emph{Commutators.} For $x,x'\in W_{12}(\auA_f)$ the commutator $[x,x']=\psi(\jmath\langle x',x
\rangle)$ is trivial when $\langle x',x\rangle\in\jmath^{-1}\mathfrak d_E^{-1}\widehat\auO_E$, in
particular for $x,x'\in L$ (self-duality). If $l=\pi_Ml'$ and $y=\aubar{\pi_M}^{\,-1}y'$ with
$l',y'\in L$, then $\langle l,y\rangle=\aubar{\pi_M}\,\aubar{\pi_M}^{\,-1}\langle l',y'\rangle=
\langle l',y'\rangle$, so $[l,y]=1$; and the pairing $(l,y)\mapsto[l,y]$ identifies $L/\pi_ML$ with
the dual of $\aubar{\pi_M}^{\,-1}L/L$, by self-duality of $L$ (the annihilator of
$\aubar{\pi_M}^{\,-1}L$ in $L$ is $\pi_ML$). The index is $\#Y=[\aubar{\pi_M}^{\,-1}L:L]=
|\auNm_{E/\Q}\aubar{\pi_M}|^{12}=M^{12f}$. Representatives can be taken in $\aubar{\pi_M}^{\,-1}
\Lambda_L$ by strong approximation ($W_{12}(E)$ is dense in $W_{12}(\auA_f)$).

\emph{The coefficients.} Let $\mathsf P=\sum_{\gamma\in L/\pi_ML}\psrho^\circ(\gamma)$; it is well
defined on the space $\mathsf V$ of vectors fixed by $\psrho^\circ(\pi_ML)$, because $\psrho^\circ$ is
a homomorphism on $L$ ({\cite[\S7.4]{P49}}). For $y\in Y$, $\psrho_f(-y)\Phi\in\mathsf V$ by
Proposition~\ref{pst:I2}(a) and $[\pi_ML,y]=1$; and $\mathsf P\psrho_f(-y)\Phi$ is fixed by
$\psrho^\circ(L)$, hence, by (W4), equal to $a_y\psvac$ for some $a_y\in\C$. Now
$\psrho_f(y)\psrho^\circ(\gamma)\psrho_f(-y)=[y,\gamma]^{\pm1}\psrho^\circ(\gamma)$ (the sign of the
exponent depends only on the convention for $[\,,]$), so
\[
\sum_{y\in Y}a_y\psrho_f(y)\psvac=\sum_{\gamma\in L/\pi_ML}\Bigl(\sum_{y\in Y}[y,\gamma]^{\pm1}\Bigr)
\psrho^\circ(\gamma)\Phi=\#Y\cdot\Phi ,
\]
since $y\mapsto[y,\gamma]$ is a character of $\aubar{\pi_M}^{\,-1}L/L$, trivial exactly for
$\gamma\in\pi_ML$, where $\psrho^\circ(\gamma)\Phi=\Phi$. This is (a). For (b):
$a_y\theta(1\otimes\psrho_f(x)\psvac)=\sum_\gamma\theta(1\otimes\psrho_f(x)\psrho^\circ(\gamma)
\psrho_f(-y)\Phi)$, and each term is a root of unity times $\theta(1\otimes\psrho_f(x+\gamma-y)
\Phi)$, an algebraic integer by Proposition~\ref{pst:I2}(b).

\emph{Indicator data.} Let $u$ with $\aubar{\pi_M}u\in L$. The ideal $\pi_M\auO_K=\mathfrak p_0
\mathfrak p_1$ is a product of primes of degree one above $p_0,p_1$, which split in $K$; so $\pi_M$
is a unit at $\aubar{\mathfrak p_0},\aubar{\mathfrak p_1}$, $\aubar{\pi_M}$ is a unit at
$\mathfrak p_0,\mathfrak p_1$, and both are units at all other places. Hence $u\in L$ at every
place not above $\aubar{\mathfrak p_0}$ or $\aubar{\mathfrak p_1}$. At
$p=p_i$ write $W_{12}(\Q_p)=W_{\mathfrak p}\oplus W_{\aubar{\mathfrak p}}$ (the components of
$E\otimes\Q_p=E_{\mathfrak p}\times E_{\aubar{\mathfrak p}}$); both are isotropic for $\langle\,,
\rangle$, because $c$ exchanges the two factors, and $L_p=L_{\mathfrak p}\oplus L_{\aubar{\mathfrak
p}}$. In the Schr\"odinger model on $\mathcal S(W_{\mathfrak p})$, $\psrho(a)$ acts by translation
for $a\in W_{\mathfrak p}$, $\psrho(a')$ by multiplication by $r\mapsto\psi(\jmath\langle r,a'
\rangle)^{\pm1}$ for $a'\in W_{\aubar{\mathfrak p}}$, the local vacuum is $\mathbf 1_{L_{\mathfrak
p}}$ and the local component of $\phi_{y_0}$ is $\mathbf 1_{y_0+\pi_ML_{\mathfrak p}}$ (with $L_{\mathfrak
p}/\pi_ML_{\mathfrak p}=(\auO_E/\mathfrak p\auO_E)^{12}$). Write $u_p=u_{\mathfrak p}+u_{\aubar{
\mathfrak p}}$; then $u_{\mathfrak p}\in L_{\mathfrak p}$ and $u_{\aubar{\mathfrak p}}\in
\aubar{\pi_M}^{\,-1}L_{\aubar{\mathfrak p}}$. Up to a root of unity (a value of $\beta$),
$\psrho(u_p)\mathbf 1_{L_{\mathfrak p}}=\psrho(u_{\aubar{\mathfrak p}})\mathbf 1_{L_{\mathfrak
p}}$ is the function $r\mapsto\psi(\jmath\langle r,u_{\aubar{\mathfrak p}}\rangle)^{\pm1}\mathbf
1_{L_{\mathfrak p}}(r)$; for $r\in\pi_ML_{\mathfrak p}$ the pairing is $\langle\pi_M^{-1}r,
\aubar{\pi_M}u_{\aubar{\mathfrak p}}\rangle$ with both arguments integral, so the character is
invariant under $\pi_ML_{\mathfrak p}$, and the function is $\sum_{y_0}\zeta_{y_0}\mathbf 1_{y_0+
\pi_ML_{\mathfrak p}}$ with roots of unity $\zeta_{y_0}$. At every other place $\psrho(u_w)$ fixes
the local vacuum up to a root of unity ($u_w\in L_w$). Taking the tensor product gives the claim.
\end{proof}

\subsubsection{(I4): the vacuum constant}\label{pst:ss-vac}

\begin{proposition}[(I4)]\label{pst:I4}
There is $y\in W_{12}(\auA_f)$ with $c_{\mathrm{vac},y}:=\theta(1\otimes\psrho_f(y)\psvac^\circ)\ne0$, and for every
such $y$ and every $\sigma\in\auAut(\C)$, $|\sigma c_{\mathrm{vac},y}|\le\exp(12C_0q_1^2)\le e^b$, where $C_0=
C_0(q_0,d)$ is the constant {\cite[(76)]{P49}}.
\end{proposition}
\begin{proof}
\emph{Non-vanishing.} Suppose $\theta(1\otimes\psrho_f(y)\psvac)=0$ for all $y\in W_{12}(\auA_f)$.
For $\gamma\in W_{12}(E)$, $\theta\circ\psrho(\gamma)=\theta$, so $\theta(\psrho_\infty(\gamma)
\psrho_f(y)(1\otimes\psvac))$ is a root of unity times $\theta(1\otimes\psrho_f(y-\gamma)\psvac)=0$.
The image of $W_{12}(E)$ in $W_{12}\otimes\R$ is dense and $w\mapsto\theta(\psrho(w)(1\otimes
\psvac))$ is continuous, so it vanishes on $W_{12}(\auA)$, and $1\otimes\psvac$ lies in the kernel of
$F\mapsto\theta(\psrho(\cdot)F)$, a proper closed invariant subspace of an irreducible
representation (as in the proof of {\cite[Lemma~5.26]{P49}}); this is absurd. Since
$\psvac^\circ=c_{\mathrm{vac}}^{12}\psvac$ with $c_{\mathrm{vac}}\ne0$, some $c_{\mathrm{vac},y}\ne0$.

\emph{Rank one.} Let $\psvac^\circ_1$ be the integral vacuum of rank one {\cite[(58)]{P49}}, a sum
of $q_1^{2f}=[L_1:q_1L_1]$ vectors $\zeta_l\psrho(l)\psrho(b)F_0$ with roots of unity $\zeta_l$, where
$F_0$ is the theta product {\cite[(56)]{P49}} with $\#\mathsf C=1$. For $x\in E(\auA_f)$,
$\theta(1\otimes\psrho_f(x)\psrho(l)\psrho(b)F_0)$ is a root of unity times $\theta(1\otimes
\psrho_f(\gamma)F_0)$ for a $\gamma\in E$ ({\cite[\S5.2.8]{P49}}), which by {\cite[(57)]{P49}} at
$z=0$ is a root of unity times $\prod_{k<q_1}\theta_{\mathsf x_k\Omega_0}(0)^{q_1}$ with $\mathsf
x_k\in K$. Let $\sigma\in\auAut(\C/K)$. By {\cite[Theorem~7.9(b)]{P49}}, comparing constant terms,
$\sigma(\theta_{\mathsf x\Omega_0}(0))=\theta^\sigma_{\mathsf x'\Omega_\sigma}(0)$ for suitable
$\mathsf x'\in K$, where $\theta^\sigma$ is the theta function of the normalised lattice
$\Lambda^\sigma$ with $\auO_K\Lambda^\sigma=\Lambda^\sigma$; so $|\sigma(\theta_{\mathsf x\Omega_0}
(0))|\le\mathsf M_d$ ({\cite[(76)]{P49}}). Roots of unity go to roots of unity. Hence
\[
\bigl|\sigma\,\theta(1\otimes\psrho_f(x)\psvac^\circ_1)\bigr|\le q_1^{2f}\mathsf M_d^{\,q_1^2}\le
\exp\bigl(2q_1\log q_1+q_1^2\log^+\mathsf M_d\bigr)\le\exp(C_0q_1^2),
\]
using $f\le q_1$, $\log q_1\le q_1$ and $C_0\ge2+\log^+\mathsf M_d$. If $\sigma|_K$ is complex
conjugation, $\sigma=c\circ\sigma'$ with $\sigma'\in\auAut(\C/K)$, and $|\sigma(\cdot)|=|\sigma'(
\cdot)|$; so the bound holds for all $\sigma\in\auAut(\C)$.

\emph{Rank twelve.} $W_{12}$ is the orthogonal sum of twelve lines, $L$ and $\psvac^\circ=
\bigotimes_b\psvac^\circ_1$ are the corresponding products ({\cite[(58)]{P49}}), and the theta
functional of an orthogonal sum is the tensor product of those of the summands; so $c_{\mathrm{vac},y}=\prod_{b=1}
^{12}\theta(1\otimes\psrho_f(y_b)\psvac^\circ_1)$ and $|\sigma c_{\mathrm{vac},y}|\le\exp(12C_0q_1^2)$.

\emph{Comparison with $b$.} By {\cite[Lemma~8.7(c)]{P49}}, $C_0\le10^{-12}\log\log b\le\log x$ (by
(D2)), and $q_1\le x^{12}$ ({\cite[Lemma~13.4(S3)]{P46}}); so $12C_0q_1^2\le12x^{24}\log x\le
x^{25}\le b$ by (S1).
\end{proof}

\emph{Correspondence with (I1)--(I4).} Proposition~\ref{pst:I1} is (I1). Proposition~\ref{pst:I2}
is (I2): its vector $\Phi$ has the stated theta function, its values $\theta(1\otimes\psrho_f(x)
\Phi)$ at all finite adelic points (in particular at the points $(0,x)$) are integral at every
prime, and (a) is the invariance of its theta function under the signed translations by $\pi_M
\mathfrak a_0^{12}$. Proposition~\ref{pst:I3} is (I3) with $u_i=y\in Y$ ($\aubar{\pi_M}u_i\in
\mathfrak a_0^{12}$) and $a_i=a_y$. Proposition~\ref{pst:I4} is (I4), the value $\psvac^\circ(0,
y)$ being $\theta(1\otimes\psrho_f(y)\psvac^\circ)$. Consequently:

\begin{corollary}\label{pst:attained}
Theorem~\ref{psm:attained} holds unconditionally.
\end{corollary}




\providecommand{\auO}{\mathcal O}
\providecommand{\auU}{\mathrm U}
\providecommand{\auGL}{\mathrm{GL}}
\providecommand{\auHom}{\operatorname{Hom}}
\providecommand{\auAut}{\operatorname{Aut}}
\providecommand{\auNm}{\operatorname{Nm}}
\providecommand{\auA}{\mathbb A}
\providecommand{\auF}{\mathbb F}
\providecommand{\auPl}{\Phi}
\providecommand{\aubud}{\mathrm{bud}}
\providecommand{\ausinf}{s_\infty}
\providecommand{\ausfin}{s_{\mathrm{fin}}}
\providecommand{\aubar}[1]{\overline{#1}}
\providecommand{\plone}{\mathbf 1}
\providecommand{\psrho}{\varrho}
\providecommand{\psvac}{\mathbf v}
\providecommand{\poA}{\mathcal A}
\providecommand{\popr}{\operatorname{pr}}
\providecommand{\podel}{\delta}
\providecommand{\pop}{\mathsf p}
\providecommand{\pog}{\mathsf g}
\providecommand{\pexK}{K^{(z)}}
\providecommand{\psuC}{\mathcal C}
\providecommand{\psuT}{\mathsf T}
\providecommand{\psuN}{\mathsf N}
\providecommand{\psug}{\hat g}

\subsection{The test form of least valuation, and the untwisted norm}\label{psu:sec}

This section completes the proof of Proposition~\ref{po:small} (in the exact form
Proposition~\ref{psm:small}) by constructing the test form from the sharpness of the count
(Theorem~\ref{psm:attained}, which holds unconditionally by Corollary~\ref{pst:attained}): parts
(ii)--(iv) of the note to Proposition~\ref{po:small}. It then treats Proposition~\ref{po:arch},
which is reduced to one precisely stated law of Galois conjugation of the lifts of characters
(\S\ref{psu:ss-arch}). Together with \S\ref{psm:sec} and \S\ref{pst:sec}, this proves Proposition~\ref{po:small}
(Theorem~\ref{psu:small}, with the constant $\exp(fmb/k+3x^{35})$, which suffices for
Proposition~\ref{po:mainheight} and~{\cite[Proposition~14.42]{P46}}) and Proposition~\ref{po:arch} (Theorem~\ref{psu:arch}).

\emph{Conventions.} We use the notation of \S\ref{po:ss-hecke} and \S\ref{pex:ss-set}. Hermitian
forms are conjugate-linear in the first variable. $\psuC=\auU(1)(E^+)\backslash\auU(1)(\auA_{E^+,f})/
K^1_M$ is the finite abelian group of classes of $\auU(1)$ at the chosen level, $N_1=\#\psuC$, and
we fix representatives $x_c$ ($c\in\psuC$) with $x_{c_0}=1$ for the trivial class $c_0$. The groups
$\auU(1)(E^+)\cap x_cK^1_Mx_c^{-1}$ are trivial (proof of Lemma~\ref{pex:points}), and we use the
Haar measure on $\auU(1)(\auA_{E^+,f})$ with $\operatorname{vol}(K^1)=1$; hence for a function $F$ on
$[\auU(1)]$ that is right invariant under $K^1_M$,
\begin{equation}\label{psu:eq-int}
\int_{[\auU(1)]}F(g)\,dg=[K^1:K^1_M]^{-1}\sum_{c\in\psuC}F(x_c).
\end{equation}
The characters of $\psuC$ are the characters $\psi$ of $\auU(1)(E^+)\backslash\auU(1)(\auA_{E^+})$
that are trivial at the infinite places and on $K^1_M$; we write $\widehat\psuC$ for the group of
them and regard each $\psi$ as a function on $\auU(1)(\auA_{E^+,f})$. For a standard space $V'$ of
rank $2z'$, a level $K'$ and a class $h$ of $\Gamma\backslash\auU(V')(\auA_f)/K'$ put
$e_h=\#(\Gamma\cap hK'h^{-1})$; $\psuN_{z'}=\#\auGL_{4fz'}(\auF_3)$ is divisible by every $e_h$
(\S\ref{pex:ss-IM}).

\subsubsection{Flat pairs, their kernels and their values}\label{psu:ss-flat}

\begin{definition}[Flat pairs over a character]\label{psu:def-flat}
For an integer $2\le z'\le t$ and $\psi\in\widehat\psuC$ let $q(\psi,z')=(\upsilon^{-z'}\psi,z')$,
where $\upsilon$ is restricted to $\auU(1)$. For a tame character $\vartheta$ of modulus $Q$
(\S\ref{pex:ss-set}) write $\Psi_{q,\vartheta}$ for the block datum $\Phi^{(1)}_\vartheta$ of the pair
$q=q(\psi,z')$ on $\auHom(E,V')$, $V'=E^{2z'}$, and
\[
k_{q,\vartheta}(g,h)=\theta\bigl(\omega(g,h)(1\otimes\Psi_{q,\vartheta})\bigr)\qquad(g\in\auU(1)(\auA_{E^+}),
\ h\in\auU(V')(\auA_{E^+,f}))
\]
for its kernel, where $\omega$ is the representation \eqref{po:eq-w5} for the line and the character
$\nu_q\vartheta$, $\nu_q=\upsilon^{-z'}\psi$, and $1$ is the constant polynomial (the weight of $q$
is trivial).
\end{definition}

\begin{lemma}\label{psu:flat}
Let $q=q(\psi,z')$ as in Definition~\ref{psu:def-flat}.
\begin{enumerate}
\item[(a)] $q$ is a critical pair; it is flat ($c(\tau)=z'$ for all $\tau$, so $l_\tau=0$), and in the
notation of \S\ref{pex:ss-set} the character on $\auU(1)$ of its Weil representation is
$\alpha_{q,\vartheta}=\nu_q\vartheta\,\upsilon^{z'}=\psi\vartheta$.
\item[(b)] For every tame $\vartheta$ of modulus $Q$ and every $h$,
$\varphi_{q,\vartheta}(h)=c_{\mathrm{vac}}^{\,2z'}\sum_{c\in\psuC}k_{q,\vartheta}(x_c,h)$.
\item[(c)] For $g\in\auU(1)(\auA_{E^+,f})$ put $G_{q,\vartheta}(g)=\sum_he_h^{-1}\,\aubar{\varphi_{q,
\vartheta}(h)}\,k_{q,\vartheta}(g,h)$, the sum over the classes $h$ of $\auU(V')$ at the level of
$q$. Then $G_{q,\vartheta}(g)=\|\varphi_{q,\vartheta}\|^2/(N_1c_{\mathrm{vac}}^{\,2z'})$ for every $g$.
\end{enumerate}
\end{lemma}
\begin{proof}
(a) $\upsilon$ has weight $1$ and trivial tame part, and $\psi$ has weight $0$ and is tamely ramified
at the chosen places only (it is trivial on $K^1_M$); so $\nu_q$ has weight $-z'$ and $z'\le c=z'\le
t\le4b$ ({\cite[Lemma~13.4]{P46}}). The formula for $\alpha$ is that of \S\ref{pex:ss-set}.

(b) By \eqref{po:eq-lift} with $y=\plone$, $\varphi_{q,\vartheta}(h)=c_{\mathrm{vac}}^{\,2z'}[K^1:K^1_M]
\int_{[\auU(1)]}k_{q,\vartheta}(g,h)\,dg$ (the weight is trivial, so $\pop_\xi$ is a constant and
$W_\lambda=\C$). The integrand is right invariant under $K^1_M$ (proof of Lemma~\ref{pex:points}: at
the chosen places $T\mapsto Tk$ fixes $\plone[T\equiv e_1]$ for $k\equiv1$, and at $w\in Q$ the factor
$\vartheta_w(k)^{-1}$ cancels $\psi_w(x_1k)/\psi_w(x_1)$). Apply \eqref{psu:eq-int}.

(c) Let $u\in\auU(1)(\auA_{E^+,f})$. The element $u$ of $\auU(1)$ and the scalar matrix $u^{-1}\cdot
1_{V'}$ of $\auU(V')$ act on $\auHom(E,V')$ by the same map $T\mapsto Tu=uT$, so their images in the
symplectic group of $\auHom(E,V')$ coincide; hence $\omega(u,1)=\epsilon(u)\,\omega(1,u^{-1})$ for a
scalar $\epsilon(u)$, which is multiplicative in $u$ and of absolute value $1$ (both operators are
unitary). The scalar $u^{-1}$ is central, so right translation by it permutes the classes $h$ and
preserves the numbers $e_h$. From $k(ug,h)=\epsilon(u)k(g,hu^{-1})$ (write $k=k_{q,\vartheta}$, $\varphi
=\varphi_{q,\vartheta}$) and (b), $\varphi(hu^{-1})=c_{\mathrm{vac}}^{2z'}\sum_ck(x_c,hu^{-1})=
\epsilon(u)^{-1}c_{\mathrm{vac}}^{2z'}\sum_ck(ux_c,h)=\epsilon(u)^{-1}\varphi(h)$, because $c\mapsto
ux_c$ runs again over representatives of $\psuC$ and $k(\cdot,h)$ is left invariant under
$\auU(1)(E^+)$ and right invariant under $K^1_M$. Hence
\[
G(ug)=\epsilon(u)\sum_he_h^{-1}\aubar{\varphi(h)}k(g,hu^{-1})=\epsilon(u)\sum_he_h^{-1}
\aubar{\varphi(hu)}k(g,h)=\epsilon(u)\aubar{\epsilon(u)}\,G(g)=G(g),
\]
using $\varphi(hu)=\epsilon(u^{-1})^{-1}\varphi(h)=\epsilon(u)\varphi(h)$. So $G$ is constant, and
$N_1G=\sum_cG(x_c)=\sum_he_h^{-1}\aubar{\varphi(h)}\sum_ck(x_c,h)=\|\varphi\|^2/c_{\mathrm{vac}}^{2z'}$ by
(b) and the definition of the norm.
\end{proof}

\begin{definition}[Integral values]\label{psu:def-I}
For $q=q(\psi,z')$ and a tame character $\vartheta$ of modulus $Q$ put, with $n_q=M^{4fz'}$,
\[
I_q(\vartheta)=\psuN_{z'}\,n_q^2\,\frac{\|\varphi_{q,\vartheta}\|^2}{(\auNm Q)^{2z'-1}} .
\]
\end{definition}

\begin{lemma}[Integral values]\label{psu:I}
Let $q=q(\psi,z')$.
\begin{enumerate}
\item[(a)] $I_q(\vartheta)>0$, and $I_q(\vartheta)=I_q(1)\,\rho_q(\vartheta)$ for every twist $\vartheta$.
\item[(b)] Let $\ell\nmid\auNm Q$ and $\jmath\colon\C\to\C_\ell$. Then $n_q\varphi_{q,\vartheta}$ is
integral at $\jmath$, the numbers $(\psuN_{z'}/e_h)\,n_q\,\varphi_{q,\vartheta}(h)$ are integral at
$\jmath$, and $I_q(\vartheta)$ is integral at $\jmath$. In particular $I_q(1)$ is integral at every
prime.
\item[(c)] $|\sigma I_q(1)|\le\exp(x^{35})$ for all $\sigma\in\auAut(\C)$.
\end{enumerate}
\end{lemma}
\begin{proof}
(a) Proposition~\ref{pex:rallis} for the pair $q$ ($z'\ge2$): $\|\varphi_{q,\vartheta}\|^2=A_q(\auNm
Q)^{2z'-1}L^{\mathsf M\cup Q}(z',(\nu_q\vartheta)^{-1})$ with $A_q>0$ and a positive $L$-value; divide
by the case $\vartheta=1$.

(b) For $l_\tau=0$ the integer of the kernel integrality of \S\ref{pex:ss-IM} may be taken to be
$n_q=M^{4fz'}$ (the factor $d^{\sum l_\tau}$ of Lemma~\ref{pex:dint} is $1$), and that statement holds at
every $\ell\nmid\auNm Q$. Indeed, by the finite Fourier analysis recalled in the proof of (c) below (and
the same at the places of $Q$, where the index $q_w$ is a unit at $\ell$), $M^{2fz'}\auNm Q\cdot
c_{\mathrm{vac}}^{2z'}\Phi^0_Q$ is a sum, with roots of unity as coefficients, of vectors
$\psrho^\circ(u)\psvac^\circ$ with $\aubar{\pi_M}u$ integral above $\ell$, that is, a standard datum of
level $\aubar{\pi_M}$ at $\ell$; the reduction in the proof of {\cite[Proposition~5.30]{P49}} (weak
approximation, and the fact of {\cite[\S7.4]{P49}} that the Weil operators of elements that are
integral with integral inverse above $\ell$ preserve standard data up to units) shows that
$M^{2fz'}c_{\mathrm{vac}}^{2z'}k_\Phi(x,h;1)$, $\Phi=\Phi^0_Q$, is a unit at $\ell$ times the value
$\theta(1\otimes f)$ of a standard datum $f$, and this value is integral at every prime by
{\cite[Lemma~7.14(a)]{P49}} (the jets are of degree $0$, so no condition on $\ell$ is needed). By
Lemma~\ref{pex:points}, $\varphi_{q,\vartheta}(h)$ is a sum of values of $\vartheta^{-1}$
(roots of unity) times kernels $k^0_Q(x,h)$, so $n_q\varphi_{q,\vartheta}(h)$ is integral at $\jmath$
(the weight is trivial, so the gauge is the absolute value). Since $e_h\mid\psuN_{z'}$, the second
assertion follows. For the third write $\|\varphi\|^2=\sum_he_h^{-1}\aubar{\varphi(h)}\varphi(h)$: the
number $n_q\aubar{\varphi(h)}$ is the complex conjugate of $n_q\varphi(h)$, which is integral at
$\jmath\circ c$ by the first assertion applied to $\jmath\circ c$, so it is integral at $\jmath$; and
$\auNm Q$ is a unit at $\ell$. (This is the proof of Proposition~\ref{pex:IM}(b), which does not use
that the pair is ramified.)

(c) Replacing $\sigma$ by $c\circ\sigma$ does not change $|\sigma I_q(1)|$, so we may assume $\sigma|_K=
1$. At each chosen place $w$, finite Fourier analysis on $(\auO_w/w)^{2z'}$ writes the indicator of
$e_1+w\auO_w^{2z'}$ as $q_w^{-2z'}$ times a sum of $q_w^{2z'}$ Heisenberg translates of the local
vacuum by vectors $u_w$ with $\varpi_wu_w$ integral, with roots of unity as coefficients. The chosen
places are the $2f$ primes of $E$ above $p_0$, $p_1$, of norms $p_0$, $p_1$; so
$c_{\mathrm{vac}}^{2z'}\Psi_{q,1}=M^{-2fz'}\sum_u\zeta_u\psrho^\circ(u)\psvac^\circ$, a sum of
$M^{2fz'}$ translates of the integral vacuum with roots of unity $\zeta_u$. By Step~1 of the proof of
{\cite[Proposition~14.23]{P46}} ({\cite[\S5.1.5]{P49}}), $\omega(g,h)\psrho^\circ(u)\psvac^\circ=\epsilon
\,\psrho^\circ(u')\omega(k)\psvac^\circ$ with $|\sigma\epsilon|=1$ and $k\in\auU(\auHom(E,V'))
(\auA_f)$; and by {\cite[Lemma~5.7(i)]{P49}} in degree $0$ and {\cite[Lemma~5.7(ii)]{P49}} at $y=0$,
$|\sigma\theta(1\otimes f)|=|\theta(1\otimes\Lambda_\sigma f)|\le\exp(C_0(2z'q_1+1)^2)$ for every such
translate $f$. Hence $|\sigma(c_{\mathrm{vac}}^{2z'}k_{q,1}(g,h))|\le\exp(C_0(2z'q_1+1)^2)$, and by
(b), $|\sigma\varphi_{q,1}(h)|\le N_1\exp(C_0(2z'q_1+1)^2)$; the same holds for $\sigma'=c\sigma c$.
Since $\sigma(\aubar a\,b)=\aubar{\sigma'(a)}\,\sigma(b)$, $|\sigma\|\varphi_{q,1}\|^2|\le h_qN_1^2
\exp(2C_0(2z'q_1+1)^2)$, with $h_q$ the number of classes of the level of $q$. By
{\cite[Proposition~7.20]{P49}} (the level contains the principal congruence subgroup modulo $M$, and
$2z'\le2t\le n_m$; also for $\auU(1)$), $h_q,N_1\le\exp(x^{30})$. Therefore
\[
\log|\sigma I_q(1)|\le(4fz')^2\log3+8fz'\log M+3x^{30}+2C_0(2z'q_1+1)^2\le x^{35},
\]
by the inequality (N5) of Lemma~\ref{psu:sizes} below.
\end{proof}

\begin{lemma}[Sizes]\label{psu:sizes}
For $1\le z'\le t$:
\textup{(N5)} $(4fz')^2\log3+8fz'\log M+3x^{30}+2C_0(2z'q_1+1)^2\le x^{35}$;
\textup{(N6)} $fmb/k+3x^{35}\le\aubud$ and $x^{35}+2f\le\aubud$, where $\aubud=fxb(\log b)^3/(mk)$.
\end{lemma}
\begin{proof}
By {\cite[Lemma~13.4]{P46}}: $t=5r\le15x^5$, $f,q_1\le x^{12}$ (S3), $\log M\le22\log b\le22x^{1/100}$
(S1), (S5), and $C_0\le\log\log b\le\log x$ ((D2) and {\cite[(35)]{P46}}). So $(4fz')^2\log3\le
3600\cdot1.1\,x^{34}$, $8fz'\log M\le x^{19}$, and $2C_0(2z'q_1+1)^2\le2\log x\cdot(31x^{17})^2\le
2000\,x^{34}\log x$. Since $\log x\ge10^{14}$ (S1), $4000+2000\log x+3+1\le x$, which gives (N5).
For (N6): $\aubud/(fmb/k)=x(\log b)^3/m^2\ge(\log b)^3/4\ge2$ (as $m^2\le4x$ and $\log b\ge128$), so
$fmb/k\le\aubud/2$; and $fmb/k\ge b\ge x^{10^6}\ge6x^{35}+2f$ (as $f\ge k$ and $m\ge1$, and (S1)),
so $fmb/k+3x^{35}\le2fmb/k\le\aubud$ and $x^{35}+2f\le\aubud$.
\end{proof}

\subsubsection{Averages over the classes of $\auU(1)$}\label{psu:ss-avg}

\begin{lemma}[Orthogonality]\label{psu:orth}
Let $s$ be a function on $\psuC$, $\ell$ a prime and $\jmath\colon\C\to\C_\ell$. For $\psi\in
\widehat\psuC$ put $A_\psi=N_1^{-1}\sum_{c\in\psuC}\psi(x_c)s(c)$. Then $\sum_\psi A_\psi=s(c_0)$, and
some $\psi$ has $|A_\psi|_\jmath\ge|s(c_0)|_\jmath$.
\end{lemma}
\begin{proof}
$\sum_\psi\psi(x_c)$ is $N_1$ for $c=c_0$ and $0$ otherwise (orthogonality of the characters of the
finite abelian group $\psuC$). If all $|A_\psi|_\jmath$ were smaller than $|s(c_0)|_\jmath$, the
ultrametric inequality would give $|\sum_\psi A_\psi|_\jmath<|s(c_0)|_\jmath$.
\end{proof}

\begin{lemma}[Multiplicative pigeonhole]\label{psu:pigeon}
Let $G$ be a finite group, $a,b\colon G\to\C$, $V\in\C$ with $V^{\#G}=\prod_{\chi\in G}a(\chi)b(\chi)$,
and $|\cdot|$ an absolute value on $\C$ (archimedean or of the form $|\jmath(\cdot)|$). For every
$\psi\in G$ there are $\phi,\chi\in G$ with $\phi\chi=\psi$ and $|V|\le|a(\phi)b(\chi)|$.
\end{lemma}
\begin{proof}
$\chi\mapsto\psi\chi^{-1}$ is a bijection of $G$, so $\prod_\chi|a(\psi\chi^{-1})b(\chi)|=|V|^{\#G}$,
and the largest factor is at least the geometric mean $|V|$.
\end{proof}

\begin{lemma}[The geometric mean]\label{psu:mean}
Let $z\ge12$. For each $\chi\in\widehat\psuC$ let $q_3(\chi)=q(\chi,2)$ and $q_1(\chi)=q(\chi,z-10)$
(Definition~\ref{psu:def-flat}). There is $L\in\C^\times$ with
\[
L^{N_1}=\prod_{\chi\in\widehat\psuC}I_{q_3(\chi)}(1)\,I_{q_1(\chi)}(1)\qquad\text{and}\qquad
|\sigma L|\le\exp(2x^{35})\quad(\sigma\in\auAut(\C)).
\]
\end{lemma}
\begin{proof}
$2\le z-10\le t$ because $z\le r$. Let $L$ be any $N_1$-th root of the product, which is non-zero by
Lemma~\ref{psu:I}(a) ($\#\widehat\psuC=N_1$). For $\sigma\in\auAut(\C)$, the number $\sigma L$ is an
$N_1$-th root of $\prod_\chi\sigma(I_{q_3(\chi)}(1))\sigma(I_{q_1(\chi)}(1))$; Lemma~\ref{psu:pigeon} with
the usual absolute value and $\psi=1$ gives $\phi,\chi$ with $|\sigma L|\le|\sigma I_{q_3(\phi)}(1)|\,
|\sigma I_{q_1(\chi)}(1)|\le\exp(2x^{35})$ by Lemma~\ref{psu:I}(c).
\end{proof}

\subsubsection{The point form and the seesaw factorisation}\label{psu:ss-seesaw}

Fix a main pair $p=(\nu,z)$ for $j$ (Definition~\ref{po:main}), with $\nu$ of weight $-c$, and put
$l_\tau=c(\tau)-z$. Then
\begin{equation}\label{psu:eq-range}
\kappa(\tau)\le l_\tau\le\kappa(\tau)+x^{42}\qquad(\tau\in\auPl):
\end{equation}
from $\kappa_j(\tau)+r-t\le c(\tau)\le\kappa_j(\tau)$ and $m<z\le r$ we get $l_\tau\ge\kappa_j(\tau)-t
=\kappa(\tau)+(j+\frac12)\Delta-t\ge\kappa(\tau)$ (as $\Delta=2r^8\ge5r=t$) and $l_\tau\le\kappa_j(\tau)
\le\kappa(\tau)+(m+\frac12)\Delta\le\kappa(\tau)+x^{42}$, since $(m+\frac12)\Delta=(2m+1)r^8\le(4\sqrt x+1)\,3^8x^{40}\le x^{42}$ by $m\le2\sqrt x$ and $r\le3x^5$ ({\cite[Definition~13.3(D3)]{P46}}) and $x$ large. Also
$z>m\ge10^4$, so $2z-20\ge4$ and $2\cdot1+2\cdot20+3<2z$.

\emph{Blocks.} Write $V=V^{(z)}=E^{2z}=V_0\perp V_2\perp V_3\perp V_1$ for the decomposition of the
standard coordinates into consecutive blocks of ranks $12$, $4$, $4$, $2z-20$; $e^{(i)}_1$ is the
first basis vector of $V_i$, and $e_{13}=e^{(2)}_1$. For $h_i\in\auU(V_i)(\auA_f)$ let
$[h_0,h_2,h_3,h_1]$ be the block diagonal element. Fix a twist $\vartheta$, of modulus $Q$, and a
class $x_0\in(\auO_E/\pi_M)^{12}$.

\emph{The point.} Let $\psug=\psug(x_0)\in\auU(V)(\auA_{E^+,f})$ be integral, equal to $1$ at the places
not below $\mathsf M\cup Q$, such that at each chosen place $w$ (where $\auU(V)(E^+_v)=\auGL_{2z}(E_w)$,
$v$ below $w$) $\psug_we_1\equiv Y:=(x_0,e^{(2)}_1,e^{(3)}_1,e^{(1)}_1)\bmod w$, and at each $w\in Q$
$\psug_w$ is the permutation matrix of the transposition of $e_1$ and $e_{13}$. It exists: $Y\bmod w$
is non-zero (its last block is), so it is the first column of an element of $\auGL_{2z}(\auO_w/w)$,
which lifts to $\auGL_{2z}(\auO_w)$.

\emph{The data.} Let $\Psi_0$ be the datum on $\auHom(E,V_0)$ that is $\plone[T\equiv x_0\bmod w]$ at
the chosen places and the vacuum at all other places (it is the indicator datum $\phi_{x_0}$ of
\S\ref{psm:ss-attained}). For a twist $\vartheta$ (or $\vartheta=1$) and $\psi_3,\psi_1\in\widehat\psuC$
let $q_2=q(\chi_0,2)=(\upsilon^{-2}\chi_0,2)$ (the pair of Proposition~\ref{po:small}), $q_3=q(\psi_3,2)$
and $q_1=q(\psi_1,z-10)$, on $V_2$, $V_3$, $V_1$, with the kernels $k_2=k_{q_2,\vartheta}$,
$k_3=k_{q_3,1}$, $k_1=k_{q_1,1}$ of Definition~\ref{psu:def-flat}. Put $\psi=\psi_3\psi_1$ and
\[
\chi'=\nu\,\upsilon^{z}\chi_0^{-1}\quad\text{on }\auU(1),\qquad s(g)=\theta\bigl(\omega_{12}(g,1)(\pop_\alpha
\otimes\Psi_0)\bigr)\quad(g\in\auU(1)(\auA_{E^+})),
\]
where $\omega_{12}$ is the Weil representation of the pair $(E,V_0)$ in the normalisation of
\eqref{po:eq-w5} with $\auU(1)$-character $\chi'$ (in the notation of \S\ref{pex:ss-set}: $\alpha=\chi'$,
$\beta=1$), and $\pop_\alpha$ is the monomial $z^\alpha=\prod_{\tau}\prod_{b\le12}z_{\tau,b}^{\alpha(\tau,b)}$
in the Fock variables of $\auHom(E,V_0)$, for an exponent $\alpha\in\N^{\auPl\times12}$ with
$\sum_b\alpha(\tau,b)=l_\tau$. The function $s$ is left invariant under $\auU(1)(E^+)$ and right
invariant under $K^1_M$ (at the chosen places $T\mapsto Tk$ fixes the indicator for $k\equiv1$; at the
other finite places $\chi'$ is unramified and the datum is the vacuum), so it is a function on
$\psuC$.

\begin{lemma}[The vector]\label{psu:vector}
Let $\alpha$ be as above and $a=\prod_\tau l_\tau!/\prod_{\tau,b}\alpha(\tau,b)!\in\Z_{>0}$. There is
$v\in W_\lambda(\Z)$ (the weight of $\poA(V;p,\vartheta)$: one row of length $l_\tau$ at $\tau$) with
$\pop_v=a\,z^\alpha$, the monomial in the variables of the first block; $v$ is fixed by every element of
$\prod_\tau\auGL_{2z}$ (over any field of characteristic $0$, and over $\auO_{\C_\ell}$) that fixes
$e_1,\dots,e_{12}$.
\end{lemma}
\begin{proof}
At $\tau$ the weight is $\operatorname{Sym}^{l_\tau}\C^{2z}\subseteq(\C^{2z})^{\otimes l_\tau}$, the span of
the vectors $\mathrm{ev}(y)=y^{\otimes l_\tau}$ of \eqref{po:eq-ev} (one column of height $1$ for each
box), and $W_\lambda(\Z)$ is its intersection with the integral tensors. Let $v_\tau$ be the sum of
the $l_\tau!/\prod_b\alpha(\tau,b)!$ distinct basis tensors $e_{b_1}\otimes\dots\otimes e_{b_{l_\tau}}$ in
which each $b\le12$ occurs $\alpha(\tau,b)$ times; it is symmetric and integral, and $v=\bigotimes_\tau
v_\tau$. By \eqref{po:eq-poly} the polynomial of a basis tensor is the product of the corresponding
Fock variables, so $\pop_v=\prod_\tau(l_\tau!/\prod_b\alpha(\tau,b)!)\prod_bz_{\tau,b}^{\alpha(\tau,b)}=
az^\alpha$. An element fixing $e_1,\dots,e_{12}$ (acting by its tensor powers, or by those of its
complex conjugate) fixes each such basis tensor.
\end{proof}

\begin{lemma}[Seesaw factorisation]\label{psu:factor}
With the notation above, let $x=[1,h_2,h_3,h_1]\,\psug$ and let $k_p(g,x;v)$ be the kernel of
\S\ref{po:ss-set} for the lift of $p$ at $\vartheta$ (the line $E$, the space $V$, the character
$\nu\vartheta$, the block datum $\Phi^{(1)}_\vartheta$). Then for $g\in\auU(1)(\auA_{E^+})$
\[
k_p(g,x;v)=a\,\psi(\det g)\,s(g)\,k_2(g,h_2)\,k_3(g,h_3)\,k_1(g,h_1).
\]
\end{lemma}
\begin{proof}
\emph{The datum.} $\psug$ is integral, so the unramified scalar of \eqref{po:eq-w5} is $1$ and
$\omega(1,\psug)\Phi(T)=\Phi(\psug^{-1}T)$ at the places of $\mathsf M\cup Q$ (at the conjugate places
$\psug$ is integral and fixes the vacuum). At a chosen place, $\plone[\psug^{-1}T\equiv e_1]=\plone[T
\equiv Y]=\plone[T_0\equiv x_0]\,\plone[T_2\equiv e^{(2)}_1]\,\plone[T_3\equiv e^{(3)}_1]\,\plone[T_1\equiv
e^{(1)}_1]$, $T_i$ the block of $T$ in $V_i$. At $w\in Q$, $\plone[T\text{ integral}]\psi_w((\psug^{-1}
T)_1)=\plone[T\text{ integral}]\psi_w(T_{13})$, and $T_{13}$ is the first coordinate of $T_2$. Hence
$\omega(1,\psug)\Phi^{(1)}_\vartheta=\Psi_0\otimes\Psi_{q_2,\vartheta}\otimes\Psi_{q_3,1}\otimes\Psi_{q_1,1}$.

\emph{The characters.} The $\auU(1)$-character of the Weil representation of $p$ at $\vartheta$ is
$\nu\vartheta\upsilon^z$ (\S\ref{pex:ss-set}), and those of the four blocks are $\chi'\psi^{-1}$ and, by
Lemma~\ref{psu:flat}(a), $\chi_0\vartheta$, $\psi_3$, $\psi_1$; their product is $\nu\vartheta\upsilon^z$.

\emph{The factorisation.} By (W7) of \S\ref{po:ss-glue}, which applies verbatim to the line $E$ and
the orthogonal decomposition $V=V_0\perp V_2\perp V_3\perp V_1$ (the characters multiply), the Weil
representation, the theta functional and the Fock vectors factor over the blocks. By
Lemma~\ref{psu:vector}, $\pop_v=a\,z^\alpha\otimes1\otimes1\otimes1$. Since $[1,h_2,h_3,h_1]$ acts on the
blocks separately and $\omega(g,hh')=\omega(g,h)\omega(1,h')$,
\[
k_p(g,x;v)=a\,\theta_{12}\bigl(\omega_{12}'(g,1)(z^\alpha\otimes\Psi_0)\bigr)\,k_2(g,h_2)\,k_3(g,h_3)\,
k_1(g,h_1),
\]
with $\omega'_{12}$ the representation of $(E,V_0)$ with character $\chi'\psi^{-1}$. By
\eqref{po:eq-w5} the action of $\auU(1)$ carries the factor $\alpha(\det g)^{-1}$ of its character
$\alpha$, so $\omega_{12}'(g,1)=\psi(\det g)\,\omega_{12}(g,1)$, which gives the factor $\psi(\det g)s(g)$.
\end{proof}

\subsubsection{The test form}\label{psu:ss-test}

Keep $p$, $\vartheta$, $x_0$, $\alpha$, $v$, $\psi_3$, $\psi_1$ and the three flat pairs. For
$i\in\{2,3,1\}$ let $\vartheta_2=\vartheta$, $\vartheta_3=\vartheta_1=1$, $z_2=z_3=2$, $z_1=z-10$, and
$\psuT_i(h)=\psuN_{z_i}n_{q_i}^2(\auNm Q_i)^{1-2z_i}\varphi_{q_i,\vartheta_i}(h)$, where $Q_2=Q$ and
$Q_3=Q_1=\emptyset$. Define
\begin{equation}\label{psu:eq-Xi}
\Xi=\sum_{h_2}\sum_{h_3}\sum_{h_1}\frac{\psuT_2(h_2)\,\psuT_3(h_3)\,\psuT_1(h_1)}{e_{h_2}e_{h_3}e_{h_1}}\,
\podel\bigl[[1,h_2,h_3,h_1]\psug;\,v\bigr]\ \in\poA(V;p,\vartheta),
\end{equation}
the sums over the classes of $\auU(V_i)$ at the levels of the pairs $q_i$ (twisted by $\vartheta_i$),
with $\podel[\cdot;\cdot]$ as in \S\ref{po:ss-glue} (for the line $E$ and the space $V$).

\emph{The summand depends only on the classes.} Let $\gamma\in\auU(V_i)(E^+)$ and $k$ in the level of
$q_i$ at $\vartheta_i$, and replace $h_i$ by $\gamma h_ik$. The element $\gamma$ (placed in the block
$V_i$) is a rational point of $\auU(V)$ that fixes $e_1,\dots,e_{12}$, hence fixes $v$
(Lemma~\ref{psu:vector}), and $\varphi_{q_i}(\gamma h_i)=\varphi_{q_i}(h_i)$ (trivial weight). Next
$[\dots,k,\dots]\psug=\psug k'$ with $k'=\psug^{-1}[\dots,k,\dots]\psug$. At a chosen place $k$ fixes
$e^{(i)}_1$ modulo $w$, so $k'e_1\equiv\psug^{-1}[\dots,k,\dots]Y\equiv\psug^{-1}Y\equiv e_1$; at $w\in Q$,
$k'$ is $k$ with the coordinates $1$ and $13$ exchanged, so for $i\ne2$ it does not touch the first
row and column, and for $i=2$ its first row is the first row of $k\bmod w$ placed in the coordinates
$1,14,15,16$, which is $((k)_{11},0,0,0)$ by the definition of the level of $q_2$ at $\vartheta$. So
$k'\in\pexK_\vartheta$, and the character of $k'$ (the value at $k'$ of the character by which the level
acts on $\Phi^{(1)}_\vartheta$: trivial at the chosen places, $\psi_w$ of the $(1,1)$ entry at $w\in Q$)
equals the value at $k$ of the character of the level of $q_i$ at $\vartheta_i$. By the rule
$\podel[\gamma xk;\gamma v]=\chi(k)^{-1}\podel[x;v]$ of \S\ref{po:ss-glue} and $\varphi_{q_i}(hk)=
\chi_{q_i}(k)\varphi_{q_i}(h)$ ({\cite[\S14.3]{P46}}), the two characters cancel; $e_{\gamma h_ik}=
e_{h_i}$. Hence \eqref{psu:eq-Xi} is well defined.

\begin{lemma}[Integrality]\label{psu:integral}
Let $\ell\nmid\auNm Q$. Then $\Xi$ is integral at every $\jmath\colon\C\to\C_\ell$.
\end{lemma}
\begin{proof}
By Lemma~\ref{psu:I}(b) (for $q_2$ at $\vartheta$, $\ell\nmid\auNm Q$, and for $q_3$, $q_1$ at $1$) the
coefficients $\psuT_i(h_i)/e_{h_i}$ are integral at $\jmath$ ($\auNm Q$ is a unit at $\ell$). By
\eqref{po:eq-deltasize} the form $\podel[x;v]$ has gauge at most the gauge of $v$ at $(\jmath,x)$, the
largest absolute value of a coordinate of $x_\ell^{-1}v$. For $x=[1,h_2,h_3,h_1]\psug$,
$x_\ell^{-1}v=\psug_\ell^{-1}[1,h_2,h_3,h_1]_\ell^{-1}v=\psug_\ell^{-1}v$ by Lemma~\ref{psu:vector}, and
$\psug_\ell$ is integral with integral inverse, so the coordinates of $\psug_\ell^{-1}v$ are integral.
Sums of integral forms with integral coefficients are integral.
\end{proof}

\begin{proposition}[Periods of false lifts in rank one]\label{psu:periods}
Let $w$ be off the level and $y\in\varrho(e_1)\poA(V;p,\vartheta)\cap\ker P_w(X_w)\cap\ker
D^{(1)}_\vartheta$ (Definition~\ref{po:R3}). Write $V=V'\perp V_1$ with $V'=V_0\perp V_2\perp V_3$ of
rank $20$. Let $q_1=q(\psi_1,z-10)$ and $k_1=k_{q_1,1}$. Let $x'\in\auU(V)(\auA_f)$ be integral with
integral inverse, equal to $1$ at the places not below $\mathsf M\cup Q$, and such that for every $k$ in
the level of $q_1$ the element $k'=x'^{-1}[1,k]\,x'$ lies in $\pexK_\vartheta$ and acts on
$\Phi^{(1)}_\vartheta$ by the scalar by which $k$ acts on $\Psi_{q_1,1}$ (for $x'=\psug$ this was shown
after \eqref{psu:eq-Xi}); under this condition the summand below depends only on the class of $h_1$.
Then for every $h'\in\auU(V')(\auA_f)$, every
$u\in W_\lambda$ whose polynomial $\pop_u$ involves only the variables of $V'$, and every
$g\in\auU(1)(\auA_{E^+,f})$,
\[
\sum_{h_1}e_{h_1}^{-1}\,\aubar{k_1(g,h_1)}\,\bigl\langle u,\,y\bigl([h',h_1]x'\bigr)\bigr\rangle=0,
\]
the sum over the classes of $\auU(V_1)$ at the level of $q_1$.
\end{proposition}
\begin{proof}
This is Theorem~\ref{ppe:main} for the line $W=E$ (rank $r=1$, definite), $V=V^{(z)}$ of rank $n=2z$
(so $t=n-r=2z-1$), the character $\nu\vartheta$, the idempotent $e_1$ in place of $e_0$, the splitting
$V=V'\perp V_1$ in place of $V''\perp V^\circ$ (with $n_1=20$), the element $x'$ in place of $\pog_s$, the
datum $\Psi_{q_1,1}$ and the trivial weight on $V_1$ in place of $\Psi^\circ$ and $\pop^\circ$, and the
level of $q_1$ in place of $J^\circ$. The steps of the proof of \S\ref{ppe:sec} are gone through one by
one in Remark~\ref{ppe:rem-rankone}, steps (1)--(8), which says for each step what replaces what, which
steps are unchanged, and proves the others (the property (F2), the generation by harmonics, and the
final step). Here we give the points where the rank, the idempotent and the data enter.
(F1) $D^{(1)}_\vartheta y=0$ is part of the hypothesis (it replaces Proposition~\ref{plf:L4}).
(F2) For $r=1$ and a chosen place $v$, $\lambda\in\auF_v^\times$ and $\psuN_\lambda=\{1+u\lambda\}=
\auF_v^\times$, so the hyperplane sum is $\eta=\sum_{s\in\auF_v^\times}\psi_v(s)^{-1}s$, where $\psi_v$ is
the tame part of $\nu\vartheta$ at $v$, namely $\omega_v^{\,j}$ ($\vartheta$ is unramified at $v$),
which is non-trivial because $1\le j\le m$ and $\omega_v$ has order greater than $8m$; since $y=\varrho
(e_1)y$ and $\eta e_1=(\sum_s\psi_v(s)^{-1})e_1=0$, $\varrho(\eta)y=0$ (this replaces
Lemma~\ref{ppe:green}). (F3) $P_w(X_w)y=0$ is part of the hypothesis, with $X_w$ of
Lemma~\ref{po:facts}(L3) for the line, $|X_w|=q^{-1/2-(t-1)}$. The tower
(Proposition~\ref{ppe:tower}) is stated for an arbitrary rank $r$ and needs $2r+2p+3<n$, here for
$p=n_1=20$: $45<2z$. Proposition~\ref{ppe:rallis} is stated for an arbitrary splitting of the standard
coordinates. In the final step of the proof of Theorem~\ref{ppe:main}, $\omega(1,x'^{-1})$ preserves
the data that are the vacuum outside $\Sigma=\mathsf S_\vartheta\cup\aubar{\mathsf S_\vartheta}$ because $x'$
is $1$ outside $\Sigma$; the factorisation of the kernel uses (W7) for $V'\perp V_1$ and
$\pop_u=\pop'\otimes1$, the characters on $\auU(1)$ being immaterial (an extra character of $\det g$ only
multiplies descents, \S\ref{ppe:ss-setup}). The function $h_1\mapsto\aubar{k_1(g,h_1)}\,\langle u,
\rho(h_{1,\infty})^{-1}y([h',h_{1,f}]x')\rangle$ on $\auU(V_1)(\auA)$ is invariant under
$\auU(V_1)(E^+\otimes\R)$ and under $\auU(V_1)(E^+)$ on the left, because $u$ is fixed by the block of
$V_1$ (its polynomial involves only the variables of $V'$) and the weight of $q_1$ is trivial; it is
right invariant under the level of $q_1$ by the condition on $x'$: with $s(k)$ the scalar by which $k$
acts on $\Psi_{q_1,1}$, $k_1(g,h_1k)=s(k)k_1(g,h_1)$ and $y(\,\cdot\,k')=s(k)\,y(\,\cdot\,)$, and
$|s(k)|=1$. Hence the weighted sum over the classes of the level, which
need not be free, is the integral over $[\auU(V_1)]$ divided by the volume of the level, as noted at the end of
Theorem~\ref{ppe:main}. The conclusion of that proof is the vanishing of the integral of $\theta^\circ
\cdot\aubar{c_\zeta y}$, whose complex conjugate is \eqref{ppe:eq-goal}, the integral form of the
displayed sum (with $\langle u,y(\cdot)\rangle=c_u y$ on finite points and $\theta^\circ=k_1$).
\end{proof}

\begin{lemma}[Orthogonality to the false lifts]\label{psu:orthfalse}
For every $w$ off the level, $\Xi$ is orthogonal to $\varrho(e_1)\poA(V;p,\vartheta)\cap\ker P_w(X_w)
\cap\ker D^{(1)}_\vartheta$.
\end{lemma}
\begin{proof}
Let $y$ be in this space. By \eqref{po:eq-delta}, $\langle\podel[x;v],y\rangle=\langle v,y(x)\rangle$,
so $\langle\Xi,y\rangle$ is a sum over $h_2,h_3$ of the conjugates of the coefficients
$\psuT_2(h_2)\psuT_3(h_3)/(e_{h_2}e_{h_3})$ times
\[
\sum_{h_1}e_{h_1}^{-1}\aubar{\psuT_1(h_1)}\,\langle v,y([h',h_1]\psug)\rangle
=\psuN_{z-10}n_{q_1}^2\,\aubar{c_{\mathrm{vac}}^{2z-20}}\sum_{c\in\psuC}\ \sum_{h_1}e_{h_1}^{-1}
\aubar{k_1(x_c,h_1)}\,\langle v,y([h',h_1]\psug)\rangle ,
\]
with $h'=[1,h_2,h_3]$, by Lemma~\ref{psu:flat}(b). Each inner sum vanishes by
Proposition~\ref{psu:periods} with $u=v$ (its polynomial involves only the variables of $V_0$,
Lemma~\ref{psu:vector}), $x'=\psug$ (which satisfies the condition of that proposition, by the
paragraph after \eqref{psu:eq-Xi} for $i=1$) and $g=x_c$. Hence $\langle\Xi,y\rangle=0$.
\end{proof}

\begin{proposition}[The pairing]\label{psu:pairing}
With $A_\psi=N_1^{-1}\sum_{c\in\psuC}\psi(x_c)s(x_c)$,
\[
\langle\Xi,\varphi_{p,\vartheta}\rangle=N_1^{-2}\,c_{\mathrm{vac}}^{\,12}(\pi_Ms_0)^{\sum_\tau l_\tau}\,a\,
I_{q_2}(\vartheta)\,I_{q_3}(1)\,I_{q_1}(1)\,A_\psi .
\]
\end{proposition}
\begin{proof}
By \eqref{po:eq-delta}, \eqref{po:eq-lift} with $y=\plone$, the definition \eqref{pex:eq-phi} of
$\varphi_{p,\vartheta}$ and \eqref{psu:eq-int} (the integrand is right invariant under $K^1_M$, proof of
Lemma~\ref{pex:points}), for every $x$
\[
\langle\podel[x;v],\varphi_{p,\vartheta}\rangle=\langle v,\varphi_{p,\vartheta}(x)\rangle=\kappa_p\sum_{c\in
\psuC}k_p(x_c,x;v),\qquad\kappa_p=c_{\mathrm{vac}}^{2z}(\pi_Ms_0)^{\sum_\tau l_\tau}.
\]
The coefficients $\psuT_i(h)/e_h$ enter conjugated (the pairing is conjugate-linear in the first
variable). Insert Lemma~\ref{psu:factor} and sum over $h_2,h_3,h_1$:
\[
\langle\Xi,\varphi_{p,\vartheta}\rangle=\kappa_p\,a\sum_c\psi(x_c)s(x_c)\prod_{i\in\{2,3,1\}}\Bigl(
\psuN_{z_i}n_{q_i}^2(\auNm Q_i)^{1-2z_i}\,G_{q_i,\vartheta_i}(x_c)\Bigr),
\]
with $G$ as in Lemma~\ref{psu:flat}(c) ($\psuN$, $n_q$, $\auNm Q$ are real). By Lemma~\ref{psu:flat}(c)
the $i$-th factor is $I_{q_i}(\vartheta_i)/(N_1c_{\mathrm{vac}}^{2z_i})$, independent of $c$. Since
$z_2+z_3+z_1=z-6$, $\kappa_pc_{\mathrm{vac}}^{-2(z-6)}=c_{\mathrm{vac}}^{12}(\pi_Ms_0)^{\sum l_\tau}$, and
$\sum_c\psi(x_c)s(x_c)=N_1A_\psi$.
\end{proof}

\subsubsection{Proof of Proposition~\ref{po:small}}\label{psu:ss-small}

\begin{theorem}[Proposition~\ref{po:small}, with the constant corrected]\label{psu:small}
Let $p$ be a main pair for $j$ and $q_2=(\upsilon^{-2}\chi_0,2)$. There is $c_0\in\C^\times$ with
\[
|\sigma c_0|\le\exp\bigl(fmb/k+3x^{35}\bigr)\qquad(\sigma\in\auAut(\C))
\]
such that for every prime $\ell$, every $\jmath\colon\C\to\C_\ell$ and every twist $\vartheta$ with
$\ell\nmid\auNm Q_\vartheta$ there is a test form $\Xi$ at $\ell$ for $(p,\vartheta)$ with
\[
\bigl|c_0\,\rho_{q_2}(\vartheta)\bigr|_\jmath\cdot\ell^{-F(\ell)}\le\bigl|\langle\Xi,\varphi_{p,\vartheta}
\rangle\bigr|_\jmath,\qquad F(\ell)=\sum_{\sigma\in\auPl}s_\ell(\kappa(\sigma),\sigma).
\]
\end{theorem}
\begin{proof}
\emph{The constant.} By \eqref{psu:eq-range}, Theorem~\ref{psm:attained} (unconditional by
Corollary~\ref{pst:attained}) applies to $l=(l_\tau)$ and gives $t\in\C^\times$ with $|\sigma t|\le
\exp(fmb/k)$. Let $L$ be as in Lemma~\ref{psu:mean} and put $c_0=t\,I_{q_2}(1)\,L$. It is non-zero
(Lemma~\ref{psu:I}(a)), it does not depend on $\ell$, $\jmath$, $\vartheta$, and by
Lemmas~\ref{psu:I}(c), \ref{psu:mean}, $|\sigma c_0|\le\exp(fmb/k+x^{35}+2x^{35})$.

\emph{The test form.} Fix $\ell$, $\jmath$ and $\vartheta$ with $\ell\nmid\auNm Q_\vartheta$. By
Theorem~\ref{psm:attained} there are $\gamma\in\N^{\auPl\times12}$ with $\sum_b\gamma(\tau,b)=l_\tau$ and
a class $y_0$ modulo $\pi_M$ with
\begin{equation}\label{psu:eq-att}
\ell^{-F(\ell)}\le\Bigl|\frac{c_{\mathrm{vac}}^{12}}{t}\prod_\tau l_\tau!\cdot[\gamma]\,\Theta_{\phi_{y_0}}
(\pi_Ms_0\,\cdot)\Bigr|_\jmath .
\end{equation}
Take $\alpha=\gamma$ and $x_0=y_0$ in \S\ref{psu:ss-seesaw}; then $\Psi_0=\phi_{y_0}$. By
{\cite[(33)]{P49}} with $Q=z^\alpha$ and $c=1$, $s(1)=\theta(z^\alpha\otimes\phi_{y_0})=\partial^\alpha
\Theta_{\phi_{y_0}}(0)=\prod_{\tau,b}\alpha(\tau,b)!\cdot[\alpha]\Theta_{\phi_{y_0}}$, and $[\alpha]\Theta_{\phi_{y_0}}
(\pi_Ms_0\,\cdot)=(\pi_Ms_0)^{\sum l_\tau}[\alpha]\Theta_{\phi_{y_0}}$. So \eqref{psu:eq-att} reads
\begin{equation}\label{psu:eq-s1}
|t|_\jmath\,\ell^{-F(\ell)}\le\bigl|c_{\mathrm{vac}}^{12}(\pi_Ms_0)^{\sum l_\tau}\,a\,s(1)\bigr|_\jmath .
\end{equation}
By Lemma~\ref{psu:orth} (applied to $s$ on $\psuC$, $x_{c_0}=1$) there is $\psi\in\widehat\psuC$ with
$|A_\psi|_\jmath\ge|s(1)|_\jmath$. By Lemma~\ref{psu:pigeon} for the group $\widehat\psuC$, the functions
$a(\chi)=I_{q(\chi,2)}(1)$, $b(\chi)=I_{q(\chi,z-10)}(1)$, $V=L$ and the absolute value $|\cdot|_\jmath$,
there are $\psi_3,\psi_1$ with $\psi_3\psi_1=\psi$ and $|L|_\jmath\le|I_{q_3}(1)I_{q_1}(1)|_\jmath$,
$q_3=q(\psi_3,2)$, $q_1=q(\psi_1,z-10)$. Let $\Xi$ be the form \eqref{psu:eq-Xi} for these data. It is
a test form at $\ell$ for $(p,\vartheta)$ (Definition~\ref{po:R3}) by Lemmas~\ref{psu:integral}
and~\ref{psu:orthfalse}.

\emph{The inequality.} By Proposition~\ref{psu:pairing}, $|N_1^{-2}|_\jmath\ge1$, \eqref{psu:eq-s1}, the
choice of $\psi$, $\psi_3$, $\psi_1$, and $I_{q_2}(\vartheta)=I_{q_2}(1)\rho_{q_2}(\vartheta)$
(Lemma~\ref{psu:I}(a)),
\[
|\langle\Xi,\varphi_{p,\vartheta}\rangle|_\jmath\ge|c_{\mathrm{vac}}^{12}(\pi_Ms_0)^{\sum l_\tau}a|_\jmath\,
|I_{q_2}(\vartheta)|_\jmath\,|L|_\jmath\,|s(1)|_\jmath\ge|t\,I_{q_2}(1)\,L\,\rho_{q_2}(\vartheta)|_\jmath\,
\ell^{-F(\ell)} .\qedhere
\]
\end{proof}

\emph{Remark.} The deduction of Proposition~\ref{po:mainheight} needs $fmb/k+3x^{35}\le\aubud$, which is (N6) of Lemma~\ref{psu:sizes}. The only property of $q_2$ used is that it is the flat pair $q(\chi_0,2)$ on a block
of rank $4$; it is ramified ($\chi_0$ has non-trivial tame part at every chosen place), which is used
afterwards through (R2).

\subsubsection{The untwisted norm at all conjugates: Proposition~\ref{po:arch}}\label{psu:ss-arch}

\begin{lemma}[The norm as an $L$-value, with the archimedean factor]\label{psu:normL}
Let $2\le z\le t$. There is $B(z)>0$ such that for every critical pair $p=(\nu,z)$, with
$l_\tau=c(\tau)-z$,
\[
\|\varphi_{p,1}\|^2=B(z)\,\Bigl(\frac M\Omega\Bigr)^{\sum_\tau l_\tau}\prod_\tau l_\tau!\cdot d_p\cdot
L^{\mathsf M}(z,\nu^{-1}),\qquad d_p=\prod_\tau\binom{l_\tau+2z-1}{2z-1}\ge1,
\]
and $d_p=1$ if $p$ is flat.
\end{lemma}
\begin{proof}
Proposition~\ref{pex:rallis} with $\vartheta=1$ gives $\|\varphi_{p,1}\|^2=A_pL^{\mathsf M}(z,\nu^{-1})$ with
$A_p=|\kappa_p|^2\prod_{w\in\mathsf M}(q_w-1)\cdot\dim W_\lambda\cdot\|v^+\|^{-2}\cdot c'\langle P,P\rangle
\operatorname{vol}([\auU(V^{(z)})])\operatorname{vol}([\auU(1)])\zeta_{E^+}(2z)^{-1}>0$. Here
$|\kappa_p|^2=|c_{\mathrm{vac}}|^{4z}|\pi_M|^{2\sum l}|s_0|^{2\sum l}=|c_{\mathrm{vac}}|^{4z}(M/\Omega)^{\sum l}$
($|\pi_M|^2=M$, $|s_0|^2=1/\Omega$, {\cite[\S5.1.1]{P49}}); $W_\lambda=\bigotimes_\tau\operatorname{Sym}^{l_\tau}
\C^{2z}$ has dimension $d_p$; the highest weight vector is $v^+=\bigotimes_\tau e_1^{\otimes l_\tau}$, of norm $1$,
and $P=\pop_{v^+}=\prod_\tau z_{\tau,1}^{l_\tau}$ (\eqref{po:eq-poly}), whose Fock norm is $\prod_\tau l_\tau!$
in the normalisation of {\cite[(33)]{P49}}, in which $\langle z^\beta,z^{\beta'}\rangle=\delta_{\beta
\beta'}\beta!$. The remaining factors depend only on $z$ (and on the data); their product with
$|c_{\mathrm{vac}}|^{4z}$ is $B(z)$, positive because $A_p>0$ for the flat pair $q(1,z)$.
\end{proof}

\begin{lemma}[Euler products at $s\ge2$]\label{psu:euler}
Let $\chi$ be a unitary Hecke character of $E$, $S$ a finite set of places, and $s\ge2$ real. Then
$|\log L^S(s,\chi)|\le2f\log\zeta(2)\le f$ (principal branch of the logarithm of the Euler product).
\end{lemma}
\begin{proof}
$\log L^S(s,\chi)=\sum_{w\notin S,\,\chi\text{ unr.\ at }w}\sum_{k\ge1}\chi(\varpi_w)^kq_w^{-ks}/k$, so
$|\log L^S(s,\chi)|\le\sum_w\sum_kq_w^{-2k}/k=\log\zeta_E(2)$. For each rational prime $p$ there are at
most $[E:\Q]=2f$ places $w\mid p$, each with $q_w\ge p$, so $\zeta_E(2)\le\zeta(2)^{2f}$; and $\log\zeta(2)
=\log(\pi^2/6)<1/2$.
\end{proof}

The proof of Proposition~\ref{po:arch} needs, besides Lemmas~\ref{psu:normL}, \ref{psu:euler}
and~\ref{psu:I}, the following law of Galois conjugation for the norms of the lifts of characters,
proved as {\cite[Theorem~9.1]{P53}} in {\cite[\S9]{P53}}.

\emph{(CL) Let $\sigma\in\auAut(\C/K)$ and $2\le z\le t$. There is $a_\sigma(z)>0$ such that for every
critical pair $p=(\nu,z)$ there is a critical pair $p^\sigma=(\nu^\sigma,z)$, whose weight $-c^\sigma$
satisfies $c^\sigma(\sigma\circ\tau)=c(\tau)$ for all $\tau\in\auPl$, with
$|\sigma(\|\varphi_{p,1}\|^2)|=a_\sigma(z)\,\|\varphi_{p^\sigma,1}\|^2$.}

\begin{theorem}[Proposition~\ref{po:arch}]\label{psu:arch}
For a main pair $p=(\nu,z)$ for $j$ the constant family $\|\varphi_{p,1}\|^2$ has height
at most $\aubud-\sum_{\tau\in\auPl}\ausinf(\kappa(\tau))$.
\end{theorem}
\begin{proof}
Let $q=q(1,z)=(\upsilon^{-z},z)$, a flat critical pair of the same rank (Lemma~\ref{psu:flat}(a);
$2\le z\le r\le t$). Take the multiplier $H=I_q(1)/\|\varphi_{p,1}\|^2\in\C^\times$, $\alpha=\aubud-
\sum_\tau\ausinf(\kappa(\tau))$ and $h_\ell=0$ in {\cite[Definition~14.1]{P46}}.

\emph{The primes.} $H\|\varphi_{p,1}\|^2=I_q(1)$ is integral at every $\iota$ (Lemma~\ref{psu:I}(b)), so
$|H\|\varphi_{p,1}\|^2|_\iota\le1=\ell^{0}$.

\emph{The conjugates.} Both $I_q(1)$ and $\|\varphi_{p,1}\|^2$ are $\ne0$, and $|(c\circ\sigma)(x)|=|\sigma(x)|$;
since one of $\sigma$, $c\circ\sigma$ fixes $K$, we may assume $\sigma\in\auAut(\C/K)$. By (CL) ({\cite[Theorem~9.1]{P53}}) for $p$ and
for $q$ (note $q^\sigma$ is again flat, and $p^\sigma$ has degrees $l^\sigma_{\sigma\tau}=l_\tau$), and
Lemma~\ref{psu:normL} for $p^\sigma$ and $q^\sigma$,
\begin{align*}
\log|\sigma H|&=\log(\psuN_zn_q^2)+\log|\sigma\|\varphi_{q,1}\|^2|-\log|\sigma\|\varphi_{p,1}\|^2|\\
&=\log(\psuN_zn_q^2)+\log L^{\mathsf M}(z,(\nu_q^\sigma)^{-1})-\sum_\tau\ausinf(l_\tau)-\log d_{p^\sigma}-\log
L^{\mathsf M}(z,(\nu^\sigma)^{-1}),
\end{align*}
because $\log(a_\sigma(z)B(z))$ cancels, $\sum_\tau\bigl(\log l_\tau!+l_\tau\log(M/\Omega)\bigr)=\sum_\tau
\ausinf(l_\tau)$ is invariant under the permutation $\tau\mapsto\sigma\circ\tau$, and the flat pair has
$l=0$, $d=1$. By Lemma~\ref{psu:euler} the two $L$-terms contribute at most $2f$, and $\log d_{p^\sigma}
\ge0$. Since $\ausinf$ is increasing ($M/\Omega\ge1$) and $l_\tau\ge\kappa(\tau)$ by \eqref{psu:eq-range},
\[
\log|\sigma H|\le\log\psuN_z+8fz\log M+2f-\sum_\tau\ausinf(\kappa(\tau))\le\aubud-\sum_\tau\ausinf(\kappa
(\tau)),
\]
because $\log\psuN_z+8fz\log M\le x^{35}$ by (N5) and $x^{35}+2f\le\aubud$ by (N6) of
Lemma~\ref{psu:sizes}. Hence the family has height at most $\alpha$.
\end{proof}

\emph{Remarks.} (i) The multiplier is $I_q(1)/\|\varphi_{p,1}\|^2$ with the integer $\psuN_zn_q^2$ in $I_q$.
(ii) The proof uses {\cite[Theorem~9.1]{P53}} (through {\cite[Theorem~7.9]{P49}}), Lemma~\ref{psu:normL} (from Proposition~\ref{pex:rallis}), Lemma~\ref{psu:euler} and Lemma~\ref{psu:I}(b).



\appendix

\section{Statements used from \texorpdfstring{\cite{P46}}{[Row D]}, \texorpdfstring{\cite{P49}}{[Theta]}, \texorpdfstring{\cite{P51}}{[Moment]}, \texorpdfstring{\cite{P53}}{[Exchange]}, \texorpdfstring{\cite{P55}}{[Galois]}}\label{C2:quoted}

For the convenience of the reader we reproduce here, without proofs, the statements of other papers that are used in this one, in the notation of the paper quoted. Each is headed by its place in that paper.

\begin{quoteddef}{\cite[Definition~13.1]{P46}}{}
\emph{Parameters} are a pair $(q_0,d)$, where $q_0$ is a positive multiple of
$18480=2^4\cdot3\cdot5\cdot7\cdot11$ and $d$ is a prime with
$d\equiv q_0/2-1\pmod{q_0}$.
\end{quoteddef}

\begin{quoteddef}{\cite[Definition~13.2]{P46}}{}
Let $L\ge1$. A set $\mathsf T\subseteq\anPl_s$ is \emph{sparse for the primes
up to $L$} if for every prime $\ell\le L$ and all integers $0<i<j<l_0(s)$
there are no $\sigma,\tau,\rho\in\mathsf T$ and signs
$\epsilon,\epsilon'\in\{\pm1\}$ with
\[
 \ancl_{q'}(\tau/\sigma)=\epsilon\,\ell^{\,i},\qquad \ancl_{q'}(\rho/\sigma)=\epsilon'\,\ell^{\,j}\qquad\text{in }(\mathbb Z/q')^{\times},
\]
where $q'$ is the largest divisor of $q_1(s)$ prime to $\ell$.
\end{quoteddef}

\begin{quoteddef}{\cite[Definition~13.3]{P46}}{ (admissible data)}
Let $(q_0,d)$ be parameters. \emph{Admissible data} $D$ for $(q_0,d)$ consist
of the following objects, subject to the conditions \textup{(D1)--(D8)}.
\begin{enumerate}
\item[(D1)] Integers $1\le x\le x_1\le b$, the \emph{scales}.
\item[(D2)] They satisfy
\begin{gather*}
 10^{12}\,\anCP\le\log\log b,\qquad 100\log\log b\le\log x,\qquad 10^{6}\log x\le\log x_1,\\
 10^{6}\log x_1\le\log b,\qquad 6\log\log x_1\le\log\log b .
\end{gather*}
\item[(D3)] Integers $m$, $r$, $s$: $m$ is a power of $2$ with
$\sqrt x\le m\le2\sqrt x$; $r$ is a power of $3$ with $x^{5}\le r\le3x^{5}$;
$2^{s}=m^{22}$; and $s\le32\log x$. The field of the data is $E=E_s$, with
$q_1=q_1(s)$, $q_2=q_2(s)$, $f=[E:K]$, $\anPl=\anPl_s$.
\item[(D4)] $q_1x_1$ divides $b$.
\item[(D5)] A set $\anPb\subseteq\anPl$ (the \emph{$b$-places}) with
$\ancl_{q_0}(\sigma)=-1$ for all $\sigma\in\anPb$, with
$T(s)\le\#\anPb\cdot\log b$, and sparse for the primes up to $5r(r+1-m)$.
\item[(D6)] The weight $\kappa\colon\anPl\to\mathbb N$: $\kappa(\tau)=b$ for
$\tau\in\anPb$ and $\kappa(\tau)=xb/x_1$ otherwise.
\item[(D7)] Two distinct primes $p_0,p_1\equiv1\pmod{q_2^{2}}$ in the interval
$[b^{10},b^{11}]$ and an integer $\pi_M$ of $K$ with
$\pi_M\overline{\pi_M}=p_0p_1=:M$.
\item[(D8)] A prime $N\equiv1\pmod{q_2^{2}}$ with
\[
 \log N_-\le\log N\le\log N_-+\sqrt{\log b},\qquad
 \log N_-:=\frac mf\sum_{\tau\in\anPl}\log\kappa(\tau)+(m-1)\log\frac{M\,\anrd_E}{2\pi}-\frac{15}2\log b ,
\]
where $\anrd_E=|\operatorname{disc}E|^{1/[E:\mathbb Q]}$ is the root
discriminant, and an integer $\pi_N$ of $E$ with $\pi_N\pi_N^{*}=N$.
\end{enumerate}
We write $a=xb$ and $k=r+1-m$; the sets $\anPa$ and $\anPb$ are disjoint.
\end{quoteddef}

\begin{quotedthm}{\cite[Lemma~13.4]{P46}}{}
For admissible data the following hold.
\begin{enumerate}
\item[(S1)] $x^{10^{6}}\le x_1$, $x_1^{10^{6}}\le b$,
$4\log x_1\le(\log x_1)^{6}\le\log b$,
$(\log b)^{100}\le x$, $\log b\ge128$, and
$\log x\ge10^{14}(q_0+2d+|\aneu|)$.
\item[(S2)] $x^{11}\le2^{s}\le2^{22}x^{11}$; $l_0(s)\le2592\log x$;
$2m\le r$, so $r/2\le k\le r$.
\item[(S3)] $x^{11}\le T(s)=\#\anPa\le f\le q_1\le x^{12}$, and
$f=\varphi(q_0)\,\#\anPa$; moreover $1\le\#\anPa/\#\anPb\le\log b$.
\item[(S4)] $\anrd_E=\sqrt d\;q_1\prod_{p\mid q_0}p^{-1/(p-1)}$, and
$x^{10}\le\anrd_E\le x^{12}$.
\item[(S5)] $1\le xb/x_1\le\kappa(\tau)\le b\le a\le b^{2}$ for all $\tau$;
$b^{20}\le M\le b^{22}$;
$20\log b\le\log(M\anrd_E/2\pi)\le23\log b$; and
$m\log b\le\log N\le25\,m\log b$.
\item[(S6)] With $\Omega:=e^{\aneu}\cdot2\pi/\anrd_E$: $0<\Omega\le1$ and
$20\log b\le\log(M/\Omega)\le23\log b$.
\item[(S7)] With $n=m/2$ and $\varrho=n-1$ (so that $r=k+2\varrho+1$): $r$ is odd; every prime
dividing $q_1$ divides $q_0$; $q_0^{\,q_0}\le x$; $\varrho\ge10^{5}$;
$k\ge512000\,n^{9}$ and $k\ge1000\,(2\varrho+2)^{9}$; $k\ge2\varrho^{2}+3$; $18\varrho f\le b^{8}$;
$f\le b/1000$; $\anrd_E\le b$; and $N\le b^{60n}$. Moreover, with $\Delta=2r^{8}$:
\begin{itemize}
\item[(a)] $\varrho\ge10^{8}$ (so $n\ge100$); $n^{7}\le k$; $12n^{3}+500\le k$; $k-3n-6>10^{9}n^{3}$;
\item[(b)] $r\ge100$ (so $2r^{8}\ge500r^{3}+10r^{5}$); $3r^{3}f\le r^{7}$; $8nr\le b$; $kf\ge40$; $f\ge4000\,n$;
\item[(c)] $2f\le2^{2\varrho}$ and $2^{-\varrho}\le b^{-6}$;
\item[(d)] $(200\,f\,k\,\anrd_E)^{1000}\le b$ and $(25\,n^{3}f\,\anrd_E)^{1000}\le b$;
\item[(e)] $\sqrt d\le\anrd_E$, and $n,r,f,\sqrt d\le b$; $M\le b^{23}$;
\item[(f)] $\binom{a+k-1}{k-1}^{f}\le b^{2kf}$ (this bounds the dimension $d_\infty=\dim W_{\lambda^X}$ of
{\cite[\S14.3]{P46}}, a product of at most $f$ such binomial coefficients);
\item[(g)] $m\Delta+r\le\kappa(\tau)$ for all $\tau$; $k+2b+2m\Delta\le b^{2}$ (this bounds every
entry of $2\varsigma_\tau$, so that $L_\tau:=\max_i|2\varsigma_{\tau,i}|\le b^{2}$); and
$(r+1)^{4}(b^{2}+r+1+2n^{2})\le b^{5}$.
\end{itemize}
\end{enumerate}
\end{quotedthm}

\begin{quoteddef}{\cite[Definition~14.1]{P46}}{}
A family $y$ \emph{has height at most $c\in\R$} if there are $H\in\C^\times$, $\alpha\in\R$ and
reals $h_\ell\ge0$, one for each prime, such that $|\sigma H|\le e^\alpha$ for all
$\sigma\in\auAut(\C)$, $|H\,y_\iota|_\iota\le\ell(\iota)^{h_{\ell(\iota)}}$ for all $\iota$, and
$\alpha+\sum_{\ell\in T}h_\ell\log\ell\le c$ for every finite set $T$ of primes.
\end{quoteddef}

\begin{quotedthm}{\cite[Lemma~14.5]{P46}}{ (Free levels)}
Let $n\le2n_m$ and let $J\subseteq K^\circ$ consist of elements congruent to $1$ modulo every
chosen place. Then $\Gamma\cap gJg^{-1}=1$ for all $g\in G_f$.
\end{quotedthm}

\begin{quoteddef}{\cite[Definition~14.6]{P46}}{ (Target forms)}
The space of \emph{target forms} is $\auT=\varrho(e_0)\auAz$, the \emph{good forms} on $Y_0$; it
is a finite-dimensional inner product space, and $\audT=\dim\auT$.
\end{quoteddef}

\begin{quoteddef}{\cite[Definition~14.8]{P46}}{ (Prime sums, gain, main amount, remainder)}
Let $y\ge0$ be real and $\sigma\in\auPl$.
\begin{enumerate}
\item[(a)] The \emph{local prime sum} at $p$ is
$s_p(y,\sigma)=v(\pi_p)\sum_{e\in I_p}\bigl(y/p^e-\kappa(u(-p^e)\sigma)\bigr)^+$, and the
\emph{prime sum} is
$\ausfin(y,\sigma)=\sum_{1\le\nu\le a}\Lambda_K(\nu)\bigl(y/\nu-\kappa(u(-\nu)\sigma)\bigr)^+$.
\item[(b)] The \emph{gain} at $p$ is
\[
\augain_p=m\Bigl(\sum_{\tau\in\auPa}s_p(a,\tau)+k\sum_{\sigma\in\auPl}s_p(\kappa(\sigma),\sigma)
\Bigr).
\]
\item[(c)] The size of a derivative of order $j$ is $\ausinf(j)=\log j!+j\log(M/\Omega)$; the
\emph{direction amount} is $\auhdir=\sum_{\sigma\in\auPl}\bigl(\ausinf(\kappa(\sigma))-
\ausfin(\kappa(\sigma),\sigma)\bigr)$; the \emph{main amount} is $\auhmain=m\,k\,\auhdir$.
\item[(d)] The \emph{expected value} is
$\ausmain=a\bigl(\log a-\tfrac1f\sum_{\tau\in\auPl}\log\kappa(\tau)+c_1\bigr)$, and the
\emph{total remainder} is $\aurem=\sum_{\tau\in\auPa}\bigl(\ausfin(a,\tau)-\ausmain\bigr)$.
\end{enumerate}
\end{quoteddef}

\begin{quotedthm}{\cite[Lemma~14.9]{P46}}{ (The gains and the prime sums)}
Let $0\le y\le a$ and $\sigma\in\auPl$.
\begin{enumerate}
\item[(a)] $\ausfin(y,\sigma)=\sum_p2\log p\cdot s_p(y,\sigma)$, and $s_p(y,\sigma)=0$ if $p$
splits in $K$ or $p\ge x_1$. Hence $\augain_p=0$ for such $p$.
\item[(b)] With $\mathcal N$ the set of primes $p<x_1$ that do not split in $K$,
\[
2\sum_{p\in\mathcal N}\augain_p\log p=m\sum_{\tau\in\auPa}\ausfin(a,\tau)+m\,k\sum_{\sigma\in\auPl}
\ausfin(\kappa(\sigma),\sigma),
\]
and
\begin{equation*}\tag*{{\cite[(42)]{P46}}}
\auhmain-m\cdot\aurem+2\sum_{p\in\mathcal N}\augain_p\log p=m\,k\sum_{\sigma\in\auPl}\ausinf
(\kappa(\sigma))+m\cdot\#\auPa\cdot\ausmain .
\end{equation*}
\item[(c)] $\ausinf(a)-\ausmain\le a\bigl(\tfrac1f\sum_\tau\log\kappa(\tau)+\log\frac{M\aurd}
{2\pi}\bigr)+1+\tfrac12\log a$, and
\begin{equation*}\tag*{{\cite[(43)]{P46}}}
m\bigl(\ausinf(a)-\ausmain\bigr)-a\log N\le32\,a\log b .
\end{equation*}
\end{enumerate}
\end{quotedthm}

\begin{quotedthm}{\cite[Proposition~14.10]{P46}}{ (The count)}
Let $F$ have the normal form. Let $\lambda$ be a weight whose columns have height less than $n'$, with $a_\tau$ columns of
height $1$, $a'_\tau$ other columns and $l_\tau$ boxes in these, and $B\ge1$ boxes in all. Let $Q$ be either
the product $Q_{T,y}$, for a filling $T$ of $\lambda$ and integral directions $y$, of the minors of
$z_\tau y_\tau$ on the rows of each column and the first $h$ columns ($h$ the height of the column),
as in {\cite[Proposition~4.1]{P49}}, or an entry $S_\lambda(z)_{T',T}$ of the Schur matrix with $T'$
standard ({\cite[Lemma~4.4]{P49}}). Then the jet
$Q(\partial)F(0)$ has valuation at least
\[
v(\pi_p)\sum_{e\in I_p}\sum_{\tau}\Bigl(\bigl(\tfrac{a_\tau}{p^e}-a_{\tau'}-a'_{\tau'}\bigr)^+
+\bigl(\tfrac{l_\tau}{p^e}-l_{\tau'}\bigr)^+-\tfrac{a'_\tau}{p^e}\Bigr)-f\,n'^3\bigl(1+\log_p
B\bigr),\qquad\tau'=u(-p^e)\tau.
\]
For $p=d$ the same bound holds after subtracting in
addition the term $\tfrac12\sum_\tau\mathrm{ded}_0(\tau)=v(\pi_d)\sum_\tau\mathrm{ded}_0(\tau)$ for $e=0$, with
$\mathrm{ded}_0$ as in {\cite[Proposition~4.1]{P49}}, which is the precise statement in all cases.
\end{quotedthm}

\begin{quotedthm}{\cite[Theorem~14.15]{P46}}{ (the Covolume Theorem)}
Let $y_1,\dots,y_n\in\auT$ be linearly independent with $n\ge(1-4/b^2)\audT$, of index at most
$\delta$ and covolume at most $c$ in $\auLat_D$. Then
\[
c+2\delta\ \ge\ \audT\cdot\bigl(\auhmain-3\,f\,a\,(\log b)^4\bigr).
\]
\end{quotedthm}

\begin{quotedthm}{\cite[Lemma~14.21]{P46}}{ (Short bases)}
Let $n\ge1$ and let $c_0$ be a positive integer in the different of $E$. For every
$g\in\auU(E^n)(\auA_{E^+,f})$ there is $Y\in\auGL_n(E)$ with $g^{-1}Y\in M_n(\widehat{\auO}_E)$
all of whose entries have all conjugates of absolute value at most $C_{\mathrm{sb}}=6c_0(fn)^2$.
\end{quotedthm}

\begin{quotedthm}{\cite[Proposition~14.23]{P46}}{ (One kernel)}
Let $0\le j<m$, let $\phi$ be a translate of the vacuum and $\sigma\in\auAut(\C)$. Then for all
$g,g'$ and all $\xi$,
\[
\|\sigma(\theta_j(\phi;g,g')\xi)\|\le\exp\bigl(x^{50}+\tfrac12\mathrm{rs}_j\bigr)\,\|\sigma(\xi)\|,
\qquad\mathrm{rs}_j=\sum_{\tau\in\auPl}\ \sum_{\text{rows }i}\ausinf(\text{length of row }i
\text{ at }\tau).
\]
\end{quotedthm}

\begin{quotedthm}{\cite[Proposition~14.37]{P46}}{ (Factorisation of the chain determinant)}
For $1\le j\le m$ and a twist $\vartheta$ let $Z_j(\vartheta)=\det(u^*u)>0$, where
$u=\Theta_{j,\vartheta}$ is the theta lift from $e_0\mathcal A(W_j)$ to algebraic modular forms on
$\auU(V_j)$ with the character $\mu_j^0\vartheta$, normalised as in {\cite[\S2.4]{P53}} (joint
harmonic at infinity; at the chosen places the indicator of the inclusion; at $w\in Q_\vartheta$
the function $x\mapsto\psi_w(\det x^{\mathrm{top}})$; values in the algebraic modular forms of level
$J^V_\vartheta$ with their inner product). There are $e\in\Z$ and the ramified ratios $\rho_{ij}$,
$0<i<j\le m$, of {\cite[(9)]{P53}}, such that
for every $m$-fold twist $\vartheta$
\[
\frac{Z(1)}{Z(\vartheta)}=\Bigl((\auNm Q_\vartheta)^{e/2}\prod_{0<i<j\le m}\rho_{ij}(\vartheta)
\Bigr)^{-\audT}\cdot\prod_{j=1}^m\frac{Z_j(1)}{Z_j(\vartheta^j)} .
\]
\end{quotedthm}

\begin{quoteddef}{\cite[Definition~14.39]{P46}}{ (Regular families of twists)}
Put $\aubud=\aueu/(mk)$. A family $\theta\colon\auI\to\{m\text{-fold twists}\}$ is \emph{regular}
if:
\begin{enumerate}
\item[(R1)] $\ell(\iota)\nmid\auNm Q_{\theta_\iota}$ for all $\iota$;
\item[(R2)] for every ramified critical pair $(\nu,z)$, $1\le n\le m$ and sign $\pm$, the family
$\iota\mapsto|L(z,\nu^{-1})|/|L(z,(\nu\theta_\iota^{\pm n})^{-1})|$ has height at most $x_1$;
\item[(R3)] for $1\le j\le m$ there are $\beta_\ell\ge0$ with $\sum_{\ell\in T}\beta_\ell\le
\aubud$ for all finite $T$ such that, for every $\iota$ with prime $\ell$: (i) the twist
$\theta_\iota^{\,j}$ \emph{has denominators at most $\beta_\ell$ at $\iota$ for the $j$-th lift}
in the sense of Definition~\ref{po:denom} (the false lifts being those of
Definition~\ref{po:false}: the forms in the image of $e_0$ that are orthogonal to the lifts and
are killed by the Hecke polynomial $P_w(X_w)$ of Lemma~\ref{po:facts}(L3) at one split place $w$
off the level; a form is orthogonal to the false lifts if it is so for every such $w$); and (ii)
for every main pair $p$ for $j$ (Definition~\ref{po:main}) and all test forms $\Xi,\Xi'$ at
$\ell$ for $(p,\theta_\iota^{\,j})$ (Definition~\ref{po:R3}), the inner product of the orthogonal
projections of $\Xi$ and $\Xi'$ to the line of the normalised lift
$\varphi_{p,\theta_\iota^{\,j}}$ has absolute value at most $e^{\beta_\ell}$ at $\iota$.
\end{enumerate}
\end{quoteddef}

\begin{quotedthm}{\cite[Proposition~14.41]{P46}}{ (One lift and Hecke values)}
Let $1\le j\le m$ and $d_j=\dim e_0\mathcal A(W_j)$. There are a constant $C_j\ne0$, for every twist
$\vartheta$ a number $\Delta_j(\vartheta)\ne0$ (the determinant of a \emph{comparison operator}
compressed to the space of lifts) and the positive number $H_j(\vartheta)$ of \eqref{po:eq-H}: for the three critical pairs $(\nu_i,z_i)$ of \S\ref{po:ss-hecke} it is the product of $(\auNm\mathsf S_\vartheta)^{z_i^2}\prod_{1\le z\le z_i}|L^{\mathsf S_\vartheta}(z,(\nu_i\vartheta^{\pm1})^{-1})|^2$, such that
\[
Z_j(\vartheta)\cdot\Delta_j(\vartheta)=C_j\cdot\bigl(H_j(\vartheta)/H_j(1)\bigr)^{d_j},
\]
and
\begin{enumerate}
\item[(a)] $|\sigma\Delta_j(1)|\le\exp(3d_j\cdot f\,m\,b\,(\log b)^2)$ for all
$\sigma\in\auAut(\C)$;
\item[(b)] if $\ell(\iota)\nmid\auNm Q_\vartheta$ and $\vartheta$ has denominators at most $\beta$
at $\iota$ for the $j$-th lift, then $|\Delta_j(\vartheta)|_\iota\le\exp(d_j\,r\,\beta)$.
\end{enumerate}
\end{quotedthm}

\begin{quotedthm}{\cite[Proposition~14.42]{P46}}{ (Heights of twisted Hecke values)}
For a regular family $\theta$ and $1\le j\le m$ the family $\iota\mapsto H_j(1)/H_j(\theta_\iota^
{\,j})$ has height at most $10\,m\,x_1+k\,(6\,\aubud-\auhdir)$.
\end{quotedthm}

\begin{quotedthm}{\cite[Proposition~14.43]{P46}}{ (The cost of one lift)}
For a regular family $\theta$ and $1\le j\le m$ the family $\iota\mapsto Z_j(1)/Z_j(\theta_\iota^
{\,j})$ has height at most $\audT\bigl(f\,a\,(\log b)^4/m-k\,\auhdir\bigr)$.
\end{quotedthm}

\begin{quotedthm}{\cite[Proposition~2.14]{P49}}{ ((W2), (W2$'$), (W2$_\infty$))}
With $\theta_W,\theta_{W_I}$ as in {\cite[\S5.1.1]{P49}} and $di$ the Haar measure of mass $1$ on
$[I]=I(k)\backslash I(\auA)$:
\begin{itemize}
\item[(a)] $\theta_{W_I}(\psZ_IF)=\int_{[I]}\theta_W(\psrho(i)F)\,di$; this property together with
\textup{(b)} characterises $\psZ_I$.
\item[(b)] $\psZ_I\psrho(w)=\psrho(\bar w)\psZ_I$ for $w\in I^\perp(\auA)$.
\item[(c)] $\psZ_I\,\omega_W(p)=\chi_I(p)\,\omega_{W_I}(\bar p)\,\psZ_I$ for $p\in P_I(\auA)$, with
$\chi_I(p)=\upsilon(\det p|_I)^{-1}|\det p|_I|^{1/2}$ ($\det p|_I\in\auA_E^\times$, $|\cdot|$ the
idele norm of $E$). Locally, $\psZ_{I,v}\omega_{W,v}(p)=\chi_{I,v}(p)\omega_{W_I,v}(\bar p)
\psZ_{I,v}$ with $\chi_{I,v}(p)=\upsilon_v(\det p|_I)^{-1}|\det p|_I|_{E_v}^{1/2}$; at a real place
and for a line $I$ on which $p$ acts by $a\in\C^\times$, $\chi_{I,\tau}(p)=\upsilon_\tau(a)^{-1}|a|$
\textup{(W2$_\infty$)}. In particular $\chi_I=1$ on $N_I(\auA)$ and on $P_I(k)$.
\item[(d)] For isotropic $I\subset I'$, $\psZ_{I'}=\psZ_{I'/I}\circ\psZ_I$, and $\psZ_I=\psZ_{I,\infty}
\otimes\psZ_{I,f}$ (for suitable normalisations of the local maps).
\end{itemize}
\end{quotedthm}

\begin{quotedthm}{\cite[Lemma~2.20]{P49}}{ (Reduction to lines)}
There are integers $e_+,e_-$, independent of $W$, with $\Delta(\gamma^+,\gamma^-)=\det(\gamma^+)^{e_+}
\det(\gamma^-)^{e_-}$ for all $W$; and $e_++e_-=-1$. Hence \eqref{ptf:eq-compact} holds for all $W$
if and only if $e_+=0$, i.e.\ if and only if, for a positive definite line $\ell$, $\omega_{\ell,\tau}(z)$
fixes the vacuum for $z\in\auU(1)(\R)$.
\end{quotedthm}

\begin{quotedthm}{\cite[Proposition~2.21]{P49}}{ ($e_+=0$; \eqref{ptf:eq-compact})}
Let $\ell$ be a positive definite hermitian line at the real place $\tau$ and $\ell^-$ the negative
line. For $(z_1,z_2)\in\auU(\ell)(\R)\times\auU(\ell^-)(\R)$ the operator $\omega_{\ell\oplus\ell^-,\tau}
(z_1,z_2)$ multiplies the Fock vacuum by $z_2^{-1}$. Hence $e_+=0$, $e_-=-1$ in
{\cite[Lemma~2.20]{P49}} and~\eqref{ptf:eq-compact} holds: $\Delta(\gamma)=\aubar{\det\gamma^-}$.
\end{quotedthm}

\begin{quotedthm}{\cite[Theorem~3.2]{P49}}{ ((A1))}
Assume that every $\Lambda^{(j)}_\tau$ satisfies \textup{(Dec$_c$)} with $c=B+N_j^2+1$, $B=x^{25}$.
Then for $0\le j\le m$ the function $c_j=\bigotimes_\tau c_{j,\tau}$ is smooth, left and right
$K_{j,\infty}$-finite, an eigenfunction of the centre of the enveloping algebra, and
$\int_{G_{j,\infty}}|\mathsf Xc_j(x)|\,\|x\|^B\,dx<\infty$ for every left- or right-invariant
differential operator $\mathsf X$.
\end{quotedthm}

\begin{quotedthm}{\cite[Corollary~3.3]{P49}}{ ((A2))}
Assume \textup{(Dec$_c$)} with $c=5N_j^3$ for every $\Lambda^{(j)}_\tau$. Let $\mathsf P=\prod_\tau
\mathsf P_\tau$ be a parabolic subgroup of $G_{j,\infty}$ with $\mathsf P_\tau$ proper for at least
one $\tau$, and $\mathsf N$ its unipotent radical. Then $\int_{\mathsf N}|c_j(anb)|dn<\infty$ and
$\int_{\mathsf N}c_j(anb)\,dn=0$ for all $a,b$.
\end{quotedthm}

\begin{quotedthm}{\cite[Lemma~3.5]{P49}}{ (Rigidity)}
Let $\sigma$ be a unitary representation of $G$ on a separable Hilbert space and $y$ a smooth
vector with: $\sigma(t)y=t^{\Lambda^K(\lambda)}y$ for $t\in T$; $d\sigma(E_{ab})y=0$ for
$\varepsilon_a=\varepsilon_b$, $\lambda_a>\lambda_b$; $d\sigma(\Omega)y=\kappa_\lambda y$. Put
$e_\lambda=d_\lambda\overline{c_\lambda}$. Then $\langle y,\sigma(g)y\rangle=\|y\|^2c_\lambda(g)$ and
$\sigma(e_\lambda)y=y$, provided $e_\lambda\in L^1(G)$.
\end{quotedthm}

\begin{quotedthm}{\cite[Lemma~3.6]{P49}}{ (The derived action)}
After conjugation by an element of $T\times T'$, the derived actions of $\omega_\infty(\cdot,1)$ and
$\omega_\infty(1,\cdot)$ at $\tau$ on polynomial vectors are {\cite[(12)]{P49}}, with
\[
\gamma=\tau_0+\tfrac{n'}2,\quad\gamma'=\tau'_0+\tfrac n2,\qquad
\tau_0=\kappa_{j+1}(\tau)+\tfrac12[j\text{ odd}],\quad\tau'_0=-\kappa_{j+1}(\tau)-\tfrac12[j\text{ even}],
\]
provided the following holds.
\begin{enumerate}
\item[(V)] At $\tau$, the operator $\omega_{\auHom(Y_j,Y_{j+1})}(g,h)$ of {\cite[\S5.1.1]{P49}} is the
genuine oscillator action of the metaplectic covers (Weyl-ordered quadratic operators) times
$\det(g)^{n'/2}\det(h)^{-n/2}$, where $\upsilon_\tau$ restricted to $\auU(1)$ is $z\mapsto z$.
\end{enumerate}
\end{quotedthm}

\begin{quoteddef}{\cite[Definition~3.9]{P49}}{}
A \emph{joint harmonic} is a polynomial $\Phi\ne0$ with $E_{aa}\Phi=\Lambda_a\Phi$,
$E'_{cc}\Phi=\Lambda'_c\Phi$, $E_{ab}\Phi=0$ ($\varepsilon_a=\varepsilon_b$, $\lambda_a>\lambda_b$),
$E'_{cd}\Phi=0$ ($\varepsilon'_c=\varepsilon'_d$, $\lambda'_c>\lambda'_d$), $\Delta_{ab}\Phi=0$
($\varepsilon_a\ne\varepsilon_b$) and $\Delta'_{cd}\Phi=0$ ($\varepsilon'_c\ne\varepsilon'_d$). A
\emph{pre-joint harmonic} satisfies the same conditions except $E'_{cc}\Phi=\Lambda'_c\Phi$.
\end{quoteddef}

\begin{quotedthm}{\cite[Lemma~3.10]{P49}}{}
\textup{(i)} $w_r\parsgn\xi_r\ge0$ for all $r$. \textup{(ii)} $w_r=-w'_{c(r)}$, and $w'_c=0$ for $c\in B$.
\textup{(iii)} If $\varepsilon_a=\varepsilon_b$ and $\lambda_a>\lambda_b$ then $\Lambda_a\ge\Lambda_b$;
likewise for $\Lambda'$.
\end{quotedthm}

\begin{quotedthm}{\cite[Proposition~3.11]{P49}}{ (Existence)}
For $r\in R$ and an integer $0\le\ell<|w_r|$, call $(r,c(r))$ a pair of \emph{level} $\ell$. For each
$\ell$ and each pair of signs $(\epsilon,\epsilon')$ let $L_{\ell,\epsilon,\epsilon'}$ be the set of
rows $r$ having a pair of level $\ell$ with $\varepsilon_r=\epsilon$, $\varepsilon'_{c(r)}=\epsilon'$, and
$M_{\ell,\epsilon,\epsilon'}=\det(z_{r\,c(r')})_{r,r'\in L_{\ell,\epsilon,\epsilon'}}$. Then
$\Phi=\prod_{\ell,\epsilon,\epsilon'}M_{\ell,\epsilon,\epsilon'}$ is a joint harmonic. Moreover no monomial
of $\Phi$ contains two variables of one column from rows of opposite signs, or of one row from columns
of opposite signs.
\end{quotedthm}

\begin{quotedthm}{\cite[Lemma~3.14]{P49}}{}
Put $\Omega=\sum E_{ab}E_{ba}$, $\Omega'=\sum E'_{cd}E'_{dc}$, $\mathcal N=\sum_{a,c}a^+_{ac}a^-_{ac}=
\sum_aE_{aa}-n\gamma$ and $s_0=\gamma+\gamma'-\tfrac{n+n'}2$. On polynomials,
$\Omega'=\Omega-2s_0\mathcal N+n^2n'-2\gamma'nn'+n'\gamma'^2-n\gamma^2$. If $(\lambda,\lambda')$ match,
then on every polynomial $\Phi$ with $\sum_aE_{aa}\Phi=(\sum_a\lambda_a)\Phi$ the scalars satisfy
$\kappa_{\lambda'}=\kappa_\lambda-2s_0(\sum_a\lambda_a-n\gamma)+n^2n'-2\gamma'nn'+n'\gamma'^2-n\gamma^2$.
\end{quotedthm}

\begin{quotedthm}{\cite[Theorem~3.15]{P49}}{ ((A3))}
Assume {\cite[Lemma~3.6]{P49}} (that is, (V)), the conclusion of {\cite[Theorem~3.2]{P49}}, and
\textup{(D1)--(D3)}. Let $0\le j<m$. A joint harmonic $\Phi_j=\bigotimes_\tau\Phi_\tau$ exists and is
unique up to a scalar, and
\[
\int_{G_{j,\infty}}e_j(x)\,\omega_\infty(x^{-1},1)\Phi_j\,dx=\int_{G_{j+1,\infty}}e_{j+1}(x)\,\omega_\infty(1,x)
\Phi_j\,dx .
\]
\end{quotedthm}

\begin{quotedthm}{\cite[Theorem~3.17]{P49}}{}
Assume $h\ge n^2+1$. Then $F=c_\lambda$ on $G$; that is, {\cite[(59)]{P51}} holds.
\end{quotedthm}

\begin{quotedthm}{\cite[Theorem~3.29]{P49}}{}
Assume \textup{(W1$_\infty$)}, \textup{(W2$_\infty$)}, \textup{(M)} for every nondegenerate hermitian
space over $E_\tau$ (in particular for $W$ and, if $W$ is definite, for $W\oplus W'$ as in Step~3),
and, if $W$ is definite, \textup{(S)}. Then \textup{(V$'$)} holds; that is, \textup{(V)} of {\cite[Lemma~3.6]{P49}} holds at
$\tau$, with $\upsilon_\tau|_{\auU(1)}(z)=z$.
\end{quotedthm}

\begin{quotedthm}{\cite[Lemma~4.5]{P49}}{}
$a_1$ is an algebraic integer and the discriminant $a_1^4-64$ of $\mathcal E$ is a unit at every
odd prime. The reduction $\widetilde{\mathcal E}$ modulo $\pam$ is a supersingular elliptic curve,
that is, it has no point of order $p$ over the algebraic closure of the residue field.
\end{quotedthm}

\begin{quotedthm}{\cite[Lemma~4.6]{P49}}{ (The parameter)}
Let $T(z)=\frac1{q-1}\sum_{\zeta\in\mu_{q-1}}\zeta^{-1}y(\zeta z)$, the part of $y$ in the degrees
$\equiv1\pmod{q-1}$.
\begin{enumerate}
\item[(a)] $T\in y+y^2R$, so $R=\pao[[T]]$; and $T(\pi z)=\pi T+\sum_{k\ge1}b_kT^{1+k(q-1)}$ with
$b_k\in\pao$ and $b_1\in\pao^\times$.
\item[(b)] Let $\pavp^{q-1}=-\pi/b_1$; then $v(\pavp)=\rho$. The series $Y(z)=T(\pavp z)/\pavp$ is
integral, and $Y\equiv\sum_{j\ge0}z^{q^j}\pmod\pi$.
\item[(c)] For every $G\in R$ the series $G(\pavp z)$ is integral, and the coefficients of $G$
itself in the degrees $<q$ in $z$ are integral. In particular this holds for $x$ and $y$.
\end{enumerate}
\end{quotedthm}

\begin{quotedthm}{\cite[Lemma~4.7]{P49}}{ (Units)}
\begin{enumerate}
\item[(a)] For $\alpha\in1+2\auO_K$, $\alpha\ne0$, the function $f_\alpha(z)=\theta(\alpha z)/
\theta(z)^{\auNm\alpha}$ is a polynomial $R_\alpha(x)$ in $x$ of degree $(\auNm\alpha-1)/2$ with
$R_\alpha\in\pao[x]$, $R_\alpha(0)=1$. In particular $f_\alpha\in1+yR$.
\item[(b)] For $\zeta\in\mu_{q-1}$, $\exp\bigl(\zeta^{p+1}g(z)-g(\zeta z)\bigr)\in1+yR$.
\item[(c)] $\pi g_2\in\pao$, and $\pi g_n\in\pao$ for all $n<q$, where $g=\sum g_nz^n$.
\end{enumerate}
\end{quotedthm}

\begin{quotedthm}{\cite[Theorem~4.9]{P49}}{ (Normal form on one curve)}
Let $\pavp$ be as in {\cite[Lemma~4.6(b)]{P49}} and $c_\pavp=\pavp^{p+1}/p$. There are integral series
$W,\Xi_1\in\pao[[z]]$, $\Xi_1(0)=0$, and $\lambda\in\pao$ such that
\[
\theta(\pavp z)=W(z)\exp\bigl(c_\pavp\,\Xi_1(z)\bigr),\qquad
\Xi_1\equiv\sum_{e\in\paI}\sum_{j\ge0}\lambda_{e,j}\bigl(z^{p^e}\cdot z\bigr)^{p^j}\pmod\pi ,
\]
where $\lambda_{e,j}=\lambda$ if $p$ is inert, and $\lambda_{0,0}=\lambda/2$,
$\lambda_{0,j}=\lambda$ $(j\ge1)$, $\lambda_{e,j}=2\lambda$ $(e\ge1)$ if $p=d$.
\end{quotedthm}

\begin{quotedthm}{\cite[Proposition~4.12]{P49}}{ (Normal form of theta products)}
Let $(u_j,m_j)$ be a resolution of the identity by vectors whose coordinates are integral at all
primes above $p$, let $\mu_j=u_j^\dagger$, and let $G$ be an integral power series in the series
$y(\mu_j(z))$. Then, for every $c\in\pao$, the series
\[
F(z)=F_0(cz),\qquad F_0(z)=\prod_j\theta(\mu_j(z))^{m_j}\cdot G(z),
\]
has the normal form {\cite[(20)]{P49}}, with $\pavp$ as in {\cite[Lemma~4.6(b)]{P49}}, the index set
$\paI$ and the partner maps $\mathrm{pr}_e(\tau)=u(-p^e)\tau$.
\end{quotedthm}

\begin{quotedthm}{\cite[Lemma~4.13]{P49}}{ (The kernels)}
Let $\ell\ne N$, let $\vartheta$ be a twist with $\ell\nmid\auNm Q_\vartheta$, $\Psi$ a chain datum
of $\vartheta$, and $\Psi'$ as in {\cite[Proposition~14.25]{P46}}, so that
$B_{\vartheta,\Psi}=\operatorname{res}\circ C_{\vartheta,\Psi'}$.
\begin{enumerate}
\item[(K0)] For $\phi\in\auS$ and $g_m\in\auU(X_m)(\auA_{E^+,f})$,
\[
(C_{\vartheta,\Psi'}\phi)(g_m)=\sum_{g_0,\dots,g_{m-1}}\ \sum_x\Psi'(x)\;\theta_{m-1}
(\phi^\circ_{m-1,x_{m-1}};g_{m-1},g_m)\cdots\theta_0(\phi^\circ_{0,x_0};g_0,g_1)\,\phi(g_0),
\]
where $g_i$ runs over representatives of the classes of the level of $X_i$ and $\theta_i$ is the
kernel of {\cite[\S14.8]{P46}}.
\item[(K1)] For every $x$ in the support of $\Psi'$, every $j<m$ and all $g,g'$ there is a power
series $F=F_{x,j,g,g'}$ in the variables $z_{\tau,b,a}$ ($\tau\in\auPl$, $b\le n_{j+1}$, $a\le n_j$)
such that the matrix of $g'^{-1}_\ell\circ\theta_j(\phi^\circ_{j,x_j};g,g')\circ g_\ell$ in the
bases $(e_T)$ of {\cite[\S6.1]{P49}} is, up to multiplication on both sides by matrices that are
invertible over $\auO_{\C_\ell}$, the matrix
$\bigl(S_{\lambda_{j+1}}(\partial)_{T',T}F(0)\bigr)_{T',T}$ of jets along the entries of the Schur
matrix. If $\ell\ne d$, all jets $Q(\partial)F(0)$ with $Q\in\auO_{\C_\ell}[z]$ are integral
({\cite[Proposition~5.30]{P49}} for $\ell\nmid2q_1M$; {\cite[Corollary~7.16(a)]{P49}} for all $\ell\ne d$).
\item[(K2)] If $\ell=p\in\mathcal N$, then $F$ is one of the series of
{\cite[Proposition~4.12]{P49}}, for $\mathsf C=[n_{j+1}]\times[n_j]$, the resolution of the identity
$u_{k,i}=(\zeta_{q_1}^k/q_1)e_i$ $(k<q_1,\ i\in\mathsf C)$ with multiplicities $q_1$, the lattice
$\Lambda=\Omega_0\auO_K$ and $c=\pi_M\varpi_{\mathrm{vac}}/\mathsf u$.
\item[(K1$_N$)] Let $\ell=N$ and $\iota\colon\C\to\C_N$. Let $z\in\operatorname{Hom}(X_j,X_{j+1})(\auA_{E^+,f})$
be integral at the primes above $N$ that are conjugate to the primes induced by the $\iota\circ\tau$,
and such that $Nz$ is integral at the primes induced by the $\iota\circ\tau$. Then for all $g,g'$ the
matrix of $g'^{-1}_N\circ\theta_j(\phi^\circ_{j,z};g,g')\circ g_N$ in the bases $(e_T)$ has entries in
$\auO_\iota$. (This is {\cite[Lemma~6.20]{P49}}.)
\end{enumerate}
\end{quotedthm}

\begin{quotedthm}{\cite[Lemma~5.2]{P49}}{ (Standard facts)}
Let $W$ be a hermitian space over $E$.
\begin{enumerate}
\item[(W1)] $\omega(g)\psrho(w)\omega(g)^{-1}=\psrho(gw)$. There is a functional $\theta=\theta_W$
on $\psHS(W)$, unique up to a scalar, with $\theta\circ\psrho(w)=\theta$ for $w\in W(E)$; it
satisfies $\theta\circ\omega(\gamma)=\theta$ for $\gamma\in\auU(W)(E^+)$.
\item[(W2)] Let $I\subset W$ be an isotropic $E$-subspace, $W_I=I^\perp/I$, and $P_I$ the stabiliser
of $I$ in $\auU(W)$. There is a continuous linear map $\psZ_I\colon\psHS(W)\to\psHS(W_I)$ (the
\emph{zero term}) with $\theta_{W_I}(\psZ_IF)=\int_{I(E)\backslash I(\auA)}\theta(\psrho(i)F)\,di$,
$\psZ_I\psrho(w)=\psrho(\bar w)\psZ_I$ for $w\in I^\perp(\auA)$ with image $\bar w$, and
$\psZ_I\omega(p)=\chi_I(p)\,\omega_{W_I}(\bar p)\psZ_I$ for $p\in P_I(\auA)$, where $\chi_I(p)=
(\upsilon|\cdot|^{-1/2})(\det p|_I)^{-1}=\upsilon(\det p|_I)^{-1}|\det p|_I|^{1/2}$, with
$\det p|_I\in\auA_E^\times$ the determinant on the $\auA_E$-module $I(\auA)$ and $|\cdot|$ the
idele norm of $E$; so $\chi_I=1$ on unipotent elements and on $P_I(E^+)$.
For isotropic $I\subset I'$ one has $\psZ_{I'}=\psZ_{I'/I}\psZ_I$, and $\psZ_I=\psZ_{I,\infty}
\otimes\psZ_{I,f}$.
\item[(W3)] For $F\in\psHS(W)$, $\theta(F)=\sum_{w\in W(E)/I^\perp(E)}\theta_{W_I}(\psZ_I\psrho(w)F)$,
the series converging absolutely. If $W=\auHom(V,V')$, $F=\Psi\otimes f$ with $\Psi$ smooth, and
$\mathcal C$ is a compact set of finite adelic points of $\auU(V)\times\auU(V')$, there are $C,A$ with
$\sum_w|\theta_{W_I}(\psZ_I\psrho(w)\omega(g,h)F)|\le C(\|g_\infty\|\|h_\infty\|)^A$ whenever
$(g_f,h_f)\in\mathcal C$.
\item[(W4)] Let $L\subset W(\auA_f)$ be a lattice that is its own dual for $[\,,]$. The vectors of
$\psHS_f(W)$ on which $\psrho(l)$, $l\in L$, acts by a fixed character form a line. If $\psvac_L$
spans the line of the character $\epsilon_L$ with $\epsilon_L(l)\epsilon_L(l')=\beta(l,l')
\epsilon_L(l+l')$ fixed below, and $\psrho^\circ(x)=\epsilon(x)\psrho(x)$ for a fixed function
$\epsilon$ with values in roots of unity extending $\epsilon_L^{-1}$, then
$\psrho^\circ(l)\psvac_L=\psvac_L$, $\psrho^\circ(x+l)\psvac_L=\zeta\psrho^\circ(x)\psvac_L$ with a
root of unity $\zeta$, the $\psrho^\circ(x)\psvac_L$, $x\in W(\auA_f)/L$, form a basis of
$\psHS_f(W)$, and $\omega(g)\psvac_L\in\C^\times\psvac_{gL}$ for $g\in\auU(W)(\auA_f)$. If
$g\in\auU(W)(\auA_f)$ stabilises $L$ then $\omega(g)\psvac_L=\psvac_L$.
\item[(W5)] At a place $v=w\aubar w$ of $E^+$ split in $E$, $\psHS_v(\auHom(V,V'))$ is the space
of Schwartz--Bruhat functions on $\auHom(V_w,V'_w)$, with $(\omega_\mu(g,h)\Psi)(x)=\mu_w(\det h_w)
\mu_w(\det g_w)^{-1}\upsilon'_w(g_w,h_w)\,\Psi(h_w^{-1}xg_w)$, where, for even ranks $n,n'$, $\mu_w\upsilon'_w$ is given by {\cite[(62)]{P49}}: the factor is
$\chi_w(\det h_w)|\det h_w|_w^{-n/2}\chi_w(\det g_w)^{-1}|\det g_w|_w^{n'/2}$ with $\chi$ the base
change of $\mu$; for arbitrary ranks the untwisted operator is
\[\Psi\mapsto\upsilon_w(\det h_w)^{-n}\upsilon_w(\det g_w)^{n'}|\det h_w|_w^{-n/2}|\det g_w|_w^{n'/2}\Psi(h_w^{-1}xg_w)\]
({\cite[Proposition~2.17]{P49}}), so that all factors multiplying $\Psi(h_w^{-1}xg_w)$ are built from $\mu_w$, $\upsilon_w$ and $|\cdot|_w$ evaluated at $\det g_w,\det h_w$; $\psrho(a_w,0)\Psi(x)=\Psi(x+a_w)$
and $\psrho(0,a_{\bar w})$ is multiplication by the character $x\mapsto[x,a_{\bar w}]$.
\item[(W0), (W6)] Two further facts are used; they are stated in full in {\cite[\S5.2.1]{P49}}.
(W0): $\theta$ is defined on the smooth vectors of the Fock space tensored with $\psHS_f(W)$ and is
continuous in the first factor. (W6): for hermitian spaces $V,V'$ of ranks $n,n'$ the map
$w\mapsto w^\dagger$ identifies $\psHS(\auHom(V,V'))$ with $\psHS(\auHom(V',V^-))$, carrying
$\omega_{\auHom(V,V')}(g,g')$ to $\upsilon(\det g)^{n'}\upsilon(\det g')^{-n}\,
\omega_{\auHom(V',V^-)}(g',g)$. It carries the self-dual lattice $\alpha\auHom(\psM,\psM')$ to
$\aubar\alpha\auHom(\psM',\psM)$, which differs from $\alpha\auHom(\psM',\psM)$ at the places above $q_1$
({\cite[Remark~7.18]{P49}}).
\end{enumerate}
\end{quotedthm}

\begin{quotedthm}{\cite[Lemma~5.7]{P49}}{ (Conjugates of jets)}
Let $X,X'$ be standard of ranks $n<n'$, $W=\auHom(X,X')$, and $\sigma\in\auAut(\C)$ with
$\sigma|_K=1$. There are $t_\sigma\in\C$ with $|t_\sigma|^2=\Omega^{-1}$, a permutation
$\varpi_\sigma$ of $\auPl$, and a $\sigma$-semilinear map $\Lambda_\sigma$ of $\psHS_f(W)$ such
that for all $f\in\psHS_f(W)$:
\begin{enumerate}
\item[(i)] the Taylor series of $\Theta_{\Lambda_\sigma f}(t_\sigma\,\cdot)$ is obtained from that
of $\Theta_f(s_0\,\cdot)$ by applying $\sigma$ to the coefficients and $\varpi_\sigma$ to the
variables $z_\tau$;
\item[(ii)] if $f=\psrho^\circ(x)\omega(g)\psvac^\circ$ with $x\in W(\auA_f)$, $g\in\auU(W)
(\auA_f)$, then $|\theta(\psrho_\infty(y)(1\otimes\Lambda_\sigma f))|\le\exp\bigl(C_0(q_1nn'+1)^2
\bigr)$ for all $y$, where $C_0$ depends only on $q_0$ and $d$.
\end{enumerate}
\end{quotedthm}

\begin{quotedthm}{\cite[Lemma~5.22]{P49}}{ (The hook number)}
Let $\langle P,Q\rangle_F=(P(\partial)\aubar Q)(0)=\sum_\alpha\alpha!\,P_\alpha\aubar{Q_\alpha}$ be the
Fischer form on polynomials in the entries of $y$, and $H_\mu=\langle f_\mu,f_\mu\rangle_F$. Then
$H_\mu$ is a positive integer, $H_\mu\le|\mu|^{3|\mu|}$ $(0^0=1)$, the invariant form of $W_\mu$ is
$H_\mu^{-1}\langle\ ,\,\rangle_F$, and $H_\mu\Sigma_\mu\subseteq\Sigma_\mu^\#$.
\end{quotedthm}

\begin{quotedthm}{\cite[Proposition~5.25]{P49}}{}
Let $(\Lambda,t,a_1)$ be a normalised lattice.
\begin{enumerate}
\item[(C1)] {\rm(Conjugation.)} For $\sigma\in\auAut(\C/K)$ there is a normalised lattice
$(\Lambda^\sigma,t^\sigma,\sigma a_1)$, unique, such that $\sigma$ maps the Taylor series at $0$ of
$x$, $y$, $\theta$ to those of $\Lambda^\sigma$. One has $g_i(\Lambda^\sigma)=\sigma g_i(\Lambda)$
$(i=2,3)$ and $s_2(\Lambda^\sigma)=\sigma s_2(\Lambda)$.
\item[(C2)] {\rm(Integrality.)} For every prime $\ell\ne d$ and every $\C\to\C_\ell$, $x$, $y$ and
$\theta$ are $\partial$-integral.
\item[(C3)] {\rm(Translates.)} Let $u\in\C$ with $\mathsf Nu\in\Lambda$, $\mathsf N\ge1$,
$u\notin t+\Lambda$.
\begin{itemize}
\item[(i)] Let $\ell$ be an odd prime not dividing $\mathsf N$ and $\C\to\C_\ell$. Then
$\theta_u(0)\in\pao$ (this holds at every prime, {\cite[Proposition~7.5]{P49}}) and $\theta_u(z)=\theta_u(0)\,\theta(z)\,\mathsf G_u(z)$ with
$\mathsf G_u=\Gamma_u(y(z))$, $\Gamma_u\in1+Y\pao[[Y]]$; and $\mathsf G_u$ is $\partial$-integral.
If $\ell\ne d$, $\theta_u$ is $\partial$-integral.
\item[(ii)] Let $\Lambda=\Omega_0\auO_K$ and $\sigma\in\auAut(\C/K)$. There are $\Omega_\sigma\in\C^\times$ and an ideal $\mathfrak c$ of $\auO_K$ with
$\Lambda^\sigma=\Omega_\sigma\mathfrak a$, $\mathfrak a=\mathfrak c\aubar{\mathfrak c}^{-1}$, and a finite idele
$\beta_\sigma$ of $K$ generating $\mathfrak a\widehat\auO_K$, such that for every $\mathsf x\in K$, every $n\ge1$
with $n\mathsf x\in\auO_K$ and every $\mathsf x'\in K$ with $\mathsf x'-\beta_\sigma\mathsf x\in2n\mathfrak a\widehat\auO_K$,
$\sigma$ maps the Taylor series of $\theta_{\mathsf x\Omega_0}$ to that of the twisted translate
$\theta^\sigma_{\mathsf x'\Omega_\sigma}$ of the theta function of $\Lambda^\sigma$.
\end{itemize}
\item[(C4)] {\rm(Chowla--Selberg.)} There is an algebraic number $\mathsf u\ne0$, a unit at every prime, with $\operatorname{area}(\Lambda^\sigma)\,|\sigma\mathsf u|^2=2\boldsymbol\pi^2e^{c_{CS}}/
\sqrt d$ for all $\sigma\in\auAut(\C/K)$, where $c_{CS}=\gamma-\Rel\frac{L'}L(1,\chi_K)$.
\end{enumerate}
\end{quotedthm}

\begin{quotedthm}{\cite[Lemma~5.26]{P49}}{ (Dictionary)}
The map $f\mapsto\Theta_f$ is a bijection from $\psHS_f(W)$ onto the space of entire functions
$\Theta$ on $\C^{\auPl\times\mathsf C}$ for which there are an integer $\mathsf N\ge1$ and a function
$\epsilon$ on $\mathsf N\Lambda_L$, with values in the roots of unity, with
\[
\Theta(z+\lambda^\natural)=\epsilon(\lambda)\,e^{\langle z,\lambda^\natural\rangle+|\lambda^\natural|^2
/2}\,\Theta(z)\qquad(\lambda\in\mathsf N\Lambda_L),\qquad\langle z,w\rangle=\sum\aubar w z .
\]
For $\gamma\in W(E)$, $\Theta_{\psrho_f(\gamma)f}(z)$ is a root of unity times
$e^{\langle z,\gamma^\natural\rangle-|\gamma^\natural|^2/2}\,\Theta_f(z-\gamma^\natural)$.
\end{quotedthm}

\begin{quotedthm}{\cite[Proposition~5.27]{P49}}{ ({\cite[Lemma~5.7]{P49}})}
{\cite[Lemma~5.7]{P49}} holds, with
$t_\sigma=\sigma(\varpi_{\mathrm{vac}}/\mathsf u)\,c_\psi/\Omega_\sigma$, $\varpi_\sigma$ the
permutation $\tau\mapsto\sigma\circ\tau$ of $\auPl$, and $C_0=C_0(q_0,d)$ given at the end of the proof.
\end{quotedthm}

\begin{quotedthm}{\cite[Proposition~5.30]{P49}}{ ({\cite[Lemma~4.13(K1)]{P49}}, (K2))}
{\rm(K1)} and {\rm(K2)} of {\cite[Lemma~4.13]{P49}} hold.
The integrality of the jets in (K1) holds for every $\ell\ne d$, by (S2) and {\cite[Theorem~7.15]{P49}}.
\end{quotedthm}

\begin{quotedthm}{\cite[Corollary~6.4]{P49}}{}
Let $X$ be a standard space, $G=\auU(X)$, $\ell$ a prime.
{\rm(a)} For $g,d\in G_f$ there is $\gamma\in\Gamma$ such that $\gamma g$ and $d\,\gamma g\,d^{-1}$
are integral above $\ell$. In particular every class of $\Gamma\backslash G_f/J$ has a
representative that is integral above $\ell$.
{\rm(b)} If $\varphi$ is a form of level $J\subseteq K^\circ$ and $(x)$ a system of representatives
that are integral above $\ell$, and if $J$ is the full group $\auGL_n(\auO_w)$ at the places above
$\ell$, then $\varphi$ is integral at $\iota$ if and only if all $\iota(\varphi(x))$ have
coordinates in $\pao$.
\end{quotedthm}

\begin{quotedthm}{\cite[Lemma~6.8]{P49}}{ (The vacuum)}
Let $\alpha=\chi_{j+1,\mathfrak l}(\varpi)q^{n/2}$ and $\beta=\chi_{j+1,\mathfrak l}(\varpi)
q^{n'/2}$. For every $v$, the vector $\plone\otimes v$ is fixed by $\omega(K,K')$, and
\[
P'(X)(\plone\otimes v)=\prod_{l<k_0}(1-q^l\alpha X)\cdot P^-\bigl(q^{k_0}\alpha\beta^{-1}X\bigr)
(\plone\otimes v).
\]
\end{quotedthm}

\begin{quotedthm}{\cite[Proposition~6.9]{P49}}{ (One lift)}
Let $\Phi=\plone\otimes v$ be a datum that is the vacuum at $\mathfrak l,\aubar{\mathfrak l}$, and
let $\varphi$ be a smooth form on $X_j$ fixed by $K$. Then $\Theta^X_j(\Phi)\varphi$ is fixed by
$K'$, and with $z=q^{1/2}$ and $\alpha$ as in {\cite[Lemma~6.8]{P49}},
\[
P_{\mathfrak l}(X)\,\Theta^X_j(\Phi)\varphi=\prod_{l<t+1}(1-q^l\alpha X)\cdot\Theta^X_j(\Phi)\bigl(
P_{\mathfrak l}(z^{\,t+1}X)\varphi\bigr).
\]
\end{quotedthm}

\begin{quotedthm}{\cite[Lemma~7.1]{P49}}{}
Let $(\Lambda,t,a_1)$ be a normalised lattice and $\chi$ the function of {\cite[(27)]{P49}}.
\begin{enumerate}
\item[(a)] For $\lambda\in\Lambda$: $\chi(\lambda)=-1$ if $\lambda/2\in t+\Lambda$, and
$\chi(\lambda)=1$ otherwise. For any choice of $t\in\frac12\Lambda\setminus\Lambda$ this holds.
\item[(b)] $\chi(\aubar\delta\lambda)=\chi(\lambda)^{\auNm\delta}$ for all $\delta\in\auO_K$. If
$\Lambda=\Omega_0\auO_K$, then $\chi(\Omega_0a)=(-1)^{\auNm a}$ for $a\in\auO_K$.
\item[(c)] For $u\in\C$ and $\lambda\in\Lambda$, $\theta_u(z+\lambda)=\chi(\lambda)e^{2\boldsymbol
\pi iE(\lambda,u)}e^{H(\lambda,z+\lambda/2)}\theta_u(z)$ and $\theta_{u+\lambda}=\chi(\lambda)
e^{\boldsymbol\pi iE(\lambda,u)}\theta_u$. In particular $\theta_{u+2n\lambda}=\theta_u$ if
$n\ge1$ and $nu\in\Lambda$, and $\theta_\lambda=\chi(\lambda)\theta$.
\item[(d)] Let $\delta\in\auO_K$, $\delta\ne0$, and $\delta u\in\Lambda$. Then $z\mapsto\theta_u
(\aubar\delta z)$ satisfies $G(z+\lambda)=\chi(\lambda)^{m}e^{mH(\lambda,z+\lambda/2)}G(z)$
$(\lambda\in\Lambda)$ with $m=\auNm\delta$.
\end{enumerate}
\end{quotedthm}

\begin{quotedthm}{\cite[Lemma~7.2]{P49}}{ (Good reduction everywhere)}
Let $(\Lambda,t,a_1)$ be a normalised lattice. Then $a_1$ is a unit at every place above $2$.
Consequently $a_1^4-64$ is a unit at every prime, the scheme $y^2+a_1xy=x^3+x$ over the ring of
integers of $\Q(a_1)$ is an elliptic curve, and at every $\iota$ the point $P_0=(0,0)$ does not
reduce to the origin.
\end{quotedthm}

\begin{quotedthm}{\cite[Corollary~7.3]{P49}}{ ({\cite[Proposition~5.25(C1)]{P49}}, (C2), (C4) at every prime)}
With normalised lattices in the sense above, {\rm(C1)} holds, {\rm(C2)} holds at every prime
$\ell\ne d$ including $\ell=2$, and the number $\mathsf u$ of {\rm(C4)} is a unit at every prime,
including $2$.
\end{quotedthm}

\begin{quotedthm}{\cite[Lemma~7.4]{P49}}{}
Let $W$ be a hermitian space and $L\subset W(\auA_f)$ a self-dual lattice. Then $\epsilon_L$,
defined by {\cite[(70)]{P49}}, has values $\pm1$. If $W=E^{\mathsf C}$ is standard and
$L=\alpha\widehat\auO_E^{\,\mathsf C}$, then $\epsilon_L(\lambda)=(-1)^{\pfT_{E^+/\Q}\langle\lambda,
\lambda\rangle}$ for $\lambda\in\Lambda_L=W(E)\cap L=\mathfrak a_0^{\mathsf C}$.
\end{quotedthm}

\begin{quotedthm}{\cite[Proposition~7.5]{P49}}{ (Values at torsion points)}
Let $u\in\Q\Lambda$. If $u\in t+\Lambda$ then $\pfbr u=0$. Otherwise $\pfbr u$ is an algebraic
number and, at every prime $\ell$ and every $\iota$,
\[
|\pfbr u|=1\ \text{ if }u-t\notin E_1,\qquad|\pfbr u|=|x(u)|^{-1/2}<1\ \text{ if }u-t\in E_1 .
\]
In particular $\theta_u(0)$ is an algebraic integer for every torsion point $u$.
\end{quotedthm}

\begin{quotedthm}{\cite[Proposition~7.6]{P49}}{ (Integrality principle)}
Let $G\in\mathcal S_m$ and suppose that $e^{-mH(z,z)/2}G(z)$ is integral for all torsion points $z$
of odd order prime to $\ell$. Then $G=\theta^m\cdot\Gamma(y)$ with $\Gamma\in\pao[[Y]]$, as power
series in $z$. In particular, if $\delta\in\auO_K$, $\delta\ne0$, and $\delta u\in\Lambda$, then
\[
\theta_u(\aubar\delta z)=\theta(z)^{\auNm\delta}\,\Gamma_{u,\delta}(y(z)),\qquad\Gamma_{u,\delta}\in
\pao[[Y]],
\]
at every prime $\ell$, and for $\ell\ne d$ the series $\theta_u(\aubar\delta z)$ is
$\partial$-integral. If moreover $\iota(\aubar\delta)$ is a unit and $\ell\ne d$, then $\theta_u$
itself is $\partial$-integral.
\end{quotedthm}

\begin{quotedthm}{\cite[Theorem~7.9]{P49}}{ ({\cite[Proposition~5.25{\rm(C3)(ii)}]{P49}})}
Let $\Lambda=\Omega_0\auO_K$ be normalised and $\sigma\in\auAut(\C/K)$. There are $\Omega_\sigma
\in\C^\times$, an ideal $\mathfrak c\ne0$ of $\auO_K$ and a finite idele $\beta_\sigma$ of $K$ such
that, with $\mathfrak a=\mathfrak c\aubar{\mathfrak c}^{-1}$ and $\widehat{\mathfrak a}=\mathfrak
a\widehat\auO_K$:
\begin{enumerate}
\item[(a)] $\Lambda^\sigma=\Omega_\sigma\mathfrak a$, $\beta_\sigma\widehat\auO_K=\widehat{
\mathfrak a}$ (so $\beta_\sigma\aubar{\beta_\sigma}$ is a unit), and $S_\sigma(\mathsf x\Omega_0)
\equiv\mathsf x'\Omega_\sigma$ modulo $\Lambda^\sigma$ whenever $\mathsf x,\mathsf x'\in K$ and
$\mathsf x'-\beta_\sigma\mathsf x\in\widehat{\mathfrak a}$;
\item[(b)] if $\mathsf x\in K$, $n\ge1$, $n\mathsf x\in\auO_K$, and $\mathsf x'\in K$ satisfies
$\mathsf x'-\beta_\sigma\mathsf x\in2n\widehat{\mathfrak a}$, then $\sigma$ maps the Taylor series
at $0$ of $\theta_{\mathsf x\Omega_0}$ to that of $\theta^\sigma_{\mathsf x'\Omega_\sigma}$.
\end{enumerate}
\end{quotedthm}

\begin{quotedthm}{\cite[Lemma~7.11]{P49}}{ (A frame of length prime to $q_1$)}
There is a tight frame of a length $\mathsf h$ prime to $q_1$.
\end{quotedthm}

\begin{quotedthm}{\cite[Lemma~7.12]{P49}}{ (The scale)}
For every prime $\ell$, every $\iota$, every $\tau\in\auPl$ and every $\mathsf v\in\mathfrak a_0$,
the number $\varpi_{\mathrm{vac}}\cdot\aubar{\tau(\mathsf v)}$ is integral. More precisely
$|\varpi_{\mathrm{vac}}|=|\alpha_{w_\tau}|_{w_\tau}^{-1}$, where $w_\tau$ is the place of $E$
induced by $\iota\circ\tau\circ c$ and $|\cdot|_{w_\tau}=|\iota\tau c(\cdot)|$.
\end{quotedthm}

\begin{quotedthm}{\cite[Lemma~7.13]{P49}}{ (The product and the vacuum have the same character)}
$\Pi^v=\Theta_{F^v}$ for a unique $F^v\in\psHS_f(W)$; $\theta(1\otimes F^v)=1$;
$\psrho_f(l)F^v=\epsilon_L(l)F^v$ for all $l\in\mathsf bL$; and for $\gamma\in W(E)$
\begin{equation*}\tag*{{\cite[(73)]{P49}}}
\Theta_{\psrho_f(\gamma)F^v}(z)=\zeta_\gamma\prod_{i,\mathsf c}\theta_{\mathsf x_{i,\mathsf c}
\Omega_0}\bigl(a_{i,\mathsf c}(z)\bigr)^{\mathsf b},\qquad\mathsf x_{i,\mathsf c}=-\pfT_{E/K}(
\aubar{u_i}\gamma_{\mathsf c}),
\end{equation*}
with a root of unity $\zeta_\gamma$.
\end{quotedthm}

\begin{quotedthm}{\cite[Lemma~7.14]{P49}}{ (Values and expansion of the integral vacuum)}
{\rm(a)} For every $x\in W(\auA_f)$ the number $\theta(1\otimes\psrho^\circ(x)\psvac^\circ)$ is an
algebraic integer.
{\rm(b)} Let $(v_i)$ be a tight frame of length $\mathsf b$. Then $\psvac^\circ=\sum_ja_j\,
\psrho_f(y_j)F^v$, a finite sum with $y_j\in\mathsf b^{-1}L$ and $\mathsf b^{2f\#\mathsf C}a_j$
algebraic integers.
\end{quotedthm}

\begin{quotedthm}{\cite[Theorem~7.15]{P49}}{ (Standard data at every prime)}
Let $\ell$ be a prime, $\delta\in\auO_K$, $\delta\ne0$, and $f$ a standard datum of level $\delta$
at $\ell$. There are a resolution of the identity $(u'_j,m_j)_j$ of $E^{\mathsf C}$ with
$\alpha^{-1}u'_j$ integral above $\ell$, and an integral power series $\mathsf H$ in the series
$y(\mu_j(z))$, $\mu_j=(\varpi_{\mathrm{vac}}/\mathsf u)\,u_j'^{\dagger}$, such that
\[
\Theta_f\bigl(\aubar\delta\,s_0\,z\bigr)=\Bigl(\prod_j\theta(\mu_j(z))^{m_j}\Bigr)^{\auNm\delta}
\cdot\mathsf H(z),
\]
and the linear forms $\mu_j$ have integral coefficients. If $\ell\ne d$, the series
$\Theta_f(\aubar\delta s_0\,\cdot)$ is $\partial$-integral, and so is $\Theta_f(c\,\aubar\delta s_0
\,\cdot)$ for every $c\in\pao$.
\end{quotedthm}

\begin{quotedthm}{\cite[Corollary~7.16]{P49}}{ ({\rm(K1)} at every prime; the prime $N$)}
\begin{enumerate}
\item[(a)] The integrality assertion of {\cite[Lemma~4.13{\rm(K1)}]{P49}} holds for every prime
$\ell\ne d,N$ with $\ell\nmid\auNm Q_\vartheta$, in particular for $\ell=2$, for the primes
dividing $q_1$, and for $p_0,p_1$.
\item[(b)] {\cite[Lemma~6.20]{P49}} $(=$ {\rm(K1$_N$)}$)$ holds as stated, without the condition
{\rm(Gen)}.
\item[(c)] {\cite[Proposition~5.29{\rm(S3)}]{P49}}, hence {\cite[Lemma~4.13{\rm(K2)}]{P49}}, holds without
appeal to \cite{KubertLang}.
\end{enumerate}
\end{quotedthm}

\begin{quotedthm}{\cite[Proposition~7.20]{P49}}{}
Let $X=E^n$ be a standard space with $n\le n_m$, $G=\auU(X)$, $\Gamma=G(E^+)$, $K^\circ$ the
stabiliser of $\widehat\auO_E^{\,n}$ in $G_f$, and for an integer $\mathsf G\ge1$ let $K(\mathsf G)$
be the group of the $k\in K^\circ$ with $k\equiv1$ modulo $\mathsf G$. Then
\[
\#\bigl(\Gamma\backslash G_f/K^\circ\bigr)\le(2\mathsf B+1)^{4fn^2},\qquad\mathsf B=n^2C_{\mathrm{sb}}^2
\cdot2f\sqrt d,\quad C_{\mathrm{sb}}=6q_1d(fn)^2,
\]
and $\#(\Gamma\backslash G_f/J)\le(2\mathsf B+1)^{4fn^2}\mathsf G^{2fn^2}$ for every subgroup
$J\supseteq K(\mathsf G)$ of $G_f$. For $\mathsf G=MN$ this is at most $\exp(x^{30})$. In
particular the levels $J_{X_j}$ $(j\le m)$, and the level $J_1$ of the target forms for the trivial
twist, have at most $\exp(x^{30})$ classes.
\end{quotedthm}

\begin{quoteddef}{\cite[Definition~8.1]{P49}}{}
Let $e=e_j$ or $\overline{e_j}$. For a function $F$ on $G_j(\auA)$ put $(R(e)F)(g)=
\int_{G_{j,\infty}}e(x)F(gx)\,dx$ when the integral converges absolutely. A \emph{form of type $e$}
on $G_j$ is a bounded continuous function $y$ on $G_j(\auA)$, left invariant under $G_j(E^+)$, right
invariant under an open subgroup of $G_j(\auA_f)$ (a \emph{level}), with $R(e)y=y$. The
\emph{forms of the archimedean type of the chain} on $G_j$ are the forms of type $e_j$; their
complex conjugates are the forms of type $\overline{e_j}$. The space $\auAz$ of
{\cite[\S14.3]{P46}} is the space of the forms of type $e_0$ and level $J_M$; this is the precise
meaning of ``the fixed vector of the discrete series with parameter $\varsigma_\tau$''. For $j=m$
the forms of type $e_m$ are the functions $g\mapsto l\bigl(g_\infty^{-1}\varphi(g_f)\bigr)$ with
$\varphi$ an algebraic modular form of weight $\lambda^\circ$ and $l$ the linear form given by
$v_m$ (Peter--Weyl); in this way the last lift has values in algebraic modular forms.
\end{quoteddef}

\begin{quotedthm}{\cite[Lemma~8.2]{P49}}{ (Growth)}
Let $s=D+1$.
\begin{enumerate}
\item[(a)] For $g\in G_{j,\infty}$, $h\in G_{j+1,\infty}$ and $\Psi\in\pzF^{\mathrm{sm}}$,
$\|\omega_\infty(g,h)\Psi\|_s\le(2D\|g\|\|h\|)^{2s}\|\Psi\|_s$. Hence $\omega_\infty$ extends to a
strongly continuous representation on $\pzF_s$ with the same bound.
\item[(b)] $(1+|z|)^{2s}\,|\Psi(z)|\,e^{-|z|^2/2}\le(D+1)^{2s}\|\Psi\|_s$ for $\Psi\in\pzF_s$ and
$z\in\C^D$.
\item[(c)] For every $f\in\psHS_f(W)$ there is $C_f$ with $|\theta(\Psi\otimes f)|\le C_f\|\Psi\|_s$
for $\Psi\in\pzF^{\mathrm{sm}}$; so $\theta(\,\cdot\otimes f)$ extends continuously to $\pzF_s$.
\end{enumerate}
\end{quotedthm}

\begin{quotedthm}{\cite[Proposition~8.4]{P49}}{ (The lower chain is well defined)}
Assume {\rm(A1)--(A3)}. Let $0\le j\le m$.
\begin{enumerate}
\item[(a)] Every form $y$ of type $e_j$ or $\overline{e_j}$ on $G_j$ is smooth, $K_{j,\infty}$-finite
and finite under the centre of the enveloping algebra; it is a cusp form; it is rapidly
decreasing on Siegel sets, in the form: for every Siegel set $\mathfrak S_\infty\subset G_{j,\infty}$, finite set $\mathsf c\subset G_j(\auA_f)$, open subgroup $K'$ of $G_j(\auA_f)$ and $N\ge0$ there is $C_N$ with $|y(g)|\le C_N\|g_\infty\|^{-N}$ for $g\in\mathfrak S_\infty\mathsf cK'$; and the forms of a given level form a finite-dimensional space.
\item[(b)] Let $j<m$, $f\in\psHS_f(W)$ and let $y$ be a form of type $e_j$ on $G_j$. Then the
integral $F(h)=\int_{[G_j]}\theta\bigl(\omega_{\mu_{j+1}}(g,h)(\Phi_j\otimes f)\bigr)y(g)\,dg$ converges
absolutely and $F$ is a form of type $e_{j+1}$ on $G_{j+1}$.
\item[(c)] Let $j<m$, $f\in\psHS_f(W)$ and let $z$ be a form of type $\overline{e_{j+1}}$ on
$G_{j+1}$. Then $F^*(g)=\int_{[G_{j+1}]}\theta\bigl(\omega_{\mu_{j+1}}(g,h)(\Phi_j\otimes f)\bigr)z(h)
\,dh$ converges absolutely and $F^*$ is a form of type $\overline{e_j}$ on $G_j$.
\item[(d)] The space $\auAz$ is finite-dimensional and stable under the right translations by
the stabiliser of the lattice at the chosen places, so $\auG$ acts on it; the Poincar\'e map
$\auP$ of {\cite[\S14.3]{P46}} is well defined; and $\iota_\vartheta$ maps $\auAz$ to algebraic
modular forms of weight $\lambda^\circ$ on $\auU(Y_m)$.
\end{enumerate}
In particular the assumption made in {\cite[\S5.1.3]{P49}} holds: the lifts $\Theta^Y_{j,\vartheta}$ and
their transposes $\Theta^{Y,*}_{j,\vartheta}$ preserve the types, and all the forms concerned are
cusp forms and rapidly decreasing.
\end{quotedthm}

\begin{quotedthm}{\cite[Lemma~8.6]{P49}}{ (The characters of the lifts exist)}
Let $S$ be a finite set of primes of $E$, no two above the same place of $E^+$, each above a place
split in $E$; let $(c_\tau)_{\tau\in\auPl}$ be integers and $\psi_w$ ($w\in S$) characters of
$\auF_w^\times$ with $\prod_\tau\tau(\xi)^{c_\tau}\prod_{w\in S}\psi_w(\xi\bmod w)=1$ for every root of
unity $\xi\in E$. There is a unitary character $\mu$ of $\auU(1)(E^+)\backslash\auU(1)(\auA_{E^+})$ with
$\mu_\tau(u)=\tau(u)^{c_\tau}$ at the infinite places, with tame part $\psi_w$ at $w\in S$, and
unramified at all other finite places. If the reduction of the roots of unity of $E$ modulo some
$w_0\in S$ is injective (for instance a chosen place $w_0$), then for all $(c_\tau)$ and all $\psi_w$,
$w\ne w_0$, there is $\psi_{w_0}$ satisfying the condition. In particular the characters $\chi_M$,
$\mu^\flat$ and $\mu_j^0$ of {\cite[\S14.3]{P46}} exist (Definition~\ref{pfq:def-chars},
Proposition~\ref{pfq:tame}).
\end{quotedthm}

\begin{quotedthm}{\cite[Lemma~8.7]{P49}}{}
\begin{enumerate}
\item[(a)] $C_0(q_0,d)$ is a finite real number that depends only on $q_0$ and $d$: not on the
scales $x,x_1,b$, nor on $s$, $q_1$, $E$ or any other part of admissible data.
\item[(b)] It is the constant $C_0$ of {\cite[Proposition~5.27]{P49}} and of {\cite[Lemma~5.7{\rm(ii)}]{P49}}.
\item[(c)] With the constant $C_P$ of {\cite[(35)]{P46}} in condition {\rm(D2)} of
{\cite[Definition~13.3]{P46}}, which contains $d$ and $C_0(q_0,d)$ as summands (all summands are
non-negative), all admissible data satisfy
$C_0\le10^{-12}\log\log b$ and $d\le10^{-12}\log\log b$.
\end{enumerate}
\end{quotedthm}

\begin{quotedthm}{\cite[Lemma~4.7]{P51}}{}
Let $e_\tau$ be as in {\cite[(42)]{P51}}, on $G_\tau\cong G_s$.
\begin{enumerate}
\item[(a)] $e_\tau$ is continuous, $e_\tau^*=e_\tau$, $|e_\tau(g)|\le C\|g\|^{-B}$,
$e_\tau\in L^1$, and $e_\tau*e_\tau=e_\tau$.
\item[(b)] There is $\alpha_\tau\in C_c^\infty(G_\tau)$ with $\alpha_\tau*e_\tau=e_\tau$. Hence
$e_\tau=e_\tau*\alpha_\tau^*$, $e_\tau$ is smooth, and $|(De_\tau)(g)|\le C_D\|g\|^{-B}$ for every
left-invariant differential operator $D$.
\item[(c)] $\int_{G_\tau}\|g\|^{-B'}dg<\infty$ for $B'>n^2$.
\end{enumerate}
The same holds for $e$ on $G_\infty$ with $\alpha=\bigotimes\alpha_\tau$, with $\pmP(g)^{-B}$ in place
of $\|g\|^{-B}$.
\end{quotedthm}

\begin{quotedthm}{\cite[Theorem~5.6]{P51}}{ (\textup{(DS2)})}
For every $g\in G$ and every integer $0\le c\le h$ (formula {\cite[(59)]{P51}} holds by {\cite[Theorem~3.17]{P49}}),
\[
|c_\tau(g)|\le|\det a(g)|^{-c/2}\le\bigl(2/\|g\|\bigr)^c .
\]
In particular $|c_\tau(g)|\le(2/\|g\|)^{j+\pbeq}$ for $0\le j\le r^7$, $\pbeq=9r^2$, and for $g\in
G_\tau$ with the norm of {\cite[\S3.2.1]{P51}}.
\end{quotedthm}

\begin{quotedthm}{\cite[Theorem~5.9]{P51}}{ (\textup{(DS3)})}
For every real $\mathsf s>4N^2$, $\pbdeg_\tau\int_G\|g\|^{-\mathsf s}dg\le4(5/\delta)^{2N^2}\mathsf s/
(\mathsf s-4N^2)$. In particular $c_\tau$ is integrable, and
\[
\pbdeg_\tau\int_{G_\tau}\bigl(32/\|g\|\bigr)^{\pbeq}dg\le8\cdot32^{\pbeq}\bigl(60N(L+2N)\bigr)^{2N^2}
\le b^{99r^2}\qquad(\pbeq=9r^2).
\]
\end{quotedthm}

\begin{quotedthm}{\cite[Lemma~5.10]{P51}}{}
Let $\zeta\in E$ be a primitive $q_1$-th root of unity and $\varpi=\gamma\,\Phi_{q_1}'(\zeta)\sqrt{-d}$,
where $\gamma\in\auO_E^\times$ satisfies $\gamma\aubar\gamma=1$ and $\aubar{\gamma\Phi'_{q_1}(\zeta)}=
\gamma\Phi'_{q_1}(\zeta)$ \textup{(}such $\gamma$ exist, e.g.\ $\gamma=\zeta^s$ with
$s=q_1/4+(q_1-\varphi(q_1))/2+1$; see the first paragraph of the proof\textup{)}. Then $\varpi$ is a purely imaginary generator of the different of $E$, and
for every $\tau\in\auPl$, $2\le|\tau(\varpi)|\le b$ and $1\le M|\tau(\varpi)|/2\pi\le b^{23}$.
\end{quotedthm}

\begin{quotedthm}{\cite[Lemma~2.9]{P53}}{}
Let $\alpha,\beta$ be unramified outside $\pcm$ and tamely ramified at the primes of $\pcm$, with
tame parts giving characters $\alpha_\pcm,\beta_\pcm$ of $\pcR_\pcm^\times$. For $k\in K_S$,
$k'\in K_T$ with reductions $\bar k,\bar k'$ modulo $\pcm$,
\[
\omega_{\alpha,\beta}(k,k')[P,\Psi]=[P,\Psi'],\qquad\Psi'(x)=\alpha_\pcm(\det\bar k)^{-1}\beta_\pcm
(\det\bar k')\,\Psi(\bar k'^{-1}x\bar k).
\]
\end{quotedthm}

\begin{quotedthm}{\cite[Lemma~2.10]{P53}}{ (Archimedean types)}
Let $(S,T)$ be one of the pairs $(W_j,V_j)$, $1\le j\le m$, or one of the pairs
$(\pcG_{i,0},\pcG_{i,1})$, $0\le i<m$, of {\cite[\S2.5]{P53}}, with the types $e_S,e_T$ of the
chain and the joint harmonic polynomial $P$ of the pair. Then:
\begin{enumerate}
\item[(A1)] the types $e_S$, $e_T$ are integrable, and every element of $\pcF(S,e_S)$ and of
$\pcF(T,e_T)$ decreases faster than every power of the height on the automorphic quotient;
\item[(A2)] for every modulus $\pcm$, every $\Psi$ and all characters $\alpha,\beta$ of the weights
of the pair, tamely ramified at $\pcm$ and unramified elsewhere, $\Theta_{\alpha,\beta}(P,\Psi)$
maps $\pcF(S,e_S)$ to $\pcF(T,e_T)$;
\item[(A3)] the integral $Z_\infty=\int_{\auU(S)_\infty}\langle P,\omega_\infty(g,1)P\rangle\,\langle
v_0,\pi_S(g)v_0\rangle\,dg$ converges absolutely and is not zero; it does not depend on the twist,
which is trivial at the infinite places.
\end{enumerate}
\end{quotedthm}

\begin{quotedthm}{\cite[Lemma~2.11]{P53}}{ (The series)}
\begin{enumerate}
\item[(a)] $\pcmc_{M',\alpha}$ does not depend on the choice of $T$, is the same for $S$ and for
the space $S^-$ with the negative form, and for all spaces $S$ of rank $n$ under the identification
of their finite adelic groups; it is bi-invariant under $K^\pcm$; and
$|\pcmc_{M',\alpha}(\gamma)|=[L_S:L_S\cap\gamma L_S]^{-M'/2}$.
\item[(b)] If $n\ge2$ and $M'\ge6n$, then
$\sum_{\gamma\in G^\pcm/K^\pcm,\ \gamma\notin K^\pcm}|\pcmc_{M',\alpha}(\gamma)|\le\exp\bigl(2f(4/5)^n
\bigr)-1$.
\item[(c)] Let $n=r$ and $M'=r+t$. The \emph{series}
$\pcRal=\pcRal_{M',\alpha}=\sum_{\gamma\in G^\pcm/K^\pcm}\pcmc_{M',\alpha}(\gamma)\rho(\gamma)$
converges absolutely in the supremum norm on the bounded $K^\pcm$-invariant functions on
$\auU(S)(\auA)$, maps them to themselves, commutes with left translations and with right
translations by $\auU(S)_\infty$ and by the components at $S(\pcm)$, and satisfies
$\|\pcRal y\|_\infty\ge\tfrac12\|y\|_\infty$. In particular $\pcRal$ is injective, and it preserves
$\pcF(S,e)^J$ for every type $e$ and every open $J\supseteq K^\pcm$.
\end{enumerate}
\end{quotedthm}

\begin{quotedthm}{\cite[Lemma~2.13]{P53}}{ (Two standard facts)}
\begin{enumerate}
\item[(W7)] \emph{(complex conjugation).} Let $S,T$ be spaces of ranks $n,M'$ and $S^-$ the space
$S$ with the negative form. There is a conjugate-linear isometry
$\mathsf C\colon\psHS(\auHom(S,T))\to\psHS(\auHom(S^-,T))$ with
\[
\mathsf C\psrho(x)=\psrho^-(x)\mathsf C,\qquad\theta^-(\mathsf CF)=\aubar{\theta(F)},\qquad
\mathsf C\,\omega_{\auHom(S,T)}(g,h)\,\mathsf C^{-1}=\lambda(g,h)\,\omega_{\auHom(S^-,T)}(g,h),
\]
for a character $\lambda$ of $\auU(S)(\auA)\times\auU(T)(\auA)$ that is trivial on the rational
points and is a product of powers of $\upsilon\circ\det$.
\item[(SW)] \emph{(Weil's Siegel formula).} Let $\pcW$ be a hermitian space of rank $2n$ over $E$
that contains an isotropic subspace of dimension $n$, and let $T$ be a hermitian space of rank
$M'$ such that (i) the theta integral converges absolutely ($T$ anisotropic over $E^+$, or
$M'-r_T>n$ with $r_T$ the Witt index of $T$) and (ii) $M'>2n$; both hold if $M'\ge6n$. Let $\auU(T)(\auA)$ act on $\psHS(\auHom(\pcW,T))$ by the Weil representation
$\omega(1,h)$ with the trivial splitting character; for $\pcW=S\oplus S^-$ and
$\psHS(\auHom(\pcW,T))=\psHS(\auHom(S,T))\otimes\psHS(\auHom(S^-,T))$ this is
$\omega_{\auHom(S,T)}(1,h)\otimes\mathsf C\,\omega_{\auHom(S,T)}(1,h)\,\mathsf C^{-1}$. For a
Lagrangian (maximal isotropic subspace) $L\subset\pcW$ let
$I_L=\{w\in\auHom(\pcW,T):w|_L=0\}$, a maximal isotropic subspace of $\auHom(\pcW,T)$, and let
$\psZ_L=\psZ_{I_L}\colon\psHS(\auHom(\pcW,T))\to\C$ be its zero term \textup{(W2)}. Then for every
smooth $\Phi$
\[
\int_{[\auU(T)]}\theta\bigl(\omega(1,h)\Phi\bigr)\,dh=\pcvol([\auU(T)])\sum_L\psZ_L(\Phi),
\]
the sum being over all Lagrangians of $\pcW$; the integral and the sum converge absolutely, and
$\sum_L|\psZ_L(\omega(g,1)\Phi)|$ is bounded by a power of the height of $g\in\auU(\pcW)(\auA)$.
Moreover, for $\pcW=S\oplus S^-$, data $F,F'$ of $\auHom(S,T)$ and characters $\alpha,\beta$ there are
$C,A$ with
\[
\int_{[\auU(T)]}\bigl|\theta(\omega_{\alpha,\beta}(g,h)F)\,\theta(\omega_{\alpha,\beta}(g',h)F')\bigr|\,dh\le
C\,(\mathrm{ht}(g)\mathrm{ht}(g'))^A\qquad(g,g'\in\auU(S)(\auA));
\]
and if $T$ is anisotropic over $E^+$, $\int|\theta(\omega(g,h)\Phi)|\,dh$ is bounded by a power of the
height of $g\in\auU(\pcW)(\auA)$.
\end{enumerate}
\end{quotedthm}

\begin{quotedthm}{\cite[Lemma~2.14]{P53}}{ (Clean data)}
Let $S,T$ be as in \textup{(SW)}, $\pcW=S\oplus S^-$, and let $F\in\psHS(\auHom(S,T))$ be a datum
whose component at one split place $v=w\aubar w$ is, in the model \textup{(W5)}, a Schwartz--Bruhat
function supported on injective matrices. Let $F'\in\psHS(\auHom(S,T))$ and
$\Phi=F\otimes\mathsf CF'\in\psHS(\auHom(S,T))\otimes\psHS(\auHom(S^-,T))=\psHS(\auHom(\pcW,T))$.
\begin{enumerate}
\item[(a)] If the Lagrangian $L$ of $\pcW$ is not of the form $L_\delta=\{(\delta s,-s):s\in S\}$
for an element $\delta\in\auU(S)(E^+)$, then $\psZ_L(\Phi)=0$.
\item[(b)] There is a constant $c'\ne0$, depending only on the normalisations, such that for all
$\delta\in\auU(S)(E^+)$ and all $F,F'$
\[
\psZ_{L_\delta}(F\otimes\mathsf CF')=c'\,\bigl\langle F',\omega_{\auHom(S,T)}(\delta^{-1},1)F\bigr
\rangle .
\]
\end{enumerate}
\end{quotedthm}

\begin{quoteddef}{\cite[Definition~2.15]{P53}}{ (Sharp functions)}
Let $S,T$ be spaces of ranks $n,M'$, $\pcm$ a modulus, $\alpha$ a character tamely ramified at
$\pcm$, and $J$ a subgroup of $K_S$ containing $J_S(\pcm)$ and the components of $K_S$ outside
$S(\pcm)$, with image $\bar J\subseteq\auGL_n(\pcR_\pcm)$. A function $\Psi\ne0$ on
$\pcMat_{M'\times n}(\pcR_\pcm)$ is \emph{sharp for $J$} if
$\alpha_\pcm(\det\bar k)^{-1}\Psi(x\bar k)=\Psi(x)$ for $\bar k\in\bar J$ and there is a matrix
$A\in\pcMat_{n\times M'}(\pcR_\pcm)$ with $Ax\in\bar J$ for all $x$ in the support of $\Psi$. We put
$\|\Psi\|^2=\sum_x|\Psi(x)|^2$.
\end{quoteddef}

\begin{quotedthm}{\cite[Proposition~2.16]{P53}}{ (Inner product formula)}
Let $(S,T)$, of ranks $r$ and $r+t$, with types $e_S,e_T$ and polynomial $P$, be a pair for which
\textup{(A1)--(A3)} of {\cite[Lemma~2.10]{P53}} hold. Let $\pcm$ be a modulus with $S(\pcm)\ne\emptyset$ \textup(i.e.\ $\pcm\ne(1)$\textup), $\alpha,\beta$ characters
as in \textup{(A2)}, and $\Psi$ sharp for $J$. Put $\Theta=\Theta_{\alpha,\beta}(P,\Psi)$ and
$\pcRal=\pcRal_{r+t,\alpha}$. Then for all $y,y'\in\pcF(S,e_S)^J$
\[
\langle\Theta y',\Theta y\rangle=A_\infty\cdot\frac{\|\Psi\|^2}{\#\pcMat_{(r+t)\times r}(\pcR_\pcm)
\cdot[K_S:J]}\cdot\langle y',\pcRal\,y\rangle,\qquad A_\infty=c'\,Z_\infty\,\pcvol([\auU(T)])\ne0,
\]
where $A_\infty$ depends only on the pair and its archimedean data. Moreover $\Theta$ is injective
on $\pcF(S,e_S)^J$.
\end{quotedthm}

\begin{quotedthm}{\cite[Lemma~2.18]{P53}}{ (Algebraic forms and the Petersson product)}
Let $S$ be a definite space, $\lambda$ a weight with highest weight vector $v^+$, $J\subseteq K_S$
open and $\chi$ a character of $J$ of finite order. For $\varphi\in\mathcal A_\lambda(J,\chi)$ let
$f_\varphi(h)=\langle v^+,h_\infty^{-1}\varphi(h_f)\rangle$. Then
$\langle\varphi',\varphi\rangle=[K_S:J]\cdot\dim W_\lambda\cdot\|v^+\|^{-2}\cdot\langle f_{\varphi'},
f_\varphi\rangle_{[\auU(S)]}$. The statement and proof hold for every unit vector in place of $v^+$.
\end{quotedthm}

\begin{quoteddef}{\cite[Definition~3.1]{P53}}{ (Matched functions)}
Let $R$ be a commutative ring, $\chi$ a character of $R^\times$, $N=b+a$, the indices of $R^N$
ordered so that the first $b$ come first. For $k\in\auGL_N(R)$ let $k_{[b]}\in\pcMat_{N\times b}(R)$
be its first $b$ columns and $k^{[a]}\in\pcMat_{a\times N}(R)$ the last $a$ rows of $k^{-1}$.
Functions $\Phi$ on $\pcMat_{a\times N}(R)$ and $\Psi$ on $\pcMat_{N\times b}(R)$ are \emph{matched
for $\chi$} if $\Phi$ vanishes off $\{k^{[a]}\}$, $\Psi$ vanishes off $\{k_{[b]}\}$, and
$\Phi(k^{[a]})=\chi(\det k)\Psi(k_{[b]})$ for all $k\in\auGL_N(R)$. Over $F_S$ the same definition is
used with $k\in\auGL_N(\auO_S)$.
\end{quoteddef}

\begin{quotedthm}{\cite[Theorem~3.2]{P53}}{ (Exchange identity)}
Let $W,V,V'$ be spaces of ranks $N=a+b$, $a$, $b$, with $a=b+2e$, $e\ge0$ (so $N$ is even). Let
$q,q_V,q_{V'}$ be their numbers of negative signs; assume $q=q_V+q_{V'}$ and $2q+2\le N$. Let
$\mu=(\mu_\tau)$ be integers with $N\le2\mu_\tau+2q$ for all $\tau$, and let $v$ be a vector of the
Fock space of $\auHom(V',W)\otimes\R$ with finite moments. Then there are $C>0$ and a vector $x$ of
the Fock space of $\auHom(W,V)\otimes\R$ with finite moments, such that: for every modulus
$\pcm\ne(1)$ and every tame character $\nu$ for $\pcm$ of weight $\mu$ there is $c\in\C$ with
\[
|c|^2=C\cdot L_*(\pcm,\nu,e)\ne0,
\]
such that for all automorphic characters $\alpha,\alpha'$ of $\auU(1)$ with
\[
\alpha\alpha'=\gamma:=\nu^{-1}\upsilon^{N/2},
\]
all functions $\Phi$ on $\pcMat_{a\times N}(\pcR_\pcm)$ and $\Psi$ on $\pcMat_{N\times b}(\pcR_\pcm)$
matched for $\chi$, and all $g\in\auU(W)(\auA)$,
\begin{equation*}\tag*{{\cite[(12)]{P53}}}
\int_{[\auU(V)]}\theta^{W,V}_{\alpha,\gamma}(g,h;x,\Phi)\,dh=c\int_{[\auU(V')]}
\theta^{V',W}_{\gamma,\alpha'}(h',g;v,\Psi)\,dh' .
\end{equation*}
Here $\theta(\cdot\,;x,\Phi)$ is the kernel of the datum $[x,\Phi]$ of {\cite[\S2.3]{P53}} (matrices
in the frames), and $C$, $x$ do not depend on $\pcm,\nu,\alpha,\alpha',\Phi,\Psi,g$.
\end{quotedthm}

\begin{quotedthm}{\cite[Corollary~3.4]{P53}}{ (The constant is non-zero for every $e$)}
In {\cite[Theorem~3.2]{P53}}, $c\ne0$ for every $e\ge0$; for $e=0$, $|c|^2=C$.
\end{quotedthm}

\begin{quotedthm}{\cite[Lemma~3.24]{P53}}{ (Data modulo $\pcm$)}
Let $\nu$ be a tame character for $\pcm$, $S$ the primes of $\pcm$, and let $\Phi$, $\Psi$ be
functions on $\pcMat_{a\times N}(\pcR_\pcm)$, $\pcMat_{N\times b}(\pcR_\pcm)$ matched for $\chi$. Let
$\Phi'(x)=\Phi(x\bmod\pcm)$ for integral $x\in\pcMat_{a\times N}(F_S)$ and $\Phi'(x)=0$ otherwise,
and define $\Psi'$ in the same way. Then $\Phi'$, $\Psi'$ are matched for $\chi_S$ over $\auO_S$,
$\Psi'$ is locally constant, and $[x,\Phi]=[x,\Phi']_S$, $[v,\Psi]=[v,\Psi']_S$ for all Fock
vectors $x$, $v$ (the left sides being the data of {\cite[\S2.3]{P53}}).
\end{quotedthm}

\begin{quotedthm}{\cite[Proposition~3.26]{P53}}{ (Decay under a rank condition)}
\textup{(a)} Suppose that every $x$ with $\Phi(x)\ne0$ is surjective at every place of $P$ (each
component has a right inverse). For every $N_0$ there are $B,M,C$ such that for all
Fock vectors $y$ with finite moment $\|y\|_M=\int(1+|z|)^M|y(z)|\,dz$, all characters
$\alpha,\beta$ and all $(g,h)$,
\[
|\theta^{S_1,S_2}_{\alpha,\beta}(g,h;[y,\Phi]_P)|\,\mathrm{ht}(h)^{N_0}\le C\,\|y\|_M\,
\mathrm{ht}(g)^B .
\]
\textup{(b)} If every $x$ with $\Phi(x)\ne0$ is injective at every place of $P$, the same holds with
the roles of $g$ and $h$ exchanged.
\end{quotedthm}

\begin{quoteddef}{\cite[Definition~3.27]{P53}}{ (Lattice and height at a point)}
Let $S_1$ be a space of rank $n$ with frame giving $S_1(\auA)\cong\auA_E^n$. For $g\in\auU(S_1)(\auA)$
the \emph{lattice at $g$} is $\Lambda_g=\{\gamma\in S_1(E): g^{-1}\gamma\text{ is integral at all
finite places}\}$, with the length $|\gamma|_g=\|(g^{-1}\gamma)_\infty\|$ (the euclidean norm of the
archimedean coordinates). Let $\lambda_1(g)=\min\{|\gamma|_g:0\ne\gamma\in\Lambda_g\}$ and
$\mathrm{ht}(g)=\max(1,\lambda_1(g)^{-1})$.
\end{quoteddef}

\begin{quotedthm}{\cite[Lemma~3.28]{P53}}{ (Reduction to a Siegel set)}
Let $S_1$ be a space of rank $n$ with a frame, $G=\auU(S_1)$, viewed as a connected reductive group over
$E^+$, and $d=[E^+:\Q]$. There are a compact set $C_f\subset G(\auA_f)$ and a constant $C$ such that every
$g\in G(\auA)$ can be written $g=\gamma^{-1}x$ with $\gamma\in G(E^+)$, $x=x_\infty x_f$, $x_f,x_f^{-1}\in C_f$
and $\|x_\infty\|,\|x_\infty^{-1}\|\le C\,\mathrm{ht}(g)$ (operator norms in the
coordinates of the frame).
\end{quotedthm}

\begin{quotedthm}{\cite[Lemma~3.29]{P53}}{ (Growth of theta functions)}
Let $S_1,S_2$ be spaces, $X=\auHom(S_1,S_2)$, $D=\frac12\dim_\R(X\otimes\R)$, $s=D+1$, $\pzF_s$ and
$\|\cdot\|_s$ as in {\cite[Lemma~8.2]{P49}} for $W=X$, and $f_0$ a finite datum. There are $Y,A$ such that for
all characters $\alpha,\beta$, $g\in\auU(S_1)(\auA)$, $h\in\auU(S_2)(\auA)$, $z\in X\otimes\R$ and
$\Psi\in\pzF_s$,
\[
|\theta(\omega_{\alpha,\beta}(g,h)(\psrho_\infty(z)\Psi\otimes f_0))|\le Y\,\|\Psi\|_s\,
(\mathrm{ht}(g)\mathrm{ht}(h))^{A};\qquad A=2s\text{ will do.}
\]
\end{quotedthm}

\begin{quotedthm}{\cite[Lemma~3.30]{P53}}{ (Fock vectors are superpositions of translated vacua)}
Let $\mathcal F$ be the Fock space of a real hermitian space $\C^m$ (real dimension $2m$), realised as in
\S\ref{psp:ss-fock} as the space of functions $v\in L^1\cap L^2$ on $\C^m$ with $v(x)=\int v(y)\beta(y,x)
e(x-y)\,dy$, where $e(x)=e^{-|x|^2/2}$ is the vacuum, $|\beta|=1$, and the Heisenberg translations are
$(\rho(y)v)(x)=\beta(y,x)v(x-y)$; the measure is the normalised Lebesgue measure. For $N\ge0$ put
$\|v\|_N=\int(1+|x|)^N|v(x)|\,dx\in[0,\infty]$. Let $s\ge0$ be an integer, and let $\|\cdot\|_s$ and
$\mathcal F_s\subseteq\mathcal F$ be the norm and the Banach space of {\cite[Lemma~8.2]{P49}} (with $D=m$),
transported by the unitary isomorphism with the model {\cite[(41)]{P49}} that intertwines the
Heisenberg translations and carries vacuum to vacuum; the operators $\partial_w$, being the derivatives
of $\rho$, correspond. Let $c_\rho$ be a constant with
$\rho(w)\rho(y)=e^{\kappa(w,y)}\rho(y)\rho(w)$, $\kappa(w,y)\in i\R$ real bilinear,
$|\kappa(w,y)|\le c_\rho|w||y|$ (in the model {\cite[(41)]{P49}}, $\kappa(w,y)=2i\Iml\operatorname{tr}(y^*w)$ by the
formula for $\psrho_\infty(x)\psrho_\infty(y)$ recalled after it, so $c_\rho=2$ will do).
\begin{enumerate}
\item[\textup{(a)}] For $y\in\C^m$, $\rho(y)$ maps $\mathcal F_s$ into itself, with
$\|\rho(y)\Psi\|_s\le(1+c_\rho|y|)^{2s}\|\Psi\|_s$, and $y\mapsto\rho(y)\Psi$ is continuous into
$\mathcal F_s$ for every $\Psi\in\mathcal F_s$.
\item[\textup{(b)}] If $\|v\|_N<\infty$ then $v=\int v(y)\rho(y)e\,dy$, a Bochner integral converging for
the norm $\|\cdot\|_N$, because $\|\rho(y)e\|_N\le(1+|y|)^N\|e\|_N$. If $\|v\|_{2s}<\infty$, then
$v\in\mathcal F_s$, the same integral converges in $\mathcal F_s$, and $\|v\|_s\le c_1\|v\|_{2s}$ with
$c_1=\max(1,c_\rho)^{2s}\|e\|_s$.
\item[\textup{(c)}] Let $\ell$ be a continuous linear functional on $\mathcal F_s$ with $|\ell(\rho(z)e)|\le
S(1+|z|)^N$ for all $z$. Then $|\ell(v)|\le S\|v\|_N$ for every $v$ with $\|v\|_{\max(N,2s)}<\infty$.
\end{enumerate}
\end{quotedthm}

\begin{quotedthm}{\cite[Lemma~3.32]{P53}}{ (Comparison with the mean over an isotropic subspace)}
Let $X$ be a space (here $X=\auHom(S_1,S_2)$), $s$, the data of class $s$ and $\|\cdot\|_{f_0}$ as in
\textup{(R2)}, $f_0$ a finite datum (a finite combination of finite Heisenberg translates of the
vacuum), $l\ge0$ and $\iota$ a finite set. There is $C$ such that: if $k\in\auU(X)(\auA)$ (acting
through the Weil representation), $\delta\ge0$, $z_i\in\Lambda_k$ $(i\in\iota)$ with $|z_i|_k\le\delta$
span an isotropic subspace $Z\subseteq X(E)$, $\Lambda$ is a linear functional on the data of class $s$,
invariant under $\rho(Z(E))$, such that $x\mapsto\Lambda(\rho(x)F)$ is continuous on $Z(\auA)$ for every
datum $F$ of class $s$ and $\|\Lambda\circ\omega(k)\|_{f_0}<\infty$, and
$\Lambda_0(F)=\int_{Z(\auA)/Z(E)}\Lambda(\rho(x)F)\,dx$ (probability Haar measure on the compact group
$Z(\auA)/Z(E)$), then for every Fock vector $v$ of $X\otimes\R$ with $\|v\|_{\max(l,2s)}<\infty$,
\[
|(\Lambda-\Lambda_0)(\omega(k)(v\otimes f_0))|\le C\,\|\Lambda\circ\omega(k)\|_{f_0}\,\delta^l\,\|v\|_l .
\]
(Here $v\in\pzF_s$ by Lemma~\textup{{\cite[Lemma~3.30(b)]{P53}}}, so $v\otimes f_0$ is a datum of class $s$.)
\end{quotedthm}

\begin{quotedthm}{\cite[Corollary~3.34]{P53}}{ (Absolute convergence in the exchange identity)}
Assume {\cite[Proposition~3.26]{P53}}. \textup{(i)} If $\Phi$ on $\pcMat_{r\times N}(F_S)$ is supported
on matrices of rank $r$ (at every place of $S$) then $h\mapsto\theta^{W,V}_{\alpha,\beta}(g,h;[v,\Phi]_S)$
is integrable on $[\auU(V)]$, uniformly for $g$ in compact sets, and the integral has moderate
growth in $g$. \textup{(ii)} The same holds for $h'\mapsto\theta^{V',W}_{\alpha,\beta}(h',g;[v,\Psi]_S)$
when $\Psi$ on $\pcMat_{N\times b}(F_S)$ is supported on matrices of rank $b$. \textup{(iii)} In
particular, for $\Phi,\Psi$ matched for $\chi$ (over $\pcR_\pcm$, or over $\auO_S$), both integrals of
{\cite[(12)]{P53}} converge absolutely.
\end{quotedthm}

\begin{quotedthm}{\cite[Lemma~3.62]{P53}}{ (Levi and radical laws)}
Let $\chi_V(u)=\alpha(u/\bar u)\upsilon(u)^{-a}$ for the source character $\alpha$. Then
\begin{align*}
(\omega_{\alpha,\beta\upsilon^b}(m(q,g),h)f)[x]&=\chi_V(\det q)^{-1}|\det q|^{a/2}\,
\omega_{\alpha,\beta}(g,h)\bigl(f[h^{-1}xq]\bigr),\\
(\omega(n(c,z),1)f)[x]&=\psi_{x^*J_Vx}(z)\,\psrho(x c^*J_{W_0})
\bigl(f[x]\bigr),\qquad\psi_{x^*J_Vx}(z)=\psi\bigl(-\tfrac12\jmath\operatorname{tr}(z\,x^*J_Vx)\bigr),
\end{align*}
with $\psi_\beta$ as in {\cite[(18)]{P53}}.
\end{quotedthm}

\begin{quotedthm}{\cite[Lemma~4.24]{P53}}{ (The vacuum coefficient does not vanish)}
In the setting of {\cite[\S4.2.1]{P53}}, let $u\in\auU(\eta)$ be block diagonal and $\omega(u)$ as
there. Then $c_u=\langle1,\omega(u)1\rangle\ne0$.
\end{quotedthm}

\begin{quotedthm}{\cite[Proposition~6.9]{P53}}{}
\begin{enumerate}
\item[(a)] $\alpha$ is a unit at every prime $w$ of $E$ with $w\nmid q_1$; in particular at every
chosen place, and at every prime of $\pcm(\vartheta)$ if $\auNm Q_\vartheta$ is prime to $q_1$.
\item[(b)] Let $p\mid q_1$ be a prime that ramifies in $\Q(\zeta_{q_1})$ (every odd $p\mid q_1$, and
$p=2$ if $4\mid q_1$). Then $p$ splits in $K$, and every prime $w$ of $E$ above $p$ is split over $E^+$
($w\ne\aubar w$) and ramified over $K$; $\alpha$ is not a unit at the primes $w\mid p$ that divide
$\omega$.
\end{enumerate}
\end{quotedthm}

\begin{quotedthm}{\cite[Lemma~6.30]{P53}}{ (The vacuum at a split place)}
The self-dual lattice $M_v=\alpha\auHom(L_S,L_T)_v$ is $M_w\oplus M_{\aubar w}$ with
$M_w=\alpha_w\pcMat_{N'\times N}(\auO_w)$ and $M_{\aubar w}$ the dual of $M_w$ for the pairing
$(x,a')\mapsto[x,a']$ of $X_w\times X_{\aubar w}$. In the model (W5) the vacuum $\psvac_{M_v}$ is,
up to a scalar, the indicator function $\mathbf 1_{M_w}$.
\end{quotedthm}

\begin{quoteddef}{\cite[Definition~6.31]{P53}}{ (Data at the primes of a modulus)}
For every $w\mid\pcm$ let $a_w$ be the generator of the ideal $\alpha_w\auO_w$ fixed in
{\cite[\S5.1.1]{P49}} ($a_w=\alpha_w$ if $\alpha_w$ is not a unit, $a_w=1$ otherwise). For a function $\Psi$ on
$\pcMat_{N'\times N}(\pcR_\pcm)$ let $\Psi^\natural$ be the function on $\prod_{w\mid\pcm}
\pcMat_{N'\times N}(E_w)$ equal to $\Psi(a^{-1}x\bmod\pcm)$ at $x\in a\,\pcMat_{N'\times N}
(\widehat\auO)$ (that is, $x_w\in a_w\pcMat(\auO_w)$ for all $w\mid\pcm$) and to $0$ elsewhere, and
use at $v\in S(\pcm)$ the $L^2$-norm for the measure giving $M_w=a_w\pcMat(\auO_w)$ the volume $1$.
The datum $[P,\Psi]$, the kernels, the lifts and the transposed lifts are then defined by the
formulas of {\cite[\S2.3]{P53}}.
\end{quoteddef}

\begin{quotedthm}{\cite[Lemma~6.32]{P53}}{ (The normalised coordinates)}
Let $a$ be the idele of {\cite[\S5.1.1]{P49}}, $S,T$ spaces of ranks $N,N'$, $v=w\aubar w$ a place of
$E^+$ split in $E$, and $X_w,X_{\aubar w}$ the components of $\auHom(S,T)_v$, written as matrices $x$
in the frames. Put $\tilde x=a_w^{-1}x$ on $X_w$ and $\tilde x=a_{\aubar w}^{-1}x$ on $X_{\aubar w}$, and
more generally $\tilde x=a^{-1}x$ on $\auHom(S,T)(\auA_f)$.
\begin{enumerate}
\item[(a)] $\alpha_w$ is not a unit if and only if $w$ lies above a prime $p\mid q_1$ and $w\mid\omega$;
then $w\ne\aubar w$, $\aubar w\nmid\omega$ and $\alpha_{\aubar w}=a_{\aubar w}=1$. Every place of $E^+$
above $p\mid q_1$ is $w\aubar w$ with exactly one of $w,\aubar w$ dividing $\omega$. No chosen place
and no prime above $N$ lies above $q_1$. A prime of $Q_\vartheta$ may lie above $q_1$ and divide
$\omega$; such primes are the only primes of a modulus $\pcm_\vartheta$, of $S(\pcm_\vartheta)$, or of
the moduli of the grids at which $a_w\ne1$.
\item[(b)] The self-dual lattice $\alpha\auHom(L_S,L_T)$ is $a\auHom(L_S,L_T)$; at $v$ its components
are $M_w=a_w\pcMat(\auO_w)$ and $M_{\aubar w}=a_{\aubar w}\pcMat(\auO_{\aubar w})$, i.e.\ the integral
matrices in the coordinates $\tilde x$, and they are dual to each other for $[\,,]$. For $x\in X_w$,
$y\in X_{\aubar w}$ one has $[x,y]=\lambda_w(\operatorname{tr}({}^t\tilde y\tilde x))^{\pm1}$ with
$\lambda_w(t)=\psi_w(\jmath(a\aubar a)_wt)$, trivial on $\auO_w$ and not on $\varpi_w^{-1}\auO_w$
($E_{\aubar w}$ identified with $E_w$ by conjugation). In the frame coordinates the dual of
$\pcMat(\auO_w)$ is $\aubar{a_w}\,a_{\aubar w}\pcMat(\auO_{\aubar w})$, which is not
$\pcMat(\auO_{\aubar w})$ when $a_w\ne1$ or $a_{\aubar w}\ne1$.
\item[(c)] On $X_w$ the coordinate $T$ of (N5) and the coordinate $a_w^{-1}x$ of
{\cite[Definition~6.31]{P53}} are both $\tilde x$. Hence for $w\mid\pcm$:
$\Psi^\natural=\Psi(\tilde x\bmod\pcm)\,\mathbf 1[\tilde x\ \text{integral}]$; the datum
$\delta_\iota^\natural$ is $\mathbf 1_{a_w(\iota+\varpi\pcMat(\auO_w))}$ in the frames, which is the
datum of {\cite[\S5.1.3]{P49}} at a level place; and a datum defined as a function of $T$ by (N5) (as in
\S\S\ref{po:ss-glue}, \ref{pte:sec} and~\ref{plf:ss-KS}) is the same function of $\tilde x$, so it is
of the form $\Psi^\natural$ exactly when, as a function of $T$, it is supported on the integral
matrices and depends only on $T$ modulo $\pcm$.
\item[(d)] For $u\in\auHom(S,T)(\auA_f)$, $u\in\alpha\auHom(L_S,L_T)$ if and only if $\tilde u$ is
integral. If $u=s_0y$ with $s_0=a$, then $\tilde u=y$; for $h\in\auU(T)(\auA_f)$,
$\widetilde{h^{-1}u}=h^{-1}y$.
\item[(e)] Let $u$ be a finite idele of $E$ with $u\aubar u=1$, acting as a central element of
$\auU(T)(\auA_f)$. Then $\omega(1,u)$ carries the vacuum of a self-dual lattice $\Lambda$ of
$\auHom(S,T)(\auA_f)$ to a multiple of modulus one of the vacuum of $u\Lambda$, and for every datum $f$,
all characters $\alpha',\beta'$ and $g\in\auU(S)(\auA)$,
\[
\int_{[\auU(T)]}\theta\bigl(\omega_{\alpha',\beta'}(g,h)\,\omega(1,u)f\bigr)\,dh=\beta'(\det u)^{-1}
\int_{[\auU(T)]}\theta\bigl(\omega_{\alpha',\beta'}(g,h)f\bigr)\,dh
\]
whenever either side converges absolutely.
\item[(f)] The map $w\mapsto w^\dagger$ carries $\alpha\auHom(L_S,L_T)$ onto $\aubar\alpha\auHom(L_T,L_S)=
u\,\alpha\auHom(L_T,L_S)$, $u=\aubar a/a$ (a finite idele with $u\aubar u=1$, central in $\auU(S)$,
equal to $1$ at the places not above $q_1$). At $v$, in the coordinates $\tilde x$ of $\auHom(S,T)$ and
$\tilde X$ of $\auHom(T,S^-)$ (frames of $S$ and $T$), it is $\tilde X=-\nu^{-1}\,{}^t\tilde x$ after
identifying $E_{\aubar w}$ with $E_w$ by conjugation, where $\nu=\aubar{a_{\aubar w}}/a_w$ is the
conjugate of $u_{\aubar w}^{-1}$.
\end{enumerate}
\end{quotedthm}

\begin{quotedthm}{\cite[Theorem~9.1]{P53}}{ (Statement (CL))}
Let $\sigma\in\auAut(\C/K)$ and $2\le z\le t$. There is $a_\sigma(z)>0$ such that for every critical
pair $p=(\nu,z)$ there is a critical pair $p^\sigma=(\nu^{\sigma},z)$, whose weight $-c^\sigma$ satisfies
$c^\sigma(\sigma\circ\tau)=c(\tau)$ for all $\tau\in\auPl$, with
\[
\bigl|\sigma\bigl(\|\varphi_{p,1}\|^2\bigr)\bigr|=a_\sigma(z)\,\|\varphi_{p^\sigma,1}\|^2 .
\]
\end{quotedthm}

\begin{quotedthm}{\cite[Theorem~2.2]{P55}}{ (A twist with few congruences at one prime)}
Let $D$ be admissible data, $\ell$ a prime and $\iota\colon\C\to\auQlb$ an isomorphism. There is an
$m$-fold twist $\tau$ such that:
\begin{enumerate}
\item[(a)] $\ell\nmid\auNm Q_\tau$;
\item[(b)] for every ramified critical pair $(\nu,z)$, every $1\le n\le m$ and both signs,
\[
\pgv\Bigl(\frac{|L(z,(\nu\tau^{\pm n})^{-1})|}{|L(z,\nu^{-1})|}\Bigr)\le e'_\ell;
\]
\item[(c)] for $1\le j\le m$, the projection to $\mathcal Z_j(\tau^j)$ has denominator at most
$\beta_\ell$ at $\iota$, for the idempotent $e_0$;
\item[(d)] for $1\le j\le m$, every main pair $(\nu,z)$ for $j$ and all $\Xi,\Xi'\in\mathcal
F(O_{\nu,z}(\tau^j))$ that are integral at every isomorphism $\C\to\auQlb$ and orthogonal to those
false lifts (for $e_1$) that lie in the kernel of the adjoint of the lift from $\auU(1)$
(\emph{test forms}),
\[
\pgv\bigl(\langle\pgpr\,\Xi,\pgpr\,\Xi'\rangle\bigr)\ge-\beta_\ell-\pgv(h_1),
\]
where $\pgpr$ is the orthogonal projection to the line $\C L_{\nu,z}(\tau^j)$.
\end{enumerate}
\end{quotedthm}

\begin{thebibliography}{L-values}
\bibitem[Row D]{P46} D.~A.~Trejo Pizzo, \emph{The sign of the radial intersection form on the idele class group}, preprint, 2026, \url{https://github.com/dtrejopizzo/riemann/tree/main/04-papers/46-row-d-sign-radial-intersection-form}.
\bibitem[Theta]{P49} D.~A.~Trejo Pizzo, \emph{The Short Forms Theorem: theta lifts, complex multiplication and divisibility}, preprint, 2026, \url{https://github.com/dtrejopizzo/riemann/tree/main/04-papers/49-short-forms-theorem-theta-lifts-complex-multiplication}.
\bibitem[Moment]{P51} D.~A.~Trejo Pizzo, \emph{The Spanning Theorem: the second moment of a Poincar\'e map}, preprint, 2026, \url{https://github.com/dtrejopizzo/riemann/tree/main/04-papers/51-spanning-theorem-second-moment-poincare-map}.
\bibitem[Exchange]{P53} D.~A.~Trejo Pizzo, \emph{The Covolume Theorem: exchange identities and the determinant of a chain of theta lifts}, preprint, 2026, \url{https://github.com/dtrejopizzo/riemann/tree/main/04-papers/53-covolume-theorem-exchange-identities-chain-determinant}.
\bibitem[Galois]{P55} D.~A.~Trejo Pizzo, \emph{The Covolume Theorem: regular families of twists, Selmer groups and Ribet's method}, preprint, 2026, \url{https://github.com/dtrejopizzo/riemann/tree/main/04-papers/55-covolume-theorem-regular-families-of-twists}.
\bibitem{MV} H.~L.~Montgomery and R.~C.~Vaughan, \emph{Multiplicative Number Theory I.
Classical Theory}, Cambridge Studies in Advanced Mathematics \textbf{97}, Cambridge
Univ.~Press, 2007.
\bibitem{Banaszczyk} W.~Banaszczyk, \emph{New bounds in some transference theorems in the geometry
of numbers}, Math.\ Ann.\ \textbf{296} (1993), 625--635.
\bibitem{Green1955} J.~A.~Green, \emph{The characters of the finite general linear groups},
Trans.\ Amer.\ Math.\ Soc.\ \textbf{80} (1955), 402--447.
\bibitem{KashiwaraVergne} M.~Kashiwara and M.~Vergne, \emph{On the Segal--Shale--Weil
representations and harmonic polynomials}, Invent.\ Math.\ \textbf{44} (1978), 1--47.
\bibitem{RallisHoweDuality} S.~Rallis, \emph{On the Howe duality conjecture}, Compositio Math.\ \textbf{51} (1984), 333--399.
\bibitem{RallisInjectivity} S.~Rallis, \emph{Injectivity properties of liftings associated to
Weil representations}, Compositio Math.\ \textbf{52} (1984), 139--169.
\bibitem{BLR} S.~Bosch, W.~L\"utkebohmert and M.~Raynaud, \emph{N\'eron Models}, Ergebnisse der
Mathematik und ihrer Grenzgebiete (3) \textbf{21}, Springer, 1990.
\bibitem{LangEF} S.~Lang, \emph{Elliptic Functions}, 2nd ed., Graduate Texts in Mathematics
\textbf{112}, Springer, 1987.
\bibitem{SilvermanAEC} J.~H.~Silverman, \emph{The Arithmetic of Elliptic Curves}, 2nd ed., Graduate
Texts in Mathematics \textbf{106}, Springer, 2009.
\bibitem{SilvermanAT} J.~H.~Silverman, \emph{Advanced Topics in the Arithmetic of Elliptic Curves},
Graduate Texts in Mathematics \textbf{151}, Springer, 1994.
\bibitem{HKS} M.~Harris, S.~S.~Kudla and W.~J.~Sweet, \emph{Theta dichotomy for unitary groups},
J.\ Amer.\ Math.\ Soc.\ \textbf{9} (1996), 941--1004.
\bibitem{KudlaSplitting} S.~S.~Kudla, \emph{Splitting metaplectic covers of dual reductive pairs},
Israel J.\ Math.\ \textbf{87} (1994), 361--401.
\bibitem{KubertLang} D.~S.~Kubert and S.~Lang, \emph{Modular Units}, Grundlehren der mathematischen
Wissenschaften \textbf{244}, Springer, 1981.
\bibitem{HoweJAMS} R.~Howe, \emph{Transcending classical invariant theory}, J.\ Amer.\ Math.\ Soc.\
\textbf{2} (1989), 535--552.
\bibitem{FollandPS} G.~B.~Folland, Harmonic Analysis in Phase Space, Ann. of Math. Stud. 122, Princeton Univ. Press, 1989.
\bibitem{BorelFin} A.~Borel, \emph{Some finiteness properties of adele groups over number
fields}, Publ.\ Math.\ Inst.\ Hautes \'Etudes Sci.\ \textbf{16} (1963), 5--30.
\bibitem{Jacobowitz} R.~Jacobowitz, \emph{Hermitian forms over local fields}, Amer.\ J.\ Math.\
\textbf{84} (1962), 441--465.
\bibitem{Neukirch} J.~Neukirch, \emph{Algebraic Number Theory}, Grundlehren Math.\ Wiss.\
\textbf{322}, Springer, 1999.
\bibitem{WeilSiegel} A.~Weil, \emph{Sur la formule de Siegel dans la th\'eorie des groupes
classiques}, Acta Math.\ \textbf{113} (1965), 1--87.
\bibitem{Hecke1920} E.~Hecke, \emph{Eine neue Art von Zetafunktionen und ihre Beziehungen zur Verteilung der Primzahlen. (Zweite Mitteilung)}, Math.\ Z.\ \textbf{6} (1920), 11--51.
\bibitem{FiteSatoTate} F.~Fit\'e, \emph{Equidistribution, $L$-functions, and Sato--Tate groups}, preprint, arXiv:1405.5162 (2014).
\bibitem{ReedSimonI} M.~Reed and B.~Simon, \emph{Methods of Modern Mathematical Physics I: Functional
Analysis}, revised ed., Academic Press, 1980.















\bibitem{ZelevinskyLNM} A.~V.~Zelevinsky, \emph{Representations of finite classical groups. A Hopf
algebra approach}, Lecture Notes in Math.\ 869, Springer, Berlin, 1981.
\bibitem{PoorYuen} C.~Poor and D.~S.~Yuen, \emph{Linear dependence among Siegel modular forms},
Math.\ Ann.\ \textbf{318} (2000), 205--234.

\end{thebibliography}

\end{document}
